\usepackage{etoolbox} % \AtEndEnvironment, used for the itemize level-drop spacing below
\usepackage{tikz} % For painting a white box over title bar
\usetikzlibrary{calc}
\usetikzlibrary{arrows.meta} % for e.g. arr/.style={-{Stealth[...]}...} on diagram slides
\usetikzlibrary{positioning} % for `right=<dist> of <node>`-style relative node placement
\usetikzlibrary{decorations.text} % text-along-a-path, for a circular photo's curved caption
\usepackage{booktabs} % \toprule/\midrule/\bottomrule for gothamtable
\usepackage{colortbl} % \arrayrulecolor{}, used by gothamtable to color those rules
\usepackage{pifont} % \ding{51}/\ding{55} checkmark/X-mark glyphs, e.g. slides/fr3-Forms-of-Capital-Market-Efficiency-Strong-Semistrong-Weak-current.tex's own validity ratings
\usepackage{pgfplots} % real-data line/bar charts (e.g. slides/3vp-Metallgesellschaft-Oil-Price-Chart-current.tex) -- confirmed safe alongside this deck's own \tableofcontents, see CONVERSION_NOTES.md's "beamer/TeX gotchas" entry
\pgfplotsset{compat=1.18}

% Restores genuine Computer Modern math fonts. Found 2026-07-15, direct
% instruction ("math equations do not have serifs, look at vcv"): despite
% fontspec being loaded with `no-math` (beamerfontthemegotham.sty), math
% still wasn't landing on Computer Modern -- checked directly against a
% compiled PDF's own per-glyph font names (pymupdf), not just eyeballed:
% digits/math-italic letters were coming out as FiraSans-Light(Italic)
% (this deck's own body-text sans font) and operators/parens as CMSS10
% (Computer Modern's *sans-serif* shape, not its serif one). Root cause:
% beamer's own base classes redefine the math alphabets to follow the
% presentation's sans-serif text font by default, specifically so math
% matches an all-sans deck -- `no-math` only stops *fontspec* from doing
% this, it has no say over beamer's own separate override. Beamer ships
% the standard fix for exactly this case: a font theme named "serif"
% with an "onlymath" option that restores ordinary (Latin
% Modern/Computer-Modern-equivalent) math fonts without touching any of
% the surrounding sans-serif text. Confirmed directly afterward: every
% glyph in a re-checked formula now reports as LMRoman*/CMR10/CMMI8 (all
% genuinely serif), not FiraSans/CMSS10.
\usefonttheme[onlymath]{serif}

% \text{} (amsmath, used throughout for plain-word operands inside math,
% e.g. \text{E}[\tilde{R}_M]) is a separate mechanism from the math
% alphabets \usefonttheme[onlymath]{serif} above just fixed: unlike
% \mathrm/\mathit/etc., \text{} is defined to typeset in whatever the
% *surrounding running text* font is at that point in the document, not a
% math alphabet -- so on this deck's all-Fira body text, \text{} inside a
% formula still comes out in Fira Sans (and picks up italic if the
% enclosing text happens to be italicized), even though the digits/
% operators around it are correctly Computer Modern. Found 2026-07-20 on
% slides/2gd-Analyzing-Reward-to-Risk-Ratio-CAPM-current.tex's own
% \text{E}/\text{Cov} operands. Fixed generally, project-wide, by wrapping
% \text{}'s argument in an explicit upright-Computer-Modern font switch --
% \usefont{OT1}{cmr}{m}{n} rather than \mathrm{} (which would pull in
% \text{}'s own math-alphabet substitution logic unnecessarily) or
% \textrm{} (which switches *family* only, so still comes out italic
% inside an italicized context, which \text{} should never do here) --
% deliberately hardcoding "not italic" (shape n) rather than inheriting
% the ambient shape, since \text{}'s whole purpose in this deck's formulas
% is plain upright labels (E, Cov, "the" CAPM, ...), regardless of
% whatever font context happens to surround the formula itself.
\let\gothamOrigText\text
\renewcommand{\text}[1]{\gothamOrigText{\usefont{OT1}{cmr}{m}{n}#1}}

% Color definition, institute logo path, and the two logo-calibration
% lengths (\gothamLogoHeight, \gothamFrametitleRightpading) all live in the
% per-university preamble that \input{}s this file, not here -- this file
% is shared machinery only, no institution-specific branding. Every color
% token below (\c_primary_*, \c_secondary_*, \c_tertiary_*, \c_custom_*) and
% \gothamInstituteLogoPath/\gothamLogoHeight/\gothamFrametitleRightpading
% are assumed already defined by the time this file is reached -- see e.g.
% Slidemaster_preamble_uni-graz.tex for the Graz values, or
% Slidemaster_preamble_mci.tex/Slidemaster_preamble_uni-ibk.tex for the
% other two institutions this deck supports. \gothamLogoWidth is the one
% logo-derived length that *does* stay computed here (via \settowidth,
% further below) rather than hardcoded per institution, since it already
% adapts automatically to whatever image \gothamInstituteLogoPath points at.

% Style hyperlinks (\href{url}{text}) consistently across all slides:
% underlined, so a link is visually recognizable as such against plain
% body text -- yellow (this project's usual accent color) was tried first
% for the *text* itself, but yellow-on-white body text is hard to read,
% unlike yellow used as an accent against dark/black text elsewhere in the
% theme. Went through a few more candidates on 2026-07-06 (all direct
% instruction) before settling on c_secondary_blue for both text
% and underline (a real color reads more clearly as "this is a link" than
% any gray does, without yellow's legibility problem).
%
% 2026-07-07, direct instruction: split text color and underline color
% instead of using one color for both -- link text back to
% c_custom_gray (the near-black dark gray also used for
% quoteframe's attribution text), underline in a separate accent color.
% Plain \underline can't do this: its rule is drawn in whatever color is
% current at that point, so a single \underline{...} call always
% underlines in its own text's color, never a second, different one.
% \gothamHrefUnderline below reimplements it instead: measure the linked
% text's own width into a savebox, print that box in customgray, then
% back up by the box's own width and draw a plain \rule of the same
% width, raised to sit just below the text's baseline, in the accent
% color. Same single-line limitation as plain \underline (and the reason
% \url{} is used instead of \href on slides/acw-Some-Topics-in-Testing-the-EMH-current.tex --
% see that slide's own comment): a \rule, like \underline, can't span a
% line break either, so this only ever wraps normal, comparatively short
% \href display text, never something that needs mid-word breaking.
%
% First tried c_primary_highlight for the underline (same day) -- confirmed
% via the compiled PDF that the rule *was* being drawn (right color,
% position, width), so not a compile-time bug, but yellow (very high
% luminance, close to white) at this rule's 0.4pt thickness reads as
% essentially invisible against the white slide background, worst at
% footnote size (see e.g. slides/za2-Startup-Financing-Categorization-current.tex's
% own long inline URL) where it's easy to miss entirely. Switched
% 2026-07-07 (direct instruction) to c_secondary_blue instead --
% far more contrast against white at the same thickness -- and the rule
% was also thickened at the same time (0.4pt -> 0.6pt), on the reasoning
% that it might need it for the same contrast reason. It didn't: blue
% already reads clearly at the original thickness, and 0.6pt instead
% just looked heavier than intended (direct feedback 2026-07-08, e.g. on
% slides/6qi-Yield-Spreads-Corporate-Bonds-current.tex). Reverted to
% 0.4pt -- only the color was ever actually asked to change.
\hypersetup{colorlinks=true, urlcolor=c_custom_gray, linkcolor=c_custom_gray}
\let\gothamOrigHref\href
\newsavebox{\gothamHrefTextBox}
%
% 2026-07-08 (direct instruction): the underline's own rule thickness now
% scales with the ambient font size instead of a single fixed 0.2pt for
% every link everywhere -- 0.2pt at this deck's footnote size, 0.4pt at its
% normal body-text size, interpolated (and extrapolated, for any other
% size a link happens to appear at) linearly in between.
%
% Deriving "current font size" from \f@size (the textbook NFSS way) was
% tried first and confirmed broken specifically *inside* a beamer frame --
% \typeout{\meaning\f@size} right after this preamble loads shows a normal,
% defined macro (`macro:->10.95`), but the exact same check inside
% `\begin{standardframe}...\end{standardframe}` reports `\f@size` as
% *undefined*, presumably from beamer's own overlay-aware font-switching
% machinery not keeping classic NFSS's own current-size bookkeeping in
% sync the way plain (non-beamer) LaTeX does. Since every real use of this
% macro happens inside a frame, that path was a dead end regardless of how
% well it worked in isolation.
%
% Measuring a reference glyph's own rendered height instead sidesteps this
% entirely: building a box and reading back its own height just measures
% whatever font is *actually* active at the call site, independent of
% whether beamer's own internal "current size" bookkeeping happens to
% agree -- and within one font
% family, every glyph's metrics scale by exactly the same ratio the point
% size does, so a measured height is just as valid a stand-in for "current
% size" as \f@size would have been, had it worked. "M" (present in every
% context this is used in) is the reference glyph; \gothamHrefUnderlineFnHeight/
% \gothamHrefUnderlineBodyHeight are its height at this deck's footnote font
% and normal body font respectively, each measured once, here, in the
% relevant font context.
% Deliberately built from raw TeX primitives (\newdimen/\setbox/\ht),
% *not* \newlength/\setlength/\settoheight -- found, while building this,
% that loading this deck's theme silently turns \setlength and
% \settoheight into something else entirely: `\meaning\setlength` reads
% `macro:->\protect \setlength` right after `\usetheme{gotham}` (still the
% ordinary kernel definition beforehand), and the same for `\settoheight`.
% Not chased to which file does this or why (bisected only as far as
% "plain beamer alone doesn't, beamer+gotham does" -- most likely
% `translator`, which beamer loads for localized strings and is known to
% robustify commands this way, but not confirmed) -- what matters
% practically is that a *global* assignment through either one silently
% no-ops instead of erroring, confirmed directly and in isolation (a
% minimal beamer+gotham file, no other content): `\global\setlength{\len}
% {\otherlen}` leaves `\len` at 0pt after the group closes, while the
% exact same pattern with plain `\newdimen`/raw `=` assignment propagates
% correctly. Every plain (non-global) `\setlength` elsewhere in this
% project's own files is unaffected -- it's specifically the *global*
% form that breaks, which is why this is the first place it's shown up
% despite `\setlength` being used constantly throughout this preamble.
\newdimen\gothamHrefUnderlineFnHeight
\newdimen\gothamHrefUnderlineBodyHeight
\newdimen\gothamHrefUnderlineNowHeight
\newdimen\gothamHrefUnderlineThickness
% A dedicated box register, not the commonly-borrowed box0 -- several
% slides in this deck (the "Long-term bonds" tables) already use box0
% themselves (`\sbox0{...}`/`\usebox0`) for their own, unrelated
% mid-table-layout purposes, and a link can appear anywhere, including
% inside one of those tables -- reusing box0 here would silently clobber
% whatever they were holding there.
\newbox\gothamHrefUnderlineBox
% The two reference heights are (re)measured on every single \href call,
% not once at preamble-load time -- tried once-at-load first and found a
% second, independent problem with it: `\usebeamerfont{footnote}` right
% here (before `\begin{document}`) silently has no effect at all (footnote-
% size "M" and normal-size "M" measured identically, confirmed directly),
% presumably because beamer's font-theme machinery isn't fully installed
% until the document starts. Since every real `\href` call happens well
% inside the document (inside some frame), measuring fresh each time
% sidesteps that timing question entirely rather than needing to find the
% right hook to measure once, early enough to be free but late enough to
% be correct -- the extra cost is two throwaway one-glyph boxes per link,
% not something worth optimizing away in a slide deck.
\newcommand{\gothamHrefUnderlineMeasureRefHeights}{%
	{\usebeamerfont{footnote}\setbox\gothamHrefUnderlineBox=\hbox{M}\global\gothamHrefUnderlineFnHeight=\ht\gothamHrefUnderlineBox}%
	{\normalsize\setbox\gothamHrefUnderlineBox=\hbox{M}\global\gothamHrefUnderlineBodyHeight=\ht\gothamHrefUnderlineBox}%
}
% A plain e-TeX \dimexpr can multiply a dimension by another dimension and
% then divide by a third, as one combined operation, as long as the
% *middle* operand of the chain is itself wrapped in its own nested
% \dimexpr...\relax -- confirmed directly (an expl3 \dim_set:Nn attempt
% first, using \seq_.../\fp_... the way \setfootnotesymbols above does,
% kept hitting "Missing ) inserted for expression"/type-mismatch errors:
% dividing one dimension by another is a *ratio*, not a dimension, and
% \dim_set:Nn's own expression grammar doesn't implicitly convert between
% the two the way this nested-\dimexpr form does). No expl3 needed here at
% all -- plain e-TeX arithmetic, already available in any current LaTeX.
\newcommand{\gothamHrefUnderlineSetThickness}{%
	\gothamHrefUnderlineThickness=\dimexpr 0.2pt +
		(\gothamHrefUnderlineNowHeight - \gothamHrefUnderlineFnHeight)
		* \dimexpr0.2pt\relax
		/ (\gothamHrefUnderlineBodyHeight - \gothamHrefUnderlineFnHeight)\relax
}
\newcommand{\gothamHrefUnderline}[1]{%
	\sbox{\gothamHrefTextBox}{#1}%
	\textcolor{c_custom_gray}{\usebox{\gothamHrefTextBox}}%
	\hspace{-\wd\gothamHrefTextBox}%
	\gothamHrefUnderlineMeasureRefHeights
	\setbox\gothamHrefUnderlineBox=\hbox{M}\gothamHrefUnderlineNowHeight=\ht\gothamHrefUnderlineBox
	\gothamHrefUnderlineSetThickness
	\raisebox{-0.35ex}[0pt][0pt]{\textcolor{c_secondary_blue}{\rule{\wd\gothamHrefTextBox}{\gothamHrefUnderlineThickness}}}%
}
\renewcommand{\href}[2]{\gothamOrigHref{#1}{\gothamHrefUnderline{#2}}}
% \url{} (the `url` package hyperref depends on/patches) defaults to a
% monospace-ish "tt" style regardless of the above -- used deliberately in
% only one place in this deck,
% slides/acw-Some-Topics-in-Testing-the-EMH-current.tex, specifically
% *for* its line-breaking-at-slashes behavior (a plain \href with the
% URL as its own display text has no break points at all -- see that
% slide's own comment). Every other same-text-as-target URL in the deck
% (including slides/bb6-A-Perpetual-Bond-current.tex's own video link,
% which used to be a \url{} too) uses plain \href{url}{url} instead, since
% \url{} doesn't pick up this deck's custom colored underline below (its
% breakable content can't be measured into the savebox the underline
% needs -- see that mechanism's own comment) and most of these URLs are
% short enough to not need the breaking \url{} exists for anyway. 2026-07-07,
% direct instruction: that one genuinely-needed \url{} (on s-acw) should
% still look like this deck's other hyperlinks, i.e. the surrounding body
% font, not \url{}'s own default monospace.
% `\urlstyle{same}` (a standard url-package switch) does exactly that,
% with no effect on \url{}'s own breaking behavior. Set deck-wide, not
% only on that one slide, so any future \url{} use matches by default
% too, rather than needing to remember this each time. Coloring already
% matches automatically (hyperref's own \url respects
% \hypersetup{urlcolor=...} the same way \href does); the underline
% \href gets above doesn't carry over, since \underline's own argument
% can't contain line breaks -- wrapping \url{}'s necessarily-breakable
% text in \underline would silently undo the breaking this exists for,
% so that one difference is accepted rather than "fixed."
\urlstyle{same}

% gotham's own footnote color (beamercolorthemegotham.sty:
% \setbeamercolor{footnote}{fg=normal~text.fg!90}) is a 90%-black gray --
% against plain black body text, customgray links contrast clearly (found
% while investigating a 2026-07-05 report that a footnote's own \href
% didn't read as a link the way body-text ones like the "Moodle" link on
% slides/hxr-Administrative-Details-current.tex do); against gotham's own
% dimmed footnote gray, the same customgray is only 25/255 lighter and
% barely distinguishable at footnote size, with only the underline giving
% it away. Overridden to plain black here so a footnote link contrasts the
% same way a body-text one does -- a one-line override of the theme's own
% color (not editing beamercolorthemegotham.sty itself, which stays
% untouched per its own upstream-file status), and a plausible general
% legibility improvement for footnote text regardless of links.
\setbeamercolor{footnote}{fg=c_primary_dark}

% Footer text color (\FooterText/slide id/frame number) -- \gothamFooterTextColor,
% set by the per-university preamble before \input{}ing this file (the
% footer box's own background color, \gothamFooterBoxColorDefault above, is
% a separate, independently configurable setting). No explicit
% \setbeamercolor{footline}{...} existed before this -- the footer text
% simply inherited whatever "normal text" foreground already was (black),
% which happened to read fine against every frame's own yellow/orange
% footer box so far, but wouldn't automatically read against an arbitrary
% future footer color (e.g. a dark blue) without this override.
\setbeamercolor{footline}{fg=\gothamFooterTextColor}

\colorlet{colorDark}{c_primary_dark} % FG in light/normal mode
% Title bar background -- \gothamTitleBoxColor, set by the per-university
% preamble before \input{}ing this file (same idiom as
% \gothamBulletColor/\gothamFooterBoxColorDefault above).
\colorlet{colorA}{\gothamTitleBoxColor} % frametitle, standin.out

\gothamset{background=light} % Needs to go after color redefinitions, so they are applied
\setbeamercolor{background canvas}{bg=white}

% gotham's own default behavior on \section{} (used for a deck's chapter
% grouping/table of contents) is to auto-insert its *own* plain, unstyled
% section-divider frame plus a small "current section highlighted in the
% outline" frame, every time -- redundant with (and visually inconsistent
% next to) this project's own hand-styled sectionframe (see README.md's
% "Special frame types" section), which every chapter in this deck already
% opens with as an explicit \slide{}. Disabled so \section{} only feeds
% \tableofcontents (see slides/nkj-Agenda-current.tex) without also
% inserting its own extra frames.
\gothamset{sectionframe default=false}

% Same reasoning, for \subsection{} (2026-07-15, direct instruction):
% gotham's default subsection-divider frame is disabled too, so
% \subsection{} calls still register in the PDF's own bookmark/outline
% structure (Acrobat's "chapters" pane -- a plain PDF sectioning feature,
% independent of whether beamer auto-inserts a frame for it) and still feed
% \tableofcontents (see slides/nkj-Agenda-current.tex's own
% `\tableofcontents[hideallsubsections]`, which then hides them from that
% one slide specifically), without also spawning their own unstyled divider
% slide in the middle of a chapter every time one is added.
\gothamset{subsectionframe default=false}

% Title text color -- \gothamTitleTextColor, set by the per-university
% preamble before \input{}ing this file (the title bar's own background
% color, \gothamTitleBoxColor/colorA above, is a separate, independently
% configurable setting).
\setbeamercolor{frametitle}{fg=\gothamTitleTextColor}

% A color leak from \titleframe/\quoteframe into whichever frame follows
% them is fixed in standardframe's own opening code below (search
% "Forces this frame's own body") rather than here -- an
% \AtBeginEnvironment{frame}{\color{c_primary_dark}} hook was tried first
% and confirmed, by isolated reproduction, to have *no effect at all*:
% beamer's own "frame" isn't a plain \newenvironment etoolbox's
% \AtBeginEnvironment can patch that way, so the hook silently never fired.

% Set content text margins
\setbeamersize{text margin left=.45cm, text margin right=.45cm}

%\colorlet{colorPale}{gPaleYell} % BG in light/normal mode
%\colorlet{colorDark}{gDarkBlack} % FG in light/normal mode
%\colorlet{colorA}{gAnthracite} % frametitle, standin.out,
%\colorlet{colorAreversed}{gLightTeal} % frametitle, standin.in,
%\colorlet{colorB}{gMidGrey} % gray BG : progress bar, blocks
%\colorlet{colorC}{gDeepYellOr} % progress bar
%\colorlet{colorD}{gLightOrange} % alert
%\colorlet{colorE}{gLightGreen} % example

% Opt-in yellow border for an image placed via a tikz node (e.g.
% \node[gothamImageBorder] {\includegraphics{...}};) -- matches the
% 3pt/accent1-yellow outline PowerPoint draws directly on a picture shape's
% own border (<a:ln w="38100"><a:solidFill><a:schemeClr val="accent1"/>,
% where accent1 turned out to already be this project's own
% c_primary_highlight, not a new color to map). A tikz *style* (a named
% bundle of node keys), not a wrapper command: adding or removing it is
% just adding or deleting one word from that one \node's own option list,
% no separate macro call/argument needed -- add `gothamImageBorder` to draw
% the border, leave it off for a plain image, same as any other tikz key.
%
% Thickness is a single named length, \gothamImageBorderWidth, not a bare
% "3pt" baked directly into the style (2026-07-15, direct instruction: "make
% sure the yellow image borders have consistent thickness throughout the
% project by defining a single variable") -- confirmed directly, by
% measuring the actual PDF content-stream "w" operator, that every
% gothamImageBorder use already rendered at the exact same 3pt regardless
% (a genuinely shared style was already in effect, no numeric drift found),
% so this is a refactor for future adjustability/discoverability, not a fix
% for an actual inconsistency -- any perceived thickness difference between
% slides traces to something else instead (e.g. a border straddling flush
% image content on one slide and not another, which reads as "heavier" even
% at an identical stroke width -- see the gothamImageBorder padding gotcha
% below).
\newlength{\gothamImageBorderWidth}
\setlength{\gothamImageBorderWidth}{2.5pt}
\tikzset{gothamImageBorder/.style={draw=c_primary_highlight, line width=\gothamImageBorderWidth}}

% The "padding gotcha" promised above, finally root-caused (2026-07-23,
% background-task follow-up to the IA2 fine-tuning batch #2 anomaly on
% slides/m2k-Exercise-CAPM-2-current.tex/
% slides/n8e-Tail-Risk-Variance-Clustering-current.tex/
% slides/fm6-Tail-Risk-Market-vs-Simulated-current.tex/
% slides/xqq-Fat-Tails-current.tex, all fixed there with a pragmatic
% `line width=1.2pt` override and an explicitly *unresolved* root cause at
% the time -- "declares an identical stroke width and CTM in the compiled
% PDF's own raw content stream, yet renders at a genuinely different pixel
% thickness", with CTM scale/clips/dash/cap/join all directly checked and
% ruled out). Confirmed via a controlled reproducer (four `gothamImageBorder`
% nodes -- plain opaque image, nested pgfplots axis, plain image at
% `inner sep=3pt`, nested plain tikzpicture at `inner sep=3pt` -- rendered
% at 600dpi and measured by a half-max cross-section profile, not just
% eyeballed): the one variable that actually predicts the rendered width is
% `inner sep`, not content type at all. At `inner sep=0pt` a *fully opaque*
% image measured 1.2pt; the exact same image at `inner sep=3pt` measured
% 2.52pt (matching the declared 2.5pt); a nested pgfplots axis and a plain
% nested tikzpicture, both at `inner sep=3pt`, also measured 2.52pt. Two
% more points confirmed it's a smooth transition, not a binary switch:
% `inner sep=0.6pt` measured 1.8pt (= the 1.2pt floor plus that 0.6pt gap,
% exactly); `inner sep=1.25pt` (= \gothamImageBorderWidth/2, the stroke's
% own half-width) already measured the full 2.52pt.
%
% The mechanism: tikz draws a node's border stroke centered on the node's
% own boundary path, then paints the node's *content* on top of it,
% afterwards. At `inner sep=0`, the content's own bounding box coincides
% exactly with that boundary path, so if the content is opaque all the way
% to its own edge, it paints directly over the stroke's inner half,
% visually leaving only the outer half (`\gothamImageBorderWidth`/2,
% ~1.2pt) -- not a rendering bug, just tikz's normal draw order interacting
% with zero padding. A gap of at least `\gothamImageBorderWidth`/2 (or
% content that's transparent right at its own edge -- see below) leaves the
% border with nothing opaque to cover it, so it renders at its true
% declared width instead.
%
% This also explains why the fix landed on exactly the 4 slides it did,
% each for a different reason: slides/m2k-Exercise-CAPM-2-current.tex
% used `inner sep=3pt` for unrelated layout reasons (room for the axis's
% own tick labels), well past the 1.25pt threshold, so nothing ever covered
% its border. slides/n8e-Tail-Risk-Variance-Clustering-current.tex/
% slides/fm6-Tail-Risk-Market-vs-Simulated-current.tex/
% slides/xqq-Fat-Tails-current.tex all use `inner sep=0pt` -- matching the
% deck's own overwhelming convention -- but their own chart PNGs
% (`chart-index-daily-returns-*.png`/`chart-index-fat-tails-histogram*.png`,
% matplotlib exports with a transparent background) measured alpha=0 at
% every edge pixel (confirmed directly, Pillow), so there was no opaque
% content to cover their own border either. Every other `gothamImageBorder`
% use in this project that was never flagged uses an opaque image (a JPG,
% or a PNG with no transparent margin) at `inner sep=0pt` -- which
% *accidentally* renders at the ~1.2pt covered-half width instead of the
% full 2.5pt this constant nominally declares, and that accidental width is
% what reads as "the deck's normal border thickness" throughout, since it's
% by far the most common case.
%
% Practical upshot for a *new* slide: a `gothamImageBorder` node with
% `inner sep=0pt` wrapping opaque content (the common case) needs no
% special handling -- it will match the rest of the deck automatically. A
% node that instead uses a nonzero `inner sep`, or wraps content with a
% transparent edge (a matplotlib export, an SVG-derived PDF with padding,
% ...), will render at the *full* \gothamImageBorderWidth and look
% noticeably heavier next to the deck's other borders -- either add an
% explicit `line width=\gothamImageBorderWidth/2` (matching the 4 fixed
% slides' own literal 1.2pt) to match the accidental deck-wide look, or
% reduce `inner sep` below \gothamImageBorderWidth/2 so the same coverage
% effect kicks in naturally. The 4 existing per-slide overrides above are
% left as-is -- they're already the correct fix for what's now a fully
% understood cause, not a mystery workaround needing revisiting.

% Reusable "mugshot" -- a circular-cropped photo, bordered like any other
% gothamImageBorder image, with a caption curving along its own bottom arc
% (2026-07-06, direct instruction, generalized from
% slides/4uq-Definition-of-Capital-Market-Efficiency-current.tex's own
% Eugene Fama photo -- see that slide's earlier comment for the source
% PPTX effect this reproduces, `<a:prstGeom prst="ellipse">` +
% `<a:prstTxWarp prst="textArchDown">`, and CONVERSION_NOTES.md's own
% "beamer/TeX gotchas" entry, search "circular photo crop", for the
% clip/border/decorations.text mechanics themselves -- unchanged here,
% just factored out so a second photo doesn't need to repeat them).
%
%   \coordinate (myMugshotCenter) at (...);
%   \gothamMugshot{myMugshotCenter}{media/photo-example.jpeg}{2.6cm}{Some Person, NP 2013}
%
% #1 the *name* of a coordinate the caller has already placed (not a raw
% point expression -- needed as a plain name, not just a point, since the
% caption's own arc path re-references it by name via `++(...)`, which
% only works with a named coordinate/node, not an arbitrary `($...$)`
% calc expression); #2 the image path; #3 the desired final diameter;
% #4 the caption text.
%
% Scales the source photo so it fully covers the circle with no gap
% anywhere -- both dimensions end up at least as big as the diameter --
% while using as much of the source photo as possible (2026-07-07,
% direct instruction: "should not include whitespace... scale the photo
% such that its edges just touch the circle's inner border"). This
% briefly (same day) scaled the *other* way -- by the photo's own
% diagonal, the largest scale at which the *entire* rectangle fits
% inside the circle with nothing clipped at all -- which does avoid
% ever clipping real content, but leaves the circle's own four corners
% uncovered (visible as plain background there, since a non-circular
% photo can only touch an inscribing circle at its corners, not along
% the whole rim) -- exactly the "whitespace" this round's feedback
% ruled out in favor of full coverage. Covering scale is
% `diameter / min(w,h)`: whichever of the photo's own two dimensions is
% *smaller* becomes exactly the diameter (that edge pair touches the
% circle's own inner border, per the direct instruction), and the other,
% larger dimension overhangs the diameter and gets clipped by the
% circular `\clip` -- for a portrait photo (every one used here so far),
% that means width becomes the diameter and a sliver clips off the top
% and bottom, which is the origin of the clipping problem the diagonal
% approach above existed to solve in the first place. Reintroduced
% deliberately: showing as much of the photo as the "no whitespace"
% constraint allows is now the stated priority, over never clipping
% anything.
%
% Measuring the photo's own native size needs `\savebox` (into
% \gothamMugshotBox, declared once below, not a fresh box register per
% call) -- `\includegraphics` with no size option at all renders at the
% image's native resolution, and `\wd`/`\ht` then read that box's actual
% dimensions. pgfmath (already loaded for tikz) converts any TeX
% dimension fed into it to a plain pt-count number automatically, so
% arithmetic on the two measured lengths just works. `min(1,ratio)`
% (ratio = height/width) picks out whichever of the two is smaller
% *relative to the other*, without needing a `sqrt`/`^` at all this time
% -- so this version doesn't hit the "! Dimension too large" pgfmath
% overflow the diagonal version's own arithmetic did (see this file's
% own comment history / CONVERSION_NOTES.md's "beamer/TeX gotchas" for
% that point (easy to misread as a harmless warning, since compilation
% continues and produces *a* number regardless). Reformulating via the
% aspect ratio sidesteps the whole problem: `scale = diameter /
% sqrt(1+ratio^2)` is the same formula algebraically (factor `w` out of
% `sqrt(w^2+h^2)` and it becomes `w*sqrt(1+(h/w)^2)`), but every
% the full story of that overflow, from when this macro's own arithmetic
% did use a `sqrt`).
%
% Caption font reverted to this deck's normal body font (2026-07-07,
% direct instruction: "change the mugshot template back to using regular
% font instead of monospace") -- `\ttfamily` (Fira Mono) was a direct
% instruction on 2026-07-06, one day earlier; \normalfont resets
% explicitly rather than just omitting \ttfamily, so this doesn't
% silently inherit whatever font happens to be active at the call site.
\newsavebox{\gothamMugshotBox}
\newcommand{\gothamMugshot}[4]{%
	\savebox{\gothamMugshotBox}{\includegraphics{#2}}%
	\pgfmathsetmacro{\gothamMugshotImgW}{\wd\gothamMugshotBox}%
	\pgfmathsetmacro{\gothamMugshotImgH}{\ht\gothamMugshotBox}%
	\pgfmathsetmacro{\gothamMugshotRatio}{\gothamMugshotImgH/\gothamMugshotImgW}%
	\pgfmathsetmacro{\gothamMugshotScaledW}{(#3)/min(1,\gothamMugshotRatio)}%
	\begin{scope}
		\clip (#1) circle (0.5*#3);
		\node[anchor=center, inner sep=0pt] at (#1)
			{\includegraphics[width=\gothamMugshotScaledW pt]{#2}};
	\end{scope}
	\draw[gothamImageBorder] (#1) circle (0.5*#3);
	{\normalfont\tiny
	\path[decoration={text along path, text={#4}, text align=center}, decorate]
		(#1) ++(200:0.5*#3+0.25cm) arc (200:340:0.5*#3+0.25cm);
	}
}

% Reusable "operating/cash conversion cycle" bar diagram (FM2_2026s_ps_deck.pptx
% slides 22/24/25/26 all embed the exact same picture -- an EMF, not a native
% shape/table -- so rather than convert that EMF and re-embed it 4 times over,
% it's redrawn once natively here and reused, matching this project's own
% "clean/geometric diagrams get redrawn, not extracted" policy). Three
% stacked bars (Inventory period+Accounts receivable period on top, Accounts
% payable period+Cash conversion cycle in the middle, Operating cycle
% spanning the full width on the bottom) plus 4 tick labels below at
% x=0/(1/3)/(2/3)/1. Geometry/colors measured directly from the source
% deck's own rendered PNG via pixel analysis (PIL run-length color
% detection), not eyeballed: the two internal split points are exactly 1/3
% and 2/3 of the total width (measured 0.3339/0.6672 against a plain
% PowerPoint illustrative diagram, not slide-23's own specific-day example --
% close enough to exact thirds, and a plausible round number for the
% presenter to have actually used, that thirds is treated as the real value
% rather than curve-fitting noise). Fill colors matched to this project's
% *existing* palette rather than the source's own literal RGB, per this
% deck's established emphasis-color policy: source yellow was already an
% exact match for `c_primary_highlight`; source green (RGB 146,208,80) and
% purple (RGB 128,100,162) were close enough to `c_secondary_green`
% and `c_secondary_violet` respectively to reuse rather than define
% two new one-off colors.
%
% Drawn at a fixed "natural" size (10cm wide) with its own top-left corner
% at the calling tikzpicture's origin -- wrap a call in
% \begin{scope}[shift={(x,y)}, scale=s]\gothamCashCycleBars\end{scope} to
% place/resize it, the same convention as this deck's other reusable
% diagram pieces (e.g. slides/bdh-Working-Capital-Example-current.tex's
% circular flow diagram, resized the same way after the fact per user
% feedback) -- resizing later without needing to re-derive any coordinates.
\newcommand{\gothamCashCycleBars}{%
	\fill[c_primary_highlight] (0,0) rectangle (6.667,-0.42);
	\fill[c_secondary_green] (6.667,0) rectangle (10,-0.42);
	\node[align=center, font=\tiny] at (3.333,-0.21) {Inventory period};
	\node[align=center, font=\tiny] at (8.333,-0.21) {Accounts receivable period};
	\fill[c_secondary_green] (0,-0.47) rectangle (3.333,-0.89);
	\fill[c_primary_highlight] (3.333,-0.47) rectangle (10,-0.89);
	\node[align=center, font=\tiny] at (1.667,-0.68) {Accounts payable period};
	\node[align=center, font=\tiny] at (6.667,-0.68) {Cash conversion cycle};
	\fill[c_secondary_violet] (0,-0.94) rectangle (10,-1.36);
	\node[align=center, font=\tiny, text=c_primary_white] at (5,-1.15) {Operating cycle};
	\foreach \gccx in {0,3.333,6.667,10} {
		\draw[c_primary_dark] (\gccx,-1.36) -- (\gccx,-1.46);
	}
	% text width on all 4, narrow enough to force an actual 2-line wrap
	% (not just cap an already-short line): without it, "Raw material
	% purchase" and "Payment for raw materials" (and, symmetrically, the
	% last two) run into each other -- confirmed directly by rendering.
	% A *wide* text width (e.g. 3.1cm, comfortably more than either label's
	% own one-line width) does NOT fix this: \node's text width only forces
	% a wrap once the text exceeds it, so a generous width just lets both
	% neighbors keep their full one-line extent and still collide -- this
	% was tried first and confirmed not to work by rendering, not assumed.
	% 1.5cm is deliberately narrower than any of the 4 would take on one
	% line, guaranteeing the wrap for all of them.
	%
	% All 4 use anchor=north (centered on their own tick, like the middle
	% two always did) -- the outer two originally used anchor=north
	% west/east instead (flush with the tick rather than centered on it),
	% which looked inconsistent with the middle two once seen side by side
	% at actual size. Centering the outer two on their own tick (rather
	% than flush against the diagram's own left/right edge) does push their
	% wrapped text a bit further left/right past the diagram's own
	% bounding box -- checked this stays within the page margin at every
	% scale/position this macro is actually called at in this deck (the
	% narrowest margin case leaves ~0.7cm of clearance), not just assumed
	% safe from the numbers alone.
	\node[anchor=north, font=\tiny, align=center, text width=1.5cm] at (0,-1.5) {Raw material purchase};
	\node[anchor=north, font=\tiny, align=center, text width=1.5cm] at (3.333,-1.5) {Payment for raw materials};
	\node[anchor=north, font=\tiny, align=center, text width=1.5cm] at (6.667,-1.5) {Sale of finished goods};
	\node[anchor=north, font=\tiny, align=center, text width=1.5cm] at (10,-1.5) {Cash collected on sales};
}

% Reusable circular flow diagram (Cash -> Raw materials -> Finished goods ->
% Receivables -> Cash), first built for
% slides/bdh-Working-Capital-Example-current.tex and reused (smaller) on
% slides/qq5-Cash-Conversion-Cycle-current.tex -- factored out here once a
% second slide needed the same picture, same reasoning as
% \gothamCashCycleBars above. Args: #1 diagram center coordinate, #2 radius
% (box-center distance from that center, a bare cm number, no unit suffix),
% #3 the box `scale` factor (a plain node `scale` transform, not a bigger
% font -- see bdh's own comment, preserved below, for why: this diagram was
% resized once already per user feedback, and `scale` was what correctly
% grew the padding/rounded-corners/arrow proportions together rather than
% just the text). Arrow length/width scale with #3 too, computed via
% \pgfmathsetmacro since tikz option values can't do arithmetic inline.
%
% [Preserved from bdh's own resize, 2026-07-06: box size grew by a uniform
% ~1.44x there (compared via the isolated "Cash" box, the one corner with no
% arrow touching it to contaminate the measurement: 0.94cm to 1.36cm wide,
% 0.50cm to 0.74cm tall), and the box-to-box radius grew by the same
% ~1.44-1.45x (1.15cm to ~1.65cm) -- both measured from the user's own
% "Clipboard 1.jpg" reference via pixel analysis, not eyeballed, and
% consistent with a single uniform resize of the whole group.]
% #1 is a *complete* coordinate spec, exactly as it would appear after
% "\coordinate (name) at " -- e.g. "(current page.north east)" or
% "($(current page.north east)+(-3.81cm,-4.61cm)$)", dollar signs and all
% when it's a calc expression. Resolved into a plain \coordinate first
% rather than substituted directly into each box's own ($...$) calc
% expression, since a calc expression can't contain another calc
% expression's $...$ delimiters nested inside it.
\newcommand{\gothamCircularFlowDiagram}[3]{%
	\pgfmathsetmacro{\gothamCFDarrowLen}{5*#3}
	\pgfmathsetmacro{\gothamCFDarrowWidth}{1.3*#3}
	\begin{scope}
		\coordinate (gothamCFDctr) at #1;
		\node[fill=c_primary_highlight, rounded corners=1pt, inner sep=4pt, align=center, font=\footnotesize, scale=#3] (gothamCFDcash) at ($(gothamCFDctr)+(0,#2)$) {Cash};
		\node[fill=c_primary_highlight, rounded corners=1pt, inner sep=4pt, align=center, font=\footnotesize, scale=#3] (gothamCFDraw) at ($(gothamCFDctr)+(#2,0)$) {Raw materials};
		\node[fill=c_primary_highlight, rounded corners=1pt, inner sep=4pt, align=center, font=\footnotesize, scale=#3] (gothamCFDfin) at ($(gothamCFDctr)+(0,-#2)$) {Finished goods};
		\node[fill=c_primary_highlight, rounded corners=1pt, inner sep=4pt, align=center, font=\footnotesize, scale=#3] (gothamCFDrec) at ($(gothamCFDctr)+(-#2,0)$) {Receivables};
		\draw[-{Stealth[length=\gothamCFDarrowLen pt]}, line width=\gothamCFDarrowWidth pt, c_primary_highlight] (gothamCFDcash) to[bend left=18] (gothamCFDraw);
		\draw[-{Stealth[length=\gothamCFDarrowLen pt]}, line width=\gothamCFDarrowWidth pt, c_primary_highlight] (gothamCFDraw) to[bend left=18] (gothamCFDfin);
		\draw[-{Stealth[length=\gothamCFDarrowLen pt]}, line width=\gothamCFDarrowWidth pt, c_primary_highlight] (gothamCFDfin) to[bend left=18] (gothamCFDrec);
		\draw[-{Stealth[length=\gothamCFDarrowLen pt]}, line width=\gothamCFDarrowWidth pt, c_primary_highlight] (gothamCFDrec) to[bend left=18] (gothamCFDcash);
	\end{scope}
}

% A \frac/\dfrac-stacked formula reads as cramped against a plain-text
% neighbor in an itemize -- it's visibly taller than a normal line (both
% above *and* below the line, since the fraction bar sits on the text
% baseline with numerator/denominator stacked around it), so both the
% fraction's own item *and* the item that follows it need a bit of extra
% headroom that a plain-text \item doesn't. Started as a manually-repeated
% \vspace{3pt} before each such \item (first on
% slides/u74-Example-Dynamic-Mattress-Company-Periods-current.tex);
% \tallitem below is that same convention factored into a single drop-in
% \item replacement, so a future fraction-heavy slide (or a change to the
% spacing amount) doesn't mean re-editing several slide files by hand.
% Takes an optional leading multiplier (default 1, i.e. 3pt): \tallitem[1.5]
% for 4.5pt, for a formula tall enough that the default doesn't read as
% enough room. Usage: replace \item with \tallitem both on any itemize line
% containing a stacked fraction *and* on the item immediately following it
% (same nesting level; skip this second replacement if that next item is
% itself already a \tallitem for its own fraction, or if the fraction item
% is the last one in its list) -- it's a plain substitute, not an
% environment, so it doesn't change how the surrounding list itself is
% structured. (Named \gothamFracItem originally; renamed 2026-07-21 to the
% more discoverable \tallitem, and given the optional multiplier and the
% following-item rule above at the same time -- no behavior change for any
% existing default-strength call site.)
% Ratio, not a frozen length (2026-07-22, direct instruction: "modify
% \tallitem and \tallitemnobullet to also use scaling relative to font
% size") -- the original \newlength+\setlength{...}{3pt} pair freezes to
% one fixed pt value the moment \setlength runs, the same issue (and same
% fix) as \gothamItemsepRatioI/II/III/IV further below: converted to a
% bare ratio macro, with the "em" unit written back on inside \tallitem's
% and \tallitemnobullet's own bodies (not here), so it's recomputed fresh
% against whatever font is actually live at each call site -- correct
% regardless of nesting level or a \footnotesize/\scriptsize override,
% unlike the old fixed 3pt (calibrated against one specific level's font,
% silently wrong at any other). 0.3 reproduces the original 3pt exactly
% at level 2's own font (10pt) -- the level \tallitem was first used at,
% on slides/u74-Example-Dynamic-Mattress-Company-Periods-current.tex (see
% this macro's own comment above) -- so no visual change there, while now
% correctly scaling up/down at every other level and font override.
\newcommand{\gothamTallItemBaseRatio}{0.3}
\newlength{\gothamTallItemBaseUnit}
\newlength{\gothamTallItemSpace}
\newcommand{\tallitem}[1][1]{%
	\setlength{\gothamTallItemBaseUnit}{\gothamTallItemBaseRatio em}%
	\setlength{\gothamTallItemSpace}{#1\gothamTallItemBaseUnit}%
	\vspace{\gothamTallItemSpace}\item}

% For a line that needs \tallitem's own extra headroom but no bullet at
% all -- a display equation continuing the previous \item's own sentence,
% rather than a new bulleted point of its own (see README.md's own
% \tallitem section, which already documents this exact combination as a
% rare, previously hand-written case: "\vspace{\gothamTallItemBaseUnit}\item[]").
% \tallitemnobullet (2026-07-21, direct instruction) factors that hand-
% written idiom into the same kind of drop-in replacement \tallitem
% already is for its own case, same optional multiplier: \tallitemnobullet[1.5].
\newcommand{\tallitemnobullet}[1][1]{%
	% Recomputes \gothamTallItemBaseUnit itself, fresh, rather than reusing
	% whatever \tallitem's own body last left it at -- since it's now a
	% per-call, live "em"-based value (see \tallitem's own comment above),
	% not the one-time preamble-load-time constant it used to be, relying
	% on call order/a preceding \tallitem to have set it correctly would
	% silently break the moment \tallitemnobullet is used on its own, or
	% at a different font size than whichever \tallitem last ran.
	\setlength{\gothamTallItemBaseUnit}{\gothamTallItemBaseRatio em}%
	\setlength{\gothamTallItemSpace}{#1\gothamTallItemBaseUnit}%
	\vspace{\gothamTallItemSpace}\item[]}

% flalign* (used to left-align a display equation inside an itemize --
% see CONVERSION_NOTES.md's "flalign* for left-aligning..." entry) inherits
% plain LaTeX's \abovedisplayskip/\belowdisplayskip, sized for a full
% document page, not a compact slide itemize -- reads as far too much blank
% space above/below the equation next to the surrounding item text (found
% 2026-07-21, direct instruction; most visible on
% slides/bwi-Pearson-Correlation-Definition-current.tex's own single-line
% equation, but the same oversized gap is present everywhere flalign* is
% used this way, e.g. slides/7se-Portfolio-Risk-Three-Assets-current.tex
% and slides/rb3-Asset-Allocation-Example-CAL-current.tex). \gothamTightDisplaySkips
% tightens both skips -- and their "short" variants, since amsmath picks
% between the two depending on how much text precedes the display, so
% setting only one pair left the gap unfixed whenever the preceding line
% happened to be short -- down to something that reads as part of the
% same list rather than a separate floating display. Usage: wrap the
% *real*, unmodified flalign* in a plain TeX group together with this
% command, `{\gothamTightDisplaySkips\begin{flalign*} ... \end{flalign*}}`,
% rather than a new same-named environment -- tried that first
% (`\newenvironment{gothamtightflalign*}{...\begin{flalign*}}{\end{flalign*}}`)
% and it broke LaTeX's own environment-nesting tracking (`\@currenvir`),
% surfacing as "\begin{flalign*} ... ended by \end{gothamtightflalign*}"
% and a fatal, PDF-less compile on the very first slide that used it. A
% plain group sidesteps that class of bug entirely since it never
% introduces a same-named wrapper environment.
\newlength{\gothamTightDisplaySkip}
\setlength{\gothamTightDisplaySkip}{3pt}
\newcommand{\gothamTightDisplaySkips}{%
	\abovedisplayskip=\gothamTightDisplaySkip
	\abovedisplayshortskip=\gothamTightDisplaySkip
	\belowdisplayskip=\gothamTightDisplaySkip
	\belowdisplayshortskip=\gothamTightDisplaySkip}

% Set logo
% \gothamInstituteLogoPath (a plain \newcommand storing an image path,
% same idiom as \gothamChapterLogoPath below) and \gothamLogoHeight (a
% length) are both defined by the per-university preamble that \input{}s
% this file, before this file is reached -- see e.g.
% Slidemaster_preamble_uni-graz.tex. Whichever source image
% \gothamInstituteLogoPath points at should have any symmetric margin of
% empty transparent pixels around the actual logo artwork trimmed off first
% (see Slidemaster_preamble_uni-graz.tex's own comment on why: an untrimmed
% margin makes \includegraphics's box bigger than the visible logo, so the
% logo's right-hand margin ends up looking bigger than its top/bottom
% margins, since only the latter get absorbed into vertical centering).
% \gothamLogoHeight is the single source of truth for the logo's rendered
% size: it drives the (invisible) in-flow logo below, the white patch's
% width further down, and the visible redrawn logo, so those three can never
% drift out of sync with each other again (that drift -- redrawing a smaller
% logo than the one the patch/position were sized for -- is what caused the
% left/right margins around the logo to come out uneven).

% The tikz node wrapper (remember picture, but not overlay, so it still
% behaves as normal inline content) records this logo's exact page position
% under the name "gothamLogoAnchor", so it can be redrawn on top of a white
% patch further down (see the "yellow box shorter than the page" block below).
\renewcommand{\gothamInstituteLogoSquare}[1][\gothamLogoHeight]{%
	\tikz[remember picture, baseline=(gothamLogoAnchor.base)]{%
		\node[inner sep=0pt, outer sep=0pt](gothamLogoAnchor){%
			\includegraphics[height=#1]{\gothamInstituteLogoPath}%
		};%
	}%
}

% \gothamFrametitleRightpading (left/right padding around the logo, chosen
% to match its top/bottom margins -- measured against the title colorbox's
% own top and bottom edges -- rather than the theme's 2.2ex default) is
% also set by the per-university preamble, not here: the top/bottom margins
% are a side effect of vertically centering a \gothamLogoHeight-tall logo in
% the title bar, not something set directly, so this value has to be
% re-measured and set per institution whenever its own logo/height differs
% from another's.

% Logo width, used below to size the white patch behind it.
\newlength{\gothamLogoWidth}
\settowidth{\gothamLogoWidth}{\includegraphics[height=\gothamLogoHeight]{\gothamInstituteLogoPath}}

% Let the yellow title bar visually stop short of the page's right edge, so
% the logo sits on white instead of on yellow: paint a white patch spanning
% the full height of the title bar behind the logo, then redraw the logo
% (already positioned correctly by \gothamFrametitleRightpading above, via
% the "gothamLogoAnchor" node captured above) on top of it. This needs two
% compilation passes: like any \tikz "remember picture", the node positions
% are only available in the .aux file written by the previous run.
% (The redraw uses \includegraphics directly rather than
% \usebeamertemplate{instituteLogo}/{includeLogo}: those templates rely on
% gotham's internal property-list lookup, which is only valid in the
% context of the original frametitle rendering and silently produces
% nothing when re-invoked here.)
% Additionally paints white across the *entire* width of the strip between
% the logo's bottom edge and the title bar's actual bottom edge, so the
% yellow ends flush with the logo's bottom (rather than extending further
% down under the title text too) while the box's width and top position
% stay unchanged. Checked against measured PDF coordinates that the title
% text's own (font-metric, not just this-string's-glyphs) lowest descender
% still clears the logo's bottom edge by about 1.3pt, so this doesn't clip
% text even on titles with descenders (e.g. "Highlighting...").
%
% Factored into its own macro, rather than written directly into the
% \addtobeamertemplate call below, so \gotham_no_inline_logo (a whole
% separate frametitle template variant, further down) can reuse it too --
% \addtobeamertemplate patches whatever template is *currently* installed
% at the time it's called (gotham's own default one here), so a template
% defined fresh afterwards doesn't inherit that patch automatically and
% needs to invoke this explicitly instead.
% Every edge-anchored node in this hook (and in every other page-corner
% overlay in this file) needs an explicit "outer sep=0pt" alongside its
% "inner sep=0pt": tikz's default outer sep is 0.5\pgflinewidth (0.2pt at
% the default 0.4pt line width), and a node's edge/corner anchors sit on
% the *outer* boundary -- outer sep past the filled rectangle -- so an
% anchor placed exactly on a page corner leaves the fill itself 0.2pt
% short of that corner on the anchored sides. Found 2026-07-16 (direct
% instruction: a hairline white line above/left of titleframe's and
% quoteframe's yellow boxes, their right edges 0.4pt off the standard
% title bar's own visible yellow edge, and their logos 0.2pt left of
% every other frame's) by measuring the compiled PDF's drawn rectangles:
% every affected fill/image was inset by exactly 0.199bp = 0.2pt from
% its anchor. This is also the same offset the footer box's own comment
% (further down) had earlier measured, worked around with a hand-nudged
% anchor, and left as "root cause not chased further" -- now chased:
% it *is* a general property of every edge/corner-anchored tikz node,
% not something specific to that one construction. (Center anchors are
% unaffected -- outer sep is symmetric around the center -- so the
% gothamLogoAnchor.center redraw below doesn't need it.)
\newcommand{\gothamFrametitleHookCode}{%
	\tikz[remember picture, overlay]\node(gothamTitleBottom){};%
	\begin{tikzpicture}[remember picture, overlay]
		\path let \p1=(current page.north east), \p2=(gothamTitleBottom) in
			node[fill=white, anchor=north east, inner sep=0pt, outer sep=0pt, name=gothamMainPatch,
			     minimum width=\gothamLogoWidth+2\gothamFrametitleRightpading,
			     minimum height=\y1-\y2] at (\x1,\y1) {};
		\ifx\gothamFrametitleFullHeightYellow\empty
			\path let \p1=(current page.north east), \p2=(gothamTitleBottom),
			          \p3=(gothamLogoAnchor.south), \p4=(current page.west) in
				node[fill=white, anchor=north west, inner sep=0pt, outer sep=0pt,
				     minimum width=\x1-\x4, minimum height=\y3-\y2] at (\x4,\y3) {};
		\fi
		% Normally redrawn at gothamLogoAnchor's own position (wherever the
		% invisible in-flow logo -- on the same line as the title text --
		% ended up). \gothamFrametitleLogoFixed instead redraws it at a fixed
		% page-corner position, independent of where the title text sits --
		% for a frame type (sectionframe) whose title is deliberately *not*
		% at the position a normal frame's title/logo line would put it.
		% gothamMainPatch already covers the full logo column from the page
		% top down to the box's true bottom edge (above), so it's already
		% blank there regardless of which position is used.
		\ifx\gothamFrametitleLogoFixed\empty
			\node[anchor=center, inner sep=0pt] at (gothamLogoAnchor.center) {\includegraphics[height=\gothamLogoHeight]{\gothamInstituteLogoPath}};
		\else
			% inner ysep is derived, not \gothamFrametitleRightpading again:
			% the inline logo this fixed redraw must line up with is
			% *vertically centered* in the title bar, and its resulting top
			% margin (8.14pt, measured) is not exactly the horizontal
			% padding (0.28cm = 7.97pt) -- close enough that the pre-outer-sep-fix
			% code's accidental extra 0.2pt masked the difference, but a
			% visible micro-jump when flipping pages once the fix landed.
			% Rather than a fresh calibrated constant:
			% \gothamTitleBarBottomAlign (the visible bar bottom) *is* the
			% inline logo's bottom edge, so subtracting \gothamLogoHeight
			% reproduces its top margin exactly, and re-measuring one
			% constant keeps both in sync if the bar ever changes.
			\node[anchor=north east, inner xsep=\gothamFrametitleRightpading,
			      inner ysep=\dimexpr\gothamTitleBarBottomAlign-\gothamLogoHeight\relax,
			      outer sep=0pt] at (current page.north east) {\includegraphics[height=\gothamLogoHeight]{\gothamInstituteLogoPath}};
		\fi
		% Optional per-frame "chapter logo" (see \chapterlogo{} below): drawn
		% immediately to the left of the main logo's white patch, on its own
		% white patch, only when \gothamChapterLogoPath is non-empty. Sized
		% with exactly one \gothamFrametitleRightpading of its own padding, on
		% its outer (left) side only -- the inner (right) side is left flush
		% against gothamMainPatch, since that patch already contributes its
		% own \gothamFrametitleRightpading of padding around the main logo.
		% This makes both gaps -- chapter logo to main logo, and chapter logo
		% to where the yellow bar resumes -- come out equal, at exactly one
		% \gothamFrametitleRightpading each, without double-counting or
		% needing a second length constant.
		\ifx\gothamChapterLogoPath\empty\else
			\settowidth{\gothamChapterLogoWidth}{\includegraphics[height=\gothamLogoHeight]{\gothamChapterLogoPath}}%
			% The patch's right edge is nudged 0.5pt *past* gothamMainPatch's own
			% west edge (i.e. slightly overlapping it) rather than exactly
			% touching it -- both are solid white, so the overlap itself is
			% invisible, but touching them exactly left a hairline yellow seam
			% at the join (a real, if sub-millimetre, gap that persisted across
			% recompiles, most likely PDF fill-antialiasing at the shared edge
			% of two separately-filled abutting shapes). The logo image itself
			% (below) is *not* shifted, so the actual margin it defines is
			% unaffected.
			\path let \p1=(current page.north east), \p2=(gothamTitleBottom) in
				node[fill=white, anchor=east, inner sep=0pt, outer sep=0pt,
				     minimum width=\gothamChapterLogoWidth+\gothamFrametitleRightpading+0.5pt,
				     minimum height=\y1-\y2] at ($(gothamMainPatch.west)+(0.5pt,0)$) {};
			\node[anchor=east, inner sep=0pt, outer sep=0pt] at (gothamMainPatch.west |- gothamLogoAnchor.center)
				{\includegraphics[height=\gothamLogoHeight]{\gothamChapterLogoPath}};
		\fi
	\end{tikzpicture}%
}
\addtobeamertemplate{frametitle}{}{\gothamFrametitleHookCode}

% \chapterlogo{path} opts a single upcoming standardframe into an extra logo
% (e.g. a topic/chapter icon) next to the main institute logo -- see the
% frametitle hook above for how it's drawn. Call it *before* \begin{standardframe},
% since the frame title (and this hook) is typeset as part of \begin{frame}
% itself, before any frame-body content would run. standardframe resets this
% back to empty on \end{frame} (below), so it only ever applies to the one
% frame that follows the \chapterlogo{} call, never leaking into the next slide.
% Both the set and the reset use \global\def (checked empirically, via
% \typeout, against the frametitle hook's actual view of this macro on a
% following frame): beamer's frame environment opens its own TeX group around
% the frame body, so a plain (local) assignment done in standardframe's
% closing code -- which runs right after \end{frame} -- was silently undone
% once that group actually closed, reviving the *previous* value instead of
% the reset. \global\renewcommand alone wasn't enough to fix it (still leaked
% into the next frame) -- only dropping to the lower-level \global\def did.
\newlength{\gothamChapterLogoWidth}
\newcommand{\gothamChapterLogoPath}{}
\newcommand{\chapterlogo}[1]{\global\def\gothamChapterLogoPath{#1}}

% \gothamFrametitleFullHeightYellow: opt-out for the frametitle hook's
% "trim yellow to the logo's bottom edge" white patch (see the hook above) --
% empty (default) keeps that trim; non-empty skips it, letting the yellow
% show all the way down to the title box's own true bottom edge instead.
% That trim exists to erase a normally-tiny (~1-2pt) sliver of yellow that
% would otherwise stick out below the logo on an ordinary short title bar --
% harmless there, but sectionframe's deliberately tall bar turns that
% "tiny sliver" into most of the box, silently wiping out nearly all of the
% intended yellow with white. sectionframe sets this (globally, same
% set/reset reasoning as \gothamChapterLogoPath above); standardframe resets
% it back to empty on \end{frame}, so it's opt-in per frame like the others.
\newcommand{\gothamFrametitleFullHeightYellow}{}

% \gothamFrametitleLogoFixed: opt-in for the frametitle hook's logo redraw
% (see the hook above) to use a fixed page-corner position instead of
% wherever the title text's own line puts it. Normally the two coincide (the
% logo sits inline with the title, gotham's own design), which is exactly
% why this is off by default -- but sectionframe's title is deliberately
% *not* where a normal frame's title sits (bottom-anchored within a much
% taller bar), and letting the logo follow it down there would put it
% somewhere no other slide in the deck ever has it. Same set/reset reasoning
% and mechanism as \gothamChapterLogoPath/\gothamFrametitleFullHeightYellow
% above.
\newcommand{\gothamFrametitleLogoFixed}{}

% \gothamFooterBoxColor: the footer box's own fill color (see the footline
% hook further down), overridable per frame the same set-before/reset-after
% way as the flags above -- default \gothamFooterBoxColorDefault (set by the
% per-university preamble before \input{}ing this file, same "value
% supplied externally, referenced lazily" idiom as \gothamBulletColor/
% \gothamInstituteLogoPath), matching every frame's title bar by default;
% highlightframe overrides it to match its own (likewise overridable) title
% bar color, \gothamHighlightColor, so both boxes change together with a
% single retargeting.
\newcommand{\gothamFooterBoxColor}{\gothamFooterBoxColorDefault}

% A frametitle template variant that drops the inline institute logo
% entirely -- for use together with \gothamFrametitleLogoFixed (which
% redraws the logo separately, at a fixed position, making the inline one
% pointless anyway). gotham's own template (beamerouterthemegotham.sty,
% "default/gotham~subsameline") puts \nolinebreak\hfill{...logo...} directly
% after \insertframetitle; this variant drops that whole part.
%
% Doing so, though, means the logo no longer reserves *any* horizontal room
% on the title's own line the way it used to -- so rightskip here widens
% \gothamFrametitleRightpading to reserve the same width gothamMainPatch
% (the frametitle hook's white patch, drawn afterward) already reserves for
% the logo column: \gothamLogoWidth + 2*\gothamFrametitleRightpading. Without
% this, the title paragraph wraps as if that whole column were free text
% space -- not an actual TeX overfull-hbox condition (no warning, the text
% genuinely fits the *box's own* width), just a title that runs visually
% underneath the white patch and logo drawn on top of it afterward, which
% looks identical to clipped/hidden text. (An earlier, wrong diagnosis of
% this same visual symptom blamed gotham's \nolinebreak forcing an overfull
% line -- checked by testing gotham's original template installed with a
% long \huge title and no changes at all, which reproduces something that
% *looks* the same; but the actual fix needed here turned out to be this
% rightskip, not removing \nolinebreak, which was really addressing a red
% herring from an isolated test that happened not to include a competing
% logo column at all.)
%
% "plus 1fil" on that rightskip keeps every line ragged-right (its own
% natural width, however short of the margin that leaves) rather than fully
% justified -- without it, a wrapped multi-line title's non-last lines get
% stretched to exactly fill the available width, which reads as oddly wide,
% uneven word-gaps on a short title (most noticeable exactly when a title is
% long enough to wrap, i.e. always, for this template).
\makeatletter
\defbeamertemplate{frametitle}{gotham_no_inline_logo}{%
	\nointerlineskip
	\begin{beamercolorbox}[wd=\paperwidth, sep=0pt,
	                        leftskip=\gothamFrametitleLeftpading,
	                        rightskip=\dimexpr\gothamLogoWidth+2\gothamFrametitleRightpading\relax plus 1fil]{frametitle}
		\gotham@frametitlestrut@start
		\insertframetitle
		\ifx\insertframesubtitle\@empty\else
			{\usebeamerfont{framesubtitle}\usebeamercolor[fg]{framesubtitle}\gothamFrameSubtitleSep \insertframesubtitle}%
			\vspace*{\gothamFramesubtitleStrutend}%
		\fi
		\gotham@frametitlestrut@end
	\end{beamercolorbox}
	\gothamFrametitleHookCode
}
\makeatother


% Define slideID as a simple value holder (initially empty)
\newcommand{\slideID}{}%

% Define slideID as a simple value holder (not a command with arguments)
\makeatletter
\newcommand{\setSlideID}[1]{%
	% Grab everything up to the first '-' and store it in \slideID
	\edef\tempa{#1}%
	\expandafter\setSlideID@aux\tempa-\relax
}
\def\setSlideID@aux#1-#2\relax{%
	\renewcommand{\slideID}{#1}%
}
\makeatother

% Define a new \slide command to automatically pick the current version of the slide
% 2026-07-08 (direct instruction): slide filenames dropped their leading
% "s-" (e.g. slides/ab3-Some-Slide-current.tex, not the old
% slides/s-ab3-Some-Slide-current.tex) so slide IDs are directly
% searchable in Explorer -- this macro (and the other three below) is the
% *only* place that ever needed to change for that: every deck's own
% \slide{}/\slideversion{}/\omitslide{}/\omitslideversion{} calls pass just
% "ID-slug", never the filename itself, so renaming every file on disk to
% match needed no changes anywhere else (not the deck files, not
% slide_ids.csv, which only ever stored bare IDs to begin with).
\newcommand{\slide}[1]{%
	\setSlideID{#1}%
	\input{slides/#1-current.tex}%
}

% \slidefinal{ID-slug} -- like \slide{}, but renders the slide at its final,
% fully-revealed overlay state in every deck version, instructor included
% (2026-07-12, direct instruction: show jww's finished build in place of a
% separate recap slide, "irrespective of which deck version is being
% compiled" -- that recap slide has since been deleted entirely, its own ID
% no longer in use). A slide
% built with \conditionalpause[instructor] already renders fully built in
% the during-course/after-course compiles -- that pause is a no-op for
% them, so nothing there ever waits on a reveal step -- the only case
% actually needing help is the instructor deck, which would otherwise still
% step through the build. Locally shadowing \conditionalpause to a no-op
% (same optional-argument signature as the real one, so it still accepts
% e.g. \conditionalpause[instructor] without erroring) for the duration of
% one \slide{} call makes every \conditionalpause in that one inclusion
% inert, regardless of the actual compiled \gothamDeckVersion -- scoped
% inside \begingroup/\endgroup, so it reverts immediately after, and has no
% effect on any *other* \slide{}/\slidefinal{} call to the same ID
% elsewhere in the same deck (e.g. an earlier, normally-building instance of
% the same slide). \deckonly/\deckuncover/\deckvisible are untouched, so
% deck-version-tagged content (as opposed to reveal-step-tagged content)
% still follows its usual per-version rules.
\newcommand{\slidefinal}[1]{%
	\begingroup
		\renewcommand{\conditionalpause}[1][instructor]{}%
		\slide{#1}%
	\endgroup
}

% Load a specific, pinned version of a slide instead of whichever is "current"
% Usage: \slideversion{ID-slug}{M-m-p} loads slides/ID-slug-vM-m-p.tex
\newcommand{\slideversion}[2]{%
	\setSlideID{#1}%
	\input{slides/#1-v#2.tex}%
}

% Student-version decks: \studentversion (called anywhere in a deck, e.g.
% right after % Slidemaster preamble variant for University of Graz decks. Institution-
% specific content only (colors, logo, logo calibration) -- every frame
% type, the deck-version system, footer/hyperlink/footnote styling, tikz
% styles, etc. live in Slidemaster_preamble_base.tex, \input{}ed at the
% bottom of this file, and are shared verbatim across every institution
% this project supports (see Slidemaster_preamble_mci.tex and
% Slidemaster_preamble_uni-ibk.tex for the other two). A deck picks its
% institution by \input{}ing this file instead of one of those, nothing
% else changes.

% Color definition (Graz palette, sourced from the original single-preamble
% version of this project -- these are the University of Graz's own official
% colors, not derived/approximated).
\definecolor{c_primary_dark}{HTML}{000000}
\definecolor{c_primary_highlight}{HTML}{ffd500}
\definecolor{c_primary_gray}{HTML}{c6c6c6}
\definecolor{c_primary_white}{HTML}{FFFFFF}
\definecolor{c_secondary_violet}{HTML}{948dc3}
\definecolor{c_tertiary_violet}{HTML}{c2bddf}
\definecolor{c_secondary_blue}{HTML}{689cce}
\definecolor{c_tertiary_blue}{HTML}{bacce8}
\definecolor{c_secondary_green}{HTML}{9fcb81}
\definecolor{c_tertiary_green}{HTML}{cce2ba}
\definecolor{c_secondary_red}{HTML}{e4684b}
\definecolor{c_tertiary_red}{HTML}{f1ae95}
\definecolor{c_secondary_orange}{HTML}{f6a968}
\definecolor{c_tertiary_orange}{HTML}{fbd0aa}
\definecolor{c_secondary_yellow}{HTML}{ffe063}
\definecolor{c_tertiary_yellow}{HTML}{ffeeaf}
\definecolor{c_tertiary_gray}{HTML}{e3e3e3}
% Not an official Uni Graz color (hence "custom", not "secondary" like the
% palette above) -- a softer near-black than c_primary_dark, used for
% quoteframe's full-bleed dark block, matching the specific dark gray (not
% pure black) ExampleSlides.pptx used for that same block. Also reused for
% hyperlink text color, for the same reason it works for quoteframe's
% attribution text: dark enough for good contrast/legibility on a light
% background.
\definecolor{c_custom_gray}{HTML}{333333}
\definecolor{c_custom_red}{HTML}{C00000}
\definecolor{c_custom_green}{HTML}{00B050}

% Logo + calibration.
% schemes/scheme_uni_graz/Logo_Uni-Graz_4c-cropped.pdf -- official vector artwork,
% user-supplied (schemes/scheme_uni_graz/Logo_Uni-Graz_4c.eps, kept alongside as the
% untouched source), converted to PDF via epstopdf and cropped to its true
% ink bounding box (computed with Ghostscript's own `-sDEVICE=bbox`, same
% treatment as MCI's and Innsbruck's own logos -- see
% Slidemaster_preamble_mci.tex's own comment for the fuller explanation of
% why the EPS's own %%BoundingBox isn't trusted directly). Vector, not
% raster, so it stays crisp at any zoom. Supersedes this file's own earlier
% raster logo (schemes/scheme_uni_graz/Logo Uni Graz 4c-cropped.png, itself cropped
% from schemes/scheme_uni_graz/Logo Uni Graz 4c.png's own symmetric transparent
% margin -- both kept alongside as historical/fallback assets, not deleted).
\newcommand{\gothamInstituteLogoPath}{schemes/scheme_uni_graz/Logo_Uni-Graz_4c-cropped.pdf}
\newlength{\gothamLogoHeight}
\setlength{\gothamLogoHeight}{3ex}
% Re-measured against the compiled PDF for this specific logo image -- see
% Slidemaster_preamble_base.tex's own comment on \gothamFrametitleRightpading
% for why this needs re-measuring per institution rather than reusing
% another one's value.
\setlength{\gothamFrametitleRightpading}{0.28cm}

% Itemize bullet color (2026-07-16, direct instruction): Graz keeps the
% project's original behavior, bullets in the primary highlight color.
\newcommand{\gothamBulletColor}{c_primary_highlight}

% Title bar and footer box/text colors (2026-07-16, direct instruction):
% Graz keeps the project's original behavior throughout -- title/footer
% boxes in the primary highlight color, title/footer text in the primary
% dark color (unlike MCI/Innsbruck, which switch these to their own
% secondary blue/white -- see Slidemaster_preamble_mci.tex's own comment).
\newcommand{\gothamTitleBoxColor}{c_primary_highlight}
\newcommand{\gothamTitleTextColor}{c_primary_dark}
\newcommand{\gothamFooterBoxColorDefault}{c_primary_highlight}
\newcommand{\gothamFooterTextColor}{c_primary_dark}

\usepackage{etoolbox} % \AtEndEnvironment, used for the itemize level-drop spacing below
\usepackage{tikz} % For painting a white box over title bar
\usetikzlibrary{calc}
\usetikzlibrary{arrows.meta} % for e.g. arr/.style={-{Stealth[...]}...} on diagram slides
\usetikzlibrary{positioning} % for `right=<dist> of <node>`-style relative node placement
\usetikzlibrary{decorations.text} % text-along-a-path, for a circular photo's curved caption
\usepackage{booktabs} % \toprule/\midrule/\bottomrule for gothamtable
\usepackage{colortbl} % \arrayrulecolor{}, used by gothamtable to color those rules
\usepackage{pifont} % \ding{51}/\ding{55} checkmark/X-mark glyphs, e.g. slides/fr3-Forms-of-Capital-Market-Efficiency-Strong-Semistrong-Weak-current.tex's own validity ratings
\usepackage{pgfplots} % real-data line/bar charts (e.g. slides/3vp-Metallgesellschaft-Oil-Price-Chart-current.tex) -- confirmed safe alongside this deck's own \tableofcontents, see CONVERSION_NOTES.md's "beamer/TeX gotchas" entry
\pgfplotsset{compat=1.18}

% Restores genuine Computer Modern math fonts. Found 2026-07-15, direct
% instruction ("math equations do not have serifs, look at vcv"): despite
% fontspec being loaded with `no-math` (beamerfontthemegotham.sty), math
% still wasn't landing on Computer Modern -- checked directly against a
% compiled PDF's own per-glyph font names (pymupdf), not just eyeballed:
% digits/math-italic letters were coming out as FiraSans-Light(Italic)
% (this deck's own body-text sans font) and operators/parens as CMSS10
% (Computer Modern's *sans-serif* shape, not its serif one). Root cause:
% beamer's own base classes redefine the math alphabets to follow the
% presentation's sans-serif text font by default, specifically so math
% matches an all-sans deck -- `no-math` only stops *fontspec* from doing
% this, it has no say over beamer's own separate override. Beamer ships
% the standard fix for exactly this case: a font theme named "serif"
% with an "onlymath" option that restores ordinary (Latin
% Modern/Computer-Modern-equivalent) math fonts without touching any of
% the surrounding sans-serif text. Confirmed directly afterward: every
% glyph in a re-checked formula now reports as LMRoman*/CMR10/CMMI8 (all
% genuinely serif), not FiraSans/CMSS10.
\usefonttheme[onlymath]{serif}

% \text{} (amsmath, used throughout for plain-word operands inside math,
% e.g. \text{E}[\tilde{R}_M]) is a separate mechanism from the math
% alphabets \usefonttheme[onlymath]{serif} above just fixed: unlike
% \mathrm/\mathit/etc., \text{} is defined to typeset in whatever the
% *surrounding running text* font is at that point in the document, not a
% math alphabet -- so on this deck's all-Fira body text, \text{} inside a
% formula still comes out in Fira Sans (and picks up italic if the
% enclosing text happens to be italicized), even though the digits/
% operators around it are correctly Computer Modern. Found 2026-07-20 on
% slides/2gd-Analyzing-Reward-to-Risk-Ratio-CAPM-current.tex's own
% \text{E}/\text{Cov} operands. Fixed generally, project-wide, by wrapping
% \text{}'s argument in an explicit upright-Computer-Modern font switch --
% \usefont{OT1}{cmr}{m}{n} rather than \mathrm{} (which would pull in
% \text{}'s own math-alphabet substitution logic unnecessarily) or
% \textrm{} (which switches *family* only, so still comes out italic
% inside an italicized context, which \text{} should never do here) --
% deliberately hardcoding "not italic" (shape n) rather than inheriting
% the ambient shape, since \text{}'s whole purpose in this deck's formulas
% is plain upright labels (E, Cov, "the" CAPM, ...), regardless of
% whatever font context happens to surround the formula itself.
\let\gothamOrigText\text
\renewcommand{\text}[1]{\gothamOrigText{\usefont{OT1}{cmr}{m}{n}#1}}

% Color definition, institute logo path, and the two logo-calibration
% lengths (\gothamLogoHeight, \gothamFrametitleRightpading) all live in the
% per-university preamble that \input{}s this file, not here -- this file
% is shared machinery only, no institution-specific branding. Every color
% token below (\c_primary_*, \c_secondary_*, \c_tertiary_*, \c_custom_*) and
% \gothamInstituteLogoPath/\gothamLogoHeight/\gothamFrametitleRightpading
% are assumed already defined by the time this file is reached -- see e.g.
% Slidemaster_preamble_uni-graz.tex for the Graz values, or
% Slidemaster_preamble_mci.tex/Slidemaster_preamble_uni-ibk.tex for the
% other two institutions this deck supports. \gothamLogoWidth is the one
% logo-derived length that *does* stay computed here (via \settowidth,
% further below) rather than hardcoded per institution, since it already
% adapts automatically to whatever image \gothamInstituteLogoPath points at.

% Style hyperlinks (\href{url}{text}) consistently across all slides:
% underlined, so a link is visually recognizable as such against plain
% body text -- yellow (this project's usual accent color) was tried first
% for the *text* itself, but yellow-on-white body text is hard to read,
% unlike yellow used as an accent against dark/black text elsewhere in the
% theme. Went through a few more candidates on 2026-07-06 (all direct
% instruction) before settling on c_secondary_blue for both text
% and underline (a real color reads more clearly as "this is a link" than
% any gray does, without yellow's legibility problem).
%
% 2026-07-07, direct instruction: split text color and underline color
% instead of using one color for both -- link text back to
% c_custom_gray (the near-black dark gray also used for
% quoteframe's attribution text), underline in a separate accent color.
% Plain \underline can't do this: its rule is drawn in whatever color is
% current at that point, so a single \underline{...} call always
% underlines in its own text's color, never a second, different one.
% \gothamHrefUnderline below reimplements it instead: measure the linked
% text's own width into a savebox, print that box in customgray, then
% back up by the box's own width and draw a plain \rule of the same
% width, raised to sit just below the text's baseline, in the accent
% color. Same single-line limitation as plain \underline (and the reason
% \url{} is used instead of \href on slides/acw-Some-Topics-in-Testing-the-EMH-current.tex --
% see that slide's own comment): a \rule, like \underline, can't span a
% line break either, so this only ever wraps normal, comparatively short
% \href display text, never something that needs mid-word breaking.
%
% First tried c_primary_highlight for the underline (same day) -- confirmed
% via the compiled PDF that the rule *was* being drawn (right color,
% position, width), so not a compile-time bug, but yellow (very high
% luminance, close to white) at this rule's 0.4pt thickness reads as
% essentially invisible against the white slide background, worst at
% footnote size (see e.g. slides/za2-Startup-Financing-Categorization-current.tex's
% own long inline URL) where it's easy to miss entirely. Switched
% 2026-07-07 (direct instruction) to c_secondary_blue instead --
% far more contrast against white at the same thickness -- and the rule
% was also thickened at the same time (0.4pt -> 0.6pt), on the reasoning
% that it might need it for the same contrast reason. It didn't: blue
% already reads clearly at the original thickness, and 0.6pt instead
% just looked heavier than intended (direct feedback 2026-07-08, e.g. on
% slides/6qi-Yield-Spreads-Corporate-Bonds-current.tex). Reverted to
% 0.4pt -- only the color was ever actually asked to change.
\hypersetup{colorlinks=true, urlcolor=c_custom_gray, linkcolor=c_custom_gray}
\let\gothamOrigHref\href
\newsavebox{\gothamHrefTextBox}
%
% 2026-07-08 (direct instruction): the underline's own rule thickness now
% scales with the ambient font size instead of a single fixed 0.2pt for
% every link everywhere -- 0.2pt at this deck's footnote size, 0.4pt at its
% normal body-text size, interpolated (and extrapolated, for any other
% size a link happens to appear at) linearly in between.
%
% Deriving "current font size" from \f@size (the textbook NFSS way) was
% tried first and confirmed broken specifically *inside* a beamer frame --
% \typeout{\meaning\f@size} right after this preamble loads shows a normal,
% defined macro (`macro:->10.95`), but the exact same check inside
% `\begin{standardframe}...\end{standardframe}` reports `\f@size` as
% *undefined*, presumably from beamer's own overlay-aware font-switching
% machinery not keeping classic NFSS's own current-size bookkeeping in
% sync the way plain (non-beamer) LaTeX does. Since every real use of this
% macro happens inside a frame, that path was a dead end regardless of how
% well it worked in isolation.
%
% Measuring a reference glyph's own rendered height instead sidesteps this
% entirely: building a box and reading back its own height just measures
% whatever font is *actually* active at the call site, independent of
% whether beamer's own internal "current size" bookkeeping happens to
% agree -- and within one font
% family, every glyph's metrics scale by exactly the same ratio the point
% size does, so a measured height is just as valid a stand-in for "current
% size" as \f@size would have been, had it worked. "M" (present in every
% context this is used in) is the reference glyph; \gothamHrefUnderlineFnHeight/
% \gothamHrefUnderlineBodyHeight are its height at this deck's footnote font
% and normal body font respectively, each measured once, here, in the
% relevant font context.
% Deliberately built from raw TeX primitives (\newdimen/\setbox/\ht),
% *not* \newlength/\setlength/\settoheight -- found, while building this,
% that loading this deck's theme silently turns \setlength and
% \settoheight into something else entirely: `\meaning\setlength` reads
% `macro:->\protect \setlength` right after `\usetheme{gotham}` (still the
% ordinary kernel definition beforehand), and the same for `\settoheight`.
% Not chased to which file does this or why (bisected only as far as
% "plain beamer alone doesn't, beamer+gotham does" -- most likely
% `translator`, which beamer loads for localized strings and is known to
% robustify commands this way, but not confirmed) -- what matters
% practically is that a *global* assignment through either one silently
% no-ops instead of erroring, confirmed directly and in isolation (a
% minimal beamer+gotham file, no other content): `\global\setlength{\len}
% {\otherlen}` leaves `\len` at 0pt after the group closes, while the
% exact same pattern with plain `\newdimen`/raw `=` assignment propagates
% correctly. Every plain (non-global) `\setlength` elsewhere in this
% project's own files is unaffected -- it's specifically the *global*
% form that breaks, which is why this is the first place it's shown up
% despite `\setlength` being used constantly throughout this preamble.
\newdimen\gothamHrefUnderlineFnHeight
\newdimen\gothamHrefUnderlineBodyHeight
\newdimen\gothamHrefUnderlineNowHeight
\newdimen\gothamHrefUnderlineThickness
% A dedicated box register, not the commonly-borrowed box0 -- several
% slides in this deck (the "Long-term bonds" tables) already use box0
% themselves (`\sbox0{...}`/`\usebox0`) for their own, unrelated
% mid-table-layout purposes, and a link can appear anywhere, including
% inside one of those tables -- reusing box0 here would silently clobber
% whatever they were holding there.
\newbox\gothamHrefUnderlineBox
% The two reference heights are (re)measured on every single \href call,
% not once at preamble-load time -- tried once-at-load first and found a
% second, independent problem with it: `\usebeamerfont{footnote}` right
% here (before `\begin{document}`) silently has no effect at all (footnote-
% size "M" and normal-size "M" measured identically, confirmed directly),
% presumably because beamer's font-theme machinery isn't fully installed
% until the document starts. Since every real `\href` call happens well
% inside the document (inside some frame), measuring fresh each time
% sidesteps that timing question entirely rather than needing to find the
% right hook to measure once, early enough to be free but late enough to
% be correct -- the extra cost is two throwaway one-glyph boxes per link,
% not something worth optimizing away in a slide deck.
\newcommand{\gothamHrefUnderlineMeasureRefHeights}{%
	{\usebeamerfont{footnote}\setbox\gothamHrefUnderlineBox=\hbox{M}\global\gothamHrefUnderlineFnHeight=\ht\gothamHrefUnderlineBox}%
	{\normalsize\setbox\gothamHrefUnderlineBox=\hbox{M}\global\gothamHrefUnderlineBodyHeight=\ht\gothamHrefUnderlineBox}%
}
% A plain e-TeX \dimexpr can multiply a dimension by another dimension and
% then divide by a third, as one combined operation, as long as the
% *middle* operand of the chain is itself wrapped in its own nested
% \dimexpr...\relax -- confirmed directly (an expl3 \dim_set:Nn attempt
% first, using \seq_.../\fp_... the way \setfootnotesymbols above does,
% kept hitting "Missing ) inserted for expression"/type-mismatch errors:
% dividing one dimension by another is a *ratio*, not a dimension, and
% \dim_set:Nn's own expression grammar doesn't implicitly convert between
% the two the way this nested-\dimexpr form does). No expl3 needed here at
% all -- plain e-TeX arithmetic, already available in any current LaTeX.
\newcommand{\gothamHrefUnderlineSetThickness}{%
	\gothamHrefUnderlineThickness=\dimexpr 0.2pt +
		(\gothamHrefUnderlineNowHeight - \gothamHrefUnderlineFnHeight)
		* \dimexpr0.2pt\relax
		/ (\gothamHrefUnderlineBodyHeight - \gothamHrefUnderlineFnHeight)\relax
}
\newcommand{\gothamHrefUnderline}[1]{%
	\sbox{\gothamHrefTextBox}{#1}%
	\textcolor{c_custom_gray}{\usebox{\gothamHrefTextBox}}%
	\hspace{-\wd\gothamHrefTextBox}%
	\gothamHrefUnderlineMeasureRefHeights
	\setbox\gothamHrefUnderlineBox=\hbox{M}\gothamHrefUnderlineNowHeight=\ht\gothamHrefUnderlineBox
	\gothamHrefUnderlineSetThickness
	\raisebox{-0.35ex}[0pt][0pt]{\textcolor{c_secondary_blue}{\rule{\wd\gothamHrefTextBox}{\gothamHrefUnderlineThickness}}}%
}
\renewcommand{\href}[2]{\gothamOrigHref{#1}{\gothamHrefUnderline{#2}}}
% \url{} (the `url` package hyperref depends on/patches) defaults to a
% monospace-ish "tt" style regardless of the above -- used deliberately in
% only one place in this deck,
% slides/acw-Some-Topics-in-Testing-the-EMH-current.tex, specifically
% *for* its line-breaking-at-slashes behavior (a plain \href with the
% URL as its own display text has no break points at all -- see that
% slide's own comment). Every other same-text-as-target URL in the deck
% (including slides/bb6-A-Perpetual-Bond-current.tex's own video link,
% which used to be a \url{} too) uses plain \href{url}{url} instead, since
% \url{} doesn't pick up this deck's custom colored underline below (its
% breakable content can't be measured into the savebox the underline
% needs -- see that mechanism's own comment) and most of these URLs are
% short enough to not need the breaking \url{} exists for anyway. 2026-07-07,
% direct instruction: that one genuinely-needed \url{} (on s-acw) should
% still look like this deck's other hyperlinks, i.e. the surrounding body
% font, not \url{}'s own default monospace.
% `\urlstyle{same}` (a standard url-package switch) does exactly that,
% with no effect on \url{}'s own breaking behavior. Set deck-wide, not
% only on that one slide, so any future \url{} use matches by default
% too, rather than needing to remember this each time. Coloring already
% matches automatically (hyperref's own \url respects
% \hypersetup{urlcolor=...} the same way \href does); the underline
% \href gets above doesn't carry over, since \underline's own argument
% can't contain line breaks -- wrapping \url{}'s necessarily-breakable
% text in \underline would silently undo the breaking this exists for,
% so that one difference is accepted rather than "fixed."
\urlstyle{same}

% gotham's own footnote color (beamercolorthemegotham.sty:
% \setbeamercolor{footnote}{fg=normal~text.fg!90}) is a 90%-black gray --
% against plain black body text, customgray links contrast clearly (found
% while investigating a 2026-07-05 report that a footnote's own \href
% didn't read as a link the way body-text ones like the "Moodle" link on
% slides/hxr-Administrative-Details-current.tex do); against gotham's own
% dimmed footnote gray, the same customgray is only 25/255 lighter and
% barely distinguishable at footnote size, with only the underline giving
% it away. Overridden to plain black here so a footnote link contrasts the
% same way a body-text one does -- a one-line override of the theme's own
% color (not editing beamercolorthemegotham.sty itself, which stays
% untouched per its own upstream-file status), and a plausible general
% legibility improvement for footnote text regardless of links.
\setbeamercolor{footnote}{fg=c_primary_dark}

% Footer text color (\FooterText/slide id/frame number) -- \gothamFooterTextColor,
% set by the per-university preamble before \input{}ing this file (the
% footer box's own background color, \gothamFooterBoxColorDefault above, is
% a separate, independently configurable setting). No explicit
% \setbeamercolor{footline}{...} existed before this -- the footer text
% simply inherited whatever "normal text" foreground already was (black),
% which happened to read fine against every frame's own yellow/orange
% footer box so far, but wouldn't automatically read against an arbitrary
% future footer color (e.g. a dark blue) without this override.
\setbeamercolor{footline}{fg=\gothamFooterTextColor}

\colorlet{colorDark}{c_primary_dark} % FG in light/normal mode
% Title bar background -- \gothamTitleBoxColor, set by the per-university
% preamble before \input{}ing this file (same idiom as
% \gothamBulletColor/\gothamFooterBoxColorDefault above).
\colorlet{colorA}{\gothamTitleBoxColor} % frametitle, standin.out

\gothamset{background=light} % Needs to go after color redefinitions, so they are applied
\setbeamercolor{background canvas}{bg=white}

% gotham's own default behavior on \section{} (used for a deck's chapter
% grouping/table of contents) is to auto-insert its *own* plain, unstyled
% section-divider frame plus a small "current section highlighted in the
% outline" frame, every time -- redundant with (and visually inconsistent
% next to) this project's own hand-styled sectionframe (see README.md's
% "Special frame types" section), which every chapter in this deck already
% opens with as an explicit \slide{}. Disabled so \section{} only feeds
% \tableofcontents (see slides/nkj-Agenda-current.tex) without also
% inserting its own extra frames.
\gothamset{sectionframe default=false}

% Same reasoning, for \subsection{} (2026-07-15, direct instruction):
% gotham's default subsection-divider frame is disabled too, so
% \subsection{} calls still register in the PDF's own bookmark/outline
% structure (Acrobat's "chapters" pane -- a plain PDF sectioning feature,
% independent of whether beamer auto-inserts a frame for it) and still feed
% \tableofcontents (see slides/nkj-Agenda-current.tex's own
% `\tableofcontents[hideallsubsections]`, which then hides them from that
% one slide specifically), without also spawning their own unstyled divider
% slide in the middle of a chapter every time one is added.
\gothamset{subsectionframe default=false}

% Title text color -- \gothamTitleTextColor, set by the per-university
% preamble before \input{}ing this file (the title bar's own background
% color, \gothamTitleBoxColor/colorA above, is a separate, independently
% configurable setting).
\setbeamercolor{frametitle}{fg=\gothamTitleTextColor}

% A color leak from \titleframe/\quoteframe into whichever frame follows
% them is fixed in standardframe's own opening code below (search
% "Forces this frame's own body") rather than here -- an
% \AtBeginEnvironment{frame}{\color{c_primary_dark}} hook was tried first
% and confirmed, by isolated reproduction, to have *no effect at all*:
% beamer's own "frame" isn't a plain \newenvironment etoolbox's
% \AtBeginEnvironment can patch that way, so the hook silently never fired.

% Set content text margins
\setbeamersize{text margin left=.45cm, text margin right=.45cm}

%\colorlet{colorPale}{gPaleYell} % BG in light/normal mode
%\colorlet{colorDark}{gDarkBlack} % FG in light/normal mode
%\colorlet{colorA}{gAnthracite} % frametitle, standin.out,
%\colorlet{colorAreversed}{gLightTeal} % frametitle, standin.in,
%\colorlet{colorB}{gMidGrey} % gray BG : progress bar, blocks
%\colorlet{colorC}{gDeepYellOr} % progress bar
%\colorlet{colorD}{gLightOrange} % alert
%\colorlet{colorE}{gLightGreen} % example

% Opt-in yellow border for an image placed via a tikz node (e.g.
% \node[gothamImageBorder] {\includegraphics{...}};) -- matches the
% 3pt/accent1-yellow outline PowerPoint draws directly on a picture shape's
% own border (<a:ln w="38100"><a:solidFill><a:schemeClr val="accent1"/>,
% where accent1 turned out to already be this project's own
% c_primary_highlight, not a new color to map). A tikz *style* (a named
% bundle of node keys), not a wrapper command: adding or removing it is
% just adding or deleting one word from that one \node's own option list,
% no separate macro call/argument needed -- add `gothamImageBorder` to draw
% the border, leave it off for a plain image, same as any other tikz key.
%
% Thickness is a single named length, \gothamImageBorderWidth, not a bare
% "3pt" baked directly into the style (2026-07-15, direct instruction: "make
% sure the yellow image borders have consistent thickness throughout the
% project by defining a single variable") -- confirmed directly, by
% measuring the actual PDF content-stream "w" operator, that every
% gothamImageBorder use already rendered at the exact same 3pt regardless
% (a genuinely shared style was already in effect, no numeric drift found),
% so this is a refactor for future adjustability/discoverability, not a fix
% for an actual inconsistency -- any perceived thickness difference between
% slides traces to something else instead (e.g. a border straddling flush
% image content on one slide and not another, which reads as "heavier" even
% at an identical stroke width -- see the gothamImageBorder padding gotcha
% below).
\newlength{\gothamImageBorderWidth}
\setlength{\gothamImageBorderWidth}{2.5pt}
\tikzset{gothamImageBorder/.style={draw=c_primary_highlight, line width=\gothamImageBorderWidth}}

% The "padding gotcha" promised above, finally root-caused (2026-07-23,
% background-task follow-up to the IA2 fine-tuning batch #2 anomaly on
% slides/m2k-Exercise-CAPM-2-current.tex/
% slides/n8e-Tail-Risk-Variance-Clustering-current.tex/
% slides/fm6-Tail-Risk-Market-vs-Simulated-current.tex/
% slides/xqq-Fat-Tails-current.tex, all fixed there with a pragmatic
% `line width=1.2pt` override and an explicitly *unresolved* root cause at
% the time -- "declares an identical stroke width and CTM in the compiled
% PDF's own raw content stream, yet renders at a genuinely different pixel
% thickness", with CTM scale/clips/dash/cap/join all directly checked and
% ruled out). Confirmed via a controlled reproducer (four `gothamImageBorder`
% nodes -- plain opaque image, nested pgfplots axis, plain image at
% `inner sep=3pt`, nested plain tikzpicture at `inner sep=3pt` -- rendered
% at 600dpi and measured by a half-max cross-section profile, not just
% eyeballed): the one variable that actually predicts the rendered width is
% `inner sep`, not content type at all. At `inner sep=0pt` a *fully opaque*
% image measured 1.2pt; the exact same image at `inner sep=3pt` measured
% 2.52pt (matching the declared 2.5pt); a nested pgfplots axis and a plain
% nested tikzpicture, both at `inner sep=3pt`, also measured 2.52pt. Two
% more points confirmed it's a smooth transition, not a binary switch:
% `inner sep=0.6pt` measured 1.8pt (= the 1.2pt floor plus that 0.6pt gap,
% exactly); `inner sep=1.25pt` (= \gothamImageBorderWidth/2, the stroke's
% own half-width) already measured the full 2.52pt.
%
% The mechanism: tikz draws a node's border stroke centered on the node's
% own boundary path, then paints the node's *content* on top of it,
% afterwards. At `inner sep=0`, the content's own bounding box coincides
% exactly with that boundary path, so if the content is opaque all the way
% to its own edge, it paints directly over the stroke's inner half,
% visually leaving only the outer half (`\gothamImageBorderWidth`/2,
% ~1.2pt) -- not a rendering bug, just tikz's normal draw order interacting
% with zero padding. A gap of at least `\gothamImageBorderWidth`/2 (or
% content that's transparent right at its own edge -- see below) leaves the
% border with nothing opaque to cover it, so it renders at its true
% declared width instead.
%
% This also explains why the fix landed on exactly the 4 slides it did,
% each for a different reason: slides/m2k-Exercise-CAPM-2-current.tex
% used `inner sep=3pt` for unrelated layout reasons (room for the axis's
% own tick labels), well past the 1.25pt threshold, so nothing ever covered
% its border. slides/n8e-Tail-Risk-Variance-Clustering-current.tex/
% slides/fm6-Tail-Risk-Market-vs-Simulated-current.tex/
% slides/xqq-Fat-Tails-current.tex all use `inner sep=0pt` -- matching the
% deck's own overwhelming convention -- but their own chart PNGs
% (`chart-index-daily-returns-*.png`/`chart-index-fat-tails-histogram*.png`,
% matplotlib exports with a transparent background) measured alpha=0 at
% every edge pixel (confirmed directly, Pillow), so there was no opaque
% content to cover their own border either. Every other `gothamImageBorder`
% use in this project that was never flagged uses an opaque image (a JPG,
% or a PNG with no transparent margin) at `inner sep=0pt` -- which
% *accidentally* renders at the ~1.2pt covered-half width instead of the
% full 2.5pt this constant nominally declares, and that accidental width is
% what reads as "the deck's normal border thickness" throughout, since it's
% by far the most common case.
%
% Practical upshot for a *new* slide: a `gothamImageBorder` node with
% `inner sep=0pt` wrapping opaque content (the common case) needs no
% special handling -- it will match the rest of the deck automatically. A
% node that instead uses a nonzero `inner sep`, or wraps content with a
% transparent edge (a matplotlib export, an SVG-derived PDF with padding,
% ...), will render at the *full* \gothamImageBorderWidth and look
% noticeably heavier next to the deck's other borders -- either add an
% explicit `line width=\gothamImageBorderWidth/2` (matching the 4 fixed
% slides' own literal 1.2pt) to match the accidental deck-wide look, or
% reduce `inner sep` below \gothamImageBorderWidth/2 so the same coverage
% effect kicks in naturally. The 4 existing per-slide overrides above are
% left as-is -- they're already the correct fix for what's now a fully
% understood cause, not a mystery workaround needing revisiting.

% Reusable "mugshot" -- a circular-cropped photo, bordered like any other
% gothamImageBorder image, with a caption curving along its own bottom arc
% (2026-07-06, direct instruction, generalized from
% slides/4uq-Definition-of-Capital-Market-Efficiency-current.tex's own
% Eugene Fama photo -- see that slide's earlier comment for the source
% PPTX effect this reproduces, `<a:prstGeom prst="ellipse">` +
% `<a:prstTxWarp prst="textArchDown">`, and CONVERSION_NOTES.md's own
% "beamer/TeX gotchas" entry, search "circular photo crop", for the
% clip/border/decorations.text mechanics themselves -- unchanged here,
% just factored out so a second photo doesn't need to repeat them).
%
%   \coordinate (myMugshotCenter) at (...);
%   \gothamMugshot{myMugshotCenter}{media/photo-example.jpeg}{2.6cm}{Some Person, NP 2013}
%
% #1 the *name* of a coordinate the caller has already placed (not a raw
% point expression -- needed as a plain name, not just a point, since the
% caption's own arc path re-references it by name via `++(...)`, which
% only works with a named coordinate/node, not an arbitrary `($...$)`
% calc expression); #2 the image path; #3 the desired final diameter;
% #4 the caption text.
%
% Scales the source photo so it fully covers the circle with no gap
% anywhere -- both dimensions end up at least as big as the diameter --
% while using as much of the source photo as possible (2026-07-07,
% direct instruction: "should not include whitespace... scale the photo
% such that its edges just touch the circle's inner border"). This
% briefly (same day) scaled the *other* way -- by the photo's own
% diagonal, the largest scale at which the *entire* rectangle fits
% inside the circle with nothing clipped at all -- which does avoid
% ever clipping real content, but leaves the circle's own four corners
% uncovered (visible as plain background there, since a non-circular
% photo can only touch an inscribing circle at its corners, not along
% the whole rim) -- exactly the "whitespace" this round's feedback
% ruled out in favor of full coverage. Covering scale is
% `diameter / min(w,h)`: whichever of the photo's own two dimensions is
% *smaller* becomes exactly the diameter (that edge pair touches the
% circle's own inner border, per the direct instruction), and the other,
% larger dimension overhangs the diameter and gets clipped by the
% circular `\clip` -- for a portrait photo (every one used here so far),
% that means width becomes the diameter and a sliver clips off the top
% and bottom, which is the origin of the clipping problem the diagonal
% approach above existed to solve in the first place. Reintroduced
% deliberately: showing as much of the photo as the "no whitespace"
% constraint allows is now the stated priority, over never clipping
% anything.
%
% Measuring the photo's own native size needs `\savebox` (into
% \gothamMugshotBox, declared once below, not a fresh box register per
% call) -- `\includegraphics` with no size option at all renders at the
% image's native resolution, and `\wd`/`\ht` then read that box's actual
% dimensions. pgfmath (already loaded for tikz) converts any TeX
% dimension fed into it to a plain pt-count number automatically, so
% arithmetic on the two measured lengths just works. `min(1,ratio)`
% (ratio = height/width) picks out whichever of the two is smaller
% *relative to the other*, without needing a `sqrt`/`^` at all this time
% -- so this version doesn't hit the "! Dimension too large" pgfmath
% overflow the diagonal version's own arithmetic did (see this file's
% own comment history / CONVERSION_NOTES.md's "beamer/TeX gotchas" for
% that point (easy to misread as a harmless warning, since compilation
% continues and produces *a* number regardless). Reformulating via the
% aspect ratio sidesteps the whole problem: `scale = diameter /
% sqrt(1+ratio^2)` is the same formula algebraically (factor `w` out of
% `sqrt(w^2+h^2)` and it becomes `w*sqrt(1+(h/w)^2)`), but every
% the full story of that overflow, from when this macro's own arithmetic
% did use a `sqrt`).
%
% Caption font reverted to this deck's normal body font (2026-07-07,
% direct instruction: "change the mugshot template back to using regular
% font instead of monospace") -- `\ttfamily` (Fira Mono) was a direct
% instruction on 2026-07-06, one day earlier; \normalfont resets
% explicitly rather than just omitting \ttfamily, so this doesn't
% silently inherit whatever font happens to be active at the call site.
\newsavebox{\gothamMugshotBox}
\newcommand{\gothamMugshot}[4]{%
	\savebox{\gothamMugshotBox}{\includegraphics{#2}}%
	\pgfmathsetmacro{\gothamMugshotImgW}{\wd\gothamMugshotBox}%
	\pgfmathsetmacro{\gothamMugshotImgH}{\ht\gothamMugshotBox}%
	\pgfmathsetmacro{\gothamMugshotRatio}{\gothamMugshotImgH/\gothamMugshotImgW}%
	\pgfmathsetmacro{\gothamMugshotScaledW}{(#3)/min(1,\gothamMugshotRatio)}%
	\begin{scope}
		\clip (#1) circle (0.5*#3);
		\node[anchor=center, inner sep=0pt] at (#1)
			{\includegraphics[width=\gothamMugshotScaledW pt]{#2}};
	\end{scope}
	\draw[gothamImageBorder] (#1) circle (0.5*#3);
	{\normalfont\tiny
	\path[decoration={text along path, text={#4}, text align=center}, decorate]
		(#1) ++(200:0.5*#3+0.25cm) arc (200:340:0.5*#3+0.25cm);
	}
}

% Reusable "operating/cash conversion cycle" bar diagram (FM2_2026s_ps_deck.pptx
% slides 22/24/25/26 all embed the exact same picture -- an EMF, not a native
% shape/table -- so rather than convert that EMF and re-embed it 4 times over,
% it's redrawn once natively here and reused, matching this project's own
% "clean/geometric diagrams get redrawn, not extracted" policy). Three
% stacked bars (Inventory period+Accounts receivable period on top, Accounts
% payable period+Cash conversion cycle in the middle, Operating cycle
% spanning the full width on the bottom) plus 4 tick labels below at
% x=0/(1/3)/(2/3)/1. Geometry/colors measured directly from the source
% deck's own rendered PNG via pixel analysis (PIL run-length color
% detection), not eyeballed: the two internal split points are exactly 1/3
% and 2/3 of the total width (measured 0.3339/0.6672 against a plain
% PowerPoint illustrative diagram, not slide-23's own specific-day example --
% close enough to exact thirds, and a plausible round number for the
% presenter to have actually used, that thirds is treated as the real value
% rather than curve-fitting noise). Fill colors matched to this project's
% *existing* palette rather than the source's own literal RGB, per this
% deck's established emphasis-color policy: source yellow was already an
% exact match for `c_primary_highlight`; source green (RGB 146,208,80) and
% purple (RGB 128,100,162) were close enough to `c_secondary_green`
% and `c_secondary_violet` respectively to reuse rather than define
% two new one-off colors.
%
% Drawn at a fixed "natural" size (10cm wide) with its own top-left corner
% at the calling tikzpicture's origin -- wrap a call in
% \begin{scope}[shift={(x,y)}, scale=s]\gothamCashCycleBars\end{scope} to
% place/resize it, the same convention as this deck's other reusable
% diagram pieces (e.g. slides/bdh-Working-Capital-Example-current.tex's
% circular flow diagram, resized the same way after the fact per user
% feedback) -- resizing later without needing to re-derive any coordinates.
\newcommand{\gothamCashCycleBars}{%
	\fill[c_primary_highlight] (0,0) rectangle (6.667,-0.42);
	\fill[c_secondary_green] (6.667,0) rectangle (10,-0.42);
	\node[align=center, font=\tiny] at (3.333,-0.21) {Inventory period};
	\node[align=center, font=\tiny] at (8.333,-0.21) {Accounts receivable period};
	\fill[c_secondary_green] (0,-0.47) rectangle (3.333,-0.89);
	\fill[c_primary_highlight] (3.333,-0.47) rectangle (10,-0.89);
	\node[align=center, font=\tiny] at (1.667,-0.68) {Accounts payable period};
	\node[align=center, font=\tiny] at (6.667,-0.68) {Cash conversion cycle};
	\fill[c_secondary_violet] (0,-0.94) rectangle (10,-1.36);
	\node[align=center, font=\tiny, text=c_primary_white] at (5,-1.15) {Operating cycle};
	\foreach \gccx in {0,3.333,6.667,10} {
		\draw[c_primary_dark] (\gccx,-1.36) -- (\gccx,-1.46);
	}
	% text width on all 4, narrow enough to force an actual 2-line wrap
	% (not just cap an already-short line): without it, "Raw material
	% purchase" and "Payment for raw materials" (and, symmetrically, the
	% last two) run into each other -- confirmed directly by rendering.
	% A *wide* text width (e.g. 3.1cm, comfortably more than either label's
	% own one-line width) does NOT fix this: \node's text width only forces
	% a wrap once the text exceeds it, so a generous width just lets both
	% neighbors keep their full one-line extent and still collide -- this
	% was tried first and confirmed not to work by rendering, not assumed.
	% 1.5cm is deliberately narrower than any of the 4 would take on one
	% line, guaranteeing the wrap for all of them.
	%
	% All 4 use anchor=north (centered on their own tick, like the middle
	% two always did) -- the outer two originally used anchor=north
	% west/east instead (flush with the tick rather than centered on it),
	% which looked inconsistent with the middle two once seen side by side
	% at actual size. Centering the outer two on their own tick (rather
	% than flush against the diagram's own left/right edge) does push their
	% wrapped text a bit further left/right past the diagram's own
	% bounding box -- checked this stays within the page margin at every
	% scale/position this macro is actually called at in this deck (the
	% narrowest margin case leaves ~0.7cm of clearance), not just assumed
	% safe from the numbers alone.
	\node[anchor=north, font=\tiny, align=center, text width=1.5cm] at (0,-1.5) {Raw material purchase};
	\node[anchor=north, font=\tiny, align=center, text width=1.5cm] at (3.333,-1.5) {Payment for raw materials};
	\node[anchor=north, font=\tiny, align=center, text width=1.5cm] at (6.667,-1.5) {Sale of finished goods};
	\node[anchor=north, font=\tiny, align=center, text width=1.5cm] at (10,-1.5) {Cash collected on sales};
}

% Reusable circular flow diagram (Cash -> Raw materials -> Finished goods ->
% Receivables -> Cash), first built for
% slides/bdh-Working-Capital-Example-current.tex and reused (smaller) on
% slides/qq5-Cash-Conversion-Cycle-current.tex -- factored out here once a
% second slide needed the same picture, same reasoning as
% \gothamCashCycleBars above. Args: #1 diagram center coordinate, #2 radius
% (box-center distance from that center, a bare cm number, no unit suffix),
% #3 the box `scale` factor (a plain node `scale` transform, not a bigger
% font -- see bdh's own comment, preserved below, for why: this diagram was
% resized once already per user feedback, and `scale` was what correctly
% grew the padding/rounded-corners/arrow proportions together rather than
% just the text). Arrow length/width scale with #3 too, computed via
% \pgfmathsetmacro since tikz option values can't do arithmetic inline.
%
% [Preserved from bdh's own resize, 2026-07-06: box size grew by a uniform
% ~1.44x there (compared via the isolated "Cash" box, the one corner with no
% arrow touching it to contaminate the measurement: 0.94cm to 1.36cm wide,
% 0.50cm to 0.74cm tall), and the box-to-box radius grew by the same
% ~1.44-1.45x (1.15cm to ~1.65cm) -- both measured from the user's own
% "Clipboard 1.jpg" reference via pixel analysis, not eyeballed, and
% consistent with a single uniform resize of the whole group.]
% #1 is a *complete* coordinate spec, exactly as it would appear after
% "\coordinate (name) at " -- e.g. "(current page.north east)" or
% "($(current page.north east)+(-3.81cm,-4.61cm)$)", dollar signs and all
% when it's a calc expression. Resolved into a plain \coordinate first
% rather than substituted directly into each box's own ($...$) calc
% expression, since a calc expression can't contain another calc
% expression's $...$ delimiters nested inside it.
\newcommand{\gothamCircularFlowDiagram}[3]{%
	\pgfmathsetmacro{\gothamCFDarrowLen}{5*#3}
	\pgfmathsetmacro{\gothamCFDarrowWidth}{1.3*#3}
	\begin{scope}
		\coordinate (gothamCFDctr) at #1;
		\node[fill=c_primary_highlight, rounded corners=1pt, inner sep=4pt, align=center, font=\footnotesize, scale=#3] (gothamCFDcash) at ($(gothamCFDctr)+(0,#2)$) {Cash};
		\node[fill=c_primary_highlight, rounded corners=1pt, inner sep=4pt, align=center, font=\footnotesize, scale=#3] (gothamCFDraw) at ($(gothamCFDctr)+(#2,0)$) {Raw materials};
		\node[fill=c_primary_highlight, rounded corners=1pt, inner sep=4pt, align=center, font=\footnotesize, scale=#3] (gothamCFDfin) at ($(gothamCFDctr)+(0,-#2)$) {Finished goods};
		\node[fill=c_primary_highlight, rounded corners=1pt, inner sep=4pt, align=center, font=\footnotesize, scale=#3] (gothamCFDrec) at ($(gothamCFDctr)+(-#2,0)$) {Receivables};
		\draw[-{Stealth[length=\gothamCFDarrowLen pt]}, line width=\gothamCFDarrowWidth pt, c_primary_highlight] (gothamCFDcash) to[bend left=18] (gothamCFDraw);
		\draw[-{Stealth[length=\gothamCFDarrowLen pt]}, line width=\gothamCFDarrowWidth pt, c_primary_highlight] (gothamCFDraw) to[bend left=18] (gothamCFDfin);
		\draw[-{Stealth[length=\gothamCFDarrowLen pt]}, line width=\gothamCFDarrowWidth pt, c_primary_highlight] (gothamCFDfin) to[bend left=18] (gothamCFDrec);
		\draw[-{Stealth[length=\gothamCFDarrowLen pt]}, line width=\gothamCFDarrowWidth pt, c_primary_highlight] (gothamCFDrec) to[bend left=18] (gothamCFDcash);
	\end{scope}
}

% A \frac/\dfrac-stacked formula reads as cramped against a plain-text
% neighbor in an itemize -- it's visibly taller than a normal line (both
% above *and* below the line, since the fraction bar sits on the text
% baseline with numerator/denominator stacked around it), so both the
% fraction's own item *and* the item that follows it need a bit of extra
% headroom that a plain-text \item doesn't. Started as a manually-repeated
% \vspace{3pt} before each such \item (first on
% slides/u74-Example-Dynamic-Mattress-Company-Periods-current.tex);
% \tallitem below is that same convention factored into a single drop-in
% \item replacement, so a future fraction-heavy slide (or a change to the
% spacing amount) doesn't mean re-editing several slide files by hand.
% Takes an optional leading multiplier (default 1, i.e. 3pt): \tallitem[1.5]
% for 4.5pt, for a formula tall enough that the default doesn't read as
% enough room. Usage: replace \item with \tallitem both on any itemize line
% containing a stacked fraction *and* on the item immediately following it
% (same nesting level; skip this second replacement if that next item is
% itself already a \tallitem for its own fraction, or if the fraction item
% is the last one in its list) -- it's a plain substitute, not an
% environment, so it doesn't change how the surrounding list itself is
% structured. (Named \gothamFracItem originally; renamed 2026-07-21 to the
% more discoverable \tallitem, and given the optional multiplier and the
% following-item rule above at the same time -- no behavior change for any
% existing default-strength call site.)
% Ratio, not a frozen length (2026-07-22, direct instruction: "modify
% \tallitem and \tallitemnobullet to also use scaling relative to font
% size") -- the original \newlength+\setlength{...}{3pt} pair freezes to
% one fixed pt value the moment \setlength runs, the same issue (and same
% fix) as \gothamItemsepRatioI/II/III/IV further below: converted to a
% bare ratio macro, with the "em" unit written back on inside \tallitem's
% and \tallitemnobullet's own bodies (not here), so it's recomputed fresh
% against whatever font is actually live at each call site -- correct
% regardless of nesting level or a \footnotesize/\scriptsize override,
% unlike the old fixed 3pt (calibrated against one specific level's font,
% silently wrong at any other). 0.3 reproduces the original 3pt exactly
% at level 2's own font (10pt) -- the level \tallitem was first used at,
% on slides/u74-Example-Dynamic-Mattress-Company-Periods-current.tex (see
% this macro's own comment above) -- so no visual change there, while now
% correctly scaling up/down at every other level and font override.
\newcommand{\gothamTallItemBaseRatio}{0.3}
\newlength{\gothamTallItemBaseUnit}
\newlength{\gothamTallItemSpace}
\newcommand{\tallitem}[1][1]{%
	\setlength{\gothamTallItemBaseUnit}{\gothamTallItemBaseRatio em}%
	\setlength{\gothamTallItemSpace}{#1\gothamTallItemBaseUnit}%
	\vspace{\gothamTallItemSpace}\item}

% For a line that needs \tallitem's own extra headroom but no bullet at
% all -- a display equation continuing the previous \item's own sentence,
% rather than a new bulleted point of its own (see README.md's own
% \tallitem section, which already documents this exact combination as a
% rare, previously hand-written case: "\vspace{\gothamTallItemBaseUnit}\item[]").
% \tallitemnobullet (2026-07-21, direct instruction) factors that hand-
% written idiom into the same kind of drop-in replacement \tallitem
% already is for its own case, same optional multiplier: \tallitemnobullet[1.5].
\newcommand{\tallitemnobullet}[1][1]{%
	% Recomputes \gothamTallItemBaseUnit itself, fresh, rather than reusing
	% whatever \tallitem's own body last left it at -- since it's now a
	% per-call, live "em"-based value (see \tallitem's own comment above),
	% not the one-time preamble-load-time constant it used to be, relying
	% on call order/a preceding \tallitem to have set it correctly would
	% silently break the moment \tallitemnobullet is used on its own, or
	% at a different font size than whichever \tallitem last ran.
	\setlength{\gothamTallItemBaseUnit}{\gothamTallItemBaseRatio em}%
	\setlength{\gothamTallItemSpace}{#1\gothamTallItemBaseUnit}%
	\vspace{\gothamTallItemSpace}\item[]}

% flalign* (used to left-align a display equation inside an itemize --
% see CONVERSION_NOTES.md's "flalign* for left-aligning..." entry) inherits
% plain LaTeX's \abovedisplayskip/\belowdisplayskip, sized for a full
% document page, not a compact slide itemize -- reads as far too much blank
% space above/below the equation next to the surrounding item text (found
% 2026-07-21, direct instruction; most visible on
% slides/bwi-Pearson-Correlation-Definition-current.tex's own single-line
% equation, but the same oversized gap is present everywhere flalign* is
% used this way, e.g. slides/7se-Portfolio-Risk-Three-Assets-current.tex
% and slides/rb3-Asset-Allocation-Example-CAL-current.tex). \gothamTightDisplaySkips
% tightens both skips -- and their "short" variants, since amsmath picks
% between the two depending on how much text precedes the display, so
% setting only one pair left the gap unfixed whenever the preceding line
% happened to be short -- down to something that reads as part of the
% same list rather than a separate floating display. Usage: wrap the
% *real*, unmodified flalign* in a plain TeX group together with this
% command, `{\gothamTightDisplaySkips\begin{flalign*} ... \end{flalign*}}`,
% rather than a new same-named environment -- tried that first
% (`\newenvironment{gothamtightflalign*}{...\begin{flalign*}}{\end{flalign*}}`)
% and it broke LaTeX's own environment-nesting tracking (`\@currenvir`),
% surfacing as "\begin{flalign*} ... ended by \end{gothamtightflalign*}"
% and a fatal, PDF-less compile on the very first slide that used it. A
% plain group sidesteps that class of bug entirely since it never
% introduces a same-named wrapper environment.
\newlength{\gothamTightDisplaySkip}
\setlength{\gothamTightDisplaySkip}{3pt}
\newcommand{\gothamTightDisplaySkips}{%
	\abovedisplayskip=\gothamTightDisplaySkip
	\abovedisplayshortskip=\gothamTightDisplaySkip
	\belowdisplayskip=\gothamTightDisplaySkip
	\belowdisplayshortskip=\gothamTightDisplaySkip}

% Set logo
% \gothamInstituteLogoPath (a plain \newcommand storing an image path,
% same idiom as \gothamChapterLogoPath below) and \gothamLogoHeight (a
% length) are both defined by the per-university preamble that \input{}s
% this file, before this file is reached -- see e.g.
% Slidemaster_preamble_uni-graz.tex. Whichever source image
% \gothamInstituteLogoPath points at should have any symmetric margin of
% empty transparent pixels around the actual logo artwork trimmed off first
% (see Slidemaster_preamble_uni-graz.tex's own comment on why: an untrimmed
% margin makes \includegraphics's box bigger than the visible logo, so the
% logo's right-hand margin ends up looking bigger than its top/bottom
% margins, since only the latter get absorbed into vertical centering).
% \gothamLogoHeight is the single source of truth for the logo's rendered
% size: it drives the (invisible) in-flow logo below, the white patch's
% width further down, and the visible redrawn logo, so those three can never
% drift out of sync with each other again (that drift -- redrawing a smaller
% logo than the one the patch/position were sized for -- is what caused the
% left/right margins around the logo to come out uneven).

% The tikz node wrapper (remember picture, but not overlay, so it still
% behaves as normal inline content) records this logo's exact page position
% under the name "gothamLogoAnchor", so it can be redrawn on top of a white
% patch further down (see the "yellow box shorter than the page" block below).
\renewcommand{\gothamInstituteLogoSquare}[1][\gothamLogoHeight]{%
	\tikz[remember picture, baseline=(gothamLogoAnchor.base)]{%
		\node[inner sep=0pt, outer sep=0pt](gothamLogoAnchor){%
			\includegraphics[height=#1]{\gothamInstituteLogoPath}%
		};%
	}%
}

% \gothamFrametitleRightpading (left/right padding around the logo, chosen
% to match its top/bottom margins -- measured against the title colorbox's
% own top and bottom edges -- rather than the theme's 2.2ex default) is
% also set by the per-university preamble, not here: the top/bottom margins
% are a side effect of vertically centering a \gothamLogoHeight-tall logo in
% the title bar, not something set directly, so this value has to be
% re-measured and set per institution whenever its own logo/height differs
% from another's.

% Logo width, used below to size the white patch behind it.
\newlength{\gothamLogoWidth}
\settowidth{\gothamLogoWidth}{\includegraphics[height=\gothamLogoHeight]{\gothamInstituteLogoPath}}

% Let the yellow title bar visually stop short of the page's right edge, so
% the logo sits on white instead of on yellow: paint a white patch spanning
% the full height of the title bar behind the logo, then redraw the logo
% (already positioned correctly by \gothamFrametitleRightpading above, via
% the "gothamLogoAnchor" node captured above) on top of it. This needs two
% compilation passes: like any \tikz "remember picture", the node positions
% are only available in the .aux file written by the previous run.
% (The redraw uses \includegraphics directly rather than
% \usebeamertemplate{instituteLogo}/{includeLogo}: those templates rely on
% gotham's internal property-list lookup, which is only valid in the
% context of the original frametitle rendering and silently produces
% nothing when re-invoked here.)
% Additionally paints white across the *entire* width of the strip between
% the logo's bottom edge and the title bar's actual bottom edge, so the
% yellow ends flush with the logo's bottom (rather than extending further
% down under the title text too) while the box's width and top position
% stay unchanged. Checked against measured PDF coordinates that the title
% text's own (font-metric, not just this-string's-glyphs) lowest descender
% still clears the logo's bottom edge by about 1.3pt, so this doesn't clip
% text even on titles with descenders (e.g. "Highlighting...").
%
% Factored into its own macro, rather than written directly into the
% \addtobeamertemplate call below, so \gotham_no_inline_logo (a whole
% separate frametitle template variant, further down) can reuse it too --
% \addtobeamertemplate patches whatever template is *currently* installed
% at the time it's called (gotham's own default one here), so a template
% defined fresh afterwards doesn't inherit that patch automatically and
% needs to invoke this explicitly instead.
% Every edge-anchored node in this hook (and in every other page-corner
% overlay in this file) needs an explicit "outer sep=0pt" alongside its
% "inner sep=0pt": tikz's default outer sep is 0.5\pgflinewidth (0.2pt at
% the default 0.4pt line width), and a node's edge/corner anchors sit on
% the *outer* boundary -- outer sep past the filled rectangle -- so an
% anchor placed exactly on a page corner leaves the fill itself 0.2pt
% short of that corner on the anchored sides. Found 2026-07-16 (direct
% instruction: a hairline white line above/left of titleframe's and
% quoteframe's yellow boxes, their right edges 0.4pt off the standard
% title bar's own visible yellow edge, and their logos 0.2pt left of
% every other frame's) by measuring the compiled PDF's drawn rectangles:
% every affected fill/image was inset by exactly 0.199bp = 0.2pt from
% its anchor. This is also the same offset the footer box's own comment
% (further down) had earlier measured, worked around with a hand-nudged
% anchor, and left as "root cause not chased further" -- now chased:
% it *is* a general property of every edge/corner-anchored tikz node,
% not something specific to that one construction. (Center anchors are
% unaffected -- outer sep is symmetric around the center -- so the
% gothamLogoAnchor.center redraw below doesn't need it.)
\newcommand{\gothamFrametitleHookCode}{%
	\tikz[remember picture, overlay]\node(gothamTitleBottom){};%
	\begin{tikzpicture}[remember picture, overlay]
		\path let \p1=(current page.north east), \p2=(gothamTitleBottom) in
			node[fill=white, anchor=north east, inner sep=0pt, outer sep=0pt, name=gothamMainPatch,
			     minimum width=\gothamLogoWidth+2\gothamFrametitleRightpading,
			     minimum height=\y1-\y2] at (\x1,\y1) {};
		\ifx\gothamFrametitleFullHeightYellow\empty
			\path let \p1=(current page.north east), \p2=(gothamTitleBottom),
			          \p3=(gothamLogoAnchor.south), \p4=(current page.west) in
				node[fill=white, anchor=north west, inner sep=0pt, outer sep=0pt,
				     minimum width=\x1-\x4, minimum height=\y3-\y2] at (\x4,\y3) {};
		\fi
		% Normally redrawn at gothamLogoAnchor's own position (wherever the
		% invisible in-flow logo -- on the same line as the title text --
		% ended up). \gothamFrametitleLogoFixed instead redraws it at a fixed
		% page-corner position, independent of where the title text sits --
		% for a frame type (sectionframe) whose title is deliberately *not*
		% at the position a normal frame's title/logo line would put it.
		% gothamMainPatch already covers the full logo column from the page
		% top down to the box's true bottom edge (above), so it's already
		% blank there regardless of which position is used.
		\ifx\gothamFrametitleLogoFixed\empty
			\node[anchor=center, inner sep=0pt] at (gothamLogoAnchor.center) {\includegraphics[height=\gothamLogoHeight]{\gothamInstituteLogoPath}};
		\else
			% inner ysep is derived, not \gothamFrametitleRightpading again:
			% the inline logo this fixed redraw must line up with is
			% *vertically centered* in the title bar, and its resulting top
			% margin (8.14pt, measured) is not exactly the horizontal
			% padding (0.28cm = 7.97pt) -- close enough that the pre-outer-sep-fix
			% code's accidental extra 0.2pt masked the difference, but a
			% visible micro-jump when flipping pages once the fix landed.
			% Rather than a fresh calibrated constant:
			% \gothamTitleBarBottomAlign (the visible bar bottom) *is* the
			% inline logo's bottom edge, so subtracting \gothamLogoHeight
			% reproduces its top margin exactly, and re-measuring one
			% constant keeps both in sync if the bar ever changes.
			\node[anchor=north east, inner xsep=\gothamFrametitleRightpading,
			      inner ysep=\dimexpr\gothamTitleBarBottomAlign-\gothamLogoHeight\relax,
			      outer sep=0pt] at (current page.north east) {\includegraphics[height=\gothamLogoHeight]{\gothamInstituteLogoPath}};
		\fi
		% Optional per-frame "chapter logo" (see \chapterlogo{} below): drawn
		% immediately to the left of the main logo's white patch, on its own
		% white patch, only when \gothamChapterLogoPath is non-empty. Sized
		% with exactly one \gothamFrametitleRightpading of its own padding, on
		% its outer (left) side only -- the inner (right) side is left flush
		% against gothamMainPatch, since that patch already contributes its
		% own \gothamFrametitleRightpading of padding around the main logo.
		% This makes both gaps -- chapter logo to main logo, and chapter logo
		% to where the yellow bar resumes -- come out equal, at exactly one
		% \gothamFrametitleRightpading each, without double-counting or
		% needing a second length constant.
		\ifx\gothamChapterLogoPath\empty\else
			\settowidth{\gothamChapterLogoWidth}{\includegraphics[height=\gothamLogoHeight]{\gothamChapterLogoPath}}%
			% The patch's right edge is nudged 0.5pt *past* gothamMainPatch's own
			% west edge (i.e. slightly overlapping it) rather than exactly
			% touching it -- both are solid white, so the overlap itself is
			% invisible, but touching them exactly left a hairline yellow seam
			% at the join (a real, if sub-millimetre, gap that persisted across
			% recompiles, most likely PDF fill-antialiasing at the shared edge
			% of two separately-filled abutting shapes). The logo image itself
			% (below) is *not* shifted, so the actual margin it defines is
			% unaffected.
			\path let \p1=(current page.north east), \p2=(gothamTitleBottom) in
				node[fill=white, anchor=east, inner sep=0pt, outer sep=0pt,
				     minimum width=\gothamChapterLogoWidth+\gothamFrametitleRightpading+0.5pt,
				     minimum height=\y1-\y2] at ($(gothamMainPatch.west)+(0.5pt,0)$) {};
			\node[anchor=east, inner sep=0pt, outer sep=0pt] at (gothamMainPatch.west |- gothamLogoAnchor.center)
				{\includegraphics[height=\gothamLogoHeight]{\gothamChapterLogoPath}};
		\fi
	\end{tikzpicture}%
}
\addtobeamertemplate{frametitle}{}{\gothamFrametitleHookCode}

% \chapterlogo{path} opts a single upcoming standardframe into an extra logo
% (e.g. a topic/chapter icon) next to the main institute logo -- see the
% frametitle hook above for how it's drawn. Call it *before* \begin{standardframe},
% since the frame title (and this hook) is typeset as part of \begin{frame}
% itself, before any frame-body content would run. standardframe resets this
% back to empty on \end{frame} (below), so it only ever applies to the one
% frame that follows the \chapterlogo{} call, never leaking into the next slide.
% Both the set and the reset use \global\def (checked empirically, via
% \typeout, against the frametitle hook's actual view of this macro on a
% following frame): beamer's frame environment opens its own TeX group around
% the frame body, so a plain (local) assignment done in standardframe's
% closing code -- which runs right after \end{frame} -- was silently undone
% once that group actually closed, reviving the *previous* value instead of
% the reset. \global\renewcommand alone wasn't enough to fix it (still leaked
% into the next frame) -- only dropping to the lower-level \global\def did.
\newlength{\gothamChapterLogoWidth}
\newcommand{\gothamChapterLogoPath}{}
\newcommand{\chapterlogo}[1]{\global\def\gothamChapterLogoPath{#1}}

% \gothamFrametitleFullHeightYellow: opt-out for the frametitle hook's
% "trim yellow to the logo's bottom edge" white patch (see the hook above) --
% empty (default) keeps that trim; non-empty skips it, letting the yellow
% show all the way down to the title box's own true bottom edge instead.
% That trim exists to erase a normally-tiny (~1-2pt) sliver of yellow that
% would otherwise stick out below the logo on an ordinary short title bar --
% harmless there, but sectionframe's deliberately tall bar turns that
% "tiny sliver" into most of the box, silently wiping out nearly all of the
% intended yellow with white. sectionframe sets this (globally, same
% set/reset reasoning as \gothamChapterLogoPath above); standardframe resets
% it back to empty on \end{frame}, so it's opt-in per frame like the others.
\newcommand{\gothamFrametitleFullHeightYellow}{}

% \gothamFrametitleLogoFixed: opt-in for the frametitle hook's logo redraw
% (see the hook above) to use a fixed page-corner position instead of
% wherever the title text's own line puts it. Normally the two coincide (the
% logo sits inline with the title, gotham's own design), which is exactly
% why this is off by default -- but sectionframe's title is deliberately
% *not* where a normal frame's title sits (bottom-anchored within a much
% taller bar), and letting the logo follow it down there would put it
% somewhere no other slide in the deck ever has it. Same set/reset reasoning
% and mechanism as \gothamChapterLogoPath/\gothamFrametitleFullHeightYellow
% above.
\newcommand{\gothamFrametitleLogoFixed}{}

% \gothamFooterBoxColor: the footer box's own fill color (see the footline
% hook further down), overridable per frame the same set-before/reset-after
% way as the flags above -- default \gothamFooterBoxColorDefault (set by the
% per-university preamble before \input{}ing this file, same "value
% supplied externally, referenced lazily" idiom as \gothamBulletColor/
% \gothamInstituteLogoPath), matching every frame's title bar by default;
% highlightframe overrides it to match its own (likewise overridable) title
% bar color, \gothamHighlightColor, so both boxes change together with a
% single retargeting.
\newcommand{\gothamFooterBoxColor}{\gothamFooterBoxColorDefault}

% A frametitle template variant that drops the inline institute logo
% entirely -- for use together with \gothamFrametitleLogoFixed (which
% redraws the logo separately, at a fixed position, making the inline one
% pointless anyway). gotham's own template (beamerouterthemegotham.sty,
% "default/gotham~subsameline") puts \nolinebreak\hfill{...logo...} directly
% after \insertframetitle; this variant drops that whole part.
%
% Doing so, though, means the logo no longer reserves *any* horizontal room
% on the title's own line the way it used to -- so rightskip here widens
% \gothamFrametitleRightpading to reserve the same width gothamMainPatch
% (the frametitle hook's white patch, drawn afterward) already reserves for
% the logo column: \gothamLogoWidth + 2*\gothamFrametitleRightpading. Without
% this, the title paragraph wraps as if that whole column were free text
% space -- not an actual TeX overfull-hbox condition (no warning, the text
% genuinely fits the *box's own* width), just a title that runs visually
% underneath the white patch and logo drawn on top of it afterward, which
% looks identical to clipped/hidden text. (An earlier, wrong diagnosis of
% this same visual symptom blamed gotham's \nolinebreak forcing an overfull
% line -- checked by testing gotham's original template installed with a
% long \huge title and no changes at all, which reproduces something that
% *looks* the same; but the actual fix needed here turned out to be this
% rightskip, not removing \nolinebreak, which was really addressing a red
% herring from an isolated test that happened not to include a competing
% logo column at all.)
%
% "plus 1fil" on that rightskip keeps every line ragged-right (its own
% natural width, however short of the margin that leaves) rather than fully
% justified -- without it, a wrapped multi-line title's non-last lines get
% stretched to exactly fill the available width, which reads as oddly wide,
% uneven word-gaps on a short title (most noticeable exactly when a title is
% long enough to wrap, i.e. always, for this template).
\makeatletter
\defbeamertemplate{frametitle}{gotham_no_inline_logo}{%
	\nointerlineskip
	\begin{beamercolorbox}[wd=\paperwidth, sep=0pt,
	                        leftskip=\gothamFrametitleLeftpading,
	                        rightskip=\dimexpr\gothamLogoWidth+2\gothamFrametitleRightpading\relax plus 1fil]{frametitle}
		\gotham@frametitlestrut@start
		\insertframetitle
		\ifx\insertframesubtitle\@empty\else
			{\usebeamerfont{framesubtitle}\usebeamercolor[fg]{framesubtitle}\gothamFrameSubtitleSep \insertframesubtitle}%
			\vspace*{\gothamFramesubtitleStrutend}%
		\fi
		\gotham@frametitlestrut@end
	\end{beamercolorbox}
	\gothamFrametitleHookCode
}
\makeatother


% Define slideID as a simple value holder (initially empty)
\newcommand{\slideID}{}%

% Define slideID as a simple value holder (not a command with arguments)
\makeatletter
\newcommand{\setSlideID}[1]{%
	% Grab everything up to the first '-' and store it in \slideID
	\edef\tempa{#1}%
	\expandafter\setSlideID@aux\tempa-\relax
}
\def\setSlideID@aux#1-#2\relax{%
	\renewcommand{\slideID}{#1}%
}
\makeatother

% Define a new \slide command to automatically pick the current version of the slide
% 2026-07-08 (direct instruction): slide filenames dropped their leading
% "s-" (e.g. slides/ab3-Some-Slide-current.tex, not the old
% slides/s-ab3-Some-Slide-current.tex) so slide IDs are directly
% searchable in Explorer -- this macro (and the other three below) is the
% *only* place that ever needed to change for that: every deck's own
% \slide{}/\slideversion{}/\omitslide{}/\omitslideversion{} calls pass just
% "ID-slug", never the filename itself, so renaming every file on disk to
% match needed no changes anywhere else (not the deck files, not
% slide_ids.csv, which only ever stored bare IDs to begin with).
\newcommand{\slide}[1]{%
	\setSlideID{#1}%
	\input{slides/#1-current.tex}%
}

% \slidefinal{ID-slug} -- like \slide{}, but renders the slide at its final,
% fully-revealed overlay state in every deck version, instructor included
% (2026-07-12, direct instruction: show jww's finished build in place of a
% separate recap slide, "irrespective of which deck version is being
% compiled" -- that recap slide has since been deleted entirely, its own ID
% no longer in use). A slide
% built with \conditionalpause[instructor] already renders fully built in
% the during-course/after-course compiles -- that pause is a no-op for
% them, so nothing there ever waits on a reveal step -- the only case
% actually needing help is the instructor deck, which would otherwise still
% step through the build. Locally shadowing \conditionalpause to a no-op
% (same optional-argument signature as the real one, so it still accepts
% e.g. \conditionalpause[instructor] without erroring) for the duration of
% one \slide{} call makes every \conditionalpause in that one inclusion
% inert, regardless of the actual compiled \gothamDeckVersion -- scoped
% inside \begingroup/\endgroup, so it reverts immediately after, and has no
% effect on any *other* \slide{}/\slidefinal{} call to the same ID
% elsewhere in the same deck (e.g. an earlier, normally-building instance of
% the same slide). \deckonly/\deckuncover/\deckvisible are untouched, so
% deck-version-tagged content (as opposed to reveal-step-tagged content)
% still follows its usual per-version rules.
\newcommand{\slidefinal}[1]{%
	\begingroup
		\renewcommand{\conditionalpause}[1][instructor]{}%
		\slide{#1}%
	\endgroup
}

% Load a specific, pinned version of a slide instead of whichever is "current"
% Usage: \slideversion{ID-slug}{M-m-p} loads slides/ID-slug-vM-m-p.tex
\newcommand{\slideversion}[2]{%
	\setSlideID{#1}%
	\input{slides/#1-v#2.tex}%
}

% Student-version decks: \studentversion (called anywhere in a deck, e.g.
% right after % Slidemaster preamble variant for University of Graz decks. Institution-
% specific content only (colors, logo, logo calibration) -- every frame
% type, the deck-version system, footer/hyperlink/footnote styling, tikz
% styles, etc. live in Slidemaster_preamble_base.tex, \input{}ed at the
% bottom of this file, and are shared verbatim across every institution
% this project supports (see Slidemaster_preamble_mci.tex and
% Slidemaster_preamble_uni-ibk.tex for the other two). A deck picks its
% institution by \input{}ing this file instead of one of those, nothing
% else changes.

% Color definition (Graz palette, sourced from the original single-preamble
% version of this project -- these are the University of Graz's own official
% colors, not derived/approximated).
\definecolor{c_primary_dark}{HTML}{000000}
\definecolor{c_primary_highlight}{HTML}{ffd500}
\definecolor{c_primary_gray}{HTML}{c6c6c6}
\definecolor{c_primary_white}{HTML}{FFFFFF}
\definecolor{c_secondary_violet}{HTML}{948dc3}
\definecolor{c_tertiary_violet}{HTML}{c2bddf}
\definecolor{c_secondary_blue}{HTML}{689cce}
\definecolor{c_tertiary_blue}{HTML}{bacce8}
\definecolor{c_secondary_green}{HTML}{9fcb81}
\definecolor{c_tertiary_green}{HTML}{cce2ba}
\definecolor{c_secondary_red}{HTML}{e4684b}
\definecolor{c_tertiary_red}{HTML}{f1ae95}
\definecolor{c_secondary_orange}{HTML}{f6a968}
\definecolor{c_tertiary_orange}{HTML}{fbd0aa}
\definecolor{c_secondary_yellow}{HTML}{ffe063}
\definecolor{c_tertiary_yellow}{HTML}{ffeeaf}
\definecolor{c_tertiary_gray}{HTML}{e3e3e3}
% Not an official Uni Graz color (hence "custom", not "secondary" like the
% palette above) -- a softer near-black than c_primary_dark, used for
% quoteframe's full-bleed dark block, matching the specific dark gray (not
% pure black) ExampleSlides.pptx used for that same block. Also reused for
% hyperlink text color, for the same reason it works for quoteframe's
% attribution text: dark enough for good contrast/legibility on a light
% background.
\definecolor{c_custom_gray}{HTML}{333333}
\definecolor{c_custom_red}{HTML}{C00000}
\definecolor{c_custom_green}{HTML}{00B050}

% Logo + calibration.
% schemes/scheme_uni_graz/Logo_Uni-Graz_4c-cropped.pdf -- official vector artwork,
% user-supplied (schemes/scheme_uni_graz/Logo_Uni-Graz_4c.eps, kept alongside as the
% untouched source), converted to PDF via epstopdf and cropped to its true
% ink bounding box (computed with Ghostscript's own `-sDEVICE=bbox`, same
% treatment as MCI's and Innsbruck's own logos -- see
% Slidemaster_preamble_mci.tex's own comment for the fuller explanation of
% why the EPS's own %%BoundingBox isn't trusted directly). Vector, not
% raster, so it stays crisp at any zoom. Supersedes this file's own earlier
% raster logo (schemes/scheme_uni_graz/Logo Uni Graz 4c-cropped.png, itself cropped
% from schemes/scheme_uni_graz/Logo Uni Graz 4c.png's own symmetric transparent
% margin -- both kept alongside as historical/fallback assets, not deleted).
\newcommand{\gothamInstituteLogoPath}{schemes/scheme_uni_graz/Logo_Uni-Graz_4c-cropped.pdf}
\newlength{\gothamLogoHeight}
\setlength{\gothamLogoHeight}{3ex}
% Re-measured against the compiled PDF for this specific logo image -- see
% Slidemaster_preamble_base.tex's own comment on \gothamFrametitleRightpading
% for why this needs re-measuring per institution rather than reusing
% another one's value.
\setlength{\gothamFrametitleRightpading}{0.28cm}

% Itemize bullet color (2026-07-16, direct instruction): Graz keeps the
% project's original behavior, bullets in the primary highlight color.
\newcommand{\gothamBulletColor}{c_primary_highlight}

% Title bar and footer box/text colors (2026-07-16, direct instruction):
% Graz keeps the project's original behavior throughout -- title/footer
% boxes in the primary highlight color, title/footer text in the primary
% dark color (unlike MCI/Innsbruck, which switch these to their own
% secondary blue/white -- see Slidemaster_preamble_mci.tex's own comment).
\newcommand{\gothamTitleBoxColor}{c_primary_highlight}
\newcommand{\gothamTitleTextColor}{c_primary_dark}
\newcommand{\gothamFooterBoxColorDefault}{c_primary_highlight}
\newcommand{\gothamFooterTextColor}{c_primary_dark}

\input{Slidemaster_preamble_base}
, similar to \sidefootnotes or
% \setfootnotesymbols) switches that compile into the "handout" variant --
% including, as of 2026-07-03, actually forcing beamer's own "handout" mode
% (\gdef\beamer@currentmode{handout}, normally only reachable via the
% \documentclass[handout]{beamer} *class option*, set once at class-loading
% time -- too early for a command like \studentversion, called later from a
% deck's own preamble, to use that route). This was checked, not assumed:
% until this point \studentversion only ever touched \omitslide{}, so
% nothing suppressed \pause/overlay specs -- true only by coincidence, since
% no slide converted so far happens to use them. Late-redefining
% \beamer@currentmode was verified (isolated minimal-beamer tests, plain
% \pause *and* inline specs like \item<2->) to reproduce the same collapse
% -to-one-page, show-the-final-revealed-state behavior as the class option,
% even set well after \documentclass{beamer} -- beamer's overlay decoder
% (beamerbasedecode.sty) reads \beamer@currentmode at frame-processing time,
% not just once at class-option time, so this is safe, not a hack that
% happens to work by luck of file ordering.
% \omitslide{}/\omitslideversion{} are drop-in replacements for
% \slide{}/\slideversion{} that tag a slide as excluded from that variant:
%
%   \omitslide{xyz-Answer-key}                 % instead of \slide{...}
%   \omitslideversion{xyz-Answer-key}{1-0-0}    % instead of \slideversion{...}{...}
%
% In a \studentversion compile, a tagged slide's \input{} is skipped
% entirely -- no frame is produced for it at all (not a blank page, not a
% numbering gap: beamer's frame/slide counters simply never advance for a
% call that never ran). In every other (i.e. normal/instructor) compile, a
% tagged slide loads exactly as it would via \slide{}/\slideversion{}, with
% one difference: \gothamOmittedSlide is set for the duration of that one
% \input{}, so the footer template below can print a marker next to that
% frame's slide ID -- a deliberately visible reminder, while reviewing the
% instructor deck, of which slides won't make it to students. The same tag
% works for both compiles because it lives at the call site in the *deck*
% file (not inside the slide's own content file): the identical slide list
% can serve both the instructor deck and the student handout, e.g.
%
%   % Slidemaster preamble variant for University of Graz decks. Institution-
% specific content only (colors, logo, logo calibration) -- every frame
% type, the deck-version system, footer/hyperlink/footnote styling, tikz
% styles, etc. live in Slidemaster_preamble_base.tex, \input{}ed at the
% bottom of this file, and are shared verbatim across every institution
% this project supports (see Slidemaster_preamble_mci.tex and
% Slidemaster_preamble_uni-ibk.tex for the other two). A deck picks its
% institution by \input{}ing this file instead of one of those, nothing
% else changes.

% Color definition (Graz palette, sourced from the original single-preamble
% version of this project -- these are the University of Graz's own official
% colors, not derived/approximated).
\definecolor{c_primary_dark}{HTML}{000000}
\definecolor{c_primary_highlight}{HTML}{ffd500}
\definecolor{c_primary_gray}{HTML}{c6c6c6}
\definecolor{c_primary_white}{HTML}{FFFFFF}
\definecolor{c_secondary_violet}{HTML}{948dc3}
\definecolor{c_tertiary_violet}{HTML}{c2bddf}
\definecolor{c_secondary_blue}{HTML}{689cce}
\definecolor{c_tertiary_blue}{HTML}{bacce8}
\definecolor{c_secondary_green}{HTML}{9fcb81}
\definecolor{c_tertiary_green}{HTML}{cce2ba}
\definecolor{c_secondary_red}{HTML}{e4684b}
\definecolor{c_tertiary_red}{HTML}{f1ae95}
\definecolor{c_secondary_orange}{HTML}{f6a968}
\definecolor{c_tertiary_orange}{HTML}{fbd0aa}
\definecolor{c_secondary_yellow}{HTML}{ffe063}
\definecolor{c_tertiary_yellow}{HTML}{ffeeaf}
\definecolor{c_tertiary_gray}{HTML}{e3e3e3}
% Not an official Uni Graz color (hence "custom", not "secondary" like the
% palette above) -- a softer near-black than c_primary_dark, used for
% quoteframe's full-bleed dark block, matching the specific dark gray (not
% pure black) ExampleSlides.pptx used for that same block. Also reused for
% hyperlink text color, for the same reason it works for quoteframe's
% attribution text: dark enough for good contrast/legibility on a light
% background.
\definecolor{c_custom_gray}{HTML}{333333}
\definecolor{c_custom_red}{HTML}{C00000}
\definecolor{c_custom_green}{HTML}{00B050}

% Logo + calibration.
% schemes/scheme_uni_graz/Logo_Uni-Graz_4c-cropped.pdf -- official vector artwork,
% user-supplied (schemes/scheme_uni_graz/Logo_Uni-Graz_4c.eps, kept alongside as the
% untouched source), converted to PDF via epstopdf and cropped to its true
% ink bounding box (computed with Ghostscript's own `-sDEVICE=bbox`, same
% treatment as MCI's and Innsbruck's own logos -- see
% Slidemaster_preamble_mci.tex's own comment for the fuller explanation of
% why the EPS's own %%BoundingBox isn't trusted directly). Vector, not
% raster, so it stays crisp at any zoom. Supersedes this file's own earlier
% raster logo (schemes/scheme_uni_graz/Logo Uni Graz 4c-cropped.png, itself cropped
% from schemes/scheme_uni_graz/Logo Uni Graz 4c.png's own symmetric transparent
% margin -- both kept alongside as historical/fallback assets, not deleted).
\newcommand{\gothamInstituteLogoPath}{schemes/scheme_uni_graz/Logo_Uni-Graz_4c-cropped.pdf}
\newlength{\gothamLogoHeight}
\setlength{\gothamLogoHeight}{3ex}
% Re-measured against the compiled PDF for this specific logo image -- see
% Slidemaster_preamble_base.tex's own comment on \gothamFrametitleRightpading
% for why this needs re-measuring per institution rather than reusing
% another one's value.
\setlength{\gothamFrametitleRightpading}{0.28cm}

% Itemize bullet color (2026-07-16, direct instruction): Graz keeps the
% project's original behavior, bullets in the primary highlight color.
\newcommand{\gothamBulletColor}{c_primary_highlight}

% Title bar and footer box/text colors (2026-07-16, direct instruction):
% Graz keeps the project's original behavior throughout -- title/footer
% boxes in the primary highlight color, title/footer text in the primary
% dark color (unlike MCI/Innsbruck, which switch these to their own
% secondary blue/white -- see Slidemaster_preamble_mci.tex's own comment).
\newcommand{\gothamTitleBoxColor}{c_primary_highlight}
\newcommand{\gothamTitleTextColor}{c_primary_dark}
\newcommand{\gothamFooterBoxColorDefault}{c_primary_highlight}
\newcommand{\gothamFooterTextColor}{c_primary_dark}

\input{Slidemaster_preamble_base}

%   \FooterText{...}
%   \makeatletter\ifdefined\HANDOUT\studentversion\fi\makeatother  % or similar
%   \input{course-content.tex}   % the shared \slide{}/\omitslide{} list
%
% with two thin wrapper decks (one plain, one calling \studentversion first)
% both \input-ing the same shared content file -- see README.md.
\newif\ifgothamStudentVersion
% \makeatletter must wrap the \newcommand itself, not sit *inside* the macro
% body: catcodes apply at tokenization time, which for a macro body happens
% once, immediately, when \newcommand reads it -- not deferred until the
% macro is later called. \makeatletter placed inside the body (tried first)
% was too late: \beamer@currentmode had already been mis-tokenized as
% "\beamer" + the six-way-otherwise-ordinary character "@" + eleven plain
% letters "currentmode", by the time that inner \makeatletter could act on
% it, so the \gdef silently created/redefined the wrong thing and never
% touched beamer's actual mode variable. Confirmed via \meaning\beamer@
% currentmode before/after \studentversion -- unchanged with the inside-the-
% body version, fixed once \makeatletter moved out here.
\makeatletter
\newcommand{\studentversion}{%
	\global\gothamStudentVersiontrue
	\gdef\beamer@currentmode{handout}%
}
\makeatother

% Deck version system - new boolean flags and string for current version
\newif\ifgothamInstructorVersion
\newif\ifgothamDuringCourseVersion
\newif\ifgothamAfterCourseVersion
\gdef\gothamDeckVersion{instructor}
\gdef\gothamDeckVersionText{}

% \deckversion[] - set the deck version(s)
% Usage: \deckversion[instructor], \deckversion[during-course], etc.
% Multiple versions: \deckversion[instructor,during-course,after-course]
% Parses #1 as a comma-separated list via etoolbox's \forcsvlist (rather
% than the single \ifstrequal{whole-list}{"instructor"} check this used to
% do, which only ever matched a single-version list -- a multi-version
% \deckversion[instructor,during-course] silently set none of the three
% booleans, since the compared string was the literal, unsplit
% "instructor,during-course"). Note \forcsvlist, not etoolbox's similarly-
% named \forlistloop: \forlistloop iterates etoolbox's own internal
% pipe-delimited list format (built via \listadd), not a plain comma string
% like this command's own argument -- confirmed directly (isolated test),
% passing a raw comma string to it doesn't error or split on the commas, it
% just treats the whole string as a single one-item "list".
\makeatletter
\newcommand{\deckversion}[1][instructor]{%
	% \xdef (not \gdef): #1 needs to be fully expanded here, not just
	% stored as whatever token(s) were passed -- \deckversion[\myvar]
	% (compile-all-versions.bat's own override mechanism, via a wrapper
	% file's \def\deckversionoverride{...}) passes a macro, not a literal
	% string, and everywhere else in this deck-version system that reads
	% \gothamDeckVersion only ever applies one \expandafter to it (see
	% \ifdeckversioninlist's own comment) -- enough to unwrap \gothamDeckVersion
	% itself, but not a further macro still sitting inside it. Confirmed
	% directly: with the old \gdef, \deckversion[\deckversionoverride]
	% left \gothamDeckVersionText empty and every \ifdeckversioninlist
	% check false, since \gothamDeckVersion's own stored content was the
	% single unexpanded token \deckversionoverride, not "during-course".
	\xdef\gothamDeckVersion{#1}%
	\global\gothamInstructorVersionfalse
	\global\gothamDuringCourseVersionfalse
	\global\gothamAfterCourseVersionfalse
	\gdef\gothamDeckVersionText{}%
	\def\do@deckversions##1{%
		\ifstrequal{##1}{instructor}{\global\gothamInstructorVersiontrue\gdef\gothamDeckVersionText{instructor deck}}{}%
		\ifstrequal{##1}{during-course}{\global\gothamDuringCourseVersiontrue\gdef\gothamDeckVersionText{student course deck}}{}%
		\ifstrequal{##1}{after-course}{\global\gothamAfterCourseVersiontrue\gdef\gothamDeckVersionText{student final deck}}{}%
	}%
	% \forcsvlist itself splits on literal commas in its own raw,
	% unexpanded argument (confirmed directly against etoolbox's own
	% source, \etb@defparser@do@aux) -- passing #1 straight through here
	% has the exact same problem as the \gdef above did, one call site
	% later. Building the whole \forcsvlist{...}{...} call inside an
	% \edef (protecting \forcsvlist/\do@deckversions themselves with
	% \noexpand so they're invoked *later*, once this constructed call is
	% run, not expanded away here) expands only \gothamDeckVersion --
	% already fully resolved by the \xdef above -- into the list argument.
	\edef\gotham@deckversion@call{\noexpand\forcsvlist{\noexpand\do@deckversions}{\gothamDeckVersion}}%
	\gotham@deckversion@call
}
\makeatother

\newif\ifgothamOmittedSlide
\newcommand{\gothamOmittedSlideMarker}{ | omitted from student deck}
% \stepcounter{framenumber} in the skipped branch below keeps the student
% deck's slide numbers matching the instructor deck's, gap and all, so "see
% slide 12" means the same physical slide to both audiences: beamer steps
% "framenumber" exactly once per \begin{frame} (checked in
% beamerbaseframe.sty -- once per *frame*, not per overlay page, so this is
% right regardless of how many overlay steps the omitted slide would have
% had), and this is the one place that increment is skipped along with the
% frame itself, so it has to be replayed by hand.
\newcommand{\omitslide}[1]{%
	\ifgothamStudentVersion
		\stepcounter{framenumber}%
	\else
		\global\gothamOmittedSlidetrue
		\setSlideID{#1}%
		\input{slides/#1-current.tex}%
		\global\gothamOmittedSlidefalse
	\fi
}
\newcommand{\omitslideversion}[2]{%
	\ifgothamStudentVersion
		\stepcounter{framenumber}%
	\else
		\global\gothamOmittedSlidetrue
		\setSlideID{#1}%
		\input{slides/#1-v#2.tex}%
		\global\gothamOmittedSlidefalse
	\fi
}

% Deck version-aware overlay commands
% Syntax: \deckonly<overlay>[versions]{content}
% If [versions] omitted, defaults to all decks (instructor,during-course,after-course)
%
% Built on etoolbox's own \forcsvlist for the comma-splitting (rather than a
% hand-rolled parser: an earlier version tried to split a list at each comma
% using \newcommand's plain, non-delimited argument grabbing
% (\newcommand{\ifdeckversioninlist@split}[3]{...} called as
% \ifdeckversioninlist@split#1\@nil{#1}) -- but a plain \newcommand doesn't do
% delimited-parameter matching against \@nil the way a \def with a matching
% pattern would; it just grabs the next 3 tokens/groups one at a time. So
% e.g. "instructor" silently split into single characters ("i", "n", "s",
% ...) instead of being read as one list item, which eventually ran the
% parser off the end of its own input and surfaced as "Undefined control
% sequence" wherever the first \conditionalpause[...]/\deckonly[...] in a
% deck happened to be used -- confirmed directly against the compile log,
% traced to slide y5z-CoCo-Bonds's own \conditionalpause[instructor].
%
% NOT etoolbox's own \ifinlist{item}{list}, despite the name similarity: that
% one tests membership in etoolbox's own internal list format, built only via
% \listadd (items joined with a literal "|", not ","); fed a plain comma
% string instead, confirmed directly (isolated test) to silently return
% false for every item, including an exact single-item match.
%
% The membership test itself needs one more subtlety: #1 here is the deck
% version to search for -- always \gothamDeckVersion at every call site
% below, i.e. a macro, not a literal string -- so comparing it via a plain
% \ifstrequal{##1}{#1} does not work either: etoolbox's \ifstrequal compares
% via \detokenize, which (confirmed directly, isolated test) stringifies an
% unexpanded macro token to its own control-sequence *name* ("gothamDeckVersion"),
% not the current value that macro holds. \expandafter forces that expansion
% before \ifstrequal ever sees it.
\makeatletter

\let\gotham@original@only\only
\let\gotham@original@uncover\uncover
\let\gotham@original@visible\visible

\newif\ifgotham@deckversion@found

% Usage: \ifdeckversioninlist{version}{list}{true-branch}{false-branch}
\newcommand{\ifdeckversioninlist}[2]{%
	\global\gotham@deckversion@foundfalse
	\def\gotham@deckversion@check##1{%
		\expandafter\ifstrequal\expandafter{#1}{##1}{\global\gotham@deckversion@foundtrue}{}%
	}%
	\forcsvlist{\gotham@deckversion@check}{#2}%
	\ifgotham@deckversion@found
		\expandafter\@firstoftwo
	\else
		\expandafter\@secondoftwo
	\fi
}

% \@ifnextchar-based parsing (the same idiom LaTeX2e's own kernel commands
% use for an optional argument that isn't a plain trailing [...]) so the
% overlay spec reads exactly like beamer's own \only<...>, with the deck-
% version list as a further, genuinely optional [...] afterwards:
%   \deckonly<2>{content}                              visible in all versions
%   \deckonly<2>[instructor]{content}                  instructor only
%   \deckonly<2>[during-course,after-course]{content}  during + after
\newcommand{\deckonly}{\@ifnextchar<{\gotham@deckonly@overlay}{\gotham@deckonly@overlay<1->}}
\def\gotham@deckonly@overlay<#1>{\@ifnextchar[{\gotham@deckonly@versions<#1>}{\gotham@deckonly@versions<#1>[instructor,during-course,after-course]}}
\def\gotham@deckonly@versions<#1>[#2]#3{%
	\ifgothamInstructorVersion
		\gotham@original@only<#1>{#3}%
	\else
		\ifdeckversioninlist{\gothamDeckVersion}{#2}{\gotham@original@only<#1>{#3}}{}%
	\fi
}

% Same interface as \deckonly, wrapping \uncover and \visible respectively
\newcommand{\deckuncover}{\@ifnextchar<{\gotham@deckuncover@overlay}{\gotham@deckuncover@overlay<1->}}
\def\gotham@deckuncover@overlay<#1>{\@ifnextchar[{\gotham@deckuncover@versions<#1>}{\gotham@deckuncover@versions<#1>[instructor,during-course,after-course]}}
\def\gotham@deckuncover@versions<#1>[#2]#3{%
	\ifgothamInstructorVersion
		\gotham@original@uncover<#1>{#3}%
	\else
		\ifdeckversioninlist{\gothamDeckVersion}{#2}{\gotham@original@uncover<#1>{#3}}{}%
	\fi
}

\newcommand{\deckvisible}{\@ifnextchar<{\gotham@deckvisible@overlay}{\gotham@deckvisible@overlay<1->}}
\def\gotham@deckvisible@overlay<#1>{\@ifnextchar[{\gotham@deckvisible@versions<#1>}{\gotham@deckvisible@versions<#1>[instructor,during-course,after-course]}}
\def\gotham@deckvisible@versions<#1>[#2]#3{%
	\ifgothamInstructorVersion
		\gotham@original@visible<#1>{#3}%
	\else
		\ifdeckversioninlist{\gothamDeckVersion}{#2}{\gotham@original@visible<#1>{#3}}{}%
	\fi
}

% \conditionalpause[versions] - \pause, but only takes effect in the listed
% deck versions (default: instructor only); a no-op, not even a skipped
% overlay step, in any other version
% Usage: \conditionalpause[instructor,during-course] or \conditionalpause[after-course]
\newcommand{\conditionalpause}[1][instructor]{%
	\ifdeckversioninlist{\gothamDeckVersion}{#1}{\pause}{}%
}

% \deckreveal[step]{content} - the "hidden answer, click-revealed for the
% instructor, absent for during-course, pre-shown for after-course" pattern
% already documented by hand on several exercise slides (e.g. mbz/x7q/37s/
% 3cv's own header comments: "a slide/reveal-step that is hidden in the
% source PowerPoint should, by default, only be made visible on the
% instructor version"). A bare, unconditional \only<N->{content} looks like
% it does this, but doesn't: \only's overlay range registers with beamer's
% own overlay-counting pass regardless of \ifgothamInstructorVersion/
% \conditionalpause -- neither of those touch \only itself, they only guard
% *other* content around it -- so a frame with nothing else establishing a
% second overlay step still gets a real second slide out of a bare
% \only<2->, in *every* deck version, not just the instructor one. Confirmed
% directly (2026-07-21, direct instruction) on 37s/3cv: both a) leaked the
% "hidden" content to during-course as a genuine second slide, the exact
% thing this pattern exists to prevent, and b) showed a blank second bullet/
% cell first on that leaked slide's own step 1 -- \only<2-> hides the
% *content*, not the space it would occupy, so during-course's own step 1
% still had nothing else deciding where the bullet/cell landed.
% \deckreveal fixes both by picking a genuinely different rendering per
% version rather than reusing the same overlay-restricted \only call
% everywhere: instructor gets the real overlay build (\only<step->,
% default step 2, matching \tallitem's own default-multiplier convention);
% during-course gets nothing at all (no token emitted, so beamer's overlay
% pass never sees an overlay spec to register); after-course gets the
% content unconditionally, with no overlay spec of its own either, so it
% doesn't silently force a second slide there either. Content-level only
% (wraps a value/expression, like \only itself) -- doesn't cover an
% itemize \item's own bullet marker, which needs its own overlay spec
% (`\item<step->`) directly on the \item, not wrapped from outside; see
% 37s's own comment for that variant, hand-written rather than factored
% into a second macro since only one call site needs it so far.
\newcommand{\deckreveal}[2][2]{%
	\ifgothamDuringCourseVersion
	\else
		\ifgothamInstructorVersion
			\only<#1->{#2}%
		\else
			#2%
		\fi
	\fi
}

% \conditionalslide[versions]{ID-slug} - \slide{}, but only included in the
% listed deck versions; keeps the frame counter in step with the omitted
% version otherwise (same \stepcounter{framenumber} technique \omitslide{}
% already uses for the student deck), so "see slide N" still points at the
% same physical slide across every deck version.
% Usage: \conditionalslide[instructor,after-course]{xyz-Answer-key}
\newcommand{\conditionalslide}[2][instructor,during-course,after-course]{%
	\ifdeckversioninlist{\gothamDeckVersion}{#1}%
		{\setSlideID{#2}%
		 \input{slides/#2-current.tex}}%
		{\stepcounter{framenumber}}%
}

\makeatother

% PDF accessibility: mark purely decorative/running content (footer, below;
% also reused on titleframe's/quoteframe's own decorative background boxes;
% logo not yet included, see README.md's own accessibility section for that
% open question) as a PDF "Artifact", so screen readers skip it -- via
% LuaTeX's raw \pdfextension literal
% primitive. \ifluatex-guarded (iftex is already loaded, via beamer.cls
% itself) because \pdfextension is LuaTeX-only -- this project's README
% documents xelatex as a supported alternative engine, and compiling with it
% hit exactly this: "Undefined control sequence: \pdfextension", hard-failing
% every frame, since the first version of this didn't guard the primitive at
% all. Under any non-LuaTeX engine these two commands now silently do
% nothing instead: decorative content just isn't marked as an Artifact under
% xelatex/pdflatex, rather than breaking the compile -- a real loss of
% accessibility under those engines specifically, not a full equivalent fix,
% but not something to solve today (see also: compiling with plain pdflatex
% already silently mis-renders the whole deck's fonts, so xelatex is really
% the only *other* engine this needs to not break under; pdflatex was never
% viable here to begin with, per "Compiling" in README.md).
% Deliberately NOT tagpdf/kernel tagging (\DocumentMetadata):
% beamer 3.77 refuses to compile at all when that's active ("The beamer
% class is not compatible with \DocumentMetadata ... switch to the new
% ltx-talk class if you require tagged output"), and ltx-talk is a from-
% scratch class (not beamer-based, no theme/template system) that this
% framework can't realistically migrate to. Bypassing beamer's check by
% force-loading the underlying package before \documentclass also failed --
% it gets past hyperref but then hits a genuine internal error in tagpdf's
% LuaTeX backend from the non-standard load order, not something to build
% real coursework on.
%
% /Artifact is a standalone PDF marked-content tag (PDF spec 14.8.2.2):
% valid in *any* PDF, with or without a structure tree, no /Properties
% resource entry or MCID required -- unlike structure-tree tagging, this
% doesn't require kernel tagging support at all. It's a real, working
% partial fix, not full PDF/UA conformance: there's no structure tree, so
% no true semantic headings, and reading order elsewhere on the slide still
% depends on plain content-stream order (verified, not enforced, for
% standardframe -- see its own comment).
% #1 is the artifact's inline PDF dictionary content, e.g.
% "/Type /Pagination /Subtype /Footer" for a running footer, or
% "/Type /Background" for a decorative fill -- not hardcoded to one /Type,
% since not everything decorative here is "Pagination" in the PDF spec's own
% sense of that word (running headers/footers/page numbers/watermarks);
% titleframe's/quoteframe's plain background color fills are semantically
% "Background" artifacts instead, a distinct standard type.
\newcommand{\gothamArtifactBegin}[1]{%
	\ifluatex\pdfextension literal direct{/Artifact <<#1>> BDC}\fi
}
\newcommand{\gothamArtifactEnd}{\ifluatex\pdfextension literal direct{EMC}\fi}

% Footer
% Smaller footer + a bit of padding under it
\setbeamerfont{frame footer}{size=\fontsize{5}{3}\selectfont} % {font size}{baseline skip} in pt

% Shrinks the gap between footnotes and the footer box below them. The
% theme positions the footer's content via \vspace*{\gothamFootlineVOffset}
% (a fixed, negative length) before it starts; making that more negative
% turns out to give footnotes (which are placed just before the footer, and
% stretch to fill whatever room remains via beamer's own
% \beamer@framebottomskip) more room to extend into, pulling them down
% towards the footer -- checked empirically against the PDF's own
% coordinates, this does *not* move the footer's own text at all (that
% still sits exactly where \gothamFooterBoxHeight below was calibrated for),
% so the footer box itself needs no compensating adjustment.
%
% Partly given back (+1.4pt) 2026-07-05: a footnote/source containing an
% underlined \href (see the hyperlink-styling comment above) sits close
% enough to the footer at the original -1.3pt that the underline visually
% touches the footer box (checked directly against the compiled PDF's own
% pixel coordinates: a plain-text footnote's own gap to the footer measured
% ~1.9pt, but an underlined-link footnote's gap measured only ~0.6pt at
% the same setting, since the strut rebuilt for footnotes just below --
% correctly sized for plain text's own natural depth -- doesn't account
% for the underline rule extending further down). +1.4pt splits the
% difference: a linked footnote's gap now measures close to the plain-
% footnote reference (~1.9pt) instead of colliding with the footer, and a
% plain footnote's own gap just becomes a little more relaxed (~3.3pt)
% rather than stretched uncomfortably tight -- applied globally rather
% than only to footnotes with links, since there's no clean per-footnote
% hook for it and the plain-footnote side effect is minor.
%
% Nudged back down by 0.4pt the next day (direct instruction: move
% footnotes down "by the width of the underline line") -- the +1.4pt above
% left a bit more headroom than necessary; 0.4pt is the underline rule's
% own measured stroke thickness (a solid ~1px core at 600dpi, i.e.
% 1/600in \approx 0.12pt, plus antialiasing bleed either side consistent
% with the common 0.4pt LaTeX default rule width -- measured directly
% against the compiled PDF's own pixel values, not assumed).
%
% Reverted the very next round (direct instruction: move back *up* by that
% same underline thickness) -- net effect is back to the plain +1.4pt this
% section started with; kept as two separate historical adjustments below
% rather than collapsing them into one, so this back-and-forth stays
% visible rather than silently disappearing.
\addtolength{\gothamFootlineVOffset}{-1.3pt}
\addtolength{\gothamFootlineVOffset}{1.4pt}
\addtolength{\gothamFootlineVOffset}{-0.4pt}
\addtolength{\gothamFootlineVOffset}{0.4pt}

% 2026-07-06 (direct instruction): move footnotes down by half the
% measured distance between a linked footnote's own underline and the
% footer box, checked on slides/3ic-Corporate-Venture-Capital-
% current.tex specifically (the underline sits at y=245.95pt, the footer
% box's own top edge at y=247.69pt, a 1.74pt gap) -- half of that is
% 0.87pt, applied as a further negative adjustment (more negative =
% footnotes get pulled further down, per this section's own comment
% above).
\addtolength{\gothamFootlineVOffset}{-0.87pt}

% Shifts the right end of the footer line (the frame number) so it lines up
% with the logo's own right edge above it, rather than wherever the theme's
% default footline margins would otherwise put it. Empirically calibrated
% against the rendered PDF coordinates of both; re-measure and adjust if
% \gothamFrametitleRightpading (the logo's own right inset) changes.
\newlength{\gothamFooterRightShift}
\setlength{\gothamFooterRightShift}{4.6pt}

% Nudges the footer text down within its own yellow box, to vertically
% center it there: it previously sat noticeably closer to the box's top
% edge than its bottom (0.14pt above vs. 1.12pt below, checked against the
% rendered PDF's own text coordinates) since \gothamFootlineVOffset's own
% positioning (further up) turns out not to move this text at all (see its
% comment above) -- this is a separate, additive leading \vspace* instead.
\newlength{\gothamFooterTextDrop}
\setlength{\gothamFooterTextDrop}{-0.49pt}

% Compile timestamp: lets you tell, from a PDF alone (no need to check file
% properties/mtime, which don't survive e.g. emailing or uploading a copy),
% whether it reflects the latest edits or a stale earlier build -- handy
% whenever a PDF's *content* is in question, not just when debugging compile
% tooling itself. LuaTeX-only, matching this project's required engine (see
% README.md's "Compiling" section) -- no \ifluatex guard needed since
% nothing else here works under another engine to begin with.
%
% Deliberately avoids os.date's own %-format directives (e.g. "%Y-%m-%d")
% entirely, building the string from os.date("*t")'s field table and plain
% Lua concatenation instead -- \directlua's argument is tokenized by plain
% TeX first, where % starts a comment, before the text ever reaches Lua as a
% string. The obvious-looking fix, escaping each one as \%, does *not* work:
% confirmed directly, it produces a Lua string containing a literal
% backslash followed by a percent sign (detokenization keeps a control
% sequence's own backslash, for \% same as any other), which Lua's own
% string-literal parser then rejects outright ("invalid escape sequence"),
% silently aborting that whole \directlua call -- with no visible LaTeX-side
% error, just a cache file ending up written with an empty value. Switching
% catcode temporarily (\catcode`\%=12 ... ) would also work but needs a
% \begingroup/\endgroup wrapper at every call site; skipping % entirely is
% simpler and needs no such wrapper.
%
% Cached to \jobname.timestamp (via plain TeX \write, so MiKTeX's
% -aux-directory flag -- already applied through .latexmkrc's %O -- routes
% it into .build/ like every other aux file, not left cluttering decks/ or
% the project root) rather than recomputed fresh on every pass: a live
% os.date() call on every pass, tried first, broke latexmk's own rerun
% convergence outright -- confirmed directly, by isolating the change, that
% an otherwise-identical clean compile converges normally without it. Each
% of latexmk's internal reruns relaunches lualatex as a brand-new process,
% and a full pass over a deck this size (font loading alone) comfortably
% crosses a one-minute boundary, so two consecutive passes embedded two
% different strings in the footer; latexmk saw the .aux checksum differ
% *again* on the very next pass and looped ("Latexmk: References changed.")
% until hitting its rerun limit and reporting a false "gave an error", with
% no actual LaTeX error anywhere in the log. Caching to a file read back on
% later passes fixes that: the value is fixed the moment it's first
% written, so the footer's rendered text -- and therefore the .aux
% checksum -- stays identical across every subsequent pass of the same
% compile.
%
% This cache is deliberately *not* self-invalidating by date/content -- it
% keeps reporting an old timestamp indefinitely once written, across
% unrelated future compiles, unless something clears it first.
% SlideManager's compiler.py does exactly that: before every job it deletes
% that job's own .build/<jobname>.* (which includes <jobname>.timestamp), so
% a run through SlideManager always gets a fresh stamp. An ad-hoc TeXstudio
% compile without going through SlideManager first may show a stale stamp if
% .build/ already has one left over from an earlier session -- delete
% .build/<jobname>.timestamp by hand (or just compile that job once through
% SlideManager) if that matters for what you're checking.
\directlua{
	function gothamZeroPad(n) if n < 10 then return "0" .. n else return "" .. n end end
}
\newwrite\gothamCompileTimestampFile
\newcommand{\gothamCompileTimestamp}{%
	\ifdefined\gothamCompileTimestampValue
	\else
		\IfFileExists{\jobname.timestamp}{%
			\input{\jobname.timestamp}%
		}{%
			\xdef\gothamCompileTimestampValue{\directlua{%
				local d = os.date("*t")%
				tex.sprint(d.year .. "-" .. gothamZeroPad(d.month) .. "-" .. gothamZeroPad(d.day) .. " " .. gothamZeroPad(d.hour) .. ":" .. gothamZeroPad(d.min))%
			}}%
			\immediate\openout\gothamCompileTimestampFile=\jobname.timestamp\relax
			\immediate\write\gothamCompileTimestampFile{\string\gdef\string\gothamCompileTimestampValue{\gothamCompileTimestampValue}}%
			\immediate\closeout\gothamCompileTimestampFile
		}%
	\fi
	\gothamCompileTimestampValue
}

% \gothamHideFooterText: opt-in flag suppressing just \FooterText (the left
% part of the footer) -- default false. Set true by the standout template's
% own hook below, so `\begin{frame}[standout]` hides only that left part,
% keeping the "(slide id: ...)" block and frame number visible (2026-07-16,
% direct instruction). Reset immediately after being read, inside the
% footer template itself, rather than relying on a per-environment closing
% hook (standout frames are bare `\begin{frame}[standout]...\end{frame}`,
% not one of this project's own environments like standardframe, so there's
% no custom closing code to reset it in) -- safe because \gothamHideFooterText
% is only ever consumed here, once, immediately before being cleared.
\newif\ifgothamHideFooterText
% gotham's own \defbeamertemplate*{standout}{gotham} (gotham.dtx) is what
% `\begin{frame}[standout]` actually activates; \addtobeamertemplate's
% "after" hook appends to it without needing to touch that upstream file.
\addtobeamertemplate{standout}{}{\global\gothamHideFooterTexttrue}

\setbeamertemplate{frame footer}{%
	\vspace*{\gothamFooterTextDrop}%
	\ifgothamHideFooterText
		\global\gothamHideFooterTextfalse
	\else
		\FooterText%
	\fi
	\hfill%
	{\fontsize{4}{4}\selectfont (slide id: \slideID\ifgothamOmittedSlide\gothamOmittedSlideMarker\fi\ | \gothamDeckVersionText\ | \gothamCompileTimestamp)}\hspace*{2em}\insertframenumber%
	\hspace*{-4em}%
	\hspace*{-\gothamFooterRightShift}%
	\vspace*{0.4ex}%
}

% Give the footer the same yellow-box treatment as the title bar, but --
% unlike the title bar, which is tall for visual weight -- only as tall as
% the footer's own text plus equal padding above and below.
%
% This paints a yellow patch *behind* the theme's own footline rather than
% replacing that template's content, because \usebeamertemplate{frame~footer}
% (which is what actually typesets \FooterText etc.) only produces its
% content when invoked from the theme's own original footline template --
% calling it again from anywhere else (a redefined footline, the frame body,
% a \settoheight measurement, ...) silently produces nothing. It's
% positioned from "current page.south" (the page's own bottom-left corner)
% rather than a \tikz "remember picture" node tracking the footer's actual
% position, since there's no per-frame content to track: the footer text
% (in its actual 5pt font, per \setbeamerfont{frame footer} above -- despite
% the "page number in head/foot" font switch inside the theme's footline
% template, this is the font that actually ends up applied) sits at a fixed
% distance from the bottom of the page on every frame.
%
% \gothamFooterBoxHeight below is calibrated to that actual measured
% position (checked via the PDF's own text coordinates, not just visually):
% the text sits about 0.84pt above the page's bottom edge and spans about
% one em of its own 5pt font, plus a bit of padding above and below. If the
% footer font size changes, this will need re-measuring.
\newlength{\gothamFooterBoxHeight}
\setlength{\gothamFooterBoxHeight}{7.26pt}
% \vbox to 0pt{...} forces this to occupy exactly zero space in the normal
% vertical list, *without* switching into horizontal mode the way an hbox
% (e.g. \raisebox) would -- inserting an hbox here forces LaTeX to open an
% implicit paragraph before the theme's own \vspace*{\gothamFootlineVOffset}
% runs, which changes how that \vspace behaves (it special-cases being at
% the very start of vertical mode) and nudges the footer text down by a few
% points. A \vbox is already vertical material, so it doesn't trigger that.
% \gothamArtifactBegin/\gothamArtifactEnd (defined above, "PDF accessibility"
% comment) wrap both this box *and* the original footline template's own
% content (the theme's \usebeamertemplate{frame footer} call, i.e.
% \FooterText/slide id/page number) -- \addtobeamertemplate's "after" slot
% runs once the wrapped-around original template content has finished, so
% \gothamArtifactEnd there closes the same marked-content span this begins.
% The footer box's own south-west corner originally measured 0.2pt short
% of the page's actual left/bottom edges on every slide checked
% (2026-07-06, direct instruction: "check whether the footer ... has a bit
% of white on its left and below") -- confirmed directly against the
% compiled PDF's own drawn rectangle (consistently ~0.199pt in from both
% edges, not a one-off), unlike the title bar (gotham's own native
% frametitle background, a completely different mechanism -- normal TeX
% box placement, not a tikz "current page" lookup), which sits exactly
% flush at (0,0). At the time, the root cause wasn't chased further and
% the box was fixed by hand-nudging the anchor 0.2pt past the corner and
% padding the size to match. Root cause since found (2026-07-16, while
% chasing the identical hairline on titleframe/quoteframe -- see the
% frametitle hook's own comment above): tikz's default outer sep
% (0.5\pgflinewidth = 0.2pt) puts every edge/corner anchor 0.2pt outside
% the filled rectangle, on *every* edge-anchored node, so the earlier
% "nothing else anchored to a page corner shows the same offset" turned
% out to be wrong. "outer sep=0pt" below is the real fix, replacing the
% nudge-and-pad workaround; the box now lands exactly flush by
% construction rather than by measured compensation.
\addtobeamertemplate{footline}{%
	\gothamArtifactBegin{/Type /Pagination /Subtype /Footer}%
	\vbox to 0pt{%
	\vss
	\begin{tikzpicture}[remember picture, overlay]
		\node[fill=\gothamFooterBoxColor, anchor=south west, inner sep=0pt, outer sep=0pt,
		      minimum width=\paperwidth, minimum height=\gothamFooterBoxHeight]
			at (current page.south west) {};
	\end{tikzpicture}%
	}%
}{\gothamArtifactEnd}

% Footnotes: small font (smaller than the footer's own 5pt), symbol markers
% (*, dagger, ...) instead of arabic numbers, right-aligned (flush with the
% logo's own right edge, like the footer's frame number), no rule above
% them, and tightened-up line spacing.
%
% Line spacing needs two different levers, not just \baselineskip, because
% at this font size (5pt) \baselineskip's normal value is *smaller* than the
% font's own natural line height, so TeX falls back to \lineskip (a flat
% add-on) between every line regardless -- \baselineskip itself turns out to
% have no effect here at all. That fallback \lineskip is what governs
% spacing *within* a single (wrapped) footnote's own lines: reduced only a
% tiny amount, from the default 1pt.
% The *between*-footnote gap is inflated by a separate mechanism entirely:
% \footnotesep, a rule LaTeX inserts before every individual footnote entry
% (measured to be 7.7pt here -- clearly sized for a much bigger font than
% ours), so that's what actually needs the bigger reduction; \parskip
% (tried first) turned out to have no visible effect on it.
\setlength{\footnotesep}{3pt}
% \reset@font\footnotesize is applied by beamer's own footnote machinery
% before this template runs (beamerbaseframecomponents.sty), so the font
% size has to be set again here to actually take effect.
\renewcommand{\footnoterule}{}
% Same marker sequence as LaTeX's own \fnsymbol (*, dagger, double dagger,
% ...), but built from plain Unicode text characters instead: \fnsymbol's
% symbols are typeset via \ensuremath{...}, i.e. in *math* mode, and the
% theme loads fontspec with the "no-math" option (see
% beamerfontthemegotham.sty) specifically so that Fira Sans -- the font
% fontspec puts everything else in -- is left in charge of text only, not
% math; math mode therefore still falls back to LaTeX's default Computer
% Modern math fonts. \fnsymbol's markers were consequently rendering in a
% different font than the surrounding text (checked by comparing rendered
% glyph shapes) -- most visible as a marker that looked shifted left of
% flush-left text next to it (e.g. a \source{} line, which has no marker of
% its own), since Computer Modern's symbol glyphs have different sidebearings
% than Fira Sans's. Plain characters typeset with no \ensuremath wrapper stay
% in text mode and pick up whatever font is already active at the call site,
% i.e. Fira Sans -- matching the rest of the footnote.
% Literal Unicode characters, not the \textdagger/\textdaggerdbl/\textbardbl
% kernel commands: those need the textcomp/TS1 setup to render as plain text,
% and silently fall back to \ensuremath{...} (i.e. math mode, the very thing
% this replaces \fnsymbol to avoid -- see the comment above) when that isn't
% loaded. LuaLaTeX plus fontspec maps raw UTF-8 input straight to the current
% font's (Fira Sans's) own glyphs with no encoding layer in between, which is
% the entire point of using them -- so writing the characters directly, with
% no wrapper command, guarantees they stay in text mode and stay in Fira Sans.
%
% The sequence itself is user-configurable via \setfootnotesymbols{...} (a
% comma-separated list, e.g. \setfootnotesymbols{*,+,#}) -- call it anywhere
% from that point in the deck onward (e.g. once, right after
% % Slidemaster preamble variant for University of Graz decks. Institution-
% specific content only (colors, logo, logo calibration) -- every frame
% type, the deck-version system, footer/hyperlink/footnote styling, tikz
% styles, etc. live in Slidemaster_preamble_base.tex, \input{}ed at the
% bottom of this file, and are shared verbatim across every institution
% this project supports (see Slidemaster_preamble_mci.tex and
% Slidemaster_preamble_uni-ibk.tex for the other two). A deck picks its
% institution by \input{}ing this file instead of one of those, nothing
% else changes.

% Color definition (Graz palette, sourced from the original single-preamble
% version of this project -- these are the University of Graz's own official
% colors, not derived/approximated).
\definecolor{c_primary_dark}{HTML}{000000}
\definecolor{c_primary_highlight}{HTML}{ffd500}
\definecolor{c_primary_gray}{HTML}{c6c6c6}
\definecolor{c_primary_white}{HTML}{FFFFFF}
\definecolor{c_secondary_violet}{HTML}{948dc3}
\definecolor{c_tertiary_violet}{HTML}{c2bddf}
\definecolor{c_secondary_blue}{HTML}{689cce}
\definecolor{c_tertiary_blue}{HTML}{bacce8}
\definecolor{c_secondary_green}{HTML}{9fcb81}
\definecolor{c_tertiary_green}{HTML}{cce2ba}
\definecolor{c_secondary_red}{HTML}{e4684b}
\definecolor{c_tertiary_red}{HTML}{f1ae95}
\definecolor{c_secondary_orange}{HTML}{f6a968}
\definecolor{c_tertiary_orange}{HTML}{fbd0aa}
\definecolor{c_secondary_yellow}{HTML}{ffe063}
\definecolor{c_tertiary_yellow}{HTML}{ffeeaf}
\definecolor{c_tertiary_gray}{HTML}{e3e3e3}
% Not an official Uni Graz color (hence "custom", not "secondary" like the
% palette above) -- a softer near-black than c_primary_dark, used for
% quoteframe's full-bleed dark block, matching the specific dark gray (not
% pure black) ExampleSlides.pptx used for that same block. Also reused for
% hyperlink text color, for the same reason it works for quoteframe's
% attribution text: dark enough for good contrast/legibility on a light
% background.
\definecolor{c_custom_gray}{HTML}{333333}
\definecolor{c_custom_red}{HTML}{C00000}
\definecolor{c_custom_green}{HTML}{00B050}

% Logo + calibration.
% schemes/scheme_uni_graz/Logo_Uni-Graz_4c-cropped.pdf -- official vector artwork,
% user-supplied (schemes/scheme_uni_graz/Logo_Uni-Graz_4c.eps, kept alongside as the
% untouched source), converted to PDF via epstopdf and cropped to its true
% ink bounding box (computed with Ghostscript's own `-sDEVICE=bbox`, same
% treatment as MCI's and Innsbruck's own logos -- see
% Slidemaster_preamble_mci.tex's own comment for the fuller explanation of
% why the EPS's own %%BoundingBox isn't trusted directly). Vector, not
% raster, so it stays crisp at any zoom. Supersedes this file's own earlier
% raster logo (schemes/scheme_uni_graz/Logo Uni Graz 4c-cropped.png, itself cropped
% from schemes/scheme_uni_graz/Logo Uni Graz 4c.png's own symmetric transparent
% margin -- both kept alongside as historical/fallback assets, not deleted).
\newcommand{\gothamInstituteLogoPath}{schemes/scheme_uni_graz/Logo_Uni-Graz_4c-cropped.pdf}
\newlength{\gothamLogoHeight}
\setlength{\gothamLogoHeight}{3ex}
% Re-measured against the compiled PDF for this specific logo image -- see
% Slidemaster_preamble_base.tex's own comment on \gothamFrametitleRightpading
% for why this needs re-measuring per institution rather than reusing
% another one's value.
\setlength{\gothamFrametitleRightpading}{0.28cm}

% Itemize bullet color (2026-07-16, direct instruction): Graz keeps the
% project's original behavior, bullets in the primary highlight color.
\newcommand{\gothamBulletColor}{c_primary_highlight}

% Title bar and footer box/text colors (2026-07-16, direct instruction):
% Graz keeps the project's original behavior throughout -- title/footer
% boxes in the primary highlight color, title/footer text in the primary
% dark color (unlike MCI/Innsbruck, which switch these to their own
% secondary blue/white -- see Slidemaster_preamble_mci.tex's own comment).
\newcommand{\gothamTitleBoxColor}{c_primary_highlight}
\newcommand{\gothamTitleTextColor}{c_primary_dark}
\newcommand{\gothamFooterBoxColorDefault}{c_primary_highlight}
\newcommand{\gothamFooterTextColor}{c_primary_dark}

\input{Slidemaster_preamble_base}
) to replace the default marker set for every
% \footnote{}/\source{} from then on. \seq_item:Nn (expl3's indexed lookup;
% expl3 is already loaded by the gotham theme itself, see e.g.
% beamerfontthemegotham.sty) is what removes the old \ifcase's hard 9-symbol
% ceiling -- a custom list can be any length. Indexing wraps around (via
% \int_mod:nn) rather than erroring past the end of the list, so an unusually
% footnote-heavy slide degrades to repeating symbols instead of a hard
% failure; slides/ZZU-Footnote-symbol-options-current.tex is a reference
% sheet of candidate Unicode symbols to build a custom list from.
%
% Project-wide default (chosen deliberately, not the classic */dagger/etc.
% set): a currency-sign sequence, in this specific order. All 14 were
% confirmed against slides/ZZU-Footnote-symbol-options-current.tex to
% actually exist in Fira Sans.
\ExplSyntaxOn
\seq_new:N \g_gotham_footnote_symbols_seq
\seq_gset_from_clist:Nn \g_gotham_footnote_symbols_seq { £,¥,₹,₽,₺,₴,₪,₫,₦,₡,₩,ƒ,¢,€ }
\newcommand{\setfootnotesymbols}[1]{%
	\seq_gset_from_clist:Nn \g_gotham_footnote_symbols_seq { #1 }
}
\newcommand{\gothamFnSymbol}[1]{%
	\seq_item:Nn \g_gotham_footnote_symbols_seq
		{ \int_eval:n { \int_mod:nn { #1 - 1 } { \seq_count:N \g_gotham_footnote_symbols_seq } + 1 } }
}
\ExplSyntaxOff
\renewcommand{\thefootnote}{\gothamFnSymbol{\value{footnote}}}
% A \footnote{} called inside a minipage doesn't use \thefootnote/\c@footnote
% at all -- LaTeX's own minipage machinery silently switches to a separate
% counter/display command, \c@mpfootnote/\thempfootnote, normally displayed
% via \@alph (a, b, c, ...). Found 2026-07-14 converting
% slides/j6h-Foreign-Exchange-Market-current.tex (its first slide combining
% \footnote{} with a minipage, since a narrowed text column was added
% alongside a corner chart): the footnote *text* still correctly landed in
% this deck's usual bottom-of-frame collection (beamer's own footnote
% storage isn't minipage-scoped), but the *marker*, both inline and at the
% bottom, showed as a plain "a"/"b"/"c" instead of this project's own
% currency-symbol sequence. Redefined the same way as \thefootnote so any
% future \footnote{} inside a minipage picks up the right symbol too --
% \c@mpfootnote resets to 0 for every new minipage on its own (unlike
% \c@footnote, which needs standardframe's own explicit per-frame reset
% below), which happens to already match this deck's "restart every slide"
% footnote convention for the common case of one minipage per frame.
\renewcommand{\thempfootnote}{\gothamFnSymbol{\value{mpfootnote}}}
\setbeamerfont{footnote}{size=\fontsize{3.2}{1.9}\selectfont}
% Beamer's own \@makefnmark (used for the marker both inline in the body
% text and again next to the footnote itself) hard-codes the "footnote
% mark" beamerfont onto it -- an *absolute* size, unrelated to whatever
% text surrounds it -- which is why the inline markers looked tinier than
% a normal superscript when that was set to this footnote's own small size:
% a fixed size that looks fine amid 3.2pt footnote text looks minuscule
% amid much bigger body text. Reverting to plain LaTeX's own \@makefnmark
% instead makes it a normal \textsuperscript of whatever font size is
% already active at each call site -- proportionally scaled either way, so
% it looks like a normal superscript inline, and still comes out small next
% to the (small) footnote text since \usebeamerfont{footnote} above is
% already active there by the time it's invoked.
\makeatletter
\renewcommand{\@makefnmark}{\hbox{\@textsuperscript{\normalfont\@thefnmark}}}
\makeatother
% \raggedleft's own \rightskip (0pt) is flush with the frame's normal body
% text margin, not the logo's edge above it; this pushes further right by
% the difference between the two, empirically calibrated (like
% \gothamFooterRightShift above) against the rendered PDF coordinates of
% both -- re-measure and adjust if the text margins or logo inset change.
\newlength{\gothamFootnoteRightShift}
\setlength{\gothamFootnoteRightShift}{4.8pt}

% Side footnotes: \sidefootnotes (called *before* \begin{standardframe}, like
% \chapterlogo{}) moves that one frame's footnotes -- including \source{} lines,
% which are footnotes underneath -- from their usual spot above the footer to a
% vertical strip along the right edge of the page, rotated 90 degrees counter-
% clockwise: the text starts just above the footer box and reads bottom-to-top,
% as if the usual bottom-right footnote block had been turned on its side (with
% one deliberate difference: the lines are flush with the strip's bottom end,
% not the top -- see the \raggedright comment in the template below).
% standardframe resets the switch
% on \end{frame} (same \global mechanics as \chapterlogo{} -- see the comment
% there for why a plain local assignment doesn't survive beamer's frame group),
% so it opts in only the single frame that follows, and any frame type derived
% from standardframe inherits the behaviour automatically.
%
% How it works: the footnote *template* below branches on this switch. In side
% mode it typesets each footnote into \gothamSideFootnotesBox (a plain box
% register accumulated across the frame's footnotes) instead of contributing
% anything to beamer's own bottom-of-frame footnote insert, and the footline
% hook further down draws that box -- rotated -- as a tikz overlay once the
% whole frame body has run. The footline is the right place to draw from: it is
% typeset at page shipout, i.e. after every \footnote{} in the body has already
% executed, whereas e.g. the frametitle would run too early and find the box
% empty. (Beamer also typesets the footline once per frame just to *measure*
% it, at \begin{frame} time -- before any footnotes exist -- which is why the
% drawing hook guards on \ifvoid: the measurement pass must not consume, or
% try to draw, a box that doesn't exist yet.)
%
% Known quirk: a footnote with an overlay spec (\footnote<2->{...}) is
% collected -- and therefore shown -- on *every* overlay slide of its frame in
% side mode. Beamer implements covered footnotes by typesetting them invisibly
% in place, which relocating the text into a box register defeats.
\newif\ifgothamSideFootnotes
\newcommand{\sidefootnotes}{\global\gothamSideFootnotestrue}
\newbox\gothamSideFootnotesBox
% The rotated block's line width, i.e. the vertical run the strip may occupy
% down the right edge of the page: from just below the title bar to just above
% the footer's yellow box. Like \gothamFooterBoxHeight above, calibrated
% against the measured PDF coordinates: the title bar's bottom edge sits
% 219.57pt above the page's bottom edge; minus \gothamSideFootnoteTopSep (3pt),
% \gothamFooterBoxHeight (7.26pt), and ~2pt of clearance above the footer box
% leaves 207pt. Re-measure if the frametitle font or its padding ever changes.
\newlength{\gothamSideFootnoteHsize}
\setlength{\gothamSideFootnoteHsize}{207pt}
% Gap between the title bar's bottom edge and the top end of the strip. The
% text itself sits at the strip's *bottom* end (it is \raggedright, i.e.
% flush with the line start, which the rotation turns into bottom-aligned),
% so this gap only matters for footnotes long enough to fill the whole run.
\newlength{\gothamSideFootnoteTopSep}
\setlength{\gothamSideFootnoteTopSep}{3pt}
% Distance from the page's right edge to the strip's outer edge -- the rotated
% counterpart of the small gap the bottom footnotes keep above the footer box.
\newlength{\gothamSideFootnoteRightInset}
\setlength{\gothamSideFootnoteRightInset}{2pt}
% A real footnote's marker (e.g. "*") sits in an \hbox of its own whose
% *ink* -- unlike \source{}'s marker, which is emptied to nothing and so has
% no ink at all -- extends slightly left of that hbox's own nominal left
% edge (ordinary optical behaviour for a small centred glyph like an
% asterisk, exaggerated further by \@makefnmark's superscript raise). With
% every footnote's paragraph starting flush at the same \noindent left edge,
% that overhang made real footnotes' markers visibly poke out to the left of
% \source{}'s flush-left text above/below them. This shift nudges real
% footnotes' *entire* line (marker, comma, and text together) right by that
% overhang, so the marker's ink -- not the text after it -- lines up with
% \source{}'s left edge; \source{} itself is untouched; empirically
% calibrated against the rendered PDF, like the other shift constants here.
\newlength{\gothamSideFootnoteMarkerShift}
\setlength{\gothamSideFootnoteMarkerShift}{1.3pt}

% A new template variant (rather than reusing beamer's own "default", which
% it has already defined and then refuses to redefine) that we immediately
% activate below.
% Wrapped in \makeatletter (unlike the rest of this section) because the side
% branch below tests \@thefnmark against \@empty to tell a real footnote
% (non-empty marker) apart from a \source{} call (marker forced empty) --
% see \gothamSideFootnoteMarkerShift above.
\makeatletter
\defbeamertemplate{footnote}{gothamfn}{%
	\ifgothamSideFootnotes
		% Side mode (see \sidefootnotes above): append this footnote to the
		% collector box instead of typesetting it here. \insertfootnotetext
		% carries beamer's own \rule\z@\footnotesep lead-in and
		% \@finalstrut\strutbox tail along with it (beamer bakes both into the
		% text it hands the template), so the between-footnote spacing and the
		% strutbox rebuild work exactly as in the bottom branch below -- and,
		% equally, nothing of the footnote (not even the spacing rule) is left
		% behind in beamer's bottom-of-frame insert, which therefore stays
		% visually empty. No \gothamFootnoteRightShift here: that constant
		% compensates the *bottom* block's \raggedleft for the text-margin-vs-
		% logo-edge difference, while the strip's position is set directly by
		% the tikz node in the footline hook.
		\global\setbox\gothamSideFootnotesBox\vbox{%
			\unvbox\gothamSideFootnotesBox
			\hsize\gothamSideFootnoteHsize
			\usebeamerfont{footnote}%
			\setlength{\lineskip}{0.6pt}%
			\settoheight{\dimen0}{Xy}%
			\settodepth{\dimen1}{Xy}%
			\setbox\strutbox\hbox{\vrule height\dimen0 depth\dimen1 width0pt}%
			\parindent 1em\noindent%
			% \raggedright, not the bottom block's \raggedleft: after the
			% rotation, the line start is the *bottom* end of the strip, so
			% this makes the text begin just above the footer box and run
			% upward (rather than hang down from the title bar).
			\raggedright
			% \source{} forces \thefootnote (hence \@thefnmark) empty, so this
			% only fires for a real footnote's marker -- see
			% \gothamSideFootnoteMarkerShift above for why it needs the shift
			% and \source{} doesn't. The "\," below (a separator between the
			% marker and the footnote text) needs the same guard: it used to
			% run unconditionally, so an empty-marker \source{} still got a
			% thin space's worth of leading indent before its own text --
			% invisible under the *bottom* block's \raggedleft (nothing there
			% depends on the left edge), but a real, visible defect here under
			% \raggedright, where the left edge is exactly what's on display:
			% a wrapped \source{}'s first line sat measurably indented from
			% its own second line, which never picks up that one-time leading
			% space (2026-07-12, direct instruction: "the first line seems
			% indented compared to the second, which it shouldn't be").
			\ifx\@thefnmark\@empty
				\insertfootnotetext\par%
			\else
				\hspace*{\gothamSideFootnoteMarkerShift}\insertfootnotemark\,\insertfootnotetext\par%
			\fi
		}%
	\else
	\usebeamerfont{footnote}%
	\setlength{\lineskip}{0.6pt}%
	% beamer's footnote machinery resets to plain \footnotesize (a much
	% bigger font) before this template runs, and switching to our own
	% smaller font via \usebeamerfont doesn't update \strutbox to match --
	% so the invisible strut beamer appends after every footnote
	% (\@finalstrut\strutbox) is still sized for \footnotesize, silently
	% padding out extra depth below the last line. Rebuilding \strutbox
	% from this font's own measured metrics removes that padding, which is
	% most of the gap between the footnotes and the footer box below them.
	\settoheight{\dimen0}{Xy}%
	\settodepth{\dimen1}{Xy}%
	\setbox\strutbox\hbox{\vrule height\dimen0 depth\dimen1 width0pt}%
	\parindent 1em\noindent%
	\raggedleft
	\advance\rightskip by -\gothamFootnoteRightShift%
	% Same \@thefnmark emptiness guard as the side-footnote branch above,
	% for the same reason (an empty marker still leaving behind the "\,"
	% separator's own leading thin space) -- invisible here under
	% \raggedleft, where the left edge is never on display, but the same
	% underlying defect, so fixed the same way for consistency.
	\ifx\@thefnmark\@empty
		\insertfootnotetext\par%
	\else
		\insertfootnotemark\,\insertfootnotetext\par%
	\fi
	\fi
}
\makeatother
\setbeamertemplate{footnote}[gothamfn]

% Draws the side-footnote strip (see \sidefootnotes above). Lives in the
% footline template because that is typeset at page shipout, after the frame
% body -- and with it every \footnote{} call -- has run, so the collector box
% is complete by the time this executes. The \ifvoid guard doubles as the
% protection against beamer's measure-the-footline pass at \begin{frame} time
% (see the \sidefootnotes comment). \box (not \copy) empties the register
% globally as a side effect, which is exactly right: each overlay slide of a
% frame re-executes the body and re-collects its footnotes from scratch, so
% whatever this slide's footline draws must not leak into the next slide.
%
% Node geometry: tikz pins the node's *unrotated* anchor corner at the given
% coordinate and rotates the box counter-clockwise around it, so anchor=south
% east (the end of the collected block's last line) ends up as the strip's
% top-east corner: from there the block extends down the page (its lines run
% bottom-to-top, last footnote nearest the page edge), with the \raggedright
% text sitting at the strip's bottom end, just above the footer box (see the
% comment in the footnote template above). The y-coordinate
% comes from the same remembered gothamTitleBottom node the title-bar patches
% already use (same two-pass .aux mechanics -- see "Title bar and logo").
% The \ifcsname guard skips the drawing while the gothamTitleBottom shape is
% entirely unknown to pgf (i.e. no frametitle has been typeset yet in this
% run *and* nothing was inherited from a previous run's .aux) -- without it,
% referencing the node is a hard error rather than the usual harmless
% first-pass misplacement. In that situation the position would have been
% garbage anyway, and the frame-start clearing in standardframe below mops up
% the box this branch leaves unconsumed.
\makeatletter
\addtobeamertemplate{footline}{%
	\ifvoid\gothamSideFootnotesBox\else
		\ifcsname pgf@sh@ns@gothamTitleBottom\endcsname
		\vbox to 0pt{% zero-space trick, as in the yellow-box footline hook above
		\vss
		\begin{tikzpicture}[remember picture, overlay]
			\node[rotate=90, anchor=south east, inner sep=0pt]
				at ($(gothamTitleBottom -| current page.east)+(-\gothamSideFootnoteRightInset,-\gothamSideFootnoteTopSep)$)
				{\box\gothamSideFootnotesBox};
		\end{tikzpicture}%
		}%
		\fi
	\fi
}{}
% Empties the side-footnote collector (used by standardframe at frame start).
% Defined here, inside the \makeatletter block, because \voidb@x needs @ to be
% a letter at tokenization time -- standardframe's own definition below sits
% outside any \makeatletter region.
\newcommand{\gothamClearSideFootnotesBox}{\global\setbox\gothamSideFootnotesBox=\box\voidb@x}
\makeatother

% \source{text} -- for a "Source: ..."/"Image source: ..."/etc. citation line
% (common at the bottom/edge of a converted PowerPoint slide) that should look
% and sit *exactly* like the footnotes above -- same font, spacing,
% right-alignment, and the same stretch-to-sit-just-above-the-footer
% behaviour beamer already gives real footnotes -- but with no marker, neither
% inline in the body text nor next to the note itself, since a source line
% isn't anchored to one specific word. Implemented as a real \footnote{} (so
% it gets that positioning for free, rather than reimplementing it) with
% \thefootnote emptied locally; since \footnote{} still advances the footnote
% counter even with an empty marker, \addtocounter reverses that so a real
% \footnote{} elsewhere on the same slide doesn't skip a symbol in the
% \fnsymbol sequence. Unlike a real footnote, no label is added automatically
% here -- the label varies ("Source:", "Image source:", ...), so write it as
% part of #1.
\newcommand{\source}[1]{%
	\begingroup
	\renewcommand{\thefootnote}{}%
	\footnote{#1}%
	\endgroup
	\addtocounter{footnote}{-1}%
}

% A plain tabular with yellow \toprule/\midrule/\bottomrule, matching the
% frametitle/footer accent color. Takes the same column-spec argument as
% tabular itself, e.g. \begin{gothamtable}{lccc} ... \end{gothamtable}. Only
% booktabs rules are colored (via \arrayrulecolor, scoped to this table by
% \begingroup/\endgroup) -- if a table ever needs a plain black \hline too,
% that's unaffected since \arrayrulecolor only recolors rules drawn *after*
% it takes effect, and booktabs' rules read it fresh each time.
%
% RULE-BLEED WARNING (2026-07-12, direct instruction: "make sure the table
% lines do not extend beyond the text"), read this before writing a new
% gothamtable: a plain column spec like `{lccc}` above leaves the default
% \tabcolsep (6pt, or whatever a slide-local \setlength last set it to) as
% *implicit* padding on the table's own outermost edges -- invisible to the
% table's own text (which starts/ends flush with each column's actual
% content), but very visible to booktabs' \toprule/\midrule/\bottomrule and
% to colortbl's \rowcolor/\cellcolor, both of which span the *declared*
% column width including that implicit padding. The result, confirmed
% directly and independently on at least three slides so far
% (slides/u5h-Long-term-Bonds-Provisions-current.tex's own v1-1-0/v1-4-0,
% slides/c4c-Macbeth-Spot-Removers-Proposed-Structure-current.tex's own
% v1-2-2, slides/tw5-Exercise-Leverage-and-ROE-Solution-current.tex's own
% v1-4-0): every rule/highlight bleeds \tabcolsep past the visible text on
% both the left and right edges, most noticeable on a left-aligned
% (non-\centering) table where the mismatch is easy to see against the
% title's own left margin.
%
% Not fixed here in the environment definition itself: doing so would
% change every existing gothamtable's own column spacing out from under it
% (this environment has no way to tell "plain default spacing was intended"
% apart from "the bleed was never noticed"). Fix per table instead, the
% same three-part pattern every slide above uses -- worth reaching for by
% default on any *new* gothamtable, not just when the bleed is spotted:
%   {\setlength{\tabcolsep}{0pt}
%   \begin{gothamtable}{@{}l<{\hspace{<gap>pt}}c<{\hspace{<gap>pt}}c@{}}
%   ...
%   \end{gothamtable}}
% -- \tabcolsep zeroed (scoped locally, e.g. by wrapping the whole call in
% a `{...}` group, so it doesn't leak into later tables), explicit `@{}` at
% both outer edges of the column spec (so booktabs/colortbl have no
% implicit padding left to bleed into), and array's `<{\hspace{<gap>pt}}>`
% column modifier appended after every column but the last, to put back an
% explicit, *content-column* gap (matching whatever `2*\tabcolsep` the
% table used to have, if preserving its exact look, e.g. `<{\hspace{12pt}}>`
% for the ordinary 6pt-\tabcolsep default) -- `<{...}>` puts the gap
% *inside* the preceding column's own content, which `\rowcolor` (if used)
% colors along with everything else, unlike a plain `@{...}` separator,
% which `\rowcolor` deliberately leaves untinted (see u5h's own v1-6-0
% comment for that exact failure mode if a highlighted table ever needs
% this too).
%
% RE-CONFIRMED 2026-07-16, direct instruction, after this exact bleed
% reappeared on essentially every gothamtable in a freshly-converted
% chapter (Capital Budgeting Techniques) despite the warning above already
% existing at the time -- the rule above is not optional/only-if-noticed,
% it's the default way to write *any* new gothamtable, full stop. Check
% every new gothamtable against it before considering a slide done, not
% just when a rendered page happens to make the bleed visually obvious.
%
% BLOCK-RULE-COLOR CONVENTION (2026-07-16, direct instruction): a
% gothamtable placed *inside* any block environment (`infoblock`,
% `exampleblock`, `referenceblock`, `definitionblock`, plain `block`, or
% their `gothamfloating*` wrappers) should use black rules, not this
% environment's own default yellow -- yellow-on-the-block's-own-tinted-
% background reads poorly, unlike yellow-on-plain-white outside a block.
% Override locally, *inside* `\begin{gothamtable}{...}`, before `\toprule`
% (not before `\begin{gothamtable}` itself -- the environment's own
% opening code sets yellow *after* its `\begingroup`, so a color set
% before that point is simply overwritten; confirmed directly, first
% attempt this same day):
%   \begin{gothamtable}{...}
%   	\arrayrulecolor{c_primary_dark}
%   	\toprule
%   	...
%   \end{gothamtable}
% Booktabs rules read `\arrayrulecolor` fresh at each rule, and
% `\endgroup` (gothamtable's own closing code) resets it back to yellow
% afterward, so this never leaks into a later table. First applied on
% slides/6ux-Example-51-BMA-Payback-Table-current.tex's own
% `gothamfloatinginfoblock`.
\newenvironment{gothamtable}[1]{%
	\begingroup
	\arrayrulecolor{c_primary_highlight}%
	\begin{tabular}{#1}%
}{%
	\end{tabular}%
	\endgroup
}

% For a side-corner image overlay (the usual \vbox to 0pt{...} tikz overlay
% pattern used throughout this deck's converted slides) whose *top edge*
% should line up with the content text's own top edge -- distance from the
% page's own top edge to the cap-height top of a standardframe's first body
% line, measured directly against the compiled PDF's own text coordinates,
% not derived from \gothamContentTopShift or guessed. Previous per-slide
% image placements each guessed their own downward shift (e.g. a bare
% "-1.3cm") independently -- close, but not exact, landing the image's top
% a couple mm above the text's own top; use this constant instead of a
% fresh guess for any new "text top edge"-aligned image.
%
% 41.58pt (2026-07-22, re-measured): the previous value, 44.0pt, went stale
% the moment \gothamContentTopShift itself was increased (7.33pt -> 11.5964pt,
% "reduce the Top Margin", see that length's own comment above) without a
% matching update here -- that change pulls plain body text up, but this
% length is a separate, hand-measured constant with no formula tying it to
% \gothamContentTopShift, so it kept anchoring images at the *old*, lower
% position. Confirmed directly (pymupdf against this same freshly-compiled
% preview deck, not assumed from the \gothamContentTopShift delta alone):
% three independent plain-first-line standardframes with no per-slide
% top-shift correction (slides/5v4-CML-vs-SML-current.tex,
% slides/u5j-Exercise-CAPM-4-current.tex,
% slides/nt4-Exercise-CAPM-5-current.tex) all land their first line's own
% top edge at y=41.58pt post-change, not 44.0pt. Re-measure and adjust
% again if \gothamContentTopShift ever changes in the future.
\newlength{\gothamContentTopAlign}
\setlength{\gothamContentTopAlign}{41.58pt}

% For a "side filler image" -- a decorative image stuck flush to the page's
% own right edge (zero margin) with its top edge flush against the title
% bar's own bottom edge (zero gap), rather than sitting in the usual
% content-aligned corner inset -- distance from the page's own top edge
% down to the standardframe title bar's bottom edge, measured directly
% against the compiled PDF (a pixel scan for the yellow-to-white
% transition, at this deck's 16cm-wide/9cm-tall page size). First used on
% slides/dfc-Angel-Investors-current.tex.
%
% 25.43pt -- the raw measured value (originally 25.44pt; re-measured
% 2026-07-07 against the current compiled PDF via the title bar's own
% white-patch rectangle, which came back 25.427pt, close enough to the
% original measurement to be the same underlying edge, not a drift).
% Briefly set to 24.94pt instead (2026-07-06-ish), a deliberate 0.5pt
% overlap *into* the title bar, after a report of a visible hairline
% white gap right at this seam in one render -- not reproduced in this
% project's own pdftoppm-based rendering checks at the time (every pixel
% scan at the raw value showed a clean cut, no gap), so assumed to be a
% renderer-specific sub-pixel rounding artifact rather than a real error
% in the constant, and "fixed" by deliberately overlapping instead of
% chasing the specific viewer. That overlap turned out to be the wrong
% call: direct feedback 2026-07-07 reported the *opposite* symptom on
% `sideimageframe` (the image's own top edge visibly overlapping the
% title box) -- exactly what a real 0.5pt overlap looks like once it's
% actually noticed, rather than a rendering artifact. Reverted back to
% the precise raw measurement: a real, small, deliberate overlap chosen
% to paper over an unreproduced gap in an unspecified other renderer
% turns out to be a worse trade than just using the measured value and
% letting a genuine future gap report (if one ever recurs, in a
% specific, identified renderer) be diagnosed on its own terms.
%
% Re-measured 2026-07-16, 25.43pt -> 25.32pt: the frametitle hook's
% outer sep fix (see its own comment) moved the visible yellow bottom
% edge up by the 0.2pt the trim patch's fill had previously been inset
% below its gothamLogoAnchor.south anchor -- the yellow now genuinely
% ends flush with the logo's bottom edge, as always intended. New raw
% measurement against the rebuilt PDF: the trim patch's fill top edge
% (= the visible yellow bottom) sits 25.228 PDF units below the page
% top, i.e. 25.3226pt (earlier measurements in this comment quoted PDF
% coordinate values directly as "pt"; PDF units are big points, 1bp =
% 1.00375pt -- a systematic ~0.1pt understatement at this magnitude,
% harmless then, but corrected now that the constant is being reset
% anyway). Note this value is also what the overlay-logo nodes
% (titleframe/quoteframe/the hook's fixed-position redraw) derive their
% vertical inset from: visible-bar-bottom minus \gothamLogoHeight is
% exactly the inline logo's own vertically-centered top margin.
\newlength{\gothamTitleBarBottomAlign}
\setlength{\gothamTitleBarBottomAlign}{25.32pt}

% Note on the [t] default below: it top-aligns content within the frame
% (beamer's own default is vertical centering, "c"), so slides don't have
% their content floating in the middle. It also sidesteps a real beamer bug:
% \begin{frame}[]{Title} (a literal empty options bracket) fails outright,
% even in plain beamer with no theme at all -- so the optional argument must
% never default to an empty string. The title is always passed as an
% explicit argument (rather than left for beamer to pick up from a trailing
% brace group), since anything placed between \begin{frame} and its title --
% like our \setbeamertemplate calls -- breaks that trailing-brace mechanism.

% Pulls the frame's content up closer to the title bar: measured (via the
% PDF's own coordinates) at about 22pt of gap by default between the title
% bar's bottom edge and the first line of body content: this is a third of
% that, shrinking the gap to about two thirds of its original size.
%
% Called "Top Margin" in conversation from 2026-07-12 onward (direct
% instruction) -- the deck-wide name for the resulting gap between the
% title bar and a standardframe's first line of body content, i.e. the
% vertical position a plain-text frame's first baseline naturally lands at
% (measured directly against a reference slide's own compiled PDF, e.g.
% slides/aii-Leverage-Effects-current.tex's first line). Any first content
% that isn't plain text (a table, a minipage, an itemize with an unusual
% first item, ...) needs its own per-slide \gothamIncreaseContentTopShift/
% \gothamReduceContentTopShift correction to land on this same baseline --
% not portable between slides even for superficially similar content (see
% slides/zgz-EBIT-Definition-current.tex's own v1-2-0 comment for a case
% where an earlier, wrong correction was carried over uncritically).
\newlength{\gothamContentTopShift}
% 11.5964pt (2026-07-22, direct instruction: "reduce the Top Margin"
% generally, tested at 20%/30%/40% smaller against
% slides/8s8-Unit-Investment-Trusts-current.tex before settling on 40% --
% up from the original 7.33pt so the *resulting* gap (this length pulls
% content up by, so a bigger value means a smaller visible margin) comes
% out to 60% of its old size: measured directly on 8s8 at 10.67pt before,
% 6.42pt after (target was 6.40pt, within 0.02pt). Since every
% standardframe applies this same pull-up before any per-slide
% \gothamIncreaseContentTopShift/\gothamReduceContentTopShift correction,
% changing it here shifts the baseline project-wide, not just on 8s8.
\setlength{\gothamContentTopShift}{11.5964pt}

% "Top Margin" (above) is one of four deck-wide margin names introduced in
% conversation from 2026-07-12 onward (direct instruction) for the plain
% page-edge/title-bar/footer gaps every slide is expected to obey by
% default. The other three, for reference (each maps to a length already
% declared elsewhere, not a new one introduced here):
% - "Left Margin": the default gap from the page's own left edge to the
%   title text and body content -- \gothamFrametitleLeftpading
%   (beamerouterthemegotham.sty/beamerthemegotham.sty, 2.2ex, measured
%   directly against the compiled PDF at 12.552pt in this deck's usual
%   fonts -- see e.g. slides/tw5-Exercise-Leverage-and-ROE-Solution-current.tex's
%   own v1-1-0 comment for that same measured value).
% - "Right Margin": the default gap from the page's own right edge to the
%   institute logo (and, by the same convention, to anything else
%   deliberately flush-right, e.g. a corner-anchored image) --
%   \gothamFrametitleRightpading (0.28cm/7.93pt, this file's own
%   \gothamInstituteLogoSquare placement above uses it directly as the
%   logo's own inner sep).
% - "Bottom Margin": the default gap between a frame's own body content
%   and the \source{}/footnote line below it. Unlike the other three, this
%   one has no single settable length behind it -- \source{}'s own
%   vertical position isn't fixed the way the footer box is; it tracks the
%   frame's actual content height once that content gets tall enough (see
%   slides/u5h-Long-term-Bonds-Provisions-current.tex's own v1-3-0 comment,
%   "past a certain table height it starts tracking the table's own growth
%   roughly 1:1"). "Respecting Bottom Margin" in practice means leaving
%   comparable clearance to whatever a reference slide with unremarkable
%   content already shows, the same empirical, measure-against-a-reference-
%   slide approach "Top Margin" itself uses, rather than a fixed constant.

% Compensates a gap that opens up specifically when a standardframe's body
% starts with a \begin{minipage}{...} (used to width-limit text next to an
% image/diagram overlay, e.g. slides/igr-Budgeting-current.tex,
% slides/bdh-Working-Capital-Example-current.tex -- see README.md's
% "Available layouts" section for why minipage instead of columns there) --
% a minipage is a *box*, and TeX's vertical glue insertion (\topskip
% specifically) treats the first box in a vertical list differently from
% the first line of plain paragraph text, so the same \gothamContentTopShift
% pull-up above lands closer to the title for minipage-wrapped content than
% for plain content. Cannot be fixed by \vspace(*) right after
% \begin{minipage}{...}: both plain and starred forms were tried and
% measurably had zero effect (the star only protects glue from page-break
% discard, not from \topskip's own separate glue-insertion calculation,
% which absorbs leading glue at the start of a vertical list regardless).
% Fixed instead by shrinking standardframe's own pull-up for the affected
% frame -- see \gothamContentTopShiftReduction below. The value below is
% calibrated empirically, not derived from first principles: an isolated
% two-frame test (single-level itemize, one plain, one minipage-wrapped)
% measured only a 4.4pt gap, but real slides (nested itemize inside the
% minipage) showed a larger ~7.5-7.6pt gap against a real plain-itemize
% slide (slides/jyn-Categorization-current.tex) -- the extra few points
% likely come from itemize's own topsep/partopsep being computed differently
% depending on whether the list opens a box or a paragraph, on top of the
% topskip/box effect alone. Tuned so slides/igr-Budgeting-current.tex and
% slides/bdh-Working-Capital-Example-current.tex's first content line
% lands within 0.1pt of slides/jyn-Categorization-current.tex's own
% (200.917pt vs. 200.871pt/200.980pt in the compiled PDF).
\newlength{\gothamMinipageTopCorrection}
% 9.38pt (2026-07-22, re-measured): the previous 7.6pt value went stale
% the moment \gothamContentTopShift itself was increased (7.33pt ->
% 11.5964pt, "reduce the Top Margin", see that length's own comment above)
% without a matching update here -- confirmed directly (pymupdf against a
% freshly-compiled preview deck, not assumed from the \gothamContentTopShift
% delta alone, which doesn't translate 1:1 into a minipage's own \topskip
% behavior): four independent slides all using this constant completely
% unmodified (slides/fr3-Forms-of-Capital-Market-Efficiency-Strong-Semistrong-Weak-current.tex,
% slides/za2-Startup-Financing-Categorization-current.tex,
% slides/yke-Implications-of-Homemade-Dividends-current.tex,
% slides/6j4-Modigliani-Miller-Assumptions-current.tex) all land their
% own first line at an identical y=39.80pt post-change, 1.78pt short of
% this deck's current 41.58pt Top Margin reference -- not slide-specific
% noise, the same deterministic amount on every one of them. Re-measure
% and adjust again if \gothamContentTopShift ever changes in the future.
\setlength{\gothamMinipageTopCorrection}{9.38pt}

% A *third* minipage-top correction, on top of \gothamMinipageTopCorrection
% above rather than replacing it, for a minipage whose own first internal
% content is a zero-height `\vbox to 0pt{...}` tikz overlay (a
% corner-anchored decorative image, placed there specifically so it stays
% unconditionally visible even when the minipage's own later content
% includes a `\pause` -- see e.g.
% slides/y5z-CoCo-Bonds-current.tex's own "v1-2-0" note for the fuller
% pause-inheritance story behind why the image has to sit there at all).
% 2026-07-08, found while chasing "the top margin of the text on y5z/22s
% is too low" -- measured directly against the compiled PDF, y5h/22s both
% landed 16.5pt too far down despite already using
% \gothamMinipageTopCorrection, with the zero-height image overlay sitting
% *before* the minipage entirely as the actual root cause (that ordering
% makes the image, not the minipage, the frame body's true first box,
% which \topskip's own special first-box glue handling reacts to
% differently than a plain first minipage). Moving the image inside the
% minipage as its own first content closed most of the gap on its own
% (restoring the minipage as the frame's true first box) but left a
% further 10.35pt uncorrected -- confirmed the same on two independent
% slides (y5z, 22s) with identical structure, not assumed transferable
% from one to the other. Call via
% `\gothamReduceContentTopShift{\dimexpr\gothamMinipageTopCorrection-\gothamMinipageWithCornerImageExtraCorrection\relax}`
% for any slide combining this exact pattern (corner image as a minipage's
% own first content, ahead of a `\pause`) -- confirmed transferable as-is
% to two more slides with the same structure
% (slides/qbj-Bank-Loans-Credit-Lines-current.tex,
% slides/jkf-Bank-Loans-Term-Loan-current.tex), neither of which
% actually uses `\pause` itself, so this isn't specific to that -- the
% image-inside-minipage-as-first-content shape alone is what triggers it.
\newlength{\gothamMinipageWithCornerImageExtraCorrection}
\setlength{\gothamMinipageWithCornerImageExtraCorrection}{10.35pt}

% A *different* minipage-top correction, specifically for a minipage+
% itemize sitting as the first content of a `sideimageframe` (2026-07-06,
% found while chasing "the top margin of the text on ezz/acw is again
% off" -- see CONVERSION_NOTES.md's own "beamer/TeX gotchas" entry for the
% full story). `\gothamMinipageTopCorrection` above is calibrated for a
% *plain* standardframe, and simply does not transfer to `sideimageframe`:
% measured directly on three real sideimageframe slides
% (slides/xmp-Implications-of-Capital-Market-Efficiency-Considerations-current.tex,
% slides/dfc-Angel-Investors-current.tex,
% slides/ezz-Examples-of-Tests-for-Market-Efficiency-current.tex --
% one flat itemize, one flat itemize, one nested-first-item itemize
% respectively, so this isn't a flat-vs-nested effect either), all three
% landed within 0.6pt of each other at their own *uncorrected* natural
% position (52.03-52.57pt against this deck's usual 44.0pt first-line
% reference) regardless of which of the two *wrong* corrections
% (`\gothamReduceContentTopShift{\gothamMinipageTopCorrection}` or
% `\gothamIncreaseContentTopShift{\gothamFlatItemizeTopCorrection}`) had
% been applied to each beforehand -- i.e. `sideimageframe` itself pushes
% its own first-content minipage down by a consistent ~8.2pt relative to
% a plain standardframe, on top of whatever generic itemize-shape
% correction would otherwise apply, most likely from the extra
% \sidefootnotes/\settowidth setup work `sideimageframe`'s own opening
% code does before `\begin{standardframe}` ever runs. Use
% `\gothamIncreaseContentTopShift{\gothamSideImageFrameMinipageCorrection}`
% (note: *increase*, not reduce) as the *only* topshift call on any
% future sideimageframe slide whose first content is a minipage+itemize --
% not the generic `\gothamMinipageTopCorrection`/
% `\gothamFlatItemizeTopCorrection` this deck uses everywhere else.
\newlength{\gothamSideImageFrameMinipageCorrection}
\setlength{\gothamSideImageFrameMinipageCorrection}{8.2pt}

% 6.26pt (2026-07-22, re-measured): stale for the same reason
% \gothamMinipageTopCorrection above went stale -- \gothamContentTopShift
% itself changed (7.33pt -> 11.5964pt) without a matching update here.
% Re-measured directly against a freshly-compiled preview deck across six
% independent sideimageframe slides using this constant unmodified: four
% (ezz, xmp, hzx, u6n) land their own first line at an identical
% y=39.64pt, 1.94pt short of this deck's current 41.58pt Top Margin
% reference (i.e. already pulled up *too far*); the other two (c7w, dfc)
% are close but not identical (39.53pt/-2.05, 40.18pt/-1.4) -- no single
% value fixes all six exactly, so this recalibrates to the four-way
% cluster (the majority case) and accepts a small (<0.6pt) residual on the
% two outliers. Since this constant feeds
% `\gothamIncreaseContentTopShift{}` (more pull-up, i.e. *smaller* y), and
% the observed y was already short of target, the fix is a *decrease*
% (8.2pt - 1.94pt), not an increase -- caught by re-measuring after a
% first, wrong-direction attempt at 10.14pt drove the four-slide cluster
% to y=37.70pt, further from target instead of closer.
\setlength{\gothamSideImageFrameMinipageCorrection}{6.26pt}

% Opt-in override: reduces standardframe's \gothamContentTopShift pull-up by
% this amount for the frame currently being opened. Empty-default-style flag
% (mirrors \gothamChapterLogoPath etc.): 0pt by default (no effect), set via
% \gothamReduceContentTopShift{<length>} called *before* \begin{standardframe},
% reset back to 0pt in standardframe's own closing code so it never leaks
% into the next frame. Used by minipage-based slides to counteract
% \gothamMinipageTopCorrection above.
\newcommand{\gothamContentTopShiftReduction}{0pt}
\newcommand{\gothamReduceContentTopShift}[1]{\global\def\gothamContentTopShiftReduction{#1}}

% A *second*, separate first-line-position quirk, found and fixed
% 2026-07-05 after a run of slides (batch 4 plus several earlier ones) were
% found to land at two visibly different heights depending purely on
% content shape, with no minipage involved at all: a plain (non-minipage)
% standardframe whose top-level itemize's very *first* \item is immediately
% followed by its own nested \begin{itemize}/\begin{enumerate} lands 1.834pt
% higher than one whose first \item has no such immediate child (checked
% against every content slide converted so far via each compiled page's own
% first-line PDF baseline coordinate, not eyeballed -- 13 "nested-first-item"
% slides all landed within 0.05pt of each other, and 9 "flat-first-item"
% slides all landed within 0.05pt of each other, at a clean +1.834pt offset
% from the first group; isolated two-frame test cases, differing only in
% presence/absence of that immediate nested list, reproduced the exact same
% 1.834pt gap independent of which text/list type/item count was used).
% Not chased to one named TeX primitive (most likely some interaction of
% \addvspace's glue-coalescing behaviour with \partopsep/\topsep at the
% \begin{itemize}/\begin{frame} boundary, given the existing
% \gothamMinipageTopCorrection comment above already flags topsep/partopsep
% as varying with box vs. paragraph context) -- calibrated empirically like
% \gothamMinipageTopCorrection, not derived from first principles.
% Interacts with minipage differently, not additively: inside a minipage,
% this nested-vs-flat distinction disappears entirely (both land within
% 0.05pt of each other) -- so this correction is for *plain* itemize-first
% standardframes only, never combined with \gothamMinipageTopCorrection.
\newlength{\gothamFlatItemizeTopCorrection}
\setlength{\gothamFlatItemizeTopCorrection}{1.834pt}

% 0.01pt (2026-07-22, re-measured): stale for the same reason
% \gothamMinipageTopCorrection above went stale -- \gothamContentTopShift
% itself changed (7.33pt -> 11.5964pt) without a matching update here.
% Re-measured directly against three independent flat-itemize slides using
% this constant unmodified (574, 5cr, c8e): all three landed their own
% first line at an identical y=39.76pt, 1.82pt short of this deck's
% current 41.58pt Top Margin reference -- coincidentally, the recalibrated
% \gothamContentTopAlign reference now already accounts for almost all of
% what this correction used to supply on its own, leaving next to nothing
% left to add. Effectively zeroed out rather than removing the call
% itself, so every flat-itemize slide's own
% `\gothamIncreaseContentTopShift{\gothamFlatItemizeTopCorrection}` line
% keeps working unchanged.
\setlength{\gothamFlatItemizeTopCorrection}{0.01pt}

% Opposite-direction counterpart to \gothamReduceContentTopShift: *increases*
% the pull-up instead of reducing it (needed for \gothamFlatItemizeTopCorrection
% above, since a flat first item needs *more* pull-up, not less). Same
% storage variable, just negated -- \gothamContentTopShiftReduction is
% allowed to go negative, since the frame-opening code only ever computes
% \gothamContentTopShift-\gothamContentTopShiftReduction, never checks its
% sign.
\newcommand{\gothamIncreaseContentTopShift}[1]{\global\def\gothamContentTopShiftReduction{-#1}}

% Define custom frame type with en-dash bullets
\newenvironment{standardframe}[2][t]{%
	% Restart footnote numbering (the *, dagger, ... symbols) fresh on every
	% slide, rather than letting it run on cumulatively across the whole deck
	% the way beamer does by default. Two reasons: \fnsymbol only defines 9
	% symbols before erroring out, which a real multi-slide deck would run
	% through in no time if every \footnote{}/\source{} call across the whole
	% talk drew from one shared counter; and each slide's footnotes are read
	% and referenced only on that one slide anyway, so there's no reason for
	% their numbering to depend on how many footnotes earlier slides happened
	% to have. This must run *before* \begin{frame}: beamer saves the
	% counter's current value as that frame's baseline as part of
	% \begin{frame}'s own internal setup (restoring it before every overlay
	% replay, so overlays within one frame don't double-count) -- resetting
	% after \begin{frame} would be reset too late, after that baseline was
	% already captured.
	\setcounter{footnote}{0}%
	\begin{frame}[#1]{#2}%
	% Forces this frame's own body to start with the correct default text
	% color, regardless of what any *previous* frame left active
	% (2026-07-16, direct instruction, root-caused via an isolated minimal
	% reproduction). \titleframe's/\quoteframe's own
	% \color{\gothamTitleTextColor} calls -- set inside content built via a
	% stored box register or a zero-space overlay, used/shipped out *later*
	% rather than immediately -- don't reliably get "undone": xcolor's
	% \color doesn't emit a PDF color operator immediately when called, only
	% when the *next* real glyph/path actually needs it, so whatever color
	% one of those frame types last set can end up as the active PDF color
	% well into the *next* frame's own first draw call. Invisible for Graz
	% (the leaked color, c_primary_dark, happens to match ordinary body text
	% anyway) but catastrophic for MCI/Innsbruck (c_primary_white, rendering
	% plain non-itemize body text invisible against its own white
	% background -- itemize items escape it, since beamer's own itemize
	% template sets color fresh per bullet). An \AtBeginEnvironment{frame}
	% hook (etoolbox) was tried first and confirmed *not* to fire at all --
	% beamer's own "frame" isn't a plain \newenvironment etoolbox can patch
	% that way -- so this resets directly in standardframe's own opening
	% code instead, exactly where this project already has full control.
	\color{c_primary_dark}%
	% Defensive: the tokens after \begin{frame} are collected into the frame
	% body and replayed once per overlay slide, so this clears the side-
	% footnote collector at the start of every slide's pass -- normally a
	% no-op (the previous slide's footline already consumed the box), but it
	% stops stale footnotes leaking in if a previous frame ever collected
	% without shipping a footline.
	\gothamClearSideFootnotesBox%
	\setbeamertemplate{itemize item}{\gothamBulletDot}%
	\setbeamertemplate{itemize subitem}{\gothamBulletDot}%
	\setbeamertemplate{itemize subsubitem}{\gothamBulletDot}%
	\setbeamertemplate{itemize subsubsubitem}{\gothamBulletDot}% 4th level, see its own comment further below (search "subsubsub")
	% Only for top-aligned frames (the default, #1=t): a [c]-centered frame
	% (2026-07-22, direct instruction, found while centering
	% slides/bw9-Pearson-Correlation-Illustration-current.tex/agc/crj) was
	% landing its content almost flush against the title bar instead of
	% truly centered -- root cause was this same pull-up firing
	% unconditionally regardless of #1. beamer's own [c] option already
	% centers the frame body within the available space; stacking this
	% deck's *own* top-alignment pull-up on top of that just drags the
	% already-centered box further upward by the same fixed amount, which
	% is exactly the asymmetry that was measured (a [c] figureframe's image
	% landing ~0pt from the title bar and ~10pt from the footer instead of
	% symmetric). Guarded with \ifstrequal (etoolbox, already loaded above)
	% rather than \ifx, matching this file's own \ifstrequal{##1}{...}
	% pattern elsewhere (\ifx would need #1 fully expanded first).
	\ifstrequal{#1}{c}{}{\vspace*{-\dimexpr\gothamContentTopShift-\gothamContentTopShiftReduction\relax}}%
}{%
	\end{frame}%
	\global\def\gothamChapterLogoPath{}% opt-in only for the frame just closed
	\global\def\gothamContentTopShiftReduction{0pt}% opt-in only for the frame just closed
	\global\gothamSideFootnotesfalse% opt-in only for the frame just closed
	\global\def\gothamFrametitleFullHeightYellow{}% opt-in only for the frame just closed
	\global\def\gothamFrametitleLogoFixed{}% opt-in only for the frame just closed
	\global\def\gothamFooterBoxColor{\gothamFooterBoxColorDefault}% opt-in only for the frame just closed
}

% The bullet itself: a small filled circle, matching the source PowerPoint
% deck's own round bullets (a plain \textbullet, tried first, is too
% small/thin in Fira Sans at any standard text size to read as a solid dot
% -- a tikz-drawn circle gives full control over the actual rendered size
% instead of depending on the font's own glyph design). baseline=-0.65ex
% roughly centers it on the surrounding text's x-height, the same vertical
% position a font's own bullet glyph would normally sit at.
%
% Its fill color is \gothamBulletColor, not a literal color token here --
% set by the per-university preamble before \input{}ing this file, same
% "value supplied externally, referenced lazily" idiom as
% \gothamInstituteLogoPath/\gothamLogoHeight above (\gothamBulletDot's own
% body only expands \gothamBulletColor when a bullet actually gets drawn,
% long after this \newcommand runs). Graz's own bullets are its primary
% highlight color (matching this project's original, single-institution
% behavior); MCI's and Innsbruck's are their own secondary blue instead
% (2026-07-16, direct instruction) -- see each institution's own
% Slidemaster_preamble_*.tex for its \gothamBulletColor setting.
\newcommand{\gothamBulletDot}{\tikz[baseline=-0.65ex]{\fill[\gothamBulletColor] circle (1.7pt);}}

% "a)"/"b)"/"c)"-labeled list, for exercise sub-tasks that read better
% lettered than bulleted (2026-07-20, direct instruction, first used on
% slides/y3a-Exercise-Diversification-2-current.tex's own bottom three
% items).
%
% First attempt (superseded, kept as a documented wrong turn -- see
% README.md's own "lualatex vs latex" section for why this project keeps
% wrong write-ups rather than deleting them): reuse itemize's own top-level
% `\llap`-based label positioning by locally overriding
% `\setbeamertemplate{itemize item}` inside a plain `\begin{itemize}` --
% the reasoning being that `\llap` draws its content flush against the
% text's own left edge "regardless of the label's width." That reasoning
% turned out to be wrong in a way that only showed up once actually
% measured against the compiled PDF (not just visually eyeballed): a plain
% bullet's own `\llap`-drawn footprint (title/Left-Margin-relative-left-
% edge to text-start) measured only ~9.05pt end-to-end on this deck's real
% slides, matching the Left-Margin-to-text-start gap exactly -- but "a)"
% rendered through the *exact same* mechanism measured a ~13.97pt
% footprint for the same gap, not the ~8.5-9pt `\settowidth` reported for
% the bare glyphs in isolation. The extra ~5pt was never fully root-caused
% (this custom itemize's own `\list` call, further below, only overrides
% `\makelabel` -- it never sets `\leftmargin`/`\labelwidth` itself, so
% those are left at whatever this deck's global itemize defaults happen to
% be, and `\llap`'s actual box-width behavior inside that inherited,
% unexamined list-parameter state is where the extra width most likely
% comes from) -- rather than keep guessing at *why*, measured the actual
% failure directly (the compiled "a)" sat ~4.9pt left of the Left Margin,
% confirmed via exact PDF text-position extraction, not a visual guess)
% and abandoned the approach once even `\tiny` plus manual per-letter
% kerning only just barely closed the gap while reading as illegibly small
% -- a real dead end, not a tuning problem.
%
% Working fix: don't reuse itemize's `\llap` label mechanism at all -- use
% a real, width-reserving `\makebox[W][r]{...}` label instead (the same
% *kind* of mechanism `\gothamEnumerateLabelWidth` above already uses for
% a nested enumerate, just with an explicit `\leftmargin`/`\labelwidth`/
% `\labelsep` this environment sets itself rather than inheriting), inside
% its own plain `\list`, not a reused `\itemize`. `\leftmargin` (label box
% + gap) is set to 11.1pt -- 9.1pt label box (fits "a)"/"b)"/"c)" at this
% deck's normal `\small` itemize size with a little clearance, confirmed
% directly: measured label-left-edge stays right of the 12.552pt Left
% Margin for all three letters, closest case (widest letter, "c)") 1.37pt
% clear) plus 2pt `\labelsep` for a visible gap before the text. Trade-off,
% accepted deliberately rather than found by accident: item text lands
% ~2.2pt right of where a plain bulleted item's text would (23.8pt vs.
% 21.6pt) -- fitting a legible "a)"/"b)"/"c)" at normal reading size plus
% a real gap simply doesn't fit in the ~9pt a plain bullet's tiny dot
% leaves, so *exact* text-position parity with plain bullets and a legible
% label are mutually exclusive here; legibility won given the two-point
% shift is visually negligible and there's no plain-bulleted item left on
% the same slide to compare against directly anyway (y3a's own first item
% moved into its infoblock). Re-derive `\leftmargin`/`\labelwidth` the same
% way (render, measure with exact text-position extraction, adjust) if
% this is ever used at a different font size than this deck's normal
% itemize `\small`.
\newcounter{gothamalphaitem}
\newcommand{\gothamAlphaItem}{\stepcounter{gothamalphaitem}\alph{gothamalphaitem})}
\newenvironment{gothamalphaitemize}{%
	\setcounter{gothamalphaitem}{0}%
	\list{\gothamAlphaItem}{%
		% 4.75pt (2026-07-20, direct instruction: "half the width of the
		% 'C' in 'Calculate'" -- \settowidth{C} at this deck's normal
		% itemize \small measured 5.5pt, so 2pt original + 2.75pt half-C).
		% Only \labelsep changes the text's own start position -- the
		% label's own box stays exactly where it was (its left edge is
		% pinned to this list's origin regardless of \labelsep, since
		% \leftmargin=\labelwidth+\labelsep always cancels \labelsep back
		% out of the label box's own position), so this widens the gap
		% without reopening the Left-Margin question already solved above.
		\setlength{\labelsep}{4.75pt}%
		\setlength{\labelwidth}{9.1pt}%
		\setlength{\leftmargin}{\labelwidth}%
		\addtolength{\leftmargin}{\labelsep}%
		\setlength{\itemindent}{0pt}%
		% \itemsep/\parsep left at plain \list's own defaults otherwise --
		% since this environment is built directly on \list rather than
		% reusing beamer's own itemize template, it never got beamer's own
		% itemize spacing hooks at all, only whatever \list falls back to.
		% Measured directly (2026-07-21, direct instruction, found on
		% slides/nt4-Exercise-CAPM-5-current.tex): that default reads as
		% *less* space between two items (13.7pt, item-to-item) than
		% between two wrapped lines within one *single* item (15.6pt,
		% pure \baselineskip) -- backwards from how a list should read
		% (item boundaries should be at least as open as a same-item line
		% wrap, not tighter). \itemsep set explicitly so the total gap at
		% an item boundary (baselineskip + itemsep) reads as clearly more
		% open than a plain wrapped line, not accidentally tighter.
		% Reduced from the original 6pt to 2.5pt (2026-07-22, found during
		% a full-project check after the itemize-rules work, not part of
		% that work itself): 6pt overshot by more than needed to clear the
		% 1.9pt gap (15.6 vs. 13.7), and slides/kmf-Exercise-CAPM-6-current.tex's
		% own 5-item list -- longer than nt4's own -- compounded that
		% excess across 4 item boundaries into a ~25pt overfull frame, its
		% own last item cut off by the footer (confirmed directly,
		% rendered and compared against a pre-regression PDF). 2.5pt
		% keeps a real, visible margin over the 1.9pt needed (not just
		% barely crossing it) while re-measuring clean (no overfull) on
		% kmf; re-verify against both nt4 and kmf directly if this is
		% ever changed again -- they're this environment's own tightest
		% (most items) and originally-tuned-against cases respectively.
		\setlength{\itemsep}{2.5pt}%
	}
	\raggedright% plain \list defaults to justified+hyphenated text, unlike
	% this deck's own itemize (which calls this too, further below) --
	% missed on the first pass, only caught by comparing a rendered a)/b)/c)
	% item side by side against a plain bulleted one ("stan-\ndard deviation"
	% hyphenating where a plain-bulleted version of the same sentence
	% elsewhere in this deck doesn't).
}{%
	\endlist
}

% Itemize indentation: level 1 bullets should sit flush with the frame title
% (no extra left margin), and each deeper level should be indented just
% enough for its bullet to line up under the *text* (not the bullet) of the
% level above. A nested list's \leftmarginN is added on top of its parent's
% (indentation is cumulative), so giving every level the same "step" --
% label width + labelsep, the width a bullet plus its gap to text takes up --
% makes level N's text sit at N*step and level N's bullet at (N-1)*step, i.e.
% exactly at level (N-1)'s text (0 for level 1, flush with the title).
\makeatletter
\newlength{\gothamBulletWidth}
\settowidth{\gothamBulletWidth}{\gothamBulletDot}
\newlength{\gothamItemStep}
\setlength{\gothamItemStep}{\labelsep}
\addtolength{\gothamItemStep}{\gothamBulletWidth}
\setlength{\leftmargini}{\gothamItemStep}
\setlength{\leftmarginii}{\gothamItemStep}
\setlength{\leftmarginiii}{\gothamItemStep}
\setlength{\leftmarginiv}{\gothamItemStep}
\makeatother

% A 4th itemize level -- beamer's own \itemize (beamerbaselocalstructure.sty)
% hard-caps nesting at 3 levels (\ifnum \@itemdepth >2\relax\@toodeep),
% independently of anything this deck sets: a 4th \begin{itemize} always
% raised "! LaTeX Error: Too deeply nested." even though plain LaTeX's own
% kernel itemize supports 4. Found 2026-07-15, direct instruction, on
% slides/hrg-Exchange-Rate-Quotation-Forward-current.tex -- genuinely
% 4 levels deep in the source pptx itself (confirmed via a PowerPoint COM
% render, not assumed), not a conversion artifact to design around.
% Since this deck already styles every existing level identically (same
% \gothamBulletDot, same \gothamItemStep margin -- see the comments right
% above/below), extending to a 4th level needs no new visual design, just
% patching beamer's own depth limit up by one and telling it what to call
% the new level ("subsubsub", the next name in beamer's own
% sub/subsub/... naming scheme) so its template/color lookups
% (\usebeamertemplate{itemize subsubsubitem}, \usebeamercolor{itemize/
% enumerate subsubsub body}, etc.) resolve to *something* instead of
% erroring -- undefined beamer templates/colors are expected to fall back
% to sensible defaults (no extra spacing hook, plain text color) rather
% than raising an error, same as any other beamer template/color that a
% given theme never bothers to set; the actual bullet symbol still needs
% its own \gothamBulletDot template below, added explicitly the same way
% as the existing 3 levels, not left to whatever fallback beamer has for
% an unset "itemize subsubsubitem" template.
% \beamer@computepref is patched wholesale (not via \patchcmd) since its
% body is one single \ifcase chain, not a stable string to match against.
% \itemize itself was *first* tried via \patchcmd{\itemize}{...}
% (matching this deck's own established practice for beamer/theme
% internals, e.g. beamerfontthemegotham.sty's own
% \patchcmd{\beamer@title}{...}) -- this failed outright, confirmed
% directly (\patchcmd's own failure branch fired, raising a clean
% "Patching failed" error rather than silently doing nothing): recent
% LaTeX2e kernels implement an optional-argument environment like
% beamer's own `\renewcommand{\itemize}[1][]{...}` via
% `\@protected@testopt`, which stores the *real* body not under `\itemize`
% itself but under a differently-named internal control sequence (`\show
% \itemize` confirms this directly: it prints as
% `\@protected@testopt \itemize \\itemize {}`, i.e. plain `\itemize` is
% only a thin dispatcher, and the actual `\ifnum \@itemdepth >2...` body
% \patchcmd was trying to match against was never there to find). Rather
% than chase that internal storage name (liable to differ across LaTeX
% kernel versions, which \patchcmd's own documented usage doesn't need to
% care about), \itemize is fully re-declared here instead, copied
% verbatim from beamer's own beamerbaselocalstructure.sty (this project's
% installed version) with only the depth bound itself changed from 2 to
% 3 -- the same declaration form beamer itself used
% (\renewcommand{\itemize}[1][]{...}), which regenerates both the
% dispatcher and the real body correctly under whatever internal name
% the current kernel expects, sidestepping the naming question entirely.
% Trade-off, stated plainly: if a future beamer release changes this
% macro's own body for an unrelated reason, this override will silently
% keep using today's (2026-07) copy instead of picking up that change --
% acceptable here since the override is one line different from upstream
% and clearly marked, but re-diff against beamerbaselocalstructure.sty's
% own \itemize if beamer is ever upgraded and itemize behavior looks off.
\makeatletter
\long\def\beamer@computepref#1{%
	\let\beameritemnestingprefix\@empty%
	\ifcase#1\or\or\def\beameritemnestingprefix{sub}\or\def\beameritemnestingprefix{subsub}\or\def\beameritemnestingprefix{subsubsub}\or\@toodeep\fi%
}
\renewcommand{\itemize}[1][]{%
	\ifblank{#1}{}{\def\beamer@defaultospec{#1}}%
	\ifnum \@itemdepth >3\relax\@toodeep\else
		\advance\@itemdepth\@ne
		\beamer@computepref\@itemdepth% sets \beameritemnestingprefix
		\usebeamerfont{itemize/enumerate \beameritemnestingprefix body}%
		\usebeamercolor[fg]{itemize/enumerate \beameritemnestingprefix body}%
		\usebeamertemplate{itemize/enumerate \beameritemnestingprefix body begin}%
		\list
			{\usebeamertemplate{itemize \beameritemnestingprefix item}}
			{\def\makelabel##1{%
					{%
						\hss\llap{{%
								\usebeamerfont*{itemize \beameritemnestingprefix item}%
								\usebeamercolor[fg]{itemize \beameritemnestingprefix item}##1}}%
					}%
				}%
			}
	\fi%
	\beamer@cramped%
	% \beamer@cramped (beamerbaselocalstructure.sty, unmodified upstream
	% code) unconditionally does `\itemsep0pt\topsep0.5pt`, every single
	% time -- called right here, in this project's own copy of beamer's
	% own \itemize body, same position as upstream. Setting \itemsep any
	% earlier (first tried as a separate \AtBeginEnvironment{itemize} hook,
	% run *before* this whole macro body even starts) gets silently wiped
	% out by this same line a moment later -- confirmed directly via
	% \typeout debugging (not assumed): the hook's own \setlength really
	% did run with the right values, \itemsep still read back as 0pt
	% immediately after. Setting it *here*, right after \beamer@cramped
	% already ran (not before), is what actually survives.
	\ifcase\@itemdepth
	\or\setlength{\itemsep}{\gothamItemsepRatioI em}% level 1
	\or\setlength{\itemsep}{\gothamItemsepRatioII em}% level 2
	\or\setlength{\itemsep}{\gothamItemsepRatioIII em}% level 3
	\or\setlength{\itemsep}{\gothamItemsepRatioIV em}% level 4
	\fi%
	\raggedright%
	\beamer@firstlineitemizeunskip%
}
\makeatother

% A nested \begin{enumerate} inherits the same \leftmarginN step as itemize
% above for *where its text starts* (matching itemize's own "one step past
% the parent level's text" rule -- confirmed directly against the compiled
% PDF: an enumerate's own body text, e.g. "Senior secured debt", already
% lands exactly one \gothamItemStep past its parent item's text), but its
% *label* ("1.", "2.", ...) doesn't follow \labelwidth the way plain LaTeX's
% list mechanism would -- beamer draws enumerate labels through its own
% template machinery, which turned out not to respect \labelwidth at all:
% \setlength{\labelwidth}{...} right after \begin{enumerate} measurably took
% effect (checked with \typeout{\the\labelwidth} immediately after) but the
% compiled numeral's own position never moved. What *does* move it: directly
% overriding \makelabel with a fixed-width \makebox -- empirically
% calibrated (render, measure the numeral's own x-position in the compiled
% PDF, adjust, repeat) since the box's own left edge turned out to sit at a
% fixed position unrelated to \labelwidth or \leftmarginN, with the
% right-aligned numeral inside it moving in lockstep with the box's width
% (confirmed linear: +1pt width -> +~1pt numeral position, across several
% test widths). 19.23pt (at this deck's default itemize font size) lands a
% single-digit numeral flush with the parent item's own text -- only valid
% at that font size and for single-digit lists (1-9); re-calibrate the same
% way (not by recomputing from first principles) for a differently-sized or
% 10+-item nested enumerate. Call right after \begin{enumerate}, before the
% first \item.
%
% 20.33pt (2026-07-08, direct instruction: "the numbers still extend a
% little farther to the left than the text above them does" --
% slides/f6m-Recovery-Rates-current.tex's own "Pecking order:" list).
% 19.23pt landed *close* but not quite flush: measured directly against
% the compiled PDF, every numeral's own left edge (20.55pt-20.95pt,
% varying slightly by digit width even though all are right-aligned to
% the same box) still sat up to ~1.05pt left of "Pecking order:"'s own
% text start (21.6pt) -- the widest digit tried, "6.", was the binding
% case. +1.1pt of box width (confirmed +1pt width -> +1pt numeral
% position, per this comment's own note above) covers that with a small
% margin, re-confirmed against the recompiled PDF.
\newcommand{\gothamEnumerateLabelWidth}{\renewcommand{\makelabel}[1]{\makebox[20.33pt][r]{##1}}}

% Automatic version of the above (2026-07-23, background-task follow-up
% flagged in slides/kfh-Asset-Allocation-Process-current.tex's own v1-1-0
% comment: "fix the enumerate behavior generally, if possible", task
% `task_bb3fbb66`). That comment's first attempt, an
% `\AtBeginEnvironment{enumerate}{...}` hook, failed because the hook fires
% *before* `\enumerate`'s own body runs, and that body's own internal
% `\list{...}{\def\makelabel...}` setup (the beamer copy directly below the
% `\itemize` override above) clobbers whatever a hook that early leaves
% behind -- the exact same "template runs its own setting, external pre-set
% gets overwritten" problem `\itemize`'s own `\itemsep` comment above
% already documents. The real fix, same idea `\itemize` itself used above:
% override beamer's own low-level piece directly rather than hooking
% \AtBeginEnvironment. Unlike `\itemize` (declared via
% `\@protected@testopt`, needing a full re-declaration to reach -- see the
% `\itemize` override's own comment for the full story), `\enumerate`'s
% relevant piece, `\beamer@enum@` (beamerbaselocalstructure.sty), is a
% plain `\def` with no optional-argument indirection, so it can be
% `\renewcommand`'d directly, verbatim from upstream except for the
% `\ifnum\@itemdepth>1` branch added below.
%
% Only fires for a *nested* enumerate (`\@itemdepth>1`, i.e. already inside
% at least one `\itemize`/`\enumerate` when this one opens) -- a top-level
% enumerate (`\@itemdepth=1`, e.g.
% slides/t9z-Outland-Steel-Hedging-Borrowing-Investing-current.tex/
% slides/g4c-Outland-Steel-Hedging-Forward-Sale-current.tex/
% slides/j84-Outland-Steel-No-Hedging-current.tex's own outermost numbered
% list, or slides/nkj-Agenda-current.tex's own \tableofcontents-driven one)
% is left at beamer's stock label position entirely untouched, matching
% this project's own established behavior for that shape (never reported
% as misaligned, unlike the nested case) -- and, for nkj specifically,
% doesn't matter either way, since that slide's own
% `\gothamTocLabelFlushLeft` call (right after its own `\begin{enumerate}`)
% overrides `\makelabel` again regardless of whatever this default sets.
%
% Reuses the exact same 20.33pt calibration as the manual
% `\gothamEnumerateLabelWidth` above, confirmed (pymupdf, compiled PDF) to
% land correctly not just at `\@itemdepth=2` (one `\itemize` above the
% enumerate -- slides/kfh-Asset-Allocation-Process-current.tex/
% slides/f6m-Recovery-Rates-current.tex/
% slides/6pz-Outland-Steel-Hedging-Strategies-Overview-current.tex, all
% previously needing the manual call) but also at `\@itemdepth=3` (two
% `\itemize`s above it -- slides/bdh-Working-Capital-Example-current.tex's
% own "Development of balance sheet:" list, also previously needing the
% manual call, with the identical 20.33pt value) -- this deck's itemize/
% enumerate templates don't shrink font size by nesting depth (only the
% bullet glyph/margin step does), so the numeral's own glyph width -- the
% only thing \gothamEnumerateLabelWidth's own box width actually
% compensates for -- stays the same at every depth, and one constant covers
% both cases without needing a depth-dependent formula. Not verified beyond
% depth 3 (no slide currently nests an enumerate that deep); if one ever
% does, re-measure before trusting this constant that far down, same as any
% other empirical calibration in this file.
% (slides/scs-Political-Risk-Stability-Index-current.tex also matches
% "`\gothamEnumerateLabelWidth`" in a grep -- but only inside its own prose
% comments describing the general technique; its actual call is the
% bespoke `\gothamEnumerateLabelFlushLeftShifted{16pt}{10.2pt}` for its own
% alphabetic, left-aligned, narrower-box labels, unaffected by this
% override either way since that call runs after `\begin{enumerate}` opens,
% same ordering as every other post-open `\makelabel` override here.)
%
% Former per-slide manual `\gothamEnumerateLabelWidth` calls removed from
% the 4 slides above now that this covers them automatically; the macro
% itself is kept defined, both for this override to reuse and as a manual
% escape hatch (e.g. a differently-calibrated width some future slide might
% need).
\makeatletter
\renewcommand{\beamer@enum@}{%
	\beamer@computepref\@itemdepth% sets \beameritemnestingprefix
	\usebeamerfont{itemize/enumerate \beameritemnestingprefix body}%
	\usebeamercolor[fg]{itemize/enumerate \beameritemnestingprefix body}%
	\usebeamertemplate{itemize/enumerate \beameritemnestingprefix body begin}%
	\expandafter
		\list
			{\usebeamertemplate{\beamer@enumtempl}}
			{\usecounter\@enumctr%
				\ifnum\@itemdepth>1
					\def\makelabel##1{\makebox[20.33pt][r]{##1}}%
				\else
					\def\makelabel##1{{\hss\llap{{%
							\usebeamerfont*{enumerate \beameritemnestingprefix item}%
							\usebeamercolor[fg]{enumerate \beameritemnestingprefix item}##1}}}}%
				\fi}%
	\beamer@cramped%
	\raggedright%
	\beamer@firstlineitemizeunskip%
}
\makeatother

% Extra breathing room when a nested itemize closes and control returns to
% its *parent* list, right before that parent's next \item -- e.g. sub-items
% under "Mid-Term Exam" followed by the next top-level "Final Exam" -- so a
% new top-level group reads as visually distinct from the one before it,
% rather than the same uniform gap used between every other item regardless
% of level. \@itemdepth (LaTeX's own nesting-level counter, already
% decremented back to the parent's level by the time \end{itemize}'s normal
% closing has run) says which level control returned to -- 0 means the
% outermost list itself just closed (nothing left to add space before,
% deliberately excluded), 1 means back to a level-1 item, 2 means back to a
% level-2 item.
%
% Level-dependent (2026-07-06, direct instruction) -- a flat, single
% constant regardless of *which* level control returned to (the original
% version of this hook) reads as "extra"/inconsistent precisely because a
% return-to-level-1 and a return-to-level-2 transition both got the exact
% same gap, rather than level 1 getting a visibly bigger gap than level 2 --
% found on slides/57r-Reasons-for-Leasing-current.tex, where a level-2
% sub-list closing back to a level-1 item looked like a random extra gap
% rather than "level 1 items are consistently more separated than level 2
% ones." \addvspace (not plain \vspace) matters too: it coalesces with
% whatever glue itemize's own \itemsep/\topsep already contributes at the
% closing boundary (taking whichever is larger) rather than *adding* on top
% of it unconditionally -- \vspace's own stacking behavior was the other
% half of why the old flat version could look like too much space.
\makeatletter
% Ratios, not lengths (2026-07-21, direct instruction: "all of these
% spacings should scale with font size") -- a \newlength+\setlength{...}{0.5em}
% pair, the way this used to be written, freezes to a fixed *pt* value the
% one time \setlength runs (at preamble-load time, using whatever font is
% active then), despite the "em" in its own definition -- it does NOT
% re-track whichever font size is actually active on a given slide/level
% later. Confirmed directly: a \typeout{\f@size} placed at the point each
% of these actually gets used (inside \AtEndEnvironment{itemize}'s own
% hook body, further below) showed the *closing* list's own font still
% active there (9pt/10pt/10.95pt for a level-4/3/2 list closing,
% respectively -- not yet restored to the parent's own font), confirming
% this hook fires early enough in `\end{itemize}`'s own processing that a
% live "0.5em" computed *there* would correctly track whichever level is
% actually involved. Fixed by storing the bare ratio as a plain macro
% instead, and writing the literal "em" back on at the point of use
% (`\addvspace{\gothamItemizeLevelDropSpaceRatioI em}`), so it's
% recomputed fresh, against whatever font is live, every time.
\newcommand{\gothamItemizeLevelDropSpaceRatioI}{0.5}
\newcommand{\gothamItemizeLevelDropSpaceRatioII}{0.3}
% Level 3's own drop-space, for when a closed *4th*-level list returns
% control to a level-3 item -- added 2026-07-15 alongside the 4th itemize
% level itself (search "subsubsub" further above). Interpolated from the
% existing I/II progression (0.5, 0.3 -- each step roughly 60% of the one
% before) rather than empirically re-measured against a compiled PDF the
% way I/II themselves were: no slide used a 4th level before
% slides/hrg-Exchange-Rate-Quotation-Forward-current.tex to measure
% against. Re-derive properly (render, measure, adjust) if this ever
% looks visibly off on a real slide.
\newcommand{\gothamItemizeLevelDropSpaceRatioIII}{0.2}
% Level 4's own drop-space, for when a closed level-4 list returns control
% to a level-3 item -- added 2026-07-21 alongside \gothamItemsepRatioIV.
% \@itemdepth at this hook's fire time is the depth of the list that *just
% closed*, confirmed directly via a temporary \typeout{\the\@itemdepth}
% (not assumed): closing the level-4 list on
% slides/ZZO-Itemize-Spacing-Reference-current.tex measured itemdepth=4
% there, closing level-3 measured 3, level-2 measured 2, and the outermost
% level-1 list itself closing measured 1 (never 0 -- \@itemdepth's own
% restore to the parent's depth happens *after* this hook runs, so the "0:
% outermost list itself closed" case below is dead code, kept only because
% removing a genuinely unreachable \ifcase branch isn't worth the risk of
% getting the numbering wrong a second time). So this is very much a real,
% reachable case -- not decorative symmetry -- and needed its own
% \addvspace branch (previously missing entirely, silently leaving the
% level-4-to-level-3 closing transition with no drop-space at all).
% Interpolated from the existing I/II/III progression as a starting guess,
% then re-measured directly against the compiled ZZO reference and bumped:
% 0.13 measured at 12.6pt, still *less* than level 4's own same-level gap
% of 13.4pt (violates rule III -- a level's own same-level gap must be
% smaller than that level's own closing-transition gap); raised until
% re-measurement cleared it with a comfortable margin again, matching the
% margin every other level's own pair already has.
\newcommand{\gothamItemizeLevelDropSpaceRatioIV}{0.5}
\AtEndEnvironment{itemize}{%
	\ifcase\@itemdepth
		% 0: dead code -- @itemdepth is never actually 0 when this hook
		% fires (see \gothamItemizeLevelDropSpaceRatioIV's own comment
		% just above); kept only for alignment with the \ifcase branches
		% below, each keyed on the depth of the list that just closed.
	\or\addvspace{\gothamItemizeLevelDropSpaceRatioI em}% level-1 (outermost) list itself closed -- a pre-existing quirk (not introduced by this comment's own edit): there's no parent item above level 1 for this to visually separate from, so this only actually matters if something follows the outermost list directly on the same frame; harmless on every slide seen so far since nothing does. Out of scope to redesign here.
	\or\addvspace{\gothamItemizeLevelDropSpaceRatioII em}% level-2 list closed, back to a level-1 item
	\or\addvspace{\gothamItemizeLevelDropSpaceRatioIII em}% level-3 list closed, back to a level-2 item
	\or\addvspace{\gothamItemizeLevelDropSpaceRatioIV em}% level-4 list closed, back to a level-3 item
	\fi
}

% Per-level \itemsep -- the item-to-item gap *within* one list, at one
% fixed nesting depth (as opposed to the drop-space mechanism just above,
% which only fires when a nested list *closes* back up to its parent).
% Ratios, not lengths -- same "scale with font size" reasoning and fix as
% \gothamItemizeLevelDropSpaceRatioI/II/III above (that comment has the
% full story on why a plain \newlength+\setlength pair here would freeze
% to one fixed pt value instead of actually tracking each level's own
% font size). Values only, defined here for discoverability alongside the
% drop-space ones above -- actually *applied* inside the `\itemize`
% redefinition itself (search "\beamer@cramped (beamerbaselocalstructure.sty"
% further above), not via a separate hook here; see that comment for why.
% Found missing entirely (2026-07-21, direct instruction, see
% slides/ZZO-Itemize-Spacing-Reference-current.tex's own header comment
% for the rules this and the drop-space mechanism together need to
% satisfy): this deck's own itemize never set \itemsep itself at any
% level -- confirmed directly, grepping this whole file -- so every depth
% shared whatever bare default beamer's own theme happens to leave
% \itemsep at, uniformly. That produced two visible problems, both
% measured directly (pymupdf against slides/2c2-Topics-Around-the-CAPM-Empirical-Validity-current.tex
% and an isolated bisection test, not guessed): (a) level-2 and level-3
% same-level item-to-item gaps came out *identical* (12.6pt each,
% confirmed against an isolated reproduction), since nothing distinguished
% them by depth at all -- only a closing transition ever got a
% level-dependent gap, never a same-level one; (b) that shared default
% \itemsep read as *smaller* than a single item's own wrapped-line
% \baselineskip (13.7-13.8pt within a wrap vs. 12.6pt between items),
% backwards from how a list should read -- an item boundary should never
% be tighter than a same-item line wrap.
% Ratios themselves are not arbitrary -- each chosen, then re-measured
% against the compiled ZZO reference slide and adjusted, specifically to
% satisfy all three of that slide's own header-comment rules together
% (extended to level 4 the same way, 2026-07-21), not just "level 1
% biggest" in isolation:
% - Rule I (level-1 same-level gap > level-2's > level-3's > level-4's)
%   falls straight out of \gothamItemsepRatioI > ...II > ...III > ...IV,
%   added on top of the same \baselineskip every level shares (itself
%   already smaller per level here, since this theme's own itemize/
%   enumerate body font shrinks with depth too -- confirmed directly,
%   10.95pt/10pt/9pt/9pt for levels 1-4 respectively, level 4 sharing
%   level 3's own font since this theme has no distinct "subsubsub" size
%   of its own).
% - Rule II (every level's own same-level gap > that level's own wrapped-
%   line gap) needs \gothamItemsepRatioIV alone (the smallest, hence the
%   binding case) to still be positive and large enough that it clearly
%   reads as more open than a plain wrapped line once rendered, not just
%   bigger on paper.
% - Rule III (a level's own same-level gap < that same level's own
%   *closing*-transition gap) is *not* automatic just because
%   \gothamItemizeLevelDropSpaceRatio* was already bigger than the old
%   shared default -- \addvspace's own coalesce-to-the-larger-value
%   behavior (documented above) means the closing gap is
%   max(natural glue, that *drop-space* ratio), not a sum of the two. Each
%   drop-space ratio's own \ifcase branch is keyed on the depth of the
%   list that *just closed* (see \gothamItemizeLevelDropSpaceRatioIV's own
%   comment for the full story on why, confirmed via \typeout, not
%   assumed), so it lines up with the same-numbered itemsep ratio despite
%   describing a different transition -- \gothamItemizeLevelDropSpaceRatioIV
%   governs the level-4-list-closing-back-to-level-3 transition and is the
%   one that has to individually stay bigger than \gothamItemsepRatioIV.
%   This branch was missing entirely at first (only I/II/III existed, a
%   leftover from before level 4 existed) -- confirmed directly on ZZO: the
%   level-4-to-level-3 closing transition measured 12.6pt with no drop-space
%   branch at all, less than level 4's own same-level gap of 13.4pt, a real
%   rule III violation. Adding the branch (\gothamItemizeLevelDropSpaceRatioIV,
%   starting at 0.5) fixed it. Re-verify rule III directly against ZZO's own
%   compiled PDF (don't assume it still holds on paper) if any of these ratios is
%   ever changed again.
\newcommand{\gothamItemsepRatioI}{0.4}
\newcommand{\gothamItemsepRatioII}{0.24}
\newcommand{\gothamItemsepRatioIII}{0.15}
\newcommand{\gothamItemsepRatioIV}{0.09}

% Per-level "opening" space -- the gap between a level-n item and the
% first item of a level-(n+1) list that immediately follows it, as
% opposed to \gothamItemizeLevelDropSpaceRatio* (a nested list *closing*
% back up) or \gothamItemsepRatio* (two items at the *same* depth).
% Added 2026-07-21, direct instruction, replacing the original three
% separate rules on slides/ZZO-Itemize-Spacing-Reference-current.tex's
% own header comment with one multi-part inequality stated per level n:
%   level-n -> level-(n-1) close  >  level-n same-level  >
%   level-n -> level-(n+1) open  >  level-n's own wrap gap
% The "open" term was entirely unhandled before this -- confirmed
% directly on ZZO: every opening transition (level 1 into level 2, level
% 2 into level 3, level 3 into level 4) measured *smaller* than that same
% level's own wrap gap, backwards from what the rule needs, since nothing
% but beamer's own bare default \topsep governed it.
% \AtBeginEnvironment{itemize}, not \topsep -- \topsep was tried first
% (set via the same \ifcase\@itemdepth block, right after \itemsep,
% inside the `\itemize` redefinition above) and measured *unreliable*:
% two of the three levels barely moved despite a sizeable ratio, while
% the third moved by roughly the right amount -- \topsep's own value is
% evidently consulted by `\list` at some point this override comes too
% late for, at least for some levels, not evenly. \addvspace via
% \AtBeginEnvironment instead measured consistently and predictably
% (each level's own change matched its own ratio's calculated pt value
% directly), the same mechanism \gothamItemizeLevelDropSpaceRatio* already
% relies on for the closing side. This hook fires with \@itemdepth
% reflecting the *parent* item's own depth (confirmed via \typeout, not
% assumed: opening the outermost level-1 list itself measured 0, opening
% a level-2 list from within a level-1 item measured 1, level-3-from-
% level-2 measured 2, level-4-from-level-3 measured 3) -- the opposite of
% \gothamItemizeLevelDropSpaceRatio*'s own hook, which keys on the depth
% of the list that *just closed*. That makes the \ifcase numbering here
% match the "n" in the rule directly: case 1 is level-1's own open-to-
% level-2 ratio, case 2 is level-2's own open-to-level-3, case 3 is
% level-3's own open-to-level-4; case 0 (opening the outermost list
% itself) is a genuine no-op, correctly -- there's no parent item above
% level 1 for that transition to separate from.
% Ratios chosen, then re-measured against the compiled ZZO reference and
% adjusted, the same way as every other ratio here: each has to
% individually land strictly between that level's own \gothamItemsepRatio*
% (the ceiling) and its own wrap gap (the floor). Re-verify directly
% against ZZO's own compiled PDF (don't assume it still holds on paper)
% if any of these, or the itemsep/drop-space ratios they're measured
% against, are ever changed again.
% \gothamItemizeLevelOpenSpaceRatioI trimmed from 0.28 to 0.02 (2026-07-21,
% direct instruction: "the level-1 to level-2 distance...seems a bit
% much"). Confirmed directly: even at 0, this transition still measured
% ~16.3pt -- \addvspace's own coalesce-to-the-larger-value behavior means
% the ratio can only ever push this *above* whatever glue the list
% mechanism inserts here regardless (that natural floor, not the ratio),
% so 0.28 down to 0 only ever moved the actual gap by ~2pt (18.5->16.3);
% 0.02 splits the difference, keeping the mechanism nominally "in effect"
% rather than reading as accidentally disabled. \gothamItemizeLevelOpenSpaceRatioIII
% (level-3-to-level-4) had to come down too, from 0.08 to 0.02, purely to
% restore its own margin -- changing RatioI shifted that unrelated
% transition's own measured gap closer to level 3's own same-level gap
% (12.1 vs. 12.0, too thin), confirmed directly on ZZO, not assumed;
% ratios elsewhere in this file have shown the same small-ripple behavior
% before (see \gothamItemizeLevelDropSpaceRatioIV's own comment) --
% re-measure the *whole* table after touching any single one of these,
% not just the transition being changed.
\newcommand{\gothamItemizeLevelOpenSpaceRatioI}{0.02}
\newcommand{\gothamItemizeLevelOpenSpaceRatioII}{0.15}
\newcommand{\gothamItemizeLevelOpenSpaceRatioIII}{0.02}
\AtBeginEnvironment{itemize}{%
	\ifcase\@itemdepth
		% 0: opening the outermost (level-1) list itself -- no parent item
		% precedes it, nothing to add.
	\or\addvspace{\gothamItemizeLevelOpenSpaceRatioI em}% opening a level-2 list from a level-1 item
	\or\addvspace{\gothamItemizeLevelOpenSpaceRatioII em}% opening a level-3 list from a level-2 item
	\or\addvspace{\gothamItemizeLevelOpenSpaceRatioIII em}% opening a level-4 list from a level-3 item
	\fi
}
\makeatother

% A \tableofcontents-driven agenda slide (see
% slides/nkj-Agenda-current.tex) needs the exact same top-margin/line-
% spacing as every other list in this deck, not whatever gotham's own
% "section in toc" templates use (both "gotham bullet" and "gotham simple"
% carry their own spacing/rule/font choices, tuned for gotham's generic
% look, not this project's own itemize metrics). Rather than re-tuning a
% *third* set of constants to match, this template just emits a plain
% \item -- wrapping the whole \tableofcontents call in a real
% \begin{enumerate}...\end{enumerate} then makes each section entry a
% genuine enumerate item, reusing this file's own \leftmargini/\itemsep/
% \gothamContentTopShift machinery exactly as any other list would, so it
% automatically stays consistent with the rest of the deck instead of
% drifting out of sync with it later.
\defbeamertemplate{section in toc}{gotham plain enumerate}{\item \inserttocsection}

% Even with the plain-enumerate template above, \tableofcontents still
% doesn't match this deck's own normal itemize spacing -- beamer's own
% \beamer@sectionintoc (beamerbasetoc.sty), which wraps *every* section
% entry regardless of which "section in toc" template is active,
% unconditionally inserts a fixed \vskip1.5em before each entry (an
% unbounded \vfill, even worse, if subsections aren't hidden -- neither is
% reachable via a public \tableofcontents option, just this show/hide
% binary), *and* wraps the entry itself in \hbox{\vbox{...}} -- box
% wrapping that (like \gothamMinipageTopCorrection's own minipage case
% elsewhere in this file) changes how surrounding baselineskip/topskip
% glue gets inserted compared to plain paragraph text. Both together added
% roughly 13pt of extra gap
% between entries beyond this deck's own normal ~18.6pt level-1 item
% spacing (measured directly against
% slides/57r-Reasons-for-Leasing-current.tex's own compiled PDF).
% Neither is controllable via a public \tableofcontents option, so this
% redefines \beamer@sectionintoc itself -- scoped safely since
% \tableofcontents is only ever called on the one Agenda slide in this
% deck -- keeping its numbering/hyperlink/part-number/hide-checking logic
% intact but dropping the vskip/vfill and the hbox/vbox wrapping, so an
% entry behaves exactly like a plain \item once printed.
\makeatletter
\def\beamer@sectionintoc#1#2#3#4#5{%
	\ifnum\c@tocdepth>0%
	\ifnum#4=\beamer@showpartnumber%
	{%
	\beamer@saveanother%
	\gdef\beamer@todo{}%
	\beamer@slideinframe=#1\relax%
	\expandafter\only\beamer@tocsections{\gdef\beamer@todo{%
			\beamer@tempcount=#5\relax%
			\advance\beamer@tempcount by\beamer@sectionadjust%
			\ifnum\beamer@tempcount>0
				\ifnum\beamer@tempcount>\beamer@toclastsection
					\def\inserttocsectionnumber{}%
				\else
					\edef\inserttocsectionnumber{\the\beamer@tempcount}%
				\fi
			\else
				\def\inserttocsectionnumber{}%
			\fi%
			\def\inserttocsection{\hyperlink{Navigation#3}{#2}}%
			\beamer@tocifnothide{\ifnum\c@section=#1\beamer@toc@cs\else\beamer@toc@os\fi}%
			{%
				\beamer@tocact{\ifnum\c@section=#1\beamer@toc@cs\else\beamer@toc@os\fi}{section in toc}%
			}%
		}%
	}%
	\beamer@restoreanother%
	}%
	\beamer@todo%
	\fi\fi%
}
\makeatother

% Left-aligns an enumerate's labels within a caller-chosen \labelwidth --
% used by the \tableofcontents Agenda slide (see
% slides/nkj-Agenda-current.tex's own comment for why its labels need
% this at 3.6pt). Defined here as a named helper rather than written as a
% literal \renewcommand{\makelabel}[1]{...#1...} inside the slide's own
% frame body, because that literal form -- specifically its raw `#1`
% parameter token -- breaks inside any frame that \tableofcontents forces
% through beamer's overlay reprocessing: beamer replays such a frame's
% collected body via \def-based machinery (\loop's own \iterate, and
% \beamer@doifinframe), and a `#` in the material of any \def must be
% doubled -- so the compile hits "! Illegal parameter number in definition
% of \iterate." (twice) plus the same for \beamer@doifinframe, once per
% compile, at that frame's own \end{frame}. TeX's error recovery ("I'll
% assume you meant ##") happens to produce the intended definition anyway,
% which is why the rendered PDF was always correct despite the errors --
% but lualatex still exits nonzero, which is exactly what kept breaking
% latexmk's automated lualatex/biber/rerun sequencing on the main deck
% (the "Known open issue" long tracked in CONVERSION_NOTES.md, originally
% misattributed to a biblatex interaction -- see its "Resolved" entry).
% Hoisting the definition out here means the frame body carries only this
% macro's *call* token, with no `#` in it, which sidesteps the whole
% mechanism. The inner \makelabel parameter is written `##1` below since
% it sits inside another \newcommand's body -- standard nesting, and also,
% incidentally, exactly the doubling rule the inline version fell foul of.
\newcommand{\gothamTocLabelFlushLeft}[1]{%
	\setlength{\labelwidth}{#1}%
	\renewcommand{\makelabel}[1]{\makebox[\labelwidth][l]{##1}}%
}

% Same idea as \gothamTocLabelFlushLeft above, plus an extra fixed
% \hspace shift before the label (for a list that needs to start further
% right than the plain nested-list indent, e.g. to line up with a
% sibling item's own text rather than its bullet -- see
% slides/scs-Political-Risk-Stability-Index-current.tex's own v1-5-0/
% v1-6-0 comments for the worked example this was hoisted out of).
% Found 2026-07-20 that the raw-`#1`-in-frame-body bug documented above
% is *not* actually \tableofcontents-specific, contrary to that comment's
% own framing -- slides/scs-Political-Risk-Stability-Index-current.tex has
% no \tableofcontents anywhere in the deck-wide overlay chain leading to
% it, yet hit the identical "Illegal parameter number in definition of
% \iterate" at its own \end{standardframe} from the exact same literal-
% \renewcommand{\makelabel}[1]{...#1...}-in-frame-body pattern. Whatever
% the precise trigger condition actually is, the same hoist-to-preamble
% fix (a bare macro *call*, no `#` token, left in the frame body) sidesteps
% it either way.
\newcommand{\gothamEnumerateLabelFlushLeftShifted}[2]{%
	\setlength{\labelwidth}{#2}%
	\renewcommand{\makelabel}[1]{\hspace{#1}\makebox[\labelwidth][l]{##1}}%
}

% Width of a single digit in this deck's normal content font, measured once
% here (safe, plain preamble-load-time \settowidth) rather than inline in
% the Agenda slide's own frame body -- \tableofcontents there puts the whole
% frame through beamer's overlay-replay machinery, which the comment above
% \gothamTocLabelFlushLeft already found chokes on raw tokens (like a bare
% \settowidth call) sitting directly in a replayed frame body, the same
% "Missing number"/"Illegal unit of measure" class of error as that `#`
% problem, confirmed directly the same way. Used by
% slides/nkj-Agenda-current.tex's own "widen the number-to-title gap"
% fix, which only ever needs to *reference* this precomputed value inside
% the frame, not measure anything itself there.
\newlength{\gothamTocLabelExtraGap}
\settowidth{\gothamTocLabelExtraGap}{0}

% For a two-column layout (e.g. bullets next to a chart/formula), nest
% beamer's own "columns" directly inside standardframe -- do NOT bake
% \begin{columns}/\end{columns} into a custom environment's begin/end code:
% beamer collects each frame's body as a token list to replay it once per
% overlay, and a columns environment that is opened in the "before" half and
% closed in the "after" half of a *different* environment confuses that
% collection (it never finds a properly matched \end{frame}). Nesting it in
% the document body, as below, avoids the issue entirely:
%
%   \begin{standardframe}{Title}
%   	\begin{columns}[t,onlytextwidth]
%   		\begin{column}{0.48\linewidth} ... \end{column}
%   		\begin{column}{0.48\linewidth} ... \end{column}
%   	\end{columns}
%   \end{standardframe}

% For automatic section dividers and single-statement "standout" slides, use
% what the gotham theme already provides instead of a custom environment:
%   \section{...}                                                (section divider, automatic)
%   \begin{frame}[standout]{Key takeaway}\end{frame}             (big single statement)
%   \begin{block}{...}/\begin{alertblock}{...}/\begin{exampleblock}{...}/\begin{referenceblock}{...}  (definitions, formulas, examples, references)
% For a course's own title slide, section-topic overview, quote, figure-only,
% and presenter-attention slides, see the five environments below.

% block/alertblock/exampleblock recolored to this project's own palette
% (2026-07-05, direct instruction) instead of gotham's own default colors
% -- gotham's own per-block-type color machinery (\gotham@block in
% beamerinnerthemegotham.sty) turns out to be dead code, never actually
% wired into any template, so these three blocks were rendering via plain
% beamer's own stock templates/colors the whole time; overriding the
% beamer color names directly here works regardless, since those stock
% templates read their colors from these same standard names (block~title/
% block~body and their ~alerted/~example variants) via beamercolorbox.
% Body background is the same c_tertiary_gray for all four
% block types (including referenceblock below) -- only each one's own
% title/header background differs. block's own header went
% gray→secondaryblue→secondaryviolet across three separate direct
% instructions in as many days -- secondaryviolet is the one actually in
% effect now.
\setbeamercolor{block title}{bg=c_secondary_violet, fg=c_primary_dark}
\setbeamercolor{block body}{bg=c_tertiary_gray, fg=c_primary_dark}
\setbeamercolor{block title example}{bg=c_secondary_green, fg=c_primary_dark}
\setbeamercolor{block body example}{bg=c_tertiary_gray, fg=c_primary_dark}
\setbeamercolor{block title alerted}{bg=c_secondary_red, fg=c_primary_dark}
\setbeamercolor{block body alerted}{bg=c_tertiary_gray, fg=c_primary_dark}

% All four block types' left/right margins widened to match a top-level
% itemize bullet's own text indent (2026-07-05, direct instruction: "the
% blocks should not have a smaller left-hand margin than the content
% text") -- a plain block spans \linewidth as-is, which lines up with the
% *frame's* own text margin, not with an itemize bullet's own text (which
% sits further right, past the bullet symbol and its gap) -- measured
% directly against a compiled slide with both an itemize and a block
% (slides/afa-Excursus-Price-Formation-current.tex): a bullet's own text
% starts 12.96pt further right than a plain block's left edge did.
%
% First attempt used changepage's \adjustwidth, on paper the "correct" tool
% for temporarily changing margins around a chunk of content -- it silently
% does *nothing* inside a beamer frame (beamer computes \linewidth/\hsize
% through its own internal frame-geometry machinery, not the page-geometry
% registers changepage adjusts, which are meant for article/report-style
% documents), and *loading the package at all* actively corrupted the
% title bar/logo on every frame (some interaction with gotham's own
% tikz/remember-picture-based title-bar drawing, not chased further since
% the package needed to go regardless) -- caught by actually compiling and
% looking, not assumed to work from the package's own documented purpose.
% Fixed instead with a plain \hspace{}+\begin{minipage}{\linewidth-2\times}
% wrapped around each block type's own begin/end template pair below --
% minipage genuinely resets \linewidth to its own local width for anything
% typeset inside it (already relied on elsewhere in this deck, e.g. the
% side-by-side text/image slides), which a bare beamercolorbox alone
% doesn't do on its own.
%
% 12.96pt matched a *second*-level bullet's own dot (which happens to
% coincide almost exactly with where a *first*-level bullet's own *text*
% starts, both ~21.6pt in from the frame's raw margin -- a coincidence of
% this deck's specific bullet-dot/gap/font metrics, not a general rule) --
% corrected the next day (direct instruction: align with the *first*-level
% bullet's own margin instead) to 4.20pt, matching a first-level bullet's
% own dot position (~12.84pt in from the frame's raw margin) instead.
\newlength{\gothamBlockSideMargin}
\setlength{\gothamBlockSideMargin}{4.20pt}

% Shared helper for a block title bar that carries a marker symbol flush
% right, opposite the title text (currently: infoblock's circled "i",
% definitionblock's $\triangleq$, exampleblock's $\in$ -- see each of
% those three's own template for the direct instruction behind its own
% marker choice). 2026-07-08, direct instruction, after infoblock's own
% "circled i nearly touching the title text" report on a two-line title
% (slides/u5h-Long-term-Bonds-Provisions-current.tex's own "Denomination,
% face value, or principal" step) -- generalized here so the same fix
% covers all three marker types at once, and any future one, rather than
% patching each block's own template separately.
%
% The marker itself is placed via a zero-width box *before* the title
% text, not via a trailing `\hfill<marker>` the way all three started out
% (definitionblock/exampleblock still do, until updated below) -- `\hfill`
% only fills space on whatever *line* it ends up typeset on, so for a
% title long enough to wrap onto a second line, `\hfill<marker>` lands on
% that second line instead of the first, printing the marker at the
% bottom right of the box instead of the top right (this is the same root
% cause infoblock's own marker already had fixed once already, see that
% template's own comment for the fuller writeup of *that* bug specifically
% -- this pushes the same fix out to the other two, and adds the part
% that bug's own fix didn't yet cover, below).
% `\makebox[0pt][l]{\hspace{\linewidth}\makebox[0pt][r]{<marker>}}` fixes
% *that* much on its own: the outer zero-width box doesn't displace the
% title text at all (still starts at the same point, still wraps however
% it needs to), but *inside* that zero-width box the marker is placed at
% a fixed horizontal offset (this box's own full \linewidth) from the
% title's own starting point, i.e. always the top right corner of the
% title bar regardless of how the title text itself wraps.
%
% What that alone still doesn't fix -- and what this round's actual
% report was about -- is the title text itself having no *awareness* that
% a marker will render on top of wherever it happens to end: a title just
% long enough to reach close to the marker's own fixed position, without
% actually needing to wrap, ends up visually touching or running under
% it. Fixed by also wrapping the title text itself in a
% `\begin{minipage}[t]{\linewidth-\gothamBlockTitleMarkerWidth}`, one
% \gothamBlockTitleMarkerWidth narrower than the full title bar, so the
% text itself wraps *before* reaching the marker's own reserved space
% instead of running into it -- confirmed directly against the compiled
% PDF on the exact case from the report. \gothamBlockTitleMarkerWidth is
% one shared reserved width for all three current marker types rather
% than a separate hand-measured width per symbol, which would need
% re-measuring by hand every time a new marker is added -- comfortably
% covers the widest of the three (checked directly) with a bit to spare.
% #1 = the marker's own content (e.g. \textcircled{\scriptsize i}, or a
% math-mode symbol like $\triangleq$); #2 = the title text
% (\insertblocktitle). Call as
% \gothamBlockTitleWithMarker{<marker>}{\insertblocktitle} in place of a
% bare \insertblocktitle, after the usual \usebeamerfont* call.
\newlength{\gothamBlockTitleMarkerWidth}
\setlength{\gothamBlockTitleMarkerWidth}{16pt}
% 2026-07-08 (direct instruction), two more fixes to this same helper:
% - A title long enough to wrap onto a second line
%   (slides/u5h-Long-term-Bonds-Provisions-current.tex's own
%   "Denomination, face value, or principal") read as *justified* rather
%   than left-aligned -- the minipage below had no explicit alignment, so
%   it fell back to LaTeX's default full justification, which only ever
%   visibly stretches a line short of the box's own right edge, i.e. only
%   shows up once a title is long enough to actually wrap. Every other
%   (single-line, shorter) title already looked left-aligned regardless,
%   since a line that's naturally shorter than the box is never stretched
%   whether or not justification is on -- confirmed directly, this is why
%   the bug went unnoticed until a title finally wrapped. `\raggedright`
%   fixes it without needing anything different for the single-line case.
% - Different titles' own title *bars* rendered at different heights
%   depending on whether that particular title's text happened to include
%   a descender (e.g. "Security"/"y", "Moody's rating at issue date"/"y")
%   or not (e.g. "Registered", "Trustee") -- the same class of bug already
%   fixed on slides/qxr-Introduction-to-R-Assignments-current.tex's own
%   colorboxes (see that slide's own comment): a box that tightly wraps its
%   own text naturally varies in height/depth with that text's own
%   ascenders/descenders. A leading `\strut` gives every title's first (and
%   for all but this one current title, only) line the same guaranteed
%   minimum height/depth, matching the qxr fix's own reasoning -- confirmed
%   directly, every current info/definition/example block title bar now
%   renders at the same height regardless of its own text.
\newcommand{\gothamBlockTitleWithMarker}[2]{%
	\makebox[0pt][l]{\hspace{\linewidth}\makebox[0pt][r]{#1}}%
	\begin{minipage}[t]{\dimexpr\linewidth-\gothamBlockTitleMarkerWidth\relax}\raggedright\strut#2\end{minipage}%
}

% A fourth block type, for a reference/citation-style callout (first used
% on slides/afa-Excursus-Price-Formation-current.tex and
% slides/4r7-IPO-Underpricing-International-current.tex, both replacing
% an earlier rotated hand-drawn callout box) -- beamer has no built-in
% "reference" block variant, so this defines one from scratch, mirroring
% exactly how beamer's own alertblock/exampleblock are built (compare
% beamerbaselocalstructure.sty's \newenvironment<>{alertblock} and
% beamerinnerthemedefault.sty's "block alerted begin/end" templates).
%
% First attempt (2026-07-05) copied these under the "rounded" style/
% beamerboxesrounded, on the wrong assumption that "rounded" was this
% deck's active block style generally -- it renders with visibly rounded
% corners, caught by the user, unlike block/alertblock/exampleblock's own
% square corners. Those three actually use beamerinnerthemedefault.sty's
% *starred* "default"-style templates (plain beamercolorbox, no rounding),
% which \defbeamertemplate*'s star selects immediately when that inner
% theme loads, regardless of beamer's own \defbeamertemplateparent{blocks}
% fallback to "rounded" -- gotham's own inner theme derives from
% beamerinnerthemedefault, so "default" (square) was the deck's real active
% style the whole time. Rewritten below to match that "default" style
% exactly (plain beamercolorbox instead of beamerboxesrounded) for
% consistent square corners across all four block types.
\setbeamercolor{block title reference}{bg=c_secondary_orange, fg=c_primary_dark}
\setbeamercolor{block body reference}{bg=c_tertiary_gray, fg=c_primary_dark}
% block/alertblock/exampleblock's own titles render bold via
% beamerfontthemegotham.sty's \setbeamerfont{block title}{series=\bfseries}
% (and an explicit ...alerted equivalent; example inherits it from
% beamer's own base font-fallback chain, verified directly -- both
% "Example"/"Caveat" titles on slides/ZZX-Definition-formula-current.tex
% render in the same heavier FiraSans-Regular as "Definition" despite
% gotham never setting "block title example" explicitly). "reference" is a
% brand new type name unknown to any of that inheritance, so it fell back
% to a plain, non-bold font with no color-analogous safety net (unlike
% \setbeamercolor{block title reference} above, colors for genuinely new
% names still resolve fine) -- caught only by comparing the actual
% rendered font name in the compiled PDF (FiraSans-Light, matching plain
% body text) against block/alertblock/exampleblock's own (FiraSans-Regular,
% this deck's de facto "bold" for Fira Sans), not by eye alone. Explicit
% setting needed, matching gotham's own two.
\setbeamerfont{block title reference}{size=\normalsize, series=\bfseries}
\defbeamertemplate*{block reference begin}{default}[1][]
{
	\hspace{\gothamBlockSideMargin}\begin{minipage}{\dimexpr\linewidth-2\gothamBlockSideMargin\relax}
	\par\vskip\medskipamount%
	\begin{beamercolorbox}[colsep*=.75ex,#1]{block title reference}
		\usebeamerfont*{block title reference}\strut\insertblocktitle%
	\end{beamercolorbox}\noframebreak%
	{\parskip0pt\par}%
	\ifbeamercolorempty[bg]{block title reference}
	{}
	{\ifbeamercolorempty[bg]{block body reference}{}{\nointerlineskip\vskip-0.5pt}}%
	\usebeamerfont{block body reference}%
	\begin{beamercolorbox}[colsep*=.75ex,vmode,#1]{block body reference}%
		\ifbeamercolorempty[bg]{block body reference}{\vskip-.25ex}{\vskip-.75ex}\vbox{}%
}
\defbeamertemplate*{block reference end}{default}
{\end{beamercolorbox}\vskip\smallskipamount\end{minipage}}
\newenvironment<>{referenceblock}[1]{%
	\begin{actionenv}#2%
		\def\insertblocktitle{#1}%
		\par%
		\usebeamertemplate{block reference begin}}
	{\par%
		\usebeamertemplate{block reference end}%
	\end{actionenv}}

% A sixth block type, for general background/context info that isn't
% itself a definition, example, or reference citation (2026-07-06, direct
% instruction: "use secondaryblue for a new type of info block") -- built
% the same from-scratch way as referenceblock just above (a brand-new
% style name needs \defbeamertemplate*'s registering form, not
% \setbeamertemplate -- see referenceblock's own comment for why). Marker
% is a circled "i" (\textcircled{\scriptsize i}, the generic LaTeX-kernel
% \textcircled -- not a Unicode "circled i" character: this deck's own
% Fira Sans/Fira Mono have no glyph for that codepoint either, the same
% missing-glyph problem \triangleq's own comment already found for
% definitionblock, and \textcircled sidesteps it entirely since it just
% draws a plain circle symbol around an ordinary letter rather than
% depending on one precomposed font glyph existing at all).
%
% The marker is placed via a zero-width box *before* \insertblocktitle,
% not via a trailing \hfill\textcircled{...} the way it first looked most
% natural to write -- \hfill only fills space on whatever line it's
% typeset on, and for a title long enough to wrap onto a second line (2026-07-08,
% found on the last reveal step of slides/ib3-Long-term-Bonds-Repayment-current.tex),
% \hfill\textcircled{...} then lands on that *second* line instead of the
% first, printing the marker at the bottom right of the box instead of the
% top right. Fixed originally with a one-off zero-width-box trick here;
% generalized the same day into the shared \gothamBlockTitleWithMarker
% helper (see its own comment, near \gothamBlockSideMargin above) once a
% *second*, related bug turned up on this exact block type (a long title
% visually touching the marker even without wrapping) -- that helper is
% what's actually called below now, not a copy of this trick inline.
\setbeamercolor{block title info}{bg=c_secondary_blue, fg=c_primary_dark}
\setbeamercolor{block body info}{bg=c_tertiary_gray, fg=c_primary_dark}
\setbeamerfont{block title info}{size=\normalsize, series=\bfseries}
\defbeamertemplate*{block info begin}{default}[1][]
{
	\hspace{\gothamBlockSideMargin}\begin{minipage}{\dimexpr\linewidth-2\gothamBlockSideMargin\relax}
	\par\vskip\medskipamount%
	\begin{beamercolorbox}[colsep*=.75ex,#1]{block title info}
		\usebeamerfont*{block title info}%
		\gothamBlockTitleWithMarker{\textcircled{\scriptsize i}}{\insertblocktitle}%
	\end{beamercolorbox}\noframebreak%
	{\parskip0pt\par}%
	\ifbeamercolorempty[bg]{block title info}
	{}
	{\ifbeamercolorempty[bg]{block body info}{}{\nointerlineskip\vskip-0.5pt}}%
	\usebeamerfont{block body info}%
	\begin{beamercolorbox}[colsep*=.75ex,vmode,#1]{block body info}%
		\ifbeamercolorempty[bg]{block body info}{\vskip-.25ex}{\vskip-.75ex}\vbox{}%
}
\defbeamertemplate*{block info end}{default}
{\end{beamercolorbox}\vskip\smallskipamount\end{minipage}}
\newenvironment<>{infoblock}[1]{%
	\begin{actionenv}#2%
		\def\insertblocktitle{#1}%
		\par%
		\usebeamertemplate{block info begin}}
	{\par%
		\usebeamertemplate{block info end}%
	\end{actionenv}}

% A fifth block type, for a "this introduces/defines a term" callout,
% visually distinguished from a plain \block (2026-07-06, direct
% instruction: "I would like to distinguish definition blocks from other
% blocks") by a $\triangleq$ symbol sitting flush right in the title bar,
% opposite the title text itself. Exact same colors/margins/font as plain
% \block (reuses "block title"/"block body", not a new color pair) --
% only the title-bar content differs, one extra `\hfill$\triangleq$`
% after `\insertblocktitle` -- so this mirrors referenceblock's own
% from-scratch pattern above (a brand new "definition" style name via
% \defbeamertemplate*, not \setbeamertemplate, for the same reason
% referenceblock's own comment gives: this style name has never been
% registered before, so the *-registration form is the correct one here,
% unlike plain block/alertblock/exampleblock below which are redefining
% already-registered "default"-style bodies). `\triangleq` is a genuine
% math symbol, so it needs `$...$` -- renders in this theme's math font
% (Computer Modern, see beamerfontthemegotham.sty's `no-math` option),
% a deliberate, visible one-glyph exception in an otherwise Fira Sans
% title bar, the same tradeoff every other math symbol in this deck
% already accepts.
\defbeamertemplate*{block definition begin}{default}[1][]
{
	\hspace{\gothamBlockSideMargin}\begin{minipage}{\dimexpr\linewidth-2\gothamBlockSideMargin\relax}
	\par\vskip\medskipamount%
	\begin{beamercolorbox}[colsep*=.75ex,#1]{block title}
		\usebeamerfont*{block title}\gothamBlockTitleWithMarker{$\triangleq$}{\insertblocktitle}%
	\end{beamercolorbox}\noframebreak%
	{\parskip0pt\par}%
	\ifbeamercolorempty[bg]{block title}
	{}
	{\ifbeamercolorempty[bg]{block body}{}{\nointerlineskip\vskip-0.5pt}}%
	\usebeamerfont{block body}%
	\begin{beamercolorbox}[colsep*=.75ex,vmode,#1]{block body}%
		\ifbeamercolorempty[bg]{block body}{\vskip-.25ex}{\vskip-.75ex}\vbox{}%
}
\defbeamertemplate*{block definition end}{default}
{\end{beamercolorbox}\vskip\smallskipamount\end{minipage}}
\newenvironment<>{definitionblock}[1]{%
	\begin{actionenv}#2%
		\def\insertblocktitle{#1}%
		\par%
		\usebeamertemplate{block definition begin}}
	{\par%
		\usebeamertemplate{block definition end}%
	\end{actionenv}}

% block/alertblock/exampleblock begin/end: same body as
% beamerinnerthemedefault.sty's own "default"-style templates (the ones
% actually in effect, per the comment above referenceblock), plus the
% minipage margin wrapper around each pair -- set via plain
% \setbeamertemplate{block begin}{...} (overriding the *body* of the
% existing "default" style directly) rather than
% \defbeamertemplate*{block begin}{default}{...} like referenceblock's own
% (new) templates above. The two look almost interchangeable but aren't:
% \defbeamertemplate* internally does a \newcommand-based *registration* of
% a style name, which errors outright ("... already defined") when that
% exact template+style pair already exists, as "block begin"/"default" and
% the other two pairs already do (registered once by
% beamerinnerthemedefault.sty itself) -- caught only by compiling and
% reading the log, not from any visible symptom in the rendered PDF (LaTeX
% recovers from the error by leaving the *original* stock template body in
% place, so these three blocks kept working and even picked up their new
% *colors* correctly via \setbeamercolor above, since color lookups happen
% at render time either way -- just silently without the margin fix, which
% *is* baked into the template body itself). \setbeamertemplate{name}{body}
% (no style-name argument) instead does a plain \def of that style's
% existing body storage, which is exactly "redefine what's already there"
% and carries no such restriction -- the correct tool whenever the target
% template+style pair is already registered, as opposed to
% \defbeamertemplate*, which is only for registering a style name for the
% very first time (as referenceblock's brand-new "reference" style name
% needs). Dropped the stock templates' own "#1" optional-argument slot
% (extra beamercolorbox options) in the process: nothing in this deck ever
% calls \usebeamertemplate{block begin}[...] with a bracket argument to
% populate it, so it was always empty in practice anyway.
% \strut before \insertblocktitle here, and in "block alerted"/"block
% reference" below: same title-bar-height-vs-descender fix as
% \gothamBlockTitleWithMarker's own comment above -- these three don't go
% through that helper at all (no marker symbol), so needed the same
% `\strut` added directly instead.
\setbeamertemplate{block begin}{
	\hspace{\gothamBlockSideMargin}\begin{minipage}{\dimexpr\linewidth-2\gothamBlockSideMargin\relax}
	\par\vskip\medskipamount%
	\begin{beamercolorbox}[colsep*=.75ex]{block title}
		\usebeamerfont*{block title}\strut\insertblocktitle%
	\end{beamercolorbox}\noframebreak%
	{\parskip0pt\par}%
	\ifbeamercolorempty[bg]{block title}
	{}
	{\ifbeamercolorempty[bg]{block body}{}{\nointerlineskip\vskip-0.5pt}}%
	\usebeamerfont{block body}%
	\begin{beamercolorbox}[colsep*=.75ex,vmode]{block body}%
		\ifbeamercolorempty[bg]{block body}{\vskip-.25ex}{\vskip-.75ex}\vbox{}%
}
\setbeamertemplate{block end}
{\end{beamercolorbox}\vskip\smallskipamount\end{minipage}}

% Warning-triangle marker for alertblock (direct instruction, 2026-07-12),
% same \gothamBlockTitleWithMarker treatment as definitionblock's
% $\triangleq$/exampleblock's $\in$/infoblock's \textcircled{\scriptsize i}
% above. No single built-in glyph renders a triangle-with-exclamation-mark
% in this deck's fonts (checked the same way \triangleq's own comment
% did), so it's hand-drawn: a plain tikz outline triangle with a bold "!"
% centered inside, sized by eye against \textcircled{\scriptsize i} until
% the two read as the same visual weight in a block title bar.
\newcommand{\gothamAlertMarker}{%
	\tikz[baseline=-0.6ex]{
		\draw[line width=1.1pt, line join=round] (-0.2,-0.11) -- (0.2,-0.11) -- (0,0.31) -- cycle;
		\node[font=\fontsize{7}{7}\selectfont\bfseries] at (0,0.01) {!};
	}%
}
\setbeamertemplate{block alerted begin}{
	\hspace{\gothamBlockSideMargin}\begin{minipage}{\dimexpr\linewidth-2\gothamBlockSideMargin\relax}
	\par\vskip\medskipamount%
	\begin{beamercolorbox}[colsep*=.75ex]{block title alerted}
		\usebeamerfont*{block title alerted}\gothamBlockTitleWithMarker{\gothamAlertMarker}{\insertblocktitle}%
	\end{beamercolorbox}\noframebreak%
	{\parskip0pt\par}%
	\ifbeamercolorempty[bg]{block title alerted}
	{}
	{\ifbeamercolorempty[bg]{block body alerted}{}{\nointerlineskip\vskip-0.5pt}}%
	\usebeamerfont{block body alerted}%
	\begin{beamercolorbox}[colsep*=.75ex,vmode]{block body alerted}%
		\ifbeamercolorempty[bg]{block body alerted}{\vskip-.25ex}{\vskip-.75ex}\vbox{}%
}
\setbeamertemplate{block alerted end}
{\end{beamercolorbox}\vskip\smallskipamount\end{minipage}}

% $\in$ ("is an instance/element of") flush right in the title bar
% (2026-07-06, direct instruction, mirroring definitionblock's own
% $\triangleq$ -- "what symbol would you propose ... to symbolize an
% example?"): a concrete example is naturally "an instance of" the general
% concept it illustrates, the same relationship $\in$ already denotes
% between an element and a set/class -- and, same as $\triangleq$, it's a
% math-mode symbol with a real Computer Modern glyph (no Fira Sans/Fira
% Mono text-mode substitute needed, unlike most Unicode symbol candidates
% -- see Slidemaster_preamble_uni-graz.tex's own "beamer/TeX gotchas" entry on
% $\triangleq$ for why that matters here). Applied to every \exampleblock/
% \gothamfloatingexampleblock in the deck via this shared template, not a
% one-off per slide, matching how $\triangleq$ applies to every
% definitionblock through its own shared template.
\setbeamertemplate{block example begin}{
	\hspace{\gothamBlockSideMargin}\begin{minipage}{\dimexpr\linewidth-2\gothamBlockSideMargin\relax}
	\par\vskip\medskipamount%
	\begin{beamercolorbox}[colsep*=.75ex]{block title example}
		\usebeamerfont*{block title example}\gothamBlockTitleWithMarker{$\in$}{\insertblocktitle}%
	\end{beamercolorbox}\noframebreak%
	{\parskip0pt\par}%
	\ifbeamercolorempty[bg]{block title example}
	{}
	{\ifbeamercolorempty[bg]{block body example}{}{\nointerlineskip\vskip-0.5pt}}%
	\usebeamerfont{block body example}%
	\begin{beamercolorbox}[colsep*=.75ex,vmode]{block body example}%
		\ifbeamercolorempty[bg]{block body example}{\vskip-.25ex}{\vskip-.75ex}\vbox{}%
}
\setbeamertemplate{block example end}
{\end{beamercolorbox}\vskip\smallskipamount\end{minipage}}

% A narrower, freely-positionable version of \begin{block}{...} (2026-07-06,
% direct instruction) -- a plain block always spans the *current*
% \linewidth (that's what block begin/end's own \hspace+minipage wrapper
% above is measured against), so it can't be dropped into a fixed-width
% spot (e.g. beside a table, inside a tikz overlay \node) without first
% narrowing \linewidth some other way. \gothamfloatingblock just wraps a
% plain block in one more, *outer* minipage of the caller's own chosen
% width -- block begin/end's own inner minipage then measures itself
% against that narrower width instead of the frame's full one, nesting
% correctly with no changes needed to the block templates themselves.
% Title/body colors, margins, and font are exactly \block's own (same
% "block title"/"block body" beamer color names) -- this only changes the
% overall width and (by taking it out of the normal document flow, left to
% the caller) the position, not the look. Font size is also left to the
% caller (wrap the whole environment in a size declaration, e.g.
% {\footnotesize \begin{gothamfloatingblock}{5cm}{Excel} ... \end{gothamfloatingblock}})
% rather than fixed here, since the right size depends entirely on how
% much room the block ends up placed in.
\newenvironment{gothamfloatingblock}[2]{%
	\begin{minipage}{#1}
	\begin{block}{#2}
}{%
	\end{block}
	\end{minipage}
}

% Same idea, but wrapping \exampleblock instead of \block (2026-07-06,
% direct instruction, first used on
% slides/fr3-Forms-of-Capital-Market-Efficiency-Strong-Semistrong-Weak-current.tex
% and slides/ha9-Monacorp-Truck-Lease-Breakeven-Rent-current.tex, both
% switching an existing \gothamfloatingblock to the "example" green header
% instead of the plain violet one) -- not worth generalizing into a single
% macro parameterized over the block name: \gothamfloatingblock's own body
% is only two lines, so a third variant is cheaper to just write out than
% to build a block-name-dispatch layer for.
\newenvironment{gothamfloatingexampleblock}[2]{%
	\begin{minipage}{#1}
	\begin{exampleblock}{#2}
}{%
	\end{exampleblock}
	\end{minipage}
}

% Same idea again, wrapping \infoblock (2026-07-06, direct instruction,
% first used on
% slides/fr3-Forms-of-Capital-Market-Efficiency-Strong-Semistrong-Weak-current.tex's
% own "Random walk" formula box, switched from the green
% \gothamfloatingexampleblock to this blue one the same day).
\newenvironment{gothamfloatinginfoblock}[2]{%
	\begin{minipage}{#1}
	\begin{infoblock}{#2}
}{%
	\end{infoblock}
	\end{minipage}
}

% Every gothamfloatinginfoblock anchored top-right against the institute
% logo (the deck's standard
% "(current page.north east)+(-\gothamFrametitleRightpading,-\gothamContentTopAlign)"
% corner pattern) lands its own visible right edge short of the logo's,
% by a fixed amount independent of the block's own chosen width (#1) --
% \infoblock's underlying template insets its colored fill by
% \gothamBlockSideMargin from *both* sides of whatever width it's given
% (that margin exists to match a first-level bullet's own left-hand
% position, see its own comment above; the fact that it also narrows the
% *right* side is just a side effect of reusing one symmetric
% hspace+minipage trick for both, not a deliberate right-margin choice),
% partly offset by beamer's own beamercolorbox "colsep*" background bleed
% (colsep* pads the colored fill outward past its own nominal box edge
% without also insetting the text, unlike unstarred "colsep" -- see
% beamerbasecolor.sty's own \beamercolorbox definition). Root-caused and
% measured directly (pymupdf against a compiled probe deck, not derived
% purely on paper) 2026-07-21, after
% slides/j8y-Asset-Allocation-Example-Weights-current.tex's own one-off
% "+2.2pt" nudge (measured that same day against just the institute
% logo's yellow chip) turned out to disagree with
% slides/fr3-Forms-of-Capital-Market-Efficiency-Strong-Semistrong-Weak-current.tex's
% own older "+3.165pt" nudge (2026-07-06, measured against the logo's
% *full* gray-bordered footprint, shadow/card included) -- both slides'
% own nudges were each individually correct for their own chosen
% reference edge, they just didn't agree on what "the logo's own edge"
% means. Direct instruction, 2026-07-21: the *full* gray-bordered
% footprint is the deck's intended reference, so 3.165pt is the value
% kept here; j8y's own newer 2.2pt nudge was the one revised to match,
% not this one.
%
% Use "-\gothamFloatingInfoblockRightAlign" (not the bare
% "-\gothamFrametitleRightpading") for a gothamfloatinginfoblock's own
% anchor node whenever its right edge is meant to align with the logo --
% every current use of this pattern across the deck has been updated to
% match; new slides using this same corner pattern should use this
% length from the start rather than re-measuring their own one-off nudge.
\newlength{\gothamFloatingInfoblockRightAlign}
\setlength{\gothamFloatingInfoblockRightAlign}{\dimexpr\gothamFrametitleRightpading-3.165pt\relax}

%%%%%%%%%%%%%%%%%%%%%%%%%%%%%%%%%%%%%%%%%%%%%%%%%%%%%%%%%%%%%%%%%%%%%%%%%%%%%
% Special frame types (converted from ExampleSlides.pptx; see README.md)
%
% All five below deliberately reuse as much of standardframe -- and of each
% other -- as possible, rather than reimplementing frame chrome from
% scratch: sectionframe, figureframe, and highlightframe are literally thin
% wrappers around standardframe itself (same title bar, logo, footer, bullet
% styling; only what's documented at each one differs). titleframe and
% quoteframe can't be standardframe wrappers -- both are title-less, and
% quoteframe replaces the whole body layout -- but still reuse its palette,
% fonts, and margins (\gothamFrametitleLeftpading, the institute logo
% drawing command, etc.) rather than inventing new ones.
%%%%%%%%%%%%%%%%%%%%%%%%%%%%%%%%%%%%%%%%%%%%%%%%%%%%%%%%%%%%%%%%%%%%%%%%%%%%%

% titleframe: the deck's own title slide (ExampleSlides.pptx slide 1) --
% typically used exactly once, at the very start of a course's first deck.
% One continuous yellow box -- like standardframe's own title bar -- flush
% with the page's top and left edges, stopping short of the institute logo
% on the right (the same margin already used around the logo elsewhere:
% \gothamLogoWidth + 2*\gothamFrametitleRightpading), holding the title and,
% below it, a second line (author/term/whatever) in the same black text as
% standardframe's own frametitle. Usage:
%
%   \titleframe{Financial Markets}{Stefan Palan\\Summer term 2026}
%
% A *command*, not an environment, taking the second (lower) line's content
% as its second argument -- plain text, use \\ for line breaks same as any
% normal LaTeX paragraph, not an itemize, since the source deck's own
% subtitle block is just stacked lines. This has to be a command rather
% than an environment with the body as its content: beamer's frame body is
% captured by a delimited-argument macro (\beamer@collect@@body, looked up
% by triggering the error below on purpose) hunting for a literal
% \end{frame}, and a beamercolorbox -- needed by an earlier version of this
% to wrap the subtitle text in its own colored band -- opened in a *custom
% environment's* begin code and closed in its end code confuses that hunt
% exactly like the \columns example above (same root cause, different
% environment): with
% \newenvironment{titleframe}{...\begin{beamercolorbox}{...}}{\end{beamercolorbox}...},
% compiling ends in "! File ended while scanning use of \beamer@collect@@body."
% every time, however short the body. Taking the content as a normal brace
% argument instead sidesteps the whole issue. No footer/frame number
% (matches beamer's own convention for title slides, e.g. the old
% \begin{frame}[plain,noframenumbering]\titlepage\end{frame} this replaces --
% see README.md) and no chapter-logo support: both would be meaningless on a
% slide that appears once, before any chapter exists.
%
% Unlike the frametitle hook further up, neither the logo nor the yellow box
% here need the white-patch-and-redraw trick: the box's own width already
% stops short of the logo (rather than being painted full-width and then
% partly painted back over), so there's no yellow left to paint over in the
% first place. Both still need \tikz's "remember picture" on their
% "current page"-anchored nodes, though, and so still need the usual second
% compile pass (see "Title bar and logo" in README.md) -- checked
% empirically: without it, "current page" doesn't resolve to the actual page
% box on the first pass, and the logo/yellow box end up positioned nowhere
% near the page corners they're meant to be anchored to.
%
% \gothamTitleframeBoxHeight is a fixed target for the box's *bottom* edge
% (matching what was actually asked for: "reaching down to where the lower
% grey bar currently ends" in the two-band design this replaces), not the
% blank space above the text -- that space is derived from it instead (box
% height minus the title+subtitle content's own measured height), so the
% bottom edge stays fixed even if a longer title ever wraps onto a second
% line; the blank space above would just shrink to compensate. Measured
% against the old two-band design's own rendered PDF coordinates (before
% this change): its second band's bottom edge sat 187.147pt below the page's
% top edge.
\newlength{\gothamTitleframeBoxHeight}
\setlength{\gothamTitleframeBoxHeight}{187.147pt}
\newlength{\gothamTitleframeBoxWidth}
\setlength{\gothamTitleframeBoxWidth}{\dimexpr\paperwidth-\gothamLogoWidth-2\gothamFrametitleRightpading\relax}
% Text inset from the box's own edges: \gothamFrametitleLeftpading, the same
% length standardframe's own title text uses, plus \gothamTitleframeTextSep
% (2.5ex, matching this command's own former per-band beamercolorbox "sep").
\newlength{\gothamTitleframeTextSep}
\setlength{\gothamTitleframeTextSep}{2.5ex}
\newlength{\gothamTitleframeBandGap}
\setlength{\gothamTitleframeBandGap}{0.8ex}
% Because the leading blank space above the text (further down) is itself
% computed as "box height minus the content's own natural height", the
% content's *bottom* padding is the only thing that determines how far the
% last line sits above the box's bottom edge -- the leading space silently
% absorbs whatever that padding is, keeping the bottom edge fixed either
% way. (The relationship is a direct 1:1 offset, checked empirically by
% comparing two rendered values, but with a nonzero intercept from beamer's
% own default top-of-frame vskip -- there's no clean closed-form for the
% right value, hence the empirical calibration here rather than a formula.)
% Calibrated (against the rendered PDF's own text coordinates, not just
% visually) to about 1.5 line-heights of the subtitle text's own \Large
% size (~20.6pt/line here) above the box's bottom edge; re-measure and
% adjust if the subtitle's font size, or beamer's own default top-of-frame
% vskip, ever change.
\newlength{\gothamTitleframeBottomPad}
\setlength{\gothamTitleframeBottomPad}{41pt}
\newbox\gothamTitleframeBox
\makeatletter
\newcommand{\titleframe}[2]{%
	\begin{frame}[plain,noframenumbering,t]%
	% A bare tikzpicture behaves like an hbox (same reasoning as the footer
	% box's own \vbox-to-0pt wrapper further up): sitting directly in the
	% frame's vertical list, it silently opens an implicit paragraph,
	% switching into horizontal mode. Everything placed after it -- the
	% \vspace* further down included -- would then run in horizontal mode
	% instead of vertical mode, which is what a first version of this command
	% got wrong: with two bare tikzpictures here, \vspace*'s glue ended up as
	% *horizontal* space on an implicit line together with them (visible as
	% the title/subtitle shifting sideways, not down, and the vspace amount
	% appearing to have no effect at all -- it was moving the wrong axis).
	% \vbox to 0pt{\vss ...} keeps each overlay at zero space in the normal
	% flow without ever leaving vertical mode.
	%
	% \gothamArtifactBegin/\gothamArtifactEnd (defined with the footer, "PDF
	% accessibility" comment): this box is purely decorative background, same
	% kind of element as the footer's own yellow box -- marked the same way,
	% for a screen reader to skip. titleframe has no footer at all to worry
	% about separately (it's a "plain" frame, no footline), so this box is
	% the only thing here that needs it; the title/subtitle text right below
	% is real content and stays unmarked.
	\gothamArtifactBegin{/Type /Background}%
	\vbox to 0pt{\vss
	\begin{tikzpicture}[overlay, remember picture]
		% \gothamTitleBoxColor, not a literal color token -- same setting the
		% standard frametitle bar (colorA, above) uses, so the title slide's
		% own box matches every other title bar in the deck (2026-07-16,
		% direct instruction).
		% outer sep=0pt here and on the logo below: without it, tikz's
		% default outer sep insets the fill/logo 0.2pt from the anchored
		% page corner -- a hairline white line above/left of this box, and
		% its right edge/the logo 0.2-0.4pt off every other frame's. See
		% the frametitle hook's own comment for the full story.
		\node[fill=\gothamTitleBoxColor, anchor=north west, inner sep=0pt, outer sep=0pt,
		      minimum width=\gothamTitleframeBoxWidth,
		      minimum height=\gothamTitleframeBoxHeight]
			at (current page.north west) {};
	\end{tikzpicture}%
	}%
	\gothamArtifactEnd
	\vbox to 0pt{\vss
	\begin{tikzpicture}[overlay, remember picture]
		% inner ysep derived, not \gothamFrametitleRightpading -- see the
		% frametitle hook's fixed-position logo redraw for why (matches the
		% inline logo's vertically-centered position on every other frame).
		\node[anchor=north east, inner xsep=\gothamFrametitleRightpading,
		      inner ysep=\dimexpr\gothamTitleBarBottomAlign-\gothamLogoHeight\relax,
		      outer sep=0pt]
			at (current page.north east) {\gothamInstituteLogoSquare};
	\end{tikzpicture}%
	}%
	\setbox\gothamTitleframeBox=\vbox{%
		\hsize=\gothamTitleframeBoxWidth
		\parindent=0pt
		\leftskip=\gothamFrametitleLeftpading
		\advance\leftskip by \gothamTitleframeTextSep
		\rightskip=\gothamFrametitleLeftpading
		\advance\rightskip by \gothamTitleframeTextSep
		% \gothamTitleTextColor, not a literal color token -- follows the
		% box color change right above (2026-07-16, direct instruction):
		% leaving this hardcoded dark would be unreadable once the box
		% itself switches to a dark blue for MCI/Innsbruck.
		\color{\gothamTitleTextColor}%
		% Beamer redefines \\ inside every frame to peek for an overlay spec
		% (\\<2->, for line-by-line reveals in itemize/enumerate); that
		% redefinition assumes it's running in beamer's own frame-body
		% machinery and chokes ("Undefined control sequence", silently
		% swallowing the following word) when \\ is used for a plain line
		% break in a raw \vbox like this one. Restoring plain LaTeX's own
		% \\ (\@normalcr) for this group only fixes it without affecting \\
		% anywhere else in the deck. (The @ here needs catcode 11, which is
		% why this whole command definition -- not just this line -- is
		% wrapped in \makeatletter/\makeatother: macro bodies are tokenized
		% at *definition* time, so a \makeatletter placed inside the body
		% itself would take effect too late, after "\@normalcr" had already
		% been mis-tokenized as "\@" followed by four ordinary-text letters.)
		\let\\\@normalcr
		\vskip\gothamTitleframeTextSep
		\noindent{\Huge\bfseries #1\par}%
		\vskip\gothamTitleframeTextSep
		\vskip\gothamTitleframeBandGap
		\vskip\gothamTitleframeTextSep
		\noindent{\Large #2\par}%
		\vskip\gothamTitleframeBottomPad
	}%
	\vspace*{\dimexpr\gothamTitleframeBoxHeight-\ht\gothamTitleframeBox-\dp\gothamTitleframeBox\relax}%
	\box\gothamTitleframeBox
	% This command's own \color{\gothamTitleTextColor} above (set inside a
	% stored box register) can leak into the *next* frame's body text --
	% see standardframe's own opening code ("Forces this frame's own body")
	% for the real fix and the full story (a trailing reset attempted
	% here first turned out to be a silent no-op, since nothing draws after
	% it within this same frame to actually force the color operator to be
	% emitted).
	\end{frame}%
}
\makeatother

% fullbleedframe: a slide that is *only* a full-screen image (FM2_2026s_ps_deck.pptx
% slide 13, a recruitment flyer pasted in as one embedded picture on an
% otherwise blank layout -- no title, no PPTX text at all) -- no title bar,
% no footer, no logo, matching that source's own blank layout intent, unlike
% figureframe (which is standardframe underneath and always keeps the usual
% title/footer chrome). The image is scaled to fit the slide while keeping
% its own aspect ratio (\includegraphics's keepaspectratio, capped by
% \paperwidth *and* \paperheight together), centered; whatever's left over
% on the binding axis (e.g. a widescreen slide holding a taller/narrower
% image) is filled by a solid black backdrop rather than left showing the
% frame's usual white canvas -- the point of a full-*bleed* image is that
% nothing reads as "empty margin" even when the aspect ratios don't match.
%
%   \fullbleedframe{media/photo-datacamp-recruitment-flyer}
%
% \gothamFullBleedBlack is a fixed HTML black, deliberately *not* a
% \colorlet of c_primary_dark: this frame type's backdrop should stay
% neutral true black regardless of whether c_primary_dark itself is ever
% restyled to something else (e.g. a softer near-black, the way
% c_custom_gray already is elsewhere in this file) for branding
% reasons -- those are two different jobs (brand color vs. letterbox
% backdrop) that only happen to want the same value today.
\definecolor{gothamFullBleedBlack}{HTML}{000000}
\newcommand{\fullbleedframe}[1]{%
	\begin{frame}[plain,noframenumbering,t]%
	% \gothamArtifactBegin/\gothamArtifactEnd (defined with the footer, "PDF
	% accessibility" comment): the black backdrop is purely decorative
	% letterboxing, marked the same way as titleframe's/quoteframe's own
	% background boxes so a screen reader skips it. The image itself is left
	% unmarked -- same as everywhere else in this file, since this project
	% has no alt-text infrastructure yet for imported images generally (a
	% separate, not-yet-made decision -- see README.md's accessibility
	% section for the same open question about the institute logo).
	\gothamArtifactBegin{/Type /Background}%
	\vbox to 0pt{\vss
	\begin{tikzpicture}[overlay, remember picture]
		% outer sep=0pt: see the frametitle hook's comment -- without it the
		% backdrop sits 0.2pt in from the anchored page corner.
		\node[fill=gothamFullBleedBlack, anchor=south west, inner sep=0pt, outer sep=0pt,
		      minimum width=\paperwidth, minimum height=\paperheight]
			at (current page.south west) {};
	\end{tikzpicture}%
	}%
	\gothamArtifactEnd
	\vbox to 0pt{\vss
	\begin{tikzpicture}[overlay, remember picture]
		\node[anchor=center] at (current page.center)
			{\includegraphics[width=\paperwidth, height=\paperheight, keepaspectratio]{#1}};
	\end{tikzpicture}%
	}%
	\end{frame}%
}

% sideimageframe: a standardframe that always carries one auto-placed
% "side filler image" (the user's own term, see
% slides/dfc-Angel-Investors-current.tex's own comment for the technique
% this formalizes into a reusable frame type) -- 2026-07-06, direct
% instruction. The image is positioned flush to three edges at once: its
% top touches the title bar's own bottom edge, its right touches the
% page's own right edge, and its bottom touches the footer box's own top
% edge -- so it's scaled by *height* to exactly
% \paperheight-\gothamTitleBarBottomAlign-\gothamFooterBoxHeight (the
% measured gap between those two), with `keepaspectratio` so the width
% simply follows from the image's own aspect ratio rather than being
% forced/distorted to some second, independent target -- anchoring the
% node at the page's own north-east corner then makes the top+right
% touches automatic, and the height exactly filling that gap makes the
% bottom touch automatic too, with no separate width constraint needed
% (only a *square* source image would ever visibly reach some further
% left boundary; this only guarantees top/right/bottom, matching what was
% actually asked for).
%
%   \begin{sideimageframe}{Title}{media/photo-example}
%   	\begin{itemize}
%   		\item ...
%   	\end{itemize}
%   	\source{...}
%   \end{sideimageframe}
%
% \sidefootnotes is switched on automatically (not left to the caller) --
% the whole point of this frame type is that its footnotes always end up
% rotated along the left edge of the image rather than colliding with it
% the way a full-width bottom footnote would (see
% slides/u29-Implications-of-Capital-Market-Efficiency-Prices-current.tex's
% own comment for exactly that collision, on a plain standardframe). Unlike
% every other \sidefootnotes use elsewhere in this deck, though, the strip
% can't sit at its usual fixed inset from the page's own right edge here --
% that would run it directly underneath the image, the same problem u29
% had. Fixed by *computing* the image's own rendered width (via
% \settowidth against the exact same \includegraphics call used to draw
% it) and using that to override \gothamSideFootnoteRightInset for just
% this one frame -- moving the whole rotated strip to sit immediately left
% of the image instead of at the page's raw edge, entirely automatically,
% with no manual x-coordinate needed from the caller (the direct
% instruction offered a manual-x-coordinate fallback "if it is possible
% [to do automatically], that would be great" -- turned out to be
% possible). \gothamSideImageFrameHeight/\gothamSideImageFrameWidth are
% declared once, globally, and merely *set* (not re-declared) inside the
% environment body, so the frame type can be used more than once per deck;
% the \setlength override of \gothamSideFootnoteRightInset is a plain,
% non-\global assignment, which -- like every other frame-local override in
% this file (e.g. slides/u29-...'s own \setbeamercolor{footnote}{fg=white})
% -- resets on its own once \begin{frame}...\end{frame}'s own TeX group
% closes, with no explicit reset needed.
%
% The extra gap added on top of the image's own measured width (2026-07-06,
% direct instruction: "replicate ... the precise margin we have on regular
% slides between footnote and footer box") was originally an unmeasured
% 4pt guess -- replaced with 3.1pt, back-solved from the deck's own real
% reference gap: measured directly on a plain bottom-footnote slide
% (slides/xhe-Active-vs-Passive-Management-2026-First-current.tex's own
% \source{} line) that the real gap between a footnote's own text and the
% footer box's top edge is 2.93pt, while the original 4pt guess here was
% actually landing the rotated text 3.8pt from the image -- 4pt - (3.8-2.93)
% = 3.13pt, rounded to 3.1pt.
\newlength{\gothamSideImageFrameHeight}
\newlength{\gothamSideImageFrameWidth}
\newenvironment{sideimageframe}[3][t]{%
	\setlength{\gothamSideImageFrameHeight}{\dimexpr\paperheight-\gothamTitleBarBottomAlign-\gothamFooterBoxHeight\relax}%
	\settowidth{\gothamSideImageFrameWidth}{\includegraphics[height=\gothamSideImageFrameHeight, keepaspectratio]{#3}}%
	\sidefootnotes
	\setlength{\gothamSideFootnoteRightInset}{\dimexpr\gothamSideImageFrameWidth+3.1pt\relax}%
	\begin{standardframe}[#1]{#2}%
	\vbox to 0pt{\vss
	\begin{tikzpicture}[overlay, remember picture]
		% outer sep=0pt: see the frametitle hook's comment -- without it the
		% image sits 0.2pt in from the page's right edge (and 0.2pt low).
		\node[anchor=north east, inner sep=0pt, outer sep=0pt] at ($(current page.north east)+(0,-\gothamTitleBarBottomAlign)$)
			{\includegraphics[height=\gothamSideImageFrameHeight, keepaspectratio]{#3}};
	\end{tikzpicture}%
	}%
}{%
	\end{standardframe}%
}

% sectionframe: a section-opening overview (ExampleSlides.pptx slide 3) --
% same layout as standardframe (footer, logo, bullet styling all completely
% unchanged), except the title bar itself is much taller (about a third of
% the page) and set in a bigger font, and the body text is a step bigger
% too, so it reads as "somewhat greater focus" and is unmistakable at a
% glance when flipping through slides -- while still keeping the deck's
% usual yellow-background/black-text look, unlike an earlier version of this
% that inverted to dark-background/yellow-text instead.
%
% The taller bar is NOT a separately-drawn box the way titleframe's is (see
% "Title slide" above): the gotham theme's own frametitle box already sizes
% itself from \gothamFrametitleToppading/\gothamFrametitleBottompading plus
% the title font's own strut height (beamerouterthemegotham.sty, the
% "frametitle" template), so overriding those two lengths locally -- the
% same trick already used below for the font/color overrides -- grows the
% *theme's own* box instead of fighting it with a second one. This also
% means the white-patch-behind-the-logo/chapter-logo machinery (tied to the
% frametitle template via \addtobeamertemplate) keeps working unmodified,
% since it measures the box's actual rendered bottom edge rather than
% assuming any particular height.
%
% The title text is bottom-anchored within the bar -- sitting a fixed "one
% line" above the bar's bottom edge, growing *upward* for a title long
% enough to wrap onto a second line -- rather than top-anchored the way
% gotham's own frametitle template flows by default (which would instead
% leave a large, growing gap *below* a wrapped title, between it and the
% bar's bottom edge). gotham's template has no built-in way to do this: its
% only two knobs are the toppading/bottompading struts, both anchored to
% *one* end of a single top-down paragraph flow. So instead of leaving
% toppading fixed (as an earlier version of this did) the whole title text
% is measured into \gothamSectionframeTitleBox first -- built with the same
% font and (approximately) the same usable width gotham's own template
% gives it, so its \ht+\dp closely track the title's real rendered size,
% including wrapping onto extra lines for a long title -- and toppading is
% then computed as "whatever's left" after that measured height and the
% fixed bottom gap are subtracted from the bar's total target height. A
% longer/wrapped title therefore automatically shrinks toppading (pushing
% the whole title block up, growing upward) rather than changing the gap
% below it. \gothamSectionframeTitleWidth is a slight underestimate of
% gotham's actual per-line width (which is only narrowed by the logo on
% whichever line ends up next to it, typically just the last one) --
% applying that narrower width to *every* line is a deliberate, harmless
% simplification rather than replicating gotham's per-line logic exactly.
%
% The institute logo moves down along with the title now (it's vertically
% centered against the title's own baseline, gotham's instituteLogo
% template), rather than staying pinned where it sits on every other frame
% type, as an earlier version of this kept it -- an intentional change from
% that version: with the title itself now sitting near the bar's bottom
% edge, keeping the logo up at the bar's top would visually separate it
% from the title it's meant to sit beside.
%
% \begingroup/\endgroup (rather than the \global-set-and-reset dance
% \chapterlogo{}/\sidefootnotes need) is enough here: this only ever touches
% \setbeamercolor/\setbeamerfont/\setlength/\setbox, all of which assign
% locally under the hood (checked against beamerbasecolor.sty/
% beamerbasefont.sty directly for the first two -- their key handlers are
% plain \def, not \gdef -- and \setlength/\setbox are plain TeX assignments,
% local by default), so a plain surrounding TeX group reverts them
% correctly on its own.
%
% \gothamSectionframeContentDrop adds a little extra room below the bar
% before the body content starts, on top of standardframe's own
% \gothamContentTopShift pull-up (applied inside \begin{standardframe}
% already, before this code regains control). All lengths here calibrated
% against the rendered PDF's own coordinates; re-measure and adjust if the
% title/body font sizes change.
\newlength{\gothamSectionframeContentDrop}
\setlength{\gothamSectionframeContentDrop}{10pt}
\newlength{\gothamSectionframeBoxHeight}
\setlength{\gothamSectionframeBoxHeight}{85.04pt}
% \gothamSectionframeBottomGap is the *visual* target -- one huge-font line
% height (baselineskip, measured directly from a deliberately 2-line test
% title: the gap between the two lines' own rendered top edges). It is NOT
% used directly as \gothamFrametitleBottompading: gotham's bottom strut
% (\gotham@frametitlestrut@end) consistently produces a visual gap about
% 5.58pt *smaller* than the padding value fed into it (checked empirically,
% same offset for both a one-line and a two-line test title, so it's a
% fixed per-box offset rather than something that scales with content) --
% \gothamSectionframeBottomPadCorrection adds that back.
\newlength{\gothamSectionframeBottomGap}
\setlength{\gothamSectionframeBottomGap}{14.32pt}
\newlength{\gothamSectionframeBottomPadCorrection}
\setlength{\gothamSectionframeBottomPadCorrection}{5.58pt}
% Similarly, the box's total rendered height comes out consistently ~11.4pt
% *less* than toppading+measuredTitleHeight+bottompading would suggest --
% the measured title box (built as a plain \vbox{...\par}) doesn't reproduce
% gotham's own leading strut (\gotham@frametitlestrut@start, a zero-width
% rule with an explicit height rather than natural interline glue) exactly,
% and this correction closes that gap. Checked the same way as the bottom
% one above: consistent across both the one-line and two-line test titles.
\newlength{\gothamSectionframeHeightCorrection}
\setlength{\gothamSectionframeHeightCorrection}{12.1pt}
\newlength{\gothamSectionframeTitleWidth}
\setlength{\gothamSectionframeTitleWidth}{\dimexpr\paperwidth-\gothamFrametitleLeftpading-\gothamLogoWidth-2\gothamFrametitleRightpading\relax}
\newbox\gothamSectionframeTitleBox
\newenvironment{sectionframe}[2][t]{%
	\begingroup
	% Disable hyphenation for the title -- both here, in the measurement
	% box below, and (since this is in effect for the rest of the group,
	% which \begin{standardframe} -- triggering the real \insertframetitle
	% -- runs inside) in the actual rendered title too, so a line break
	% never lands mid-word. Setting the penalties alone is the standard,
	% sufficient way to do this (TeX treats a penalty of 10000 at a
	% breakpoint as forbidden, not just heavily discouraged, so it's never
	% chosen even when the alternative is an overfull line); the hyphenmin
	% lengths are redundant insurance, not strictly needed on top of that.
	\hyphenpenalty=10000
	\exhyphenpenalty=10000
	\lefthyphenmin=100
	\righthyphenmin=100
	\setbeamerfont{frametitle}{size=\huge}
	\setbox\gothamSectionframeTitleBox=\vbox{%
		\hsize=\gothamSectionframeTitleWidth
		\usebeamerfont{frametitle}\usebeamercolor[fg]{frametitle}%
		\noindent #2\par
	}%
	\setlength{\gothamFrametitleBottompading}{\dimexpr\gothamSectionframeBottomGap+\gothamSectionframeBottomPadCorrection\relax}
	\setlength{\gothamFrametitleToppading}{\dimexpr\gothamSectionframeBoxHeight-\gothamFrametitleBottompading-\ht\gothamSectionframeTitleBox-\dp\gothamSectionframeTitleBox+\gothamSectionframeHeightCorrection\relax}
	\global\def\gothamFrametitleFullHeightYellow{1}%
	\global\def\gothamFrametitleLogoFixed{1}%
	\setbeamertemplate{frametitle}[gotham_no_inline_logo]%
	\begin{standardframe}[#1]{#2}%
	\vspace*{\gothamSectionframeContentDrop}%
	\large
}{%
	\end{standardframe}%
	\endgroup
}

% quoteframe: a quote slide (ExampleSlides.pptx slide 4) -- institute logo
% as usual, and the quote itself set large/italic/black on a yellow
% full-bleed block, same style as titleframe/sectionframe's own title boxes
% (flush top-left, stopping short of the logo on the right, same margin as
% elsewhere) rather than the dark block an earlier version of this used. The
% usual footer still shows (unlike titleframe, which drops it -- a title
% slide is a one-off outside any particular course content, while a quote
% is still one slide *among* a deck's others, worth being able to reference
% by its footer slide ID like any other). Opening/closing curly quotes are
% added automatically (write the quote text with none of its own); the
% second argument is the person it's from, set smaller/plain, below the
% yellow box and right-aligned to its own right edge (not the page's usual
% text margin) -- pass an empty second argument for an unattributed quote.
%
%   \quoteframe{One of the funny things about the stock market is that
%   	every time one person buys, another sells, and both think they
%   	are astute.}{William Feather, 1889-1981}
%
% \gothamQuoteBoxHeight is a plain percentage of the page, unlike
% titleframe's/sectionframe's own (calibrated-to-land-on-an-exact-fraction)
% box heights, since nothing here needs to line up with another element's
% measured position the way those do. \gothamTitleframeBoxWidth (defined
% above) is reused as-is for the box's width -- the same
% stop-short-of-the-logo formula, so no reason to duplicate it under a new
% name.
%
% Every piece here -- box, logo, quote text, attribution -- is drawn as a
% zero-space tikz overlay anchored to absolute page coordinates (the same
% \vbox to 0pt{\vss ...} pattern used elsewhere in this file for exactly
% this reason), rather than any of it flowing through the frame's normal
% vertical list. This sidesteps needing to know/counteract beamer's own
% default top-of-frame spacing (which standardframe explicitly measures and
% cancels via \gothamContentTopShift) -- there's no such frame-relative
% positioning to get right here, since literally everything is placed by
% explicit coordinate instead. Like the other page-corner overlays in this
% file, "current page" doesn't resolve correctly on a document's very first
% compile pass, so (like those) this needs the usual second pass too.
%
% A *command*, not an environment -- see the identically-reasoned comment on
% titleframe just above for why: the same beamercolorbox-spanning-a-custom-
% environment's-begin/end problem would apply here too, for the same reason.
\newlength{\gothamQuoteBoxHeight}
\setlength{\gothamQuoteBoxHeight}{0.8\paperheight}
\newlength{\gothamQuoteAttributionGap}
\setlength{\gothamQuoteAttributionGap}{1.8ex}
% \@empty (the empty-attribution check below) needs \makeatletter.
% \quotesource{text} -- opt-in \source{} for \quoteframe (call *before*
% \quoteframe{}{}, same set-then-consume pattern as \chapterlogo{}
% above), needed because \quoteframe is a *command* that auto-closes its
% own \end{frame} -- a bare \source{...} typed after \quoteframe{}{} in
% the calling slide file (as if it were still inside the frame, the way
% \source{} reads everywhere else in this deck) is actually issued
% *outside* any frame, and \source{}'s own \footnote{} then has no page to
% attach to yet. Confirmed directly (2026-07-12, direct instruction:
% traced a spurious extra blank leading overlay step on
% slides/c4c-Macbeth-Spot-Removers-Proposed-Structure-current.tex -- the
% very next multi-overlay frame after
% slides/w8w-Pie-Model-of-Capital-Structure-current.tex's own bare
% \source{} -- down to that orphaned footnote specifically, via a reduced
% two-slide test deck): the dangling footnote insert doesn't render
% anywhere visible on w8w itself, but silently corrupts beamer's own
% overlay-counting pre-pass on whichever frame happens to come right
% after -- reproduced with w8w+c4c alone (bug present) vs. c4c alone (bug
% absent), and confirmed it's specifically about the footnote being
% orphaned, not about w8w's content otherwise, by removing just its
% trailing \source{} line (bug also disappears). A plain \clearpage
% between the two frames also "fixes" c4c but only by giving the orphaned
% footnote a page of its own to attach to -- confirmed directly this
% splits w8w itself into two pages (real content, then a second page with
% nothing but the footnote/footer), not an acceptable trade. The only
% correct fix is keeping the footnote inside the frame that's actually
% citing something, i.e. before \quoteframe's own \end{frame} -- exactly
% what \quotesource{} below does.
\newcommand{\gothamQuoteSource}{}
\newcommand{\quotesource}[1]{\gdef\gothamQuoteSource{#1}}
\makeatletter
\newcommand{\quoteframe}[2]{%
	\begin{frame}[t]%
	% \gothamArtifactBegin/\gothamArtifactEnd (defined with the footer, "PDF
	% accessibility" comment): same reasoning as titleframe's own background
	% box above -- purely decorative, marked so a screen reader skips it. The
	% quote text and attribution below stay unmarked (real content); the
	% footer (this frame keeps one, unlike titleframe) is already marked via
	% the global footline hook.
	\gothamArtifactBegin{/Type /Background}%
	\vbox to 0pt{\vss
	\begin{tikzpicture}[overlay, remember picture]
		% \gothamTitleBoxColor, not a literal color token -- same setting
		% titleframe's own box and the standard frametitle bar (colorA) use
		% (2026-07-16, direct instruction).
		% outer sep=0pt here and on the logo/attribution below: same hairline-
		% inset-from-the-anchored-corner issue as titleframe's own box -- see
		% the frametitle hook's comment for the full story.
		\node[fill=\gothamTitleBoxColor, anchor=north west, inner sep=0pt, outer sep=0pt,
		      minimum width=\gothamTitleframeBoxWidth,
		      minimum height=\gothamQuoteBoxHeight]
			at (current page.north west) {};
	\end{tikzpicture}%
	}%
	\gothamArtifactEnd
	\vbox to 0pt{\vss
	\begin{tikzpicture}[overlay, remember picture]
		% inner ysep derived, not \gothamFrametitleRightpading -- see the
		% frametitle hook's fixed-position logo redraw for why (matches the
		% inline logo's vertically-centered position on every other frame).
		\node[anchor=north east, inner xsep=\gothamFrametitleRightpading,
		      inner ysep=\dimexpr\gothamTitleBarBottomAlign-\gothamLogoHeight\relax,
		      outer sep=0pt]
			at (current page.north east) {\gothamInstituteLogoSquare};
	\end{tikzpicture}%
	}%
	\vbox to 0pt{\vss
	\begin{tikzpicture}[overlay, remember picture]
		% \gothamTitleTextColor, not a literal color token -- follows the box
		% color change right above (2026-07-16, direct instruction). Set
		% *inside* the node's own {...} content group (not a stored box
		% register, unlike titleframe's own vbox above), so this doesn't hit
		% the same "\color leaks out of a stored \vbox" pitfall documented on
		% \titleframe's own closing code -- confirmed directly by checking
		% the frame right after this one for the same symptom.
		\node[anchor=center, inner sep=0pt,
		      text width=\dimexpr\gothamTitleframeBoxWidth-2\gothamFrametitleLeftpading\relax]
			at ($(current page.north west)+(0.5\gothamTitleframeBoxWidth,-0.5\gothamQuoteBoxHeight)$)
			{\color{\gothamTitleTextColor}\Large\itshape\textquotedblleft #1\textquotedblright};
	\end{tikzpicture}%
	}%
	\def\gothamQuoteAttributionTemp{#2}%
	\ifx\gothamQuoteAttributionTemp\@empty\else
		\vbox to 0pt{\vss
		\begin{tikzpicture}[overlay, remember picture]
			\node[anchor=north east, inner sep=0pt, outer sep=0pt]
				at ($(current page.north west)+(\gothamTitleframeBoxWidth,-\gothamQuoteBoxHeight-\gothamQuoteAttributionGap)$)
				{\normalfont\itshape\color{c_custom_gray} #2};
		\end{tikzpicture}%
		}%
	\fi
	\ifx\gothamQuoteSource\@empty\else
		\source{\gothamQuoteSource}%
		\global\let\gothamQuoteSource\@empty
	\fi
	\end{frame}%
}
\makeatother

% figureframe: a slide that is just a figure (ExampleSlides.pptx slide 5) --
% same standardframe layout (title, footer, \source{} for the usual
% "Source: ..."/"Image source: ..." line), but with the figure positioned
% via \figurecenter (call before the figure, inside the frame body) rather
% than left-aligned, which -- like the source deck's own figure slide -- is
% the default:
%
%   \begin{figureframe}{Correlation between DJIA stock returns over 15 years}
%   	\includegraphics[width=0.5\linewidth]{correlation-matrix}
%   	\source{Source: Own illustration, data from Yahoo! Finance, 24.02.2026.}
%   \end{figureframe}
%
%   \begin{figureframe}{Title}
%   	\figurecenter
%   	\includegraphics[width=0.6\linewidth]{diagram}
%   \end{figureframe}
%
% Vertical placement is standardframe's own [t]/[c] (its existing optional
% argument, passed straight through as figureframe's own -- so
% \begin{figureframe}[c]{Title} vertically centers exactly like
% \begin{standardframe}[c]{Title} already does); \figurecenter only ever
% affects the horizontal axis.
%
% \figurecenter is plain \centering under a more discoverable name, nothing
% more -- *not* an opt-in-before-the-frame flag like \chapterlogo{}/
% \sidefootnotes (those must run before \begin{standardframe}, since they
% configure something the frame's own opening code reads). This is called
% *inside* the body instead, exactly where it takes effect, so it can just
% take effect immediately like any other formatting declaration -- no flag,
% no state to reset afterwards, and (checked by first writing it as a flag
% read in figureframe's own opening code, mirroring \sidefootnotes) no way
% for the frame's opening code to see it anyway, since that code runs
% before the body -- and \figurecenter -- does.
\newcommand{\figurecenter}{\centering}
\newenvironment{figureframe}[2][t]{%
	\begin{standardframe}[#1]{#2}%
}{%
	\end{standardframe}%
}

% highlightframe: a standardframe with two things about it made unmistakably
% different -- the title bar's *and* the footer's background color, both
% together -- so a presenter can never flip to it without immediately
% noticing (ExampleSlides.pptx slide 6 was the inspiration for wanting a
% standardframe variant here, though the source slide itself has no such
% color cue; this is a new addition, not a conversion). Everything else --
% body, logo, bullet styling -- is exactly standardframe's own, via the same
% \begingroup/\endgroup wrapper sectionframe above uses for the title-color
% half of this. \gothamHighlightColor is a \colorlet, not a color baked
% directly into the environment, so retargeting every highlightframe in a
% deck to a different color -- title bar and footer both -- is still the
% same one-line change it was when this only touched the title bar.
%
% The footer box needs the \global\def dance (\gothamFooterBoxColor, defined
% above with the frametitle hook's other opt-in flags, reset in
% standardframe's own closing code) that \setbeamercolor{frametitle} doesn't:
% that assigns through beamer's own color-management system, which handles
% the frame-closing group correctly on its own (checked -- see
% \gothamChapterLogoPath's comment for why plain macros need \global\def but
% beamer's own set-commands don't), but the footer box's fill is read from a
% plain macro inside a tikz key, same category as \gothamChapterLogoPath.
\colorlet{gothamHighlightColor}{c_primary_gray}
\newenvironment{highlightframe}[2][t]{%
	\begingroup
	\setbeamercolor{frametitle}{bg=gothamHighlightColor}
	\global\def\gothamFooterBoxColor{gothamHighlightColor}%
	\begin{standardframe}[#1]{#2}%
}{%
	\end{standardframe}%
	\endgroup
}


, similar to \sidefootnotes or
% \setfootnotesymbols) switches that compile into the "handout" variant --
% including, as of 2026-07-03, actually forcing beamer's own "handout" mode
% (\gdef\beamer@currentmode{handout}, normally only reachable via the
% \documentclass[handout]{beamer} *class option*, set once at class-loading
% time -- too early for a command like \studentversion, called later from a
% deck's own preamble, to use that route). This was checked, not assumed:
% until this point \studentversion only ever touched \omitslide{}, so
% nothing suppressed \pause/overlay specs -- true only by coincidence, since
% no slide converted so far happens to use them. Late-redefining
% \beamer@currentmode was verified (isolated minimal-beamer tests, plain
% \pause *and* inline specs like \item<2->) to reproduce the same collapse
% -to-one-page, show-the-final-revealed-state behavior as the class option,
% even set well after \documentclass{beamer} -- beamer's overlay decoder
% (beamerbasedecode.sty) reads \beamer@currentmode at frame-processing time,
% not just once at class-option time, so this is safe, not a hack that
% happens to work by luck of file ordering.
% \omitslide{}/\omitslideversion{} are drop-in replacements for
% \slide{}/\slideversion{} that tag a slide as excluded from that variant:
%
%   \omitslide{xyz-Answer-key}                 % instead of \slide{...}
%   \omitslideversion{xyz-Answer-key}{1-0-0}    % instead of \slideversion{...}{...}
%
% In a \studentversion compile, a tagged slide's \input{} is skipped
% entirely -- no frame is produced for it at all (not a blank page, not a
% numbering gap: beamer's frame/slide counters simply never advance for a
% call that never ran). In every other (i.e. normal/instructor) compile, a
% tagged slide loads exactly as it would via \slide{}/\slideversion{}, with
% one difference: \gothamOmittedSlide is set for the duration of that one
% \input{}, so the footer template below can print a marker next to that
% frame's slide ID -- a deliberately visible reminder, while reviewing the
% instructor deck, of which slides won't make it to students. The same tag
% works for both compiles because it lives at the call site in the *deck*
% file (not inside the slide's own content file): the identical slide list
% can serve both the instructor deck and the student handout, e.g.
%
%   % Slidemaster preamble variant for University of Graz decks. Institution-
% specific content only (colors, logo, logo calibration) -- every frame
% type, the deck-version system, footer/hyperlink/footnote styling, tikz
% styles, etc. live in Slidemaster_preamble_base.tex, \input{}ed at the
% bottom of this file, and are shared verbatim across every institution
% this project supports (see Slidemaster_preamble_mci.tex and
% Slidemaster_preamble_uni-ibk.tex for the other two). A deck picks its
% institution by \input{}ing this file instead of one of those, nothing
% else changes.

% Color definition (Graz palette, sourced from the original single-preamble
% version of this project -- these are the University of Graz's own official
% colors, not derived/approximated).
\definecolor{c_primary_dark}{HTML}{000000}
\definecolor{c_primary_highlight}{HTML}{ffd500}
\definecolor{c_primary_gray}{HTML}{c6c6c6}
\definecolor{c_primary_white}{HTML}{FFFFFF}
\definecolor{c_secondary_violet}{HTML}{948dc3}
\definecolor{c_tertiary_violet}{HTML}{c2bddf}
\definecolor{c_secondary_blue}{HTML}{689cce}
\definecolor{c_tertiary_blue}{HTML}{bacce8}
\definecolor{c_secondary_green}{HTML}{9fcb81}
\definecolor{c_tertiary_green}{HTML}{cce2ba}
\definecolor{c_secondary_red}{HTML}{e4684b}
\definecolor{c_tertiary_red}{HTML}{f1ae95}
\definecolor{c_secondary_orange}{HTML}{f6a968}
\definecolor{c_tertiary_orange}{HTML}{fbd0aa}
\definecolor{c_secondary_yellow}{HTML}{ffe063}
\definecolor{c_tertiary_yellow}{HTML}{ffeeaf}
\definecolor{c_tertiary_gray}{HTML}{e3e3e3}
% Not an official Uni Graz color (hence "custom", not "secondary" like the
% palette above) -- a softer near-black than c_primary_dark, used for
% quoteframe's full-bleed dark block, matching the specific dark gray (not
% pure black) ExampleSlides.pptx used for that same block. Also reused for
% hyperlink text color, for the same reason it works for quoteframe's
% attribution text: dark enough for good contrast/legibility on a light
% background.
\definecolor{c_custom_gray}{HTML}{333333}
\definecolor{c_custom_red}{HTML}{C00000}
\definecolor{c_custom_green}{HTML}{00B050}

% Logo + calibration.
% schemes/scheme_uni_graz/Logo_Uni-Graz_4c-cropped.pdf -- official vector artwork,
% user-supplied (schemes/scheme_uni_graz/Logo_Uni-Graz_4c.eps, kept alongside as the
% untouched source), converted to PDF via epstopdf and cropped to its true
% ink bounding box (computed with Ghostscript's own `-sDEVICE=bbox`, same
% treatment as MCI's and Innsbruck's own logos -- see
% Slidemaster_preamble_mci.tex's own comment for the fuller explanation of
% why the EPS's own %%BoundingBox isn't trusted directly). Vector, not
% raster, so it stays crisp at any zoom. Supersedes this file's own earlier
% raster logo (schemes/scheme_uni_graz/Logo Uni Graz 4c-cropped.png, itself cropped
% from schemes/scheme_uni_graz/Logo Uni Graz 4c.png's own symmetric transparent
% margin -- both kept alongside as historical/fallback assets, not deleted).
\newcommand{\gothamInstituteLogoPath}{schemes/scheme_uni_graz/Logo_Uni-Graz_4c-cropped.pdf}
\newlength{\gothamLogoHeight}
\setlength{\gothamLogoHeight}{3ex}
% Re-measured against the compiled PDF for this specific logo image -- see
% Slidemaster_preamble_base.tex's own comment on \gothamFrametitleRightpading
% for why this needs re-measuring per institution rather than reusing
% another one's value.
\setlength{\gothamFrametitleRightpading}{0.28cm}

% Itemize bullet color (2026-07-16, direct instruction): Graz keeps the
% project's original behavior, bullets in the primary highlight color.
\newcommand{\gothamBulletColor}{c_primary_highlight}

% Title bar and footer box/text colors (2026-07-16, direct instruction):
% Graz keeps the project's original behavior throughout -- title/footer
% boxes in the primary highlight color, title/footer text in the primary
% dark color (unlike MCI/Innsbruck, which switch these to their own
% secondary blue/white -- see Slidemaster_preamble_mci.tex's own comment).
\newcommand{\gothamTitleBoxColor}{c_primary_highlight}
\newcommand{\gothamTitleTextColor}{c_primary_dark}
\newcommand{\gothamFooterBoxColorDefault}{c_primary_highlight}
\newcommand{\gothamFooterTextColor}{c_primary_dark}

\usepackage{etoolbox} % \AtEndEnvironment, used for the itemize level-drop spacing below
\usepackage{tikz} % For painting a white box over title bar
\usetikzlibrary{calc}
\usetikzlibrary{arrows.meta} % for e.g. arr/.style={-{Stealth[...]}...} on diagram slides
\usetikzlibrary{positioning} % for `right=<dist> of <node>`-style relative node placement
\usetikzlibrary{decorations.text} % text-along-a-path, for a circular photo's curved caption
\usepackage{booktabs} % \toprule/\midrule/\bottomrule for gothamtable
\usepackage{colortbl} % \arrayrulecolor{}, used by gothamtable to color those rules
\usepackage{pifont} % \ding{51}/\ding{55} checkmark/X-mark glyphs, e.g. slides/fr3-Forms-of-Capital-Market-Efficiency-Strong-Semistrong-Weak-current.tex's own validity ratings
\usepackage{pgfplots} % real-data line/bar charts (e.g. slides/3vp-Metallgesellschaft-Oil-Price-Chart-current.tex) -- confirmed safe alongside this deck's own \tableofcontents, see CONVERSION_NOTES.md's "beamer/TeX gotchas" entry
\pgfplotsset{compat=1.18}

% Restores genuine Computer Modern math fonts. Found 2026-07-15, direct
% instruction ("math equations do not have serifs, look at vcv"): despite
% fontspec being loaded with `no-math` (beamerfontthemegotham.sty), math
% still wasn't landing on Computer Modern -- checked directly against a
% compiled PDF's own per-glyph font names (pymupdf), not just eyeballed:
% digits/math-italic letters were coming out as FiraSans-Light(Italic)
% (this deck's own body-text sans font) and operators/parens as CMSS10
% (Computer Modern's *sans-serif* shape, not its serif one). Root cause:
% beamer's own base classes redefine the math alphabets to follow the
% presentation's sans-serif text font by default, specifically so math
% matches an all-sans deck -- `no-math` only stops *fontspec* from doing
% this, it has no say over beamer's own separate override. Beamer ships
% the standard fix for exactly this case: a font theme named "serif"
% with an "onlymath" option that restores ordinary (Latin
% Modern/Computer-Modern-equivalent) math fonts without touching any of
% the surrounding sans-serif text. Confirmed directly afterward: every
% glyph in a re-checked formula now reports as LMRoman*/CMR10/CMMI8 (all
% genuinely serif), not FiraSans/CMSS10.
\usefonttheme[onlymath]{serif}

% \text{} (amsmath, used throughout for plain-word operands inside math,
% e.g. \text{E}[\tilde{R}_M]) is a separate mechanism from the math
% alphabets \usefonttheme[onlymath]{serif} above just fixed: unlike
% \mathrm/\mathit/etc., \text{} is defined to typeset in whatever the
% *surrounding running text* font is at that point in the document, not a
% math alphabet -- so on this deck's all-Fira body text, \text{} inside a
% formula still comes out in Fira Sans (and picks up italic if the
% enclosing text happens to be italicized), even though the digits/
% operators around it are correctly Computer Modern. Found 2026-07-20 on
% slides/2gd-Analyzing-Reward-to-Risk-Ratio-CAPM-current.tex's own
% \text{E}/\text{Cov} operands. Fixed generally, project-wide, by wrapping
% \text{}'s argument in an explicit upright-Computer-Modern font switch --
% \usefont{OT1}{cmr}{m}{n} rather than \mathrm{} (which would pull in
% \text{}'s own math-alphabet substitution logic unnecessarily) or
% \textrm{} (which switches *family* only, so still comes out italic
% inside an italicized context, which \text{} should never do here) --
% deliberately hardcoding "not italic" (shape n) rather than inheriting
% the ambient shape, since \text{}'s whole purpose in this deck's formulas
% is plain upright labels (E, Cov, "the" CAPM, ...), regardless of
% whatever font context happens to surround the formula itself.
\let\gothamOrigText\text
\renewcommand{\text}[1]{\gothamOrigText{\usefont{OT1}{cmr}{m}{n}#1}}

% Color definition, institute logo path, and the two logo-calibration
% lengths (\gothamLogoHeight, \gothamFrametitleRightpading) all live in the
% per-university preamble that \input{}s this file, not here -- this file
% is shared machinery only, no institution-specific branding. Every color
% token below (\c_primary_*, \c_secondary_*, \c_tertiary_*, \c_custom_*) and
% \gothamInstituteLogoPath/\gothamLogoHeight/\gothamFrametitleRightpading
% are assumed already defined by the time this file is reached -- see e.g.
% Slidemaster_preamble_uni-graz.tex for the Graz values, or
% Slidemaster_preamble_mci.tex/Slidemaster_preamble_uni-ibk.tex for the
% other two institutions this deck supports. \gothamLogoWidth is the one
% logo-derived length that *does* stay computed here (via \settowidth,
% further below) rather than hardcoded per institution, since it already
% adapts automatically to whatever image \gothamInstituteLogoPath points at.

% Style hyperlinks (\href{url}{text}) consistently across all slides:
% underlined, so a link is visually recognizable as such against plain
% body text -- yellow (this project's usual accent color) was tried first
% for the *text* itself, but yellow-on-white body text is hard to read,
% unlike yellow used as an accent against dark/black text elsewhere in the
% theme. Went through a few more candidates on 2026-07-06 (all direct
% instruction) before settling on c_secondary_blue for both text
% and underline (a real color reads more clearly as "this is a link" than
% any gray does, without yellow's legibility problem).
%
% 2026-07-07, direct instruction: split text color and underline color
% instead of using one color for both -- link text back to
% c_custom_gray (the near-black dark gray also used for
% quoteframe's attribution text), underline in a separate accent color.
% Plain \underline can't do this: its rule is drawn in whatever color is
% current at that point, so a single \underline{...} call always
% underlines in its own text's color, never a second, different one.
% \gothamHrefUnderline below reimplements it instead: measure the linked
% text's own width into a savebox, print that box in customgray, then
% back up by the box's own width and draw a plain \rule of the same
% width, raised to sit just below the text's baseline, in the accent
% color. Same single-line limitation as plain \underline (and the reason
% \url{} is used instead of \href on slides/acw-Some-Topics-in-Testing-the-EMH-current.tex --
% see that slide's own comment): a \rule, like \underline, can't span a
% line break either, so this only ever wraps normal, comparatively short
% \href display text, never something that needs mid-word breaking.
%
% First tried c_primary_highlight for the underline (same day) -- confirmed
% via the compiled PDF that the rule *was* being drawn (right color,
% position, width), so not a compile-time bug, but yellow (very high
% luminance, close to white) at this rule's 0.4pt thickness reads as
% essentially invisible against the white slide background, worst at
% footnote size (see e.g. slides/za2-Startup-Financing-Categorization-current.tex's
% own long inline URL) where it's easy to miss entirely. Switched
% 2026-07-07 (direct instruction) to c_secondary_blue instead --
% far more contrast against white at the same thickness -- and the rule
% was also thickened at the same time (0.4pt -> 0.6pt), on the reasoning
% that it might need it for the same contrast reason. It didn't: blue
% already reads clearly at the original thickness, and 0.6pt instead
% just looked heavier than intended (direct feedback 2026-07-08, e.g. on
% slides/6qi-Yield-Spreads-Corporate-Bonds-current.tex). Reverted to
% 0.4pt -- only the color was ever actually asked to change.
\hypersetup{colorlinks=true, urlcolor=c_custom_gray, linkcolor=c_custom_gray}
\let\gothamOrigHref\href
\newsavebox{\gothamHrefTextBox}
%
% 2026-07-08 (direct instruction): the underline's own rule thickness now
% scales with the ambient font size instead of a single fixed 0.2pt for
% every link everywhere -- 0.2pt at this deck's footnote size, 0.4pt at its
% normal body-text size, interpolated (and extrapolated, for any other
% size a link happens to appear at) linearly in between.
%
% Deriving "current font size" from \f@size (the textbook NFSS way) was
% tried first and confirmed broken specifically *inside* a beamer frame --
% \typeout{\meaning\f@size} right after this preamble loads shows a normal,
% defined macro (`macro:->10.95`), but the exact same check inside
% `\begin{standardframe}...\end{standardframe}` reports `\f@size` as
% *undefined*, presumably from beamer's own overlay-aware font-switching
% machinery not keeping classic NFSS's own current-size bookkeeping in
% sync the way plain (non-beamer) LaTeX does. Since every real use of this
% macro happens inside a frame, that path was a dead end regardless of how
% well it worked in isolation.
%
% Measuring a reference glyph's own rendered height instead sidesteps this
% entirely: building a box and reading back its own height just measures
% whatever font is *actually* active at the call site, independent of
% whether beamer's own internal "current size" bookkeeping happens to
% agree -- and within one font
% family, every glyph's metrics scale by exactly the same ratio the point
% size does, so a measured height is just as valid a stand-in for "current
% size" as \f@size would have been, had it worked. "M" (present in every
% context this is used in) is the reference glyph; \gothamHrefUnderlineFnHeight/
% \gothamHrefUnderlineBodyHeight are its height at this deck's footnote font
% and normal body font respectively, each measured once, here, in the
% relevant font context.
% Deliberately built from raw TeX primitives (\newdimen/\setbox/\ht),
% *not* \newlength/\setlength/\settoheight -- found, while building this,
% that loading this deck's theme silently turns \setlength and
% \settoheight into something else entirely: `\meaning\setlength` reads
% `macro:->\protect \setlength` right after `\usetheme{gotham}` (still the
% ordinary kernel definition beforehand), and the same for `\settoheight`.
% Not chased to which file does this or why (bisected only as far as
% "plain beamer alone doesn't, beamer+gotham does" -- most likely
% `translator`, which beamer loads for localized strings and is known to
% robustify commands this way, but not confirmed) -- what matters
% practically is that a *global* assignment through either one silently
% no-ops instead of erroring, confirmed directly and in isolation (a
% minimal beamer+gotham file, no other content): `\global\setlength{\len}
% {\otherlen}` leaves `\len` at 0pt after the group closes, while the
% exact same pattern with plain `\newdimen`/raw `=` assignment propagates
% correctly. Every plain (non-global) `\setlength` elsewhere in this
% project's own files is unaffected -- it's specifically the *global*
% form that breaks, which is why this is the first place it's shown up
% despite `\setlength` being used constantly throughout this preamble.
\newdimen\gothamHrefUnderlineFnHeight
\newdimen\gothamHrefUnderlineBodyHeight
\newdimen\gothamHrefUnderlineNowHeight
\newdimen\gothamHrefUnderlineThickness
% A dedicated box register, not the commonly-borrowed box0 -- several
% slides in this deck (the "Long-term bonds" tables) already use box0
% themselves (`\sbox0{...}`/`\usebox0`) for their own, unrelated
% mid-table-layout purposes, and a link can appear anywhere, including
% inside one of those tables -- reusing box0 here would silently clobber
% whatever they were holding there.
\newbox\gothamHrefUnderlineBox
% The two reference heights are (re)measured on every single \href call,
% not once at preamble-load time -- tried once-at-load first and found a
% second, independent problem with it: `\usebeamerfont{footnote}` right
% here (before `\begin{document}`) silently has no effect at all (footnote-
% size "M" and normal-size "M" measured identically, confirmed directly),
% presumably because beamer's font-theme machinery isn't fully installed
% until the document starts. Since every real `\href` call happens well
% inside the document (inside some frame), measuring fresh each time
% sidesteps that timing question entirely rather than needing to find the
% right hook to measure once, early enough to be free but late enough to
% be correct -- the extra cost is two throwaway one-glyph boxes per link,
% not something worth optimizing away in a slide deck.
\newcommand{\gothamHrefUnderlineMeasureRefHeights}{%
	{\usebeamerfont{footnote}\setbox\gothamHrefUnderlineBox=\hbox{M}\global\gothamHrefUnderlineFnHeight=\ht\gothamHrefUnderlineBox}%
	{\normalsize\setbox\gothamHrefUnderlineBox=\hbox{M}\global\gothamHrefUnderlineBodyHeight=\ht\gothamHrefUnderlineBox}%
}
% A plain e-TeX \dimexpr can multiply a dimension by another dimension and
% then divide by a third, as one combined operation, as long as the
% *middle* operand of the chain is itself wrapped in its own nested
% \dimexpr...\relax -- confirmed directly (an expl3 \dim_set:Nn attempt
% first, using \seq_.../\fp_... the way \setfootnotesymbols above does,
% kept hitting "Missing ) inserted for expression"/type-mismatch errors:
% dividing one dimension by another is a *ratio*, not a dimension, and
% \dim_set:Nn's own expression grammar doesn't implicitly convert between
% the two the way this nested-\dimexpr form does). No expl3 needed here at
% all -- plain e-TeX arithmetic, already available in any current LaTeX.
\newcommand{\gothamHrefUnderlineSetThickness}{%
	\gothamHrefUnderlineThickness=\dimexpr 0.2pt +
		(\gothamHrefUnderlineNowHeight - \gothamHrefUnderlineFnHeight)
		* \dimexpr0.2pt\relax
		/ (\gothamHrefUnderlineBodyHeight - \gothamHrefUnderlineFnHeight)\relax
}
\newcommand{\gothamHrefUnderline}[1]{%
	\sbox{\gothamHrefTextBox}{#1}%
	\textcolor{c_custom_gray}{\usebox{\gothamHrefTextBox}}%
	\hspace{-\wd\gothamHrefTextBox}%
	\gothamHrefUnderlineMeasureRefHeights
	\setbox\gothamHrefUnderlineBox=\hbox{M}\gothamHrefUnderlineNowHeight=\ht\gothamHrefUnderlineBox
	\gothamHrefUnderlineSetThickness
	\raisebox{-0.35ex}[0pt][0pt]{\textcolor{c_secondary_blue}{\rule{\wd\gothamHrefTextBox}{\gothamHrefUnderlineThickness}}}%
}
\renewcommand{\href}[2]{\gothamOrigHref{#1}{\gothamHrefUnderline{#2}}}
% \url{} (the `url` package hyperref depends on/patches) defaults to a
% monospace-ish "tt" style regardless of the above -- used deliberately in
% only one place in this deck,
% slides/acw-Some-Topics-in-Testing-the-EMH-current.tex, specifically
% *for* its line-breaking-at-slashes behavior (a plain \href with the
% URL as its own display text has no break points at all -- see that
% slide's own comment). Every other same-text-as-target URL in the deck
% (including slides/bb6-A-Perpetual-Bond-current.tex's own video link,
% which used to be a \url{} too) uses plain \href{url}{url} instead, since
% \url{} doesn't pick up this deck's custom colored underline below (its
% breakable content can't be measured into the savebox the underline
% needs -- see that mechanism's own comment) and most of these URLs are
% short enough to not need the breaking \url{} exists for anyway. 2026-07-07,
% direct instruction: that one genuinely-needed \url{} (on s-acw) should
% still look like this deck's other hyperlinks, i.e. the surrounding body
% font, not \url{}'s own default monospace.
% `\urlstyle{same}` (a standard url-package switch) does exactly that,
% with no effect on \url{}'s own breaking behavior. Set deck-wide, not
% only on that one slide, so any future \url{} use matches by default
% too, rather than needing to remember this each time. Coloring already
% matches automatically (hyperref's own \url respects
% \hypersetup{urlcolor=...} the same way \href does); the underline
% \href gets above doesn't carry over, since \underline's own argument
% can't contain line breaks -- wrapping \url{}'s necessarily-breakable
% text in \underline would silently undo the breaking this exists for,
% so that one difference is accepted rather than "fixed."
\urlstyle{same}

% gotham's own footnote color (beamercolorthemegotham.sty:
% \setbeamercolor{footnote}{fg=normal~text.fg!90}) is a 90%-black gray --
% against plain black body text, customgray links contrast clearly (found
% while investigating a 2026-07-05 report that a footnote's own \href
% didn't read as a link the way body-text ones like the "Moodle" link on
% slides/hxr-Administrative-Details-current.tex do); against gotham's own
% dimmed footnote gray, the same customgray is only 25/255 lighter and
% barely distinguishable at footnote size, with only the underline giving
% it away. Overridden to plain black here so a footnote link contrasts the
% same way a body-text one does -- a one-line override of the theme's own
% color (not editing beamercolorthemegotham.sty itself, which stays
% untouched per its own upstream-file status), and a plausible general
% legibility improvement for footnote text regardless of links.
\setbeamercolor{footnote}{fg=c_primary_dark}

% Footer text color (\FooterText/slide id/frame number) -- \gothamFooterTextColor,
% set by the per-university preamble before \input{}ing this file (the
% footer box's own background color, \gothamFooterBoxColorDefault above, is
% a separate, independently configurable setting). No explicit
% \setbeamercolor{footline}{...} existed before this -- the footer text
% simply inherited whatever "normal text" foreground already was (black),
% which happened to read fine against every frame's own yellow/orange
% footer box so far, but wouldn't automatically read against an arbitrary
% future footer color (e.g. a dark blue) without this override.
\setbeamercolor{footline}{fg=\gothamFooterTextColor}

\colorlet{colorDark}{c_primary_dark} % FG in light/normal mode
% Title bar background -- \gothamTitleBoxColor, set by the per-university
% preamble before \input{}ing this file (same idiom as
% \gothamBulletColor/\gothamFooterBoxColorDefault above).
\colorlet{colorA}{\gothamTitleBoxColor} % frametitle, standin.out

\gothamset{background=light} % Needs to go after color redefinitions, so they are applied
\setbeamercolor{background canvas}{bg=white}

% gotham's own default behavior on \section{} (used for a deck's chapter
% grouping/table of contents) is to auto-insert its *own* plain, unstyled
% section-divider frame plus a small "current section highlighted in the
% outline" frame, every time -- redundant with (and visually inconsistent
% next to) this project's own hand-styled sectionframe (see README.md's
% "Special frame types" section), which every chapter in this deck already
% opens with as an explicit \slide{}. Disabled so \section{} only feeds
% \tableofcontents (see slides/nkj-Agenda-current.tex) without also
% inserting its own extra frames.
\gothamset{sectionframe default=false}

% Same reasoning, for \subsection{} (2026-07-15, direct instruction):
% gotham's default subsection-divider frame is disabled too, so
% \subsection{} calls still register in the PDF's own bookmark/outline
% structure (Acrobat's "chapters" pane -- a plain PDF sectioning feature,
% independent of whether beamer auto-inserts a frame for it) and still feed
% \tableofcontents (see slides/nkj-Agenda-current.tex's own
% `\tableofcontents[hideallsubsections]`, which then hides them from that
% one slide specifically), without also spawning their own unstyled divider
% slide in the middle of a chapter every time one is added.
\gothamset{subsectionframe default=false}

% Title text color -- \gothamTitleTextColor, set by the per-university
% preamble before \input{}ing this file (the title bar's own background
% color, \gothamTitleBoxColor/colorA above, is a separate, independently
% configurable setting).
\setbeamercolor{frametitle}{fg=\gothamTitleTextColor}

% A color leak from \titleframe/\quoteframe into whichever frame follows
% them is fixed in standardframe's own opening code below (search
% "Forces this frame's own body") rather than here -- an
% \AtBeginEnvironment{frame}{\color{c_primary_dark}} hook was tried first
% and confirmed, by isolated reproduction, to have *no effect at all*:
% beamer's own "frame" isn't a plain \newenvironment etoolbox's
% \AtBeginEnvironment can patch that way, so the hook silently never fired.

% Set content text margins
\setbeamersize{text margin left=.45cm, text margin right=.45cm}

%\colorlet{colorPale}{gPaleYell} % BG in light/normal mode
%\colorlet{colorDark}{gDarkBlack} % FG in light/normal mode
%\colorlet{colorA}{gAnthracite} % frametitle, standin.out,
%\colorlet{colorAreversed}{gLightTeal} % frametitle, standin.in,
%\colorlet{colorB}{gMidGrey} % gray BG : progress bar, blocks
%\colorlet{colorC}{gDeepYellOr} % progress bar
%\colorlet{colorD}{gLightOrange} % alert
%\colorlet{colorE}{gLightGreen} % example

% Opt-in yellow border for an image placed via a tikz node (e.g.
% \node[gothamImageBorder] {\includegraphics{...}};) -- matches the
% 3pt/accent1-yellow outline PowerPoint draws directly on a picture shape's
% own border (<a:ln w="38100"><a:solidFill><a:schemeClr val="accent1"/>,
% where accent1 turned out to already be this project's own
% c_primary_highlight, not a new color to map). A tikz *style* (a named
% bundle of node keys), not a wrapper command: adding or removing it is
% just adding or deleting one word from that one \node's own option list,
% no separate macro call/argument needed -- add `gothamImageBorder` to draw
% the border, leave it off for a plain image, same as any other tikz key.
%
% Thickness is a single named length, \gothamImageBorderWidth, not a bare
% "3pt" baked directly into the style (2026-07-15, direct instruction: "make
% sure the yellow image borders have consistent thickness throughout the
% project by defining a single variable") -- confirmed directly, by
% measuring the actual PDF content-stream "w" operator, that every
% gothamImageBorder use already rendered at the exact same 3pt regardless
% (a genuinely shared style was already in effect, no numeric drift found),
% so this is a refactor for future adjustability/discoverability, not a fix
% for an actual inconsistency -- any perceived thickness difference between
% slides traces to something else instead (e.g. a border straddling flush
% image content on one slide and not another, which reads as "heavier" even
% at an identical stroke width -- see the gothamImageBorder padding gotcha
% below).
\newlength{\gothamImageBorderWidth}
\setlength{\gothamImageBorderWidth}{2.5pt}
\tikzset{gothamImageBorder/.style={draw=c_primary_highlight, line width=\gothamImageBorderWidth}}

% The "padding gotcha" promised above, finally root-caused (2026-07-23,
% background-task follow-up to the IA2 fine-tuning batch #2 anomaly on
% slides/m2k-Exercise-CAPM-2-current.tex/
% slides/n8e-Tail-Risk-Variance-Clustering-current.tex/
% slides/fm6-Tail-Risk-Market-vs-Simulated-current.tex/
% slides/xqq-Fat-Tails-current.tex, all fixed there with a pragmatic
% `line width=1.2pt` override and an explicitly *unresolved* root cause at
% the time -- "declares an identical stroke width and CTM in the compiled
% PDF's own raw content stream, yet renders at a genuinely different pixel
% thickness", with CTM scale/clips/dash/cap/join all directly checked and
% ruled out). Confirmed via a controlled reproducer (four `gothamImageBorder`
% nodes -- plain opaque image, nested pgfplots axis, plain image at
% `inner sep=3pt`, nested plain tikzpicture at `inner sep=3pt` -- rendered
% at 600dpi and measured by a half-max cross-section profile, not just
% eyeballed): the one variable that actually predicts the rendered width is
% `inner sep`, not content type at all. At `inner sep=0pt` a *fully opaque*
% image measured 1.2pt; the exact same image at `inner sep=3pt` measured
% 2.52pt (matching the declared 2.5pt); a nested pgfplots axis and a plain
% nested tikzpicture, both at `inner sep=3pt`, also measured 2.52pt. Two
% more points confirmed it's a smooth transition, not a binary switch:
% `inner sep=0.6pt` measured 1.8pt (= the 1.2pt floor plus that 0.6pt gap,
% exactly); `inner sep=1.25pt` (= \gothamImageBorderWidth/2, the stroke's
% own half-width) already measured the full 2.52pt.
%
% The mechanism: tikz draws a node's border stroke centered on the node's
% own boundary path, then paints the node's *content* on top of it,
% afterwards. At `inner sep=0`, the content's own bounding box coincides
% exactly with that boundary path, so if the content is opaque all the way
% to its own edge, it paints directly over the stroke's inner half,
% visually leaving only the outer half (`\gothamImageBorderWidth`/2,
% ~1.2pt) -- not a rendering bug, just tikz's normal draw order interacting
% with zero padding. A gap of at least `\gothamImageBorderWidth`/2 (or
% content that's transparent right at its own edge -- see below) leaves the
% border with nothing opaque to cover it, so it renders at its true
% declared width instead.
%
% This also explains why the fix landed on exactly the 4 slides it did,
% each for a different reason: slides/m2k-Exercise-CAPM-2-current.tex
% used `inner sep=3pt` for unrelated layout reasons (room for the axis's
% own tick labels), well past the 1.25pt threshold, so nothing ever covered
% its border. slides/n8e-Tail-Risk-Variance-Clustering-current.tex/
% slides/fm6-Tail-Risk-Market-vs-Simulated-current.tex/
% slides/xqq-Fat-Tails-current.tex all use `inner sep=0pt` -- matching the
% deck's own overwhelming convention -- but their own chart PNGs
% (`chart-index-daily-returns-*.png`/`chart-index-fat-tails-histogram*.png`,
% matplotlib exports with a transparent background) measured alpha=0 at
% every edge pixel (confirmed directly, Pillow), so there was no opaque
% content to cover their own border either. Every other `gothamImageBorder`
% use in this project that was never flagged uses an opaque image (a JPG,
% or a PNG with no transparent margin) at `inner sep=0pt` -- which
% *accidentally* renders at the ~1.2pt covered-half width instead of the
% full 2.5pt this constant nominally declares, and that accidental width is
% what reads as "the deck's normal border thickness" throughout, since it's
% by far the most common case.
%
% Practical upshot for a *new* slide: a `gothamImageBorder` node with
% `inner sep=0pt` wrapping opaque content (the common case) needs no
% special handling -- it will match the rest of the deck automatically. A
% node that instead uses a nonzero `inner sep`, or wraps content with a
% transparent edge (a matplotlib export, an SVG-derived PDF with padding,
% ...), will render at the *full* \gothamImageBorderWidth and look
% noticeably heavier next to the deck's other borders -- either add an
% explicit `line width=\gothamImageBorderWidth/2` (matching the 4 fixed
% slides' own literal 1.2pt) to match the accidental deck-wide look, or
% reduce `inner sep` below \gothamImageBorderWidth/2 so the same coverage
% effect kicks in naturally. The 4 existing per-slide overrides above are
% left as-is -- they're already the correct fix for what's now a fully
% understood cause, not a mystery workaround needing revisiting.

% Reusable "mugshot" -- a circular-cropped photo, bordered like any other
% gothamImageBorder image, with a caption curving along its own bottom arc
% (2026-07-06, direct instruction, generalized from
% slides/4uq-Definition-of-Capital-Market-Efficiency-current.tex's own
% Eugene Fama photo -- see that slide's earlier comment for the source
% PPTX effect this reproduces, `<a:prstGeom prst="ellipse">` +
% `<a:prstTxWarp prst="textArchDown">`, and CONVERSION_NOTES.md's own
% "beamer/TeX gotchas" entry, search "circular photo crop", for the
% clip/border/decorations.text mechanics themselves -- unchanged here,
% just factored out so a second photo doesn't need to repeat them).
%
%   \coordinate (myMugshotCenter) at (...);
%   \gothamMugshot{myMugshotCenter}{media/photo-example.jpeg}{2.6cm}{Some Person, NP 2013}
%
% #1 the *name* of a coordinate the caller has already placed (not a raw
% point expression -- needed as a plain name, not just a point, since the
% caption's own arc path re-references it by name via `++(...)`, which
% only works with a named coordinate/node, not an arbitrary `($...$)`
% calc expression); #2 the image path; #3 the desired final diameter;
% #4 the caption text.
%
% Scales the source photo so it fully covers the circle with no gap
% anywhere -- both dimensions end up at least as big as the diameter --
% while using as much of the source photo as possible (2026-07-07,
% direct instruction: "should not include whitespace... scale the photo
% such that its edges just touch the circle's inner border"). This
% briefly (same day) scaled the *other* way -- by the photo's own
% diagonal, the largest scale at which the *entire* rectangle fits
% inside the circle with nothing clipped at all -- which does avoid
% ever clipping real content, but leaves the circle's own four corners
% uncovered (visible as plain background there, since a non-circular
% photo can only touch an inscribing circle at its corners, not along
% the whole rim) -- exactly the "whitespace" this round's feedback
% ruled out in favor of full coverage. Covering scale is
% `diameter / min(w,h)`: whichever of the photo's own two dimensions is
% *smaller* becomes exactly the diameter (that edge pair touches the
% circle's own inner border, per the direct instruction), and the other,
% larger dimension overhangs the diameter and gets clipped by the
% circular `\clip` -- for a portrait photo (every one used here so far),
% that means width becomes the diameter and a sliver clips off the top
% and bottom, which is the origin of the clipping problem the diagonal
% approach above existed to solve in the first place. Reintroduced
% deliberately: showing as much of the photo as the "no whitespace"
% constraint allows is now the stated priority, over never clipping
% anything.
%
% Measuring the photo's own native size needs `\savebox` (into
% \gothamMugshotBox, declared once below, not a fresh box register per
% call) -- `\includegraphics` with no size option at all renders at the
% image's native resolution, and `\wd`/`\ht` then read that box's actual
% dimensions. pgfmath (already loaded for tikz) converts any TeX
% dimension fed into it to a plain pt-count number automatically, so
% arithmetic on the two measured lengths just works. `min(1,ratio)`
% (ratio = height/width) picks out whichever of the two is smaller
% *relative to the other*, without needing a `sqrt`/`^` at all this time
% -- so this version doesn't hit the "! Dimension too large" pgfmath
% overflow the diagonal version's own arithmetic did (see this file's
% own comment history / CONVERSION_NOTES.md's "beamer/TeX gotchas" for
% that point (easy to misread as a harmless warning, since compilation
% continues and produces *a* number regardless). Reformulating via the
% aspect ratio sidesteps the whole problem: `scale = diameter /
% sqrt(1+ratio^2)` is the same formula algebraically (factor `w` out of
% `sqrt(w^2+h^2)` and it becomes `w*sqrt(1+(h/w)^2)`), but every
% the full story of that overflow, from when this macro's own arithmetic
% did use a `sqrt`).
%
% Caption font reverted to this deck's normal body font (2026-07-07,
% direct instruction: "change the mugshot template back to using regular
% font instead of monospace") -- `\ttfamily` (Fira Mono) was a direct
% instruction on 2026-07-06, one day earlier; \normalfont resets
% explicitly rather than just omitting \ttfamily, so this doesn't
% silently inherit whatever font happens to be active at the call site.
\newsavebox{\gothamMugshotBox}
\newcommand{\gothamMugshot}[4]{%
	\savebox{\gothamMugshotBox}{\includegraphics{#2}}%
	\pgfmathsetmacro{\gothamMugshotImgW}{\wd\gothamMugshotBox}%
	\pgfmathsetmacro{\gothamMugshotImgH}{\ht\gothamMugshotBox}%
	\pgfmathsetmacro{\gothamMugshotRatio}{\gothamMugshotImgH/\gothamMugshotImgW}%
	\pgfmathsetmacro{\gothamMugshotScaledW}{(#3)/min(1,\gothamMugshotRatio)}%
	\begin{scope}
		\clip (#1) circle (0.5*#3);
		\node[anchor=center, inner sep=0pt] at (#1)
			{\includegraphics[width=\gothamMugshotScaledW pt]{#2}};
	\end{scope}
	\draw[gothamImageBorder] (#1) circle (0.5*#3);
	{\normalfont\tiny
	\path[decoration={text along path, text={#4}, text align=center}, decorate]
		(#1) ++(200:0.5*#3+0.25cm) arc (200:340:0.5*#3+0.25cm);
	}
}

% Reusable "operating/cash conversion cycle" bar diagram (FM2_2026s_ps_deck.pptx
% slides 22/24/25/26 all embed the exact same picture -- an EMF, not a native
% shape/table -- so rather than convert that EMF and re-embed it 4 times over,
% it's redrawn once natively here and reused, matching this project's own
% "clean/geometric diagrams get redrawn, not extracted" policy). Three
% stacked bars (Inventory period+Accounts receivable period on top, Accounts
% payable period+Cash conversion cycle in the middle, Operating cycle
% spanning the full width on the bottom) plus 4 tick labels below at
% x=0/(1/3)/(2/3)/1. Geometry/colors measured directly from the source
% deck's own rendered PNG via pixel analysis (PIL run-length color
% detection), not eyeballed: the two internal split points are exactly 1/3
% and 2/3 of the total width (measured 0.3339/0.6672 against a plain
% PowerPoint illustrative diagram, not slide-23's own specific-day example --
% close enough to exact thirds, and a plausible round number for the
% presenter to have actually used, that thirds is treated as the real value
% rather than curve-fitting noise). Fill colors matched to this project's
% *existing* palette rather than the source's own literal RGB, per this
% deck's established emphasis-color policy: source yellow was already an
% exact match for `c_primary_highlight`; source green (RGB 146,208,80) and
% purple (RGB 128,100,162) were close enough to `c_secondary_green`
% and `c_secondary_violet` respectively to reuse rather than define
% two new one-off colors.
%
% Drawn at a fixed "natural" size (10cm wide) with its own top-left corner
% at the calling tikzpicture's origin -- wrap a call in
% \begin{scope}[shift={(x,y)}, scale=s]\gothamCashCycleBars\end{scope} to
% place/resize it, the same convention as this deck's other reusable
% diagram pieces (e.g. slides/bdh-Working-Capital-Example-current.tex's
% circular flow diagram, resized the same way after the fact per user
% feedback) -- resizing later without needing to re-derive any coordinates.
\newcommand{\gothamCashCycleBars}{%
	\fill[c_primary_highlight] (0,0) rectangle (6.667,-0.42);
	\fill[c_secondary_green] (6.667,0) rectangle (10,-0.42);
	\node[align=center, font=\tiny] at (3.333,-0.21) {Inventory period};
	\node[align=center, font=\tiny] at (8.333,-0.21) {Accounts receivable period};
	\fill[c_secondary_green] (0,-0.47) rectangle (3.333,-0.89);
	\fill[c_primary_highlight] (3.333,-0.47) rectangle (10,-0.89);
	\node[align=center, font=\tiny] at (1.667,-0.68) {Accounts payable period};
	\node[align=center, font=\tiny] at (6.667,-0.68) {Cash conversion cycle};
	\fill[c_secondary_violet] (0,-0.94) rectangle (10,-1.36);
	\node[align=center, font=\tiny, text=c_primary_white] at (5,-1.15) {Operating cycle};
	\foreach \gccx in {0,3.333,6.667,10} {
		\draw[c_primary_dark] (\gccx,-1.36) -- (\gccx,-1.46);
	}
	% text width on all 4, narrow enough to force an actual 2-line wrap
	% (not just cap an already-short line): without it, "Raw material
	% purchase" and "Payment for raw materials" (and, symmetrically, the
	% last two) run into each other -- confirmed directly by rendering.
	% A *wide* text width (e.g. 3.1cm, comfortably more than either label's
	% own one-line width) does NOT fix this: \node's text width only forces
	% a wrap once the text exceeds it, so a generous width just lets both
	% neighbors keep their full one-line extent and still collide -- this
	% was tried first and confirmed not to work by rendering, not assumed.
	% 1.5cm is deliberately narrower than any of the 4 would take on one
	% line, guaranteeing the wrap for all of them.
	%
	% All 4 use anchor=north (centered on their own tick, like the middle
	% two always did) -- the outer two originally used anchor=north
	% west/east instead (flush with the tick rather than centered on it),
	% which looked inconsistent with the middle two once seen side by side
	% at actual size. Centering the outer two on their own tick (rather
	% than flush against the diagram's own left/right edge) does push their
	% wrapped text a bit further left/right past the diagram's own
	% bounding box -- checked this stays within the page margin at every
	% scale/position this macro is actually called at in this deck (the
	% narrowest margin case leaves ~0.7cm of clearance), not just assumed
	% safe from the numbers alone.
	\node[anchor=north, font=\tiny, align=center, text width=1.5cm] at (0,-1.5) {Raw material purchase};
	\node[anchor=north, font=\tiny, align=center, text width=1.5cm] at (3.333,-1.5) {Payment for raw materials};
	\node[anchor=north, font=\tiny, align=center, text width=1.5cm] at (6.667,-1.5) {Sale of finished goods};
	\node[anchor=north, font=\tiny, align=center, text width=1.5cm] at (10,-1.5) {Cash collected on sales};
}

% Reusable circular flow diagram (Cash -> Raw materials -> Finished goods ->
% Receivables -> Cash), first built for
% slides/bdh-Working-Capital-Example-current.tex and reused (smaller) on
% slides/qq5-Cash-Conversion-Cycle-current.tex -- factored out here once a
% second slide needed the same picture, same reasoning as
% \gothamCashCycleBars above. Args: #1 diagram center coordinate, #2 radius
% (box-center distance from that center, a bare cm number, no unit suffix),
% #3 the box `scale` factor (a plain node `scale` transform, not a bigger
% font -- see bdh's own comment, preserved below, for why: this diagram was
% resized once already per user feedback, and `scale` was what correctly
% grew the padding/rounded-corners/arrow proportions together rather than
% just the text). Arrow length/width scale with #3 too, computed via
% \pgfmathsetmacro since tikz option values can't do arithmetic inline.
%
% [Preserved from bdh's own resize, 2026-07-06: box size grew by a uniform
% ~1.44x there (compared via the isolated "Cash" box, the one corner with no
% arrow touching it to contaminate the measurement: 0.94cm to 1.36cm wide,
% 0.50cm to 0.74cm tall), and the box-to-box radius grew by the same
% ~1.44-1.45x (1.15cm to ~1.65cm) -- both measured from the user's own
% "Clipboard 1.jpg" reference via pixel analysis, not eyeballed, and
% consistent with a single uniform resize of the whole group.]
% #1 is a *complete* coordinate spec, exactly as it would appear after
% "\coordinate (name) at " -- e.g. "(current page.north east)" or
% "($(current page.north east)+(-3.81cm,-4.61cm)$)", dollar signs and all
% when it's a calc expression. Resolved into a plain \coordinate first
% rather than substituted directly into each box's own ($...$) calc
% expression, since a calc expression can't contain another calc
% expression's $...$ delimiters nested inside it.
\newcommand{\gothamCircularFlowDiagram}[3]{%
	\pgfmathsetmacro{\gothamCFDarrowLen}{5*#3}
	\pgfmathsetmacro{\gothamCFDarrowWidth}{1.3*#3}
	\begin{scope}
		\coordinate (gothamCFDctr) at #1;
		\node[fill=c_primary_highlight, rounded corners=1pt, inner sep=4pt, align=center, font=\footnotesize, scale=#3] (gothamCFDcash) at ($(gothamCFDctr)+(0,#2)$) {Cash};
		\node[fill=c_primary_highlight, rounded corners=1pt, inner sep=4pt, align=center, font=\footnotesize, scale=#3] (gothamCFDraw) at ($(gothamCFDctr)+(#2,0)$) {Raw materials};
		\node[fill=c_primary_highlight, rounded corners=1pt, inner sep=4pt, align=center, font=\footnotesize, scale=#3] (gothamCFDfin) at ($(gothamCFDctr)+(0,-#2)$) {Finished goods};
		\node[fill=c_primary_highlight, rounded corners=1pt, inner sep=4pt, align=center, font=\footnotesize, scale=#3] (gothamCFDrec) at ($(gothamCFDctr)+(-#2,0)$) {Receivables};
		\draw[-{Stealth[length=\gothamCFDarrowLen pt]}, line width=\gothamCFDarrowWidth pt, c_primary_highlight] (gothamCFDcash) to[bend left=18] (gothamCFDraw);
		\draw[-{Stealth[length=\gothamCFDarrowLen pt]}, line width=\gothamCFDarrowWidth pt, c_primary_highlight] (gothamCFDraw) to[bend left=18] (gothamCFDfin);
		\draw[-{Stealth[length=\gothamCFDarrowLen pt]}, line width=\gothamCFDarrowWidth pt, c_primary_highlight] (gothamCFDfin) to[bend left=18] (gothamCFDrec);
		\draw[-{Stealth[length=\gothamCFDarrowLen pt]}, line width=\gothamCFDarrowWidth pt, c_primary_highlight] (gothamCFDrec) to[bend left=18] (gothamCFDcash);
	\end{scope}
}

% A \frac/\dfrac-stacked formula reads as cramped against a plain-text
% neighbor in an itemize -- it's visibly taller than a normal line (both
% above *and* below the line, since the fraction bar sits on the text
% baseline with numerator/denominator stacked around it), so both the
% fraction's own item *and* the item that follows it need a bit of extra
% headroom that a plain-text \item doesn't. Started as a manually-repeated
% \vspace{3pt} before each such \item (first on
% slides/u74-Example-Dynamic-Mattress-Company-Periods-current.tex);
% \tallitem below is that same convention factored into a single drop-in
% \item replacement, so a future fraction-heavy slide (or a change to the
% spacing amount) doesn't mean re-editing several slide files by hand.
% Takes an optional leading multiplier (default 1, i.e. 3pt): \tallitem[1.5]
% for 4.5pt, for a formula tall enough that the default doesn't read as
% enough room. Usage: replace \item with \tallitem both on any itemize line
% containing a stacked fraction *and* on the item immediately following it
% (same nesting level; skip this second replacement if that next item is
% itself already a \tallitem for its own fraction, or if the fraction item
% is the last one in its list) -- it's a plain substitute, not an
% environment, so it doesn't change how the surrounding list itself is
% structured. (Named \gothamFracItem originally; renamed 2026-07-21 to the
% more discoverable \tallitem, and given the optional multiplier and the
% following-item rule above at the same time -- no behavior change for any
% existing default-strength call site.)
% Ratio, not a frozen length (2026-07-22, direct instruction: "modify
% \tallitem and \tallitemnobullet to also use scaling relative to font
% size") -- the original \newlength+\setlength{...}{3pt} pair freezes to
% one fixed pt value the moment \setlength runs, the same issue (and same
% fix) as \gothamItemsepRatioI/II/III/IV further below: converted to a
% bare ratio macro, with the "em" unit written back on inside \tallitem's
% and \tallitemnobullet's own bodies (not here), so it's recomputed fresh
% against whatever font is actually live at each call site -- correct
% regardless of nesting level or a \footnotesize/\scriptsize override,
% unlike the old fixed 3pt (calibrated against one specific level's font,
% silently wrong at any other). 0.3 reproduces the original 3pt exactly
% at level 2's own font (10pt) -- the level \tallitem was first used at,
% on slides/u74-Example-Dynamic-Mattress-Company-Periods-current.tex (see
% this macro's own comment above) -- so no visual change there, while now
% correctly scaling up/down at every other level and font override.
\newcommand{\gothamTallItemBaseRatio}{0.3}
\newlength{\gothamTallItemBaseUnit}
\newlength{\gothamTallItemSpace}
\newcommand{\tallitem}[1][1]{%
	\setlength{\gothamTallItemBaseUnit}{\gothamTallItemBaseRatio em}%
	\setlength{\gothamTallItemSpace}{#1\gothamTallItemBaseUnit}%
	\vspace{\gothamTallItemSpace}\item}

% For a line that needs \tallitem's own extra headroom but no bullet at
% all -- a display equation continuing the previous \item's own sentence,
% rather than a new bulleted point of its own (see README.md's own
% \tallitem section, which already documents this exact combination as a
% rare, previously hand-written case: "\vspace{\gothamTallItemBaseUnit}\item[]").
% \tallitemnobullet (2026-07-21, direct instruction) factors that hand-
% written idiom into the same kind of drop-in replacement \tallitem
% already is for its own case, same optional multiplier: \tallitemnobullet[1.5].
\newcommand{\tallitemnobullet}[1][1]{%
	% Recomputes \gothamTallItemBaseUnit itself, fresh, rather than reusing
	% whatever \tallitem's own body last left it at -- since it's now a
	% per-call, live "em"-based value (see \tallitem's own comment above),
	% not the one-time preamble-load-time constant it used to be, relying
	% on call order/a preceding \tallitem to have set it correctly would
	% silently break the moment \tallitemnobullet is used on its own, or
	% at a different font size than whichever \tallitem last ran.
	\setlength{\gothamTallItemBaseUnit}{\gothamTallItemBaseRatio em}%
	\setlength{\gothamTallItemSpace}{#1\gothamTallItemBaseUnit}%
	\vspace{\gothamTallItemSpace}\item[]}

% flalign* (used to left-align a display equation inside an itemize --
% see CONVERSION_NOTES.md's "flalign* for left-aligning..." entry) inherits
% plain LaTeX's \abovedisplayskip/\belowdisplayskip, sized for a full
% document page, not a compact slide itemize -- reads as far too much blank
% space above/below the equation next to the surrounding item text (found
% 2026-07-21, direct instruction; most visible on
% slides/bwi-Pearson-Correlation-Definition-current.tex's own single-line
% equation, but the same oversized gap is present everywhere flalign* is
% used this way, e.g. slides/7se-Portfolio-Risk-Three-Assets-current.tex
% and slides/rb3-Asset-Allocation-Example-CAL-current.tex). \gothamTightDisplaySkips
% tightens both skips -- and their "short" variants, since amsmath picks
% between the two depending on how much text precedes the display, so
% setting only one pair left the gap unfixed whenever the preceding line
% happened to be short -- down to something that reads as part of the
% same list rather than a separate floating display. Usage: wrap the
% *real*, unmodified flalign* in a plain TeX group together with this
% command, `{\gothamTightDisplaySkips\begin{flalign*} ... \end{flalign*}}`,
% rather than a new same-named environment -- tried that first
% (`\newenvironment{gothamtightflalign*}{...\begin{flalign*}}{\end{flalign*}}`)
% and it broke LaTeX's own environment-nesting tracking (`\@currenvir`),
% surfacing as "\begin{flalign*} ... ended by \end{gothamtightflalign*}"
% and a fatal, PDF-less compile on the very first slide that used it. A
% plain group sidesteps that class of bug entirely since it never
% introduces a same-named wrapper environment.
\newlength{\gothamTightDisplaySkip}
\setlength{\gothamTightDisplaySkip}{3pt}
\newcommand{\gothamTightDisplaySkips}{%
	\abovedisplayskip=\gothamTightDisplaySkip
	\abovedisplayshortskip=\gothamTightDisplaySkip
	\belowdisplayskip=\gothamTightDisplaySkip
	\belowdisplayshortskip=\gothamTightDisplaySkip}

% Set logo
% \gothamInstituteLogoPath (a plain \newcommand storing an image path,
% same idiom as \gothamChapterLogoPath below) and \gothamLogoHeight (a
% length) are both defined by the per-university preamble that \input{}s
% this file, before this file is reached -- see e.g.
% Slidemaster_preamble_uni-graz.tex. Whichever source image
% \gothamInstituteLogoPath points at should have any symmetric margin of
% empty transparent pixels around the actual logo artwork trimmed off first
% (see Slidemaster_preamble_uni-graz.tex's own comment on why: an untrimmed
% margin makes \includegraphics's box bigger than the visible logo, so the
% logo's right-hand margin ends up looking bigger than its top/bottom
% margins, since only the latter get absorbed into vertical centering).
% \gothamLogoHeight is the single source of truth for the logo's rendered
% size: it drives the (invisible) in-flow logo below, the white patch's
% width further down, and the visible redrawn logo, so those three can never
% drift out of sync with each other again (that drift -- redrawing a smaller
% logo than the one the patch/position were sized for -- is what caused the
% left/right margins around the logo to come out uneven).

% The tikz node wrapper (remember picture, but not overlay, so it still
% behaves as normal inline content) records this logo's exact page position
% under the name "gothamLogoAnchor", so it can be redrawn on top of a white
% patch further down (see the "yellow box shorter than the page" block below).
\renewcommand{\gothamInstituteLogoSquare}[1][\gothamLogoHeight]{%
	\tikz[remember picture, baseline=(gothamLogoAnchor.base)]{%
		\node[inner sep=0pt, outer sep=0pt](gothamLogoAnchor){%
			\includegraphics[height=#1]{\gothamInstituteLogoPath}%
		};%
	}%
}

% \gothamFrametitleRightpading (left/right padding around the logo, chosen
% to match its top/bottom margins -- measured against the title colorbox's
% own top and bottom edges -- rather than the theme's 2.2ex default) is
% also set by the per-university preamble, not here: the top/bottom margins
% are a side effect of vertically centering a \gothamLogoHeight-tall logo in
% the title bar, not something set directly, so this value has to be
% re-measured and set per institution whenever its own logo/height differs
% from another's.

% Logo width, used below to size the white patch behind it.
\newlength{\gothamLogoWidth}
\settowidth{\gothamLogoWidth}{\includegraphics[height=\gothamLogoHeight]{\gothamInstituteLogoPath}}

% Let the yellow title bar visually stop short of the page's right edge, so
% the logo sits on white instead of on yellow: paint a white patch spanning
% the full height of the title bar behind the logo, then redraw the logo
% (already positioned correctly by \gothamFrametitleRightpading above, via
% the "gothamLogoAnchor" node captured above) on top of it. This needs two
% compilation passes: like any \tikz "remember picture", the node positions
% are only available in the .aux file written by the previous run.
% (The redraw uses \includegraphics directly rather than
% \usebeamertemplate{instituteLogo}/{includeLogo}: those templates rely on
% gotham's internal property-list lookup, which is only valid in the
% context of the original frametitle rendering and silently produces
% nothing when re-invoked here.)
% Additionally paints white across the *entire* width of the strip between
% the logo's bottom edge and the title bar's actual bottom edge, so the
% yellow ends flush with the logo's bottom (rather than extending further
% down under the title text too) while the box's width and top position
% stay unchanged. Checked against measured PDF coordinates that the title
% text's own (font-metric, not just this-string's-glyphs) lowest descender
% still clears the logo's bottom edge by about 1.3pt, so this doesn't clip
% text even on titles with descenders (e.g. "Highlighting...").
%
% Factored into its own macro, rather than written directly into the
% \addtobeamertemplate call below, so \gotham_no_inline_logo (a whole
% separate frametitle template variant, further down) can reuse it too --
% \addtobeamertemplate patches whatever template is *currently* installed
% at the time it's called (gotham's own default one here), so a template
% defined fresh afterwards doesn't inherit that patch automatically and
% needs to invoke this explicitly instead.
% Every edge-anchored node in this hook (and in every other page-corner
% overlay in this file) needs an explicit "outer sep=0pt" alongside its
% "inner sep=0pt": tikz's default outer sep is 0.5\pgflinewidth (0.2pt at
% the default 0.4pt line width), and a node's edge/corner anchors sit on
% the *outer* boundary -- outer sep past the filled rectangle -- so an
% anchor placed exactly on a page corner leaves the fill itself 0.2pt
% short of that corner on the anchored sides. Found 2026-07-16 (direct
% instruction: a hairline white line above/left of titleframe's and
% quoteframe's yellow boxes, their right edges 0.4pt off the standard
% title bar's own visible yellow edge, and their logos 0.2pt left of
% every other frame's) by measuring the compiled PDF's drawn rectangles:
% every affected fill/image was inset by exactly 0.199bp = 0.2pt from
% its anchor. This is also the same offset the footer box's own comment
% (further down) had earlier measured, worked around with a hand-nudged
% anchor, and left as "root cause not chased further" -- now chased:
% it *is* a general property of every edge/corner-anchored tikz node,
% not something specific to that one construction. (Center anchors are
% unaffected -- outer sep is symmetric around the center -- so the
% gothamLogoAnchor.center redraw below doesn't need it.)
\newcommand{\gothamFrametitleHookCode}{%
	\tikz[remember picture, overlay]\node(gothamTitleBottom){};%
	\begin{tikzpicture}[remember picture, overlay]
		\path let \p1=(current page.north east), \p2=(gothamTitleBottom) in
			node[fill=white, anchor=north east, inner sep=0pt, outer sep=0pt, name=gothamMainPatch,
			     minimum width=\gothamLogoWidth+2\gothamFrametitleRightpading,
			     minimum height=\y1-\y2] at (\x1,\y1) {};
		\ifx\gothamFrametitleFullHeightYellow\empty
			\path let \p1=(current page.north east), \p2=(gothamTitleBottom),
			          \p3=(gothamLogoAnchor.south), \p4=(current page.west) in
				node[fill=white, anchor=north west, inner sep=0pt, outer sep=0pt,
				     minimum width=\x1-\x4, minimum height=\y3-\y2] at (\x4,\y3) {};
		\fi
		% Normally redrawn at gothamLogoAnchor's own position (wherever the
		% invisible in-flow logo -- on the same line as the title text --
		% ended up). \gothamFrametitleLogoFixed instead redraws it at a fixed
		% page-corner position, independent of where the title text sits --
		% for a frame type (sectionframe) whose title is deliberately *not*
		% at the position a normal frame's title/logo line would put it.
		% gothamMainPatch already covers the full logo column from the page
		% top down to the box's true bottom edge (above), so it's already
		% blank there regardless of which position is used.
		\ifx\gothamFrametitleLogoFixed\empty
			\node[anchor=center, inner sep=0pt] at (gothamLogoAnchor.center) {\includegraphics[height=\gothamLogoHeight]{\gothamInstituteLogoPath}};
		\else
			% inner ysep is derived, not \gothamFrametitleRightpading again:
			% the inline logo this fixed redraw must line up with is
			% *vertically centered* in the title bar, and its resulting top
			% margin (8.14pt, measured) is not exactly the horizontal
			% padding (0.28cm = 7.97pt) -- close enough that the pre-outer-sep-fix
			% code's accidental extra 0.2pt masked the difference, but a
			% visible micro-jump when flipping pages once the fix landed.
			% Rather than a fresh calibrated constant:
			% \gothamTitleBarBottomAlign (the visible bar bottom) *is* the
			% inline logo's bottom edge, so subtracting \gothamLogoHeight
			% reproduces its top margin exactly, and re-measuring one
			% constant keeps both in sync if the bar ever changes.
			\node[anchor=north east, inner xsep=\gothamFrametitleRightpading,
			      inner ysep=\dimexpr\gothamTitleBarBottomAlign-\gothamLogoHeight\relax,
			      outer sep=0pt] at (current page.north east) {\includegraphics[height=\gothamLogoHeight]{\gothamInstituteLogoPath}};
		\fi
		% Optional per-frame "chapter logo" (see \chapterlogo{} below): drawn
		% immediately to the left of the main logo's white patch, on its own
		% white patch, only when \gothamChapterLogoPath is non-empty. Sized
		% with exactly one \gothamFrametitleRightpading of its own padding, on
		% its outer (left) side only -- the inner (right) side is left flush
		% against gothamMainPatch, since that patch already contributes its
		% own \gothamFrametitleRightpading of padding around the main logo.
		% This makes both gaps -- chapter logo to main logo, and chapter logo
		% to where the yellow bar resumes -- come out equal, at exactly one
		% \gothamFrametitleRightpading each, without double-counting or
		% needing a second length constant.
		\ifx\gothamChapterLogoPath\empty\else
			\settowidth{\gothamChapterLogoWidth}{\includegraphics[height=\gothamLogoHeight]{\gothamChapterLogoPath}}%
			% The patch's right edge is nudged 0.5pt *past* gothamMainPatch's own
			% west edge (i.e. slightly overlapping it) rather than exactly
			% touching it -- both are solid white, so the overlap itself is
			% invisible, but touching them exactly left a hairline yellow seam
			% at the join (a real, if sub-millimetre, gap that persisted across
			% recompiles, most likely PDF fill-antialiasing at the shared edge
			% of two separately-filled abutting shapes). The logo image itself
			% (below) is *not* shifted, so the actual margin it defines is
			% unaffected.
			\path let \p1=(current page.north east), \p2=(gothamTitleBottom) in
				node[fill=white, anchor=east, inner sep=0pt, outer sep=0pt,
				     minimum width=\gothamChapterLogoWidth+\gothamFrametitleRightpading+0.5pt,
				     minimum height=\y1-\y2] at ($(gothamMainPatch.west)+(0.5pt,0)$) {};
			\node[anchor=east, inner sep=0pt, outer sep=0pt] at (gothamMainPatch.west |- gothamLogoAnchor.center)
				{\includegraphics[height=\gothamLogoHeight]{\gothamChapterLogoPath}};
		\fi
	\end{tikzpicture}%
}
\addtobeamertemplate{frametitle}{}{\gothamFrametitleHookCode}

% \chapterlogo{path} opts a single upcoming standardframe into an extra logo
% (e.g. a topic/chapter icon) next to the main institute logo -- see the
% frametitle hook above for how it's drawn. Call it *before* \begin{standardframe},
% since the frame title (and this hook) is typeset as part of \begin{frame}
% itself, before any frame-body content would run. standardframe resets this
% back to empty on \end{frame} (below), so it only ever applies to the one
% frame that follows the \chapterlogo{} call, never leaking into the next slide.
% Both the set and the reset use \global\def (checked empirically, via
% \typeout, against the frametitle hook's actual view of this macro on a
% following frame): beamer's frame environment opens its own TeX group around
% the frame body, so a plain (local) assignment done in standardframe's
% closing code -- which runs right after \end{frame} -- was silently undone
% once that group actually closed, reviving the *previous* value instead of
% the reset. \global\renewcommand alone wasn't enough to fix it (still leaked
% into the next frame) -- only dropping to the lower-level \global\def did.
\newlength{\gothamChapterLogoWidth}
\newcommand{\gothamChapterLogoPath}{}
\newcommand{\chapterlogo}[1]{\global\def\gothamChapterLogoPath{#1}}

% \gothamFrametitleFullHeightYellow: opt-out for the frametitle hook's
% "trim yellow to the logo's bottom edge" white patch (see the hook above) --
% empty (default) keeps that trim; non-empty skips it, letting the yellow
% show all the way down to the title box's own true bottom edge instead.
% That trim exists to erase a normally-tiny (~1-2pt) sliver of yellow that
% would otherwise stick out below the logo on an ordinary short title bar --
% harmless there, but sectionframe's deliberately tall bar turns that
% "tiny sliver" into most of the box, silently wiping out nearly all of the
% intended yellow with white. sectionframe sets this (globally, same
% set/reset reasoning as \gothamChapterLogoPath above); standardframe resets
% it back to empty on \end{frame}, so it's opt-in per frame like the others.
\newcommand{\gothamFrametitleFullHeightYellow}{}

% \gothamFrametitleLogoFixed: opt-in for the frametitle hook's logo redraw
% (see the hook above) to use a fixed page-corner position instead of
% wherever the title text's own line puts it. Normally the two coincide (the
% logo sits inline with the title, gotham's own design), which is exactly
% why this is off by default -- but sectionframe's title is deliberately
% *not* where a normal frame's title sits (bottom-anchored within a much
% taller bar), and letting the logo follow it down there would put it
% somewhere no other slide in the deck ever has it. Same set/reset reasoning
% and mechanism as \gothamChapterLogoPath/\gothamFrametitleFullHeightYellow
% above.
\newcommand{\gothamFrametitleLogoFixed}{}

% \gothamFooterBoxColor: the footer box's own fill color (see the footline
% hook further down), overridable per frame the same set-before/reset-after
% way as the flags above -- default \gothamFooterBoxColorDefault (set by the
% per-university preamble before \input{}ing this file, same "value
% supplied externally, referenced lazily" idiom as \gothamBulletColor/
% \gothamInstituteLogoPath), matching every frame's title bar by default;
% highlightframe overrides it to match its own (likewise overridable) title
% bar color, \gothamHighlightColor, so both boxes change together with a
% single retargeting.
\newcommand{\gothamFooterBoxColor}{\gothamFooterBoxColorDefault}

% A frametitle template variant that drops the inline institute logo
% entirely -- for use together with \gothamFrametitleLogoFixed (which
% redraws the logo separately, at a fixed position, making the inline one
% pointless anyway). gotham's own template (beamerouterthemegotham.sty,
% "default/gotham~subsameline") puts \nolinebreak\hfill{...logo...} directly
% after \insertframetitle; this variant drops that whole part.
%
% Doing so, though, means the logo no longer reserves *any* horizontal room
% on the title's own line the way it used to -- so rightskip here widens
% \gothamFrametitleRightpading to reserve the same width gothamMainPatch
% (the frametitle hook's white patch, drawn afterward) already reserves for
% the logo column: \gothamLogoWidth + 2*\gothamFrametitleRightpading. Without
% this, the title paragraph wraps as if that whole column were free text
% space -- not an actual TeX overfull-hbox condition (no warning, the text
% genuinely fits the *box's own* width), just a title that runs visually
% underneath the white patch and logo drawn on top of it afterward, which
% looks identical to clipped/hidden text. (An earlier, wrong diagnosis of
% this same visual symptom blamed gotham's \nolinebreak forcing an overfull
% line -- checked by testing gotham's original template installed with a
% long \huge title and no changes at all, which reproduces something that
% *looks* the same; but the actual fix needed here turned out to be this
% rightskip, not removing \nolinebreak, which was really addressing a red
% herring from an isolated test that happened not to include a competing
% logo column at all.)
%
% "plus 1fil" on that rightskip keeps every line ragged-right (its own
% natural width, however short of the margin that leaves) rather than fully
% justified -- without it, a wrapped multi-line title's non-last lines get
% stretched to exactly fill the available width, which reads as oddly wide,
% uneven word-gaps on a short title (most noticeable exactly when a title is
% long enough to wrap, i.e. always, for this template).
\makeatletter
\defbeamertemplate{frametitle}{gotham_no_inline_logo}{%
	\nointerlineskip
	\begin{beamercolorbox}[wd=\paperwidth, sep=0pt,
	                        leftskip=\gothamFrametitleLeftpading,
	                        rightskip=\dimexpr\gothamLogoWidth+2\gothamFrametitleRightpading\relax plus 1fil]{frametitle}
		\gotham@frametitlestrut@start
		\insertframetitle
		\ifx\insertframesubtitle\@empty\else
			{\usebeamerfont{framesubtitle}\usebeamercolor[fg]{framesubtitle}\gothamFrameSubtitleSep \insertframesubtitle}%
			\vspace*{\gothamFramesubtitleStrutend}%
		\fi
		\gotham@frametitlestrut@end
	\end{beamercolorbox}
	\gothamFrametitleHookCode
}
\makeatother


% Define slideID as a simple value holder (initially empty)
\newcommand{\slideID}{}%

% Define slideID as a simple value holder (not a command with arguments)
\makeatletter
\newcommand{\setSlideID}[1]{%
	% Grab everything up to the first '-' and store it in \slideID
	\edef\tempa{#1}%
	\expandafter\setSlideID@aux\tempa-\relax
}
\def\setSlideID@aux#1-#2\relax{%
	\renewcommand{\slideID}{#1}%
}
\makeatother

% Define a new \slide command to automatically pick the current version of the slide
% 2026-07-08 (direct instruction): slide filenames dropped their leading
% "s-" (e.g. slides/ab3-Some-Slide-current.tex, not the old
% slides/s-ab3-Some-Slide-current.tex) so slide IDs are directly
% searchable in Explorer -- this macro (and the other three below) is the
% *only* place that ever needed to change for that: every deck's own
% \slide{}/\slideversion{}/\omitslide{}/\omitslideversion{} calls pass just
% "ID-slug", never the filename itself, so renaming every file on disk to
% match needed no changes anywhere else (not the deck files, not
% slide_ids.csv, which only ever stored bare IDs to begin with).
\newcommand{\slide}[1]{%
	\setSlideID{#1}%
	\input{slides/#1-current.tex}%
}

% \slidefinal{ID-slug} -- like \slide{}, but renders the slide at its final,
% fully-revealed overlay state in every deck version, instructor included
% (2026-07-12, direct instruction: show jww's finished build in place of a
% separate recap slide, "irrespective of which deck version is being
% compiled" -- that recap slide has since been deleted entirely, its own ID
% no longer in use). A slide
% built with \conditionalpause[instructor] already renders fully built in
% the during-course/after-course compiles -- that pause is a no-op for
% them, so nothing there ever waits on a reveal step -- the only case
% actually needing help is the instructor deck, which would otherwise still
% step through the build. Locally shadowing \conditionalpause to a no-op
% (same optional-argument signature as the real one, so it still accepts
% e.g. \conditionalpause[instructor] without erroring) for the duration of
% one \slide{} call makes every \conditionalpause in that one inclusion
% inert, regardless of the actual compiled \gothamDeckVersion -- scoped
% inside \begingroup/\endgroup, so it reverts immediately after, and has no
% effect on any *other* \slide{}/\slidefinal{} call to the same ID
% elsewhere in the same deck (e.g. an earlier, normally-building instance of
% the same slide). \deckonly/\deckuncover/\deckvisible are untouched, so
% deck-version-tagged content (as opposed to reveal-step-tagged content)
% still follows its usual per-version rules.
\newcommand{\slidefinal}[1]{%
	\begingroup
		\renewcommand{\conditionalpause}[1][instructor]{}%
		\slide{#1}%
	\endgroup
}

% Load a specific, pinned version of a slide instead of whichever is "current"
% Usage: \slideversion{ID-slug}{M-m-p} loads slides/ID-slug-vM-m-p.tex
\newcommand{\slideversion}[2]{%
	\setSlideID{#1}%
	\input{slides/#1-v#2.tex}%
}

% Student-version decks: \studentversion (called anywhere in a deck, e.g.
% right after % Slidemaster preamble variant for University of Graz decks. Institution-
% specific content only (colors, logo, logo calibration) -- every frame
% type, the deck-version system, footer/hyperlink/footnote styling, tikz
% styles, etc. live in Slidemaster_preamble_base.tex, \input{}ed at the
% bottom of this file, and are shared verbatim across every institution
% this project supports (see Slidemaster_preamble_mci.tex and
% Slidemaster_preamble_uni-ibk.tex for the other two). A deck picks its
% institution by \input{}ing this file instead of one of those, nothing
% else changes.

% Color definition (Graz palette, sourced from the original single-preamble
% version of this project -- these are the University of Graz's own official
% colors, not derived/approximated).
\definecolor{c_primary_dark}{HTML}{000000}
\definecolor{c_primary_highlight}{HTML}{ffd500}
\definecolor{c_primary_gray}{HTML}{c6c6c6}
\definecolor{c_primary_white}{HTML}{FFFFFF}
\definecolor{c_secondary_violet}{HTML}{948dc3}
\definecolor{c_tertiary_violet}{HTML}{c2bddf}
\definecolor{c_secondary_blue}{HTML}{689cce}
\definecolor{c_tertiary_blue}{HTML}{bacce8}
\definecolor{c_secondary_green}{HTML}{9fcb81}
\definecolor{c_tertiary_green}{HTML}{cce2ba}
\definecolor{c_secondary_red}{HTML}{e4684b}
\definecolor{c_tertiary_red}{HTML}{f1ae95}
\definecolor{c_secondary_orange}{HTML}{f6a968}
\definecolor{c_tertiary_orange}{HTML}{fbd0aa}
\definecolor{c_secondary_yellow}{HTML}{ffe063}
\definecolor{c_tertiary_yellow}{HTML}{ffeeaf}
\definecolor{c_tertiary_gray}{HTML}{e3e3e3}
% Not an official Uni Graz color (hence "custom", not "secondary" like the
% palette above) -- a softer near-black than c_primary_dark, used for
% quoteframe's full-bleed dark block, matching the specific dark gray (not
% pure black) ExampleSlides.pptx used for that same block. Also reused for
% hyperlink text color, for the same reason it works for quoteframe's
% attribution text: dark enough for good contrast/legibility on a light
% background.
\definecolor{c_custom_gray}{HTML}{333333}
\definecolor{c_custom_red}{HTML}{C00000}
\definecolor{c_custom_green}{HTML}{00B050}

% Logo + calibration.
% schemes/scheme_uni_graz/Logo_Uni-Graz_4c-cropped.pdf -- official vector artwork,
% user-supplied (schemes/scheme_uni_graz/Logo_Uni-Graz_4c.eps, kept alongside as the
% untouched source), converted to PDF via epstopdf and cropped to its true
% ink bounding box (computed with Ghostscript's own `-sDEVICE=bbox`, same
% treatment as MCI's and Innsbruck's own logos -- see
% Slidemaster_preamble_mci.tex's own comment for the fuller explanation of
% why the EPS's own %%BoundingBox isn't trusted directly). Vector, not
% raster, so it stays crisp at any zoom. Supersedes this file's own earlier
% raster logo (schemes/scheme_uni_graz/Logo Uni Graz 4c-cropped.png, itself cropped
% from schemes/scheme_uni_graz/Logo Uni Graz 4c.png's own symmetric transparent
% margin -- both kept alongside as historical/fallback assets, not deleted).
\newcommand{\gothamInstituteLogoPath}{schemes/scheme_uni_graz/Logo_Uni-Graz_4c-cropped.pdf}
\newlength{\gothamLogoHeight}
\setlength{\gothamLogoHeight}{3ex}
% Re-measured against the compiled PDF for this specific logo image -- see
% Slidemaster_preamble_base.tex's own comment on \gothamFrametitleRightpading
% for why this needs re-measuring per institution rather than reusing
% another one's value.
\setlength{\gothamFrametitleRightpading}{0.28cm}

% Itemize bullet color (2026-07-16, direct instruction): Graz keeps the
% project's original behavior, bullets in the primary highlight color.
\newcommand{\gothamBulletColor}{c_primary_highlight}

% Title bar and footer box/text colors (2026-07-16, direct instruction):
% Graz keeps the project's original behavior throughout -- title/footer
% boxes in the primary highlight color, title/footer text in the primary
% dark color (unlike MCI/Innsbruck, which switch these to their own
% secondary blue/white -- see Slidemaster_preamble_mci.tex's own comment).
\newcommand{\gothamTitleBoxColor}{c_primary_highlight}
\newcommand{\gothamTitleTextColor}{c_primary_dark}
\newcommand{\gothamFooterBoxColorDefault}{c_primary_highlight}
\newcommand{\gothamFooterTextColor}{c_primary_dark}

\input{Slidemaster_preamble_base}
, similar to \sidefootnotes or
% \setfootnotesymbols) switches that compile into the "handout" variant --
% including, as of 2026-07-03, actually forcing beamer's own "handout" mode
% (\gdef\beamer@currentmode{handout}, normally only reachable via the
% \documentclass[handout]{beamer} *class option*, set once at class-loading
% time -- too early for a command like \studentversion, called later from a
% deck's own preamble, to use that route). This was checked, not assumed:
% until this point \studentversion only ever touched \omitslide{}, so
% nothing suppressed \pause/overlay specs -- true only by coincidence, since
% no slide converted so far happens to use them. Late-redefining
% \beamer@currentmode was verified (isolated minimal-beamer tests, plain
% \pause *and* inline specs like \item<2->) to reproduce the same collapse
% -to-one-page, show-the-final-revealed-state behavior as the class option,
% even set well after \documentclass{beamer} -- beamer's overlay decoder
% (beamerbasedecode.sty) reads \beamer@currentmode at frame-processing time,
% not just once at class-option time, so this is safe, not a hack that
% happens to work by luck of file ordering.
% \omitslide{}/\omitslideversion{} are drop-in replacements for
% \slide{}/\slideversion{} that tag a slide as excluded from that variant:
%
%   \omitslide{xyz-Answer-key}                 % instead of \slide{...}
%   \omitslideversion{xyz-Answer-key}{1-0-0}    % instead of \slideversion{...}{...}
%
% In a \studentversion compile, a tagged slide's \input{} is skipped
% entirely -- no frame is produced for it at all (not a blank page, not a
% numbering gap: beamer's frame/slide counters simply never advance for a
% call that never ran). In every other (i.e. normal/instructor) compile, a
% tagged slide loads exactly as it would via \slide{}/\slideversion{}, with
% one difference: \gothamOmittedSlide is set for the duration of that one
% \input{}, so the footer template below can print a marker next to that
% frame's slide ID -- a deliberately visible reminder, while reviewing the
% instructor deck, of which slides won't make it to students. The same tag
% works for both compiles because it lives at the call site in the *deck*
% file (not inside the slide's own content file): the identical slide list
% can serve both the instructor deck and the student handout, e.g.
%
%   % Slidemaster preamble variant for University of Graz decks. Institution-
% specific content only (colors, logo, logo calibration) -- every frame
% type, the deck-version system, footer/hyperlink/footnote styling, tikz
% styles, etc. live in Slidemaster_preamble_base.tex, \input{}ed at the
% bottom of this file, and are shared verbatim across every institution
% this project supports (see Slidemaster_preamble_mci.tex and
% Slidemaster_preamble_uni-ibk.tex for the other two). A deck picks its
% institution by \input{}ing this file instead of one of those, nothing
% else changes.

% Color definition (Graz palette, sourced from the original single-preamble
% version of this project -- these are the University of Graz's own official
% colors, not derived/approximated).
\definecolor{c_primary_dark}{HTML}{000000}
\definecolor{c_primary_highlight}{HTML}{ffd500}
\definecolor{c_primary_gray}{HTML}{c6c6c6}
\definecolor{c_primary_white}{HTML}{FFFFFF}
\definecolor{c_secondary_violet}{HTML}{948dc3}
\definecolor{c_tertiary_violet}{HTML}{c2bddf}
\definecolor{c_secondary_blue}{HTML}{689cce}
\definecolor{c_tertiary_blue}{HTML}{bacce8}
\definecolor{c_secondary_green}{HTML}{9fcb81}
\definecolor{c_tertiary_green}{HTML}{cce2ba}
\definecolor{c_secondary_red}{HTML}{e4684b}
\definecolor{c_tertiary_red}{HTML}{f1ae95}
\definecolor{c_secondary_orange}{HTML}{f6a968}
\definecolor{c_tertiary_orange}{HTML}{fbd0aa}
\definecolor{c_secondary_yellow}{HTML}{ffe063}
\definecolor{c_tertiary_yellow}{HTML}{ffeeaf}
\definecolor{c_tertiary_gray}{HTML}{e3e3e3}
% Not an official Uni Graz color (hence "custom", not "secondary" like the
% palette above) -- a softer near-black than c_primary_dark, used for
% quoteframe's full-bleed dark block, matching the specific dark gray (not
% pure black) ExampleSlides.pptx used for that same block. Also reused for
% hyperlink text color, for the same reason it works for quoteframe's
% attribution text: dark enough for good contrast/legibility on a light
% background.
\definecolor{c_custom_gray}{HTML}{333333}
\definecolor{c_custom_red}{HTML}{C00000}
\definecolor{c_custom_green}{HTML}{00B050}

% Logo + calibration.
% schemes/scheme_uni_graz/Logo_Uni-Graz_4c-cropped.pdf -- official vector artwork,
% user-supplied (schemes/scheme_uni_graz/Logo_Uni-Graz_4c.eps, kept alongside as the
% untouched source), converted to PDF via epstopdf and cropped to its true
% ink bounding box (computed with Ghostscript's own `-sDEVICE=bbox`, same
% treatment as MCI's and Innsbruck's own logos -- see
% Slidemaster_preamble_mci.tex's own comment for the fuller explanation of
% why the EPS's own %%BoundingBox isn't trusted directly). Vector, not
% raster, so it stays crisp at any zoom. Supersedes this file's own earlier
% raster logo (schemes/scheme_uni_graz/Logo Uni Graz 4c-cropped.png, itself cropped
% from schemes/scheme_uni_graz/Logo Uni Graz 4c.png's own symmetric transparent
% margin -- both kept alongside as historical/fallback assets, not deleted).
\newcommand{\gothamInstituteLogoPath}{schemes/scheme_uni_graz/Logo_Uni-Graz_4c-cropped.pdf}
\newlength{\gothamLogoHeight}
\setlength{\gothamLogoHeight}{3ex}
% Re-measured against the compiled PDF for this specific logo image -- see
% Slidemaster_preamble_base.tex's own comment on \gothamFrametitleRightpading
% for why this needs re-measuring per institution rather than reusing
% another one's value.
\setlength{\gothamFrametitleRightpading}{0.28cm}

% Itemize bullet color (2026-07-16, direct instruction): Graz keeps the
% project's original behavior, bullets in the primary highlight color.
\newcommand{\gothamBulletColor}{c_primary_highlight}

% Title bar and footer box/text colors (2026-07-16, direct instruction):
% Graz keeps the project's original behavior throughout -- title/footer
% boxes in the primary highlight color, title/footer text in the primary
% dark color (unlike MCI/Innsbruck, which switch these to their own
% secondary blue/white -- see Slidemaster_preamble_mci.tex's own comment).
\newcommand{\gothamTitleBoxColor}{c_primary_highlight}
\newcommand{\gothamTitleTextColor}{c_primary_dark}
\newcommand{\gothamFooterBoxColorDefault}{c_primary_highlight}
\newcommand{\gothamFooterTextColor}{c_primary_dark}

\input{Slidemaster_preamble_base}

%   \FooterText{...}
%   \makeatletter\ifdefined\HANDOUT\studentversion\fi\makeatother  % or similar
%   \input{course-content.tex}   % the shared \slide{}/\omitslide{} list
%
% with two thin wrapper decks (one plain, one calling \studentversion first)
% both \input-ing the same shared content file -- see README.md.
\newif\ifgothamStudentVersion
% \makeatletter must wrap the \newcommand itself, not sit *inside* the macro
% body: catcodes apply at tokenization time, which for a macro body happens
% once, immediately, when \newcommand reads it -- not deferred until the
% macro is later called. \makeatletter placed inside the body (tried first)
% was too late: \beamer@currentmode had already been mis-tokenized as
% "\beamer" + the six-way-otherwise-ordinary character "@" + eleven plain
% letters "currentmode", by the time that inner \makeatletter could act on
% it, so the \gdef silently created/redefined the wrong thing and never
% touched beamer's actual mode variable. Confirmed via \meaning\beamer@
% currentmode before/after \studentversion -- unchanged with the inside-the-
% body version, fixed once \makeatletter moved out here.
\makeatletter
\newcommand{\studentversion}{%
	\global\gothamStudentVersiontrue
	\gdef\beamer@currentmode{handout}%
}
\makeatother

% Deck version system - new boolean flags and string for current version
\newif\ifgothamInstructorVersion
\newif\ifgothamDuringCourseVersion
\newif\ifgothamAfterCourseVersion
\gdef\gothamDeckVersion{instructor}
\gdef\gothamDeckVersionText{}

% \deckversion[] - set the deck version(s)
% Usage: \deckversion[instructor], \deckversion[during-course], etc.
% Multiple versions: \deckversion[instructor,during-course,after-course]
% Parses #1 as a comma-separated list via etoolbox's \forcsvlist (rather
% than the single \ifstrequal{whole-list}{"instructor"} check this used to
% do, which only ever matched a single-version list -- a multi-version
% \deckversion[instructor,during-course] silently set none of the three
% booleans, since the compared string was the literal, unsplit
% "instructor,during-course"). Note \forcsvlist, not etoolbox's similarly-
% named \forlistloop: \forlistloop iterates etoolbox's own internal
% pipe-delimited list format (built via \listadd), not a plain comma string
% like this command's own argument -- confirmed directly (isolated test),
% passing a raw comma string to it doesn't error or split on the commas, it
% just treats the whole string as a single one-item "list".
\makeatletter
\newcommand{\deckversion}[1][instructor]{%
	% \xdef (not \gdef): #1 needs to be fully expanded here, not just
	% stored as whatever token(s) were passed -- \deckversion[\myvar]
	% (compile-all-versions.bat's own override mechanism, via a wrapper
	% file's \def\deckversionoverride{...}) passes a macro, not a literal
	% string, and everywhere else in this deck-version system that reads
	% \gothamDeckVersion only ever applies one \expandafter to it (see
	% \ifdeckversioninlist's own comment) -- enough to unwrap \gothamDeckVersion
	% itself, but not a further macro still sitting inside it. Confirmed
	% directly: with the old \gdef, \deckversion[\deckversionoverride]
	% left \gothamDeckVersionText empty and every \ifdeckversioninlist
	% check false, since \gothamDeckVersion's own stored content was the
	% single unexpanded token \deckversionoverride, not "during-course".
	\xdef\gothamDeckVersion{#1}%
	\global\gothamInstructorVersionfalse
	\global\gothamDuringCourseVersionfalse
	\global\gothamAfterCourseVersionfalse
	\gdef\gothamDeckVersionText{}%
	\def\do@deckversions##1{%
		\ifstrequal{##1}{instructor}{\global\gothamInstructorVersiontrue\gdef\gothamDeckVersionText{instructor deck}}{}%
		\ifstrequal{##1}{during-course}{\global\gothamDuringCourseVersiontrue\gdef\gothamDeckVersionText{student course deck}}{}%
		\ifstrequal{##1}{after-course}{\global\gothamAfterCourseVersiontrue\gdef\gothamDeckVersionText{student final deck}}{}%
	}%
	% \forcsvlist itself splits on literal commas in its own raw,
	% unexpanded argument (confirmed directly against etoolbox's own
	% source, \etb@defparser@do@aux) -- passing #1 straight through here
	% has the exact same problem as the \gdef above did, one call site
	% later. Building the whole \forcsvlist{...}{...} call inside an
	% \edef (protecting \forcsvlist/\do@deckversions themselves with
	% \noexpand so they're invoked *later*, once this constructed call is
	% run, not expanded away here) expands only \gothamDeckVersion --
	% already fully resolved by the \xdef above -- into the list argument.
	\edef\gotham@deckversion@call{\noexpand\forcsvlist{\noexpand\do@deckversions}{\gothamDeckVersion}}%
	\gotham@deckversion@call
}
\makeatother

\newif\ifgothamOmittedSlide
\newcommand{\gothamOmittedSlideMarker}{ | omitted from student deck}
% \stepcounter{framenumber} in the skipped branch below keeps the student
% deck's slide numbers matching the instructor deck's, gap and all, so "see
% slide 12" means the same physical slide to both audiences: beamer steps
% "framenumber" exactly once per \begin{frame} (checked in
% beamerbaseframe.sty -- once per *frame*, not per overlay page, so this is
% right regardless of how many overlay steps the omitted slide would have
% had), and this is the one place that increment is skipped along with the
% frame itself, so it has to be replayed by hand.
\newcommand{\omitslide}[1]{%
	\ifgothamStudentVersion
		\stepcounter{framenumber}%
	\else
		\global\gothamOmittedSlidetrue
		\setSlideID{#1}%
		\input{slides/#1-current.tex}%
		\global\gothamOmittedSlidefalse
	\fi
}
\newcommand{\omitslideversion}[2]{%
	\ifgothamStudentVersion
		\stepcounter{framenumber}%
	\else
		\global\gothamOmittedSlidetrue
		\setSlideID{#1}%
		\input{slides/#1-v#2.tex}%
		\global\gothamOmittedSlidefalse
	\fi
}

% Deck version-aware overlay commands
% Syntax: \deckonly<overlay>[versions]{content}
% If [versions] omitted, defaults to all decks (instructor,during-course,after-course)
%
% Built on etoolbox's own \forcsvlist for the comma-splitting (rather than a
% hand-rolled parser: an earlier version tried to split a list at each comma
% using \newcommand's plain, non-delimited argument grabbing
% (\newcommand{\ifdeckversioninlist@split}[3]{...} called as
% \ifdeckversioninlist@split#1\@nil{#1}) -- but a plain \newcommand doesn't do
% delimited-parameter matching against \@nil the way a \def with a matching
% pattern would; it just grabs the next 3 tokens/groups one at a time. So
% e.g. "instructor" silently split into single characters ("i", "n", "s",
% ...) instead of being read as one list item, which eventually ran the
% parser off the end of its own input and surfaced as "Undefined control
% sequence" wherever the first \conditionalpause[...]/\deckonly[...] in a
% deck happened to be used -- confirmed directly against the compile log,
% traced to slide y5z-CoCo-Bonds's own \conditionalpause[instructor].
%
% NOT etoolbox's own \ifinlist{item}{list}, despite the name similarity: that
% one tests membership in etoolbox's own internal list format, built only via
% \listadd (items joined with a literal "|", not ","); fed a plain comma
% string instead, confirmed directly (isolated test) to silently return
% false for every item, including an exact single-item match.
%
% The membership test itself needs one more subtlety: #1 here is the deck
% version to search for -- always \gothamDeckVersion at every call site
% below, i.e. a macro, not a literal string -- so comparing it via a plain
% \ifstrequal{##1}{#1} does not work either: etoolbox's \ifstrequal compares
% via \detokenize, which (confirmed directly, isolated test) stringifies an
% unexpanded macro token to its own control-sequence *name* ("gothamDeckVersion"),
% not the current value that macro holds. \expandafter forces that expansion
% before \ifstrequal ever sees it.
\makeatletter

\let\gotham@original@only\only
\let\gotham@original@uncover\uncover
\let\gotham@original@visible\visible

\newif\ifgotham@deckversion@found

% Usage: \ifdeckversioninlist{version}{list}{true-branch}{false-branch}
\newcommand{\ifdeckversioninlist}[2]{%
	\global\gotham@deckversion@foundfalse
	\def\gotham@deckversion@check##1{%
		\expandafter\ifstrequal\expandafter{#1}{##1}{\global\gotham@deckversion@foundtrue}{}%
	}%
	\forcsvlist{\gotham@deckversion@check}{#2}%
	\ifgotham@deckversion@found
		\expandafter\@firstoftwo
	\else
		\expandafter\@secondoftwo
	\fi
}

% \@ifnextchar-based parsing (the same idiom LaTeX2e's own kernel commands
% use for an optional argument that isn't a plain trailing [...]) so the
% overlay spec reads exactly like beamer's own \only<...>, with the deck-
% version list as a further, genuinely optional [...] afterwards:
%   \deckonly<2>{content}                              visible in all versions
%   \deckonly<2>[instructor]{content}                  instructor only
%   \deckonly<2>[during-course,after-course]{content}  during + after
\newcommand{\deckonly}{\@ifnextchar<{\gotham@deckonly@overlay}{\gotham@deckonly@overlay<1->}}
\def\gotham@deckonly@overlay<#1>{\@ifnextchar[{\gotham@deckonly@versions<#1>}{\gotham@deckonly@versions<#1>[instructor,during-course,after-course]}}
\def\gotham@deckonly@versions<#1>[#2]#3{%
	\ifgothamInstructorVersion
		\gotham@original@only<#1>{#3}%
	\else
		\ifdeckversioninlist{\gothamDeckVersion}{#2}{\gotham@original@only<#1>{#3}}{}%
	\fi
}

% Same interface as \deckonly, wrapping \uncover and \visible respectively
\newcommand{\deckuncover}{\@ifnextchar<{\gotham@deckuncover@overlay}{\gotham@deckuncover@overlay<1->}}
\def\gotham@deckuncover@overlay<#1>{\@ifnextchar[{\gotham@deckuncover@versions<#1>}{\gotham@deckuncover@versions<#1>[instructor,during-course,after-course]}}
\def\gotham@deckuncover@versions<#1>[#2]#3{%
	\ifgothamInstructorVersion
		\gotham@original@uncover<#1>{#3}%
	\else
		\ifdeckversioninlist{\gothamDeckVersion}{#2}{\gotham@original@uncover<#1>{#3}}{}%
	\fi
}

\newcommand{\deckvisible}{\@ifnextchar<{\gotham@deckvisible@overlay}{\gotham@deckvisible@overlay<1->}}
\def\gotham@deckvisible@overlay<#1>{\@ifnextchar[{\gotham@deckvisible@versions<#1>}{\gotham@deckvisible@versions<#1>[instructor,during-course,after-course]}}
\def\gotham@deckvisible@versions<#1>[#2]#3{%
	\ifgothamInstructorVersion
		\gotham@original@visible<#1>{#3}%
	\else
		\ifdeckversioninlist{\gothamDeckVersion}{#2}{\gotham@original@visible<#1>{#3}}{}%
	\fi
}

% \conditionalpause[versions] - \pause, but only takes effect in the listed
% deck versions (default: instructor only); a no-op, not even a skipped
% overlay step, in any other version
% Usage: \conditionalpause[instructor,during-course] or \conditionalpause[after-course]
\newcommand{\conditionalpause}[1][instructor]{%
	\ifdeckversioninlist{\gothamDeckVersion}{#1}{\pause}{}%
}

% \deckreveal[step]{content} - the "hidden answer, click-revealed for the
% instructor, absent for during-course, pre-shown for after-course" pattern
% already documented by hand on several exercise slides (e.g. mbz/x7q/37s/
% 3cv's own header comments: "a slide/reveal-step that is hidden in the
% source PowerPoint should, by default, only be made visible on the
% instructor version"). A bare, unconditional \only<N->{content} looks like
% it does this, but doesn't: \only's overlay range registers with beamer's
% own overlay-counting pass regardless of \ifgothamInstructorVersion/
% \conditionalpause -- neither of those touch \only itself, they only guard
% *other* content around it -- so a frame with nothing else establishing a
% second overlay step still gets a real second slide out of a bare
% \only<2->, in *every* deck version, not just the instructor one. Confirmed
% directly (2026-07-21, direct instruction) on 37s/3cv: both a) leaked the
% "hidden" content to during-course as a genuine second slide, the exact
% thing this pattern exists to prevent, and b) showed a blank second bullet/
% cell first on that leaked slide's own step 1 -- \only<2-> hides the
% *content*, not the space it would occupy, so during-course's own step 1
% still had nothing else deciding where the bullet/cell landed.
% \deckreveal fixes both by picking a genuinely different rendering per
% version rather than reusing the same overlay-restricted \only call
% everywhere: instructor gets the real overlay build (\only<step->,
% default step 2, matching \tallitem's own default-multiplier convention);
% during-course gets nothing at all (no token emitted, so beamer's overlay
% pass never sees an overlay spec to register); after-course gets the
% content unconditionally, with no overlay spec of its own either, so it
% doesn't silently force a second slide there either. Content-level only
% (wraps a value/expression, like \only itself) -- doesn't cover an
% itemize \item's own bullet marker, which needs its own overlay spec
% (`\item<step->`) directly on the \item, not wrapped from outside; see
% 37s's own comment for that variant, hand-written rather than factored
% into a second macro since only one call site needs it so far.
\newcommand{\deckreveal}[2][2]{%
	\ifgothamDuringCourseVersion
	\else
		\ifgothamInstructorVersion
			\only<#1->{#2}%
		\else
			#2%
		\fi
	\fi
}

% \conditionalslide[versions]{ID-slug} - \slide{}, but only included in the
% listed deck versions; keeps the frame counter in step with the omitted
% version otherwise (same \stepcounter{framenumber} technique \omitslide{}
% already uses for the student deck), so "see slide N" still points at the
% same physical slide across every deck version.
% Usage: \conditionalslide[instructor,after-course]{xyz-Answer-key}
\newcommand{\conditionalslide}[2][instructor,during-course,after-course]{%
	\ifdeckversioninlist{\gothamDeckVersion}{#1}%
		{\setSlideID{#2}%
		 \input{slides/#2-current.tex}}%
		{\stepcounter{framenumber}}%
}

\makeatother

% PDF accessibility: mark purely decorative/running content (footer, below;
% also reused on titleframe's/quoteframe's own decorative background boxes;
% logo not yet included, see README.md's own accessibility section for that
% open question) as a PDF "Artifact", so screen readers skip it -- via
% LuaTeX's raw \pdfextension literal
% primitive. \ifluatex-guarded (iftex is already loaded, via beamer.cls
% itself) because \pdfextension is LuaTeX-only -- this project's README
% documents xelatex as a supported alternative engine, and compiling with it
% hit exactly this: "Undefined control sequence: \pdfextension", hard-failing
% every frame, since the first version of this didn't guard the primitive at
% all. Under any non-LuaTeX engine these two commands now silently do
% nothing instead: decorative content just isn't marked as an Artifact under
% xelatex/pdflatex, rather than breaking the compile -- a real loss of
% accessibility under those engines specifically, not a full equivalent fix,
% but not something to solve today (see also: compiling with plain pdflatex
% already silently mis-renders the whole deck's fonts, so xelatex is really
% the only *other* engine this needs to not break under; pdflatex was never
% viable here to begin with, per "Compiling" in README.md).
% Deliberately NOT tagpdf/kernel tagging (\DocumentMetadata):
% beamer 3.77 refuses to compile at all when that's active ("The beamer
% class is not compatible with \DocumentMetadata ... switch to the new
% ltx-talk class if you require tagged output"), and ltx-talk is a from-
% scratch class (not beamer-based, no theme/template system) that this
% framework can't realistically migrate to. Bypassing beamer's check by
% force-loading the underlying package before \documentclass also failed --
% it gets past hyperref but then hits a genuine internal error in tagpdf's
% LuaTeX backend from the non-standard load order, not something to build
% real coursework on.
%
% /Artifact is a standalone PDF marked-content tag (PDF spec 14.8.2.2):
% valid in *any* PDF, with or without a structure tree, no /Properties
% resource entry or MCID required -- unlike structure-tree tagging, this
% doesn't require kernel tagging support at all. It's a real, working
% partial fix, not full PDF/UA conformance: there's no structure tree, so
% no true semantic headings, and reading order elsewhere on the slide still
% depends on plain content-stream order (verified, not enforced, for
% standardframe -- see its own comment).
% #1 is the artifact's inline PDF dictionary content, e.g.
% "/Type /Pagination /Subtype /Footer" for a running footer, or
% "/Type /Background" for a decorative fill -- not hardcoded to one /Type,
% since not everything decorative here is "Pagination" in the PDF spec's own
% sense of that word (running headers/footers/page numbers/watermarks);
% titleframe's/quoteframe's plain background color fills are semantically
% "Background" artifacts instead, a distinct standard type.
\newcommand{\gothamArtifactBegin}[1]{%
	\ifluatex\pdfextension literal direct{/Artifact <<#1>> BDC}\fi
}
\newcommand{\gothamArtifactEnd}{\ifluatex\pdfextension literal direct{EMC}\fi}

% Footer
% Smaller footer + a bit of padding under it
\setbeamerfont{frame footer}{size=\fontsize{5}{3}\selectfont} % {font size}{baseline skip} in pt

% Shrinks the gap between footnotes and the footer box below them. The
% theme positions the footer's content via \vspace*{\gothamFootlineVOffset}
% (a fixed, negative length) before it starts; making that more negative
% turns out to give footnotes (which are placed just before the footer, and
% stretch to fill whatever room remains via beamer's own
% \beamer@framebottomskip) more room to extend into, pulling them down
% towards the footer -- checked empirically against the PDF's own
% coordinates, this does *not* move the footer's own text at all (that
% still sits exactly where \gothamFooterBoxHeight below was calibrated for),
% so the footer box itself needs no compensating adjustment.
%
% Partly given back (+1.4pt) 2026-07-05: a footnote/source containing an
% underlined \href (see the hyperlink-styling comment above) sits close
% enough to the footer at the original -1.3pt that the underline visually
% touches the footer box (checked directly against the compiled PDF's own
% pixel coordinates: a plain-text footnote's own gap to the footer measured
% ~1.9pt, but an underlined-link footnote's gap measured only ~0.6pt at
% the same setting, since the strut rebuilt for footnotes just below --
% correctly sized for plain text's own natural depth -- doesn't account
% for the underline rule extending further down). +1.4pt splits the
% difference: a linked footnote's gap now measures close to the plain-
% footnote reference (~1.9pt) instead of colliding with the footer, and a
% plain footnote's own gap just becomes a little more relaxed (~3.3pt)
% rather than stretched uncomfortably tight -- applied globally rather
% than only to footnotes with links, since there's no clean per-footnote
% hook for it and the plain-footnote side effect is minor.
%
% Nudged back down by 0.4pt the next day (direct instruction: move
% footnotes down "by the width of the underline line") -- the +1.4pt above
% left a bit more headroom than necessary; 0.4pt is the underline rule's
% own measured stroke thickness (a solid ~1px core at 600dpi, i.e.
% 1/600in \approx 0.12pt, plus antialiasing bleed either side consistent
% with the common 0.4pt LaTeX default rule width -- measured directly
% against the compiled PDF's own pixel values, not assumed).
%
% Reverted the very next round (direct instruction: move back *up* by that
% same underline thickness) -- net effect is back to the plain +1.4pt this
% section started with; kept as two separate historical adjustments below
% rather than collapsing them into one, so this back-and-forth stays
% visible rather than silently disappearing.
\addtolength{\gothamFootlineVOffset}{-1.3pt}
\addtolength{\gothamFootlineVOffset}{1.4pt}
\addtolength{\gothamFootlineVOffset}{-0.4pt}
\addtolength{\gothamFootlineVOffset}{0.4pt}

% 2026-07-06 (direct instruction): move footnotes down by half the
% measured distance between a linked footnote's own underline and the
% footer box, checked on slides/3ic-Corporate-Venture-Capital-
% current.tex specifically (the underline sits at y=245.95pt, the footer
% box's own top edge at y=247.69pt, a 1.74pt gap) -- half of that is
% 0.87pt, applied as a further negative adjustment (more negative =
% footnotes get pulled further down, per this section's own comment
% above).
\addtolength{\gothamFootlineVOffset}{-0.87pt}

% Shifts the right end of the footer line (the frame number) so it lines up
% with the logo's own right edge above it, rather than wherever the theme's
% default footline margins would otherwise put it. Empirically calibrated
% against the rendered PDF coordinates of both; re-measure and adjust if
% \gothamFrametitleRightpading (the logo's own right inset) changes.
\newlength{\gothamFooterRightShift}
\setlength{\gothamFooterRightShift}{4.6pt}

% Nudges the footer text down within its own yellow box, to vertically
% center it there: it previously sat noticeably closer to the box's top
% edge than its bottom (0.14pt above vs. 1.12pt below, checked against the
% rendered PDF's own text coordinates) since \gothamFootlineVOffset's own
% positioning (further up) turns out not to move this text at all (see its
% comment above) -- this is a separate, additive leading \vspace* instead.
\newlength{\gothamFooterTextDrop}
\setlength{\gothamFooterTextDrop}{-0.49pt}

% Compile timestamp: lets you tell, from a PDF alone (no need to check file
% properties/mtime, which don't survive e.g. emailing or uploading a copy),
% whether it reflects the latest edits or a stale earlier build -- handy
% whenever a PDF's *content* is in question, not just when debugging compile
% tooling itself. LuaTeX-only, matching this project's required engine (see
% README.md's "Compiling" section) -- no \ifluatex guard needed since
% nothing else here works under another engine to begin with.
%
% Deliberately avoids os.date's own %-format directives (e.g. "%Y-%m-%d")
% entirely, building the string from os.date("*t")'s field table and plain
% Lua concatenation instead -- \directlua's argument is tokenized by plain
% TeX first, where % starts a comment, before the text ever reaches Lua as a
% string. The obvious-looking fix, escaping each one as \%, does *not* work:
% confirmed directly, it produces a Lua string containing a literal
% backslash followed by a percent sign (detokenization keeps a control
% sequence's own backslash, for \% same as any other), which Lua's own
% string-literal parser then rejects outright ("invalid escape sequence"),
% silently aborting that whole \directlua call -- with no visible LaTeX-side
% error, just a cache file ending up written with an empty value. Switching
% catcode temporarily (\catcode`\%=12 ... ) would also work but needs a
% \begingroup/\endgroup wrapper at every call site; skipping % entirely is
% simpler and needs no such wrapper.
%
% Cached to \jobname.timestamp (via plain TeX \write, so MiKTeX's
% -aux-directory flag -- already applied through .latexmkrc's %O -- routes
% it into .build/ like every other aux file, not left cluttering decks/ or
% the project root) rather than recomputed fresh on every pass: a live
% os.date() call on every pass, tried first, broke latexmk's own rerun
% convergence outright -- confirmed directly, by isolating the change, that
% an otherwise-identical clean compile converges normally without it. Each
% of latexmk's internal reruns relaunches lualatex as a brand-new process,
% and a full pass over a deck this size (font loading alone) comfortably
% crosses a one-minute boundary, so two consecutive passes embedded two
% different strings in the footer; latexmk saw the .aux checksum differ
% *again* on the very next pass and looped ("Latexmk: References changed.")
% until hitting its rerun limit and reporting a false "gave an error", with
% no actual LaTeX error anywhere in the log. Caching to a file read back on
% later passes fixes that: the value is fixed the moment it's first
% written, so the footer's rendered text -- and therefore the .aux
% checksum -- stays identical across every subsequent pass of the same
% compile.
%
% This cache is deliberately *not* self-invalidating by date/content -- it
% keeps reporting an old timestamp indefinitely once written, across
% unrelated future compiles, unless something clears it first.
% SlideManager's compiler.py does exactly that: before every job it deletes
% that job's own .build/<jobname>.* (which includes <jobname>.timestamp), so
% a run through SlideManager always gets a fresh stamp. An ad-hoc TeXstudio
% compile without going through SlideManager first may show a stale stamp if
% .build/ already has one left over from an earlier session -- delete
% .build/<jobname>.timestamp by hand (or just compile that job once through
% SlideManager) if that matters for what you're checking.
\directlua{
	function gothamZeroPad(n) if n < 10 then return "0" .. n else return "" .. n end end
}
\newwrite\gothamCompileTimestampFile
\newcommand{\gothamCompileTimestamp}{%
	\ifdefined\gothamCompileTimestampValue
	\else
		\IfFileExists{\jobname.timestamp}{%
			\input{\jobname.timestamp}%
		}{%
			\xdef\gothamCompileTimestampValue{\directlua{%
				local d = os.date("*t")%
				tex.sprint(d.year .. "-" .. gothamZeroPad(d.month) .. "-" .. gothamZeroPad(d.day) .. " " .. gothamZeroPad(d.hour) .. ":" .. gothamZeroPad(d.min))%
			}}%
			\immediate\openout\gothamCompileTimestampFile=\jobname.timestamp\relax
			\immediate\write\gothamCompileTimestampFile{\string\gdef\string\gothamCompileTimestampValue{\gothamCompileTimestampValue}}%
			\immediate\closeout\gothamCompileTimestampFile
		}%
	\fi
	\gothamCompileTimestampValue
}

% \gothamHideFooterText: opt-in flag suppressing just \FooterText (the left
% part of the footer) -- default false. Set true by the standout template's
% own hook below, so `\begin{frame}[standout]` hides only that left part,
% keeping the "(slide id: ...)" block and frame number visible (2026-07-16,
% direct instruction). Reset immediately after being read, inside the
% footer template itself, rather than relying on a per-environment closing
% hook (standout frames are bare `\begin{frame}[standout]...\end{frame}`,
% not one of this project's own environments like standardframe, so there's
% no custom closing code to reset it in) -- safe because \gothamHideFooterText
% is only ever consumed here, once, immediately before being cleared.
\newif\ifgothamHideFooterText
% gotham's own \defbeamertemplate*{standout}{gotham} (gotham.dtx) is what
% `\begin{frame}[standout]` actually activates; \addtobeamertemplate's
% "after" hook appends to it without needing to touch that upstream file.
\addtobeamertemplate{standout}{}{\global\gothamHideFooterTexttrue}

\setbeamertemplate{frame footer}{%
	\vspace*{\gothamFooterTextDrop}%
	\ifgothamHideFooterText
		\global\gothamHideFooterTextfalse
	\else
		\FooterText%
	\fi
	\hfill%
	{\fontsize{4}{4}\selectfont (slide id: \slideID\ifgothamOmittedSlide\gothamOmittedSlideMarker\fi\ | \gothamDeckVersionText\ | \gothamCompileTimestamp)}\hspace*{2em}\insertframenumber%
	\hspace*{-4em}%
	\hspace*{-\gothamFooterRightShift}%
	\vspace*{0.4ex}%
}

% Give the footer the same yellow-box treatment as the title bar, but --
% unlike the title bar, which is tall for visual weight -- only as tall as
% the footer's own text plus equal padding above and below.
%
% This paints a yellow patch *behind* the theme's own footline rather than
% replacing that template's content, because \usebeamertemplate{frame~footer}
% (which is what actually typesets \FooterText etc.) only produces its
% content when invoked from the theme's own original footline template --
% calling it again from anywhere else (a redefined footline, the frame body,
% a \settoheight measurement, ...) silently produces nothing. It's
% positioned from "current page.south" (the page's own bottom-left corner)
% rather than a \tikz "remember picture" node tracking the footer's actual
% position, since there's no per-frame content to track: the footer text
% (in its actual 5pt font, per \setbeamerfont{frame footer} above -- despite
% the "page number in head/foot" font switch inside the theme's footline
% template, this is the font that actually ends up applied) sits at a fixed
% distance from the bottom of the page on every frame.
%
% \gothamFooterBoxHeight below is calibrated to that actual measured
% position (checked via the PDF's own text coordinates, not just visually):
% the text sits about 0.84pt above the page's bottom edge and spans about
% one em of its own 5pt font, plus a bit of padding above and below. If the
% footer font size changes, this will need re-measuring.
\newlength{\gothamFooterBoxHeight}
\setlength{\gothamFooterBoxHeight}{7.26pt}
% \vbox to 0pt{...} forces this to occupy exactly zero space in the normal
% vertical list, *without* switching into horizontal mode the way an hbox
% (e.g. \raisebox) would -- inserting an hbox here forces LaTeX to open an
% implicit paragraph before the theme's own \vspace*{\gothamFootlineVOffset}
% runs, which changes how that \vspace behaves (it special-cases being at
% the very start of vertical mode) and nudges the footer text down by a few
% points. A \vbox is already vertical material, so it doesn't trigger that.
% \gothamArtifactBegin/\gothamArtifactEnd (defined above, "PDF accessibility"
% comment) wrap both this box *and* the original footline template's own
% content (the theme's \usebeamertemplate{frame footer} call, i.e.
% \FooterText/slide id/page number) -- \addtobeamertemplate's "after" slot
% runs once the wrapped-around original template content has finished, so
% \gothamArtifactEnd there closes the same marked-content span this begins.
% The footer box's own south-west corner originally measured 0.2pt short
% of the page's actual left/bottom edges on every slide checked
% (2026-07-06, direct instruction: "check whether the footer ... has a bit
% of white on its left and below") -- confirmed directly against the
% compiled PDF's own drawn rectangle (consistently ~0.199pt in from both
% edges, not a one-off), unlike the title bar (gotham's own native
% frametitle background, a completely different mechanism -- normal TeX
% box placement, not a tikz "current page" lookup), which sits exactly
% flush at (0,0). At the time, the root cause wasn't chased further and
% the box was fixed by hand-nudging the anchor 0.2pt past the corner and
% padding the size to match. Root cause since found (2026-07-16, while
% chasing the identical hairline on titleframe/quoteframe -- see the
% frametitle hook's own comment above): tikz's default outer sep
% (0.5\pgflinewidth = 0.2pt) puts every edge/corner anchor 0.2pt outside
% the filled rectangle, on *every* edge-anchored node, so the earlier
% "nothing else anchored to a page corner shows the same offset" turned
% out to be wrong. "outer sep=0pt" below is the real fix, replacing the
% nudge-and-pad workaround; the box now lands exactly flush by
% construction rather than by measured compensation.
\addtobeamertemplate{footline}{%
	\gothamArtifactBegin{/Type /Pagination /Subtype /Footer}%
	\vbox to 0pt{%
	\vss
	\begin{tikzpicture}[remember picture, overlay]
		\node[fill=\gothamFooterBoxColor, anchor=south west, inner sep=0pt, outer sep=0pt,
		      minimum width=\paperwidth, minimum height=\gothamFooterBoxHeight]
			at (current page.south west) {};
	\end{tikzpicture}%
	}%
}{\gothamArtifactEnd}

% Footnotes: small font (smaller than the footer's own 5pt), symbol markers
% (*, dagger, ...) instead of arabic numbers, right-aligned (flush with the
% logo's own right edge, like the footer's frame number), no rule above
% them, and tightened-up line spacing.
%
% Line spacing needs two different levers, not just \baselineskip, because
% at this font size (5pt) \baselineskip's normal value is *smaller* than the
% font's own natural line height, so TeX falls back to \lineskip (a flat
% add-on) between every line regardless -- \baselineskip itself turns out to
% have no effect here at all. That fallback \lineskip is what governs
% spacing *within* a single (wrapped) footnote's own lines: reduced only a
% tiny amount, from the default 1pt.
% The *between*-footnote gap is inflated by a separate mechanism entirely:
% \footnotesep, a rule LaTeX inserts before every individual footnote entry
% (measured to be 7.7pt here -- clearly sized for a much bigger font than
% ours), so that's what actually needs the bigger reduction; \parskip
% (tried first) turned out to have no visible effect on it.
\setlength{\footnotesep}{3pt}
% \reset@font\footnotesize is applied by beamer's own footnote machinery
% before this template runs (beamerbaseframecomponents.sty), so the font
% size has to be set again here to actually take effect.
\renewcommand{\footnoterule}{}
% Same marker sequence as LaTeX's own \fnsymbol (*, dagger, double dagger,
% ...), but built from plain Unicode text characters instead: \fnsymbol's
% symbols are typeset via \ensuremath{...}, i.e. in *math* mode, and the
% theme loads fontspec with the "no-math" option (see
% beamerfontthemegotham.sty) specifically so that Fira Sans -- the font
% fontspec puts everything else in -- is left in charge of text only, not
% math; math mode therefore still falls back to LaTeX's default Computer
% Modern math fonts. \fnsymbol's markers were consequently rendering in a
% different font than the surrounding text (checked by comparing rendered
% glyph shapes) -- most visible as a marker that looked shifted left of
% flush-left text next to it (e.g. a \source{} line, which has no marker of
% its own), since Computer Modern's symbol glyphs have different sidebearings
% than Fira Sans's. Plain characters typeset with no \ensuremath wrapper stay
% in text mode and pick up whatever font is already active at the call site,
% i.e. Fira Sans -- matching the rest of the footnote.
% Literal Unicode characters, not the \textdagger/\textdaggerdbl/\textbardbl
% kernel commands: those need the textcomp/TS1 setup to render as plain text,
% and silently fall back to \ensuremath{...} (i.e. math mode, the very thing
% this replaces \fnsymbol to avoid -- see the comment above) when that isn't
% loaded. LuaLaTeX plus fontspec maps raw UTF-8 input straight to the current
% font's (Fira Sans's) own glyphs with no encoding layer in between, which is
% the entire point of using them -- so writing the characters directly, with
% no wrapper command, guarantees they stay in text mode and stay in Fira Sans.
%
% The sequence itself is user-configurable via \setfootnotesymbols{...} (a
% comma-separated list, e.g. \setfootnotesymbols{*,+,#}) -- call it anywhere
% from that point in the deck onward (e.g. once, right after
% % Slidemaster preamble variant for University of Graz decks. Institution-
% specific content only (colors, logo, logo calibration) -- every frame
% type, the deck-version system, footer/hyperlink/footnote styling, tikz
% styles, etc. live in Slidemaster_preamble_base.tex, \input{}ed at the
% bottom of this file, and are shared verbatim across every institution
% this project supports (see Slidemaster_preamble_mci.tex and
% Slidemaster_preamble_uni-ibk.tex for the other two). A deck picks its
% institution by \input{}ing this file instead of one of those, nothing
% else changes.

% Color definition (Graz palette, sourced from the original single-preamble
% version of this project -- these are the University of Graz's own official
% colors, not derived/approximated).
\definecolor{c_primary_dark}{HTML}{000000}
\definecolor{c_primary_highlight}{HTML}{ffd500}
\definecolor{c_primary_gray}{HTML}{c6c6c6}
\definecolor{c_primary_white}{HTML}{FFFFFF}
\definecolor{c_secondary_violet}{HTML}{948dc3}
\definecolor{c_tertiary_violet}{HTML}{c2bddf}
\definecolor{c_secondary_blue}{HTML}{689cce}
\definecolor{c_tertiary_blue}{HTML}{bacce8}
\definecolor{c_secondary_green}{HTML}{9fcb81}
\definecolor{c_tertiary_green}{HTML}{cce2ba}
\definecolor{c_secondary_red}{HTML}{e4684b}
\definecolor{c_tertiary_red}{HTML}{f1ae95}
\definecolor{c_secondary_orange}{HTML}{f6a968}
\definecolor{c_tertiary_orange}{HTML}{fbd0aa}
\definecolor{c_secondary_yellow}{HTML}{ffe063}
\definecolor{c_tertiary_yellow}{HTML}{ffeeaf}
\definecolor{c_tertiary_gray}{HTML}{e3e3e3}
% Not an official Uni Graz color (hence "custom", not "secondary" like the
% palette above) -- a softer near-black than c_primary_dark, used for
% quoteframe's full-bleed dark block, matching the specific dark gray (not
% pure black) ExampleSlides.pptx used for that same block. Also reused for
% hyperlink text color, for the same reason it works for quoteframe's
% attribution text: dark enough for good contrast/legibility on a light
% background.
\definecolor{c_custom_gray}{HTML}{333333}
\definecolor{c_custom_red}{HTML}{C00000}
\definecolor{c_custom_green}{HTML}{00B050}

% Logo + calibration.
% schemes/scheme_uni_graz/Logo_Uni-Graz_4c-cropped.pdf -- official vector artwork,
% user-supplied (schemes/scheme_uni_graz/Logo_Uni-Graz_4c.eps, kept alongside as the
% untouched source), converted to PDF via epstopdf and cropped to its true
% ink bounding box (computed with Ghostscript's own `-sDEVICE=bbox`, same
% treatment as MCI's and Innsbruck's own logos -- see
% Slidemaster_preamble_mci.tex's own comment for the fuller explanation of
% why the EPS's own %%BoundingBox isn't trusted directly). Vector, not
% raster, so it stays crisp at any zoom. Supersedes this file's own earlier
% raster logo (schemes/scheme_uni_graz/Logo Uni Graz 4c-cropped.png, itself cropped
% from schemes/scheme_uni_graz/Logo Uni Graz 4c.png's own symmetric transparent
% margin -- both kept alongside as historical/fallback assets, not deleted).
\newcommand{\gothamInstituteLogoPath}{schemes/scheme_uni_graz/Logo_Uni-Graz_4c-cropped.pdf}
\newlength{\gothamLogoHeight}
\setlength{\gothamLogoHeight}{3ex}
% Re-measured against the compiled PDF for this specific logo image -- see
% Slidemaster_preamble_base.tex's own comment on \gothamFrametitleRightpading
% for why this needs re-measuring per institution rather than reusing
% another one's value.
\setlength{\gothamFrametitleRightpading}{0.28cm}

% Itemize bullet color (2026-07-16, direct instruction): Graz keeps the
% project's original behavior, bullets in the primary highlight color.
\newcommand{\gothamBulletColor}{c_primary_highlight}

% Title bar and footer box/text colors (2026-07-16, direct instruction):
% Graz keeps the project's original behavior throughout -- title/footer
% boxes in the primary highlight color, title/footer text in the primary
% dark color (unlike MCI/Innsbruck, which switch these to their own
% secondary blue/white -- see Slidemaster_preamble_mci.tex's own comment).
\newcommand{\gothamTitleBoxColor}{c_primary_highlight}
\newcommand{\gothamTitleTextColor}{c_primary_dark}
\newcommand{\gothamFooterBoxColorDefault}{c_primary_highlight}
\newcommand{\gothamFooterTextColor}{c_primary_dark}

\input{Slidemaster_preamble_base}
) to replace the default marker set for every
% \footnote{}/\source{} from then on. \seq_item:Nn (expl3's indexed lookup;
% expl3 is already loaded by the gotham theme itself, see e.g.
% beamerfontthemegotham.sty) is what removes the old \ifcase's hard 9-symbol
% ceiling -- a custom list can be any length. Indexing wraps around (via
% \int_mod:nn) rather than erroring past the end of the list, so an unusually
% footnote-heavy slide degrades to repeating symbols instead of a hard
% failure; slides/ZZU-Footnote-symbol-options-current.tex is a reference
% sheet of candidate Unicode symbols to build a custom list from.
%
% Project-wide default (chosen deliberately, not the classic */dagger/etc.
% set): a currency-sign sequence, in this specific order. All 14 were
% confirmed against slides/ZZU-Footnote-symbol-options-current.tex to
% actually exist in Fira Sans.
\ExplSyntaxOn
\seq_new:N \g_gotham_footnote_symbols_seq
\seq_gset_from_clist:Nn \g_gotham_footnote_symbols_seq { £,¥,₹,₽,₺,₴,₪,₫,₦,₡,₩,ƒ,¢,€ }
\newcommand{\setfootnotesymbols}[1]{%
	\seq_gset_from_clist:Nn \g_gotham_footnote_symbols_seq { #1 }
}
\newcommand{\gothamFnSymbol}[1]{%
	\seq_item:Nn \g_gotham_footnote_symbols_seq
		{ \int_eval:n { \int_mod:nn { #1 - 1 } { \seq_count:N \g_gotham_footnote_symbols_seq } + 1 } }
}
\ExplSyntaxOff
\renewcommand{\thefootnote}{\gothamFnSymbol{\value{footnote}}}
% A \footnote{} called inside a minipage doesn't use \thefootnote/\c@footnote
% at all -- LaTeX's own minipage machinery silently switches to a separate
% counter/display command, \c@mpfootnote/\thempfootnote, normally displayed
% via \@alph (a, b, c, ...). Found 2026-07-14 converting
% slides/j6h-Foreign-Exchange-Market-current.tex (its first slide combining
% \footnote{} with a minipage, since a narrowed text column was added
% alongside a corner chart): the footnote *text* still correctly landed in
% this deck's usual bottom-of-frame collection (beamer's own footnote
% storage isn't minipage-scoped), but the *marker*, both inline and at the
% bottom, showed as a plain "a"/"b"/"c" instead of this project's own
% currency-symbol sequence. Redefined the same way as \thefootnote so any
% future \footnote{} inside a minipage picks up the right symbol too --
% \c@mpfootnote resets to 0 for every new minipage on its own (unlike
% \c@footnote, which needs standardframe's own explicit per-frame reset
% below), which happens to already match this deck's "restart every slide"
% footnote convention for the common case of one minipage per frame.
\renewcommand{\thempfootnote}{\gothamFnSymbol{\value{mpfootnote}}}
\setbeamerfont{footnote}{size=\fontsize{3.2}{1.9}\selectfont}
% Beamer's own \@makefnmark (used for the marker both inline in the body
% text and again next to the footnote itself) hard-codes the "footnote
% mark" beamerfont onto it -- an *absolute* size, unrelated to whatever
% text surrounds it -- which is why the inline markers looked tinier than
% a normal superscript when that was set to this footnote's own small size:
% a fixed size that looks fine amid 3.2pt footnote text looks minuscule
% amid much bigger body text. Reverting to plain LaTeX's own \@makefnmark
% instead makes it a normal \textsuperscript of whatever font size is
% already active at each call site -- proportionally scaled either way, so
% it looks like a normal superscript inline, and still comes out small next
% to the (small) footnote text since \usebeamerfont{footnote} above is
% already active there by the time it's invoked.
\makeatletter
\renewcommand{\@makefnmark}{\hbox{\@textsuperscript{\normalfont\@thefnmark}}}
\makeatother
% \raggedleft's own \rightskip (0pt) is flush with the frame's normal body
% text margin, not the logo's edge above it; this pushes further right by
% the difference between the two, empirically calibrated (like
% \gothamFooterRightShift above) against the rendered PDF coordinates of
% both -- re-measure and adjust if the text margins or logo inset change.
\newlength{\gothamFootnoteRightShift}
\setlength{\gothamFootnoteRightShift}{4.8pt}

% Side footnotes: \sidefootnotes (called *before* \begin{standardframe}, like
% \chapterlogo{}) moves that one frame's footnotes -- including \source{} lines,
% which are footnotes underneath -- from their usual spot above the footer to a
% vertical strip along the right edge of the page, rotated 90 degrees counter-
% clockwise: the text starts just above the footer box and reads bottom-to-top,
% as if the usual bottom-right footnote block had been turned on its side (with
% one deliberate difference: the lines are flush with the strip's bottom end,
% not the top -- see the \raggedright comment in the template below).
% standardframe resets the switch
% on \end{frame} (same \global mechanics as \chapterlogo{} -- see the comment
% there for why a plain local assignment doesn't survive beamer's frame group),
% so it opts in only the single frame that follows, and any frame type derived
% from standardframe inherits the behaviour automatically.
%
% How it works: the footnote *template* below branches on this switch. In side
% mode it typesets each footnote into \gothamSideFootnotesBox (a plain box
% register accumulated across the frame's footnotes) instead of contributing
% anything to beamer's own bottom-of-frame footnote insert, and the footline
% hook further down draws that box -- rotated -- as a tikz overlay once the
% whole frame body has run. The footline is the right place to draw from: it is
% typeset at page shipout, i.e. after every \footnote{} in the body has already
% executed, whereas e.g. the frametitle would run too early and find the box
% empty. (Beamer also typesets the footline once per frame just to *measure*
% it, at \begin{frame} time -- before any footnotes exist -- which is why the
% drawing hook guards on \ifvoid: the measurement pass must not consume, or
% try to draw, a box that doesn't exist yet.)
%
% Known quirk: a footnote with an overlay spec (\footnote<2->{...}) is
% collected -- and therefore shown -- on *every* overlay slide of its frame in
% side mode. Beamer implements covered footnotes by typesetting them invisibly
% in place, which relocating the text into a box register defeats.
\newif\ifgothamSideFootnotes
\newcommand{\sidefootnotes}{\global\gothamSideFootnotestrue}
\newbox\gothamSideFootnotesBox
% The rotated block's line width, i.e. the vertical run the strip may occupy
% down the right edge of the page: from just below the title bar to just above
% the footer's yellow box. Like \gothamFooterBoxHeight above, calibrated
% against the measured PDF coordinates: the title bar's bottom edge sits
% 219.57pt above the page's bottom edge; minus \gothamSideFootnoteTopSep (3pt),
% \gothamFooterBoxHeight (7.26pt), and ~2pt of clearance above the footer box
% leaves 207pt. Re-measure if the frametitle font or its padding ever changes.
\newlength{\gothamSideFootnoteHsize}
\setlength{\gothamSideFootnoteHsize}{207pt}
% Gap between the title bar's bottom edge and the top end of the strip. The
% text itself sits at the strip's *bottom* end (it is \raggedright, i.e.
% flush with the line start, which the rotation turns into bottom-aligned),
% so this gap only matters for footnotes long enough to fill the whole run.
\newlength{\gothamSideFootnoteTopSep}
\setlength{\gothamSideFootnoteTopSep}{3pt}
% Distance from the page's right edge to the strip's outer edge -- the rotated
% counterpart of the small gap the bottom footnotes keep above the footer box.
\newlength{\gothamSideFootnoteRightInset}
\setlength{\gothamSideFootnoteRightInset}{2pt}
% A real footnote's marker (e.g. "*") sits in an \hbox of its own whose
% *ink* -- unlike \source{}'s marker, which is emptied to nothing and so has
% no ink at all -- extends slightly left of that hbox's own nominal left
% edge (ordinary optical behaviour for a small centred glyph like an
% asterisk, exaggerated further by \@makefnmark's superscript raise). With
% every footnote's paragraph starting flush at the same \noindent left edge,
% that overhang made real footnotes' markers visibly poke out to the left of
% \source{}'s flush-left text above/below them. This shift nudges real
% footnotes' *entire* line (marker, comma, and text together) right by that
% overhang, so the marker's ink -- not the text after it -- lines up with
% \source{}'s left edge; \source{} itself is untouched; empirically
% calibrated against the rendered PDF, like the other shift constants here.
\newlength{\gothamSideFootnoteMarkerShift}
\setlength{\gothamSideFootnoteMarkerShift}{1.3pt}

% A new template variant (rather than reusing beamer's own "default", which
% it has already defined and then refuses to redefine) that we immediately
% activate below.
% Wrapped in \makeatletter (unlike the rest of this section) because the side
% branch below tests \@thefnmark against \@empty to tell a real footnote
% (non-empty marker) apart from a \source{} call (marker forced empty) --
% see \gothamSideFootnoteMarkerShift above.
\makeatletter
\defbeamertemplate{footnote}{gothamfn}{%
	\ifgothamSideFootnotes
		% Side mode (see \sidefootnotes above): append this footnote to the
		% collector box instead of typesetting it here. \insertfootnotetext
		% carries beamer's own \rule\z@\footnotesep lead-in and
		% \@finalstrut\strutbox tail along with it (beamer bakes both into the
		% text it hands the template), so the between-footnote spacing and the
		% strutbox rebuild work exactly as in the bottom branch below -- and,
		% equally, nothing of the footnote (not even the spacing rule) is left
		% behind in beamer's bottom-of-frame insert, which therefore stays
		% visually empty. No \gothamFootnoteRightShift here: that constant
		% compensates the *bottom* block's \raggedleft for the text-margin-vs-
		% logo-edge difference, while the strip's position is set directly by
		% the tikz node in the footline hook.
		\global\setbox\gothamSideFootnotesBox\vbox{%
			\unvbox\gothamSideFootnotesBox
			\hsize\gothamSideFootnoteHsize
			\usebeamerfont{footnote}%
			\setlength{\lineskip}{0.6pt}%
			\settoheight{\dimen0}{Xy}%
			\settodepth{\dimen1}{Xy}%
			\setbox\strutbox\hbox{\vrule height\dimen0 depth\dimen1 width0pt}%
			\parindent 1em\noindent%
			% \raggedright, not the bottom block's \raggedleft: after the
			% rotation, the line start is the *bottom* end of the strip, so
			% this makes the text begin just above the footer box and run
			% upward (rather than hang down from the title bar).
			\raggedright
			% \source{} forces \thefootnote (hence \@thefnmark) empty, so this
			% only fires for a real footnote's marker -- see
			% \gothamSideFootnoteMarkerShift above for why it needs the shift
			% and \source{} doesn't. The "\," below (a separator between the
			% marker and the footnote text) needs the same guard: it used to
			% run unconditionally, so an empty-marker \source{} still got a
			% thin space's worth of leading indent before its own text --
			% invisible under the *bottom* block's \raggedleft (nothing there
			% depends on the left edge), but a real, visible defect here under
			% \raggedright, where the left edge is exactly what's on display:
			% a wrapped \source{}'s first line sat measurably indented from
			% its own second line, which never picks up that one-time leading
			% space (2026-07-12, direct instruction: "the first line seems
			% indented compared to the second, which it shouldn't be").
			\ifx\@thefnmark\@empty
				\insertfootnotetext\par%
			\else
				\hspace*{\gothamSideFootnoteMarkerShift}\insertfootnotemark\,\insertfootnotetext\par%
			\fi
		}%
	\else
	\usebeamerfont{footnote}%
	\setlength{\lineskip}{0.6pt}%
	% beamer's footnote machinery resets to plain \footnotesize (a much
	% bigger font) before this template runs, and switching to our own
	% smaller font via \usebeamerfont doesn't update \strutbox to match --
	% so the invisible strut beamer appends after every footnote
	% (\@finalstrut\strutbox) is still sized for \footnotesize, silently
	% padding out extra depth below the last line. Rebuilding \strutbox
	% from this font's own measured metrics removes that padding, which is
	% most of the gap between the footnotes and the footer box below them.
	\settoheight{\dimen0}{Xy}%
	\settodepth{\dimen1}{Xy}%
	\setbox\strutbox\hbox{\vrule height\dimen0 depth\dimen1 width0pt}%
	\parindent 1em\noindent%
	\raggedleft
	\advance\rightskip by -\gothamFootnoteRightShift%
	% Same \@thefnmark emptiness guard as the side-footnote branch above,
	% for the same reason (an empty marker still leaving behind the "\,"
	% separator's own leading thin space) -- invisible here under
	% \raggedleft, where the left edge is never on display, but the same
	% underlying defect, so fixed the same way for consistency.
	\ifx\@thefnmark\@empty
		\insertfootnotetext\par%
	\else
		\insertfootnotemark\,\insertfootnotetext\par%
	\fi
	\fi
}
\makeatother
\setbeamertemplate{footnote}[gothamfn]

% Draws the side-footnote strip (see \sidefootnotes above). Lives in the
% footline template because that is typeset at page shipout, after the frame
% body -- and with it every \footnote{} call -- has run, so the collector box
% is complete by the time this executes. The \ifvoid guard doubles as the
% protection against beamer's measure-the-footline pass at \begin{frame} time
% (see the \sidefootnotes comment). \box (not \copy) empties the register
% globally as a side effect, which is exactly right: each overlay slide of a
% frame re-executes the body and re-collects its footnotes from scratch, so
% whatever this slide's footline draws must not leak into the next slide.
%
% Node geometry: tikz pins the node's *unrotated* anchor corner at the given
% coordinate and rotates the box counter-clockwise around it, so anchor=south
% east (the end of the collected block's last line) ends up as the strip's
% top-east corner: from there the block extends down the page (its lines run
% bottom-to-top, last footnote nearest the page edge), with the \raggedright
% text sitting at the strip's bottom end, just above the footer box (see the
% comment in the footnote template above). The y-coordinate
% comes from the same remembered gothamTitleBottom node the title-bar patches
% already use (same two-pass .aux mechanics -- see "Title bar and logo").
% The \ifcsname guard skips the drawing while the gothamTitleBottom shape is
% entirely unknown to pgf (i.e. no frametitle has been typeset yet in this
% run *and* nothing was inherited from a previous run's .aux) -- without it,
% referencing the node is a hard error rather than the usual harmless
% first-pass misplacement. In that situation the position would have been
% garbage anyway, and the frame-start clearing in standardframe below mops up
% the box this branch leaves unconsumed.
\makeatletter
\addtobeamertemplate{footline}{%
	\ifvoid\gothamSideFootnotesBox\else
		\ifcsname pgf@sh@ns@gothamTitleBottom\endcsname
		\vbox to 0pt{% zero-space trick, as in the yellow-box footline hook above
		\vss
		\begin{tikzpicture}[remember picture, overlay]
			\node[rotate=90, anchor=south east, inner sep=0pt]
				at ($(gothamTitleBottom -| current page.east)+(-\gothamSideFootnoteRightInset,-\gothamSideFootnoteTopSep)$)
				{\box\gothamSideFootnotesBox};
		\end{tikzpicture}%
		}%
		\fi
	\fi
}{}
% Empties the side-footnote collector (used by standardframe at frame start).
% Defined here, inside the \makeatletter block, because \voidb@x needs @ to be
% a letter at tokenization time -- standardframe's own definition below sits
% outside any \makeatletter region.
\newcommand{\gothamClearSideFootnotesBox}{\global\setbox\gothamSideFootnotesBox=\box\voidb@x}
\makeatother

% \source{text} -- for a "Source: ..."/"Image source: ..."/etc. citation line
% (common at the bottom/edge of a converted PowerPoint slide) that should look
% and sit *exactly* like the footnotes above -- same font, spacing,
% right-alignment, and the same stretch-to-sit-just-above-the-footer
% behaviour beamer already gives real footnotes -- but with no marker, neither
% inline in the body text nor next to the note itself, since a source line
% isn't anchored to one specific word. Implemented as a real \footnote{} (so
% it gets that positioning for free, rather than reimplementing it) with
% \thefootnote emptied locally; since \footnote{} still advances the footnote
% counter even with an empty marker, \addtocounter reverses that so a real
% \footnote{} elsewhere on the same slide doesn't skip a symbol in the
% \fnsymbol sequence. Unlike a real footnote, no label is added automatically
% here -- the label varies ("Source:", "Image source:", ...), so write it as
% part of #1.
\newcommand{\source}[1]{%
	\begingroup
	\renewcommand{\thefootnote}{}%
	\footnote{#1}%
	\endgroup
	\addtocounter{footnote}{-1}%
}

% A plain tabular with yellow \toprule/\midrule/\bottomrule, matching the
% frametitle/footer accent color. Takes the same column-spec argument as
% tabular itself, e.g. \begin{gothamtable}{lccc} ... \end{gothamtable}. Only
% booktabs rules are colored (via \arrayrulecolor, scoped to this table by
% \begingroup/\endgroup) -- if a table ever needs a plain black \hline too,
% that's unaffected since \arrayrulecolor only recolors rules drawn *after*
% it takes effect, and booktabs' rules read it fresh each time.
%
% RULE-BLEED WARNING (2026-07-12, direct instruction: "make sure the table
% lines do not extend beyond the text"), read this before writing a new
% gothamtable: a plain column spec like `{lccc}` above leaves the default
% \tabcolsep (6pt, or whatever a slide-local \setlength last set it to) as
% *implicit* padding on the table's own outermost edges -- invisible to the
% table's own text (which starts/ends flush with each column's actual
% content), but very visible to booktabs' \toprule/\midrule/\bottomrule and
% to colortbl's \rowcolor/\cellcolor, both of which span the *declared*
% column width including that implicit padding. The result, confirmed
% directly and independently on at least three slides so far
% (slides/u5h-Long-term-Bonds-Provisions-current.tex's own v1-1-0/v1-4-0,
% slides/c4c-Macbeth-Spot-Removers-Proposed-Structure-current.tex's own
% v1-2-2, slides/tw5-Exercise-Leverage-and-ROE-Solution-current.tex's own
% v1-4-0): every rule/highlight bleeds \tabcolsep past the visible text on
% both the left and right edges, most noticeable on a left-aligned
% (non-\centering) table where the mismatch is easy to see against the
% title's own left margin.
%
% Not fixed here in the environment definition itself: doing so would
% change every existing gothamtable's own column spacing out from under it
% (this environment has no way to tell "plain default spacing was intended"
% apart from "the bleed was never noticed"). Fix per table instead, the
% same three-part pattern every slide above uses -- worth reaching for by
% default on any *new* gothamtable, not just when the bleed is spotted:
%   {\setlength{\tabcolsep}{0pt}
%   \begin{gothamtable}{@{}l<{\hspace{<gap>pt}}c<{\hspace{<gap>pt}}c@{}}
%   ...
%   \end{gothamtable}}
% -- \tabcolsep zeroed (scoped locally, e.g. by wrapping the whole call in
% a `{...}` group, so it doesn't leak into later tables), explicit `@{}` at
% both outer edges of the column spec (so booktabs/colortbl have no
% implicit padding left to bleed into), and array's `<{\hspace{<gap>pt}}>`
% column modifier appended after every column but the last, to put back an
% explicit, *content-column* gap (matching whatever `2*\tabcolsep` the
% table used to have, if preserving its exact look, e.g. `<{\hspace{12pt}}>`
% for the ordinary 6pt-\tabcolsep default) -- `<{...}>` puts the gap
% *inside* the preceding column's own content, which `\rowcolor` (if used)
% colors along with everything else, unlike a plain `@{...}` separator,
% which `\rowcolor` deliberately leaves untinted (see u5h's own v1-6-0
% comment for that exact failure mode if a highlighted table ever needs
% this too).
%
% RE-CONFIRMED 2026-07-16, direct instruction, after this exact bleed
% reappeared on essentially every gothamtable in a freshly-converted
% chapter (Capital Budgeting Techniques) despite the warning above already
% existing at the time -- the rule above is not optional/only-if-noticed,
% it's the default way to write *any* new gothamtable, full stop. Check
% every new gothamtable against it before considering a slide done, not
% just when a rendered page happens to make the bleed visually obvious.
%
% BLOCK-RULE-COLOR CONVENTION (2026-07-16, direct instruction): a
% gothamtable placed *inside* any block environment (`infoblock`,
% `exampleblock`, `referenceblock`, `definitionblock`, plain `block`, or
% their `gothamfloating*` wrappers) should use black rules, not this
% environment's own default yellow -- yellow-on-the-block's-own-tinted-
% background reads poorly, unlike yellow-on-plain-white outside a block.
% Override locally, *inside* `\begin{gothamtable}{...}`, before `\toprule`
% (not before `\begin{gothamtable}` itself -- the environment's own
% opening code sets yellow *after* its `\begingroup`, so a color set
% before that point is simply overwritten; confirmed directly, first
% attempt this same day):
%   \begin{gothamtable}{...}
%   	\arrayrulecolor{c_primary_dark}
%   	\toprule
%   	...
%   \end{gothamtable}
% Booktabs rules read `\arrayrulecolor` fresh at each rule, and
% `\endgroup` (gothamtable's own closing code) resets it back to yellow
% afterward, so this never leaks into a later table. First applied on
% slides/6ux-Example-51-BMA-Payback-Table-current.tex's own
% `gothamfloatinginfoblock`.
\newenvironment{gothamtable}[1]{%
	\begingroup
	\arrayrulecolor{c_primary_highlight}%
	\begin{tabular}{#1}%
}{%
	\end{tabular}%
	\endgroup
}

% For a side-corner image overlay (the usual \vbox to 0pt{...} tikz overlay
% pattern used throughout this deck's converted slides) whose *top edge*
% should line up with the content text's own top edge -- distance from the
% page's own top edge to the cap-height top of a standardframe's first body
% line, measured directly against the compiled PDF's own text coordinates,
% not derived from \gothamContentTopShift or guessed. Previous per-slide
% image placements each guessed their own downward shift (e.g. a bare
% "-1.3cm") independently -- close, but not exact, landing the image's top
% a couple mm above the text's own top; use this constant instead of a
% fresh guess for any new "text top edge"-aligned image.
%
% 41.58pt (2026-07-22, re-measured): the previous value, 44.0pt, went stale
% the moment \gothamContentTopShift itself was increased (7.33pt -> 11.5964pt,
% "reduce the Top Margin", see that length's own comment above) without a
% matching update here -- that change pulls plain body text up, but this
% length is a separate, hand-measured constant with no formula tying it to
% \gothamContentTopShift, so it kept anchoring images at the *old*, lower
% position. Confirmed directly (pymupdf against this same freshly-compiled
% preview deck, not assumed from the \gothamContentTopShift delta alone):
% three independent plain-first-line standardframes with no per-slide
% top-shift correction (slides/5v4-CML-vs-SML-current.tex,
% slides/u5j-Exercise-CAPM-4-current.tex,
% slides/nt4-Exercise-CAPM-5-current.tex) all land their first line's own
% top edge at y=41.58pt post-change, not 44.0pt. Re-measure and adjust
% again if \gothamContentTopShift ever changes in the future.
\newlength{\gothamContentTopAlign}
\setlength{\gothamContentTopAlign}{41.58pt}

% For a "side filler image" -- a decorative image stuck flush to the page's
% own right edge (zero margin) with its top edge flush against the title
% bar's own bottom edge (zero gap), rather than sitting in the usual
% content-aligned corner inset -- distance from the page's own top edge
% down to the standardframe title bar's bottom edge, measured directly
% against the compiled PDF (a pixel scan for the yellow-to-white
% transition, at this deck's 16cm-wide/9cm-tall page size). First used on
% slides/dfc-Angel-Investors-current.tex.
%
% 25.43pt -- the raw measured value (originally 25.44pt; re-measured
% 2026-07-07 against the current compiled PDF via the title bar's own
% white-patch rectangle, which came back 25.427pt, close enough to the
% original measurement to be the same underlying edge, not a drift).
% Briefly set to 24.94pt instead (2026-07-06-ish), a deliberate 0.5pt
% overlap *into* the title bar, after a report of a visible hairline
% white gap right at this seam in one render -- not reproduced in this
% project's own pdftoppm-based rendering checks at the time (every pixel
% scan at the raw value showed a clean cut, no gap), so assumed to be a
% renderer-specific sub-pixel rounding artifact rather than a real error
% in the constant, and "fixed" by deliberately overlapping instead of
% chasing the specific viewer. That overlap turned out to be the wrong
% call: direct feedback 2026-07-07 reported the *opposite* symptom on
% `sideimageframe` (the image's own top edge visibly overlapping the
% title box) -- exactly what a real 0.5pt overlap looks like once it's
% actually noticed, rather than a rendering artifact. Reverted back to
% the precise raw measurement: a real, small, deliberate overlap chosen
% to paper over an unreproduced gap in an unspecified other renderer
% turns out to be a worse trade than just using the measured value and
% letting a genuine future gap report (if one ever recurs, in a
% specific, identified renderer) be diagnosed on its own terms.
%
% Re-measured 2026-07-16, 25.43pt -> 25.32pt: the frametitle hook's
% outer sep fix (see its own comment) moved the visible yellow bottom
% edge up by the 0.2pt the trim patch's fill had previously been inset
% below its gothamLogoAnchor.south anchor -- the yellow now genuinely
% ends flush with the logo's bottom edge, as always intended. New raw
% measurement against the rebuilt PDF: the trim patch's fill top edge
% (= the visible yellow bottom) sits 25.228 PDF units below the page
% top, i.e. 25.3226pt (earlier measurements in this comment quoted PDF
% coordinate values directly as "pt"; PDF units are big points, 1bp =
% 1.00375pt -- a systematic ~0.1pt understatement at this magnitude,
% harmless then, but corrected now that the constant is being reset
% anyway). Note this value is also what the overlay-logo nodes
% (titleframe/quoteframe/the hook's fixed-position redraw) derive their
% vertical inset from: visible-bar-bottom minus \gothamLogoHeight is
% exactly the inline logo's own vertically-centered top margin.
\newlength{\gothamTitleBarBottomAlign}
\setlength{\gothamTitleBarBottomAlign}{25.32pt}

% Note on the [t] default below: it top-aligns content within the frame
% (beamer's own default is vertical centering, "c"), so slides don't have
% their content floating in the middle. It also sidesteps a real beamer bug:
% \begin{frame}[]{Title} (a literal empty options bracket) fails outright,
% even in plain beamer with no theme at all -- so the optional argument must
% never default to an empty string. The title is always passed as an
% explicit argument (rather than left for beamer to pick up from a trailing
% brace group), since anything placed between \begin{frame} and its title --
% like our \setbeamertemplate calls -- breaks that trailing-brace mechanism.

% Pulls the frame's content up closer to the title bar: measured (via the
% PDF's own coordinates) at about 22pt of gap by default between the title
% bar's bottom edge and the first line of body content: this is a third of
% that, shrinking the gap to about two thirds of its original size.
%
% Called "Top Margin" in conversation from 2026-07-12 onward (direct
% instruction) -- the deck-wide name for the resulting gap between the
% title bar and a standardframe's first line of body content, i.e. the
% vertical position a plain-text frame's first baseline naturally lands at
% (measured directly against a reference slide's own compiled PDF, e.g.
% slides/aii-Leverage-Effects-current.tex's first line). Any first content
% that isn't plain text (a table, a minipage, an itemize with an unusual
% first item, ...) needs its own per-slide \gothamIncreaseContentTopShift/
% \gothamReduceContentTopShift correction to land on this same baseline --
% not portable between slides even for superficially similar content (see
% slides/zgz-EBIT-Definition-current.tex's own v1-2-0 comment for a case
% where an earlier, wrong correction was carried over uncritically).
\newlength{\gothamContentTopShift}
% 11.5964pt (2026-07-22, direct instruction: "reduce the Top Margin"
% generally, tested at 20%/30%/40% smaller against
% slides/8s8-Unit-Investment-Trusts-current.tex before settling on 40% --
% up from the original 7.33pt so the *resulting* gap (this length pulls
% content up by, so a bigger value means a smaller visible margin) comes
% out to 60% of its old size: measured directly on 8s8 at 10.67pt before,
% 6.42pt after (target was 6.40pt, within 0.02pt). Since every
% standardframe applies this same pull-up before any per-slide
% \gothamIncreaseContentTopShift/\gothamReduceContentTopShift correction,
% changing it here shifts the baseline project-wide, not just on 8s8.
\setlength{\gothamContentTopShift}{11.5964pt}

% "Top Margin" (above) is one of four deck-wide margin names introduced in
% conversation from 2026-07-12 onward (direct instruction) for the plain
% page-edge/title-bar/footer gaps every slide is expected to obey by
% default. The other three, for reference (each maps to a length already
% declared elsewhere, not a new one introduced here):
% - "Left Margin": the default gap from the page's own left edge to the
%   title text and body content -- \gothamFrametitleLeftpading
%   (beamerouterthemegotham.sty/beamerthemegotham.sty, 2.2ex, measured
%   directly against the compiled PDF at 12.552pt in this deck's usual
%   fonts -- see e.g. slides/tw5-Exercise-Leverage-and-ROE-Solution-current.tex's
%   own v1-1-0 comment for that same measured value).
% - "Right Margin": the default gap from the page's own right edge to the
%   institute logo (and, by the same convention, to anything else
%   deliberately flush-right, e.g. a corner-anchored image) --
%   \gothamFrametitleRightpading (0.28cm/7.93pt, this file's own
%   \gothamInstituteLogoSquare placement above uses it directly as the
%   logo's own inner sep).
% - "Bottom Margin": the default gap between a frame's own body content
%   and the \source{}/footnote line below it. Unlike the other three, this
%   one has no single settable length behind it -- \source{}'s own
%   vertical position isn't fixed the way the footer box is; it tracks the
%   frame's actual content height once that content gets tall enough (see
%   slides/u5h-Long-term-Bonds-Provisions-current.tex's own v1-3-0 comment,
%   "past a certain table height it starts tracking the table's own growth
%   roughly 1:1"). "Respecting Bottom Margin" in practice means leaving
%   comparable clearance to whatever a reference slide with unremarkable
%   content already shows, the same empirical, measure-against-a-reference-
%   slide approach "Top Margin" itself uses, rather than a fixed constant.

% Compensates a gap that opens up specifically when a standardframe's body
% starts with a \begin{minipage}{...} (used to width-limit text next to an
% image/diagram overlay, e.g. slides/igr-Budgeting-current.tex,
% slides/bdh-Working-Capital-Example-current.tex -- see README.md's
% "Available layouts" section for why minipage instead of columns there) --
% a minipage is a *box*, and TeX's vertical glue insertion (\topskip
% specifically) treats the first box in a vertical list differently from
% the first line of plain paragraph text, so the same \gothamContentTopShift
% pull-up above lands closer to the title for minipage-wrapped content than
% for plain content. Cannot be fixed by \vspace(*) right after
% \begin{minipage}{...}: both plain and starred forms were tried and
% measurably had zero effect (the star only protects glue from page-break
% discard, not from \topskip's own separate glue-insertion calculation,
% which absorbs leading glue at the start of a vertical list regardless).
% Fixed instead by shrinking standardframe's own pull-up for the affected
% frame -- see \gothamContentTopShiftReduction below. The value below is
% calibrated empirically, not derived from first principles: an isolated
% two-frame test (single-level itemize, one plain, one minipage-wrapped)
% measured only a 4.4pt gap, but real slides (nested itemize inside the
% minipage) showed a larger ~7.5-7.6pt gap against a real plain-itemize
% slide (slides/jyn-Categorization-current.tex) -- the extra few points
% likely come from itemize's own topsep/partopsep being computed differently
% depending on whether the list opens a box or a paragraph, on top of the
% topskip/box effect alone. Tuned so slides/igr-Budgeting-current.tex and
% slides/bdh-Working-Capital-Example-current.tex's first content line
% lands within 0.1pt of slides/jyn-Categorization-current.tex's own
% (200.917pt vs. 200.871pt/200.980pt in the compiled PDF).
\newlength{\gothamMinipageTopCorrection}
% 9.38pt (2026-07-22, re-measured): the previous 7.6pt value went stale
% the moment \gothamContentTopShift itself was increased (7.33pt ->
% 11.5964pt, "reduce the Top Margin", see that length's own comment above)
% without a matching update here -- confirmed directly (pymupdf against a
% freshly-compiled preview deck, not assumed from the \gothamContentTopShift
% delta alone, which doesn't translate 1:1 into a minipage's own \topskip
% behavior): four independent slides all using this constant completely
% unmodified (slides/fr3-Forms-of-Capital-Market-Efficiency-Strong-Semistrong-Weak-current.tex,
% slides/za2-Startup-Financing-Categorization-current.tex,
% slides/yke-Implications-of-Homemade-Dividends-current.tex,
% slides/6j4-Modigliani-Miller-Assumptions-current.tex) all land their
% own first line at an identical y=39.80pt post-change, 1.78pt short of
% this deck's current 41.58pt Top Margin reference -- not slide-specific
% noise, the same deterministic amount on every one of them. Re-measure
% and adjust again if \gothamContentTopShift ever changes in the future.
\setlength{\gothamMinipageTopCorrection}{9.38pt}

% A *third* minipage-top correction, on top of \gothamMinipageTopCorrection
% above rather than replacing it, for a minipage whose own first internal
% content is a zero-height `\vbox to 0pt{...}` tikz overlay (a
% corner-anchored decorative image, placed there specifically so it stays
% unconditionally visible even when the minipage's own later content
% includes a `\pause` -- see e.g.
% slides/y5z-CoCo-Bonds-current.tex's own "v1-2-0" note for the fuller
% pause-inheritance story behind why the image has to sit there at all).
% 2026-07-08, found while chasing "the top margin of the text on y5z/22s
% is too low" -- measured directly against the compiled PDF, y5h/22s both
% landed 16.5pt too far down despite already using
% \gothamMinipageTopCorrection, with the zero-height image overlay sitting
% *before* the minipage entirely as the actual root cause (that ordering
% makes the image, not the minipage, the frame body's true first box,
% which \topskip's own special first-box glue handling reacts to
% differently than a plain first minipage). Moving the image inside the
% minipage as its own first content closed most of the gap on its own
% (restoring the minipage as the frame's true first box) but left a
% further 10.35pt uncorrected -- confirmed the same on two independent
% slides (y5z, 22s) with identical structure, not assumed transferable
% from one to the other. Call via
% `\gothamReduceContentTopShift{\dimexpr\gothamMinipageTopCorrection-\gothamMinipageWithCornerImageExtraCorrection\relax}`
% for any slide combining this exact pattern (corner image as a minipage's
% own first content, ahead of a `\pause`) -- confirmed transferable as-is
% to two more slides with the same structure
% (slides/qbj-Bank-Loans-Credit-Lines-current.tex,
% slides/jkf-Bank-Loans-Term-Loan-current.tex), neither of which
% actually uses `\pause` itself, so this isn't specific to that -- the
% image-inside-minipage-as-first-content shape alone is what triggers it.
\newlength{\gothamMinipageWithCornerImageExtraCorrection}
\setlength{\gothamMinipageWithCornerImageExtraCorrection}{10.35pt}

% A *different* minipage-top correction, specifically for a minipage+
% itemize sitting as the first content of a `sideimageframe` (2026-07-06,
% found while chasing "the top margin of the text on ezz/acw is again
% off" -- see CONVERSION_NOTES.md's own "beamer/TeX gotchas" entry for the
% full story). `\gothamMinipageTopCorrection` above is calibrated for a
% *plain* standardframe, and simply does not transfer to `sideimageframe`:
% measured directly on three real sideimageframe slides
% (slides/xmp-Implications-of-Capital-Market-Efficiency-Considerations-current.tex,
% slides/dfc-Angel-Investors-current.tex,
% slides/ezz-Examples-of-Tests-for-Market-Efficiency-current.tex --
% one flat itemize, one flat itemize, one nested-first-item itemize
% respectively, so this isn't a flat-vs-nested effect either), all three
% landed within 0.6pt of each other at their own *uncorrected* natural
% position (52.03-52.57pt against this deck's usual 44.0pt first-line
% reference) regardless of which of the two *wrong* corrections
% (`\gothamReduceContentTopShift{\gothamMinipageTopCorrection}` or
% `\gothamIncreaseContentTopShift{\gothamFlatItemizeTopCorrection}`) had
% been applied to each beforehand -- i.e. `sideimageframe` itself pushes
% its own first-content minipage down by a consistent ~8.2pt relative to
% a plain standardframe, on top of whatever generic itemize-shape
% correction would otherwise apply, most likely from the extra
% \sidefootnotes/\settowidth setup work `sideimageframe`'s own opening
% code does before `\begin{standardframe}` ever runs. Use
% `\gothamIncreaseContentTopShift{\gothamSideImageFrameMinipageCorrection}`
% (note: *increase*, not reduce) as the *only* topshift call on any
% future sideimageframe slide whose first content is a minipage+itemize --
% not the generic `\gothamMinipageTopCorrection`/
% `\gothamFlatItemizeTopCorrection` this deck uses everywhere else.
\newlength{\gothamSideImageFrameMinipageCorrection}
\setlength{\gothamSideImageFrameMinipageCorrection}{8.2pt}

% 6.26pt (2026-07-22, re-measured): stale for the same reason
% \gothamMinipageTopCorrection above went stale -- \gothamContentTopShift
% itself changed (7.33pt -> 11.5964pt) without a matching update here.
% Re-measured directly against a freshly-compiled preview deck across six
% independent sideimageframe slides using this constant unmodified: four
% (ezz, xmp, hzx, u6n) land their own first line at an identical
% y=39.64pt, 1.94pt short of this deck's current 41.58pt Top Margin
% reference (i.e. already pulled up *too far*); the other two (c7w, dfc)
% are close but not identical (39.53pt/-2.05, 40.18pt/-1.4) -- no single
% value fixes all six exactly, so this recalibrates to the four-way
% cluster (the majority case) and accepts a small (<0.6pt) residual on the
% two outliers. Since this constant feeds
% `\gothamIncreaseContentTopShift{}` (more pull-up, i.e. *smaller* y), and
% the observed y was already short of target, the fix is a *decrease*
% (8.2pt - 1.94pt), not an increase -- caught by re-measuring after a
% first, wrong-direction attempt at 10.14pt drove the four-slide cluster
% to y=37.70pt, further from target instead of closer.
\setlength{\gothamSideImageFrameMinipageCorrection}{6.26pt}

% Opt-in override: reduces standardframe's \gothamContentTopShift pull-up by
% this amount for the frame currently being opened. Empty-default-style flag
% (mirrors \gothamChapterLogoPath etc.): 0pt by default (no effect), set via
% \gothamReduceContentTopShift{<length>} called *before* \begin{standardframe},
% reset back to 0pt in standardframe's own closing code so it never leaks
% into the next frame. Used by minipage-based slides to counteract
% \gothamMinipageTopCorrection above.
\newcommand{\gothamContentTopShiftReduction}{0pt}
\newcommand{\gothamReduceContentTopShift}[1]{\global\def\gothamContentTopShiftReduction{#1}}

% A *second*, separate first-line-position quirk, found and fixed
% 2026-07-05 after a run of slides (batch 4 plus several earlier ones) were
% found to land at two visibly different heights depending purely on
% content shape, with no minipage involved at all: a plain (non-minipage)
% standardframe whose top-level itemize's very *first* \item is immediately
% followed by its own nested \begin{itemize}/\begin{enumerate} lands 1.834pt
% higher than one whose first \item has no such immediate child (checked
% against every content slide converted so far via each compiled page's own
% first-line PDF baseline coordinate, not eyeballed -- 13 "nested-first-item"
% slides all landed within 0.05pt of each other, and 9 "flat-first-item"
% slides all landed within 0.05pt of each other, at a clean +1.834pt offset
% from the first group; isolated two-frame test cases, differing only in
% presence/absence of that immediate nested list, reproduced the exact same
% 1.834pt gap independent of which text/list type/item count was used).
% Not chased to one named TeX primitive (most likely some interaction of
% \addvspace's glue-coalescing behaviour with \partopsep/\topsep at the
% \begin{itemize}/\begin{frame} boundary, given the existing
% \gothamMinipageTopCorrection comment above already flags topsep/partopsep
% as varying with box vs. paragraph context) -- calibrated empirically like
% \gothamMinipageTopCorrection, not derived from first principles.
% Interacts with minipage differently, not additively: inside a minipage,
% this nested-vs-flat distinction disappears entirely (both land within
% 0.05pt of each other) -- so this correction is for *plain* itemize-first
% standardframes only, never combined with \gothamMinipageTopCorrection.
\newlength{\gothamFlatItemizeTopCorrection}
\setlength{\gothamFlatItemizeTopCorrection}{1.834pt}

% 0.01pt (2026-07-22, re-measured): stale for the same reason
% \gothamMinipageTopCorrection above went stale -- \gothamContentTopShift
% itself changed (7.33pt -> 11.5964pt) without a matching update here.
% Re-measured directly against three independent flat-itemize slides using
% this constant unmodified (574, 5cr, c8e): all three landed their own
% first line at an identical y=39.76pt, 1.82pt short of this deck's
% current 41.58pt Top Margin reference -- coincidentally, the recalibrated
% \gothamContentTopAlign reference now already accounts for almost all of
% what this correction used to supply on its own, leaving next to nothing
% left to add. Effectively zeroed out rather than removing the call
% itself, so every flat-itemize slide's own
% `\gothamIncreaseContentTopShift{\gothamFlatItemizeTopCorrection}` line
% keeps working unchanged.
\setlength{\gothamFlatItemizeTopCorrection}{0.01pt}

% Opposite-direction counterpart to \gothamReduceContentTopShift: *increases*
% the pull-up instead of reducing it (needed for \gothamFlatItemizeTopCorrection
% above, since a flat first item needs *more* pull-up, not less). Same
% storage variable, just negated -- \gothamContentTopShiftReduction is
% allowed to go negative, since the frame-opening code only ever computes
% \gothamContentTopShift-\gothamContentTopShiftReduction, never checks its
% sign.
\newcommand{\gothamIncreaseContentTopShift}[1]{\global\def\gothamContentTopShiftReduction{-#1}}

% Define custom frame type with en-dash bullets
\newenvironment{standardframe}[2][t]{%
	% Restart footnote numbering (the *, dagger, ... symbols) fresh on every
	% slide, rather than letting it run on cumulatively across the whole deck
	% the way beamer does by default. Two reasons: \fnsymbol only defines 9
	% symbols before erroring out, which a real multi-slide deck would run
	% through in no time if every \footnote{}/\source{} call across the whole
	% talk drew from one shared counter; and each slide's footnotes are read
	% and referenced only on that one slide anyway, so there's no reason for
	% their numbering to depend on how many footnotes earlier slides happened
	% to have. This must run *before* \begin{frame}: beamer saves the
	% counter's current value as that frame's baseline as part of
	% \begin{frame}'s own internal setup (restoring it before every overlay
	% replay, so overlays within one frame don't double-count) -- resetting
	% after \begin{frame} would be reset too late, after that baseline was
	% already captured.
	\setcounter{footnote}{0}%
	\begin{frame}[#1]{#2}%
	% Forces this frame's own body to start with the correct default text
	% color, regardless of what any *previous* frame left active
	% (2026-07-16, direct instruction, root-caused via an isolated minimal
	% reproduction). \titleframe's/\quoteframe's own
	% \color{\gothamTitleTextColor} calls -- set inside content built via a
	% stored box register or a zero-space overlay, used/shipped out *later*
	% rather than immediately -- don't reliably get "undone": xcolor's
	% \color doesn't emit a PDF color operator immediately when called, only
	% when the *next* real glyph/path actually needs it, so whatever color
	% one of those frame types last set can end up as the active PDF color
	% well into the *next* frame's own first draw call. Invisible for Graz
	% (the leaked color, c_primary_dark, happens to match ordinary body text
	% anyway) but catastrophic for MCI/Innsbruck (c_primary_white, rendering
	% plain non-itemize body text invisible against its own white
	% background -- itemize items escape it, since beamer's own itemize
	% template sets color fresh per bullet). An \AtBeginEnvironment{frame}
	% hook (etoolbox) was tried first and confirmed *not* to fire at all --
	% beamer's own "frame" isn't a plain \newenvironment etoolbox can patch
	% that way -- so this resets directly in standardframe's own opening
	% code instead, exactly where this project already has full control.
	\color{c_primary_dark}%
	% Defensive: the tokens after \begin{frame} are collected into the frame
	% body and replayed once per overlay slide, so this clears the side-
	% footnote collector at the start of every slide's pass -- normally a
	% no-op (the previous slide's footline already consumed the box), but it
	% stops stale footnotes leaking in if a previous frame ever collected
	% without shipping a footline.
	\gothamClearSideFootnotesBox%
	\setbeamertemplate{itemize item}{\gothamBulletDot}%
	\setbeamertemplate{itemize subitem}{\gothamBulletDot}%
	\setbeamertemplate{itemize subsubitem}{\gothamBulletDot}%
	\setbeamertemplate{itemize subsubsubitem}{\gothamBulletDot}% 4th level, see its own comment further below (search "subsubsub")
	% Only for top-aligned frames (the default, #1=t): a [c]-centered frame
	% (2026-07-22, direct instruction, found while centering
	% slides/bw9-Pearson-Correlation-Illustration-current.tex/agc/crj) was
	% landing its content almost flush against the title bar instead of
	% truly centered -- root cause was this same pull-up firing
	% unconditionally regardless of #1. beamer's own [c] option already
	% centers the frame body within the available space; stacking this
	% deck's *own* top-alignment pull-up on top of that just drags the
	% already-centered box further upward by the same fixed amount, which
	% is exactly the asymmetry that was measured (a [c] figureframe's image
	% landing ~0pt from the title bar and ~10pt from the footer instead of
	% symmetric). Guarded with \ifstrequal (etoolbox, already loaded above)
	% rather than \ifx, matching this file's own \ifstrequal{##1}{...}
	% pattern elsewhere (\ifx would need #1 fully expanded first).
	\ifstrequal{#1}{c}{}{\vspace*{-\dimexpr\gothamContentTopShift-\gothamContentTopShiftReduction\relax}}%
}{%
	\end{frame}%
	\global\def\gothamChapterLogoPath{}% opt-in only for the frame just closed
	\global\def\gothamContentTopShiftReduction{0pt}% opt-in only for the frame just closed
	\global\gothamSideFootnotesfalse% opt-in only for the frame just closed
	\global\def\gothamFrametitleFullHeightYellow{}% opt-in only for the frame just closed
	\global\def\gothamFrametitleLogoFixed{}% opt-in only for the frame just closed
	\global\def\gothamFooterBoxColor{\gothamFooterBoxColorDefault}% opt-in only for the frame just closed
}

% The bullet itself: a small filled circle, matching the source PowerPoint
% deck's own round bullets (a plain \textbullet, tried first, is too
% small/thin in Fira Sans at any standard text size to read as a solid dot
% -- a tikz-drawn circle gives full control over the actual rendered size
% instead of depending on the font's own glyph design). baseline=-0.65ex
% roughly centers it on the surrounding text's x-height, the same vertical
% position a font's own bullet glyph would normally sit at.
%
% Its fill color is \gothamBulletColor, not a literal color token here --
% set by the per-university preamble before \input{}ing this file, same
% "value supplied externally, referenced lazily" idiom as
% \gothamInstituteLogoPath/\gothamLogoHeight above (\gothamBulletDot's own
% body only expands \gothamBulletColor when a bullet actually gets drawn,
% long after this \newcommand runs). Graz's own bullets are its primary
% highlight color (matching this project's original, single-institution
% behavior); MCI's and Innsbruck's are their own secondary blue instead
% (2026-07-16, direct instruction) -- see each institution's own
% Slidemaster_preamble_*.tex for its \gothamBulletColor setting.
\newcommand{\gothamBulletDot}{\tikz[baseline=-0.65ex]{\fill[\gothamBulletColor] circle (1.7pt);}}

% "a)"/"b)"/"c)"-labeled list, for exercise sub-tasks that read better
% lettered than bulleted (2026-07-20, direct instruction, first used on
% slides/y3a-Exercise-Diversification-2-current.tex's own bottom three
% items).
%
% First attempt (superseded, kept as a documented wrong turn -- see
% README.md's own "lualatex vs latex" section for why this project keeps
% wrong write-ups rather than deleting them): reuse itemize's own top-level
% `\llap`-based label positioning by locally overriding
% `\setbeamertemplate{itemize item}` inside a plain `\begin{itemize}` --
% the reasoning being that `\llap` draws its content flush against the
% text's own left edge "regardless of the label's width." That reasoning
% turned out to be wrong in a way that only showed up once actually
% measured against the compiled PDF (not just visually eyeballed): a plain
% bullet's own `\llap`-drawn footprint (title/Left-Margin-relative-left-
% edge to text-start) measured only ~9.05pt end-to-end on this deck's real
% slides, matching the Left-Margin-to-text-start gap exactly -- but "a)"
% rendered through the *exact same* mechanism measured a ~13.97pt
% footprint for the same gap, not the ~8.5-9pt `\settowidth` reported for
% the bare glyphs in isolation. The extra ~5pt was never fully root-caused
% (this custom itemize's own `\list` call, further below, only overrides
% `\makelabel` -- it never sets `\leftmargin`/`\labelwidth` itself, so
% those are left at whatever this deck's global itemize defaults happen to
% be, and `\llap`'s actual box-width behavior inside that inherited,
% unexamined list-parameter state is where the extra width most likely
% comes from) -- rather than keep guessing at *why*, measured the actual
% failure directly (the compiled "a)" sat ~4.9pt left of the Left Margin,
% confirmed via exact PDF text-position extraction, not a visual guess)
% and abandoned the approach once even `\tiny` plus manual per-letter
% kerning only just barely closed the gap while reading as illegibly small
% -- a real dead end, not a tuning problem.
%
% Working fix: don't reuse itemize's `\llap` label mechanism at all -- use
% a real, width-reserving `\makebox[W][r]{...}` label instead (the same
% *kind* of mechanism `\gothamEnumerateLabelWidth` above already uses for
% a nested enumerate, just with an explicit `\leftmargin`/`\labelwidth`/
% `\labelsep` this environment sets itself rather than inheriting), inside
% its own plain `\list`, not a reused `\itemize`. `\leftmargin` (label box
% + gap) is set to 11.1pt -- 9.1pt label box (fits "a)"/"b)"/"c)" at this
% deck's normal `\small` itemize size with a little clearance, confirmed
% directly: measured label-left-edge stays right of the 12.552pt Left
% Margin for all three letters, closest case (widest letter, "c)") 1.37pt
% clear) plus 2pt `\labelsep` for a visible gap before the text. Trade-off,
% accepted deliberately rather than found by accident: item text lands
% ~2.2pt right of where a plain bulleted item's text would (23.8pt vs.
% 21.6pt) -- fitting a legible "a)"/"b)"/"c)" at normal reading size plus
% a real gap simply doesn't fit in the ~9pt a plain bullet's tiny dot
% leaves, so *exact* text-position parity with plain bullets and a legible
% label are mutually exclusive here; legibility won given the two-point
% shift is visually negligible and there's no plain-bulleted item left on
% the same slide to compare against directly anyway (y3a's own first item
% moved into its infoblock). Re-derive `\leftmargin`/`\labelwidth` the same
% way (render, measure with exact text-position extraction, adjust) if
% this is ever used at a different font size than this deck's normal
% itemize `\small`.
\newcounter{gothamalphaitem}
\newcommand{\gothamAlphaItem}{\stepcounter{gothamalphaitem}\alph{gothamalphaitem})}
\newenvironment{gothamalphaitemize}{%
	\setcounter{gothamalphaitem}{0}%
	\list{\gothamAlphaItem}{%
		% 4.75pt (2026-07-20, direct instruction: "half the width of the
		% 'C' in 'Calculate'" -- \settowidth{C} at this deck's normal
		% itemize \small measured 5.5pt, so 2pt original + 2.75pt half-C).
		% Only \labelsep changes the text's own start position -- the
		% label's own box stays exactly where it was (its left edge is
		% pinned to this list's origin regardless of \labelsep, since
		% \leftmargin=\labelwidth+\labelsep always cancels \labelsep back
		% out of the label box's own position), so this widens the gap
		% without reopening the Left-Margin question already solved above.
		\setlength{\labelsep}{4.75pt}%
		\setlength{\labelwidth}{9.1pt}%
		\setlength{\leftmargin}{\labelwidth}%
		\addtolength{\leftmargin}{\labelsep}%
		\setlength{\itemindent}{0pt}%
		% \itemsep/\parsep left at plain \list's own defaults otherwise --
		% since this environment is built directly on \list rather than
		% reusing beamer's own itemize template, it never got beamer's own
		% itemize spacing hooks at all, only whatever \list falls back to.
		% Measured directly (2026-07-21, direct instruction, found on
		% slides/nt4-Exercise-CAPM-5-current.tex): that default reads as
		% *less* space between two items (13.7pt, item-to-item) than
		% between two wrapped lines within one *single* item (15.6pt,
		% pure \baselineskip) -- backwards from how a list should read
		% (item boundaries should be at least as open as a same-item line
		% wrap, not tighter). \itemsep set explicitly so the total gap at
		% an item boundary (baselineskip + itemsep) reads as clearly more
		% open than a plain wrapped line, not accidentally tighter.
		% Reduced from the original 6pt to 2.5pt (2026-07-22, found during
		% a full-project check after the itemize-rules work, not part of
		% that work itself): 6pt overshot by more than needed to clear the
		% 1.9pt gap (15.6 vs. 13.7), and slides/kmf-Exercise-CAPM-6-current.tex's
		% own 5-item list -- longer than nt4's own -- compounded that
		% excess across 4 item boundaries into a ~25pt overfull frame, its
		% own last item cut off by the footer (confirmed directly,
		% rendered and compared against a pre-regression PDF). 2.5pt
		% keeps a real, visible margin over the 1.9pt needed (not just
		% barely crossing it) while re-measuring clean (no overfull) on
		% kmf; re-verify against both nt4 and kmf directly if this is
		% ever changed again -- they're this environment's own tightest
		% (most items) and originally-tuned-against cases respectively.
		\setlength{\itemsep}{2.5pt}%
	}
	\raggedright% plain \list defaults to justified+hyphenated text, unlike
	% this deck's own itemize (which calls this too, further below) --
	% missed on the first pass, only caught by comparing a rendered a)/b)/c)
	% item side by side against a plain bulleted one ("stan-\ndard deviation"
	% hyphenating where a plain-bulleted version of the same sentence
	% elsewhere in this deck doesn't).
}{%
	\endlist
}

% Itemize indentation: level 1 bullets should sit flush with the frame title
% (no extra left margin), and each deeper level should be indented just
% enough for its bullet to line up under the *text* (not the bullet) of the
% level above. A nested list's \leftmarginN is added on top of its parent's
% (indentation is cumulative), so giving every level the same "step" --
% label width + labelsep, the width a bullet plus its gap to text takes up --
% makes level N's text sit at N*step and level N's bullet at (N-1)*step, i.e.
% exactly at level (N-1)'s text (0 for level 1, flush with the title).
\makeatletter
\newlength{\gothamBulletWidth}
\settowidth{\gothamBulletWidth}{\gothamBulletDot}
\newlength{\gothamItemStep}
\setlength{\gothamItemStep}{\labelsep}
\addtolength{\gothamItemStep}{\gothamBulletWidth}
\setlength{\leftmargini}{\gothamItemStep}
\setlength{\leftmarginii}{\gothamItemStep}
\setlength{\leftmarginiii}{\gothamItemStep}
\setlength{\leftmarginiv}{\gothamItemStep}
\makeatother

% A 4th itemize level -- beamer's own \itemize (beamerbaselocalstructure.sty)
% hard-caps nesting at 3 levels (\ifnum \@itemdepth >2\relax\@toodeep),
% independently of anything this deck sets: a 4th \begin{itemize} always
% raised "! LaTeX Error: Too deeply nested." even though plain LaTeX's own
% kernel itemize supports 4. Found 2026-07-15, direct instruction, on
% slides/hrg-Exchange-Rate-Quotation-Forward-current.tex -- genuinely
% 4 levels deep in the source pptx itself (confirmed via a PowerPoint COM
% render, not assumed), not a conversion artifact to design around.
% Since this deck already styles every existing level identically (same
% \gothamBulletDot, same \gothamItemStep margin -- see the comments right
% above/below), extending to a 4th level needs no new visual design, just
% patching beamer's own depth limit up by one and telling it what to call
% the new level ("subsubsub", the next name in beamer's own
% sub/subsub/... naming scheme) so its template/color lookups
% (\usebeamertemplate{itemize subsubsubitem}, \usebeamercolor{itemize/
% enumerate subsubsub body}, etc.) resolve to *something* instead of
% erroring -- undefined beamer templates/colors are expected to fall back
% to sensible defaults (no extra spacing hook, plain text color) rather
% than raising an error, same as any other beamer template/color that a
% given theme never bothers to set; the actual bullet symbol still needs
% its own \gothamBulletDot template below, added explicitly the same way
% as the existing 3 levels, not left to whatever fallback beamer has for
% an unset "itemize subsubsubitem" template.
% \beamer@computepref is patched wholesale (not via \patchcmd) since its
% body is one single \ifcase chain, not a stable string to match against.
% \itemize itself was *first* tried via \patchcmd{\itemize}{...}
% (matching this deck's own established practice for beamer/theme
% internals, e.g. beamerfontthemegotham.sty's own
% \patchcmd{\beamer@title}{...}) -- this failed outright, confirmed
% directly (\patchcmd's own failure branch fired, raising a clean
% "Patching failed" error rather than silently doing nothing): recent
% LaTeX2e kernels implement an optional-argument environment like
% beamer's own `\renewcommand{\itemize}[1][]{...}` via
% `\@protected@testopt`, which stores the *real* body not under `\itemize`
% itself but under a differently-named internal control sequence (`\show
% \itemize` confirms this directly: it prints as
% `\@protected@testopt \itemize \\itemize {}`, i.e. plain `\itemize` is
% only a thin dispatcher, and the actual `\ifnum \@itemdepth >2...` body
% \patchcmd was trying to match against was never there to find). Rather
% than chase that internal storage name (liable to differ across LaTeX
% kernel versions, which \patchcmd's own documented usage doesn't need to
% care about), \itemize is fully re-declared here instead, copied
% verbatim from beamer's own beamerbaselocalstructure.sty (this project's
% installed version) with only the depth bound itself changed from 2 to
% 3 -- the same declaration form beamer itself used
% (\renewcommand{\itemize}[1][]{...}), which regenerates both the
% dispatcher and the real body correctly under whatever internal name
% the current kernel expects, sidestepping the naming question entirely.
% Trade-off, stated plainly: if a future beamer release changes this
% macro's own body for an unrelated reason, this override will silently
% keep using today's (2026-07) copy instead of picking up that change --
% acceptable here since the override is one line different from upstream
% and clearly marked, but re-diff against beamerbaselocalstructure.sty's
% own \itemize if beamer is ever upgraded and itemize behavior looks off.
\makeatletter
\long\def\beamer@computepref#1{%
	\let\beameritemnestingprefix\@empty%
	\ifcase#1\or\or\def\beameritemnestingprefix{sub}\or\def\beameritemnestingprefix{subsub}\or\def\beameritemnestingprefix{subsubsub}\or\@toodeep\fi%
}
\renewcommand{\itemize}[1][]{%
	\ifblank{#1}{}{\def\beamer@defaultospec{#1}}%
	\ifnum \@itemdepth >3\relax\@toodeep\else
		\advance\@itemdepth\@ne
		\beamer@computepref\@itemdepth% sets \beameritemnestingprefix
		\usebeamerfont{itemize/enumerate \beameritemnestingprefix body}%
		\usebeamercolor[fg]{itemize/enumerate \beameritemnestingprefix body}%
		\usebeamertemplate{itemize/enumerate \beameritemnestingprefix body begin}%
		\list
			{\usebeamertemplate{itemize \beameritemnestingprefix item}}
			{\def\makelabel##1{%
					{%
						\hss\llap{{%
								\usebeamerfont*{itemize \beameritemnestingprefix item}%
								\usebeamercolor[fg]{itemize \beameritemnestingprefix item}##1}}%
					}%
				}%
			}
	\fi%
	\beamer@cramped%
	% \beamer@cramped (beamerbaselocalstructure.sty, unmodified upstream
	% code) unconditionally does `\itemsep0pt\topsep0.5pt`, every single
	% time -- called right here, in this project's own copy of beamer's
	% own \itemize body, same position as upstream. Setting \itemsep any
	% earlier (first tried as a separate \AtBeginEnvironment{itemize} hook,
	% run *before* this whole macro body even starts) gets silently wiped
	% out by this same line a moment later -- confirmed directly via
	% \typeout debugging (not assumed): the hook's own \setlength really
	% did run with the right values, \itemsep still read back as 0pt
	% immediately after. Setting it *here*, right after \beamer@cramped
	% already ran (not before), is what actually survives.
	\ifcase\@itemdepth
	\or\setlength{\itemsep}{\gothamItemsepRatioI em}% level 1
	\or\setlength{\itemsep}{\gothamItemsepRatioII em}% level 2
	\or\setlength{\itemsep}{\gothamItemsepRatioIII em}% level 3
	\or\setlength{\itemsep}{\gothamItemsepRatioIV em}% level 4
	\fi%
	\raggedright%
	\beamer@firstlineitemizeunskip%
}
\makeatother

% A nested \begin{enumerate} inherits the same \leftmarginN step as itemize
% above for *where its text starts* (matching itemize's own "one step past
% the parent level's text" rule -- confirmed directly against the compiled
% PDF: an enumerate's own body text, e.g. "Senior secured debt", already
% lands exactly one \gothamItemStep past its parent item's text), but its
% *label* ("1.", "2.", ...) doesn't follow \labelwidth the way plain LaTeX's
% list mechanism would -- beamer draws enumerate labels through its own
% template machinery, which turned out not to respect \labelwidth at all:
% \setlength{\labelwidth}{...} right after \begin{enumerate} measurably took
% effect (checked with \typeout{\the\labelwidth} immediately after) but the
% compiled numeral's own position never moved. What *does* move it: directly
% overriding \makelabel with a fixed-width \makebox -- empirically
% calibrated (render, measure the numeral's own x-position in the compiled
% PDF, adjust, repeat) since the box's own left edge turned out to sit at a
% fixed position unrelated to \labelwidth or \leftmarginN, with the
% right-aligned numeral inside it moving in lockstep with the box's width
% (confirmed linear: +1pt width -> +~1pt numeral position, across several
% test widths). 19.23pt (at this deck's default itemize font size) lands a
% single-digit numeral flush with the parent item's own text -- only valid
% at that font size and for single-digit lists (1-9); re-calibrate the same
% way (not by recomputing from first principles) for a differently-sized or
% 10+-item nested enumerate. Call right after \begin{enumerate}, before the
% first \item.
%
% 20.33pt (2026-07-08, direct instruction: "the numbers still extend a
% little farther to the left than the text above them does" --
% slides/f6m-Recovery-Rates-current.tex's own "Pecking order:" list).
% 19.23pt landed *close* but not quite flush: measured directly against
% the compiled PDF, every numeral's own left edge (20.55pt-20.95pt,
% varying slightly by digit width even though all are right-aligned to
% the same box) still sat up to ~1.05pt left of "Pecking order:"'s own
% text start (21.6pt) -- the widest digit tried, "6.", was the binding
% case. +1.1pt of box width (confirmed +1pt width -> +1pt numeral
% position, per this comment's own note above) covers that with a small
% margin, re-confirmed against the recompiled PDF.
\newcommand{\gothamEnumerateLabelWidth}{\renewcommand{\makelabel}[1]{\makebox[20.33pt][r]{##1}}}

% Automatic version of the above (2026-07-23, background-task follow-up
% flagged in slides/kfh-Asset-Allocation-Process-current.tex's own v1-1-0
% comment: "fix the enumerate behavior generally, if possible", task
% `task_bb3fbb66`). That comment's first attempt, an
% `\AtBeginEnvironment{enumerate}{...}` hook, failed because the hook fires
% *before* `\enumerate`'s own body runs, and that body's own internal
% `\list{...}{\def\makelabel...}` setup (the beamer copy directly below the
% `\itemize` override above) clobbers whatever a hook that early leaves
% behind -- the exact same "template runs its own setting, external pre-set
% gets overwritten" problem `\itemize`'s own `\itemsep` comment above
% already documents. The real fix, same idea `\itemize` itself used above:
% override beamer's own low-level piece directly rather than hooking
% \AtBeginEnvironment. Unlike `\itemize` (declared via
% `\@protected@testopt`, needing a full re-declaration to reach -- see the
% `\itemize` override's own comment for the full story), `\enumerate`'s
% relevant piece, `\beamer@enum@` (beamerbaselocalstructure.sty), is a
% plain `\def` with no optional-argument indirection, so it can be
% `\renewcommand`'d directly, verbatim from upstream except for the
% `\ifnum\@itemdepth>1` branch added below.
%
% Only fires for a *nested* enumerate (`\@itemdepth>1`, i.e. already inside
% at least one `\itemize`/`\enumerate` when this one opens) -- a top-level
% enumerate (`\@itemdepth=1`, e.g.
% slides/t9z-Outland-Steel-Hedging-Borrowing-Investing-current.tex/
% slides/g4c-Outland-Steel-Hedging-Forward-Sale-current.tex/
% slides/j84-Outland-Steel-No-Hedging-current.tex's own outermost numbered
% list, or slides/nkj-Agenda-current.tex's own \tableofcontents-driven one)
% is left at beamer's stock label position entirely untouched, matching
% this project's own established behavior for that shape (never reported
% as misaligned, unlike the nested case) -- and, for nkj specifically,
% doesn't matter either way, since that slide's own
% `\gothamTocLabelFlushLeft` call (right after its own `\begin{enumerate}`)
% overrides `\makelabel` again regardless of whatever this default sets.
%
% Reuses the exact same 20.33pt calibration as the manual
% `\gothamEnumerateLabelWidth` above, confirmed (pymupdf, compiled PDF) to
% land correctly not just at `\@itemdepth=2` (one `\itemize` above the
% enumerate -- slides/kfh-Asset-Allocation-Process-current.tex/
% slides/f6m-Recovery-Rates-current.tex/
% slides/6pz-Outland-Steel-Hedging-Strategies-Overview-current.tex, all
% previously needing the manual call) but also at `\@itemdepth=3` (two
% `\itemize`s above it -- slides/bdh-Working-Capital-Example-current.tex's
% own "Development of balance sheet:" list, also previously needing the
% manual call, with the identical 20.33pt value) -- this deck's itemize/
% enumerate templates don't shrink font size by nesting depth (only the
% bullet glyph/margin step does), so the numeral's own glyph width -- the
% only thing \gothamEnumerateLabelWidth's own box width actually
% compensates for -- stays the same at every depth, and one constant covers
% both cases without needing a depth-dependent formula. Not verified beyond
% depth 3 (no slide currently nests an enumerate that deep); if one ever
% does, re-measure before trusting this constant that far down, same as any
% other empirical calibration in this file.
% (slides/scs-Political-Risk-Stability-Index-current.tex also matches
% "`\gothamEnumerateLabelWidth`" in a grep -- but only inside its own prose
% comments describing the general technique; its actual call is the
% bespoke `\gothamEnumerateLabelFlushLeftShifted{16pt}{10.2pt}` for its own
% alphabetic, left-aligned, narrower-box labels, unaffected by this
% override either way since that call runs after `\begin{enumerate}` opens,
% same ordering as every other post-open `\makelabel` override here.)
%
% Former per-slide manual `\gothamEnumerateLabelWidth` calls removed from
% the 4 slides above now that this covers them automatically; the macro
% itself is kept defined, both for this override to reuse and as a manual
% escape hatch (e.g. a differently-calibrated width some future slide might
% need).
\makeatletter
\renewcommand{\beamer@enum@}{%
	\beamer@computepref\@itemdepth% sets \beameritemnestingprefix
	\usebeamerfont{itemize/enumerate \beameritemnestingprefix body}%
	\usebeamercolor[fg]{itemize/enumerate \beameritemnestingprefix body}%
	\usebeamertemplate{itemize/enumerate \beameritemnestingprefix body begin}%
	\expandafter
		\list
			{\usebeamertemplate{\beamer@enumtempl}}
			{\usecounter\@enumctr%
				\ifnum\@itemdepth>1
					\def\makelabel##1{\makebox[20.33pt][r]{##1}}%
				\else
					\def\makelabel##1{{\hss\llap{{%
							\usebeamerfont*{enumerate \beameritemnestingprefix item}%
							\usebeamercolor[fg]{enumerate \beameritemnestingprefix item}##1}}}}%
				\fi}%
	\beamer@cramped%
	\raggedright%
	\beamer@firstlineitemizeunskip%
}
\makeatother

% Extra breathing room when a nested itemize closes and control returns to
% its *parent* list, right before that parent's next \item -- e.g. sub-items
% under "Mid-Term Exam" followed by the next top-level "Final Exam" -- so a
% new top-level group reads as visually distinct from the one before it,
% rather than the same uniform gap used between every other item regardless
% of level. \@itemdepth (LaTeX's own nesting-level counter, already
% decremented back to the parent's level by the time \end{itemize}'s normal
% closing has run) says which level control returned to -- 0 means the
% outermost list itself just closed (nothing left to add space before,
% deliberately excluded), 1 means back to a level-1 item, 2 means back to a
% level-2 item.
%
% Level-dependent (2026-07-06, direct instruction) -- a flat, single
% constant regardless of *which* level control returned to (the original
% version of this hook) reads as "extra"/inconsistent precisely because a
% return-to-level-1 and a return-to-level-2 transition both got the exact
% same gap, rather than level 1 getting a visibly bigger gap than level 2 --
% found on slides/57r-Reasons-for-Leasing-current.tex, where a level-2
% sub-list closing back to a level-1 item looked like a random extra gap
% rather than "level 1 items are consistently more separated than level 2
% ones." \addvspace (not plain \vspace) matters too: it coalesces with
% whatever glue itemize's own \itemsep/\topsep already contributes at the
% closing boundary (taking whichever is larger) rather than *adding* on top
% of it unconditionally -- \vspace's own stacking behavior was the other
% half of why the old flat version could look like too much space.
\makeatletter
% Ratios, not lengths (2026-07-21, direct instruction: "all of these
% spacings should scale with font size") -- a \newlength+\setlength{...}{0.5em}
% pair, the way this used to be written, freezes to a fixed *pt* value the
% one time \setlength runs (at preamble-load time, using whatever font is
% active then), despite the "em" in its own definition -- it does NOT
% re-track whichever font size is actually active on a given slide/level
% later. Confirmed directly: a \typeout{\f@size} placed at the point each
% of these actually gets used (inside \AtEndEnvironment{itemize}'s own
% hook body, further below) showed the *closing* list's own font still
% active there (9pt/10pt/10.95pt for a level-4/3/2 list closing,
% respectively -- not yet restored to the parent's own font), confirming
% this hook fires early enough in `\end{itemize}`'s own processing that a
% live "0.5em" computed *there* would correctly track whichever level is
% actually involved. Fixed by storing the bare ratio as a plain macro
% instead, and writing the literal "em" back on at the point of use
% (`\addvspace{\gothamItemizeLevelDropSpaceRatioI em}`), so it's
% recomputed fresh, against whatever font is live, every time.
\newcommand{\gothamItemizeLevelDropSpaceRatioI}{0.5}
\newcommand{\gothamItemizeLevelDropSpaceRatioII}{0.3}
% Level 3's own drop-space, for when a closed *4th*-level list returns
% control to a level-3 item -- added 2026-07-15 alongside the 4th itemize
% level itself (search "subsubsub" further above). Interpolated from the
% existing I/II progression (0.5, 0.3 -- each step roughly 60% of the one
% before) rather than empirically re-measured against a compiled PDF the
% way I/II themselves were: no slide used a 4th level before
% slides/hrg-Exchange-Rate-Quotation-Forward-current.tex to measure
% against. Re-derive properly (render, measure, adjust) if this ever
% looks visibly off on a real slide.
\newcommand{\gothamItemizeLevelDropSpaceRatioIII}{0.2}
% Level 4's own drop-space, for when a closed level-4 list returns control
% to a level-3 item -- added 2026-07-21 alongside \gothamItemsepRatioIV.
% \@itemdepth at this hook's fire time is the depth of the list that *just
% closed*, confirmed directly via a temporary \typeout{\the\@itemdepth}
% (not assumed): closing the level-4 list on
% slides/ZZO-Itemize-Spacing-Reference-current.tex measured itemdepth=4
% there, closing level-3 measured 3, level-2 measured 2, and the outermost
% level-1 list itself closing measured 1 (never 0 -- \@itemdepth's own
% restore to the parent's depth happens *after* this hook runs, so the "0:
% outermost list itself closed" case below is dead code, kept only because
% removing a genuinely unreachable \ifcase branch isn't worth the risk of
% getting the numbering wrong a second time). So this is very much a real,
% reachable case -- not decorative symmetry -- and needed its own
% \addvspace branch (previously missing entirely, silently leaving the
% level-4-to-level-3 closing transition with no drop-space at all).
% Interpolated from the existing I/II/III progression as a starting guess,
% then re-measured directly against the compiled ZZO reference and bumped:
% 0.13 measured at 12.6pt, still *less* than level 4's own same-level gap
% of 13.4pt (violates rule III -- a level's own same-level gap must be
% smaller than that level's own closing-transition gap); raised until
% re-measurement cleared it with a comfortable margin again, matching the
% margin every other level's own pair already has.
\newcommand{\gothamItemizeLevelDropSpaceRatioIV}{0.5}
\AtEndEnvironment{itemize}{%
	\ifcase\@itemdepth
		% 0: dead code -- @itemdepth is never actually 0 when this hook
		% fires (see \gothamItemizeLevelDropSpaceRatioIV's own comment
		% just above); kept only for alignment with the \ifcase branches
		% below, each keyed on the depth of the list that just closed.
	\or\addvspace{\gothamItemizeLevelDropSpaceRatioI em}% level-1 (outermost) list itself closed -- a pre-existing quirk (not introduced by this comment's own edit): there's no parent item above level 1 for this to visually separate from, so this only actually matters if something follows the outermost list directly on the same frame; harmless on every slide seen so far since nothing does. Out of scope to redesign here.
	\or\addvspace{\gothamItemizeLevelDropSpaceRatioII em}% level-2 list closed, back to a level-1 item
	\or\addvspace{\gothamItemizeLevelDropSpaceRatioIII em}% level-3 list closed, back to a level-2 item
	\or\addvspace{\gothamItemizeLevelDropSpaceRatioIV em}% level-4 list closed, back to a level-3 item
	\fi
}

% Per-level \itemsep -- the item-to-item gap *within* one list, at one
% fixed nesting depth (as opposed to the drop-space mechanism just above,
% which only fires when a nested list *closes* back up to its parent).
% Ratios, not lengths -- same "scale with font size" reasoning and fix as
% \gothamItemizeLevelDropSpaceRatioI/II/III above (that comment has the
% full story on why a plain \newlength+\setlength pair here would freeze
% to one fixed pt value instead of actually tracking each level's own
% font size). Values only, defined here for discoverability alongside the
% drop-space ones above -- actually *applied* inside the `\itemize`
% redefinition itself (search "\beamer@cramped (beamerbaselocalstructure.sty"
% further above), not via a separate hook here; see that comment for why.
% Found missing entirely (2026-07-21, direct instruction, see
% slides/ZZO-Itemize-Spacing-Reference-current.tex's own header comment
% for the rules this and the drop-space mechanism together need to
% satisfy): this deck's own itemize never set \itemsep itself at any
% level -- confirmed directly, grepping this whole file -- so every depth
% shared whatever bare default beamer's own theme happens to leave
% \itemsep at, uniformly. That produced two visible problems, both
% measured directly (pymupdf against slides/2c2-Topics-Around-the-CAPM-Empirical-Validity-current.tex
% and an isolated bisection test, not guessed): (a) level-2 and level-3
% same-level item-to-item gaps came out *identical* (12.6pt each,
% confirmed against an isolated reproduction), since nothing distinguished
% them by depth at all -- only a closing transition ever got a
% level-dependent gap, never a same-level one; (b) that shared default
% \itemsep read as *smaller* than a single item's own wrapped-line
% \baselineskip (13.7-13.8pt within a wrap vs. 12.6pt between items),
% backwards from how a list should read -- an item boundary should never
% be tighter than a same-item line wrap.
% Ratios themselves are not arbitrary -- each chosen, then re-measured
% against the compiled ZZO reference slide and adjusted, specifically to
% satisfy all three of that slide's own header-comment rules together
% (extended to level 4 the same way, 2026-07-21), not just "level 1
% biggest" in isolation:
% - Rule I (level-1 same-level gap > level-2's > level-3's > level-4's)
%   falls straight out of \gothamItemsepRatioI > ...II > ...III > ...IV,
%   added on top of the same \baselineskip every level shares (itself
%   already smaller per level here, since this theme's own itemize/
%   enumerate body font shrinks with depth too -- confirmed directly,
%   10.95pt/10pt/9pt/9pt for levels 1-4 respectively, level 4 sharing
%   level 3's own font since this theme has no distinct "subsubsub" size
%   of its own).
% - Rule II (every level's own same-level gap > that level's own wrapped-
%   line gap) needs \gothamItemsepRatioIV alone (the smallest, hence the
%   binding case) to still be positive and large enough that it clearly
%   reads as more open than a plain wrapped line once rendered, not just
%   bigger on paper.
% - Rule III (a level's own same-level gap < that same level's own
%   *closing*-transition gap) is *not* automatic just because
%   \gothamItemizeLevelDropSpaceRatio* was already bigger than the old
%   shared default -- \addvspace's own coalesce-to-the-larger-value
%   behavior (documented above) means the closing gap is
%   max(natural glue, that *drop-space* ratio), not a sum of the two. Each
%   drop-space ratio's own \ifcase branch is keyed on the depth of the
%   list that *just closed* (see \gothamItemizeLevelDropSpaceRatioIV's own
%   comment for the full story on why, confirmed via \typeout, not
%   assumed), so it lines up with the same-numbered itemsep ratio despite
%   describing a different transition -- \gothamItemizeLevelDropSpaceRatioIV
%   governs the level-4-list-closing-back-to-level-3 transition and is the
%   one that has to individually stay bigger than \gothamItemsepRatioIV.
%   This branch was missing entirely at first (only I/II/III existed, a
%   leftover from before level 4 existed) -- confirmed directly on ZZO: the
%   level-4-to-level-3 closing transition measured 12.6pt with no drop-space
%   branch at all, less than level 4's own same-level gap of 13.4pt, a real
%   rule III violation. Adding the branch (\gothamItemizeLevelDropSpaceRatioIV,
%   starting at 0.5) fixed it. Re-verify rule III directly against ZZO's own
%   compiled PDF (don't assume it still holds on paper) if any of these ratios is
%   ever changed again.
\newcommand{\gothamItemsepRatioI}{0.4}
\newcommand{\gothamItemsepRatioII}{0.24}
\newcommand{\gothamItemsepRatioIII}{0.15}
\newcommand{\gothamItemsepRatioIV}{0.09}

% Per-level "opening" space -- the gap between a level-n item and the
% first item of a level-(n+1) list that immediately follows it, as
% opposed to \gothamItemizeLevelDropSpaceRatio* (a nested list *closing*
% back up) or \gothamItemsepRatio* (two items at the *same* depth).
% Added 2026-07-21, direct instruction, replacing the original three
% separate rules on slides/ZZO-Itemize-Spacing-Reference-current.tex's
% own header comment with one multi-part inequality stated per level n:
%   level-n -> level-(n-1) close  >  level-n same-level  >
%   level-n -> level-(n+1) open  >  level-n's own wrap gap
% The "open" term was entirely unhandled before this -- confirmed
% directly on ZZO: every opening transition (level 1 into level 2, level
% 2 into level 3, level 3 into level 4) measured *smaller* than that same
% level's own wrap gap, backwards from what the rule needs, since nothing
% but beamer's own bare default \topsep governed it.
% \AtBeginEnvironment{itemize}, not \topsep -- \topsep was tried first
% (set via the same \ifcase\@itemdepth block, right after \itemsep,
% inside the `\itemize` redefinition above) and measured *unreliable*:
% two of the three levels barely moved despite a sizeable ratio, while
% the third moved by roughly the right amount -- \topsep's own value is
% evidently consulted by `\list` at some point this override comes too
% late for, at least for some levels, not evenly. \addvspace via
% \AtBeginEnvironment instead measured consistently and predictably
% (each level's own change matched its own ratio's calculated pt value
% directly), the same mechanism \gothamItemizeLevelDropSpaceRatio* already
% relies on for the closing side. This hook fires with \@itemdepth
% reflecting the *parent* item's own depth (confirmed via \typeout, not
% assumed: opening the outermost level-1 list itself measured 0, opening
% a level-2 list from within a level-1 item measured 1, level-3-from-
% level-2 measured 2, level-4-from-level-3 measured 3) -- the opposite of
% \gothamItemizeLevelDropSpaceRatio*'s own hook, which keys on the depth
% of the list that *just closed*. That makes the \ifcase numbering here
% match the "n" in the rule directly: case 1 is level-1's own open-to-
% level-2 ratio, case 2 is level-2's own open-to-level-3, case 3 is
% level-3's own open-to-level-4; case 0 (opening the outermost list
% itself) is a genuine no-op, correctly -- there's no parent item above
% level 1 for that transition to separate from.
% Ratios chosen, then re-measured against the compiled ZZO reference and
% adjusted, the same way as every other ratio here: each has to
% individually land strictly between that level's own \gothamItemsepRatio*
% (the ceiling) and its own wrap gap (the floor). Re-verify directly
% against ZZO's own compiled PDF (don't assume it still holds on paper)
% if any of these, or the itemsep/drop-space ratios they're measured
% against, are ever changed again.
% \gothamItemizeLevelOpenSpaceRatioI trimmed from 0.28 to 0.02 (2026-07-21,
% direct instruction: "the level-1 to level-2 distance...seems a bit
% much"). Confirmed directly: even at 0, this transition still measured
% ~16.3pt -- \addvspace's own coalesce-to-the-larger-value behavior means
% the ratio can only ever push this *above* whatever glue the list
% mechanism inserts here regardless (that natural floor, not the ratio),
% so 0.28 down to 0 only ever moved the actual gap by ~2pt (18.5->16.3);
% 0.02 splits the difference, keeping the mechanism nominally "in effect"
% rather than reading as accidentally disabled. \gothamItemizeLevelOpenSpaceRatioIII
% (level-3-to-level-4) had to come down too, from 0.08 to 0.02, purely to
% restore its own margin -- changing RatioI shifted that unrelated
% transition's own measured gap closer to level 3's own same-level gap
% (12.1 vs. 12.0, too thin), confirmed directly on ZZO, not assumed;
% ratios elsewhere in this file have shown the same small-ripple behavior
% before (see \gothamItemizeLevelDropSpaceRatioIV's own comment) --
% re-measure the *whole* table after touching any single one of these,
% not just the transition being changed.
\newcommand{\gothamItemizeLevelOpenSpaceRatioI}{0.02}
\newcommand{\gothamItemizeLevelOpenSpaceRatioII}{0.15}
\newcommand{\gothamItemizeLevelOpenSpaceRatioIII}{0.02}
\AtBeginEnvironment{itemize}{%
	\ifcase\@itemdepth
		% 0: opening the outermost (level-1) list itself -- no parent item
		% precedes it, nothing to add.
	\or\addvspace{\gothamItemizeLevelOpenSpaceRatioI em}% opening a level-2 list from a level-1 item
	\or\addvspace{\gothamItemizeLevelOpenSpaceRatioII em}% opening a level-3 list from a level-2 item
	\or\addvspace{\gothamItemizeLevelOpenSpaceRatioIII em}% opening a level-4 list from a level-3 item
	\fi
}
\makeatother

% A \tableofcontents-driven agenda slide (see
% slides/nkj-Agenda-current.tex) needs the exact same top-margin/line-
% spacing as every other list in this deck, not whatever gotham's own
% "section in toc" templates use (both "gotham bullet" and "gotham simple"
% carry their own spacing/rule/font choices, tuned for gotham's generic
% look, not this project's own itemize metrics). Rather than re-tuning a
% *third* set of constants to match, this template just emits a plain
% \item -- wrapping the whole \tableofcontents call in a real
% \begin{enumerate}...\end{enumerate} then makes each section entry a
% genuine enumerate item, reusing this file's own \leftmargini/\itemsep/
% \gothamContentTopShift machinery exactly as any other list would, so it
% automatically stays consistent with the rest of the deck instead of
% drifting out of sync with it later.
\defbeamertemplate{section in toc}{gotham plain enumerate}{\item \inserttocsection}

% Even with the plain-enumerate template above, \tableofcontents still
% doesn't match this deck's own normal itemize spacing -- beamer's own
% \beamer@sectionintoc (beamerbasetoc.sty), which wraps *every* section
% entry regardless of which "section in toc" template is active,
% unconditionally inserts a fixed \vskip1.5em before each entry (an
% unbounded \vfill, even worse, if subsections aren't hidden -- neither is
% reachable via a public \tableofcontents option, just this show/hide
% binary), *and* wraps the entry itself in \hbox{\vbox{...}} -- box
% wrapping that (like \gothamMinipageTopCorrection's own minipage case
% elsewhere in this file) changes how surrounding baselineskip/topskip
% glue gets inserted compared to plain paragraph text. Both together added
% roughly 13pt of extra gap
% between entries beyond this deck's own normal ~18.6pt level-1 item
% spacing (measured directly against
% slides/57r-Reasons-for-Leasing-current.tex's own compiled PDF).
% Neither is controllable via a public \tableofcontents option, so this
% redefines \beamer@sectionintoc itself -- scoped safely since
% \tableofcontents is only ever called on the one Agenda slide in this
% deck -- keeping its numbering/hyperlink/part-number/hide-checking logic
% intact but dropping the vskip/vfill and the hbox/vbox wrapping, so an
% entry behaves exactly like a plain \item once printed.
\makeatletter
\def\beamer@sectionintoc#1#2#3#4#5{%
	\ifnum\c@tocdepth>0%
	\ifnum#4=\beamer@showpartnumber%
	{%
	\beamer@saveanother%
	\gdef\beamer@todo{}%
	\beamer@slideinframe=#1\relax%
	\expandafter\only\beamer@tocsections{\gdef\beamer@todo{%
			\beamer@tempcount=#5\relax%
			\advance\beamer@tempcount by\beamer@sectionadjust%
			\ifnum\beamer@tempcount>0
				\ifnum\beamer@tempcount>\beamer@toclastsection
					\def\inserttocsectionnumber{}%
				\else
					\edef\inserttocsectionnumber{\the\beamer@tempcount}%
				\fi
			\else
				\def\inserttocsectionnumber{}%
			\fi%
			\def\inserttocsection{\hyperlink{Navigation#3}{#2}}%
			\beamer@tocifnothide{\ifnum\c@section=#1\beamer@toc@cs\else\beamer@toc@os\fi}%
			{%
				\beamer@tocact{\ifnum\c@section=#1\beamer@toc@cs\else\beamer@toc@os\fi}{section in toc}%
			}%
		}%
	}%
	\beamer@restoreanother%
	}%
	\beamer@todo%
	\fi\fi%
}
\makeatother

% Left-aligns an enumerate's labels within a caller-chosen \labelwidth --
% used by the \tableofcontents Agenda slide (see
% slides/nkj-Agenda-current.tex's own comment for why its labels need
% this at 3.6pt). Defined here as a named helper rather than written as a
% literal \renewcommand{\makelabel}[1]{...#1...} inside the slide's own
% frame body, because that literal form -- specifically its raw `#1`
% parameter token -- breaks inside any frame that \tableofcontents forces
% through beamer's overlay reprocessing: beamer replays such a frame's
% collected body via \def-based machinery (\loop's own \iterate, and
% \beamer@doifinframe), and a `#` in the material of any \def must be
% doubled -- so the compile hits "! Illegal parameter number in definition
% of \iterate." (twice) plus the same for \beamer@doifinframe, once per
% compile, at that frame's own \end{frame}. TeX's error recovery ("I'll
% assume you meant ##") happens to produce the intended definition anyway,
% which is why the rendered PDF was always correct despite the errors --
% but lualatex still exits nonzero, which is exactly what kept breaking
% latexmk's automated lualatex/biber/rerun sequencing on the main deck
% (the "Known open issue" long tracked in CONVERSION_NOTES.md, originally
% misattributed to a biblatex interaction -- see its "Resolved" entry).
% Hoisting the definition out here means the frame body carries only this
% macro's *call* token, with no `#` in it, which sidesteps the whole
% mechanism. The inner \makelabel parameter is written `##1` below since
% it sits inside another \newcommand's body -- standard nesting, and also,
% incidentally, exactly the doubling rule the inline version fell foul of.
\newcommand{\gothamTocLabelFlushLeft}[1]{%
	\setlength{\labelwidth}{#1}%
	\renewcommand{\makelabel}[1]{\makebox[\labelwidth][l]{##1}}%
}

% Same idea as \gothamTocLabelFlushLeft above, plus an extra fixed
% \hspace shift before the label (for a list that needs to start further
% right than the plain nested-list indent, e.g. to line up with a
% sibling item's own text rather than its bullet -- see
% slides/scs-Political-Risk-Stability-Index-current.tex's own v1-5-0/
% v1-6-0 comments for the worked example this was hoisted out of).
% Found 2026-07-20 that the raw-`#1`-in-frame-body bug documented above
% is *not* actually \tableofcontents-specific, contrary to that comment's
% own framing -- slides/scs-Political-Risk-Stability-Index-current.tex has
% no \tableofcontents anywhere in the deck-wide overlay chain leading to
% it, yet hit the identical "Illegal parameter number in definition of
% \iterate" at its own \end{standardframe} from the exact same literal-
% \renewcommand{\makelabel}[1]{...#1...}-in-frame-body pattern. Whatever
% the precise trigger condition actually is, the same hoist-to-preamble
% fix (a bare macro *call*, no `#` token, left in the frame body) sidesteps
% it either way.
\newcommand{\gothamEnumerateLabelFlushLeftShifted}[2]{%
	\setlength{\labelwidth}{#2}%
	\renewcommand{\makelabel}[1]{\hspace{#1}\makebox[\labelwidth][l]{##1}}%
}

% Width of a single digit in this deck's normal content font, measured once
% here (safe, plain preamble-load-time \settowidth) rather than inline in
% the Agenda slide's own frame body -- \tableofcontents there puts the whole
% frame through beamer's overlay-replay machinery, which the comment above
% \gothamTocLabelFlushLeft already found chokes on raw tokens (like a bare
% \settowidth call) sitting directly in a replayed frame body, the same
% "Missing number"/"Illegal unit of measure" class of error as that `#`
% problem, confirmed directly the same way. Used by
% slides/nkj-Agenda-current.tex's own "widen the number-to-title gap"
% fix, which only ever needs to *reference* this precomputed value inside
% the frame, not measure anything itself there.
\newlength{\gothamTocLabelExtraGap}
\settowidth{\gothamTocLabelExtraGap}{0}

% For a two-column layout (e.g. bullets next to a chart/formula), nest
% beamer's own "columns" directly inside standardframe -- do NOT bake
% \begin{columns}/\end{columns} into a custom environment's begin/end code:
% beamer collects each frame's body as a token list to replay it once per
% overlay, and a columns environment that is opened in the "before" half and
% closed in the "after" half of a *different* environment confuses that
% collection (it never finds a properly matched \end{frame}). Nesting it in
% the document body, as below, avoids the issue entirely:
%
%   \begin{standardframe}{Title}
%   	\begin{columns}[t,onlytextwidth]
%   		\begin{column}{0.48\linewidth} ... \end{column}
%   		\begin{column}{0.48\linewidth} ... \end{column}
%   	\end{columns}
%   \end{standardframe}

% For automatic section dividers and single-statement "standout" slides, use
% what the gotham theme already provides instead of a custom environment:
%   \section{...}                                                (section divider, automatic)
%   \begin{frame}[standout]{Key takeaway}\end{frame}             (big single statement)
%   \begin{block}{...}/\begin{alertblock}{...}/\begin{exampleblock}{...}/\begin{referenceblock}{...}  (definitions, formulas, examples, references)
% For a course's own title slide, section-topic overview, quote, figure-only,
% and presenter-attention slides, see the five environments below.

% block/alertblock/exampleblock recolored to this project's own palette
% (2026-07-05, direct instruction) instead of gotham's own default colors
% -- gotham's own per-block-type color machinery (\gotham@block in
% beamerinnerthemegotham.sty) turns out to be dead code, never actually
% wired into any template, so these three blocks were rendering via plain
% beamer's own stock templates/colors the whole time; overriding the
% beamer color names directly here works regardless, since those stock
% templates read their colors from these same standard names (block~title/
% block~body and their ~alerted/~example variants) via beamercolorbox.
% Body background is the same c_tertiary_gray for all four
% block types (including referenceblock below) -- only each one's own
% title/header background differs. block's own header went
% gray→secondaryblue→secondaryviolet across three separate direct
% instructions in as many days -- secondaryviolet is the one actually in
% effect now.
\setbeamercolor{block title}{bg=c_secondary_violet, fg=c_primary_dark}
\setbeamercolor{block body}{bg=c_tertiary_gray, fg=c_primary_dark}
\setbeamercolor{block title example}{bg=c_secondary_green, fg=c_primary_dark}
\setbeamercolor{block body example}{bg=c_tertiary_gray, fg=c_primary_dark}
\setbeamercolor{block title alerted}{bg=c_secondary_red, fg=c_primary_dark}
\setbeamercolor{block body alerted}{bg=c_tertiary_gray, fg=c_primary_dark}

% All four block types' left/right margins widened to match a top-level
% itemize bullet's own text indent (2026-07-05, direct instruction: "the
% blocks should not have a smaller left-hand margin than the content
% text") -- a plain block spans \linewidth as-is, which lines up with the
% *frame's* own text margin, not with an itemize bullet's own text (which
% sits further right, past the bullet symbol and its gap) -- measured
% directly against a compiled slide with both an itemize and a block
% (slides/afa-Excursus-Price-Formation-current.tex): a bullet's own text
% starts 12.96pt further right than a plain block's left edge did.
%
% First attempt used changepage's \adjustwidth, on paper the "correct" tool
% for temporarily changing margins around a chunk of content -- it silently
% does *nothing* inside a beamer frame (beamer computes \linewidth/\hsize
% through its own internal frame-geometry machinery, not the page-geometry
% registers changepage adjusts, which are meant for article/report-style
% documents), and *loading the package at all* actively corrupted the
% title bar/logo on every frame (some interaction with gotham's own
% tikz/remember-picture-based title-bar drawing, not chased further since
% the package needed to go regardless) -- caught by actually compiling and
% looking, not assumed to work from the package's own documented purpose.
% Fixed instead with a plain \hspace{}+\begin{minipage}{\linewidth-2\times}
% wrapped around each block type's own begin/end template pair below --
% minipage genuinely resets \linewidth to its own local width for anything
% typeset inside it (already relied on elsewhere in this deck, e.g. the
% side-by-side text/image slides), which a bare beamercolorbox alone
% doesn't do on its own.
%
% 12.96pt matched a *second*-level bullet's own dot (which happens to
% coincide almost exactly with where a *first*-level bullet's own *text*
% starts, both ~21.6pt in from the frame's raw margin -- a coincidence of
% this deck's specific bullet-dot/gap/font metrics, not a general rule) --
% corrected the next day (direct instruction: align with the *first*-level
% bullet's own margin instead) to 4.20pt, matching a first-level bullet's
% own dot position (~12.84pt in from the frame's raw margin) instead.
\newlength{\gothamBlockSideMargin}
\setlength{\gothamBlockSideMargin}{4.20pt}

% Shared helper for a block title bar that carries a marker symbol flush
% right, opposite the title text (currently: infoblock's circled "i",
% definitionblock's $\triangleq$, exampleblock's $\in$ -- see each of
% those three's own template for the direct instruction behind its own
% marker choice). 2026-07-08, direct instruction, after infoblock's own
% "circled i nearly touching the title text" report on a two-line title
% (slides/u5h-Long-term-Bonds-Provisions-current.tex's own "Denomination,
% face value, or principal" step) -- generalized here so the same fix
% covers all three marker types at once, and any future one, rather than
% patching each block's own template separately.
%
% The marker itself is placed via a zero-width box *before* the title
% text, not via a trailing `\hfill<marker>` the way all three started out
% (definitionblock/exampleblock still do, until updated below) -- `\hfill`
% only fills space on whatever *line* it ends up typeset on, so for a
% title long enough to wrap onto a second line, `\hfill<marker>` lands on
% that second line instead of the first, printing the marker at the
% bottom right of the box instead of the top right (this is the same root
% cause infoblock's own marker already had fixed once already, see that
% template's own comment for the fuller writeup of *that* bug specifically
% -- this pushes the same fix out to the other two, and adds the part
% that bug's own fix didn't yet cover, below).
% `\makebox[0pt][l]{\hspace{\linewidth}\makebox[0pt][r]{<marker>}}` fixes
% *that* much on its own: the outer zero-width box doesn't displace the
% title text at all (still starts at the same point, still wraps however
% it needs to), but *inside* that zero-width box the marker is placed at
% a fixed horizontal offset (this box's own full \linewidth) from the
% title's own starting point, i.e. always the top right corner of the
% title bar regardless of how the title text itself wraps.
%
% What that alone still doesn't fix -- and what this round's actual
% report was about -- is the title text itself having no *awareness* that
% a marker will render on top of wherever it happens to end: a title just
% long enough to reach close to the marker's own fixed position, without
% actually needing to wrap, ends up visually touching or running under
% it. Fixed by also wrapping the title text itself in a
% `\begin{minipage}[t]{\linewidth-\gothamBlockTitleMarkerWidth}`, one
% \gothamBlockTitleMarkerWidth narrower than the full title bar, so the
% text itself wraps *before* reaching the marker's own reserved space
% instead of running into it -- confirmed directly against the compiled
% PDF on the exact case from the report. \gothamBlockTitleMarkerWidth is
% one shared reserved width for all three current marker types rather
% than a separate hand-measured width per symbol, which would need
% re-measuring by hand every time a new marker is added -- comfortably
% covers the widest of the three (checked directly) with a bit to spare.
% #1 = the marker's own content (e.g. \textcircled{\scriptsize i}, or a
% math-mode symbol like $\triangleq$); #2 = the title text
% (\insertblocktitle). Call as
% \gothamBlockTitleWithMarker{<marker>}{\insertblocktitle} in place of a
% bare \insertblocktitle, after the usual \usebeamerfont* call.
\newlength{\gothamBlockTitleMarkerWidth}
\setlength{\gothamBlockTitleMarkerWidth}{16pt}
% 2026-07-08 (direct instruction), two more fixes to this same helper:
% - A title long enough to wrap onto a second line
%   (slides/u5h-Long-term-Bonds-Provisions-current.tex's own
%   "Denomination, face value, or principal") read as *justified* rather
%   than left-aligned -- the minipage below had no explicit alignment, so
%   it fell back to LaTeX's default full justification, which only ever
%   visibly stretches a line short of the box's own right edge, i.e. only
%   shows up once a title is long enough to actually wrap. Every other
%   (single-line, shorter) title already looked left-aligned regardless,
%   since a line that's naturally shorter than the box is never stretched
%   whether or not justification is on -- confirmed directly, this is why
%   the bug went unnoticed until a title finally wrapped. `\raggedright`
%   fixes it without needing anything different for the single-line case.
% - Different titles' own title *bars* rendered at different heights
%   depending on whether that particular title's text happened to include
%   a descender (e.g. "Security"/"y", "Moody's rating at issue date"/"y")
%   or not (e.g. "Registered", "Trustee") -- the same class of bug already
%   fixed on slides/qxr-Introduction-to-R-Assignments-current.tex's own
%   colorboxes (see that slide's own comment): a box that tightly wraps its
%   own text naturally varies in height/depth with that text's own
%   ascenders/descenders. A leading `\strut` gives every title's first (and
%   for all but this one current title, only) line the same guaranteed
%   minimum height/depth, matching the qxr fix's own reasoning -- confirmed
%   directly, every current info/definition/example block title bar now
%   renders at the same height regardless of its own text.
\newcommand{\gothamBlockTitleWithMarker}[2]{%
	\makebox[0pt][l]{\hspace{\linewidth}\makebox[0pt][r]{#1}}%
	\begin{minipage}[t]{\dimexpr\linewidth-\gothamBlockTitleMarkerWidth\relax}\raggedright\strut#2\end{minipage}%
}

% A fourth block type, for a reference/citation-style callout (first used
% on slides/afa-Excursus-Price-Formation-current.tex and
% slides/4r7-IPO-Underpricing-International-current.tex, both replacing
% an earlier rotated hand-drawn callout box) -- beamer has no built-in
% "reference" block variant, so this defines one from scratch, mirroring
% exactly how beamer's own alertblock/exampleblock are built (compare
% beamerbaselocalstructure.sty's \newenvironment<>{alertblock} and
% beamerinnerthemedefault.sty's "block alerted begin/end" templates).
%
% First attempt (2026-07-05) copied these under the "rounded" style/
% beamerboxesrounded, on the wrong assumption that "rounded" was this
% deck's active block style generally -- it renders with visibly rounded
% corners, caught by the user, unlike block/alertblock/exampleblock's own
% square corners. Those three actually use beamerinnerthemedefault.sty's
% *starred* "default"-style templates (plain beamercolorbox, no rounding),
% which \defbeamertemplate*'s star selects immediately when that inner
% theme loads, regardless of beamer's own \defbeamertemplateparent{blocks}
% fallback to "rounded" -- gotham's own inner theme derives from
% beamerinnerthemedefault, so "default" (square) was the deck's real active
% style the whole time. Rewritten below to match that "default" style
% exactly (plain beamercolorbox instead of beamerboxesrounded) for
% consistent square corners across all four block types.
\setbeamercolor{block title reference}{bg=c_secondary_orange, fg=c_primary_dark}
\setbeamercolor{block body reference}{bg=c_tertiary_gray, fg=c_primary_dark}
% block/alertblock/exampleblock's own titles render bold via
% beamerfontthemegotham.sty's \setbeamerfont{block title}{series=\bfseries}
% (and an explicit ...alerted equivalent; example inherits it from
% beamer's own base font-fallback chain, verified directly -- both
% "Example"/"Caveat" titles on slides/ZZX-Definition-formula-current.tex
% render in the same heavier FiraSans-Regular as "Definition" despite
% gotham never setting "block title example" explicitly). "reference" is a
% brand new type name unknown to any of that inheritance, so it fell back
% to a plain, non-bold font with no color-analogous safety net (unlike
% \setbeamercolor{block title reference} above, colors for genuinely new
% names still resolve fine) -- caught only by comparing the actual
% rendered font name in the compiled PDF (FiraSans-Light, matching plain
% body text) against block/alertblock/exampleblock's own (FiraSans-Regular,
% this deck's de facto "bold" for Fira Sans), not by eye alone. Explicit
% setting needed, matching gotham's own two.
\setbeamerfont{block title reference}{size=\normalsize, series=\bfseries}
\defbeamertemplate*{block reference begin}{default}[1][]
{
	\hspace{\gothamBlockSideMargin}\begin{minipage}{\dimexpr\linewidth-2\gothamBlockSideMargin\relax}
	\par\vskip\medskipamount%
	\begin{beamercolorbox}[colsep*=.75ex,#1]{block title reference}
		\usebeamerfont*{block title reference}\strut\insertblocktitle%
	\end{beamercolorbox}\noframebreak%
	{\parskip0pt\par}%
	\ifbeamercolorempty[bg]{block title reference}
	{}
	{\ifbeamercolorempty[bg]{block body reference}{}{\nointerlineskip\vskip-0.5pt}}%
	\usebeamerfont{block body reference}%
	\begin{beamercolorbox}[colsep*=.75ex,vmode,#1]{block body reference}%
		\ifbeamercolorempty[bg]{block body reference}{\vskip-.25ex}{\vskip-.75ex}\vbox{}%
}
\defbeamertemplate*{block reference end}{default}
{\end{beamercolorbox}\vskip\smallskipamount\end{minipage}}
\newenvironment<>{referenceblock}[1]{%
	\begin{actionenv}#2%
		\def\insertblocktitle{#1}%
		\par%
		\usebeamertemplate{block reference begin}}
	{\par%
		\usebeamertemplate{block reference end}%
	\end{actionenv}}

% A sixth block type, for general background/context info that isn't
% itself a definition, example, or reference citation (2026-07-06, direct
% instruction: "use secondaryblue for a new type of info block") -- built
% the same from-scratch way as referenceblock just above (a brand-new
% style name needs \defbeamertemplate*'s registering form, not
% \setbeamertemplate -- see referenceblock's own comment for why). Marker
% is a circled "i" (\textcircled{\scriptsize i}, the generic LaTeX-kernel
% \textcircled -- not a Unicode "circled i" character: this deck's own
% Fira Sans/Fira Mono have no glyph for that codepoint either, the same
% missing-glyph problem \triangleq's own comment already found for
% definitionblock, and \textcircled sidesteps it entirely since it just
% draws a plain circle symbol around an ordinary letter rather than
% depending on one precomposed font glyph existing at all).
%
% The marker is placed via a zero-width box *before* \insertblocktitle,
% not via a trailing \hfill\textcircled{...} the way it first looked most
% natural to write -- \hfill only fills space on whatever line it's
% typeset on, and for a title long enough to wrap onto a second line (2026-07-08,
% found on the last reveal step of slides/ib3-Long-term-Bonds-Repayment-current.tex),
% \hfill\textcircled{...} then lands on that *second* line instead of the
% first, printing the marker at the bottom right of the box instead of the
% top right. Fixed originally with a one-off zero-width-box trick here;
% generalized the same day into the shared \gothamBlockTitleWithMarker
% helper (see its own comment, near \gothamBlockSideMargin above) once a
% *second*, related bug turned up on this exact block type (a long title
% visually touching the marker even without wrapping) -- that helper is
% what's actually called below now, not a copy of this trick inline.
\setbeamercolor{block title info}{bg=c_secondary_blue, fg=c_primary_dark}
\setbeamercolor{block body info}{bg=c_tertiary_gray, fg=c_primary_dark}
\setbeamerfont{block title info}{size=\normalsize, series=\bfseries}
\defbeamertemplate*{block info begin}{default}[1][]
{
	\hspace{\gothamBlockSideMargin}\begin{minipage}{\dimexpr\linewidth-2\gothamBlockSideMargin\relax}
	\par\vskip\medskipamount%
	\begin{beamercolorbox}[colsep*=.75ex,#1]{block title info}
		\usebeamerfont*{block title info}%
		\gothamBlockTitleWithMarker{\textcircled{\scriptsize i}}{\insertblocktitle}%
	\end{beamercolorbox}\noframebreak%
	{\parskip0pt\par}%
	\ifbeamercolorempty[bg]{block title info}
	{}
	{\ifbeamercolorempty[bg]{block body info}{}{\nointerlineskip\vskip-0.5pt}}%
	\usebeamerfont{block body info}%
	\begin{beamercolorbox}[colsep*=.75ex,vmode,#1]{block body info}%
		\ifbeamercolorempty[bg]{block body info}{\vskip-.25ex}{\vskip-.75ex}\vbox{}%
}
\defbeamertemplate*{block info end}{default}
{\end{beamercolorbox}\vskip\smallskipamount\end{minipage}}
\newenvironment<>{infoblock}[1]{%
	\begin{actionenv}#2%
		\def\insertblocktitle{#1}%
		\par%
		\usebeamertemplate{block info begin}}
	{\par%
		\usebeamertemplate{block info end}%
	\end{actionenv}}

% A fifth block type, for a "this introduces/defines a term" callout,
% visually distinguished from a plain \block (2026-07-06, direct
% instruction: "I would like to distinguish definition blocks from other
% blocks") by a $\triangleq$ symbol sitting flush right in the title bar,
% opposite the title text itself. Exact same colors/margins/font as plain
% \block (reuses "block title"/"block body", not a new color pair) --
% only the title-bar content differs, one extra `\hfill$\triangleq$`
% after `\insertblocktitle` -- so this mirrors referenceblock's own
% from-scratch pattern above (a brand new "definition" style name via
% \defbeamertemplate*, not \setbeamertemplate, for the same reason
% referenceblock's own comment gives: this style name has never been
% registered before, so the *-registration form is the correct one here,
% unlike plain block/alertblock/exampleblock below which are redefining
% already-registered "default"-style bodies). `\triangleq` is a genuine
% math symbol, so it needs `$...$` -- renders in this theme's math font
% (Computer Modern, see beamerfontthemegotham.sty's `no-math` option),
% a deliberate, visible one-glyph exception in an otherwise Fira Sans
% title bar, the same tradeoff every other math symbol in this deck
% already accepts.
\defbeamertemplate*{block definition begin}{default}[1][]
{
	\hspace{\gothamBlockSideMargin}\begin{minipage}{\dimexpr\linewidth-2\gothamBlockSideMargin\relax}
	\par\vskip\medskipamount%
	\begin{beamercolorbox}[colsep*=.75ex,#1]{block title}
		\usebeamerfont*{block title}\gothamBlockTitleWithMarker{$\triangleq$}{\insertblocktitle}%
	\end{beamercolorbox}\noframebreak%
	{\parskip0pt\par}%
	\ifbeamercolorempty[bg]{block title}
	{}
	{\ifbeamercolorempty[bg]{block body}{}{\nointerlineskip\vskip-0.5pt}}%
	\usebeamerfont{block body}%
	\begin{beamercolorbox}[colsep*=.75ex,vmode,#1]{block body}%
		\ifbeamercolorempty[bg]{block body}{\vskip-.25ex}{\vskip-.75ex}\vbox{}%
}
\defbeamertemplate*{block definition end}{default}
{\end{beamercolorbox}\vskip\smallskipamount\end{minipage}}
\newenvironment<>{definitionblock}[1]{%
	\begin{actionenv}#2%
		\def\insertblocktitle{#1}%
		\par%
		\usebeamertemplate{block definition begin}}
	{\par%
		\usebeamertemplate{block definition end}%
	\end{actionenv}}

% block/alertblock/exampleblock begin/end: same body as
% beamerinnerthemedefault.sty's own "default"-style templates (the ones
% actually in effect, per the comment above referenceblock), plus the
% minipage margin wrapper around each pair -- set via plain
% \setbeamertemplate{block begin}{...} (overriding the *body* of the
% existing "default" style directly) rather than
% \defbeamertemplate*{block begin}{default}{...} like referenceblock's own
% (new) templates above. The two look almost interchangeable but aren't:
% \defbeamertemplate* internally does a \newcommand-based *registration* of
% a style name, which errors outright ("... already defined") when that
% exact template+style pair already exists, as "block begin"/"default" and
% the other two pairs already do (registered once by
% beamerinnerthemedefault.sty itself) -- caught only by compiling and
% reading the log, not from any visible symptom in the rendered PDF (LaTeX
% recovers from the error by leaving the *original* stock template body in
% place, so these three blocks kept working and even picked up their new
% *colors* correctly via \setbeamercolor above, since color lookups happen
% at render time either way -- just silently without the margin fix, which
% *is* baked into the template body itself). \setbeamertemplate{name}{body}
% (no style-name argument) instead does a plain \def of that style's
% existing body storage, which is exactly "redefine what's already there"
% and carries no such restriction -- the correct tool whenever the target
% template+style pair is already registered, as opposed to
% \defbeamertemplate*, which is only for registering a style name for the
% very first time (as referenceblock's brand-new "reference" style name
% needs). Dropped the stock templates' own "#1" optional-argument slot
% (extra beamercolorbox options) in the process: nothing in this deck ever
% calls \usebeamertemplate{block begin}[...] with a bracket argument to
% populate it, so it was always empty in practice anyway.
% \strut before \insertblocktitle here, and in "block alerted"/"block
% reference" below: same title-bar-height-vs-descender fix as
% \gothamBlockTitleWithMarker's own comment above -- these three don't go
% through that helper at all (no marker symbol), so needed the same
% `\strut` added directly instead.
\setbeamertemplate{block begin}{
	\hspace{\gothamBlockSideMargin}\begin{minipage}{\dimexpr\linewidth-2\gothamBlockSideMargin\relax}
	\par\vskip\medskipamount%
	\begin{beamercolorbox}[colsep*=.75ex]{block title}
		\usebeamerfont*{block title}\strut\insertblocktitle%
	\end{beamercolorbox}\noframebreak%
	{\parskip0pt\par}%
	\ifbeamercolorempty[bg]{block title}
	{}
	{\ifbeamercolorempty[bg]{block body}{}{\nointerlineskip\vskip-0.5pt}}%
	\usebeamerfont{block body}%
	\begin{beamercolorbox}[colsep*=.75ex,vmode]{block body}%
		\ifbeamercolorempty[bg]{block body}{\vskip-.25ex}{\vskip-.75ex}\vbox{}%
}
\setbeamertemplate{block end}
{\end{beamercolorbox}\vskip\smallskipamount\end{minipage}}

% Warning-triangle marker for alertblock (direct instruction, 2026-07-12),
% same \gothamBlockTitleWithMarker treatment as definitionblock's
% $\triangleq$/exampleblock's $\in$/infoblock's \textcircled{\scriptsize i}
% above. No single built-in glyph renders a triangle-with-exclamation-mark
% in this deck's fonts (checked the same way \triangleq's own comment
% did), so it's hand-drawn: a plain tikz outline triangle with a bold "!"
% centered inside, sized by eye against \textcircled{\scriptsize i} until
% the two read as the same visual weight in a block title bar.
\newcommand{\gothamAlertMarker}{%
	\tikz[baseline=-0.6ex]{
		\draw[line width=1.1pt, line join=round] (-0.2,-0.11) -- (0.2,-0.11) -- (0,0.31) -- cycle;
		\node[font=\fontsize{7}{7}\selectfont\bfseries] at (0,0.01) {!};
	}%
}
\setbeamertemplate{block alerted begin}{
	\hspace{\gothamBlockSideMargin}\begin{minipage}{\dimexpr\linewidth-2\gothamBlockSideMargin\relax}
	\par\vskip\medskipamount%
	\begin{beamercolorbox}[colsep*=.75ex]{block title alerted}
		\usebeamerfont*{block title alerted}\gothamBlockTitleWithMarker{\gothamAlertMarker}{\insertblocktitle}%
	\end{beamercolorbox}\noframebreak%
	{\parskip0pt\par}%
	\ifbeamercolorempty[bg]{block title alerted}
	{}
	{\ifbeamercolorempty[bg]{block body alerted}{}{\nointerlineskip\vskip-0.5pt}}%
	\usebeamerfont{block body alerted}%
	\begin{beamercolorbox}[colsep*=.75ex,vmode]{block body alerted}%
		\ifbeamercolorempty[bg]{block body alerted}{\vskip-.25ex}{\vskip-.75ex}\vbox{}%
}
\setbeamertemplate{block alerted end}
{\end{beamercolorbox}\vskip\smallskipamount\end{minipage}}

% $\in$ ("is an instance/element of") flush right in the title bar
% (2026-07-06, direct instruction, mirroring definitionblock's own
% $\triangleq$ -- "what symbol would you propose ... to symbolize an
% example?"): a concrete example is naturally "an instance of" the general
% concept it illustrates, the same relationship $\in$ already denotes
% between an element and a set/class -- and, same as $\triangleq$, it's a
% math-mode symbol with a real Computer Modern glyph (no Fira Sans/Fira
% Mono text-mode substitute needed, unlike most Unicode symbol candidates
% -- see Slidemaster_preamble_uni-graz.tex's own "beamer/TeX gotchas" entry on
% $\triangleq$ for why that matters here). Applied to every \exampleblock/
% \gothamfloatingexampleblock in the deck via this shared template, not a
% one-off per slide, matching how $\triangleq$ applies to every
% definitionblock through its own shared template.
\setbeamertemplate{block example begin}{
	\hspace{\gothamBlockSideMargin}\begin{minipage}{\dimexpr\linewidth-2\gothamBlockSideMargin\relax}
	\par\vskip\medskipamount%
	\begin{beamercolorbox}[colsep*=.75ex]{block title example}
		\usebeamerfont*{block title example}\gothamBlockTitleWithMarker{$\in$}{\insertblocktitle}%
	\end{beamercolorbox}\noframebreak%
	{\parskip0pt\par}%
	\ifbeamercolorempty[bg]{block title example}
	{}
	{\ifbeamercolorempty[bg]{block body example}{}{\nointerlineskip\vskip-0.5pt}}%
	\usebeamerfont{block body example}%
	\begin{beamercolorbox}[colsep*=.75ex,vmode]{block body example}%
		\ifbeamercolorempty[bg]{block body example}{\vskip-.25ex}{\vskip-.75ex}\vbox{}%
}
\setbeamertemplate{block example end}
{\end{beamercolorbox}\vskip\smallskipamount\end{minipage}}

% A narrower, freely-positionable version of \begin{block}{...} (2026-07-06,
% direct instruction) -- a plain block always spans the *current*
% \linewidth (that's what block begin/end's own \hspace+minipage wrapper
% above is measured against), so it can't be dropped into a fixed-width
% spot (e.g. beside a table, inside a tikz overlay \node) without first
% narrowing \linewidth some other way. \gothamfloatingblock just wraps a
% plain block in one more, *outer* minipage of the caller's own chosen
% width -- block begin/end's own inner minipage then measures itself
% against that narrower width instead of the frame's full one, nesting
% correctly with no changes needed to the block templates themselves.
% Title/body colors, margins, and font are exactly \block's own (same
% "block title"/"block body" beamer color names) -- this only changes the
% overall width and (by taking it out of the normal document flow, left to
% the caller) the position, not the look. Font size is also left to the
% caller (wrap the whole environment in a size declaration, e.g.
% {\footnotesize \begin{gothamfloatingblock}{5cm}{Excel} ... \end{gothamfloatingblock}})
% rather than fixed here, since the right size depends entirely on how
% much room the block ends up placed in.
\newenvironment{gothamfloatingblock}[2]{%
	\begin{minipage}{#1}
	\begin{block}{#2}
}{%
	\end{block}
	\end{minipage}
}

% Same idea, but wrapping \exampleblock instead of \block (2026-07-06,
% direct instruction, first used on
% slides/fr3-Forms-of-Capital-Market-Efficiency-Strong-Semistrong-Weak-current.tex
% and slides/ha9-Monacorp-Truck-Lease-Breakeven-Rent-current.tex, both
% switching an existing \gothamfloatingblock to the "example" green header
% instead of the plain violet one) -- not worth generalizing into a single
% macro parameterized over the block name: \gothamfloatingblock's own body
% is only two lines, so a third variant is cheaper to just write out than
% to build a block-name-dispatch layer for.
\newenvironment{gothamfloatingexampleblock}[2]{%
	\begin{minipage}{#1}
	\begin{exampleblock}{#2}
}{%
	\end{exampleblock}
	\end{minipage}
}

% Same idea again, wrapping \infoblock (2026-07-06, direct instruction,
% first used on
% slides/fr3-Forms-of-Capital-Market-Efficiency-Strong-Semistrong-Weak-current.tex's
% own "Random walk" formula box, switched from the green
% \gothamfloatingexampleblock to this blue one the same day).
\newenvironment{gothamfloatinginfoblock}[2]{%
	\begin{minipage}{#1}
	\begin{infoblock}{#2}
}{%
	\end{infoblock}
	\end{minipage}
}

% Every gothamfloatinginfoblock anchored top-right against the institute
% logo (the deck's standard
% "(current page.north east)+(-\gothamFrametitleRightpading,-\gothamContentTopAlign)"
% corner pattern) lands its own visible right edge short of the logo's,
% by a fixed amount independent of the block's own chosen width (#1) --
% \infoblock's underlying template insets its colored fill by
% \gothamBlockSideMargin from *both* sides of whatever width it's given
% (that margin exists to match a first-level bullet's own left-hand
% position, see its own comment above; the fact that it also narrows the
% *right* side is just a side effect of reusing one symmetric
% hspace+minipage trick for both, not a deliberate right-margin choice),
% partly offset by beamer's own beamercolorbox "colsep*" background bleed
% (colsep* pads the colored fill outward past its own nominal box edge
% without also insetting the text, unlike unstarred "colsep" -- see
% beamerbasecolor.sty's own \beamercolorbox definition). Root-caused and
% measured directly (pymupdf against a compiled probe deck, not derived
% purely on paper) 2026-07-21, after
% slides/j8y-Asset-Allocation-Example-Weights-current.tex's own one-off
% "+2.2pt" nudge (measured that same day against just the institute
% logo's yellow chip) turned out to disagree with
% slides/fr3-Forms-of-Capital-Market-Efficiency-Strong-Semistrong-Weak-current.tex's
% own older "+3.165pt" nudge (2026-07-06, measured against the logo's
% *full* gray-bordered footprint, shadow/card included) -- both slides'
% own nudges were each individually correct for their own chosen
% reference edge, they just didn't agree on what "the logo's own edge"
% means. Direct instruction, 2026-07-21: the *full* gray-bordered
% footprint is the deck's intended reference, so 3.165pt is the value
% kept here; j8y's own newer 2.2pt nudge was the one revised to match,
% not this one.
%
% Use "-\gothamFloatingInfoblockRightAlign" (not the bare
% "-\gothamFrametitleRightpading") for a gothamfloatinginfoblock's own
% anchor node whenever its right edge is meant to align with the logo --
% every current use of this pattern across the deck has been updated to
% match; new slides using this same corner pattern should use this
% length from the start rather than re-measuring their own one-off nudge.
\newlength{\gothamFloatingInfoblockRightAlign}
\setlength{\gothamFloatingInfoblockRightAlign}{\dimexpr\gothamFrametitleRightpading-3.165pt\relax}

%%%%%%%%%%%%%%%%%%%%%%%%%%%%%%%%%%%%%%%%%%%%%%%%%%%%%%%%%%%%%%%%%%%%%%%%%%%%%
% Special frame types (converted from ExampleSlides.pptx; see README.md)
%
% All five below deliberately reuse as much of standardframe -- and of each
% other -- as possible, rather than reimplementing frame chrome from
% scratch: sectionframe, figureframe, and highlightframe are literally thin
% wrappers around standardframe itself (same title bar, logo, footer, bullet
% styling; only what's documented at each one differs). titleframe and
% quoteframe can't be standardframe wrappers -- both are title-less, and
% quoteframe replaces the whole body layout -- but still reuse its palette,
% fonts, and margins (\gothamFrametitleLeftpading, the institute logo
% drawing command, etc.) rather than inventing new ones.
%%%%%%%%%%%%%%%%%%%%%%%%%%%%%%%%%%%%%%%%%%%%%%%%%%%%%%%%%%%%%%%%%%%%%%%%%%%%%

% titleframe: the deck's own title slide (ExampleSlides.pptx slide 1) --
% typically used exactly once, at the very start of a course's first deck.
% One continuous yellow box -- like standardframe's own title bar -- flush
% with the page's top and left edges, stopping short of the institute logo
% on the right (the same margin already used around the logo elsewhere:
% \gothamLogoWidth + 2*\gothamFrametitleRightpading), holding the title and,
% below it, a second line (author/term/whatever) in the same black text as
% standardframe's own frametitle. Usage:
%
%   \titleframe{Financial Markets}{Stefan Palan\\Summer term 2026}
%
% A *command*, not an environment, taking the second (lower) line's content
% as its second argument -- plain text, use \\ for line breaks same as any
% normal LaTeX paragraph, not an itemize, since the source deck's own
% subtitle block is just stacked lines. This has to be a command rather
% than an environment with the body as its content: beamer's frame body is
% captured by a delimited-argument macro (\beamer@collect@@body, looked up
% by triggering the error below on purpose) hunting for a literal
% \end{frame}, and a beamercolorbox -- needed by an earlier version of this
% to wrap the subtitle text in its own colored band -- opened in a *custom
% environment's* begin code and closed in its end code confuses that hunt
% exactly like the \columns example above (same root cause, different
% environment): with
% \newenvironment{titleframe}{...\begin{beamercolorbox}{...}}{\end{beamercolorbox}...},
% compiling ends in "! File ended while scanning use of \beamer@collect@@body."
% every time, however short the body. Taking the content as a normal brace
% argument instead sidesteps the whole issue. No footer/frame number
% (matches beamer's own convention for title slides, e.g. the old
% \begin{frame}[plain,noframenumbering]\titlepage\end{frame} this replaces --
% see README.md) and no chapter-logo support: both would be meaningless on a
% slide that appears once, before any chapter exists.
%
% Unlike the frametitle hook further up, neither the logo nor the yellow box
% here need the white-patch-and-redraw trick: the box's own width already
% stops short of the logo (rather than being painted full-width and then
% partly painted back over), so there's no yellow left to paint over in the
% first place. Both still need \tikz's "remember picture" on their
% "current page"-anchored nodes, though, and so still need the usual second
% compile pass (see "Title bar and logo" in README.md) -- checked
% empirically: without it, "current page" doesn't resolve to the actual page
% box on the first pass, and the logo/yellow box end up positioned nowhere
% near the page corners they're meant to be anchored to.
%
% \gothamTitleframeBoxHeight is a fixed target for the box's *bottom* edge
% (matching what was actually asked for: "reaching down to where the lower
% grey bar currently ends" in the two-band design this replaces), not the
% blank space above the text -- that space is derived from it instead (box
% height minus the title+subtitle content's own measured height), so the
% bottom edge stays fixed even if a longer title ever wraps onto a second
% line; the blank space above would just shrink to compensate. Measured
% against the old two-band design's own rendered PDF coordinates (before
% this change): its second band's bottom edge sat 187.147pt below the page's
% top edge.
\newlength{\gothamTitleframeBoxHeight}
\setlength{\gothamTitleframeBoxHeight}{187.147pt}
\newlength{\gothamTitleframeBoxWidth}
\setlength{\gothamTitleframeBoxWidth}{\dimexpr\paperwidth-\gothamLogoWidth-2\gothamFrametitleRightpading\relax}
% Text inset from the box's own edges: \gothamFrametitleLeftpading, the same
% length standardframe's own title text uses, plus \gothamTitleframeTextSep
% (2.5ex, matching this command's own former per-band beamercolorbox "sep").
\newlength{\gothamTitleframeTextSep}
\setlength{\gothamTitleframeTextSep}{2.5ex}
\newlength{\gothamTitleframeBandGap}
\setlength{\gothamTitleframeBandGap}{0.8ex}
% Because the leading blank space above the text (further down) is itself
% computed as "box height minus the content's own natural height", the
% content's *bottom* padding is the only thing that determines how far the
% last line sits above the box's bottom edge -- the leading space silently
% absorbs whatever that padding is, keeping the bottom edge fixed either
% way. (The relationship is a direct 1:1 offset, checked empirically by
% comparing two rendered values, but with a nonzero intercept from beamer's
% own default top-of-frame vskip -- there's no clean closed-form for the
% right value, hence the empirical calibration here rather than a formula.)
% Calibrated (against the rendered PDF's own text coordinates, not just
% visually) to about 1.5 line-heights of the subtitle text's own \Large
% size (~20.6pt/line here) above the box's bottom edge; re-measure and
% adjust if the subtitle's font size, or beamer's own default top-of-frame
% vskip, ever change.
\newlength{\gothamTitleframeBottomPad}
\setlength{\gothamTitleframeBottomPad}{41pt}
\newbox\gothamTitleframeBox
\makeatletter
\newcommand{\titleframe}[2]{%
	\begin{frame}[plain,noframenumbering,t]%
	% A bare tikzpicture behaves like an hbox (same reasoning as the footer
	% box's own \vbox-to-0pt wrapper further up): sitting directly in the
	% frame's vertical list, it silently opens an implicit paragraph,
	% switching into horizontal mode. Everything placed after it -- the
	% \vspace* further down included -- would then run in horizontal mode
	% instead of vertical mode, which is what a first version of this command
	% got wrong: with two bare tikzpictures here, \vspace*'s glue ended up as
	% *horizontal* space on an implicit line together with them (visible as
	% the title/subtitle shifting sideways, not down, and the vspace amount
	% appearing to have no effect at all -- it was moving the wrong axis).
	% \vbox to 0pt{\vss ...} keeps each overlay at zero space in the normal
	% flow without ever leaving vertical mode.
	%
	% \gothamArtifactBegin/\gothamArtifactEnd (defined with the footer, "PDF
	% accessibility" comment): this box is purely decorative background, same
	% kind of element as the footer's own yellow box -- marked the same way,
	% for a screen reader to skip. titleframe has no footer at all to worry
	% about separately (it's a "plain" frame, no footline), so this box is
	% the only thing here that needs it; the title/subtitle text right below
	% is real content and stays unmarked.
	\gothamArtifactBegin{/Type /Background}%
	\vbox to 0pt{\vss
	\begin{tikzpicture}[overlay, remember picture]
		% \gothamTitleBoxColor, not a literal color token -- same setting the
		% standard frametitle bar (colorA, above) uses, so the title slide's
		% own box matches every other title bar in the deck (2026-07-16,
		% direct instruction).
		% outer sep=0pt here and on the logo below: without it, tikz's
		% default outer sep insets the fill/logo 0.2pt from the anchored
		% page corner -- a hairline white line above/left of this box, and
		% its right edge/the logo 0.2-0.4pt off every other frame's. See
		% the frametitle hook's own comment for the full story.
		\node[fill=\gothamTitleBoxColor, anchor=north west, inner sep=0pt, outer sep=0pt,
		      minimum width=\gothamTitleframeBoxWidth,
		      minimum height=\gothamTitleframeBoxHeight]
			at (current page.north west) {};
	\end{tikzpicture}%
	}%
	\gothamArtifactEnd
	\vbox to 0pt{\vss
	\begin{tikzpicture}[overlay, remember picture]
		% inner ysep derived, not \gothamFrametitleRightpading -- see the
		% frametitle hook's fixed-position logo redraw for why (matches the
		% inline logo's vertically-centered position on every other frame).
		\node[anchor=north east, inner xsep=\gothamFrametitleRightpading,
		      inner ysep=\dimexpr\gothamTitleBarBottomAlign-\gothamLogoHeight\relax,
		      outer sep=0pt]
			at (current page.north east) {\gothamInstituteLogoSquare};
	\end{tikzpicture}%
	}%
	\setbox\gothamTitleframeBox=\vbox{%
		\hsize=\gothamTitleframeBoxWidth
		\parindent=0pt
		\leftskip=\gothamFrametitleLeftpading
		\advance\leftskip by \gothamTitleframeTextSep
		\rightskip=\gothamFrametitleLeftpading
		\advance\rightskip by \gothamTitleframeTextSep
		% \gothamTitleTextColor, not a literal color token -- follows the
		% box color change right above (2026-07-16, direct instruction):
		% leaving this hardcoded dark would be unreadable once the box
		% itself switches to a dark blue for MCI/Innsbruck.
		\color{\gothamTitleTextColor}%
		% Beamer redefines \\ inside every frame to peek for an overlay spec
		% (\\<2->, for line-by-line reveals in itemize/enumerate); that
		% redefinition assumes it's running in beamer's own frame-body
		% machinery and chokes ("Undefined control sequence", silently
		% swallowing the following word) when \\ is used for a plain line
		% break in a raw \vbox like this one. Restoring plain LaTeX's own
		% \\ (\@normalcr) for this group only fixes it without affecting \\
		% anywhere else in the deck. (The @ here needs catcode 11, which is
		% why this whole command definition -- not just this line -- is
		% wrapped in \makeatletter/\makeatother: macro bodies are tokenized
		% at *definition* time, so a \makeatletter placed inside the body
		% itself would take effect too late, after "\@normalcr" had already
		% been mis-tokenized as "\@" followed by four ordinary-text letters.)
		\let\\\@normalcr
		\vskip\gothamTitleframeTextSep
		\noindent{\Huge\bfseries #1\par}%
		\vskip\gothamTitleframeTextSep
		\vskip\gothamTitleframeBandGap
		\vskip\gothamTitleframeTextSep
		\noindent{\Large #2\par}%
		\vskip\gothamTitleframeBottomPad
	}%
	\vspace*{\dimexpr\gothamTitleframeBoxHeight-\ht\gothamTitleframeBox-\dp\gothamTitleframeBox\relax}%
	\box\gothamTitleframeBox
	% This command's own \color{\gothamTitleTextColor} above (set inside a
	% stored box register) can leak into the *next* frame's body text --
	% see standardframe's own opening code ("Forces this frame's own body")
	% for the real fix and the full story (a trailing reset attempted
	% here first turned out to be a silent no-op, since nothing draws after
	% it within this same frame to actually force the color operator to be
	% emitted).
	\end{frame}%
}
\makeatother

% fullbleedframe: a slide that is *only* a full-screen image (FM2_2026s_ps_deck.pptx
% slide 13, a recruitment flyer pasted in as one embedded picture on an
% otherwise blank layout -- no title, no PPTX text at all) -- no title bar,
% no footer, no logo, matching that source's own blank layout intent, unlike
% figureframe (which is standardframe underneath and always keeps the usual
% title/footer chrome). The image is scaled to fit the slide while keeping
% its own aspect ratio (\includegraphics's keepaspectratio, capped by
% \paperwidth *and* \paperheight together), centered; whatever's left over
% on the binding axis (e.g. a widescreen slide holding a taller/narrower
% image) is filled by a solid black backdrop rather than left showing the
% frame's usual white canvas -- the point of a full-*bleed* image is that
% nothing reads as "empty margin" even when the aspect ratios don't match.
%
%   \fullbleedframe{media/photo-datacamp-recruitment-flyer}
%
% \gothamFullBleedBlack is a fixed HTML black, deliberately *not* a
% \colorlet of c_primary_dark: this frame type's backdrop should stay
% neutral true black regardless of whether c_primary_dark itself is ever
% restyled to something else (e.g. a softer near-black, the way
% c_custom_gray already is elsewhere in this file) for branding
% reasons -- those are two different jobs (brand color vs. letterbox
% backdrop) that only happen to want the same value today.
\definecolor{gothamFullBleedBlack}{HTML}{000000}
\newcommand{\fullbleedframe}[1]{%
	\begin{frame}[plain,noframenumbering,t]%
	% \gothamArtifactBegin/\gothamArtifactEnd (defined with the footer, "PDF
	% accessibility" comment): the black backdrop is purely decorative
	% letterboxing, marked the same way as titleframe's/quoteframe's own
	% background boxes so a screen reader skips it. The image itself is left
	% unmarked -- same as everywhere else in this file, since this project
	% has no alt-text infrastructure yet for imported images generally (a
	% separate, not-yet-made decision -- see README.md's accessibility
	% section for the same open question about the institute logo).
	\gothamArtifactBegin{/Type /Background}%
	\vbox to 0pt{\vss
	\begin{tikzpicture}[overlay, remember picture]
		% outer sep=0pt: see the frametitle hook's comment -- without it the
		% backdrop sits 0.2pt in from the anchored page corner.
		\node[fill=gothamFullBleedBlack, anchor=south west, inner sep=0pt, outer sep=0pt,
		      minimum width=\paperwidth, minimum height=\paperheight]
			at (current page.south west) {};
	\end{tikzpicture}%
	}%
	\gothamArtifactEnd
	\vbox to 0pt{\vss
	\begin{tikzpicture}[overlay, remember picture]
		\node[anchor=center] at (current page.center)
			{\includegraphics[width=\paperwidth, height=\paperheight, keepaspectratio]{#1}};
	\end{tikzpicture}%
	}%
	\end{frame}%
}

% sideimageframe: a standardframe that always carries one auto-placed
% "side filler image" (the user's own term, see
% slides/dfc-Angel-Investors-current.tex's own comment for the technique
% this formalizes into a reusable frame type) -- 2026-07-06, direct
% instruction. The image is positioned flush to three edges at once: its
% top touches the title bar's own bottom edge, its right touches the
% page's own right edge, and its bottom touches the footer box's own top
% edge -- so it's scaled by *height* to exactly
% \paperheight-\gothamTitleBarBottomAlign-\gothamFooterBoxHeight (the
% measured gap between those two), with `keepaspectratio` so the width
% simply follows from the image's own aspect ratio rather than being
% forced/distorted to some second, independent target -- anchoring the
% node at the page's own north-east corner then makes the top+right
% touches automatic, and the height exactly filling that gap makes the
% bottom touch automatic too, with no separate width constraint needed
% (only a *square* source image would ever visibly reach some further
% left boundary; this only guarantees top/right/bottom, matching what was
% actually asked for).
%
%   \begin{sideimageframe}{Title}{media/photo-example}
%   	\begin{itemize}
%   		\item ...
%   	\end{itemize}
%   	\source{...}
%   \end{sideimageframe}
%
% \sidefootnotes is switched on automatically (not left to the caller) --
% the whole point of this frame type is that its footnotes always end up
% rotated along the left edge of the image rather than colliding with it
% the way a full-width bottom footnote would (see
% slides/u29-Implications-of-Capital-Market-Efficiency-Prices-current.tex's
% own comment for exactly that collision, on a plain standardframe). Unlike
% every other \sidefootnotes use elsewhere in this deck, though, the strip
% can't sit at its usual fixed inset from the page's own right edge here --
% that would run it directly underneath the image, the same problem u29
% had. Fixed by *computing* the image's own rendered width (via
% \settowidth against the exact same \includegraphics call used to draw
% it) and using that to override \gothamSideFootnoteRightInset for just
% this one frame -- moving the whole rotated strip to sit immediately left
% of the image instead of at the page's raw edge, entirely automatically,
% with no manual x-coordinate needed from the caller (the direct
% instruction offered a manual-x-coordinate fallback "if it is possible
% [to do automatically], that would be great" -- turned out to be
% possible). \gothamSideImageFrameHeight/\gothamSideImageFrameWidth are
% declared once, globally, and merely *set* (not re-declared) inside the
% environment body, so the frame type can be used more than once per deck;
% the \setlength override of \gothamSideFootnoteRightInset is a plain,
% non-\global assignment, which -- like every other frame-local override in
% this file (e.g. slides/u29-...'s own \setbeamercolor{footnote}{fg=white})
% -- resets on its own once \begin{frame}...\end{frame}'s own TeX group
% closes, with no explicit reset needed.
%
% The extra gap added on top of the image's own measured width (2026-07-06,
% direct instruction: "replicate ... the precise margin we have on regular
% slides between footnote and footer box") was originally an unmeasured
% 4pt guess -- replaced with 3.1pt, back-solved from the deck's own real
% reference gap: measured directly on a plain bottom-footnote slide
% (slides/xhe-Active-vs-Passive-Management-2026-First-current.tex's own
% \source{} line) that the real gap between a footnote's own text and the
% footer box's top edge is 2.93pt, while the original 4pt guess here was
% actually landing the rotated text 3.8pt from the image -- 4pt - (3.8-2.93)
% = 3.13pt, rounded to 3.1pt.
\newlength{\gothamSideImageFrameHeight}
\newlength{\gothamSideImageFrameWidth}
\newenvironment{sideimageframe}[3][t]{%
	\setlength{\gothamSideImageFrameHeight}{\dimexpr\paperheight-\gothamTitleBarBottomAlign-\gothamFooterBoxHeight\relax}%
	\settowidth{\gothamSideImageFrameWidth}{\includegraphics[height=\gothamSideImageFrameHeight, keepaspectratio]{#3}}%
	\sidefootnotes
	\setlength{\gothamSideFootnoteRightInset}{\dimexpr\gothamSideImageFrameWidth+3.1pt\relax}%
	\begin{standardframe}[#1]{#2}%
	\vbox to 0pt{\vss
	\begin{tikzpicture}[overlay, remember picture]
		% outer sep=0pt: see the frametitle hook's comment -- without it the
		% image sits 0.2pt in from the page's right edge (and 0.2pt low).
		\node[anchor=north east, inner sep=0pt, outer sep=0pt] at ($(current page.north east)+(0,-\gothamTitleBarBottomAlign)$)
			{\includegraphics[height=\gothamSideImageFrameHeight, keepaspectratio]{#3}};
	\end{tikzpicture}%
	}%
}{%
	\end{standardframe}%
}

% sectionframe: a section-opening overview (ExampleSlides.pptx slide 3) --
% same layout as standardframe (footer, logo, bullet styling all completely
% unchanged), except the title bar itself is much taller (about a third of
% the page) and set in a bigger font, and the body text is a step bigger
% too, so it reads as "somewhat greater focus" and is unmistakable at a
% glance when flipping through slides -- while still keeping the deck's
% usual yellow-background/black-text look, unlike an earlier version of this
% that inverted to dark-background/yellow-text instead.
%
% The taller bar is NOT a separately-drawn box the way titleframe's is (see
% "Title slide" above): the gotham theme's own frametitle box already sizes
% itself from \gothamFrametitleToppading/\gothamFrametitleBottompading plus
% the title font's own strut height (beamerouterthemegotham.sty, the
% "frametitle" template), so overriding those two lengths locally -- the
% same trick already used below for the font/color overrides -- grows the
% *theme's own* box instead of fighting it with a second one. This also
% means the white-patch-behind-the-logo/chapter-logo machinery (tied to the
% frametitle template via \addtobeamertemplate) keeps working unmodified,
% since it measures the box's actual rendered bottom edge rather than
% assuming any particular height.
%
% The title text is bottom-anchored within the bar -- sitting a fixed "one
% line" above the bar's bottom edge, growing *upward* for a title long
% enough to wrap onto a second line -- rather than top-anchored the way
% gotham's own frametitle template flows by default (which would instead
% leave a large, growing gap *below* a wrapped title, between it and the
% bar's bottom edge). gotham's template has no built-in way to do this: its
% only two knobs are the toppading/bottompading struts, both anchored to
% *one* end of a single top-down paragraph flow. So instead of leaving
% toppading fixed (as an earlier version of this did) the whole title text
% is measured into \gothamSectionframeTitleBox first -- built with the same
% font and (approximately) the same usable width gotham's own template
% gives it, so its \ht+\dp closely track the title's real rendered size,
% including wrapping onto extra lines for a long title -- and toppading is
% then computed as "whatever's left" after that measured height and the
% fixed bottom gap are subtracted from the bar's total target height. A
% longer/wrapped title therefore automatically shrinks toppading (pushing
% the whole title block up, growing upward) rather than changing the gap
% below it. \gothamSectionframeTitleWidth is a slight underestimate of
% gotham's actual per-line width (which is only narrowed by the logo on
% whichever line ends up next to it, typically just the last one) --
% applying that narrower width to *every* line is a deliberate, harmless
% simplification rather than replicating gotham's per-line logic exactly.
%
% The institute logo moves down along with the title now (it's vertically
% centered against the title's own baseline, gotham's instituteLogo
% template), rather than staying pinned where it sits on every other frame
% type, as an earlier version of this kept it -- an intentional change from
% that version: with the title itself now sitting near the bar's bottom
% edge, keeping the logo up at the bar's top would visually separate it
% from the title it's meant to sit beside.
%
% \begingroup/\endgroup (rather than the \global-set-and-reset dance
% \chapterlogo{}/\sidefootnotes need) is enough here: this only ever touches
% \setbeamercolor/\setbeamerfont/\setlength/\setbox, all of which assign
% locally under the hood (checked against beamerbasecolor.sty/
% beamerbasefont.sty directly for the first two -- their key handlers are
% plain \def, not \gdef -- and \setlength/\setbox are plain TeX assignments,
% local by default), so a plain surrounding TeX group reverts them
% correctly on its own.
%
% \gothamSectionframeContentDrop adds a little extra room below the bar
% before the body content starts, on top of standardframe's own
% \gothamContentTopShift pull-up (applied inside \begin{standardframe}
% already, before this code regains control). All lengths here calibrated
% against the rendered PDF's own coordinates; re-measure and adjust if the
% title/body font sizes change.
\newlength{\gothamSectionframeContentDrop}
\setlength{\gothamSectionframeContentDrop}{10pt}
\newlength{\gothamSectionframeBoxHeight}
\setlength{\gothamSectionframeBoxHeight}{85.04pt}
% \gothamSectionframeBottomGap is the *visual* target -- one huge-font line
% height (baselineskip, measured directly from a deliberately 2-line test
% title: the gap between the two lines' own rendered top edges). It is NOT
% used directly as \gothamFrametitleBottompading: gotham's bottom strut
% (\gotham@frametitlestrut@end) consistently produces a visual gap about
% 5.58pt *smaller* than the padding value fed into it (checked empirically,
% same offset for both a one-line and a two-line test title, so it's a
% fixed per-box offset rather than something that scales with content) --
% \gothamSectionframeBottomPadCorrection adds that back.
\newlength{\gothamSectionframeBottomGap}
\setlength{\gothamSectionframeBottomGap}{14.32pt}
\newlength{\gothamSectionframeBottomPadCorrection}
\setlength{\gothamSectionframeBottomPadCorrection}{5.58pt}
% Similarly, the box's total rendered height comes out consistently ~11.4pt
% *less* than toppading+measuredTitleHeight+bottompading would suggest --
% the measured title box (built as a plain \vbox{...\par}) doesn't reproduce
% gotham's own leading strut (\gotham@frametitlestrut@start, a zero-width
% rule with an explicit height rather than natural interline glue) exactly,
% and this correction closes that gap. Checked the same way as the bottom
% one above: consistent across both the one-line and two-line test titles.
\newlength{\gothamSectionframeHeightCorrection}
\setlength{\gothamSectionframeHeightCorrection}{12.1pt}
\newlength{\gothamSectionframeTitleWidth}
\setlength{\gothamSectionframeTitleWidth}{\dimexpr\paperwidth-\gothamFrametitleLeftpading-\gothamLogoWidth-2\gothamFrametitleRightpading\relax}
\newbox\gothamSectionframeTitleBox
\newenvironment{sectionframe}[2][t]{%
	\begingroup
	% Disable hyphenation for the title -- both here, in the measurement
	% box below, and (since this is in effect for the rest of the group,
	% which \begin{standardframe} -- triggering the real \insertframetitle
	% -- runs inside) in the actual rendered title too, so a line break
	% never lands mid-word. Setting the penalties alone is the standard,
	% sufficient way to do this (TeX treats a penalty of 10000 at a
	% breakpoint as forbidden, not just heavily discouraged, so it's never
	% chosen even when the alternative is an overfull line); the hyphenmin
	% lengths are redundant insurance, not strictly needed on top of that.
	\hyphenpenalty=10000
	\exhyphenpenalty=10000
	\lefthyphenmin=100
	\righthyphenmin=100
	\setbeamerfont{frametitle}{size=\huge}
	\setbox\gothamSectionframeTitleBox=\vbox{%
		\hsize=\gothamSectionframeTitleWidth
		\usebeamerfont{frametitle}\usebeamercolor[fg]{frametitle}%
		\noindent #2\par
	}%
	\setlength{\gothamFrametitleBottompading}{\dimexpr\gothamSectionframeBottomGap+\gothamSectionframeBottomPadCorrection\relax}
	\setlength{\gothamFrametitleToppading}{\dimexpr\gothamSectionframeBoxHeight-\gothamFrametitleBottompading-\ht\gothamSectionframeTitleBox-\dp\gothamSectionframeTitleBox+\gothamSectionframeHeightCorrection\relax}
	\global\def\gothamFrametitleFullHeightYellow{1}%
	\global\def\gothamFrametitleLogoFixed{1}%
	\setbeamertemplate{frametitle}[gotham_no_inline_logo]%
	\begin{standardframe}[#1]{#2}%
	\vspace*{\gothamSectionframeContentDrop}%
	\large
}{%
	\end{standardframe}%
	\endgroup
}

% quoteframe: a quote slide (ExampleSlides.pptx slide 4) -- institute logo
% as usual, and the quote itself set large/italic/black on a yellow
% full-bleed block, same style as titleframe/sectionframe's own title boxes
% (flush top-left, stopping short of the logo on the right, same margin as
% elsewhere) rather than the dark block an earlier version of this used. The
% usual footer still shows (unlike titleframe, which drops it -- a title
% slide is a one-off outside any particular course content, while a quote
% is still one slide *among* a deck's others, worth being able to reference
% by its footer slide ID like any other). Opening/closing curly quotes are
% added automatically (write the quote text with none of its own); the
% second argument is the person it's from, set smaller/plain, below the
% yellow box and right-aligned to its own right edge (not the page's usual
% text margin) -- pass an empty second argument for an unattributed quote.
%
%   \quoteframe{One of the funny things about the stock market is that
%   	every time one person buys, another sells, and both think they
%   	are astute.}{William Feather, 1889-1981}
%
% \gothamQuoteBoxHeight is a plain percentage of the page, unlike
% titleframe's/sectionframe's own (calibrated-to-land-on-an-exact-fraction)
% box heights, since nothing here needs to line up with another element's
% measured position the way those do. \gothamTitleframeBoxWidth (defined
% above) is reused as-is for the box's width -- the same
% stop-short-of-the-logo formula, so no reason to duplicate it under a new
% name.
%
% Every piece here -- box, logo, quote text, attribution -- is drawn as a
% zero-space tikz overlay anchored to absolute page coordinates (the same
% \vbox to 0pt{\vss ...} pattern used elsewhere in this file for exactly
% this reason), rather than any of it flowing through the frame's normal
% vertical list. This sidesteps needing to know/counteract beamer's own
% default top-of-frame spacing (which standardframe explicitly measures and
% cancels via \gothamContentTopShift) -- there's no such frame-relative
% positioning to get right here, since literally everything is placed by
% explicit coordinate instead. Like the other page-corner overlays in this
% file, "current page" doesn't resolve correctly on a document's very first
% compile pass, so (like those) this needs the usual second pass too.
%
% A *command*, not an environment -- see the identically-reasoned comment on
% titleframe just above for why: the same beamercolorbox-spanning-a-custom-
% environment's-begin/end problem would apply here too, for the same reason.
\newlength{\gothamQuoteBoxHeight}
\setlength{\gothamQuoteBoxHeight}{0.8\paperheight}
\newlength{\gothamQuoteAttributionGap}
\setlength{\gothamQuoteAttributionGap}{1.8ex}
% \@empty (the empty-attribution check below) needs \makeatletter.
% \quotesource{text} -- opt-in \source{} for \quoteframe (call *before*
% \quoteframe{}{}, same set-then-consume pattern as \chapterlogo{}
% above), needed because \quoteframe is a *command* that auto-closes its
% own \end{frame} -- a bare \source{...} typed after \quoteframe{}{} in
% the calling slide file (as if it were still inside the frame, the way
% \source{} reads everywhere else in this deck) is actually issued
% *outside* any frame, and \source{}'s own \footnote{} then has no page to
% attach to yet. Confirmed directly (2026-07-12, direct instruction:
% traced a spurious extra blank leading overlay step on
% slides/c4c-Macbeth-Spot-Removers-Proposed-Structure-current.tex -- the
% very next multi-overlay frame after
% slides/w8w-Pie-Model-of-Capital-Structure-current.tex's own bare
% \source{} -- down to that orphaned footnote specifically, via a reduced
% two-slide test deck): the dangling footnote insert doesn't render
% anywhere visible on w8w itself, but silently corrupts beamer's own
% overlay-counting pre-pass on whichever frame happens to come right
% after -- reproduced with w8w+c4c alone (bug present) vs. c4c alone (bug
% absent), and confirmed it's specifically about the footnote being
% orphaned, not about w8w's content otherwise, by removing just its
% trailing \source{} line (bug also disappears). A plain \clearpage
% between the two frames also "fixes" c4c but only by giving the orphaned
% footnote a page of its own to attach to -- confirmed directly this
% splits w8w itself into two pages (real content, then a second page with
% nothing but the footnote/footer), not an acceptable trade. The only
% correct fix is keeping the footnote inside the frame that's actually
% citing something, i.e. before \quoteframe's own \end{frame} -- exactly
% what \quotesource{} below does.
\newcommand{\gothamQuoteSource}{}
\newcommand{\quotesource}[1]{\gdef\gothamQuoteSource{#1}}
\makeatletter
\newcommand{\quoteframe}[2]{%
	\begin{frame}[t]%
	% \gothamArtifactBegin/\gothamArtifactEnd (defined with the footer, "PDF
	% accessibility" comment): same reasoning as titleframe's own background
	% box above -- purely decorative, marked so a screen reader skips it. The
	% quote text and attribution below stay unmarked (real content); the
	% footer (this frame keeps one, unlike titleframe) is already marked via
	% the global footline hook.
	\gothamArtifactBegin{/Type /Background}%
	\vbox to 0pt{\vss
	\begin{tikzpicture}[overlay, remember picture]
		% \gothamTitleBoxColor, not a literal color token -- same setting
		% titleframe's own box and the standard frametitle bar (colorA) use
		% (2026-07-16, direct instruction).
		% outer sep=0pt here and on the logo/attribution below: same hairline-
		% inset-from-the-anchored-corner issue as titleframe's own box -- see
		% the frametitle hook's comment for the full story.
		\node[fill=\gothamTitleBoxColor, anchor=north west, inner sep=0pt, outer sep=0pt,
		      minimum width=\gothamTitleframeBoxWidth,
		      minimum height=\gothamQuoteBoxHeight]
			at (current page.north west) {};
	\end{tikzpicture}%
	}%
	\gothamArtifactEnd
	\vbox to 0pt{\vss
	\begin{tikzpicture}[overlay, remember picture]
		% inner ysep derived, not \gothamFrametitleRightpading -- see the
		% frametitle hook's fixed-position logo redraw for why (matches the
		% inline logo's vertically-centered position on every other frame).
		\node[anchor=north east, inner xsep=\gothamFrametitleRightpading,
		      inner ysep=\dimexpr\gothamTitleBarBottomAlign-\gothamLogoHeight\relax,
		      outer sep=0pt]
			at (current page.north east) {\gothamInstituteLogoSquare};
	\end{tikzpicture}%
	}%
	\vbox to 0pt{\vss
	\begin{tikzpicture}[overlay, remember picture]
		% \gothamTitleTextColor, not a literal color token -- follows the box
		% color change right above (2026-07-16, direct instruction). Set
		% *inside* the node's own {...} content group (not a stored box
		% register, unlike titleframe's own vbox above), so this doesn't hit
		% the same "\color leaks out of a stored \vbox" pitfall documented on
		% \titleframe's own closing code -- confirmed directly by checking
		% the frame right after this one for the same symptom.
		\node[anchor=center, inner sep=0pt,
		      text width=\dimexpr\gothamTitleframeBoxWidth-2\gothamFrametitleLeftpading\relax]
			at ($(current page.north west)+(0.5\gothamTitleframeBoxWidth,-0.5\gothamQuoteBoxHeight)$)
			{\color{\gothamTitleTextColor}\Large\itshape\textquotedblleft #1\textquotedblright};
	\end{tikzpicture}%
	}%
	\def\gothamQuoteAttributionTemp{#2}%
	\ifx\gothamQuoteAttributionTemp\@empty\else
		\vbox to 0pt{\vss
		\begin{tikzpicture}[overlay, remember picture]
			\node[anchor=north east, inner sep=0pt, outer sep=0pt]
				at ($(current page.north west)+(\gothamTitleframeBoxWidth,-\gothamQuoteBoxHeight-\gothamQuoteAttributionGap)$)
				{\normalfont\itshape\color{c_custom_gray} #2};
		\end{tikzpicture}%
		}%
	\fi
	\ifx\gothamQuoteSource\@empty\else
		\source{\gothamQuoteSource}%
		\global\let\gothamQuoteSource\@empty
	\fi
	\end{frame}%
}
\makeatother

% figureframe: a slide that is just a figure (ExampleSlides.pptx slide 5) --
% same standardframe layout (title, footer, \source{} for the usual
% "Source: ..."/"Image source: ..." line), but with the figure positioned
% via \figurecenter (call before the figure, inside the frame body) rather
% than left-aligned, which -- like the source deck's own figure slide -- is
% the default:
%
%   \begin{figureframe}{Correlation between DJIA stock returns over 15 years}
%   	\includegraphics[width=0.5\linewidth]{correlation-matrix}
%   	\source{Source: Own illustration, data from Yahoo! Finance, 24.02.2026.}
%   \end{figureframe}
%
%   \begin{figureframe}{Title}
%   	\figurecenter
%   	\includegraphics[width=0.6\linewidth]{diagram}
%   \end{figureframe}
%
% Vertical placement is standardframe's own [t]/[c] (its existing optional
% argument, passed straight through as figureframe's own -- so
% \begin{figureframe}[c]{Title} vertically centers exactly like
% \begin{standardframe}[c]{Title} already does); \figurecenter only ever
% affects the horizontal axis.
%
% \figurecenter is plain \centering under a more discoverable name, nothing
% more -- *not* an opt-in-before-the-frame flag like \chapterlogo{}/
% \sidefootnotes (those must run before \begin{standardframe}, since they
% configure something the frame's own opening code reads). This is called
% *inside* the body instead, exactly where it takes effect, so it can just
% take effect immediately like any other formatting declaration -- no flag,
% no state to reset afterwards, and (checked by first writing it as a flag
% read in figureframe's own opening code, mirroring \sidefootnotes) no way
% for the frame's opening code to see it anyway, since that code runs
% before the body -- and \figurecenter -- does.
\newcommand{\figurecenter}{\centering}
\newenvironment{figureframe}[2][t]{%
	\begin{standardframe}[#1]{#2}%
}{%
	\end{standardframe}%
}

% highlightframe: a standardframe with two things about it made unmistakably
% different -- the title bar's *and* the footer's background color, both
% together -- so a presenter can never flip to it without immediately
% noticing (ExampleSlides.pptx slide 6 was the inspiration for wanting a
% standardframe variant here, though the source slide itself has no such
% color cue; this is a new addition, not a conversion). Everything else --
% body, logo, bullet styling -- is exactly standardframe's own, via the same
% \begingroup/\endgroup wrapper sectionframe above uses for the title-color
% half of this. \gothamHighlightColor is a \colorlet, not a color baked
% directly into the environment, so retargeting every highlightframe in a
% deck to a different color -- title bar and footer both -- is still the
% same one-line change it was when this only touched the title bar.
%
% The footer box needs the \global\def dance (\gothamFooterBoxColor, defined
% above with the frametitle hook's other opt-in flags, reset in
% standardframe's own closing code) that \setbeamercolor{frametitle} doesn't:
% that assigns through beamer's own color-management system, which handles
% the frame-closing group correctly on its own (checked -- see
% \gothamChapterLogoPath's comment for why plain macros need \global\def but
% beamer's own set-commands don't), but the footer box's fill is read from a
% plain macro inside a tikz key, same category as \gothamChapterLogoPath.
\colorlet{gothamHighlightColor}{c_primary_gray}
\newenvironment{highlightframe}[2][t]{%
	\begingroup
	\setbeamercolor{frametitle}{bg=gothamHighlightColor}
	\global\def\gothamFooterBoxColor{gothamHighlightColor}%
	\begin{standardframe}[#1]{#2}%
}{%
	\end{standardframe}%
	\endgroup
}



%   \FooterText{...}
%   \makeatletter\ifdefined\HANDOUT\studentversion\fi\makeatother  % or similar
%   \input{course-content.tex}   % the shared \slide{}/\omitslide{} list
%
% with two thin wrapper decks (one plain, one calling \studentversion first)
% both \input-ing the same shared content file -- see README.md.
\newif\ifgothamStudentVersion
% \makeatletter must wrap the \newcommand itself, not sit *inside* the macro
% body: catcodes apply at tokenization time, which for a macro body happens
% once, immediately, when \newcommand reads it -- not deferred until the
% macro is later called. \makeatletter placed inside the body (tried first)
% was too late: \beamer@currentmode had already been mis-tokenized as
% "\beamer" + the six-way-otherwise-ordinary character "@" + eleven plain
% letters "currentmode", by the time that inner \makeatletter could act on
% it, so the \gdef silently created/redefined the wrong thing and never
% touched beamer's actual mode variable. Confirmed via \meaning\beamer@
% currentmode before/after \studentversion -- unchanged with the inside-the-
% body version, fixed once \makeatletter moved out here.
\makeatletter
\newcommand{\studentversion}{%
	\global\gothamStudentVersiontrue
	\gdef\beamer@currentmode{handout}%
}
\makeatother

% Deck version system - new boolean flags and string for current version
\newif\ifgothamInstructorVersion
\newif\ifgothamDuringCourseVersion
\newif\ifgothamAfterCourseVersion
\gdef\gothamDeckVersion{instructor}
\gdef\gothamDeckVersionText{}

% \deckversion[] - set the deck version(s)
% Usage: \deckversion[instructor], \deckversion[during-course], etc.
% Multiple versions: \deckversion[instructor,during-course,after-course]
% Parses #1 as a comma-separated list via etoolbox's \forcsvlist (rather
% than the single \ifstrequal{whole-list}{"instructor"} check this used to
% do, which only ever matched a single-version list -- a multi-version
% \deckversion[instructor,during-course] silently set none of the three
% booleans, since the compared string was the literal, unsplit
% "instructor,during-course"). Note \forcsvlist, not etoolbox's similarly-
% named \forlistloop: \forlistloop iterates etoolbox's own internal
% pipe-delimited list format (built via \listadd), not a plain comma string
% like this command's own argument -- confirmed directly (isolated test),
% passing a raw comma string to it doesn't error or split on the commas, it
% just treats the whole string as a single one-item "list".
\makeatletter
\newcommand{\deckversion}[1][instructor]{%
	% \xdef (not \gdef): #1 needs to be fully expanded here, not just
	% stored as whatever token(s) were passed -- \deckversion[\myvar]
	% (compile-all-versions.bat's own override mechanism, via a wrapper
	% file's \def\deckversionoverride{...}) passes a macro, not a literal
	% string, and everywhere else in this deck-version system that reads
	% \gothamDeckVersion only ever applies one \expandafter to it (see
	% \ifdeckversioninlist's own comment) -- enough to unwrap \gothamDeckVersion
	% itself, but not a further macro still sitting inside it. Confirmed
	% directly: with the old \gdef, \deckversion[\deckversionoverride]
	% left \gothamDeckVersionText empty and every \ifdeckversioninlist
	% check false, since \gothamDeckVersion's own stored content was the
	% single unexpanded token \deckversionoverride, not "during-course".
	\xdef\gothamDeckVersion{#1}%
	\global\gothamInstructorVersionfalse
	\global\gothamDuringCourseVersionfalse
	\global\gothamAfterCourseVersionfalse
	\gdef\gothamDeckVersionText{}%
	\def\do@deckversions##1{%
		\ifstrequal{##1}{instructor}{\global\gothamInstructorVersiontrue\gdef\gothamDeckVersionText{instructor deck}}{}%
		\ifstrequal{##1}{during-course}{\global\gothamDuringCourseVersiontrue\gdef\gothamDeckVersionText{student course deck}}{}%
		\ifstrequal{##1}{after-course}{\global\gothamAfterCourseVersiontrue\gdef\gothamDeckVersionText{student final deck}}{}%
	}%
	% \forcsvlist itself splits on literal commas in its own raw,
	% unexpanded argument (confirmed directly against etoolbox's own
	% source, \etb@defparser@do@aux) -- passing #1 straight through here
	% has the exact same problem as the \gdef above did, one call site
	% later. Building the whole \forcsvlist{...}{...} call inside an
	% \edef (protecting \forcsvlist/\do@deckversions themselves with
	% \noexpand so they're invoked *later*, once this constructed call is
	% run, not expanded away here) expands only \gothamDeckVersion --
	% already fully resolved by the \xdef above -- into the list argument.
	\edef\gotham@deckversion@call{\noexpand\forcsvlist{\noexpand\do@deckversions}{\gothamDeckVersion}}%
	\gotham@deckversion@call
}
\makeatother

\newif\ifgothamOmittedSlide
\newcommand{\gothamOmittedSlideMarker}{ | omitted from student deck}
% \stepcounter{framenumber} in the skipped branch below keeps the student
% deck's slide numbers matching the instructor deck's, gap and all, so "see
% slide 12" means the same physical slide to both audiences: beamer steps
% "framenumber" exactly once per \begin{frame} (checked in
% beamerbaseframe.sty -- once per *frame*, not per overlay page, so this is
% right regardless of how many overlay steps the omitted slide would have
% had), and this is the one place that increment is skipped along with the
% frame itself, so it has to be replayed by hand.
\newcommand{\omitslide}[1]{%
	\ifgothamStudentVersion
		\stepcounter{framenumber}%
	\else
		\global\gothamOmittedSlidetrue
		\setSlideID{#1}%
		\input{slides/#1-current.tex}%
		\global\gothamOmittedSlidefalse
	\fi
}
\newcommand{\omitslideversion}[2]{%
	\ifgothamStudentVersion
		\stepcounter{framenumber}%
	\else
		\global\gothamOmittedSlidetrue
		\setSlideID{#1}%
		\input{slides/#1-v#2.tex}%
		\global\gothamOmittedSlidefalse
	\fi
}

% Deck version-aware overlay commands
% Syntax: \deckonly<overlay>[versions]{content}
% If [versions] omitted, defaults to all decks (instructor,during-course,after-course)
%
% Built on etoolbox's own \forcsvlist for the comma-splitting (rather than a
% hand-rolled parser: an earlier version tried to split a list at each comma
% using \newcommand's plain, non-delimited argument grabbing
% (\newcommand{\ifdeckversioninlist@split}[3]{...} called as
% \ifdeckversioninlist@split#1\@nil{#1}) -- but a plain \newcommand doesn't do
% delimited-parameter matching against \@nil the way a \def with a matching
% pattern would; it just grabs the next 3 tokens/groups one at a time. So
% e.g. "instructor" silently split into single characters ("i", "n", "s",
% ...) instead of being read as one list item, which eventually ran the
% parser off the end of its own input and surfaced as "Undefined control
% sequence" wherever the first \conditionalpause[...]/\deckonly[...] in a
% deck happened to be used -- confirmed directly against the compile log,
% traced to slide y5z-CoCo-Bonds's own \conditionalpause[instructor].
%
% NOT etoolbox's own \ifinlist{item}{list}, despite the name similarity: that
% one tests membership in etoolbox's own internal list format, built only via
% \listadd (items joined with a literal "|", not ","); fed a plain comma
% string instead, confirmed directly (isolated test) to silently return
% false for every item, including an exact single-item match.
%
% The membership test itself needs one more subtlety: #1 here is the deck
% version to search for -- always \gothamDeckVersion at every call site
% below, i.e. a macro, not a literal string -- so comparing it via a plain
% \ifstrequal{##1}{#1} does not work either: etoolbox's \ifstrequal compares
% via \detokenize, which (confirmed directly, isolated test) stringifies an
% unexpanded macro token to its own control-sequence *name* ("gothamDeckVersion"),
% not the current value that macro holds. \expandafter forces that expansion
% before \ifstrequal ever sees it.
\makeatletter

\let\gotham@original@only\only
\let\gotham@original@uncover\uncover
\let\gotham@original@visible\visible

\newif\ifgotham@deckversion@found

% Usage: \ifdeckversioninlist{version}{list}{true-branch}{false-branch}
\newcommand{\ifdeckversioninlist}[2]{%
	\global\gotham@deckversion@foundfalse
	\def\gotham@deckversion@check##1{%
		\expandafter\ifstrequal\expandafter{#1}{##1}{\global\gotham@deckversion@foundtrue}{}%
	}%
	\forcsvlist{\gotham@deckversion@check}{#2}%
	\ifgotham@deckversion@found
		\expandafter\@firstoftwo
	\else
		\expandafter\@secondoftwo
	\fi
}

% \@ifnextchar-based parsing (the same idiom LaTeX2e's own kernel commands
% use for an optional argument that isn't a plain trailing [...]) so the
% overlay spec reads exactly like beamer's own \only<...>, with the deck-
% version list as a further, genuinely optional [...] afterwards:
%   \deckonly<2>{content}                              visible in all versions
%   \deckonly<2>[instructor]{content}                  instructor only
%   \deckonly<2>[during-course,after-course]{content}  during + after
\newcommand{\deckonly}{\@ifnextchar<{\gotham@deckonly@overlay}{\gotham@deckonly@overlay<1->}}
\def\gotham@deckonly@overlay<#1>{\@ifnextchar[{\gotham@deckonly@versions<#1>}{\gotham@deckonly@versions<#1>[instructor,during-course,after-course]}}
\def\gotham@deckonly@versions<#1>[#2]#3{%
	\ifgothamInstructorVersion
		\gotham@original@only<#1>{#3}%
	\else
		\ifdeckversioninlist{\gothamDeckVersion}{#2}{\gotham@original@only<#1>{#3}}{}%
	\fi
}

% Same interface as \deckonly, wrapping \uncover and \visible respectively
\newcommand{\deckuncover}{\@ifnextchar<{\gotham@deckuncover@overlay}{\gotham@deckuncover@overlay<1->}}
\def\gotham@deckuncover@overlay<#1>{\@ifnextchar[{\gotham@deckuncover@versions<#1>}{\gotham@deckuncover@versions<#1>[instructor,during-course,after-course]}}
\def\gotham@deckuncover@versions<#1>[#2]#3{%
	\ifgothamInstructorVersion
		\gotham@original@uncover<#1>{#3}%
	\else
		\ifdeckversioninlist{\gothamDeckVersion}{#2}{\gotham@original@uncover<#1>{#3}}{}%
	\fi
}

\newcommand{\deckvisible}{\@ifnextchar<{\gotham@deckvisible@overlay}{\gotham@deckvisible@overlay<1->}}
\def\gotham@deckvisible@overlay<#1>{\@ifnextchar[{\gotham@deckvisible@versions<#1>}{\gotham@deckvisible@versions<#1>[instructor,during-course,after-course]}}
\def\gotham@deckvisible@versions<#1>[#2]#3{%
	\ifgothamInstructorVersion
		\gotham@original@visible<#1>{#3}%
	\else
		\ifdeckversioninlist{\gothamDeckVersion}{#2}{\gotham@original@visible<#1>{#3}}{}%
	\fi
}

% \conditionalpause[versions] - \pause, but only takes effect in the listed
% deck versions (default: instructor only); a no-op, not even a skipped
% overlay step, in any other version
% Usage: \conditionalpause[instructor,during-course] or \conditionalpause[after-course]
\newcommand{\conditionalpause}[1][instructor]{%
	\ifdeckversioninlist{\gothamDeckVersion}{#1}{\pause}{}%
}

% \deckreveal[step]{content} - the "hidden answer, click-revealed for the
% instructor, absent for during-course, pre-shown for after-course" pattern
% already documented by hand on several exercise slides (e.g. mbz/x7q/37s/
% 3cv's own header comments: "a slide/reveal-step that is hidden in the
% source PowerPoint should, by default, only be made visible on the
% instructor version"). A bare, unconditional \only<N->{content} looks like
% it does this, but doesn't: \only's overlay range registers with beamer's
% own overlay-counting pass regardless of \ifgothamInstructorVersion/
% \conditionalpause -- neither of those touch \only itself, they only guard
% *other* content around it -- so a frame with nothing else establishing a
% second overlay step still gets a real second slide out of a bare
% \only<2->, in *every* deck version, not just the instructor one. Confirmed
% directly (2026-07-21, direct instruction) on 37s/3cv: both a) leaked the
% "hidden" content to during-course as a genuine second slide, the exact
% thing this pattern exists to prevent, and b) showed a blank second bullet/
% cell first on that leaked slide's own step 1 -- \only<2-> hides the
% *content*, not the space it would occupy, so during-course's own step 1
% still had nothing else deciding where the bullet/cell landed.
% \deckreveal fixes both by picking a genuinely different rendering per
% version rather than reusing the same overlay-restricted \only call
% everywhere: instructor gets the real overlay build (\only<step->,
% default step 2, matching \tallitem's own default-multiplier convention);
% during-course gets nothing at all (no token emitted, so beamer's overlay
% pass never sees an overlay spec to register); after-course gets the
% content unconditionally, with no overlay spec of its own either, so it
% doesn't silently force a second slide there either. Content-level only
% (wraps a value/expression, like \only itself) -- doesn't cover an
% itemize \item's own bullet marker, which needs its own overlay spec
% (`\item<step->`) directly on the \item, not wrapped from outside; see
% 37s's own comment for that variant, hand-written rather than factored
% into a second macro since only one call site needs it so far.
\newcommand{\deckreveal}[2][2]{%
	\ifgothamDuringCourseVersion
	\else
		\ifgothamInstructorVersion
			\only<#1->{#2}%
		\else
			#2%
		\fi
	\fi
}

% \conditionalslide[versions]{ID-slug} - \slide{}, but only included in the
% listed deck versions; keeps the frame counter in step with the omitted
% version otherwise (same \stepcounter{framenumber} technique \omitslide{}
% already uses for the student deck), so "see slide N" still points at the
% same physical slide across every deck version.
% Usage: \conditionalslide[instructor,after-course]{xyz-Answer-key}
\newcommand{\conditionalslide}[2][instructor,during-course,after-course]{%
	\ifdeckversioninlist{\gothamDeckVersion}{#1}%
		{\setSlideID{#2}%
		 \input{slides/#2-current.tex}}%
		{\stepcounter{framenumber}}%
}

\makeatother

% PDF accessibility: mark purely decorative/running content (footer, below;
% also reused on titleframe's/quoteframe's own decorative background boxes;
% logo not yet included, see README.md's own accessibility section for that
% open question) as a PDF "Artifact", so screen readers skip it -- via
% LuaTeX's raw \pdfextension literal
% primitive. \ifluatex-guarded (iftex is already loaded, via beamer.cls
% itself) because \pdfextension is LuaTeX-only -- this project's README
% documents xelatex as a supported alternative engine, and compiling with it
% hit exactly this: "Undefined control sequence: \pdfextension", hard-failing
% every frame, since the first version of this didn't guard the primitive at
% all. Under any non-LuaTeX engine these two commands now silently do
% nothing instead: decorative content just isn't marked as an Artifact under
% xelatex/pdflatex, rather than breaking the compile -- a real loss of
% accessibility under those engines specifically, not a full equivalent fix,
% but not something to solve today (see also: compiling with plain pdflatex
% already silently mis-renders the whole deck's fonts, so xelatex is really
% the only *other* engine this needs to not break under; pdflatex was never
% viable here to begin with, per "Compiling" in README.md).
% Deliberately NOT tagpdf/kernel tagging (\DocumentMetadata):
% beamer 3.77 refuses to compile at all when that's active ("The beamer
% class is not compatible with \DocumentMetadata ... switch to the new
% ltx-talk class if you require tagged output"), and ltx-talk is a from-
% scratch class (not beamer-based, no theme/template system) that this
% framework can't realistically migrate to. Bypassing beamer's check by
% force-loading the underlying package before \documentclass also failed --
% it gets past hyperref but then hits a genuine internal error in tagpdf's
% LuaTeX backend from the non-standard load order, not something to build
% real coursework on.
%
% /Artifact is a standalone PDF marked-content tag (PDF spec 14.8.2.2):
% valid in *any* PDF, with or without a structure tree, no /Properties
% resource entry or MCID required -- unlike structure-tree tagging, this
% doesn't require kernel tagging support at all. It's a real, working
% partial fix, not full PDF/UA conformance: there's no structure tree, so
% no true semantic headings, and reading order elsewhere on the slide still
% depends on plain content-stream order (verified, not enforced, for
% standardframe -- see its own comment).
% #1 is the artifact's inline PDF dictionary content, e.g.
% "/Type /Pagination /Subtype /Footer" for a running footer, or
% "/Type /Background" for a decorative fill -- not hardcoded to one /Type,
% since not everything decorative here is "Pagination" in the PDF spec's own
% sense of that word (running headers/footers/page numbers/watermarks);
% titleframe's/quoteframe's plain background color fills are semantically
% "Background" artifacts instead, a distinct standard type.
\newcommand{\gothamArtifactBegin}[1]{%
	\ifluatex\pdfextension literal direct{/Artifact <<#1>> BDC}\fi
}
\newcommand{\gothamArtifactEnd}{\ifluatex\pdfextension literal direct{EMC}\fi}

% Footer
% Smaller footer + a bit of padding under it
\setbeamerfont{frame footer}{size=\fontsize{5}{3}\selectfont} % {font size}{baseline skip} in pt

% Shrinks the gap between footnotes and the footer box below them. The
% theme positions the footer's content via \vspace*{\gothamFootlineVOffset}
% (a fixed, negative length) before it starts; making that more negative
% turns out to give footnotes (which are placed just before the footer, and
% stretch to fill whatever room remains via beamer's own
% \beamer@framebottomskip) more room to extend into, pulling them down
% towards the footer -- checked empirically against the PDF's own
% coordinates, this does *not* move the footer's own text at all (that
% still sits exactly where \gothamFooterBoxHeight below was calibrated for),
% so the footer box itself needs no compensating adjustment.
%
% Partly given back (+1.4pt) 2026-07-05: a footnote/source containing an
% underlined \href (see the hyperlink-styling comment above) sits close
% enough to the footer at the original -1.3pt that the underline visually
% touches the footer box (checked directly against the compiled PDF's own
% pixel coordinates: a plain-text footnote's own gap to the footer measured
% ~1.9pt, but an underlined-link footnote's gap measured only ~0.6pt at
% the same setting, since the strut rebuilt for footnotes just below --
% correctly sized for plain text's own natural depth -- doesn't account
% for the underline rule extending further down). +1.4pt splits the
% difference: a linked footnote's gap now measures close to the plain-
% footnote reference (~1.9pt) instead of colliding with the footer, and a
% plain footnote's own gap just becomes a little more relaxed (~3.3pt)
% rather than stretched uncomfortably tight -- applied globally rather
% than only to footnotes with links, since there's no clean per-footnote
% hook for it and the plain-footnote side effect is minor.
%
% Nudged back down by 0.4pt the next day (direct instruction: move
% footnotes down "by the width of the underline line") -- the +1.4pt above
% left a bit more headroom than necessary; 0.4pt is the underline rule's
% own measured stroke thickness (a solid ~1px core at 600dpi, i.e.
% 1/600in \approx 0.12pt, plus antialiasing bleed either side consistent
% with the common 0.4pt LaTeX default rule width -- measured directly
% against the compiled PDF's own pixel values, not assumed).
%
% Reverted the very next round (direct instruction: move back *up* by that
% same underline thickness) -- net effect is back to the plain +1.4pt this
% section started with; kept as two separate historical adjustments below
% rather than collapsing them into one, so this back-and-forth stays
% visible rather than silently disappearing.
\addtolength{\gothamFootlineVOffset}{-1.3pt}
\addtolength{\gothamFootlineVOffset}{1.4pt}
\addtolength{\gothamFootlineVOffset}{-0.4pt}
\addtolength{\gothamFootlineVOffset}{0.4pt}

% 2026-07-06 (direct instruction): move footnotes down by half the
% measured distance between a linked footnote's own underline and the
% footer box, checked on slides/3ic-Corporate-Venture-Capital-
% current.tex specifically (the underline sits at y=245.95pt, the footer
% box's own top edge at y=247.69pt, a 1.74pt gap) -- half of that is
% 0.87pt, applied as a further negative adjustment (more negative =
% footnotes get pulled further down, per this section's own comment
% above).
\addtolength{\gothamFootlineVOffset}{-0.87pt}

% Shifts the right end of the footer line (the frame number) so it lines up
% with the logo's own right edge above it, rather than wherever the theme's
% default footline margins would otherwise put it. Empirically calibrated
% against the rendered PDF coordinates of both; re-measure and adjust if
% \gothamFrametitleRightpading (the logo's own right inset) changes.
\newlength{\gothamFooterRightShift}
\setlength{\gothamFooterRightShift}{4.6pt}

% Nudges the footer text down within its own yellow box, to vertically
% center it there: it previously sat noticeably closer to the box's top
% edge than its bottom (0.14pt above vs. 1.12pt below, checked against the
% rendered PDF's own text coordinates) since \gothamFootlineVOffset's own
% positioning (further up) turns out not to move this text at all (see its
% comment above) -- this is a separate, additive leading \vspace* instead.
\newlength{\gothamFooterTextDrop}
\setlength{\gothamFooterTextDrop}{-0.49pt}

% Compile timestamp: lets you tell, from a PDF alone (no need to check file
% properties/mtime, which don't survive e.g. emailing or uploading a copy),
% whether it reflects the latest edits or a stale earlier build -- handy
% whenever a PDF's *content* is in question, not just when debugging compile
% tooling itself. LuaTeX-only, matching this project's required engine (see
% README.md's "Compiling" section) -- no \ifluatex guard needed since
% nothing else here works under another engine to begin with.
%
% Deliberately avoids os.date's own %-format directives (e.g. "%Y-%m-%d")
% entirely, building the string from os.date("*t")'s field table and plain
% Lua concatenation instead -- \directlua's argument is tokenized by plain
% TeX first, where % starts a comment, before the text ever reaches Lua as a
% string. The obvious-looking fix, escaping each one as \%, does *not* work:
% confirmed directly, it produces a Lua string containing a literal
% backslash followed by a percent sign (detokenization keeps a control
% sequence's own backslash, for \% same as any other), which Lua's own
% string-literal parser then rejects outright ("invalid escape sequence"),
% silently aborting that whole \directlua call -- with no visible LaTeX-side
% error, just a cache file ending up written with an empty value. Switching
% catcode temporarily (\catcode`\%=12 ... ) would also work but needs a
% \begingroup/\endgroup wrapper at every call site; skipping % entirely is
% simpler and needs no such wrapper.
%
% Cached to \jobname.timestamp (via plain TeX \write, so MiKTeX's
% -aux-directory flag -- already applied through .latexmkrc's %O -- routes
% it into .build/ like every other aux file, not left cluttering decks/ or
% the project root) rather than recomputed fresh on every pass: a live
% os.date() call on every pass, tried first, broke latexmk's own rerun
% convergence outright -- confirmed directly, by isolating the change, that
% an otherwise-identical clean compile converges normally without it. Each
% of latexmk's internal reruns relaunches lualatex as a brand-new process,
% and a full pass over a deck this size (font loading alone) comfortably
% crosses a one-minute boundary, so two consecutive passes embedded two
% different strings in the footer; latexmk saw the .aux checksum differ
% *again* on the very next pass and looped ("Latexmk: References changed.")
% until hitting its rerun limit and reporting a false "gave an error", with
% no actual LaTeX error anywhere in the log. Caching to a file read back on
% later passes fixes that: the value is fixed the moment it's first
% written, so the footer's rendered text -- and therefore the .aux
% checksum -- stays identical across every subsequent pass of the same
% compile.
%
% This cache is deliberately *not* self-invalidating by date/content -- it
% keeps reporting an old timestamp indefinitely once written, across
% unrelated future compiles, unless something clears it first.
% SlideManager's compiler.py does exactly that: before every job it deletes
% that job's own .build/<jobname>.* (which includes <jobname>.timestamp), so
% a run through SlideManager always gets a fresh stamp. An ad-hoc TeXstudio
% compile without going through SlideManager first may show a stale stamp if
% .build/ already has one left over from an earlier session -- delete
% .build/<jobname>.timestamp by hand (or just compile that job once through
% SlideManager) if that matters for what you're checking.
\directlua{
	function gothamZeroPad(n) if n < 10 then return "0" .. n else return "" .. n end end
}
\newwrite\gothamCompileTimestampFile
\newcommand{\gothamCompileTimestamp}{%
	\ifdefined\gothamCompileTimestampValue
	\else
		\IfFileExists{\jobname.timestamp}{%
			\input{\jobname.timestamp}%
		}{%
			\xdef\gothamCompileTimestampValue{\directlua{%
				local d = os.date("*t")%
				tex.sprint(d.year .. "-" .. gothamZeroPad(d.month) .. "-" .. gothamZeroPad(d.day) .. " " .. gothamZeroPad(d.hour) .. ":" .. gothamZeroPad(d.min))%
			}}%
			\immediate\openout\gothamCompileTimestampFile=\jobname.timestamp\relax
			\immediate\write\gothamCompileTimestampFile{\string\gdef\string\gothamCompileTimestampValue{\gothamCompileTimestampValue}}%
			\immediate\closeout\gothamCompileTimestampFile
		}%
	\fi
	\gothamCompileTimestampValue
}

% \gothamHideFooterText: opt-in flag suppressing just \FooterText (the left
% part of the footer) -- default false. Set true by the standout template's
% own hook below, so `\begin{frame}[standout]` hides only that left part,
% keeping the "(slide id: ...)" block and frame number visible (2026-07-16,
% direct instruction). Reset immediately after being read, inside the
% footer template itself, rather than relying on a per-environment closing
% hook (standout frames are bare `\begin{frame}[standout]...\end{frame}`,
% not one of this project's own environments like standardframe, so there's
% no custom closing code to reset it in) -- safe because \gothamHideFooterText
% is only ever consumed here, once, immediately before being cleared.
\newif\ifgothamHideFooterText
% gotham's own \defbeamertemplate*{standout}{gotham} (gotham.dtx) is what
% `\begin{frame}[standout]` actually activates; \addtobeamertemplate's
% "after" hook appends to it without needing to touch that upstream file.
\addtobeamertemplate{standout}{}{\global\gothamHideFooterTexttrue}

\setbeamertemplate{frame footer}{%
	\vspace*{\gothamFooterTextDrop}%
	\ifgothamHideFooterText
		\global\gothamHideFooterTextfalse
	\else
		\FooterText%
	\fi
	\hfill%
	{\fontsize{4}{4}\selectfont (slide id: \slideID\ifgothamOmittedSlide\gothamOmittedSlideMarker\fi\ | \gothamDeckVersionText\ | \gothamCompileTimestamp)}\hspace*{2em}\insertframenumber%
	\hspace*{-4em}%
	\hspace*{-\gothamFooterRightShift}%
	\vspace*{0.4ex}%
}

% Give the footer the same yellow-box treatment as the title bar, but --
% unlike the title bar, which is tall for visual weight -- only as tall as
% the footer's own text plus equal padding above and below.
%
% This paints a yellow patch *behind* the theme's own footline rather than
% replacing that template's content, because \usebeamertemplate{frame~footer}
% (which is what actually typesets \FooterText etc.) only produces its
% content when invoked from the theme's own original footline template --
% calling it again from anywhere else (a redefined footline, the frame body,
% a \settoheight measurement, ...) silently produces nothing. It's
% positioned from "current page.south" (the page's own bottom-left corner)
% rather than a \tikz "remember picture" node tracking the footer's actual
% position, since there's no per-frame content to track: the footer text
% (in its actual 5pt font, per \setbeamerfont{frame footer} above -- despite
% the "page number in head/foot" font switch inside the theme's footline
% template, this is the font that actually ends up applied) sits at a fixed
% distance from the bottom of the page on every frame.
%
% \gothamFooterBoxHeight below is calibrated to that actual measured
% position (checked via the PDF's own text coordinates, not just visually):
% the text sits about 0.84pt above the page's bottom edge and spans about
% one em of its own 5pt font, plus a bit of padding above and below. If the
% footer font size changes, this will need re-measuring.
\newlength{\gothamFooterBoxHeight}
\setlength{\gothamFooterBoxHeight}{7.26pt}
% \vbox to 0pt{...} forces this to occupy exactly zero space in the normal
% vertical list, *without* switching into horizontal mode the way an hbox
% (e.g. \raisebox) would -- inserting an hbox here forces LaTeX to open an
% implicit paragraph before the theme's own \vspace*{\gothamFootlineVOffset}
% runs, which changes how that \vspace behaves (it special-cases being at
% the very start of vertical mode) and nudges the footer text down by a few
% points. A \vbox is already vertical material, so it doesn't trigger that.
% \gothamArtifactBegin/\gothamArtifactEnd (defined above, "PDF accessibility"
% comment) wrap both this box *and* the original footline template's own
% content (the theme's \usebeamertemplate{frame footer} call, i.e.
% \FooterText/slide id/page number) -- \addtobeamertemplate's "after" slot
% runs once the wrapped-around original template content has finished, so
% \gothamArtifactEnd there closes the same marked-content span this begins.
% The footer box's own south-west corner originally measured 0.2pt short
% of the page's actual left/bottom edges on every slide checked
% (2026-07-06, direct instruction: "check whether the footer ... has a bit
% of white on its left and below") -- confirmed directly against the
% compiled PDF's own drawn rectangle (consistently ~0.199pt in from both
% edges, not a one-off), unlike the title bar (gotham's own native
% frametitle background, a completely different mechanism -- normal TeX
% box placement, not a tikz "current page" lookup), which sits exactly
% flush at (0,0). At the time, the root cause wasn't chased further and
% the box was fixed by hand-nudging the anchor 0.2pt past the corner and
% padding the size to match. Root cause since found (2026-07-16, while
% chasing the identical hairline on titleframe/quoteframe -- see the
% frametitle hook's own comment above): tikz's default outer sep
% (0.5\pgflinewidth = 0.2pt) puts every edge/corner anchor 0.2pt outside
% the filled rectangle, on *every* edge-anchored node, so the earlier
% "nothing else anchored to a page corner shows the same offset" turned
% out to be wrong. "outer sep=0pt" below is the real fix, replacing the
% nudge-and-pad workaround; the box now lands exactly flush by
% construction rather than by measured compensation.
\addtobeamertemplate{footline}{%
	\gothamArtifactBegin{/Type /Pagination /Subtype /Footer}%
	\vbox to 0pt{%
	\vss
	\begin{tikzpicture}[remember picture, overlay]
		\node[fill=\gothamFooterBoxColor, anchor=south west, inner sep=0pt, outer sep=0pt,
		      minimum width=\paperwidth, minimum height=\gothamFooterBoxHeight]
			at (current page.south west) {};
	\end{tikzpicture}%
	}%
}{\gothamArtifactEnd}

% Footnotes: small font (smaller than the footer's own 5pt), symbol markers
% (*, dagger, ...) instead of arabic numbers, right-aligned (flush with the
% logo's own right edge, like the footer's frame number), no rule above
% them, and tightened-up line spacing.
%
% Line spacing needs two different levers, not just \baselineskip, because
% at this font size (5pt) \baselineskip's normal value is *smaller* than the
% font's own natural line height, so TeX falls back to \lineskip (a flat
% add-on) between every line regardless -- \baselineskip itself turns out to
% have no effect here at all. That fallback \lineskip is what governs
% spacing *within* a single (wrapped) footnote's own lines: reduced only a
% tiny amount, from the default 1pt.
% The *between*-footnote gap is inflated by a separate mechanism entirely:
% \footnotesep, a rule LaTeX inserts before every individual footnote entry
% (measured to be 7.7pt here -- clearly sized for a much bigger font than
% ours), so that's what actually needs the bigger reduction; \parskip
% (tried first) turned out to have no visible effect on it.
\setlength{\footnotesep}{3pt}
% \reset@font\footnotesize is applied by beamer's own footnote machinery
% before this template runs (beamerbaseframecomponents.sty), so the font
% size has to be set again here to actually take effect.
\renewcommand{\footnoterule}{}
% Same marker sequence as LaTeX's own \fnsymbol (*, dagger, double dagger,
% ...), but built from plain Unicode text characters instead: \fnsymbol's
% symbols are typeset via \ensuremath{...}, i.e. in *math* mode, and the
% theme loads fontspec with the "no-math" option (see
% beamerfontthemegotham.sty) specifically so that Fira Sans -- the font
% fontspec puts everything else in -- is left in charge of text only, not
% math; math mode therefore still falls back to LaTeX's default Computer
% Modern math fonts. \fnsymbol's markers were consequently rendering in a
% different font than the surrounding text (checked by comparing rendered
% glyph shapes) -- most visible as a marker that looked shifted left of
% flush-left text next to it (e.g. a \source{} line, which has no marker of
% its own), since Computer Modern's symbol glyphs have different sidebearings
% than Fira Sans's. Plain characters typeset with no \ensuremath wrapper stay
% in text mode and pick up whatever font is already active at the call site,
% i.e. Fira Sans -- matching the rest of the footnote.
% Literal Unicode characters, not the \textdagger/\textdaggerdbl/\textbardbl
% kernel commands: those need the textcomp/TS1 setup to render as plain text,
% and silently fall back to \ensuremath{...} (i.e. math mode, the very thing
% this replaces \fnsymbol to avoid -- see the comment above) when that isn't
% loaded. LuaLaTeX plus fontspec maps raw UTF-8 input straight to the current
% font's (Fira Sans's) own glyphs with no encoding layer in between, which is
% the entire point of using them -- so writing the characters directly, with
% no wrapper command, guarantees they stay in text mode and stay in Fira Sans.
%
% The sequence itself is user-configurable via \setfootnotesymbols{...} (a
% comma-separated list, e.g. \setfootnotesymbols{*,+,#}) -- call it anywhere
% from that point in the deck onward (e.g. once, right after
% % Slidemaster preamble variant for University of Graz decks. Institution-
% specific content only (colors, logo, logo calibration) -- every frame
% type, the deck-version system, footer/hyperlink/footnote styling, tikz
% styles, etc. live in Slidemaster_preamble_base.tex, \input{}ed at the
% bottom of this file, and are shared verbatim across every institution
% this project supports (see Slidemaster_preamble_mci.tex and
% Slidemaster_preamble_uni-ibk.tex for the other two). A deck picks its
% institution by \input{}ing this file instead of one of those, nothing
% else changes.

% Color definition (Graz palette, sourced from the original single-preamble
% version of this project -- these are the University of Graz's own official
% colors, not derived/approximated).
\definecolor{c_primary_dark}{HTML}{000000}
\definecolor{c_primary_highlight}{HTML}{ffd500}
\definecolor{c_primary_gray}{HTML}{c6c6c6}
\definecolor{c_primary_white}{HTML}{FFFFFF}
\definecolor{c_secondary_violet}{HTML}{948dc3}
\definecolor{c_tertiary_violet}{HTML}{c2bddf}
\definecolor{c_secondary_blue}{HTML}{689cce}
\definecolor{c_tertiary_blue}{HTML}{bacce8}
\definecolor{c_secondary_green}{HTML}{9fcb81}
\definecolor{c_tertiary_green}{HTML}{cce2ba}
\definecolor{c_secondary_red}{HTML}{e4684b}
\definecolor{c_tertiary_red}{HTML}{f1ae95}
\definecolor{c_secondary_orange}{HTML}{f6a968}
\definecolor{c_tertiary_orange}{HTML}{fbd0aa}
\definecolor{c_secondary_yellow}{HTML}{ffe063}
\definecolor{c_tertiary_yellow}{HTML}{ffeeaf}
\definecolor{c_tertiary_gray}{HTML}{e3e3e3}
% Not an official Uni Graz color (hence "custom", not "secondary" like the
% palette above) -- a softer near-black than c_primary_dark, used for
% quoteframe's full-bleed dark block, matching the specific dark gray (not
% pure black) ExampleSlides.pptx used for that same block. Also reused for
% hyperlink text color, for the same reason it works for quoteframe's
% attribution text: dark enough for good contrast/legibility on a light
% background.
\definecolor{c_custom_gray}{HTML}{333333}
\definecolor{c_custom_red}{HTML}{C00000}
\definecolor{c_custom_green}{HTML}{00B050}

% Logo + calibration.
% schemes/scheme_uni_graz/Logo_Uni-Graz_4c-cropped.pdf -- official vector artwork,
% user-supplied (schemes/scheme_uni_graz/Logo_Uni-Graz_4c.eps, kept alongside as the
% untouched source), converted to PDF via epstopdf and cropped to its true
% ink bounding box (computed with Ghostscript's own `-sDEVICE=bbox`, same
% treatment as MCI's and Innsbruck's own logos -- see
% Slidemaster_preamble_mci.tex's own comment for the fuller explanation of
% why the EPS's own %%BoundingBox isn't trusted directly). Vector, not
% raster, so it stays crisp at any zoom. Supersedes this file's own earlier
% raster logo (schemes/scheme_uni_graz/Logo Uni Graz 4c-cropped.png, itself cropped
% from schemes/scheme_uni_graz/Logo Uni Graz 4c.png's own symmetric transparent
% margin -- both kept alongside as historical/fallback assets, not deleted).
\newcommand{\gothamInstituteLogoPath}{schemes/scheme_uni_graz/Logo_Uni-Graz_4c-cropped.pdf}
\newlength{\gothamLogoHeight}
\setlength{\gothamLogoHeight}{3ex}
% Re-measured against the compiled PDF for this specific logo image -- see
% Slidemaster_preamble_base.tex's own comment on \gothamFrametitleRightpading
% for why this needs re-measuring per institution rather than reusing
% another one's value.
\setlength{\gothamFrametitleRightpading}{0.28cm}

% Itemize bullet color (2026-07-16, direct instruction): Graz keeps the
% project's original behavior, bullets in the primary highlight color.
\newcommand{\gothamBulletColor}{c_primary_highlight}

% Title bar and footer box/text colors (2026-07-16, direct instruction):
% Graz keeps the project's original behavior throughout -- title/footer
% boxes in the primary highlight color, title/footer text in the primary
% dark color (unlike MCI/Innsbruck, which switch these to their own
% secondary blue/white -- see Slidemaster_preamble_mci.tex's own comment).
\newcommand{\gothamTitleBoxColor}{c_primary_highlight}
\newcommand{\gothamTitleTextColor}{c_primary_dark}
\newcommand{\gothamFooterBoxColorDefault}{c_primary_highlight}
\newcommand{\gothamFooterTextColor}{c_primary_dark}

\usepackage{etoolbox} % \AtEndEnvironment, used for the itemize level-drop spacing below
\usepackage{tikz} % For painting a white box over title bar
\usetikzlibrary{calc}
\usetikzlibrary{arrows.meta} % for e.g. arr/.style={-{Stealth[...]}...} on diagram slides
\usetikzlibrary{positioning} % for `right=<dist> of <node>`-style relative node placement
\usetikzlibrary{decorations.text} % text-along-a-path, for a circular photo's curved caption
\usepackage{booktabs} % \toprule/\midrule/\bottomrule for gothamtable
\usepackage{colortbl} % \arrayrulecolor{}, used by gothamtable to color those rules
\usepackage{pifont} % \ding{51}/\ding{55} checkmark/X-mark glyphs, e.g. slides/fr3-Forms-of-Capital-Market-Efficiency-Strong-Semistrong-Weak-current.tex's own validity ratings
\usepackage{pgfplots} % real-data line/bar charts (e.g. slides/3vp-Metallgesellschaft-Oil-Price-Chart-current.tex) -- confirmed safe alongside this deck's own \tableofcontents, see CONVERSION_NOTES.md's "beamer/TeX gotchas" entry
\pgfplotsset{compat=1.18}

% Restores genuine Computer Modern math fonts. Found 2026-07-15, direct
% instruction ("math equations do not have serifs, look at vcv"): despite
% fontspec being loaded with `no-math` (beamerfontthemegotham.sty), math
% still wasn't landing on Computer Modern -- checked directly against a
% compiled PDF's own per-glyph font names (pymupdf), not just eyeballed:
% digits/math-italic letters were coming out as FiraSans-Light(Italic)
% (this deck's own body-text sans font) and operators/parens as CMSS10
% (Computer Modern's *sans-serif* shape, not its serif one). Root cause:
% beamer's own base classes redefine the math alphabets to follow the
% presentation's sans-serif text font by default, specifically so math
% matches an all-sans deck -- `no-math` only stops *fontspec* from doing
% this, it has no say over beamer's own separate override. Beamer ships
% the standard fix for exactly this case: a font theme named "serif"
% with an "onlymath" option that restores ordinary (Latin
% Modern/Computer-Modern-equivalent) math fonts without touching any of
% the surrounding sans-serif text. Confirmed directly afterward: every
% glyph in a re-checked formula now reports as LMRoman*/CMR10/CMMI8 (all
% genuinely serif), not FiraSans/CMSS10.
\usefonttheme[onlymath]{serif}

% \text{} (amsmath, used throughout for plain-word operands inside math,
% e.g. \text{E}[\tilde{R}_M]) is a separate mechanism from the math
% alphabets \usefonttheme[onlymath]{serif} above just fixed: unlike
% \mathrm/\mathit/etc., \text{} is defined to typeset in whatever the
% *surrounding running text* font is at that point in the document, not a
% math alphabet -- so on this deck's all-Fira body text, \text{} inside a
% formula still comes out in Fira Sans (and picks up italic if the
% enclosing text happens to be italicized), even though the digits/
% operators around it are correctly Computer Modern. Found 2026-07-20 on
% slides/2gd-Analyzing-Reward-to-Risk-Ratio-CAPM-current.tex's own
% \text{E}/\text{Cov} operands. Fixed generally, project-wide, by wrapping
% \text{}'s argument in an explicit upright-Computer-Modern font switch --
% \usefont{OT1}{cmr}{m}{n} rather than \mathrm{} (which would pull in
% \text{}'s own math-alphabet substitution logic unnecessarily) or
% \textrm{} (which switches *family* only, so still comes out italic
% inside an italicized context, which \text{} should never do here) --
% deliberately hardcoding "not italic" (shape n) rather than inheriting
% the ambient shape, since \text{}'s whole purpose in this deck's formulas
% is plain upright labels (E, Cov, "the" CAPM, ...), regardless of
% whatever font context happens to surround the formula itself.
\let\gothamOrigText\text
\renewcommand{\text}[1]{\gothamOrigText{\usefont{OT1}{cmr}{m}{n}#1}}

% Color definition, institute logo path, and the two logo-calibration
% lengths (\gothamLogoHeight, \gothamFrametitleRightpading) all live in the
% per-university preamble that \input{}s this file, not here -- this file
% is shared machinery only, no institution-specific branding. Every color
% token below (\c_primary_*, \c_secondary_*, \c_tertiary_*, \c_custom_*) and
% \gothamInstituteLogoPath/\gothamLogoHeight/\gothamFrametitleRightpading
% are assumed already defined by the time this file is reached -- see e.g.
% Slidemaster_preamble_uni-graz.tex for the Graz values, or
% Slidemaster_preamble_mci.tex/Slidemaster_preamble_uni-ibk.tex for the
% other two institutions this deck supports. \gothamLogoWidth is the one
% logo-derived length that *does* stay computed here (via \settowidth,
% further below) rather than hardcoded per institution, since it already
% adapts automatically to whatever image \gothamInstituteLogoPath points at.

% Style hyperlinks (\href{url}{text}) consistently across all slides:
% underlined, so a link is visually recognizable as such against plain
% body text -- yellow (this project's usual accent color) was tried first
% for the *text* itself, but yellow-on-white body text is hard to read,
% unlike yellow used as an accent against dark/black text elsewhere in the
% theme. Went through a few more candidates on 2026-07-06 (all direct
% instruction) before settling on c_secondary_blue for both text
% and underline (a real color reads more clearly as "this is a link" than
% any gray does, without yellow's legibility problem).
%
% 2026-07-07, direct instruction: split text color and underline color
% instead of using one color for both -- link text back to
% c_custom_gray (the near-black dark gray also used for
% quoteframe's attribution text), underline in a separate accent color.
% Plain \underline can't do this: its rule is drawn in whatever color is
% current at that point, so a single \underline{...} call always
% underlines in its own text's color, never a second, different one.
% \gothamHrefUnderline below reimplements it instead: measure the linked
% text's own width into a savebox, print that box in customgray, then
% back up by the box's own width and draw a plain \rule of the same
% width, raised to sit just below the text's baseline, in the accent
% color. Same single-line limitation as plain \underline (and the reason
% \url{} is used instead of \href on slides/acw-Some-Topics-in-Testing-the-EMH-current.tex --
% see that slide's own comment): a \rule, like \underline, can't span a
% line break either, so this only ever wraps normal, comparatively short
% \href display text, never something that needs mid-word breaking.
%
% First tried c_primary_highlight for the underline (same day) -- confirmed
% via the compiled PDF that the rule *was* being drawn (right color,
% position, width), so not a compile-time bug, but yellow (very high
% luminance, close to white) at this rule's 0.4pt thickness reads as
% essentially invisible against the white slide background, worst at
% footnote size (see e.g. slides/za2-Startup-Financing-Categorization-current.tex's
% own long inline URL) where it's easy to miss entirely. Switched
% 2026-07-07 (direct instruction) to c_secondary_blue instead --
% far more contrast against white at the same thickness -- and the rule
% was also thickened at the same time (0.4pt -> 0.6pt), on the reasoning
% that it might need it for the same contrast reason. It didn't: blue
% already reads clearly at the original thickness, and 0.6pt instead
% just looked heavier than intended (direct feedback 2026-07-08, e.g. on
% slides/6qi-Yield-Spreads-Corporate-Bonds-current.tex). Reverted to
% 0.4pt -- only the color was ever actually asked to change.
\hypersetup{colorlinks=true, urlcolor=c_custom_gray, linkcolor=c_custom_gray}
\let\gothamOrigHref\href
\newsavebox{\gothamHrefTextBox}
%
% 2026-07-08 (direct instruction): the underline's own rule thickness now
% scales with the ambient font size instead of a single fixed 0.2pt for
% every link everywhere -- 0.2pt at this deck's footnote size, 0.4pt at its
% normal body-text size, interpolated (and extrapolated, for any other
% size a link happens to appear at) linearly in between.
%
% Deriving "current font size" from \f@size (the textbook NFSS way) was
% tried first and confirmed broken specifically *inside* a beamer frame --
% \typeout{\meaning\f@size} right after this preamble loads shows a normal,
% defined macro (`macro:->10.95`), but the exact same check inside
% `\begin{standardframe}...\end{standardframe}` reports `\f@size` as
% *undefined*, presumably from beamer's own overlay-aware font-switching
% machinery not keeping classic NFSS's own current-size bookkeeping in
% sync the way plain (non-beamer) LaTeX does. Since every real use of this
% macro happens inside a frame, that path was a dead end regardless of how
% well it worked in isolation.
%
% Measuring a reference glyph's own rendered height instead sidesteps this
% entirely: building a box and reading back its own height just measures
% whatever font is *actually* active at the call site, independent of
% whether beamer's own internal "current size" bookkeeping happens to
% agree -- and within one font
% family, every glyph's metrics scale by exactly the same ratio the point
% size does, so a measured height is just as valid a stand-in for "current
% size" as \f@size would have been, had it worked. "M" (present in every
% context this is used in) is the reference glyph; \gothamHrefUnderlineFnHeight/
% \gothamHrefUnderlineBodyHeight are its height at this deck's footnote font
% and normal body font respectively, each measured once, here, in the
% relevant font context.
% Deliberately built from raw TeX primitives (\newdimen/\setbox/\ht),
% *not* \newlength/\setlength/\settoheight -- found, while building this,
% that loading this deck's theme silently turns \setlength and
% \settoheight into something else entirely: `\meaning\setlength` reads
% `macro:->\protect \setlength` right after `\usetheme{gotham}` (still the
% ordinary kernel definition beforehand), and the same for `\settoheight`.
% Not chased to which file does this or why (bisected only as far as
% "plain beamer alone doesn't, beamer+gotham does" -- most likely
% `translator`, which beamer loads for localized strings and is known to
% robustify commands this way, but not confirmed) -- what matters
% practically is that a *global* assignment through either one silently
% no-ops instead of erroring, confirmed directly and in isolation (a
% minimal beamer+gotham file, no other content): `\global\setlength{\len}
% {\otherlen}` leaves `\len` at 0pt after the group closes, while the
% exact same pattern with plain `\newdimen`/raw `=` assignment propagates
% correctly. Every plain (non-global) `\setlength` elsewhere in this
% project's own files is unaffected -- it's specifically the *global*
% form that breaks, which is why this is the first place it's shown up
% despite `\setlength` being used constantly throughout this preamble.
\newdimen\gothamHrefUnderlineFnHeight
\newdimen\gothamHrefUnderlineBodyHeight
\newdimen\gothamHrefUnderlineNowHeight
\newdimen\gothamHrefUnderlineThickness
% A dedicated box register, not the commonly-borrowed box0 -- several
% slides in this deck (the "Long-term bonds" tables) already use box0
% themselves (`\sbox0{...}`/`\usebox0`) for their own, unrelated
% mid-table-layout purposes, and a link can appear anywhere, including
% inside one of those tables -- reusing box0 here would silently clobber
% whatever they were holding there.
\newbox\gothamHrefUnderlineBox
% The two reference heights are (re)measured on every single \href call,
% not once at preamble-load time -- tried once-at-load first and found a
% second, independent problem with it: `\usebeamerfont{footnote}` right
% here (before `\begin{document}`) silently has no effect at all (footnote-
% size "M" and normal-size "M" measured identically, confirmed directly),
% presumably because beamer's font-theme machinery isn't fully installed
% until the document starts. Since every real `\href` call happens well
% inside the document (inside some frame), measuring fresh each time
% sidesteps that timing question entirely rather than needing to find the
% right hook to measure once, early enough to be free but late enough to
% be correct -- the extra cost is two throwaway one-glyph boxes per link,
% not something worth optimizing away in a slide deck.
\newcommand{\gothamHrefUnderlineMeasureRefHeights}{%
	{\usebeamerfont{footnote}\setbox\gothamHrefUnderlineBox=\hbox{M}\global\gothamHrefUnderlineFnHeight=\ht\gothamHrefUnderlineBox}%
	{\normalsize\setbox\gothamHrefUnderlineBox=\hbox{M}\global\gothamHrefUnderlineBodyHeight=\ht\gothamHrefUnderlineBox}%
}
% A plain e-TeX \dimexpr can multiply a dimension by another dimension and
% then divide by a third, as one combined operation, as long as the
% *middle* operand of the chain is itself wrapped in its own nested
% \dimexpr...\relax -- confirmed directly (an expl3 \dim_set:Nn attempt
% first, using \seq_.../\fp_... the way \setfootnotesymbols above does,
% kept hitting "Missing ) inserted for expression"/type-mismatch errors:
% dividing one dimension by another is a *ratio*, not a dimension, and
% \dim_set:Nn's own expression grammar doesn't implicitly convert between
% the two the way this nested-\dimexpr form does). No expl3 needed here at
% all -- plain e-TeX arithmetic, already available in any current LaTeX.
\newcommand{\gothamHrefUnderlineSetThickness}{%
	\gothamHrefUnderlineThickness=\dimexpr 0.2pt +
		(\gothamHrefUnderlineNowHeight - \gothamHrefUnderlineFnHeight)
		* \dimexpr0.2pt\relax
		/ (\gothamHrefUnderlineBodyHeight - \gothamHrefUnderlineFnHeight)\relax
}
\newcommand{\gothamHrefUnderline}[1]{%
	\sbox{\gothamHrefTextBox}{#1}%
	\textcolor{c_custom_gray}{\usebox{\gothamHrefTextBox}}%
	\hspace{-\wd\gothamHrefTextBox}%
	\gothamHrefUnderlineMeasureRefHeights
	\setbox\gothamHrefUnderlineBox=\hbox{M}\gothamHrefUnderlineNowHeight=\ht\gothamHrefUnderlineBox
	\gothamHrefUnderlineSetThickness
	\raisebox{-0.35ex}[0pt][0pt]{\textcolor{c_secondary_blue}{\rule{\wd\gothamHrefTextBox}{\gothamHrefUnderlineThickness}}}%
}
\renewcommand{\href}[2]{\gothamOrigHref{#1}{\gothamHrefUnderline{#2}}}
% \url{} (the `url` package hyperref depends on/patches) defaults to a
% monospace-ish "tt" style regardless of the above -- used deliberately in
% only one place in this deck,
% slides/acw-Some-Topics-in-Testing-the-EMH-current.tex, specifically
% *for* its line-breaking-at-slashes behavior (a plain \href with the
% URL as its own display text has no break points at all -- see that
% slide's own comment). Every other same-text-as-target URL in the deck
% (including slides/bb6-A-Perpetual-Bond-current.tex's own video link,
% which used to be a \url{} too) uses plain \href{url}{url} instead, since
% \url{} doesn't pick up this deck's custom colored underline below (its
% breakable content can't be measured into the savebox the underline
% needs -- see that mechanism's own comment) and most of these URLs are
% short enough to not need the breaking \url{} exists for anyway. 2026-07-07,
% direct instruction: that one genuinely-needed \url{} (on s-acw) should
% still look like this deck's other hyperlinks, i.e. the surrounding body
% font, not \url{}'s own default monospace.
% `\urlstyle{same}` (a standard url-package switch) does exactly that,
% with no effect on \url{}'s own breaking behavior. Set deck-wide, not
% only on that one slide, so any future \url{} use matches by default
% too, rather than needing to remember this each time. Coloring already
% matches automatically (hyperref's own \url respects
% \hypersetup{urlcolor=...} the same way \href does); the underline
% \href gets above doesn't carry over, since \underline's own argument
% can't contain line breaks -- wrapping \url{}'s necessarily-breakable
% text in \underline would silently undo the breaking this exists for,
% so that one difference is accepted rather than "fixed."
\urlstyle{same}

% gotham's own footnote color (beamercolorthemegotham.sty:
% \setbeamercolor{footnote}{fg=normal~text.fg!90}) is a 90%-black gray --
% against plain black body text, customgray links contrast clearly (found
% while investigating a 2026-07-05 report that a footnote's own \href
% didn't read as a link the way body-text ones like the "Moodle" link on
% slides/hxr-Administrative-Details-current.tex do); against gotham's own
% dimmed footnote gray, the same customgray is only 25/255 lighter and
% barely distinguishable at footnote size, with only the underline giving
% it away. Overridden to plain black here so a footnote link contrasts the
% same way a body-text one does -- a one-line override of the theme's own
% color (not editing beamercolorthemegotham.sty itself, which stays
% untouched per its own upstream-file status), and a plausible general
% legibility improvement for footnote text regardless of links.
\setbeamercolor{footnote}{fg=c_primary_dark}

% Footer text color (\FooterText/slide id/frame number) -- \gothamFooterTextColor,
% set by the per-university preamble before \input{}ing this file (the
% footer box's own background color, \gothamFooterBoxColorDefault above, is
% a separate, independently configurable setting). No explicit
% \setbeamercolor{footline}{...} existed before this -- the footer text
% simply inherited whatever "normal text" foreground already was (black),
% which happened to read fine against every frame's own yellow/orange
% footer box so far, but wouldn't automatically read against an arbitrary
% future footer color (e.g. a dark blue) without this override.
\setbeamercolor{footline}{fg=\gothamFooterTextColor}

\colorlet{colorDark}{c_primary_dark} % FG in light/normal mode
% Title bar background -- \gothamTitleBoxColor, set by the per-university
% preamble before \input{}ing this file (same idiom as
% \gothamBulletColor/\gothamFooterBoxColorDefault above).
\colorlet{colorA}{\gothamTitleBoxColor} % frametitle, standin.out

\gothamset{background=light} % Needs to go after color redefinitions, so they are applied
\setbeamercolor{background canvas}{bg=white}

% gotham's own default behavior on \section{} (used for a deck's chapter
% grouping/table of contents) is to auto-insert its *own* plain, unstyled
% section-divider frame plus a small "current section highlighted in the
% outline" frame, every time -- redundant with (and visually inconsistent
% next to) this project's own hand-styled sectionframe (see README.md's
% "Special frame types" section), which every chapter in this deck already
% opens with as an explicit \slide{}. Disabled so \section{} only feeds
% \tableofcontents (see slides/nkj-Agenda-current.tex) without also
% inserting its own extra frames.
\gothamset{sectionframe default=false}

% Same reasoning, for \subsection{} (2026-07-15, direct instruction):
% gotham's default subsection-divider frame is disabled too, so
% \subsection{} calls still register in the PDF's own bookmark/outline
% structure (Acrobat's "chapters" pane -- a plain PDF sectioning feature,
% independent of whether beamer auto-inserts a frame for it) and still feed
% \tableofcontents (see slides/nkj-Agenda-current.tex's own
% `\tableofcontents[hideallsubsections]`, which then hides them from that
% one slide specifically), without also spawning their own unstyled divider
% slide in the middle of a chapter every time one is added.
\gothamset{subsectionframe default=false}

% Title text color -- \gothamTitleTextColor, set by the per-university
% preamble before \input{}ing this file (the title bar's own background
% color, \gothamTitleBoxColor/colorA above, is a separate, independently
% configurable setting).
\setbeamercolor{frametitle}{fg=\gothamTitleTextColor}

% A color leak from \titleframe/\quoteframe into whichever frame follows
% them is fixed in standardframe's own opening code below (search
% "Forces this frame's own body") rather than here -- an
% \AtBeginEnvironment{frame}{\color{c_primary_dark}} hook was tried first
% and confirmed, by isolated reproduction, to have *no effect at all*:
% beamer's own "frame" isn't a plain \newenvironment etoolbox's
% \AtBeginEnvironment can patch that way, so the hook silently never fired.

% Set content text margins
\setbeamersize{text margin left=.45cm, text margin right=.45cm}

%\colorlet{colorPale}{gPaleYell} % BG in light/normal mode
%\colorlet{colorDark}{gDarkBlack} % FG in light/normal mode
%\colorlet{colorA}{gAnthracite} % frametitle, standin.out,
%\colorlet{colorAreversed}{gLightTeal} % frametitle, standin.in,
%\colorlet{colorB}{gMidGrey} % gray BG : progress bar, blocks
%\colorlet{colorC}{gDeepYellOr} % progress bar
%\colorlet{colorD}{gLightOrange} % alert
%\colorlet{colorE}{gLightGreen} % example

% Opt-in yellow border for an image placed via a tikz node (e.g.
% \node[gothamImageBorder] {\includegraphics{...}};) -- matches the
% 3pt/accent1-yellow outline PowerPoint draws directly on a picture shape's
% own border (<a:ln w="38100"><a:solidFill><a:schemeClr val="accent1"/>,
% where accent1 turned out to already be this project's own
% c_primary_highlight, not a new color to map). A tikz *style* (a named
% bundle of node keys), not a wrapper command: adding or removing it is
% just adding or deleting one word from that one \node's own option list,
% no separate macro call/argument needed -- add `gothamImageBorder` to draw
% the border, leave it off for a plain image, same as any other tikz key.
%
% Thickness is a single named length, \gothamImageBorderWidth, not a bare
% "3pt" baked directly into the style (2026-07-15, direct instruction: "make
% sure the yellow image borders have consistent thickness throughout the
% project by defining a single variable") -- confirmed directly, by
% measuring the actual PDF content-stream "w" operator, that every
% gothamImageBorder use already rendered at the exact same 3pt regardless
% (a genuinely shared style was already in effect, no numeric drift found),
% so this is a refactor for future adjustability/discoverability, not a fix
% for an actual inconsistency -- any perceived thickness difference between
% slides traces to something else instead (e.g. a border straddling flush
% image content on one slide and not another, which reads as "heavier" even
% at an identical stroke width -- see the gothamImageBorder padding gotcha
% below).
\newlength{\gothamImageBorderWidth}
\setlength{\gothamImageBorderWidth}{2.5pt}
\tikzset{gothamImageBorder/.style={draw=c_primary_highlight, line width=\gothamImageBorderWidth}}

% The "padding gotcha" promised above, finally root-caused (2026-07-23,
% background-task follow-up to the IA2 fine-tuning batch #2 anomaly on
% slides/m2k-Exercise-CAPM-2-current.tex/
% slides/n8e-Tail-Risk-Variance-Clustering-current.tex/
% slides/fm6-Tail-Risk-Market-vs-Simulated-current.tex/
% slides/xqq-Fat-Tails-current.tex, all fixed there with a pragmatic
% `line width=1.2pt` override and an explicitly *unresolved* root cause at
% the time -- "declares an identical stroke width and CTM in the compiled
% PDF's own raw content stream, yet renders at a genuinely different pixel
% thickness", with CTM scale/clips/dash/cap/join all directly checked and
% ruled out). Confirmed via a controlled reproducer (four `gothamImageBorder`
% nodes -- plain opaque image, nested pgfplots axis, plain image at
% `inner sep=3pt`, nested plain tikzpicture at `inner sep=3pt` -- rendered
% at 600dpi and measured by a half-max cross-section profile, not just
% eyeballed): the one variable that actually predicts the rendered width is
% `inner sep`, not content type at all. At `inner sep=0pt` a *fully opaque*
% image measured 1.2pt; the exact same image at `inner sep=3pt` measured
% 2.52pt (matching the declared 2.5pt); a nested pgfplots axis and a plain
% nested tikzpicture, both at `inner sep=3pt`, also measured 2.52pt. Two
% more points confirmed it's a smooth transition, not a binary switch:
% `inner sep=0.6pt` measured 1.8pt (= the 1.2pt floor plus that 0.6pt gap,
% exactly); `inner sep=1.25pt` (= \gothamImageBorderWidth/2, the stroke's
% own half-width) already measured the full 2.52pt.
%
% The mechanism: tikz draws a node's border stroke centered on the node's
% own boundary path, then paints the node's *content* on top of it,
% afterwards. At `inner sep=0`, the content's own bounding box coincides
% exactly with that boundary path, so if the content is opaque all the way
% to its own edge, it paints directly over the stroke's inner half,
% visually leaving only the outer half (`\gothamImageBorderWidth`/2,
% ~1.2pt) -- not a rendering bug, just tikz's normal draw order interacting
% with zero padding. A gap of at least `\gothamImageBorderWidth`/2 (or
% content that's transparent right at its own edge -- see below) leaves the
% border with nothing opaque to cover it, so it renders at its true
% declared width instead.
%
% This also explains why the fix landed on exactly the 4 slides it did,
% each for a different reason: slides/m2k-Exercise-CAPM-2-current.tex
% used `inner sep=3pt` for unrelated layout reasons (room for the axis's
% own tick labels), well past the 1.25pt threshold, so nothing ever covered
% its border. slides/n8e-Tail-Risk-Variance-Clustering-current.tex/
% slides/fm6-Tail-Risk-Market-vs-Simulated-current.tex/
% slides/xqq-Fat-Tails-current.tex all use `inner sep=0pt` -- matching the
% deck's own overwhelming convention -- but their own chart PNGs
% (`chart-index-daily-returns-*.png`/`chart-index-fat-tails-histogram*.png`,
% matplotlib exports with a transparent background) measured alpha=0 at
% every edge pixel (confirmed directly, Pillow), so there was no opaque
% content to cover their own border either. Every other `gothamImageBorder`
% use in this project that was never flagged uses an opaque image (a JPG,
% or a PNG with no transparent margin) at `inner sep=0pt` -- which
% *accidentally* renders at the ~1.2pt covered-half width instead of the
% full 2.5pt this constant nominally declares, and that accidental width is
% what reads as "the deck's normal border thickness" throughout, since it's
% by far the most common case.
%
% Practical upshot for a *new* slide: a `gothamImageBorder` node with
% `inner sep=0pt` wrapping opaque content (the common case) needs no
% special handling -- it will match the rest of the deck automatically. A
% node that instead uses a nonzero `inner sep`, or wraps content with a
% transparent edge (a matplotlib export, an SVG-derived PDF with padding,
% ...), will render at the *full* \gothamImageBorderWidth and look
% noticeably heavier next to the deck's other borders -- either add an
% explicit `line width=\gothamImageBorderWidth/2` (matching the 4 fixed
% slides' own literal 1.2pt) to match the accidental deck-wide look, or
% reduce `inner sep` below \gothamImageBorderWidth/2 so the same coverage
% effect kicks in naturally. The 4 existing per-slide overrides above are
% left as-is -- they're already the correct fix for what's now a fully
% understood cause, not a mystery workaround needing revisiting.

% Reusable "mugshot" -- a circular-cropped photo, bordered like any other
% gothamImageBorder image, with a caption curving along its own bottom arc
% (2026-07-06, direct instruction, generalized from
% slides/4uq-Definition-of-Capital-Market-Efficiency-current.tex's own
% Eugene Fama photo -- see that slide's earlier comment for the source
% PPTX effect this reproduces, `<a:prstGeom prst="ellipse">` +
% `<a:prstTxWarp prst="textArchDown">`, and CONVERSION_NOTES.md's own
% "beamer/TeX gotchas" entry, search "circular photo crop", for the
% clip/border/decorations.text mechanics themselves -- unchanged here,
% just factored out so a second photo doesn't need to repeat them).
%
%   \coordinate (myMugshotCenter) at (...);
%   \gothamMugshot{myMugshotCenter}{media/photo-example.jpeg}{2.6cm}{Some Person, NP 2013}
%
% #1 the *name* of a coordinate the caller has already placed (not a raw
% point expression -- needed as a plain name, not just a point, since the
% caption's own arc path re-references it by name via `++(...)`, which
% only works with a named coordinate/node, not an arbitrary `($...$)`
% calc expression); #2 the image path; #3 the desired final diameter;
% #4 the caption text.
%
% Scales the source photo so it fully covers the circle with no gap
% anywhere -- both dimensions end up at least as big as the diameter --
% while using as much of the source photo as possible (2026-07-07,
% direct instruction: "should not include whitespace... scale the photo
% such that its edges just touch the circle's inner border"). This
% briefly (same day) scaled the *other* way -- by the photo's own
% diagonal, the largest scale at which the *entire* rectangle fits
% inside the circle with nothing clipped at all -- which does avoid
% ever clipping real content, but leaves the circle's own four corners
% uncovered (visible as plain background there, since a non-circular
% photo can only touch an inscribing circle at its corners, not along
% the whole rim) -- exactly the "whitespace" this round's feedback
% ruled out in favor of full coverage. Covering scale is
% `diameter / min(w,h)`: whichever of the photo's own two dimensions is
% *smaller* becomes exactly the diameter (that edge pair touches the
% circle's own inner border, per the direct instruction), and the other,
% larger dimension overhangs the diameter and gets clipped by the
% circular `\clip` -- for a portrait photo (every one used here so far),
% that means width becomes the diameter and a sliver clips off the top
% and bottom, which is the origin of the clipping problem the diagonal
% approach above existed to solve in the first place. Reintroduced
% deliberately: showing as much of the photo as the "no whitespace"
% constraint allows is now the stated priority, over never clipping
% anything.
%
% Measuring the photo's own native size needs `\savebox` (into
% \gothamMugshotBox, declared once below, not a fresh box register per
% call) -- `\includegraphics` with no size option at all renders at the
% image's native resolution, and `\wd`/`\ht` then read that box's actual
% dimensions. pgfmath (already loaded for tikz) converts any TeX
% dimension fed into it to a plain pt-count number automatically, so
% arithmetic on the two measured lengths just works. `min(1,ratio)`
% (ratio = height/width) picks out whichever of the two is smaller
% *relative to the other*, without needing a `sqrt`/`^` at all this time
% -- so this version doesn't hit the "! Dimension too large" pgfmath
% overflow the diagonal version's own arithmetic did (see this file's
% own comment history / CONVERSION_NOTES.md's "beamer/TeX gotchas" for
% that point (easy to misread as a harmless warning, since compilation
% continues and produces *a* number regardless). Reformulating via the
% aspect ratio sidesteps the whole problem: `scale = diameter /
% sqrt(1+ratio^2)` is the same formula algebraically (factor `w` out of
% `sqrt(w^2+h^2)` and it becomes `w*sqrt(1+(h/w)^2)`), but every
% the full story of that overflow, from when this macro's own arithmetic
% did use a `sqrt`).
%
% Caption font reverted to this deck's normal body font (2026-07-07,
% direct instruction: "change the mugshot template back to using regular
% font instead of monospace") -- `\ttfamily` (Fira Mono) was a direct
% instruction on 2026-07-06, one day earlier; \normalfont resets
% explicitly rather than just omitting \ttfamily, so this doesn't
% silently inherit whatever font happens to be active at the call site.
\newsavebox{\gothamMugshotBox}
\newcommand{\gothamMugshot}[4]{%
	\savebox{\gothamMugshotBox}{\includegraphics{#2}}%
	\pgfmathsetmacro{\gothamMugshotImgW}{\wd\gothamMugshotBox}%
	\pgfmathsetmacro{\gothamMugshotImgH}{\ht\gothamMugshotBox}%
	\pgfmathsetmacro{\gothamMugshotRatio}{\gothamMugshotImgH/\gothamMugshotImgW}%
	\pgfmathsetmacro{\gothamMugshotScaledW}{(#3)/min(1,\gothamMugshotRatio)}%
	\begin{scope}
		\clip (#1) circle (0.5*#3);
		\node[anchor=center, inner sep=0pt] at (#1)
			{\includegraphics[width=\gothamMugshotScaledW pt]{#2}};
	\end{scope}
	\draw[gothamImageBorder] (#1) circle (0.5*#3);
	{\normalfont\tiny
	\path[decoration={text along path, text={#4}, text align=center}, decorate]
		(#1) ++(200:0.5*#3+0.25cm) arc (200:340:0.5*#3+0.25cm);
	}
}

% Reusable "operating/cash conversion cycle" bar diagram (FM2_2026s_ps_deck.pptx
% slides 22/24/25/26 all embed the exact same picture -- an EMF, not a native
% shape/table -- so rather than convert that EMF and re-embed it 4 times over,
% it's redrawn once natively here and reused, matching this project's own
% "clean/geometric diagrams get redrawn, not extracted" policy). Three
% stacked bars (Inventory period+Accounts receivable period on top, Accounts
% payable period+Cash conversion cycle in the middle, Operating cycle
% spanning the full width on the bottom) plus 4 tick labels below at
% x=0/(1/3)/(2/3)/1. Geometry/colors measured directly from the source
% deck's own rendered PNG via pixel analysis (PIL run-length color
% detection), not eyeballed: the two internal split points are exactly 1/3
% and 2/3 of the total width (measured 0.3339/0.6672 against a plain
% PowerPoint illustrative diagram, not slide-23's own specific-day example --
% close enough to exact thirds, and a plausible round number for the
% presenter to have actually used, that thirds is treated as the real value
% rather than curve-fitting noise). Fill colors matched to this project's
% *existing* palette rather than the source's own literal RGB, per this
% deck's established emphasis-color policy: source yellow was already an
% exact match for `c_primary_highlight`; source green (RGB 146,208,80) and
% purple (RGB 128,100,162) were close enough to `c_secondary_green`
% and `c_secondary_violet` respectively to reuse rather than define
% two new one-off colors.
%
% Drawn at a fixed "natural" size (10cm wide) with its own top-left corner
% at the calling tikzpicture's origin -- wrap a call in
% \begin{scope}[shift={(x,y)}, scale=s]\gothamCashCycleBars\end{scope} to
% place/resize it, the same convention as this deck's other reusable
% diagram pieces (e.g. slides/bdh-Working-Capital-Example-current.tex's
% circular flow diagram, resized the same way after the fact per user
% feedback) -- resizing later without needing to re-derive any coordinates.
\newcommand{\gothamCashCycleBars}{%
	\fill[c_primary_highlight] (0,0) rectangle (6.667,-0.42);
	\fill[c_secondary_green] (6.667,0) rectangle (10,-0.42);
	\node[align=center, font=\tiny] at (3.333,-0.21) {Inventory period};
	\node[align=center, font=\tiny] at (8.333,-0.21) {Accounts receivable period};
	\fill[c_secondary_green] (0,-0.47) rectangle (3.333,-0.89);
	\fill[c_primary_highlight] (3.333,-0.47) rectangle (10,-0.89);
	\node[align=center, font=\tiny] at (1.667,-0.68) {Accounts payable period};
	\node[align=center, font=\tiny] at (6.667,-0.68) {Cash conversion cycle};
	\fill[c_secondary_violet] (0,-0.94) rectangle (10,-1.36);
	\node[align=center, font=\tiny, text=c_primary_white] at (5,-1.15) {Operating cycle};
	\foreach \gccx in {0,3.333,6.667,10} {
		\draw[c_primary_dark] (\gccx,-1.36) -- (\gccx,-1.46);
	}
	% text width on all 4, narrow enough to force an actual 2-line wrap
	% (not just cap an already-short line): without it, "Raw material
	% purchase" and "Payment for raw materials" (and, symmetrically, the
	% last two) run into each other -- confirmed directly by rendering.
	% A *wide* text width (e.g. 3.1cm, comfortably more than either label's
	% own one-line width) does NOT fix this: \node's text width only forces
	% a wrap once the text exceeds it, so a generous width just lets both
	% neighbors keep their full one-line extent and still collide -- this
	% was tried first and confirmed not to work by rendering, not assumed.
	% 1.5cm is deliberately narrower than any of the 4 would take on one
	% line, guaranteeing the wrap for all of them.
	%
	% All 4 use anchor=north (centered on their own tick, like the middle
	% two always did) -- the outer two originally used anchor=north
	% west/east instead (flush with the tick rather than centered on it),
	% which looked inconsistent with the middle two once seen side by side
	% at actual size. Centering the outer two on their own tick (rather
	% than flush against the diagram's own left/right edge) does push their
	% wrapped text a bit further left/right past the diagram's own
	% bounding box -- checked this stays within the page margin at every
	% scale/position this macro is actually called at in this deck (the
	% narrowest margin case leaves ~0.7cm of clearance), not just assumed
	% safe from the numbers alone.
	\node[anchor=north, font=\tiny, align=center, text width=1.5cm] at (0,-1.5) {Raw material purchase};
	\node[anchor=north, font=\tiny, align=center, text width=1.5cm] at (3.333,-1.5) {Payment for raw materials};
	\node[anchor=north, font=\tiny, align=center, text width=1.5cm] at (6.667,-1.5) {Sale of finished goods};
	\node[anchor=north, font=\tiny, align=center, text width=1.5cm] at (10,-1.5) {Cash collected on sales};
}

% Reusable circular flow diagram (Cash -> Raw materials -> Finished goods ->
% Receivables -> Cash), first built for
% slides/bdh-Working-Capital-Example-current.tex and reused (smaller) on
% slides/qq5-Cash-Conversion-Cycle-current.tex -- factored out here once a
% second slide needed the same picture, same reasoning as
% \gothamCashCycleBars above. Args: #1 diagram center coordinate, #2 radius
% (box-center distance from that center, a bare cm number, no unit suffix),
% #3 the box `scale` factor (a plain node `scale` transform, not a bigger
% font -- see bdh's own comment, preserved below, for why: this diagram was
% resized once already per user feedback, and `scale` was what correctly
% grew the padding/rounded-corners/arrow proportions together rather than
% just the text). Arrow length/width scale with #3 too, computed via
% \pgfmathsetmacro since tikz option values can't do arithmetic inline.
%
% [Preserved from bdh's own resize, 2026-07-06: box size grew by a uniform
% ~1.44x there (compared via the isolated "Cash" box, the one corner with no
% arrow touching it to contaminate the measurement: 0.94cm to 1.36cm wide,
% 0.50cm to 0.74cm tall), and the box-to-box radius grew by the same
% ~1.44-1.45x (1.15cm to ~1.65cm) -- both measured from the user's own
% "Clipboard 1.jpg" reference via pixel analysis, not eyeballed, and
% consistent with a single uniform resize of the whole group.]
% #1 is a *complete* coordinate spec, exactly as it would appear after
% "\coordinate (name) at " -- e.g. "(current page.north east)" or
% "($(current page.north east)+(-3.81cm,-4.61cm)$)", dollar signs and all
% when it's a calc expression. Resolved into a plain \coordinate first
% rather than substituted directly into each box's own ($...$) calc
% expression, since a calc expression can't contain another calc
% expression's $...$ delimiters nested inside it.
\newcommand{\gothamCircularFlowDiagram}[3]{%
	\pgfmathsetmacro{\gothamCFDarrowLen}{5*#3}
	\pgfmathsetmacro{\gothamCFDarrowWidth}{1.3*#3}
	\begin{scope}
		\coordinate (gothamCFDctr) at #1;
		\node[fill=c_primary_highlight, rounded corners=1pt, inner sep=4pt, align=center, font=\footnotesize, scale=#3] (gothamCFDcash) at ($(gothamCFDctr)+(0,#2)$) {Cash};
		\node[fill=c_primary_highlight, rounded corners=1pt, inner sep=4pt, align=center, font=\footnotesize, scale=#3] (gothamCFDraw) at ($(gothamCFDctr)+(#2,0)$) {Raw materials};
		\node[fill=c_primary_highlight, rounded corners=1pt, inner sep=4pt, align=center, font=\footnotesize, scale=#3] (gothamCFDfin) at ($(gothamCFDctr)+(0,-#2)$) {Finished goods};
		\node[fill=c_primary_highlight, rounded corners=1pt, inner sep=4pt, align=center, font=\footnotesize, scale=#3] (gothamCFDrec) at ($(gothamCFDctr)+(-#2,0)$) {Receivables};
		\draw[-{Stealth[length=\gothamCFDarrowLen pt]}, line width=\gothamCFDarrowWidth pt, c_primary_highlight] (gothamCFDcash) to[bend left=18] (gothamCFDraw);
		\draw[-{Stealth[length=\gothamCFDarrowLen pt]}, line width=\gothamCFDarrowWidth pt, c_primary_highlight] (gothamCFDraw) to[bend left=18] (gothamCFDfin);
		\draw[-{Stealth[length=\gothamCFDarrowLen pt]}, line width=\gothamCFDarrowWidth pt, c_primary_highlight] (gothamCFDfin) to[bend left=18] (gothamCFDrec);
		\draw[-{Stealth[length=\gothamCFDarrowLen pt]}, line width=\gothamCFDarrowWidth pt, c_primary_highlight] (gothamCFDrec) to[bend left=18] (gothamCFDcash);
	\end{scope}
}

% A \frac/\dfrac-stacked formula reads as cramped against a plain-text
% neighbor in an itemize -- it's visibly taller than a normal line (both
% above *and* below the line, since the fraction bar sits on the text
% baseline with numerator/denominator stacked around it), so both the
% fraction's own item *and* the item that follows it need a bit of extra
% headroom that a plain-text \item doesn't. Started as a manually-repeated
% \vspace{3pt} before each such \item (first on
% slides/u74-Example-Dynamic-Mattress-Company-Periods-current.tex);
% \tallitem below is that same convention factored into a single drop-in
% \item replacement, so a future fraction-heavy slide (or a change to the
% spacing amount) doesn't mean re-editing several slide files by hand.
% Takes an optional leading multiplier (default 1, i.e. 3pt): \tallitem[1.5]
% for 4.5pt, for a formula tall enough that the default doesn't read as
% enough room. Usage: replace \item with \tallitem both on any itemize line
% containing a stacked fraction *and* on the item immediately following it
% (same nesting level; skip this second replacement if that next item is
% itself already a \tallitem for its own fraction, or if the fraction item
% is the last one in its list) -- it's a plain substitute, not an
% environment, so it doesn't change how the surrounding list itself is
% structured. (Named \gothamFracItem originally; renamed 2026-07-21 to the
% more discoverable \tallitem, and given the optional multiplier and the
% following-item rule above at the same time -- no behavior change for any
% existing default-strength call site.)
% Ratio, not a frozen length (2026-07-22, direct instruction: "modify
% \tallitem and \tallitemnobullet to also use scaling relative to font
% size") -- the original \newlength+\setlength{...}{3pt} pair freezes to
% one fixed pt value the moment \setlength runs, the same issue (and same
% fix) as \gothamItemsepRatioI/II/III/IV further below: converted to a
% bare ratio macro, with the "em" unit written back on inside \tallitem's
% and \tallitemnobullet's own bodies (not here), so it's recomputed fresh
% against whatever font is actually live at each call site -- correct
% regardless of nesting level or a \footnotesize/\scriptsize override,
% unlike the old fixed 3pt (calibrated against one specific level's font,
% silently wrong at any other). 0.3 reproduces the original 3pt exactly
% at level 2's own font (10pt) -- the level \tallitem was first used at,
% on slides/u74-Example-Dynamic-Mattress-Company-Periods-current.tex (see
% this macro's own comment above) -- so no visual change there, while now
% correctly scaling up/down at every other level and font override.
\newcommand{\gothamTallItemBaseRatio}{0.3}
\newlength{\gothamTallItemBaseUnit}
\newlength{\gothamTallItemSpace}
\newcommand{\tallitem}[1][1]{%
	\setlength{\gothamTallItemBaseUnit}{\gothamTallItemBaseRatio em}%
	\setlength{\gothamTallItemSpace}{#1\gothamTallItemBaseUnit}%
	\vspace{\gothamTallItemSpace}\item}

% For a line that needs \tallitem's own extra headroom but no bullet at
% all -- a display equation continuing the previous \item's own sentence,
% rather than a new bulleted point of its own (see README.md's own
% \tallitem section, which already documents this exact combination as a
% rare, previously hand-written case: "\vspace{\gothamTallItemBaseUnit}\item[]").
% \tallitemnobullet (2026-07-21, direct instruction) factors that hand-
% written idiom into the same kind of drop-in replacement \tallitem
% already is for its own case, same optional multiplier: \tallitemnobullet[1.5].
\newcommand{\tallitemnobullet}[1][1]{%
	% Recomputes \gothamTallItemBaseUnit itself, fresh, rather than reusing
	% whatever \tallitem's own body last left it at -- since it's now a
	% per-call, live "em"-based value (see \tallitem's own comment above),
	% not the one-time preamble-load-time constant it used to be, relying
	% on call order/a preceding \tallitem to have set it correctly would
	% silently break the moment \tallitemnobullet is used on its own, or
	% at a different font size than whichever \tallitem last ran.
	\setlength{\gothamTallItemBaseUnit}{\gothamTallItemBaseRatio em}%
	\setlength{\gothamTallItemSpace}{#1\gothamTallItemBaseUnit}%
	\vspace{\gothamTallItemSpace}\item[]}

% flalign* (used to left-align a display equation inside an itemize --
% see CONVERSION_NOTES.md's "flalign* for left-aligning..." entry) inherits
% plain LaTeX's \abovedisplayskip/\belowdisplayskip, sized for a full
% document page, not a compact slide itemize -- reads as far too much blank
% space above/below the equation next to the surrounding item text (found
% 2026-07-21, direct instruction; most visible on
% slides/bwi-Pearson-Correlation-Definition-current.tex's own single-line
% equation, but the same oversized gap is present everywhere flalign* is
% used this way, e.g. slides/7se-Portfolio-Risk-Three-Assets-current.tex
% and slides/rb3-Asset-Allocation-Example-CAL-current.tex). \gothamTightDisplaySkips
% tightens both skips -- and their "short" variants, since amsmath picks
% between the two depending on how much text precedes the display, so
% setting only one pair left the gap unfixed whenever the preceding line
% happened to be short -- down to something that reads as part of the
% same list rather than a separate floating display. Usage: wrap the
% *real*, unmodified flalign* in a plain TeX group together with this
% command, `{\gothamTightDisplaySkips\begin{flalign*} ... \end{flalign*}}`,
% rather than a new same-named environment -- tried that first
% (`\newenvironment{gothamtightflalign*}{...\begin{flalign*}}{\end{flalign*}}`)
% and it broke LaTeX's own environment-nesting tracking (`\@currenvir`),
% surfacing as "\begin{flalign*} ... ended by \end{gothamtightflalign*}"
% and a fatal, PDF-less compile on the very first slide that used it. A
% plain group sidesteps that class of bug entirely since it never
% introduces a same-named wrapper environment.
\newlength{\gothamTightDisplaySkip}
\setlength{\gothamTightDisplaySkip}{3pt}
\newcommand{\gothamTightDisplaySkips}{%
	\abovedisplayskip=\gothamTightDisplaySkip
	\abovedisplayshortskip=\gothamTightDisplaySkip
	\belowdisplayskip=\gothamTightDisplaySkip
	\belowdisplayshortskip=\gothamTightDisplaySkip}

% Set logo
% \gothamInstituteLogoPath (a plain \newcommand storing an image path,
% same idiom as \gothamChapterLogoPath below) and \gothamLogoHeight (a
% length) are both defined by the per-university preamble that \input{}s
% this file, before this file is reached -- see e.g.
% Slidemaster_preamble_uni-graz.tex. Whichever source image
% \gothamInstituteLogoPath points at should have any symmetric margin of
% empty transparent pixels around the actual logo artwork trimmed off first
% (see Slidemaster_preamble_uni-graz.tex's own comment on why: an untrimmed
% margin makes \includegraphics's box bigger than the visible logo, so the
% logo's right-hand margin ends up looking bigger than its top/bottom
% margins, since only the latter get absorbed into vertical centering).
% \gothamLogoHeight is the single source of truth for the logo's rendered
% size: it drives the (invisible) in-flow logo below, the white patch's
% width further down, and the visible redrawn logo, so those three can never
% drift out of sync with each other again (that drift -- redrawing a smaller
% logo than the one the patch/position were sized for -- is what caused the
% left/right margins around the logo to come out uneven).

% The tikz node wrapper (remember picture, but not overlay, so it still
% behaves as normal inline content) records this logo's exact page position
% under the name "gothamLogoAnchor", so it can be redrawn on top of a white
% patch further down (see the "yellow box shorter than the page" block below).
\renewcommand{\gothamInstituteLogoSquare}[1][\gothamLogoHeight]{%
	\tikz[remember picture, baseline=(gothamLogoAnchor.base)]{%
		\node[inner sep=0pt, outer sep=0pt](gothamLogoAnchor){%
			\includegraphics[height=#1]{\gothamInstituteLogoPath}%
		};%
	}%
}

% \gothamFrametitleRightpading (left/right padding around the logo, chosen
% to match its top/bottom margins -- measured against the title colorbox's
% own top and bottom edges -- rather than the theme's 2.2ex default) is
% also set by the per-university preamble, not here: the top/bottom margins
% are a side effect of vertically centering a \gothamLogoHeight-tall logo in
% the title bar, not something set directly, so this value has to be
% re-measured and set per institution whenever its own logo/height differs
% from another's.

% Logo width, used below to size the white patch behind it.
\newlength{\gothamLogoWidth}
\settowidth{\gothamLogoWidth}{\includegraphics[height=\gothamLogoHeight]{\gothamInstituteLogoPath}}

% Let the yellow title bar visually stop short of the page's right edge, so
% the logo sits on white instead of on yellow: paint a white patch spanning
% the full height of the title bar behind the logo, then redraw the logo
% (already positioned correctly by \gothamFrametitleRightpading above, via
% the "gothamLogoAnchor" node captured above) on top of it. This needs two
% compilation passes: like any \tikz "remember picture", the node positions
% are only available in the .aux file written by the previous run.
% (The redraw uses \includegraphics directly rather than
% \usebeamertemplate{instituteLogo}/{includeLogo}: those templates rely on
% gotham's internal property-list lookup, which is only valid in the
% context of the original frametitle rendering and silently produces
% nothing when re-invoked here.)
% Additionally paints white across the *entire* width of the strip between
% the logo's bottom edge and the title bar's actual bottom edge, so the
% yellow ends flush with the logo's bottom (rather than extending further
% down under the title text too) while the box's width and top position
% stay unchanged. Checked against measured PDF coordinates that the title
% text's own (font-metric, not just this-string's-glyphs) lowest descender
% still clears the logo's bottom edge by about 1.3pt, so this doesn't clip
% text even on titles with descenders (e.g. "Highlighting...").
%
% Factored into its own macro, rather than written directly into the
% \addtobeamertemplate call below, so \gotham_no_inline_logo (a whole
% separate frametitle template variant, further down) can reuse it too --
% \addtobeamertemplate patches whatever template is *currently* installed
% at the time it's called (gotham's own default one here), so a template
% defined fresh afterwards doesn't inherit that patch automatically and
% needs to invoke this explicitly instead.
% Every edge-anchored node in this hook (and in every other page-corner
% overlay in this file) needs an explicit "outer sep=0pt" alongside its
% "inner sep=0pt": tikz's default outer sep is 0.5\pgflinewidth (0.2pt at
% the default 0.4pt line width), and a node's edge/corner anchors sit on
% the *outer* boundary -- outer sep past the filled rectangle -- so an
% anchor placed exactly on a page corner leaves the fill itself 0.2pt
% short of that corner on the anchored sides. Found 2026-07-16 (direct
% instruction: a hairline white line above/left of titleframe's and
% quoteframe's yellow boxes, their right edges 0.4pt off the standard
% title bar's own visible yellow edge, and their logos 0.2pt left of
% every other frame's) by measuring the compiled PDF's drawn rectangles:
% every affected fill/image was inset by exactly 0.199bp = 0.2pt from
% its anchor. This is also the same offset the footer box's own comment
% (further down) had earlier measured, worked around with a hand-nudged
% anchor, and left as "root cause not chased further" -- now chased:
% it *is* a general property of every edge/corner-anchored tikz node,
% not something specific to that one construction. (Center anchors are
% unaffected -- outer sep is symmetric around the center -- so the
% gothamLogoAnchor.center redraw below doesn't need it.)
\newcommand{\gothamFrametitleHookCode}{%
	\tikz[remember picture, overlay]\node(gothamTitleBottom){};%
	\begin{tikzpicture}[remember picture, overlay]
		\path let \p1=(current page.north east), \p2=(gothamTitleBottom) in
			node[fill=white, anchor=north east, inner sep=0pt, outer sep=0pt, name=gothamMainPatch,
			     minimum width=\gothamLogoWidth+2\gothamFrametitleRightpading,
			     minimum height=\y1-\y2] at (\x1,\y1) {};
		\ifx\gothamFrametitleFullHeightYellow\empty
			\path let \p1=(current page.north east), \p2=(gothamTitleBottom),
			          \p3=(gothamLogoAnchor.south), \p4=(current page.west) in
				node[fill=white, anchor=north west, inner sep=0pt, outer sep=0pt,
				     minimum width=\x1-\x4, minimum height=\y3-\y2] at (\x4,\y3) {};
		\fi
		% Normally redrawn at gothamLogoAnchor's own position (wherever the
		% invisible in-flow logo -- on the same line as the title text --
		% ended up). \gothamFrametitleLogoFixed instead redraws it at a fixed
		% page-corner position, independent of where the title text sits --
		% for a frame type (sectionframe) whose title is deliberately *not*
		% at the position a normal frame's title/logo line would put it.
		% gothamMainPatch already covers the full logo column from the page
		% top down to the box's true bottom edge (above), so it's already
		% blank there regardless of which position is used.
		\ifx\gothamFrametitleLogoFixed\empty
			\node[anchor=center, inner sep=0pt] at (gothamLogoAnchor.center) {\includegraphics[height=\gothamLogoHeight]{\gothamInstituteLogoPath}};
		\else
			% inner ysep is derived, not \gothamFrametitleRightpading again:
			% the inline logo this fixed redraw must line up with is
			% *vertically centered* in the title bar, and its resulting top
			% margin (8.14pt, measured) is not exactly the horizontal
			% padding (0.28cm = 7.97pt) -- close enough that the pre-outer-sep-fix
			% code's accidental extra 0.2pt masked the difference, but a
			% visible micro-jump when flipping pages once the fix landed.
			% Rather than a fresh calibrated constant:
			% \gothamTitleBarBottomAlign (the visible bar bottom) *is* the
			% inline logo's bottom edge, so subtracting \gothamLogoHeight
			% reproduces its top margin exactly, and re-measuring one
			% constant keeps both in sync if the bar ever changes.
			\node[anchor=north east, inner xsep=\gothamFrametitleRightpading,
			      inner ysep=\dimexpr\gothamTitleBarBottomAlign-\gothamLogoHeight\relax,
			      outer sep=0pt] at (current page.north east) {\includegraphics[height=\gothamLogoHeight]{\gothamInstituteLogoPath}};
		\fi
		% Optional per-frame "chapter logo" (see \chapterlogo{} below): drawn
		% immediately to the left of the main logo's white patch, on its own
		% white patch, only when \gothamChapterLogoPath is non-empty. Sized
		% with exactly one \gothamFrametitleRightpading of its own padding, on
		% its outer (left) side only -- the inner (right) side is left flush
		% against gothamMainPatch, since that patch already contributes its
		% own \gothamFrametitleRightpading of padding around the main logo.
		% This makes both gaps -- chapter logo to main logo, and chapter logo
		% to where the yellow bar resumes -- come out equal, at exactly one
		% \gothamFrametitleRightpading each, without double-counting or
		% needing a second length constant.
		\ifx\gothamChapterLogoPath\empty\else
			\settowidth{\gothamChapterLogoWidth}{\includegraphics[height=\gothamLogoHeight]{\gothamChapterLogoPath}}%
			% The patch's right edge is nudged 0.5pt *past* gothamMainPatch's own
			% west edge (i.e. slightly overlapping it) rather than exactly
			% touching it -- both are solid white, so the overlap itself is
			% invisible, but touching them exactly left a hairline yellow seam
			% at the join (a real, if sub-millimetre, gap that persisted across
			% recompiles, most likely PDF fill-antialiasing at the shared edge
			% of two separately-filled abutting shapes). The logo image itself
			% (below) is *not* shifted, so the actual margin it defines is
			% unaffected.
			\path let \p1=(current page.north east), \p2=(gothamTitleBottom) in
				node[fill=white, anchor=east, inner sep=0pt, outer sep=0pt,
				     minimum width=\gothamChapterLogoWidth+\gothamFrametitleRightpading+0.5pt,
				     minimum height=\y1-\y2] at ($(gothamMainPatch.west)+(0.5pt,0)$) {};
			\node[anchor=east, inner sep=0pt, outer sep=0pt] at (gothamMainPatch.west |- gothamLogoAnchor.center)
				{\includegraphics[height=\gothamLogoHeight]{\gothamChapterLogoPath}};
		\fi
	\end{tikzpicture}%
}
\addtobeamertemplate{frametitle}{}{\gothamFrametitleHookCode}

% \chapterlogo{path} opts a single upcoming standardframe into an extra logo
% (e.g. a topic/chapter icon) next to the main institute logo -- see the
% frametitle hook above for how it's drawn. Call it *before* \begin{standardframe},
% since the frame title (and this hook) is typeset as part of \begin{frame}
% itself, before any frame-body content would run. standardframe resets this
% back to empty on \end{frame} (below), so it only ever applies to the one
% frame that follows the \chapterlogo{} call, never leaking into the next slide.
% Both the set and the reset use \global\def (checked empirically, via
% \typeout, against the frametitle hook's actual view of this macro on a
% following frame): beamer's frame environment opens its own TeX group around
% the frame body, so a plain (local) assignment done in standardframe's
% closing code -- which runs right after \end{frame} -- was silently undone
% once that group actually closed, reviving the *previous* value instead of
% the reset. \global\renewcommand alone wasn't enough to fix it (still leaked
% into the next frame) -- only dropping to the lower-level \global\def did.
\newlength{\gothamChapterLogoWidth}
\newcommand{\gothamChapterLogoPath}{}
\newcommand{\chapterlogo}[1]{\global\def\gothamChapterLogoPath{#1}}

% \gothamFrametitleFullHeightYellow: opt-out for the frametitle hook's
% "trim yellow to the logo's bottom edge" white patch (see the hook above) --
% empty (default) keeps that trim; non-empty skips it, letting the yellow
% show all the way down to the title box's own true bottom edge instead.
% That trim exists to erase a normally-tiny (~1-2pt) sliver of yellow that
% would otherwise stick out below the logo on an ordinary short title bar --
% harmless there, but sectionframe's deliberately tall bar turns that
% "tiny sliver" into most of the box, silently wiping out nearly all of the
% intended yellow with white. sectionframe sets this (globally, same
% set/reset reasoning as \gothamChapterLogoPath above); standardframe resets
% it back to empty on \end{frame}, so it's opt-in per frame like the others.
\newcommand{\gothamFrametitleFullHeightYellow}{}

% \gothamFrametitleLogoFixed: opt-in for the frametitle hook's logo redraw
% (see the hook above) to use a fixed page-corner position instead of
% wherever the title text's own line puts it. Normally the two coincide (the
% logo sits inline with the title, gotham's own design), which is exactly
% why this is off by default -- but sectionframe's title is deliberately
% *not* where a normal frame's title sits (bottom-anchored within a much
% taller bar), and letting the logo follow it down there would put it
% somewhere no other slide in the deck ever has it. Same set/reset reasoning
% and mechanism as \gothamChapterLogoPath/\gothamFrametitleFullHeightYellow
% above.
\newcommand{\gothamFrametitleLogoFixed}{}

% \gothamFooterBoxColor: the footer box's own fill color (see the footline
% hook further down), overridable per frame the same set-before/reset-after
% way as the flags above -- default \gothamFooterBoxColorDefault (set by the
% per-university preamble before \input{}ing this file, same "value
% supplied externally, referenced lazily" idiom as \gothamBulletColor/
% \gothamInstituteLogoPath), matching every frame's title bar by default;
% highlightframe overrides it to match its own (likewise overridable) title
% bar color, \gothamHighlightColor, so both boxes change together with a
% single retargeting.
\newcommand{\gothamFooterBoxColor}{\gothamFooterBoxColorDefault}

% A frametitle template variant that drops the inline institute logo
% entirely -- for use together with \gothamFrametitleLogoFixed (which
% redraws the logo separately, at a fixed position, making the inline one
% pointless anyway). gotham's own template (beamerouterthemegotham.sty,
% "default/gotham~subsameline") puts \nolinebreak\hfill{...logo...} directly
% after \insertframetitle; this variant drops that whole part.
%
% Doing so, though, means the logo no longer reserves *any* horizontal room
% on the title's own line the way it used to -- so rightskip here widens
% \gothamFrametitleRightpading to reserve the same width gothamMainPatch
% (the frametitle hook's white patch, drawn afterward) already reserves for
% the logo column: \gothamLogoWidth + 2*\gothamFrametitleRightpading. Without
% this, the title paragraph wraps as if that whole column were free text
% space -- not an actual TeX overfull-hbox condition (no warning, the text
% genuinely fits the *box's own* width), just a title that runs visually
% underneath the white patch and logo drawn on top of it afterward, which
% looks identical to clipped/hidden text. (An earlier, wrong diagnosis of
% this same visual symptom blamed gotham's \nolinebreak forcing an overfull
% line -- checked by testing gotham's original template installed with a
% long \huge title and no changes at all, which reproduces something that
% *looks* the same; but the actual fix needed here turned out to be this
% rightskip, not removing \nolinebreak, which was really addressing a red
% herring from an isolated test that happened not to include a competing
% logo column at all.)
%
% "plus 1fil" on that rightskip keeps every line ragged-right (its own
% natural width, however short of the margin that leaves) rather than fully
% justified -- without it, a wrapped multi-line title's non-last lines get
% stretched to exactly fill the available width, which reads as oddly wide,
% uneven word-gaps on a short title (most noticeable exactly when a title is
% long enough to wrap, i.e. always, for this template).
\makeatletter
\defbeamertemplate{frametitle}{gotham_no_inline_logo}{%
	\nointerlineskip
	\begin{beamercolorbox}[wd=\paperwidth, sep=0pt,
	                        leftskip=\gothamFrametitleLeftpading,
	                        rightskip=\dimexpr\gothamLogoWidth+2\gothamFrametitleRightpading\relax plus 1fil]{frametitle}
		\gotham@frametitlestrut@start
		\insertframetitle
		\ifx\insertframesubtitle\@empty\else
			{\usebeamerfont{framesubtitle}\usebeamercolor[fg]{framesubtitle}\gothamFrameSubtitleSep \insertframesubtitle}%
			\vspace*{\gothamFramesubtitleStrutend}%
		\fi
		\gotham@frametitlestrut@end
	\end{beamercolorbox}
	\gothamFrametitleHookCode
}
\makeatother


% Define slideID as a simple value holder (initially empty)
\newcommand{\slideID}{}%

% Define slideID as a simple value holder (not a command with arguments)
\makeatletter
\newcommand{\setSlideID}[1]{%
	% Grab everything up to the first '-' and store it in \slideID
	\edef\tempa{#1}%
	\expandafter\setSlideID@aux\tempa-\relax
}
\def\setSlideID@aux#1-#2\relax{%
	\renewcommand{\slideID}{#1}%
}
\makeatother

% Define a new \slide command to automatically pick the current version of the slide
% 2026-07-08 (direct instruction): slide filenames dropped their leading
% "s-" (e.g. slides/ab3-Some-Slide-current.tex, not the old
% slides/s-ab3-Some-Slide-current.tex) so slide IDs are directly
% searchable in Explorer -- this macro (and the other three below) is the
% *only* place that ever needed to change for that: every deck's own
% \slide{}/\slideversion{}/\omitslide{}/\omitslideversion{} calls pass just
% "ID-slug", never the filename itself, so renaming every file on disk to
% match needed no changes anywhere else (not the deck files, not
% slide_ids.csv, which only ever stored bare IDs to begin with).
\newcommand{\slide}[1]{%
	\setSlideID{#1}%
	\input{slides/#1-current.tex}%
}

% \slidefinal{ID-slug} -- like \slide{}, but renders the slide at its final,
% fully-revealed overlay state in every deck version, instructor included
% (2026-07-12, direct instruction: show jww's finished build in place of a
% separate recap slide, "irrespective of which deck version is being
% compiled" -- that recap slide has since been deleted entirely, its own ID
% no longer in use). A slide
% built with \conditionalpause[instructor] already renders fully built in
% the during-course/after-course compiles -- that pause is a no-op for
% them, so nothing there ever waits on a reveal step -- the only case
% actually needing help is the instructor deck, which would otherwise still
% step through the build. Locally shadowing \conditionalpause to a no-op
% (same optional-argument signature as the real one, so it still accepts
% e.g. \conditionalpause[instructor] without erroring) for the duration of
% one \slide{} call makes every \conditionalpause in that one inclusion
% inert, regardless of the actual compiled \gothamDeckVersion -- scoped
% inside \begingroup/\endgroup, so it reverts immediately after, and has no
% effect on any *other* \slide{}/\slidefinal{} call to the same ID
% elsewhere in the same deck (e.g. an earlier, normally-building instance of
% the same slide). \deckonly/\deckuncover/\deckvisible are untouched, so
% deck-version-tagged content (as opposed to reveal-step-tagged content)
% still follows its usual per-version rules.
\newcommand{\slidefinal}[1]{%
	\begingroup
		\renewcommand{\conditionalpause}[1][instructor]{}%
		\slide{#1}%
	\endgroup
}

% Load a specific, pinned version of a slide instead of whichever is "current"
% Usage: \slideversion{ID-slug}{M-m-p} loads slides/ID-slug-vM-m-p.tex
\newcommand{\slideversion}[2]{%
	\setSlideID{#1}%
	\input{slides/#1-v#2.tex}%
}

% Student-version decks: \studentversion (called anywhere in a deck, e.g.
% right after % Slidemaster preamble variant for University of Graz decks. Institution-
% specific content only (colors, logo, logo calibration) -- every frame
% type, the deck-version system, footer/hyperlink/footnote styling, tikz
% styles, etc. live in Slidemaster_preamble_base.tex, \input{}ed at the
% bottom of this file, and are shared verbatim across every institution
% this project supports (see Slidemaster_preamble_mci.tex and
% Slidemaster_preamble_uni-ibk.tex for the other two). A deck picks its
% institution by \input{}ing this file instead of one of those, nothing
% else changes.

% Color definition (Graz palette, sourced from the original single-preamble
% version of this project -- these are the University of Graz's own official
% colors, not derived/approximated).
\definecolor{c_primary_dark}{HTML}{000000}
\definecolor{c_primary_highlight}{HTML}{ffd500}
\definecolor{c_primary_gray}{HTML}{c6c6c6}
\definecolor{c_primary_white}{HTML}{FFFFFF}
\definecolor{c_secondary_violet}{HTML}{948dc3}
\definecolor{c_tertiary_violet}{HTML}{c2bddf}
\definecolor{c_secondary_blue}{HTML}{689cce}
\definecolor{c_tertiary_blue}{HTML}{bacce8}
\definecolor{c_secondary_green}{HTML}{9fcb81}
\definecolor{c_tertiary_green}{HTML}{cce2ba}
\definecolor{c_secondary_red}{HTML}{e4684b}
\definecolor{c_tertiary_red}{HTML}{f1ae95}
\definecolor{c_secondary_orange}{HTML}{f6a968}
\definecolor{c_tertiary_orange}{HTML}{fbd0aa}
\definecolor{c_secondary_yellow}{HTML}{ffe063}
\definecolor{c_tertiary_yellow}{HTML}{ffeeaf}
\definecolor{c_tertiary_gray}{HTML}{e3e3e3}
% Not an official Uni Graz color (hence "custom", not "secondary" like the
% palette above) -- a softer near-black than c_primary_dark, used for
% quoteframe's full-bleed dark block, matching the specific dark gray (not
% pure black) ExampleSlides.pptx used for that same block. Also reused for
% hyperlink text color, for the same reason it works for quoteframe's
% attribution text: dark enough for good contrast/legibility on a light
% background.
\definecolor{c_custom_gray}{HTML}{333333}
\definecolor{c_custom_red}{HTML}{C00000}
\definecolor{c_custom_green}{HTML}{00B050}

% Logo + calibration.
% schemes/scheme_uni_graz/Logo_Uni-Graz_4c-cropped.pdf -- official vector artwork,
% user-supplied (schemes/scheme_uni_graz/Logo_Uni-Graz_4c.eps, kept alongside as the
% untouched source), converted to PDF via epstopdf and cropped to its true
% ink bounding box (computed with Ghostscript's own `-sDEVICE=bbox`, same
% treatment as MCI's and Innsbruck's own logos -- see
% Slidemaster_preamble_mci.tex's own comment for the fuller explanation of
% why the EPS's own %%BoundingBox isn't trusted directly). Vector, not
% raster, so it stays crisp at any zoom. Supersedes this file's own earlier
% raster logo (schemes/scheme_uni_graz/Logo Uni Graz 4c-cropped.png, itself cropped
% from schemes/scheme_uni_graz/Logo Uni Graz 4c.png's own symmetric transparent
% margin -- both kept alongside as historical/fallback assets, not deleted).
\newcommand{\gothamInstituteLogoPath}{schemes/scheme_uni_graz/Logo_Uni-Graz_4c-cropped.pdf}
\newlength{\gothamLogoHeight}
\setlength{\gothamLogoHeight}{3ex}
% Re-measured against the compiled PDF for this specific logo image -- see
% Slidemaster_preamble_base.tex's own comment on \gothamFrametitleRightpading
% for why this needs re-measuring per institution rather than reusing
% another one's value.
\setlength{\gothamFrametitleRightpading}{0.28cm}

% Itemize bullet color (2026-07-16, direct instruction): Graz keeps the
% project's original behavior, bullets in the primary highlight color.
\newcommand{\gothamBulletColor}{c_primary_highlight}

% Title bar and footer box/text colors (2026-07-16, direct instruction):
% Graz keeps the project's original behavior throughout -- title/footer
% boxes in the primary highlight color, title/footer text in the primary
% dark color (unlike MCI/Innsbruck, which switch these to their own
% secondary blue/white -- see Slidemaster_preamble_mci.tex's own comment).
\newcommand{\gothamTitleBoxColor}{c_primary_highlight}
\newcommand{\gothamTitleTextColor}{c_primary_dark}
\newcommand{\gothamFooterBoxColorDefault}{c_primary_highlight}
\newcommand{\gothamFooterTextColor}{c_primary_dark}

\input{Slidemaster_preamble_base}
, similar to \sidefootnotes or
% \setfootnotesymbols) switches that compile into the "handout" variant --
% including, as of 2026-07-03, actually forcing beamer's own "handout" mode
% (\gdef\beamer@currentmode{handout}, normally only reachable via the
% \documentclass[handout]{beamer} *class option*, set once at class-loading
% time -- too early for a command like \studentversion, called later from a
% deck's own preamble, to use that route). This was checked, not assumed:
% until this point \studentversion only ever touched \omitslide{}, so
% nothing suppressed \pause/overlay specs -- true only by coincidence, since
% no slide converted so far happens to use them. Late-redefining
% \beamer@currentmode was verified (isolated minimal-beamer tests, plain
% \pause *and* inline specs like \item<2->) to reproduce the same collapse
% -to-one-page, show-the-final-revealed-state behavior as the class option,
% even set well after \documentclass{beamer} -- beamer's overlay decoder
% (beamerbasedecode.sty) reads \beamer@currentmode at frame-processing time,
% not just once at class-option time, so this is safe, not a hack that
% happens to work by luck of file ordering.
% \omitslide{}/\omitslideversion{} are drop-in replacements for
% \slide{}/\slideversion{} that tag a slide as excluded from that variant:
%
%   \omitslide{xyz-Answer-key}                 % instead of \slide{...}
%   \omitslideversion{xyz-Answer-key}{1-0-0}    % instead of \slideversion{...}{...}
%
% In a \studentversion compile, a tagged slide's \input{} is skipped
% entirely -- no frame is produced for it at all (not a blank page, not a
% numbering gap: beamer's frame/slide counters simply never advance for a
% call that never ran). In every other (i.e. normal/instructor) compile, a
% tagged slide loads exactly as it would via \slide{}/\slideversion{}, with
% one difference: \gothamOmittedSlide is set for the duration of that one
% \input{}, so the footer template below can print a marker next to that
% frame's slide ID -- a deliberately visible reminder, while reviewing the
% instructor deck, of which slides won't make it to students. The same tag
% works for both compiles because it lives at the call site in the *deck*
% file (not inside the slide's own content file): the identical slide list
% can serve both the instructor deck and the student handout, e.g.
%
%   % Slidemaster preamble variant for University of Graz decks. Institution-
% specific content only (colors, logo, logo calibration) -- every frame
% type, the deck-version system, footer/hyperlink/footnote styling, tikz
% styles, etc. live in Slidemaster_preamble_base.tex, \input{}ed at the
% bottom of this file, and are shared verbatim across every institution
% this project supports (see Slidemaster_preamble_mci.tex and
% Slidemaster_preamble_uni-ibk.tex for the other two). A deck picks its
% institution by \input{}ing this file instead of one of those, nothing
% else changes.

% Color definition (Graz palette, sourced from the original single-preamble
% version of this project -- these are the University of Graz's own official
% colors, not derived/approximated).
\definecolor{c_primary_dark}{HTML}{000000}
\definecolor{c_primary_highlight}{HTML}{ffd500}
\definecolor{c_primary_gray}{HTML}{c6c6c6}
\definecolor{c_primary_white}{HTML}{FFFFFF}
\definecolor{c_secondary_violet}{HTML}{948dc3}
\definecolor{c_tertiary_violet}{HTML}{c2bddf}
\definecolor{c_secondary_blue}{HTML}{689cce}
\definecolor{c_tertiary_blue}{HTML}{bacce8}
\definecolor{c_secondary_green}{HTML}{9fcb81}
\definecolor{c_tertiary_green}{HTML}{cce2ba}
\definecolor{c_secondary_red}{HTML}{e4684b}
\definecolor{c_tertiary_red}{HTML}{f1ae95}
\definecolor{c_secondary_orange}{HTML}{f6a968}
\definecolor{c_tertiary_orange}{HTML}{fbd0aa}
\definecolor{c_secondary_yellow}{HTML}{ffe063}
\definecolor{c_tertiary_yellow}{HTML}{ffeeaf}
\definecolor{c_tertiary_gray}{HTML}{e3e3e3}
% Not an official Uni Graz color (hence "custom", not "secondary" like the
% palette above) -- a softer near-black than c_primary_dark, used for
% quoteframe's full-bleed dark block, matching the specific dark gray (not
% pure black) ExampleSlides.pptx used for that same block. Also reused for
% hyperlink text color, for the same reason it works for quoteframe's
% attribution text: dark enough for good contrast/legibility on a light
% background.
\definecolor{c_custom_gray}{HTML}{333333}
\definecolor{c_custom_red}{HTML}{C00000}
\definecolor{c_custom_green}{HTML}{00B050}

% Logo + calibration.
% schemes/scheme_uni_graz/Logo_Uni-Graz_4c-cropped.pdf -- official vector artwork,
% user-supplied (schemes/scheme_uni_graz/Logo_Uni-Graz_4c.eps, kept alongside as the
% untouched source), converted to PDF via epstopdf and cropped to its true
% ink bounding box (computed with Ghostscript's own `-sDEVICE=bbox`, same
% treatment as MCI's and Innsbruck's own logos -- see
% Slidemaster_preamble_mci.tex's own comment for the fuller explanation of
% why the EPS's own %%BoundingBox isn't trusted directly). Vector, not
% raster, so it stays crisp at any zoom. Supersedes this file's own earlier
% raster logo (schemes/scheme_uni_graz/Logo Uni Graz 4c-cropped.png, itself cropped
% from schemes/scheme_uni_graz/Logo Uni Graz 4c.png's own symmetric transparent
% margin -- both kept alongside as historical/fallback assets, not deleted).
\newcommand{\gothamInstituteLogoPath}{schemes/scheme_uni_graz/Logo_Uni-Graz_4c-cropped.pdf}
\newlength{\gothamLogoHeight}
\setlength{\gothamLogoHeight}{3ex}
% Re-measured against the compiled PDF for this specific logo image -- see
% Slidemaster_preamble_base.tex's own comment on \gothamFrametitleRightpading
% for why this needs re-measuring per institution rather than reusing
% another one's value.
\setlength{\gothamFrametitleRightpading}{0.28cm}

% Itemize bullet color (2026-07-16, direct instruction): Graz keeps the
% project's original behavior, bullets in the primary highlight color.
\newcommand{\gothamBulletColor}{c_primary_highlight}

% Title bar and footer box/text colors (2026-07-16, direct instruction):
% Graz keeps the project's original behavior throughout -- title/footer
% boxes in the primary highlight color, title/footer text in the primary
% dark color (unlike MCI/Innsbruck, which switch these to their own
% secondary blue/white -- see Slidemaster_preamble_mci.tex's own comment).
\newcommand{\gothamTitleBoxColor}{c_primary_highlight}
\newcommand{\gothamTitleTextColor}{c_primary_dark}
\newcommand{\gothamFooterBoxColorDefault}{c_primary_highlight}
\newcommand{\gothamFooterTextColor}{c_primary_dark}

\input{Slidemaster_preamble_base}

%   \FooterText{...}
%   \makeatletter\ifdefined\HANDOUT\studentversion\fi\makeatother  % or similar
%   \input{course-content.tex}   % the shared \slide{}/\omitslide{} list
%
% with two thin wrapper decks (one plain, one calling \studentversion first)
% both \input-ing the same shared content file -- see README.md.
\newif\ifgothamStudentVersion
% \makeatletter must wrap the \newcommand itself, not sit *inside* the macro
% body: catcodes apply at tokenization time, which for a macro body happens
% once, immediately, when \newcommand reads it -- not deferred until the
% macro is later called. \makeatletter placed inside the body (tried first)
% was too late: \beamer@currentmode had already been mis-tokenized as
% "\beamer" + the six-way-otherwise-ordinary character "@" + eleven plain
% letters "currentmode", by the time that inner \makeatletter could act on
% it, so the \gdef silently created/redefined the wrong thing and never
% touched beamer's actual mode variable. Confirmed via \meaning\beamer@
% currentmode before/after \studentversion -- unchanged with the inside-the-
% body version, fixed once \makeatletter moved out here.
\makeatletter
\newcommand{\studentversion}{%
	\global\gothamStudentVersiontrue
	\gdef\beamer@currentmode{handout}%
}
\makeatother

% Deck version system - new boolean flags and string for current version
\newif\ifgothamInstructorVersion
\newif\ifgothamDuringCourseVersion
\newif\ifgothamAfterCourseVersion
\gdef\gothamDeckVersion{instructor}
\gdef\gothamDeckVersionText{}

% \deckversion[] - set the deck version(s)
% Usage: \deckversion[instructor], \deckversion[during-course], etc.
% Multiple versions: \deckversion[instructor,during-course,after-course]
% Parses #1 as a comma-separated list via etoolbox's \forcsvlist (rather
% than the single \ifstrequal{whole-list}{"instructor"} check this used to
% do, which only ever matched a single-version list -- a multi-version
% \deckversion[instructor,during-course] silently set none of the three
% booleans, since the compared string was the literal, unsplit
% "instructor,during-course"). Note \forcsvlist, not etoolbox's similarly-
% named \forlistloop: \forlistloop iterates etoolbox's own internal
% pipe-delimited list format (built via \listadd), not a plain comma string
% like this command's own argument -- confirmed directly (isolated test),
% passing a raw comma string to it doesn't error or split on the commas, it
% just treats the whole string as a single one-item "list".
\makeatletter
\newcommand{\deckversion}[1][instructor]{%
	% \xdef (not \gdef): #1 needs to be fully expanded here, not just
	% stored as whatever token(s) were passed -- \deckversion[\myvar]
	% (compile-all-versions.bat's own override mechanism, via a wrapper
	% file's \def\deckversionoverride{...}) passes a macro, not a literal
	% string, and everywhere else in this deck-version system that reads
	% \gothamDeckVersion only ever applies one \expandafter to it (see
	% \ifdeckversioninlist's own comment) -- enough to unwrap \gothamDeckVersion
	% itself, but not a further macro still sitting inside it. Confirmed
	% directly: with the old \gdef, \deckversion[\deckversionoverride]
	% left \gothamDeckVersionText empty and every \ifdeckversioninlist
	% check false, since \gothamDeckVersion's own stored content was the
	% single unexpanded token \deckversionoverride, not "during-course".
	\xdef\gothamDeckVersion{#1}%
	\global\gothamInstructorVersionfalse
	\global\gothamDuringCourseVersionfalse
	\global\gothamAfterCourseVersionfalse
	\gdef\gothamDeckVersionText{}%
	\def\do@deckversions##1{%
		\ifstrequal{##1}{instructor}{\global\gothamInstructorVersiontrue\gdef\gothamDeckVersionText{instructor deck}}{}%
		\ifstrequal{##1}{during-course}{\global\gothamDuringCourseVersiontrue\gdef\gothamDeckVersionText{student course deck}}{}%
		\ifstrequal{##1}{after-course}{\global\gothamAfterCourseVersiontrue\gdef\gothamDeckVersionText{student final deck}}{}%
	}%
	% \forcsvlist itself splits on literal commas in its own raw,
	% unexpanded argument (confirmed directly against etoolbox's own
	% source, \etb@defparser@do@aux) -- passing #1 straight through here
	% has the exact same problem as the \gdef above did, one call site
	% later. Building the whole \forcsvlist{...}{...} call inside an
	% \edef (protecting \forcsvlist/\do@deckversions themselves with
	% \noexpand so they're invoked *later*, once this constructed call is
	% run, not expanded away here) expands only \gothamDeckVersion --
	% already fully resolved by the \xdef above -- into the list argument.
	\edef\gotham@deckversion@call{\noexpand\forcsvlist{\noexpand\do@deckversions}{\gothamDeckVersion}}%
	\gotham@deckversion@call
}
\makeatother

\newif\ifgothamOmittedSlide
\newcommand{\gothamOmittedSlideMarker}{ | omitted from student deck}
% \stepcounter{framenumber} in the skipped branch below keeps the student
% deck's slide numbers matching the instructor deck's, gap and all, so "see
% slide 12" means the same physical slide to both audiences: beamer steps
% "framenumber" exactly once per \begin{frame} (checked in
% beamerbaseframe.sty -- once per *frame*, not per overlay page, so this is
% right regardless of how many overlay steps the omitted slide would have
% had), and this is the one place that increment is skipped along with the
% frame itself, so it has to be replayed by hand.
\newcommand{\omitslide}[1]{%
	\ifgothamStudentVersion
		\stepcounter{framenumber}%
	\else
		\global\gothamOmittedSlidetrue
		\setSlideID{#1}%
		\input{slides/#1-current.tex}%
		\global\gothamOmittedSlidefalse
	\fi
}
\newcommand{\omitslideversion}[2]{%
	\ifgothamStudentVersion
		\stepcounter{framenumber}%
	\else
		\global\gothamOmittedSlidetrue
		\setSlideID{#1}%
		\input{slides/#1-v#2.tex}%
		\global\gothamOmittedSlidefalse
	\fi
}

% Deck version-aware overlay commands
% Syntax: \deckonly<overlay>[versions]{content}
% If [versions] omitted, defaults to all decks (instructor,during-course,after-course)
%
% Built on etoolbox's own \forcsvlist for the comma-splitting (rather than a
% hand-rolled parser: an earlier version tried to split a list at each comma
% using \newcommand's plain, non-delimited argument grabbing
% (\newcommand{\ifdeckversioninlist@split}[3]{...} called as
% \ifdeckversioninlist@split#1\@nil{#1}) -- but a plain \newcommand doesn't do
% delimited-parameter matching against \@nil the way a \def with a matching
% pattern would; it just grabs the next 3 tokens/groups one at a time. So
% e.g. "instructor" silently split into single characters ("i", "n", "s",
% ...) instead of being read as one list item, which eventually ran the
% parser off the end of its own input and surfaced as "Undefined control
% sequence" wherever the first \conditionalpause[...]/\deckonly[...] in a
% deck happened to be used -- confirmed directly against the compile log,
% traced to slide y5z-CoCo-Bonds's own \conditionalpause[instructor].
%
% NOT etoolbox's own \ifinlist{item}{list}, despite the name similarity: that
% one tests membership in etoolbox's own internal list format, built only via
% \listadd (items joined with a literal "|", not ","); fed a plain comma
% string instead, confirmed directly (isolated test) to silently return
% false for every item, including an exact single-item match.
%
% The membership test itself needs one more subtlety: #1 here is the deck
% version to search for -- always \gothamDeckVersion at every call site
% below, i.e. a macro, not a literal string -- so comparing it via a plain
% \ifstrequal{##1}{#1} does not work either: etoolbox's \ifstrequal compares
% via \detokenize, which (confirmed directly, isolated test) stringifies an
% unexpanded macro token to its own control-sequence *name* ("gothamDeckVersion"),
% not the current value that macro holds. \expandafter forces that expansion
% before \ifstrequal ever sees it.
\makeatletter

\let\gotham@original@only\only
\let\gotham@original@uncover\uncover
\let\gotham@original@visible\visible

\newif\ifgotham@deckversion@found

% Usage: \ifdeckversioninlist{version}{list}{true-branch}{false-branch}
\newcommand{\ifdeckversioninlist}[2]{%
	\global\gotham@deckversion@foundfalse
	\def\gotham@deckversion@check##1{%
		\expandafter\ifstrequal\expandafter{#1}{##1}{\global\gotham@deckversion@foundtrue}{}%
	}%
	\forcsvlist{\gotham@deckversion@check}{#2}%
	\ifgotham@deckversion@found
		\expandafter\@firstoftwo
	\else
		\expandafter\@secondoftwo
	\fi
}

% \@ifnextchar-based parsing (the same idiom LaTeX2e's own kernel commands
% use for an optional argument that isn't a plain trailing [...]) so the
% overlay spec reads exactly like beamer's own \only<...>, with the deck-
% version list as a further, genuinely optional [...] afterwards:
%   \deckonly<2>{content}                              visible in all versions
%   \deckonly<2>[instructor]{content}                  instructor only
%   \deckonly<2>[during-course,after-course]{content}  during + after
\newcommand{\deckonly}{\@ifnextchar<{\gotham@deckonly@overlay}{\gotham@deckonly@overlay<1->}}
\def\gotham@deckonly@overlay<#1>{\@ifnextchar[{\gotham@deckonly@versions<#1>}{\gotham@deckonly@versions<#1>[instructor,during-course,after-course]}}
\def\gotham@deckonly@versions<#1>[#2]#3{%
	\ifgothamInstructorVersion
		\gotham@original@only<#1>{#3}%
	\else
		\ifdeckversioninlist{\gothamDeckVersion}{#2}{\gotham@original@only<#1>{#3}}{}%
	\fi
}

% Same interface as \deckonly, wrapping \uncover and \visible respectively
\newcommand{\deckuncover}{\@ifnextchar<{\gotham@deckuncover@overlay}{\gotham@deckuncover@overlay<1->}}
\def\gotham@deckuncover@overlay<#1>{\@ifnextchar[{\gotham@deckuncover@versions<#1>}{\gotham@deckuncover@versions<#1>[instructor,during-course,after-course]}}
\def\gotham@deckuncover@versions<#1>[#2]#3{%
	\ifgothamInstructorVersion
		\gotham@original@uncover<#1>{#3}%
	\else
		\ifdeckversioninlist{\gothamDeckVersion}{#2}{\gotham@original@uncover<#1>{#3}}{}%
	\fi
}

\newcommand{\deckvisible}{\@ifnextchar<{\gotham@deckvisible@overlay}{\gotham@deckvisible@overlay<1->}}
\def\gotham@deckvisible@overlay<#1>{\@ifnextchar[{\gotham@deckvisible@versions<#1>}{\gotham@deckvisible@versions<#1>[instructor,during-course,after-course]}}
\def\gotham@deckvisible@versions<#1>[#2]#3{%
	\ifgothamInstructorVersion
		\gotham@original@visible<#1>{#3}%
	\else
		\ifdeckversioninlist{\gothamDeckVersion}{#2}{\gotham@original@visible<#1>{#3}}{}%
	\fi
}

% \conditionalpause[versions] - \pause, but only takes effect in the listed
% deck versions (default: instructor only); a no-op, not even a skipped
% overlay step, in any other version
% Usage: \conditionalpause[instructor,during-course] or \conditionalpause[after-course]
\newcommand{\conditionalpause}[1][instructor]{%
	\ifdeckversioninlist{\gothamDeckVersion}{#1}{\pause}{}%
}

% \deckreveal[step]{content} - the "hidden answer, click-revealed for the
% instructor, absent for during-course, pre-shown for after-course" pattern
% already documented by hand on several exercise slides (e.g. mbz/x7q/37s/
% 3cv's own header comments: "a slide/reveal-step that is hidden in the
% source PowerPoint should, by default, only be made visible on the
% instructor version"). A bare, unconditional \only<N->{content} looks like
% it does this, but doesn't: \only's overlay range registers with beamer's
% own overlay-counting pass regardless of \ifgothamInstructorVersion/
% \conditionalpause -- neither of those touch \only itself, they only guard
% *other* content around it -- so a frame with nothing else establishing a
% second overlay step still gets a real second slide out of a bare
% \only<2->, in *every* deck version, not just the instructor one. Confirmed
% directly (2026-07-21, direct instruction) on 37s/3cv: both a) leaked the
% "hidden" content to during-course as a genuine second slide, the exact
% thing this pattern exists to prevent, and b) showed a blank second bullet/
% cell first on that leaked slide's own step 1 -- \only<2-> hides the
% *content*, not the space it would occupy, so during-course's own step 1
% still had nothing else deciding where the bullet/cell landed.
% \deckreveal fixes both by picking a genuinely different rendering per
% version rather than reusing the same overlay-restricted \only call
% everywhere: instructor gets the real overlay build (\only<step->,
% default step 2, matching \tallitem's own default-multiplier convention);
% during-course gets nothing at all (no token emitted, so beamer's overlay
% pass never sees an overlay spec to register); after-course gets the
% content unconditionally, with no overlay spec of its own either, so it
% doesn't silently force a second slide there either. Content-level only
% (wraps a value/expression, like \only itself) -- doesn't cover an
% itemize \item's own bullet marker, which needs its own overlay spec
% (`\item<step->`) directly on the \item, not wrapped from outside; see
% 37s's own comment for that variant, hand-written rather than factored
% into a second macro since only one call site needs it so far.
\newcommand{\deckreveal}[2][2]{%
	\ifgothamDuringCourseVersion
	\else
		\ifgothamInstructorVersion
			\only<#1->{#2}%
		\else
			#2%
		\fi
	\fi
}

% \conditionalslide[versions]{ID-slug} - \slide{}, but only included in the
% listed deck versions; keeps the frame counter in step with the omitted
% version otherwise (same \stepcounter{framenumber} technique \omitslide{}
% already uses for the student deck), so "see slide N" still points at the
% same physical slide across every deck version.
% Usage: \conditionalslide[instructor,after-course]{xyz-Answer-key}
\newcommand{\conditionalslide}[2][instructor,during-course,after-course]{%
	\ifdeckversioninlist{\gothamDeckVersion}{#1}%
		{\setSlideID{#2}%
		 \input{slides/#2-current.tex}}%
		{\stepcounter{framenumber}}%
}

\makeatother

% PDF accessibility: mark purely decorative/running content (footer, below;
% also reused on titleframe's/quoteframe's own decorative background boxes;
% logo not yet included, see README.md's own accessibility section for that
% open question) as a PDF "Artifact", so screen readers skip it -- via
% LuaTeX's raw \pdfextension literal
% primitive. \ifluatex-guarded (iftex is already loaded, via beamer.cls
% itself) because \pdfextension is LuaTeX-only -- this project's README
% documents xelatex as a supported alternative engine, and compiling with it
% hit exactly this: "Undefined control sequence: \pdfextension", hard-failing
% every frame, since the first version of this didn't guard the primitive at
% all. Under any non-LuaTeX engine these two commands now silently do
% nothing instead: decorative content just isn't marked as an Artifact under
% xelatex/pdflatex, rather than breaking the compile -- a real loss of
% accessibility under those engines specifically, not a full equivalent fix,
% but not something to solve today (see also: compiling with plain pdflatex
% already silently mis-renders the whole deck's fonts, so xelatex is really
% the only *other* engine this needs to not break under; pdflatex was never
% viable here to begin with, per "Compiling" in README.md).
% Deliberately NOT tagpdf/kernel tagging (\DocumentMetadata):
% beamer 3.77 refuses to compile at all when that's active ("The beamer
% class is not compatible with \DocumentMetadata ... switch to the new
% ltx-talk class if you require tagged output"), and ltx-talk is a from-
% scratch class (not beamer-based, no theme/template system) that this
% framework can't realistically migrate to. Bypassing beamer's check by
% force-loading the underlying package before \documentclass also failed --
% it gets past hyperref but then hits a genuine internal error in tagpdf's
% LuaTeX backend from the non-standard load order, not something to build
% real coursework on.
%
% /Artifact is a standalone PDF marked-content tag (PDF spec 14.8.2.2):
% valid in *any* PDF, with or without a structure tree, no /Properties
% resource entry or MCID required -- unlike structure-tree tagging, this
% doesn't require kernel tagging support at all. It's a real, working
% partial fix, not full PDF/UA conformance: there's no structure tree, so
% no true semantic headings, and reading order elsewhere on the slide still
% depends on plain content-stream order (verified, not enforced, for
% standardframe -- see its own comment).
% #1 is the artifact's inline PDF dictionary content, e.g.
% "/Type /Pagination /Subtype /Footer" for a running footer, or
% "/Type /Background" for a decorative fill -- not hardcoded to one /Type,
% since not everything decorative here is "Pagination" in the PDF spec's own
% sense of that word (running headers/footers/page numbers/watermarks);
% titleframe's/quoteframe's plain background color fills are semantically
% "Background" artifacts instead, a distinct standard type.
\newcommand{\gothamArtifactBegin}[1]{%
	\ifluatex\pdfextension literal direct{/Artifact <<#1>> BDC}\fi
}
\newcommand{\gothamArtifactEnd}{\ifluatex\pdfextension literal direct{EMC}\fi}

% Footer
% Smaller footer + a bit of padding under it
\setbeamerfont{frame footer}{size=\fontsize{5}{3}\selectfont} % {font size}{baseline skip} in pt

% Shrinks the gap between footnotes and the footer box below them. The
% theme positions the footer's content via \vspace*{\gothamFootlineVOffset}
% (a fixed, negative length) before it starts; making that more negative
% turns out to give footnotes (which are placed just before the footer, and
% stretch to fill whatever room remains via beamer's own
% \beamer@framebottomskip) more room to extend into, pulling them down
% towards the footer -- checked empirically against the PDF's own
% coordinates, this does *not* move the footer's own text at all (that
% still sits exactly where \gothamFooterBoxHeight below was calibrated for),
% so the footer box itself needs no compensating adjustment.
%
% Partly given back (+1.4pt) 2026-07-05: a footnote/source containing an
% underlined \href (see the hyperlink-styling comment above) sits close
% enough to the footer at the original -1.3pt that the underline visually
% touches the footer box (checked directly against the compiled PDF's own
% pixel coordinates: a plain-text footnote's own gap to the footer measured
% ~1.9pt, but an underlined-link footnote's gap measured only ~0.6pt at
% the same setting, since the strut rebuilt for footnotes just below --
% correctly sized for plain text's own natural depth -- doesn't account
% for the underline rule extending further down). +1.4pt splits the
% difference: a linked footnote's gap now measures close to the plain-
% footnote reference (~1.9pt) instead of colliding with the footer, and a
% plain footnote's own gap just becomes a little more relaxed (~3.3pt)
% rather than stretched uncomfortably tight -- applied globally rather
% than only to footnotes with links, since there's no clean per-footnote
% hook for it and the plain-footnote side effect is minor.
%
% Nudged back down by 0.4pt the next day (direct instruction: move
% footnotes down "by the width of the underline line") -- the +1.4pt above
% left a bit more headroom than necessary; 0.4pt is the underline rule's
% own measured stroke thickness (a solid ~1px core at 600dpi, i.e.
% 1/600in \approx 0.12pt, plus antialiasing bleed either side consistent
% with the common 0.4pt LaTeX default rule width -- measured directly
% against the compiled PDF's own pixel values, not assumed).
%
% Reverted the very next round (direct instruction: move back *up* by that
% same underline thickness) -- net effect is back to the plain +1.4pt this
% section started with; kept as two separate historical adjustments below
% rather than collapsing them into one, so this back-and-forth stays
% visible rather than silently disappearing.
\addtolength{\gothamFootlineVOffset}{-1.3pt}
\addtolength{\gothamFootlineVOffset}{1.4pt}
\addtolength{\gothamFootlineVOffset}{-0.4pt}
\addtolength{\gothamFootlineVOffset}{0.4pt}

% 2026-07-06 (direct instruction): move footnotes down by half the
% measured distance between a linked footnote's own underline and the
% footer box, checked on slides/3ic-Corporate-Venture-Capital-
% current.tex specifically (the underline sits at y=245.95pt, the footer
% box's own top edge at y=247.69pt, a 1.74pt gap) -- half of that is
% 0.87pt, applied as a further negative adjustment (more negative =
% footnotes get pulled further down, per this section's own comment
% above).
\addtolength{\gothamFootlineVOffset}{-0.87pt}

% Shifts the right end of the footer line (the frame number) so it lines up
% with the logo's own right edge above it, rather than wherever the theme's
% default footline margins would otherwise put it. Empirically calibrated
% against the rendered PDF coordinates of both; re-measure and adjust if
% \gothamFrametitleRightpading (the logo's own right inset) changes.
\newlength{\gothamFooterRightShift}
\setlength{\gothamFooterRightShift}{4.6pt}

% Nudges the footer text down within its own yellow box, to vertically
% center it there: it previously sat noticeably closer to the box's top
% edge than its bottom (0.14pt above vs. 1.12pt below, checked against the
% rendered PDF's own text coordinates) since \gothamFootlineVOffset's own
% positioning (further up) turns out not to move this text at all (see its
% comment above) -- this is a separate, additive leading \vspace* instead.
\newlength{\gothamFooterTextDrop}
\setlength{\gothamFooterTextDrop}{-0.49pt}

% Compile timestamp: lets you tell, from a PDF alone (no need to check file
% properties/mtime, which don't survive e.g. emailing or uploading a copy),
% whether it reflects the latest edits or a stale earlier build -- handy
% whenever a PDF's *content* is in question, not just when debugging compile
% tooling itself. LuaTeX-only, matching this project's required engine (see
% README.md's "Compiling" section) -- no \ifluatex guard needed since
% nothing else here works under another engine to begin with.
%
% Deliberately avoids os.date's own %-format directives (e.g. "%Y-%m-%d")
% entirely, building the string from os.date("*t")'s field table and plain
% Lua concatenation instead -- \directlua's argument is tokenized by plain
% TeX first, where % starts a comment, before the text ever reaches Lua as a
% string. The obvious-looking fix, escaping each one as \%, does *not* work:
% confirmed directly, it produces a Lua string containing a literal
% backslash followed by a percent sign (detokenization keeps a control
% sequence's own backslash, for \% same as any other), which Lua's own
% string-literal parser then rejects outright ("invalid escape sequence"),
% silently aborting that whole \directlua call -- with no visible LaTeX-side
% error, just a cache file ending up written with an empty value. Switching
% catcode temporarily (\catcode`\%=12 ... ) would also work but needs a
% \begingroup/\endgroup wrapper at every call site; skipping % entirely is
% simpler and needs no such wrapper.
%
% Cached to \jobname.timestamp (via plain TeX \write, so MiKTeX's
% -aux-directory flag -- already applied through .latexmkrc's %O -- routes
% it into .build/ like every other aux file, not left cluttering decks/ or
% the project root) rather than recomputed fresh on every pass: a live
% os.date() call on every pass, tried first, broke latexmk's own rerun
% convergence outright -- confirmed directly, by isolating the change, that
% an otherwise-identical clean compile converges normally without it. Each
% of latexmk's internal reruns relaunches lualatex as a brand-new process,
% and a full pass over a deck this size (font loading alone) comfortably
% crosses a one-minute boundary, so two consecutive passes embedded two
% different strings in the footer; latexmk saw the .aux checksum differ
% *again* on the very next pass and looped ("Latexmk: References changed.")
% until hitting its rerun limit and reporting a false "gave an error", with
% no actual LaTeX error anywhere in the log. Caching to a file read back on
% later passes fixes that: the value is fixed the moment it's first
% written, so the footer's rendered text -- and therefore the .aux
% checksum -- stays identical across every subsequent pass of the same
% compile.
%
% This cache is deliberately *not* self-invalidating by date/content -- it
% keeps reporting an old timestamp indefinitely once written, across
% unrelated future compiles, unless something clears it first.
% SlideManager's compiler.py does exactly that: before every job it deletes
% that job's own .build/<jobname>.* (which includes <jobname>.timestamp), so
% a run through SlideManager always gets a fresh stamp. An ad-hoc TeXstudio
% compile without going through SlideManager first may show a stale stamp if
% .build/ already has one left over from an earlier session -- delete
% .build/<jobname>.timestamp by hand (or just compile that job once through
% SlideManager) if that matters for what you're checking.
\directlua{
	function gothamZeroPad(n) if n < 10 then return "0" .. n else return "" .. n end end
}
\newwrite\gothamCompileTimestampFile
\newcommand{\gothamCompileTimestamp}{%
	\ifdefined\gothamCompileTimestampValue
	\else
		\IfFileExists{\jobname.timestamp}{%
			\input{\jobname.timestamp}%
		}{%
			\xdef\gothamCompileTimestampValue{\directlua{%
				local d = os.date("*t")%
				tex.sprint(d.year .. "-" .. gothamZeroPad(d.month) .. "-" .. gothamZeroPad(d.day) .. " " .. gothamZeroPad(d.hour) .. ":" .. gothamZeroPad(d.min))%
			}}%
			\immediate\openout\gothamCompileTimestampFile=\jobname.timestamp\relax
			\immediate\write\gothamCompileTimestampFile{\string\gdef\string\gothamCompileTimestampValue{\gothamCompileTimestampValue}}%
			\immediate\closeout\gothamCompileTimestampFile
		}%
	\fi
	\gothamCompileTimestampValue
}

% \gothamHideFooterText: opt-in flag suppressing just \FooterText (the left
% part of the footer) -- default false. Set true by the standout template's
% own hook below, so `\begin{frame}[standout]` hides only that left part,
% keeping the "(slide id: ...)" block and frame number visible (2026-07-16,
% direct instruction). Reset immediately after being read, inside the
% footer template itself, rather than relying on a per-environment closing
% hook (standout frames are bare `\begin{frame}[standout]...\end{frame}`,
% not one of this project's own environments like standardframe, so there's
% no custom closing code to reset it in) -- safe because \gothamHideFooterText
% is only ever consumed here, once, immediately before being cleared.
\newif\ifgothamHideFooterText
% gotham's own \defbeamertemplate*{standout}{gotham} (gotham.dtx) is what
% `\begin{frame}[standout]` actually activates; \addtobeamertemplate's
% "after" hook appends to it without needing to touch that upstream file.
\addtobeamertemplate{standout}{}{\global\gothamHideFooterTexttrue}

\setbeamertemplate{frame footer}{%
	\vspace*{\gothamFooterTextDrop}%
	\ifgothamHideFooterText
		\global\gothamHideFooterTextfalse
	\else
		\FooterText%
	\fi
	\hfill%
	{\fontsize{4}{4}\selectfont (slide id: \slideID\ifgothamOmittedSlide\gothamOmittedSlideMarker\fi\ | \gothamDeckVersionText\ | \gothamCompileTimestamp)}\hspace*{2em}\insertframenumber%
	\hspace*{-4em}%
	\hspace*{-\gothamFooterRightShift}%
	\vspace*{0.4ex}%
}

% Give the footer the same yellow-box treatment as the title bar, but --
% unlike the title bar, which is tall for visual weight -- only as tall as
% the footer's own text plus equal padding above and below.
%
% This paints a yellow patch *behind* the theme's own footline rather than
% replacing that template's content, because \usebeamertemplate{frame~footer}
% (which is what actually typesets \FooterText etc.) only produces its
% content when invoked from the theme's own original footline template --
% calling it again from anywhere else (a redefined footline, the frame body,
% a \settoheight measurement, ...) silently produces nothing. It's
% positioned from "current page.south" (the page's own bottom-left corner)
% rather than a \tikz "remember picture" node tracking the footer's actual
% position, since there's no per-frame content to track: the footer text
% (in its actual 5pt font, per \setbeamerfont{frame footer} above -- despite
% the "page number in head/foot" font switch inside the theme's footline
% template, this is the font that actually ends up applied) sits at a fixed
% distance from the bottom of the page on every frame.
%
% \gothamFooterBoxHeight below is calibrated to that actual measured
% position (checked via the PDF's own text coordinates, not just visually):
% the text sits about 0.84pt above the page's bottom edge and spans about
% one em of its own 5pt font, plus a bit of padding above and below. If the
% footer font size changes, this will need re-measuring.
\newlength{\gothamFooterBoxHeight}
\setlength{\gothamFooterBoxHeight}{7.26pt}
% \vbox to 0pt{...} forces this to occupy exactly zero space in the normal
% vertical list, *without* switching into horizontal mode the way an hbox
% (e.g. \raisebox) would -- inserting an hbox here forces LaTeX to open an
% implicit paragraph before the theme's own \vspace*{\gothamFootlineVOffset}
% runs, which changes how that \vspace behaves (it special-cases being at
% the very start of vertical mode) and nudges the footer text down by a few
% points. A \vbox is already vertical material, so it doesn't trigger that.
% \gothamArtifactBegin/\gothamArtifactEnd (defined above, "PDF accessibility"
% comment) wrap both this box *and* the original footline template's own
% content (the theme's \usebeamertemplate{frame footer} call, i.e.
% \FooterText/slide id/page number) -- \addtobeamertemplate's "after" slot
% runs once the wrapped-around original template content has finished, so
% \gothamArtifactEnd there closes the same marked-content span this begins.
% The footer box's own south-west corner originally measured 0.2pt short
% of the page's actual left/bottom edges on every slide checked
% (2026-07-06, direct instruction: "check whether the footer ... has a bit
% of white on its left and below") -- confirmed directly against the
% compiled PDF's own drawn rectangle (consistently ~0.199pt in from both
% edges, not a one-off), unlike the title bar (gotham's own native
% frametitle background, a completely different mechanism -- normal TeX
% box placement, not a tikz "current page" lookup), which sits exactly
% flush at (0,0). At the time, the root cause wasn't chased further and
% the box was fixed by hand-nudging the anchor 0.2pt past the corner and
% padding the size to match. Root cause since found (2026-07-16, while
% chasing the identical hairline on titleframe/quoteframe -- see the
% frametitle hook's own comment above): tikz's default outer sep
% (0.5\pgflinewidth = 0.2pt) puts every edge/corner anchor 0.2pt outside
% the filled rectangle, on *every* edge-anchored node, so the earlier
% "nothing else anchored to a page corner shows the same offset" turned
% out to be wrong. "outer sep=0pt" below is the real fix, replacing the
% nudge-and-pad workaround; the box now lands exactly flush by
% construction rather than by measured compensation.
\addtobeamertemplate{footline}{%
	\gothamArtifactBegin{/Type /Pagination /Subtype /Footer}%
	\vbox to 0pt{%
	\vss
	\begin{tikzpicture}[remember picture, overlay]
		\node[fill=\gothamFooterBoxColor, anchor=south west, inner sep=0pt, outer sep=0pt,
		      minimum width=\paperwidth, minimum height=\gothamFooterBoxHeight]
			at (current page.south west) {};
	\end{tikzpicture}%
	}%
}{\gothamArtifactEnd}

% Footnotes: small font (smaller than the footer's own 5pt), symbol markers
% (*, dagger, ...) instead of arabic numbers, right-aligned (flush with the
% logo's own right edge, like the footer's frame number), no rule above
% them, and tightened-up line spacing.
%
% Line spacing needs two different levers, not just \baselineskip, because
% at this font size (5pt) \baselineskip's normal value is *smaller* than the
% font's own natural line height, so TeX falls back to \lineskip (a flat
% add-on) between every line regardless -- \baselineskip itself turns out to
% have no effect here at all. That fallback \lineskip is what governs
% spacing *within* a single (wrapped) footnote's own lines: reduced only a
% tiny amount, from the default 1pt.
% The *between*-footnote gap is inflated by a separate mechanism entirely:
% \footnotesep, a rule LaTeX inserts before every individual footnote entry
% (measured to be 7.7pt here -- clearly sized for a much bigger font than
% ours), so that's what actually needs the bigger reduction; \parskip
% (tried first) turned out to have no visible effect on it.
\setlength{\footnotesep}{3pt}
% \reset@font\footnotesize is applied by beamer's own footnote machinery
% before this template runs (beamerbaseframecomponents.sty), so the font
% size has to be set again here to actually take effect.
\renewcommand{\footnoterule}{}
% Same marker sequence as LaTeX's own \fnsymbol (*, dagger, double dagger,
% ...), but built from plain Unicode text characters instead: \fnsymbol's
% symbols are typeset via \ensuremath{...}, i.e. in *math* mode, and the
% theme loads fontspec with the "no-math" option (see
% beamerfontthemegotham.sty) specifically so that Fira Sans -- the font
% fontspec puts everything else in -- is left in charge of text only, not
% math; math mode therefore still falls back to LaTeX's default Computer
% Modern math fonts. \fnsymbol's markers were consequently rendering in a
% different font than the surrounding text (checked by comparing rendered
% glyph shapes) -- most visible as a marker that looked shifted left of
% flush-left text next to it (e.g. a \source{} line, which has no marker of
% its own), since Computer Modern's symbol glyphs have different sidebearings
% than Fira Sans's. Plain characters typeset with no \ensuremath wrapper stay
% in text mode and pick up whatever font is already active at the call site,
% i.e. Fira Sans -- matching the rest of the footnote.
% Literal Unicode characters, not the \textdagger/\textdaggerdbl/\textbardbl
% kernel commands: those need the textcomp/TS1 setup to render as plain text,
% and silently fall back to \ensuremath{...} (i.e. math mode, the very thing
% this replaces \fnsymbol to avoid -- see the comment above) when that isn't
% loaded. LuaLaTeX plus fontspec maps raw UTF-8 input straight to the current
% font's (Fira Sans's) own glyphs with no encoding layer in between, which is
% the entire point of using them -- so writing the characters directly, with
% no wrapper command, guarantees they stay in text mode and stay in Fira Sans.
%
% The sequence itself is user-configurable via \setfootnotesymbols{...} (a
% comma-separated list, e.g. \setfootnotesymbols{*,+,#}) -- call it anywhere
% from that point in the deck onward (e.g. once, right after
% % Slidemaster preamble variant for University of Graz decks. Institution-
% specific content only (colors, logo, logo calibration) -- every frame
% type, the deck-version system, footer/hyperlink/footnote styling, tikz
% styles, etc. live in Slidemaster_preamble_base.tex, \input{}ed at the
% bottom of this file, and are shared verbatim across every institution
% this project supports (see Slidemaster_preamble_mci.tex and
% Slidemaster_preamble_uni-ibk.tex for the other two). A deck picks its
% institution by \input{}ing this file instead of one of those, nothing
% else changes.

% Color definition (Graz palette, sourced from the original single-preamble
% version of this project -- these are the University of Graz's own official
% colors, not derived/approximated).
\definecolor{c_primary_dark}{HTML}{000000}
\definecolor{c_primary_highlight}{HTML}{ffd500}
\definecolor{c_primary_gray}{HTML}{c6c6c6}
\definecolor{c_primary_white}{HTML}{FFFFFF}
\definecolor{c_secondary_violet}{HTML}{948dc3}
\definecolor{c_tertiary_violet}{HTML}{c2bddf}
\definecolor{c_secondary_blue}{HTML}{689cce}
\definecolor{c_tertiary_blue}{HTML}{bacce8}
\definecolor{c_secondary_green}{HTML}{9fcb81}
\definecolor{c_tertiary_green}{HTML}{cce2ba}
\definecolor{c_secondary_red}{HTML}{e4684b}
\definecolor{c_tertiary_red}{HTML}{f1ae95}
\definecolor{c_secondary_orange}{HTML}{f6a968}
\definecolor{c_tertiary_orange}{HTML}{fbd0aa}
\definecolor{c_secondary_yellow}{HTML}{ffe063}
\definecolor{c_tertiary_yellow}{HTML}{ffeeaf}
\definecolor{c_tertiary_gray}{HTML}{e3e3e3}
% Not an official Uni Graz color (hence "custom", not "secondary" like the
% palette above) -- a softer near-black than c_primary_dark, used for
% quoteframe's full-bleed dark block, matching the specific dark gray (not
% pure black) ExampleSlides.pptx used for that same block. Also reused for
% hyperlink text color, for the same reason it works for quoteframe's
% attribution text: dark enough for good contrast/legibility on a light
% background.
\definecolor{c_custom_gray}{HTML}{333333}
\definecolor{c_custom_red}{HTML}{C00000}
\definecolor{c_custom_green}{HTML}{00B050}

% Logo + calibration.
% schemes/scheme_uni_graz/Logo_Uni-Graz_4c-cropped.pdf -- official vector artwork,
% user-supplied (schemes/scheme_uni_graz/Logo_Uni-Graz_4c.eps, kept alongside as the
% untouched source), converted to PDF via epstopdf and cropped to its true
% ink bounding box (computed with Ghostscript's own `-sDEVICE=bbox`, same
% treatment as MCI's and Innsbruck's own logos -- see
% Slidemaster_preamble_mci.tex's own comment for the fuller explanation of
% why the EPS's own %%BoundingBox isn't trusted directly). Vector, not
% raster, so it stays crisp at any zoom. Supersedes this file's own earlier
% raster logo (schemes/scheme_uni_graz/Logo Uni Graz 4c-cropped.png, itself cropped
% from schemes/scheme_uni_graz/Logo Uni Graz 4c.png's own symmetric transparent
% margin -- both kept alongside as historical/fallback assets, not deleted).
\newcommand{\gothamInstituteLogoPath}{schemes/scheme_uni_graz/Logo_Uni-Graz_4c-cropped.pdf}
\newlength{\gothamLogoHeight}
\setlength{\gothamLogoHeight}{3ex}
% Re-measured against the compiled PDF for this specific logo image -- see
% Slidemaster_preamble_base.tex's own comment on \gothamFrametitleRightpading
% for why this needs re-measuring per institution rather than reusing
% another one's value.
\setlength{\gothamFrametitleRightpading}{0.28cm}

% Itemize bullet color (2026-07-16, direct instruction): Graz keeps the
% project's original behavior, bullets in the primary highlight color.
\newcommand{\gothamBulletColor}{c_primary_highlight}

% Title bar and footer box/text colors (2026-07-16, direct instruction):
% Graz keeps the project's original behavior throughout -- title/footer
% boxes in the primary highlight color, title/footer text in the primary
% dark color (unlike MCI/Innsbruck, which switch these to their own
% secondary blue/white -- see Slidemaster_preamble_mci.tex's own comment).
\newcommand{\gothamTitleBoxColor}{c_primary_highlight}
\newcommand{\gothamTitleTextColor}{c_primary_dark}
\newcommand{\gothamFooterBoxColorDefault}{c_primary_highlight}
\newcommand{\gothamFooterTextColor}{c_primary_dark}

\input{Slidemaster_preamble_base}
) to replace the default marker set for every
% \footnote{}/\source{} from then on. \seq_item:Nn (expl3's indexed lookup;
% expl3 is already loaded by the gotham theme itself, see e.g.
% beamerfontthemegotham.sty) is what removes the old \ifcase's hard 9-symbol
% ceiling -- a custom list can be any length. Indexing wraps around (via
% \int_mod:nn) rather than erroring past the end of the list, so an unusually
% footnote-heavy slide degrades to repeating symbols instead of a hard
% failure; slides/ZZU-Footnote-symbol-options-current.tex is a reference
% sheet of candidate Unicode symbols to build a custom list from.
%
% Project-wide default (chosen deliberately, not the classic */dagger/etc.
% set): a currency-sign sequence, in this specific order. All 14 were
% confirmed against slides/ZZU-Footnote-symbol-options-current.tex to
% actually exist in Fira Sans.
\ExplSyntaxOn
\seq_new:N \g_gotham_footnote_symbols_seq
\seq_gset_from_clist:Nn \g_gotham_footnote_symbols_seq { £,¥,₹,₽,₺,₴,₪,₫,₦,₡,₩,ƒ,¢,€ }
\newcommand{\setfootnotesymbols}[1]{%
	\seq_gset_from_clist:Nn \g_gotham_footnote_symbols_seq { #1 }
}
\newcommand{\gothamFnSymbol}[1]{%
	\seq_item:Nn \g_gotham_footnote_symbols_seq
		{ \int_eval:n { \int_mod:nn { #1 - 1 } { \seq_count:N \g_gotham_footnote_symbols_seq } + 1 } }
}
\ExplSyntaxOff
\renewcommand{\thefootnote}{\gothamFnSymbol{\value{footnote}}}
% A \footnote{} called inside a minipage doesn't use \thefootnote/\c@footnote
% at all -- LaTeX's own minipage machinery silently switches to a separate
% counter/display command, \c@mpfootnote/\thempfootnote, normally displayed
% via \@alph (a, b, c, ...). Found 2026-07-14 converting
% slides/j6h-Foreign-Exchange-Market-current.tex (its first slide combining
% \footnote{} with a minipage, since a narrowed text column was added
% alongside a corner chart): the footnote *text* still correctly landed in
% this deck's usual bottom-of-frame collection (beamer's own footnote
% storage isn't minipage-scoped), but the *marker*, both inline and at the
% bottom, showed as a plain "a"/"b"/"c" instead of this project's own
% currency-symbol sequence. Redefined the same way as \thefootnote so any
% future \footnote{} inside a minipage picks up the right symbol too --
% \c@mpfootnote resets to 0 for every new minipage on its own (unlike
% \c@footnote, which needs standardframe's own explicit per-frame reset
% below), which happens to already match this deck's "restart every slide"
% footnote convention for the common case of one minipage per frame.
\renewcommand{\thempfootnote}{\gothamFnSymbol{\value{mpfootnote}}}
\setbeamerfont{footnote}{size=\fontsize{3.2}{1.9}\selectfont}
% Beamer's own \@makefnmark (used for the marker both inline in the body
% text and again next to the footnote itself) hard-codes the "footnote
% mark" beamerfont onto it -- an *absolute* size, unrelated to whatever
% text surrounds it -- which is why the inline markers looked tinier than
% a normal superscript when that was set to this footnote's own small size:
% a fixed size that looks fine amid 3.2pt footnote text looks minuscule
% amid much bigger body text. Reverting to plain LaTeX's own \@makefnmark
% instead makes it a normal \textsuperscript of whatever font size is
% already active at each call site -- proportionally scaled either way, so
% it looks like a normal superscript inline, and still comes out small next
% to the (small) footnote text since \usebeamerfont{footnote} above is
% already active there by the time it's invoked.
\makeatletter
\renewcommand{\@makefnmark}{\hbox{\@textsuperscript{\normalfont\@thefnmark}}}
\makeatother
% \raggedleft's own \rightskip (0pt) is flush with the frame's normal body
% text margin, not the logo's edge above it; this pushes further right by
% the difference between the two, empirically calibrated (like
% \gothamFooterRightShift above) against the rendered PDF coordinates of
% both -- re-measure and adjust if the text margins or logo inset change.
\newlength{\gothamFootnoteRightShift}
\setlength{\gothamFootnoteRightShift}{4.8pt}

% Side footnotes: \sidefootnotes (called *before* \begin{standardframe}, like
% \chapterlogo{}) moves that one frame's footnotes -- including \source{} lines,
% which are footnotes underneath -- from their usual spot above the footer to a
% vertical strip along the right edge of the page, rotated 90 degrees counter-
% clockwise: the text starts just above the footer box and reads bottom-to-top,
% as if the usual bottom-right footnote block had been turned on its side (with
% one deliberate difference: the lines are flush with the strip's bottom end,
% not the top -- see the \raggedright comment in the template below).
% standardframe resets the switch
% on \end{frame} (same \global mechanics as \chapterlogo{} -- see the comment
% there for why a plain local assignment doesn't survive beamer's frame group),
% so it opts in only the single frame that follows, and any frame type derived
% from standardframe inherits the behaviour automatically.
%
% How it works: the footnote *template* below branches on this switch. In side
% mode it typesets each footnote into \gothamSideFootnotesBox (a plain box
% register accumulated across the frame's footnotes) instead of contributing
% anything to beamer's own bottom-of-frame footnote insert, and the footline
% hook further down draws that box -- rotated -- as a tikz overlay once the
% whole frame body has run. The footline is the right place to draw from: it is
% typeset at page shipout, i.e. after every \footnote{} in the body has already
% executed, whereas e.g. the frametitle would run too early and find the box
% empty. (Beamer also typesets the footline once per frame just to *measure*
% it, at \begin{frame} time -- before any footnotes exist -- which is why the
% drawing hook guards on \ifvoid: the measurement pass must not consume, or
% try to draw, a box that doesn't exist yet.)
%
% Known quirk: a footnote with an overlay spec (\footnote<2->{...}) is
% collected -- and therefore shown -- on *every* overlay slide of its frame in
% side mode. Beamer implements covered footnotes by typesetting them invisibly
% in place, which relocating the text into a box register defeats.
\newif\ifgothamSideFootnotes
\newcommand{\sidefootnotes}{\global\gothamSideFootnotestrue}
\newbox\gothamSideFootnotesBox
% The rotated block's line width, i.e. the vertical run the strip may occupy
% down the right edge of the page: from just below the title bar to just above
% the footer's yellow box. Like \gothamFooterBoxHeight above, calibrated
% against the measured PDF coordinates: the title bar's bottom edge sits
% 219.57pt above the page's bottom edge; minus \gothamSideFootnoteTopSep (3pt),
% \gothamFooterBoxHeight (7.26pt), and ~2pt of clearance above the footer box
% leaves 207pt. Re-measure if the frametitle font or its padding ever changes.
\newlength{\gothamSideFootnoteHsize}
\setlength{\gothamSideFootnoteHsize}{207pt}
% Gap between the title bar's bottom edge and the top end of the strip. The
% text itself sits at the strip's *bottom* end (it is \raggedright, i.e.
% flush with the line start, which the rotation turns into bottom-aligned),
% so this gap only matters for footnotes long enough to fill the whole run.
\newlength{\gothamSideFootnoteTopSep}
\setlength{\gothamSideFootnoteTopSep}{3pt}
% Distance from the page's right edge to the strip's outer edge -- the rotated
% counterpart of the small gap the bottom footnotes keep above the footer box.
\newlength{\gothamSideFootnoteRightInset}
\setlength{\gothamSideFootnoteRightInset}{2pt}
% A real footnote's marker (e.g. "*") sits in an \hbox of its own whose
% *ink* -- unlike \source{}'s marker, which is emptied to nothing and so has
% no ink at all -- extends slightly left of that hbox's own nominal left
% edge (ordinary optical behaviour for a small centred glyph like an
% asterisk, exaggerated further by \@makefnmark's superscript raise). With
% every footnote's paragraph starting flush at the same \noindent left edge,
% that overhang made real footnotes' markers visibly poke out to the left of
% \source{}'s flush-left text above/below them. This shift nudges real
% footnotes' *entire* line (marker, comma, and text together) right by that
% overhang, so the marker's ink -- not the text after it -- lines up with
% \source{}'s left edge; \source{} itself is untouched; empirically
% calibrated against the rendered PDF, like the other shift constants here.
\newlength{\gothamSideFootnoteMarkerShift}
\setlength{\gothamSideFootnoteMarkerShift}{1.3pt}

% A new template variant (rather than reusing beamer's own "default", which
% it has already defined and then refuses to redefine) that we immediately
% activate below.
% Wrapped in \makeatletter (unlike the rest of this section) because the side
% branch below tests \@thefnmark against \@empty to tell a real footnote
% (non-empty marker) apart from a \source{} call (marker forced empty) --
% see \gothamSideFootnoteMarkerShift above.
\makeatletter
\defbeamertemplate{footnote}{gothamfn}{%
	\ifgothamSideFootnotes
		% Side mode (see \sidefootnotes above): append this footnote to the
		% collector box instead of typesetting it here. \insertfootnotetext
		% carries beamer's own \rule\z@\footnotesep lead-in and
		% \@finalstrut\strutbox tail along with it (beamer bakes both into the
		% text it hands the template), so the between-footnote spacing and the
		% strutbox rebuild work exactly as in the bottom branch below -- and,
		% equally, nothing of the footnote (not even the spacing rule) is left
		% behind in beamer's bottom-of-frame insert, which therefore stays
		% visually empty. No \gothamFootnoteRightShift here: that constant
		% compensates the *bottom* block's \raggedleft for the text-margin-vs-
		% logo-edge difference, while the strip's position is set directly by
		% the tikz node in the footline hook.
		\global\setbox\gothamSideFootnotesBox\vbox{%
			\unvbox\gothamSideFootnotesBox
			\hsize\gothamSideFootnoteHsize
			\usebeamerfont{footnote}%
			\setlength{\lineskip}{0.6pt}%
			\settoheight{\dimen0}{Xy}%
			\settodepth{\dimen1}{Xy}%
			\setbox\strutbox\hbox{\vrule height\dimen0 depth\dimen1 width0pt}%
			\parindent 1em\noindent%
			% \raggedright, not the bottom block's \raggedleft: after the
			% rotation, the line start is the *bottom* end of the strip, so
			% this makes the text begin just above the footer box and run
			% upward (rather than hang down from the title bar).
			\raggedright
			% \source{} forces \thefootnote (hence \@thefnmark) empty, so this
			% only fires for a real footnote's marker -- see
			% \gothamSideFootnoteMarkerShift above for why it needs the shift
			% and \source{} doesn't. The "\," below (a separator between the
			% marker and the footnote text) needs the same guard: it used to
			% run unconditionally, so an empty-marker \source{} still got a
			% thin space's worth of leading indent before its own text --
			% invisible under the *bottom* block's \raggedleft (nothing there
			% depends on the left edge), but a real, visible defect here under
			% \raggedright, where the left edge is exactly what's on display:
			% a wrapped \source{}'s first line sat measurably indented from
			% its own second line, which never picks up that one-time leading
			% space (2026-07-12, direct instruction: "the first line seems
			% indented compared to the second, which it shouldn't be").
			\ifx\@thefnmark\@empty
				\insertfootnotetext\par%
			\else
				\hspace*{\gothamSideFootnoteMarkerShift}\insertfootnotemark\,\insertfootnotetext\par%
			\fi
		}%
	\else
	\usebeamerfont{footnote}%
	\setlength{\lineskip}{0.6pt}%
	% beamer's footnote machinery resets to plain \footnotesize (a much
	% bigger font) before this template runs, and switching to our own
	% smaller font via \usebeamerfont doesn't update \strutbox to match --
	% so the invisible strut beamer appends after every footnote
	% (\@finalstrut\strutbox) is still sized for \footnotesize, silently
	% padding out extra depth below the last line. Rebuilding \strutbox
	% from this font's own measured metrics removes that padding, which is
	% most of the gap between the footnotes and the footer box below them.
	\settoheight{\dimen0}{Xy}%
	\settodepth{\dimen1}{Xy}%
	\setbox\strutbox\hbox{\vrule height\dimen0 depth\dimen1 width0pt}%
	\parindent 1em\noindent%
	\raggedleft
	\advance\rightskip by -\gothamFootnoteRightShift%
	% Same \@thefnmark emptiness guard as the side-footnote branch above,
	% for the same reason (an empty marker still leaving behind the "\,"
	% separator's own leading thin space) -- invisible here under
	% \raggedleft, where the left edge is never on display, but the same
	% underlying defect, so fixed the same way for consistency.
	\ifx\@thefnmark\@empty
		\insertfootnotetext\par%
	\else
		\insertfootnotemark\,\insertfootnotetext\par%
	\fi
	\fi
}
\makeatother
\setbeamertemplate{footnote}[gothamfn]

% Draws the side-footnote strip (see \sidefootnotes above). Lives in the
% footline template because that is typeset at page shipout, after the frame
% body -- and with it every \footnote{} call -- has run, so the collector box
% is complete by the time this executes. The \ifvoid guard doubles as the
% protection against beamer's measure-the-footline pass at \begin{frame} time
% (see the \sidefootnotes comment). \box (not \copy) empties the register
% globally as a side effect, which is exactly right: each overlay slide of a
% frame re-executes the body and re-collects its footnotes from scratch, so
% whatever this slide's footline draws must not leak into the next slide.
%
% Node geometry: tikz pins the node's *unrotated* anchor corner at the given
% coordinate and rotates the box counter-clockwise around it, so anchor=south
% east (the end of the collected block's last line) ends up as the strip's
% top-east corner: from there the block extends down the page (its lines run
% bottom-to-top, last footnote nearest the page edge), with the \raggedright
% text sitting at the strip's bottom end, just above the footer box (see the
% comment in the footnote template above). The y-coordinate
% comes from the same remembered gothamTitleBottom node the title-bar patches
% already use (same two-pass .aux mechanics -- see "Title bar and logo").
% The \ifcsname guard skips the drawing while the gothamTitleBottom shape is
% entirely unknown to pgf (i.e. no frametitle has been typeset yet in this
% run *and* nothing was inherited from a previous run's .aux) -- without it,
% referencing the node is a hard error rather than the usual harmless
% first-pass misplacement. In that situation the position would have been
% garbage anyway, and the frame-start clearing in standardframe below mops up
% the box this branch leaves unconsumed.
\makeatletter
\addtobeamertemplate{footline}{%
	\ifvoid\gothamSideFootnotesBox\else
		\ifcsname pgf@sh@ns@gothamTitleBottom\endcsname
		\vbox to 0pt{% zero-space trick, as in the yellow-box footline hook above
		\vss
		\begin{tikzpicture}[remember picture, overlay]
			\node[rotate=90, anchor=south east, inner sep=0pt]
				at ($(gothamTitleBottom -| current page.east)+(-\gothamSideFootnoteRightInset,-\gothamSideFootnoteTopSep)$)
				{\box\gothamSideFootnotesBox};
		\end{tikzpicture}%
		}%
		\fi
	\fi
}{}
% Empties the side-footnote collector (used by standardframe at frame start).
% Defined here, inside the \makeatletter block, because \voidb@x needs @ to be
% a letter at tokenization time -- standardframe's own definition below sits
% outside any \makeatletter region.
\newcommand{\gothamClearSideFootnotesBox}{\global\setbox\gothamSideFootnotesBox=\box\voidb@x}
\makeatother

% \source{text} -- for a "Source: ..."/"Image source: ..."/etc. citation line
% (common at the bottom/edge of a converted PowerPoint slide) that should look
% and sit *exactly* like the footnotes above -- same font, spacing,
% right-alignment, and the same stretch-to-sit-just-above-the-footer
% behaviour beamer already gives real footnotes -- but with no marker, neither
% inline in the body text nor next to the note itself, since a source line
% isn't anchored to one specific word. Implemented as a real \footnote{} (so
% it gets that positioning for free, rather than reimplementing it) with
% \thefootnote emptied locally; since \footnote{} still advances the footnote
% counter even with an empty marker, \addtocounter reverses that so a real
% \footnote{} elsewhere on the same slide doesn't skip a symbol in the
% \fnsymbol sequence. Unlike a real footnote, no label is added automatically
% here -- the label varies ("Source:", "Image source:", ...), so write it as
% part of #1.
\newcommand{\source}[1]{%
	\begingroup
	\renewcommand{\thefootnote}{}%
	\footnote{#1}%
	\endgroup
	\addtocounter{footnote}{-1}%
}

% A plain tabular with yellow \toprule/\midrule/\bottomrule, matching the
% frametitle/footer accent color. Takes the same column-spec argument as
% tabular itself, e.g. \begin{gothamtable}{lccc} ... \end{gothamtable}. Only
% booktabs rules are colored (via \arrayrulecolor, scoped to this table by
% \begingroup/\endgroup) -- if a table ever needs a plain black \hline too,
% that's unaffected since \arrayrulecolor only recolors rules drawn *after*
% it takes effect, and booktabs' rules read it fresh each time.
%
% RULE-BLEED WARNING (2026-07-12, direct instruction: "make sure the table
% lines do not extend beyond the text"), read this before writing a new
% gothamtable: a plain column spec like `{lccc}` above leaves the default
% \tabcolsep (6pt, or whatever a slide-local \setlength last set it to) as
% *implicit* padding on the table's own outermost edges -- invisible to the
% table's own text (which starts/ends flush with each column's actual
% content), but very visible to booktabs' \toprule/\midrule/\bottomrule and
% to colortbl's \rowcolor/\cellcolor, both of which span the *declared*
% column width including that implicit padding. The result, confirmed
% directly and independently on at least three slides so far
% (slides/u5h-Long-term-Bonds-Provisions-current.tex's own v1-1-0/v1-4-0,
% slides/c4c-Macbeth-Spot-Removers-Proposed-Structure-current.tex's own
% v1-2-2, slides/tw5-Exercise-Leverage-and-ROE-Solution-current.tex's own
% v1-4-0): every rule/highlight bleeds \tabcolsep past the visible text on
% both the left and right edges, most noticeable on a left-aligned
% (non-\centering) table where the mismatch is easy to see against the
% title's own left margin.
%
% Not fixed here in the environment definition itself: doing so would
% change every existing gothamtable's own column spacing out from under it
% (this environment has no way to tell "plain default spacing was intended"
% apart from "the bleed was never noticed"). Fix per table instead, the
% same three-part pattern every slide above uses -- worth reaching for by
% default on any *new* gothamtable, not just when the bleed is spotted:
%   {\setlength{\tabcolsep}{0pt}
%   \begin{gothamtable}{@{}l<{\hspace{<gap>pt}}c<{\hspace{<gap>pt}}c@{}}
%   ...
%   \end{gothamtable}}
% -- \tabcolsep zeroed (scoped locally, e.g. by wrapping the whole call in
% a `{...}` group, so it doesn't leak into later tables), explicit `@{}` at
% both outer edges of the column spec (so booktabs/colortbl have no
% implicit padding left to bleed into), and array's `<{\hspace{<gap>pt}}>`
% column modifier appended after every column but the last, to put back an
% explicit, *content-column* gap (matching whatever `2*\tabcolsep` the
% table used to have, if preserving its exact look, e.g. `<{\hspace{12pt}}>`
% for the ordinary 6pt-\tabcolsep default) -- `<{...}>` puts the gap
% *inside* the preceding column's own content, which `\rowcolor` (if used)
% colors along with everything else, unlike a plain `@{...}` separator,
% which `\rowcolor` deliberately leaves untinted (see u5h's own v1-6-0
% comment for that exact failure mode if a highlighted table ever needs
% this too).
%
% RE-CONFIRMED 2026-07-16, direct instruction, after this exact bleed
% reappeared on essentially every gothamtable in a freshly-converted
% chapter (Capital Budgeting Techniques) despite the warning above already
% existing at the time -- the rule above is not optional/only-if-noticed,
% it's the default way to write *any* new gothamtable, full stop. Check
% every new gothamtable against it before considering a slide done, not
% just when a rendered page happens to make the bleed visually obvious.
%
% BLOCK-RULE-COLOR CONVENTION (2026-07-16, direct instruction): a
% gothamtable placed *inside* any block environment (`infoblock`,
% `exampleblock`, `referenceblock`, `definitionblock`, plain `block`, or
% their `gothamfloating*` wrappers) should use black rules, not this
% environment's own default yellow -- yellow-on-the-block's-own-tinted-
% background reads poorly, unlike yellow-on-plain-white outside a block.
% Override locally, *inside* `\begin{gothamtable}{...}`, before `\toprule`
% (not before `\begin{gothamtable}` itself -- the environment's own
% opening code sets yellow *after* its `\begingroup`, so a color set
% before that point is simply overwritten; confirmed directly, first
% attempt this same day):
%   \begin{gothamtable}{...}
%   	\arrayrulecolor{c_primary_dark}
%   	\toprule
%   	...
%   \end{gothamtable}
% Booktabs rules read `\arrayrulecolor` fresh at each rule, and
% `\endgroup` (gothamtable's own closing code) resets it back to yellow
% afterward, so this never leaks into a later table. First applied on
% slides/6ux-Example-51-BMA-Payback-Table-current.tex's own
% `gothamfloatinginfoblock`.
\newenvironment{gothamtable}[1]{%
	\begingroup
	\arrayrulecolor{c_primary_highlight}%
	\begin{tabular}{#1}%
}{%
	\end{tabular}%
	\endgroup
}

% For a side-corner image overlay (the usual \vbox to 0pt{...} tikz overlay
% pattern used throughout this deck's converted slides) whose *top edge*
% should line up with the content text's own top edge -- distance from the
% page's own top edge to the cap-height top of a standardframe's first body
% line, measured directly against the compiled PDF's own text coordinates,
% not derived from \gothamContentTopShift or guessed. Previous per-slide
% image placements each guessed their own downward shift (e.g. a bare
% "-1.3cm") independently -- close, but not exact, landing the image's top
% a couple mm above the text's own top; use this constant instead of a
% fresh guess for any new "text top edge"-aligned image.
%
% 41.58pt (2026-07-22, re-measured): the previous value, 44.0pt, went stale
% the moment \gothamContentTopShift itself was increased (7.33pt -> 11.5964pt,
% "reduce the Top Margin", see that length's own comment above) without a
% matching update here -- that change pulls plain body text up, but this
% length is a separate, hand-measured constant with no formula tying it to
% \gothamContentTopShift, so it kept anchoring images at the *old*, lower
% position. Confirmed directly (pymupdf against this same freshly-compiled
% preview deck, not assumed from the \gothamContentTopShift delta alone):
% three independent plain-first-line standardframes with no per-slide
% top-shift correction (slides/5v4-CML-vs-SML-current.tex,
% slides/u5j-Exercise-CAPM-4-current.tex,
% slides/nt4-Exercise-CAPM-5-current.tex) all land their first line's own
% top edge at y=41.58pt post-change, not 44.0pt. Re-measure and adjust
% again if \gothamContentTopShift ever changes in the future.
\newlength{\gothamContentTopAlign}
\setlength{\gothamContentTopAlign}{41.58pt}

% For a "side filler image" -- a decorative image stuck flush to the page's
% own right edge (zero margin) with its top edge flush against the title
% bar's own bottom edge (zero gap), rather than sitting in the usual
% content-aligned corner inset -- distance from the page's own top edge
% down to the standardframe title bar's bottom edge, measured directly
% against the compiled PDF (a pixel scan for the yellow-to-white
% transition, at this deck's 16cm-wide/9cm-tall page size). First used on
% slides/dfc-Angel-Investors-current.tex.
%
% 25.43pt -- the raw measured value (originally 25.44pt; re-measured
% 2026-07-07 against the current compiled PDF via the title bar's own
% white-patch rectangle, which came back 25.427pt, close enough to the
% original measurement to be the same underlying edge, not a drift).
% Briefly set to 24.94pt instead (2026-07-06-ish), a deliberate 0.5pt
% overlap *into* the title bar, after a report of a visible hairline
% white gap right at this seam in one render -- not reproduced in this
% project's own pdftoppm-based rendering checks at the time (every pixel
% scan at the raw value showed a clean cut, no gap), so assumed to be a
% renderer-specific sub-pixel rounding artifact rather than a real error
% in the constant, and "fixed" by deliberately overlapping instead of
% chasing the specific viewer. That overlap turned out to be the wrong
% call: direct feedback 2026-07-07 reported the *opposite* symptom on
% `sideimageframe` (the image's own top edge visibly overlapping the
% title box) -- exactly what a real 0.5pt overlap looks like once it's
% actually noticed, rather than a rendering artifact. Reverted back to
% the precise raw measurement: a real, small, deliberate overlap chosen
% to paper over an unreproduced gap in an unspecified other renderer
% turns out to be a worse trade than just using the measured value and
% letting a genuine future gap report (if one ever recurs, in a
% specific, identified renderer) be diagnosed on its own terms.
%
% Re-measured 2026-07-16, 25.43pt -> 25.32pt: the frametitle hook's
% outer sep fix (see its own comment) moved the visible yellow bottom
% edge up by the 0.2pt the trim patch's fill had previously been inset
% below its gothamLogoAnchor.south anchor -- the yellow now genuinely
% ends flush with the logo's bottom edge, as always intended. New raw
% measurement against the rebuilt PDF: the trim patch's fill top edge
% (= the visible yellow bottom) sits 25.228 PDF units below the page
% top, i.e. 25.3226pt (earlier measurements in this comment quoted PDF
% coordinate values directly as "pt"; PDF units are big points, 1bp =
% 1.00375pt -- a systematic ~0.1pt understatement at this magnitude,
% harmless then, but corrected now that the constant is being reset
% anyway). Note this value is also what the overlay-logo nodes
% (titleframe/quoteframe/the hook's fixed-position redraw) derive their
% vertical inset from: visible-bar-bottom minus \gothamLogoHeight is
% exactly the inline logo's own vertically-centered top margin.
\newlength{\gothamTitleBarBottomAlign}
\setlength{\gothamTitleBarBottomAlign}{25.32pt}

% Note on the [t] default below: it top-aligns content within the frame
% (beamer's own default is vertical centering, "c"), so slides don't have
% their content floating in the middle. It also sidesteps a real beamer bug:
% \begin{frame}[]{Title} (a literal empty options bracket) fails outright,
% even in plain beamer with no theme at all -- so the optional argument must
% never default to an empty string. The title is always passed as an
% explicit argument (rather than left for beamer to pick up from a trailing
% brace group), since anything placed between \begin{frame} and its title --
% like our \setbeamertemplate calls -- breaks that trailing-brace mechanism.

% Pulls the frame's content up closer to the title bar: measured (via the
% PDF's own coordinates) at about 22pt of gap by default between the title
% bar's bottom edge and the first line of body content: this is a third of
% that, shrinking the gap to about two thirds of its original size.
%
% Called "Top Margin" in conversation from 2026-07-12 onward (direct
% instruction) -- the deck-wide name for the resulting gap between the
% title bar and a standardframe's first line of body content, i.e. the
% vertical position a plain-text frame's first baseline naturally lands at
% (measured directly against a reference slide's own compiled PDF, e.g.
% slides/aii-Leverage-Effects-current.tex's first line). Any first content
% that isn't plain text (a table, a minipage, an itemize with an unusual
% first item, ...) needs its own per-slide \gothamIncreaseContentTopShift/
% \gothamReduceContentTopShift correction to land on this same baseline --
% not portable between slides even for superficially similar content (see
% slides/zgz-EBIT-Definition-current.tex's own v1-2-0 comment for a case
% where an earlier, wrong correction was carried over uncritically).
\newlength{\gothamContentTopShift}
% 11.5964pt (2026-07-22, direct instruction: "reduce the Top Margin"
% generally, tested at 20%/30%/40% smaller against
% slides/8s8-Unit-Investment-Trusts-current.tex before settling on 40% --
% up from the original 7.33pt so the *resulting* gap (this length pulls
% content up by, so a bigger value means a smaller visible margin) comes
% out to 60% of its old size: measured directly on 8s8 at 10.67pt before,
% 6.42pt after (target was 6.40pt, within 0.02pt). Since every
% standardframe applies this same pull-up before any per-slide
% \gothamIncreaseContentTopShift/\gothamReduceContentTopShift correction,
% changing it here shifts the baseline project-wide, not just on 8s8.
\setlength{\gothamContentTopShift}{11.5964pt}

% "Top Margin" (above) is one of four deck-wide margin names introduced in
% conversation from 2026-07-12 onward (direct instruction) for the plain
% page-edge/title-bar/footer gaps every slide is expected to obey by
% default. The other three, for reference (each maps to a length already
% declared elsewhere, not a new one introduced here):
% - "Left Margin": the default gap from the page's own left edge to the
%   title text and body content -- \gothamFrametitleLeftpading
%   (beamerouterthemegotham.sty/beamerthemegotham.sty, 2.2ex, measured
%   directly against the compiled PDF at 12.552pt in this deck's usual
%   fonts -- see e.g. slides/tw5-Exercise-Leverage-and-ROE-Solution-current.tex's
%   own v1-1-0 comment for that same measured value).
% - "Right Margin": the default gap from the page's own right edge to the
%   institute logo (and, by the same convention, to anything else
%   deliberately flush-right, e.g. a corner-anchored image) --
%   \gothamFrametitleRightpading (0.28cm/7.93pt, this file's own
%   \gothamInstituteLogoSquare placement above uses it directly as the
%   logo's own inner sep).
% - "Bottom Margin": the default gap between a frame's own body content
%   and the \source{}/footnote line below it. Unlike the other three, this
%   one has no single settable length behind it -- \source{}'s own
%   vertical position isn't fixed the way the footer box is; it tracks the
%   frame's actual content height once that content gets tall enough (see
%   slides/u5h-Long-term-Bonds-Provisions-current.tex's own v1-3-0 comment,
%   "past a certain table height it starts tracking the table's own growth
%   roughly 1:1"). "Respecting Bottom Margin" in practice means leaving
%   comparable clearance to whatever a reference slide with unremarkable
%   content already shows, the same empirical, measure-against-a-reference-
%   slide approach "Top Margin" itself uses, rather than a fixed constant.

% Compensates a gap that opens up specifically when a standardframe's body
% starts with a \begin{minipage}{...} (used to width-limit text next to an
% image/diagram overlay, e.g. slides/igr-Budgeting-current.tex,
% slides/bdh-Working-Capital-Example-current.tex -- see README.md's
% "Available layouts" section for why minipage instead of columns there) --
% a minipage is a *box*, and TeX's vertical glue insertion (\topskip
% specifically) treats the first box in a vertical list differently from
% the first line of plain paragraph text, so the same \gothamContentTopShift
% pull-up above lands closer to the title for minipage-wrapped content than
% for plain content. Cannot be fixed by \vspace(*) right after
% \begin{minipage}{...}: both plain and starred forms were tried and
% measurably had zero effect (the star only protects glue from page-break
% discard, not from \topskip's own separate glue-insertion calculation,
% which absorbs leading glue at the start of a vertical list regardless).
% Fixed instead by shrinking standardframe's own pull-up for the affected
% frame -- see \gothamContentTopShiftReduction below. The value below is
% calibrated empirically, not derived from first principles: an isolated
% two-frame test (single-level itemize, one plain, one minipage-wrapped)
% measured only a 4.4pt gap, but real slides (nested itemize inside the
% minipage) showed a larger ~7.5-7.6pt gap against a real plain-itemize
% slide (slides/jyn-Categorization-current.tex) -- the extra few points
% likely come from itemize's own topsep/partopsep being computed differently
% depending on whether the list opens a box or a paragraph, on top of the
% topskip/box effect alone. Tuned so slides/igr-Budgeting-current.tex and
% slides/bdh-Working-Capital-Example-current.tex's first content line
% lands within 0.1pt of slides/jyn-Categorization-current.tex's own
% (200.917pt vs. 200.871pt/200.980pt in the compiled PDF).
\newlength{\gothamMinipageTopCorrection}
% 9.38pt (2026-07-22, re-measured): the previous 7.6pt value went stale
% the moment \gothamContentTopShift itself was increased (7.33pt ->
% 11.5964pt, "reduce the Top Margin", see that length's own comment above)
% without a matching update here -- confirmed directly (pymupdf against a
% freshly-compiled preview deck, not assumed from the \gothamContentTopShift
% delta alone, which doesn't translate 1:1 into a minipage's own \topskip
% behavior): four independent slides all using this constant completely
% unmodified (slides/fr3-Forms-of-Capital-Market-Efficiency-Strong-Semistrong-Weak-current.tex,
% slides/za2-Startup-Financing-Categorization-current.tex,
% slides/yke-Implications-of-Homemade-Dividends-current.tex,
% slides/6j4-Modigliani-Miller-Assumptions-current.tex) all land their
% own first line at an identical y=39.80pt post-change, 1.78pt short of
% this deck's current 41.58pt Top Margin reference -- not slide-specific
% noise, the same deterministic amount on every one of them. Re-measure
% and adjust again if \gothamContentTopShift ever changes in the future.
\setlength{\gothamMinipageTopCorrection}{9.38pt}

% A *third* minipage-top correction, on top of \gothamMinipageTopCorrection
% above rather than replacing it, for a minipage whose own first internal
% content is a zero-height `\vbox to 0pt{...}` tikz overlay (a
% corner-anchored decorative image, placed there specifically so it stays
% unconditionally visible even when the minipage's own later content
% includes a `\pause` -- see e.g.
% slides/y5z-CoCo-Bonds-current.tex's own "v1-2-0" note for the fuller
% pause-inheritance story behind why the image has to sit there at all).
% 2026-07-08, found while chasing "the top margin of the text on y5z/22s
% is too low" -- measured directly against the compiled PDF, y5h/22s both
% landed 16.5pt too far down despite already using
% \gothamMinipageTopCorrection, with the zero-height image overlay sitting
% *before* the minipage entirely as the actual root cause (that ordering
% makes the image, not the minipage, the frame body's true first box,
% which \topskip's own special first-box glue handling reacts to
% differently than a plain first minipage). Moving the image inside the
% minipage as its own first content closed most of the gap on its own
% (restoring the minipage as the frame's true first box) but left a
% further 10.35pt uncorrected -- confirmed the same on two independent
% slides (y5z, 22s) with identical structure, not assumed transferable
% from one to the other. Call via
% `\gothamReduceContentTopShift{\dimexpr\gothamMinipageTopCorrection-\gothamMinipageWithCornerImageExtraCorrection\relax}`
% for any slide combining this exact pattern (corner image as a minipage's
% own first content, ahead of a `\pause`) -- confirmed transferable as-is
% to two more slides with the same structure
% (slides/qbj-Bank-Loans-Credit-Lines-current.tex,
% slides/jkf-Bank-Loans-Term-Loan-current.tex), neither of which
% actually uses `\pause` itself, so this isn't specific to that -- the
% image-inside-minipage-as-first-content shape alone is what triggers it.
\newlength{\gothamMinipageWithCornerImageExtraCorrection}
\setlength{\gothamMinipageWithCornerImageExtraCorrection}{10.35pt}

% A *different* minipage-top correction, specifically for a minipage+
% itemize sitting as the first content of a `sideimageframe` (2026-07-06,
% found while chasing "the top margin of the text on ezz/acw is again
% off" -- see CONVERSION_NOTES.md's own "beamer/TeX gotchas" entry for the
% full story). `\gothamMinipageTopCorrection` above is calibrated for a
% *plain* standardframe, and simply does not transfer to `sideimageframe`:
% measured directly on three real sideimageframe slides
% (slides/xmp-Implications-of-Capital-Market-Efficiency-Considerations-current.tex,
% slides/dfc-Angel-Investors-current.tex,
% slides/ezz-Examples-of-Tests-for-Market-Efficiency-current.tex --
% one flat itemize, one flat itemize, one nested-first-item itemize
% respectively, so this isn't a flat-vs-nested effect either), all three
% landed within 0.6pt of each other at their own *uncorrected* natural
% position (52.03-52.57pt against this deck's usual 44.0pt first-line
% reference) regardless of which of the two *wrong* corrections
% (`\gothamReduceContentTopShift{\gothamMinipageTopCorrection}` or
% `\gothamIncreaseContentTopShift{\gothamFlatItemizeTopCorrection}`) had
% been applied to each beforehand -- i.e. `sideimageframe` itself pushes
% its own first-content minipage down by a consistent ~8.2pt relative to
% a plain standardframe, on top of whatever generic itemize-shape
% correction would otherwise apply, most likely from the extra
% \sidefootnotes/\settowidth setup work `sideimageframe`'s own opening
% code does before `\begin{standardframe}` ever runs. Use
% `\gothamIncreaseContentTopShift{\gothamSideImageFrameMinipageCorrection}`
% (note: *increase*, not reduce) as the *only* topshift call on any
% future sideimageframe slide whose first content is a minipage+itemize --
% not the generic `\gothamMinipageTopCorrection`/
% `\gothamFlatItemizeTopCorrection` this deck uses everywhere else.
\newlength{\gothamSideImageFrameMinipageCorrection}
\setlength{\gothamSideImageFrameMinipageCorrection}{8.2pt}

% 6.26pt (2026-07-22, re-measured): stale for the same reason
% \gothamMinipageTopCorrection above went stale -- \gothamContentTopShift
% itself changed (7.33pt -> 11.5964pt) without a matching update here.
% Re-measured directly against a freshly-compiled preview deck across six
% independent sideimageframe slides using this constant unmodified: four
% (ezz, xmp, hzx, u6n) land their own first line at an identical
% y=39.64pt, 1.94pt short of this deck's current 41.58pt Top Margin
% reference (i.e. already pulled up *too far*); the other two (c7w, dfc)
% are close but not identical (39.53pt/-2.05, 40.18pt/-1.4) -- no single
% value fixes all six exactly, so this recalibrates to the four-way
% cluster (the majority case) and accepts a small (<0.6pt) residual on the
% two outliers. Since this constant feeds
% `\gothamIncreaseContentTopShift{}` (more pull-up, i.e. *smaller* y), and
% the observed y was already short of target, the fix is a *decrease*
% (8.2pt - 1.94pt), not an increase -- caught by re-measuring after a
% first, wrong-direction attempt at 10.14pt drove the four-slide cluster
% to y=37.70pt, further from target instead of closer.
\setlength{\gothamSideImageFrameMinipageCorrection}{6.26pt}

% Opt-in override: reduces standardframe's \gothamContentTopShift pull-up by
% this amount for the frame currently being opened. Empty-default-style flag
% (mirrors \gothamChapterLogoPath etc.): 0pt by default (no effect), set via
% \gothamReduceContentTopShift{<length>} called *before* \begin{standardframe},
% reset back to 0pt in standardframe's own closing code so it never leaks
% into the next frame. Used by minipage-based slides to counteract
% \gothamMinipageTopCorrection above.
\newcommand{\gothamContentTopShiftReduction}{0pt}
\newcommand{\gothamReduceContentTopShift}[1]{\global\def\gothamContentTopShiftReduction{#1}}

% A *second*, separate first-line-position quirk, found and fixed
% 2026-07-05 after a run of slides (batch 4 plus several earlier ones) were
% found to land at two visibly different heights depending purely on
% content shape, with no minipage involved at all: a plain (non-minipage)
% standardframe whose top-level itemize's very *first* \item is immediately
% followed by its own nested \begin{itemize}/\begin{enumerate} lands 1.834pt
% higher than one whose first \item has no such immediate child (checked
% against every content slide converted so far via each compiled page's own
% first-line PDF baseline coordinate, not eyeballed -- 13 "nested-first-item"
% slides all landed within 0.05pt of each other, and 9 "flat-first-item"
% slides all landed within 0.05pt of each other, at a clean +1.834pt offset
% from the first group; isolated two-frame test cases, differing only in
% presence/absence of that immediate nested list, reproduced the exact same
% 1.834pt gap independent of which text/list type/item count was used).
% Not chased to one named TeX primitive (most likely some interaction of
% \addvspace's glue-coalescing behaviour with \partopsep/\topsep at the
% \begin{itemize}/\begin{frame} boundary, given the existing
% \gothamMinipageTopCorrection comment above already flags topsep/partopsep
% as varying with box vs. paragraph context) -- calibrated empirically like
% \gothamMinipageTopCorrection, not derived from first principles.
% Interacts with minipage differently, not additively: inside a minipage,
% this nested-vs-flat distinction disappears entirely (both land within
% 0.05pt of each other) -- so this correction is for *plain* itemize-first
% standardframes only, never combined with \gothamMinipageTopCorrection.
\newlength{\gothamFlatItemizeTopCorrection}
\setlength{\gothamFlatItemizeTopCorrection}{1.834pt}

% 0.01pt (2026-07-22, re-measured): stale for the same reason
% \gothamMinipageTopCorrection above went stale -- \gothamContentTopShift
% itself changed (7.33pt -> 11.5964pt) without a matching update here.
% Re-measured directly against three independent flat-itemize slides using
% this constant unmodified (574, 5cr, c8e): all three landed their own
% first line at an identical y=39.76pt, 1.82pt short of this deck's
% current 41.58pt Top Margin reference -- coincidentally, the recalibrated
% \gothamContentTopAlign reference now already accounts for almost all of
% what this correction used to supply on its own, leaving next to nothing
% left to add. Effectively zeroed out rather than removing the call
% itself, so every flat-itemize slide's own
% `\gothamIncreaseContentTopShift{\gothamFlatItemizeTopCorrection}` line
% keeps working unchanged.
\setlength{\gothamFlatItemizeTopCorrection}{0.01pt}

% Opposite-direction counterpart to \gothamReduceContentTopShift: *increases*
% the pull-up instead of reducing it (needed for \gothamFlatItemizeTopCorrection
% above, since a flat first item needs *more* pull-up, not less). Same
% storage variable, just negated -- \gothamContentTopShiftReduction is
% allowed to go negative, since the frame-opening code only ever computes
% \gothamContentTopShift-\gothamContentTopShiftReduction, never checks its
% sign.
\newcommand{\gothamIncreaseContentTopShift}[1]{\global\def\gothamContentTopShiftReduction{-#1}}

% Define custom frame type with en-dash bullets
\newenvironment{standardframe}[2][t]{%
	% Restart footnote numbering (the *, dagger, ... symbols) fresh on every
	% slide, rather than letting it run on cumulatively across the whole deck
	% the way beamer does by default. Two reasons: \fnsymbol only defines 9
	% symbols before erroring out, which a real multi-slide deck would run
	% through in no time if every \footnote{}/\source{} call across the whole
	% talk drew from one shared counter; and each slide's footnotes are read
	% and referenced only on that one slide anyway, so there's no reason for
	% their numbering to depend on how many footnotes earlier slides happened
	% to have. This must run *before* \begin{frame}: beamer saves the
	% counter's current value as that frame's baseline as part of
	% \begin{frame}'s own internal setup (restoring it before every overlay
	% replay, so overlays within one frame don't double-count) -- resetting
	% after \begin{frame} would be reset too late, after that baseline was
	% already captured.
	\setcounter{footnote}{0}%
	\begin{frame}[#1]{#2}%
	% Forces this frame's own body to start with the correct default text
	% color, regardless of what any *previous* frame left active
	% (2026-07-16, direct instruction, root-caused via an isolated minimal
	% reproduction). \titleframe's/\quoteframe's own
	% \color{\gothamTitleTextColor} calls -- set inside content built via a
	% stored box register or a zero-space overlay, used/shipped out *later*
	% rather than immediately -- don't reliably get "undone": xcolor's
	% \color doesn't emit a PDF color operator immediately when called, only
	% when the *next* real glyph/path actually needs it, so whatever color
	% one of those frame types last set can end up as the active PDF color
	% well into the *next* frame's own first draw call. Invisible for Graz
	% (the leaked color, c_primary_dark, happens to match ordinary body text
	% anyway) but catastrophic for MCI/Innsbruck (c_primary_white, rendering
	% plain non-itemize body text invisible against its own white
	% background -- itemize items escape it, since beamer's own itemize
	% template sets color fresh per bullet). An \AtBeginEnvironment{frame}
	% hook (etoolbox) was tried first and confirmed *not* to fire at all --
	% beamer's own "frame" isn't a plain \newenvironment etoolbox can patch
	% that way -- so this resets directly in standardframe's own opening
	% code instead, exactly where this project already has full control.
	\color{c_primary_dark}%
	% Defensive: the tokens after \begin{frame} are collected into the frame
	% body and replayed once per overlay slide, so this clears the side-
	% footnote collector at the start of every slide's pass -- normally a
	% no-op (the previous slide's footline already consumed the box), but it
	% stops stale footnotes leaking in if a previous frame ever collected
	% without shipping a footline.
	\gothamClearSideFootnotesBox%
	\setbeamertemplate{itemize item}{\gothamBulletDot}%
	\setbeamertemplate{itemize subitem}{\gothamBulletDot}%
	\setbeamertemplate{itemize subsubitem}{\gothamBulletDot}%
	\setbeamertemplate{itemize subsubsubitem}{\gothamBulletDot}% 4th level, see its own comment further below (search "subsubsub")
	% Only for top-aligned frames (the default, #1=t): a [c]-centered frame
	% (2026-07-22, direct instruction, found while centering
	% slides/bw9-Pearson-Correlation-Illustration-current.tex/agc/crj) was
	% landing its content almost flush against the title bar instead of
	% truly centered -- root cause was this same pull-up firing
	% unconditionally regardless of #1. beamer's own [c] option already
	% centers the frame body within the available space; stacking this
	% deck's *own* top-alignment pull-up on top of that just drags the
	% already-centered box further upward by the same fixed amount, which
	% is exactly the asymmetry that was measured (a [c] figureframe's image
	% landing ~0pt from the title bar and ~10pt from the footer instead of
	% symmetric). Guarded with \ifstrequal (etoolbox, already loaded above)
	% rather than \ifx, matching this file's own \ifstrequal{##1}{...}
	% pattern elsewhere (\ifx would need #1 fully expanded first).
	\ifstrequal{#1}{c}{}{\vspace*{-\dimexpr\gothamContentTopShift-\gothamContentTopShiftReduction\relax}}%
}{%
	\end{frame}%
	\global\def\gothamChapterLogoPath{}% opt-in only for the frame just closed
	\global\def\gothamContentTopShiftReduction{0pt}% opt-in only for the frame just closed
	\global\gothamSideFootnotesfalse% opt-in only for the frame just closed
	\global\def\gothamFrametitleFullHeightYellow{}% opt-in only for the frame just closed
	\global\def\gothamFrametitleLogoFixed{}% opt-in only for the frame just closed
	\global\def\gothamFooterBoxColor{\gothamFooterBoxColorDefault}% opt-in only for the frame just closed
}

% The bullet itself: a small filled circle, matching the source PowerPoint
% deck's own round bullets (a plain \textbullet, tried first, is too
% small/thin in Fira Sans at any standard text size to read as a solid dot
% -- a tikz-drawn circle gives full control over the actual rendered size
% instead of depending on the font's own glyph design). baseline=-0.65ex
% roughly centers it on the surrounding text's x-height, the same vertical
% position a font's own bullet glyph would normally sit at.
%
% Its fill color is \gothamBulletColor, not a literal color token here --
% set by the per-university preamble before \input{}ing this file, same
% "value supplied externally, referenced lazily" idiom as
% \gothamInstituteLogoPath/\gothamLogoHeight above (\gothamBulletDot's own
% body only expands \gothamBulletColor when a bullet actually gets drawn,
% long after this \newcommand runs). Graz's own bullets are its primary
% highlight color (matching this project's original, single-institution
% behavior); MCI's and Innsbruck's are their own secondary blue instead
% (2026-07-16, direct instruction) -- see each institution's own
% Slidemaster_preamble_*.tex for its \gothamBulletColor setting.
\newcommand{\gothamBulletDot}{\tikz[baseline=-0.65ex]{\fill[\gothamBulletColor] circle (1.7pt);}}

% "a)"/"b)"/"c)"-labeled list, for exercise sub-tasks that read better
% lettered than bulleted (2026-07-20, direct instruction, first used on
% slides/y3a-Exercise-Diversification-2-current.tex's own bottom three
% items).
%
% First attempt (superseded, kept as a documented wrong turn -- see
% README.md's own "lualatex vs latex" section for why this project keeps
% wrong write-ups rather than deleting them): reuse itemize's own top-level
% `\llap`-based label positioning by locally overriding
% `\setbeamertemplate{itemize item}` inside a plain `\begin{itemize}` --
% the reasoning being that `\llap` draws its content flush against the
% text's own left edge "regardless of the label's width." That reasoning
% turned out to be wrong in a way that only showed up once actually
% measured against the compiled PDF (not just visually eyeballed): a plain
% bullet's own `\llap`-drawn footprint (title/Left-Margin-relative-left-
% edge to text-start) measured only ~9.05pt end-to-end on this deck's real
% slides, matching the Left-Margin-to-text-start gap exactly -- but "a)"
% rendered through the *exact same* mechanism measured a ~13.97pt
% footprint for the same gap, not the ~8.5-9pt `\settowidth` reported for
% the bare glyphs in isolation. The extra ~5pt was never fully root-caused
% (this custom itemize's own `\list` call, further below, only overrides
% `\makelabel` -- it never sets `\leftmargin`/`\labelwidth` itself, so
% those are left at whatever this deck's global itemize defaults happen to
% be, and `\llap`'s actual box-width behavior inside that inherited,
% unexamined list-parameter state is where the extra width most likely
% comes from) -- rather than keep guessing at *why*, measured the actual
% failure directly (the compiled "a)" sat ~4.9pt left of the Left Margin,
% confirmed via exact PDF text-position extraction, not a visual guess)
% and abandoned the approach once even `\tiny` plus manual per-letter
% kerning only just barely closed the gap while reading as illegibly small
% -- a real dead end, not a tuning problem.
%
% Working fix: don't reuse itemize's `\llap` label mechanism at all -- use
% a real, width-reserving `\makebox[W][r]{...}` label instead (the same
% *kind* of mechanism `\gothamEnumerateLabelWidth` above already uses for
% a nested enumerate, just with an explicit `\leftmargin`/`\labelwidth`/
% `\labelsep` this environment sets itself rather than inheriting), inside
% its own plain `\list`, not a reused `\itemize`. `\leftmargin` (label box
% + gap) is set to 11.1pt -- 9.1pt label box (fits "a)"/"b)"/"c)" at this
% deck's normal `\small` itemize size with a little clearance, confirmed
% directly: measured label-left-edge stays right of the 12.552pt Left
% Margin for all three letters, closest case (widest letter, "c)") 1.37pt
% clear) plus 2pt `\labelsep` for a visible gap before the text. Trade-off,
% accepted deliberately rather than found by accident: item text lands
% ~2.2pt right of where a plain bulleted item's text would (23.8pt vs.
% 21.6pt) -- fitting a legible "a)"/"b)"/"c)" at normal reading size plus
% a real gap simply doesn't fit in the ~9pt a plain bullet's tiny dot
% leaves, so *exact* text-position parity with plain bullets and a legible
% label are mutually exclusive here; legibility won given the two-point
% shift is visually negligible and there's no plain-bulleted item left on
% the same slide to compare against directly anyway (y3a's own first item
% moved into its infoblock). Re-derive `\leftmargin`/`\labelwidth` the same
% way (render, measure with exact text-position extraction, adjust) if
% this is ever used at a different font size than this deck's normal
% itemize `\small`.
\newcounter{gothamalphaitem}
\newcommand{\gothamAlphaItem}{\stepcounter{gothamalphaitem}\alph{gothamalphaitem})}
\newenvironment{gothamalphaitemize}{%
	\setcounter{gothamalphaitem}{0}%
	\list{\gothamAlphaItem}{%
		% 4.75pt (2026-07-20, direct instruction: "half the width of the
		% 'C' in 'Calculate'" -- \settowidth{C} at this deck's normal
		% itemize \small measured 5.5pt, so 2pt original + 2.75pt half-C).
		% Only \labelsep changes the text's own start position -- the
		% label's own box stays exactly where it was (its left edge is
		% pinned to this list's origin regardless of \labelsep, since
		% \leftmargin=\labelwidth+\labelsep always cancels \labelsep back
		% out of the label box's own position), so this widens the gap
		% without reopening the Left-Margin question already solved above.
		\setlength{\labelsep}{4.75pt}%
		\setlength{\labelwidth}{9.1pt}%
		\setlength{\leftmargin}{\labelwidth}%
		\addtolength{\leftmargin}{\labelsep}%
		\setlength{\itemindent}{0pt}%
		% \itemsep/\parsep left at plain \list's own defaults otherwise --
		% since this environment is built directly on \list rather than
		% reusing beamer's own itemize template, it never got beamer's own
		% itemize spacing hooks at all, only whatever \list falls back to.
		% Measured directly (2026-07-21, direct instruction, found on
		% slides/nt4-Exercise-CAPM-5-current.tex): that default reads as
		% *less* space between two items (13.7pt, item-to-item) than
		% between two wrapped lines within one *single* item (15.6pt,
		% pure \baselineskip) -- backwards from how a list should read
		% (item boundaries should be at least as open as a same-item line
		% wrap, not tighter). \itemsep set explicitly so the total gap at
		% an item boundary (baselineskip + itemsep) reads as clearly more
		% open than a plain wrapped line, not accidentally tighter.
		% Reduced from the original 6pt to 2.5pt (2026-07-22, found during
		% a full-project check after the itemize-rules work, not part of
		% that work itself): 6pt overshot by more than needed to clear the
		% 1.9pt gap (15.6 vs. 13.7), and slides/kmf-Exercise-CAPM-6-current.tex's
		% own 5-item list -- longer than nt4's own -- compounded that
		% excess across 4 item boundaries into a ~25pt overfull frame, its
		% own last item cut off by the footer (confirmed directly,
		% rendered and compared against a pre-regression PDF). 2.5pt
		% keeps a real, visible margin over the 1.9pt needed (not just
		% barely crossing it) while re-measuring clean (no overfull) on
		% kmf; re-verify against both nt4 and kmf directly if this is
		% ever changed again -- they're this environment's own tightest
		% (most items) and originally-tuned-against cases respectively.
		\setlength{\itemsep}{2.5pt}%
	}
	\raggedright% plain \list defaults to justified+hyphenated text, unlike
	% this deck's own itemize (which calls this too, further below) --
	% missed on the first pass, only caught by comparing a rendered a)/b)/c)
	% item side by side against a plain bulleted one ("stan-\ndard deviation"
	% hyphenating where a plain-bulleted version of the same sentence
	% elsewhere in this deck doesn't).
}{%
	\endlist
}

% Itemize indentation: level 1 bullets should sit flush with the frame title
% (no extra left margin), and each deeper level should be indented just
% enough for its bullet to line up under the *text* (not the bullet) of the
% level above. A nested list's \leftmarginN is added on top of its parent's
% (indentation is cumulative), so giving every level the same "step" --
% label width + labelsep, the width a bullet plus its gap to text takes up --
% makes level N's text sit at N*step and level N's bullet at (N-1)*step, i.e.
% exactly at level (N-1)'s text (0 for level 1, flush with the title).
\makeatletter
\newlength{\gothamBulletWidth}
\settowidth{\gothamBulletWidth}{\gothamBulletDot}
\newlength{\gothamItemStep}
\setlength{\gothamItemStep}{\labelsep}
\addtolength{\gothamItemStep}{\gothamBulletWidth}
\setlength{\leftmargini}{\gothamItemStep}
\setlength{\leftmarginii}{\gothamItemStep}
\setlength{\leftmarginiii}{\gothamItemStep}
\setlength{\leftmarginiv}{\gothamItemStep}
\makeatother

% A 4th itemize level -- beamer's own \itemize (beamerbaselocalstructure.sty)
% hard-caps nesting at 3 levels (\ifnum \@itemdepth >2\relax\@toodeep),
% independently of anything this deck sets: a 4th \begin{itemize} always
% raised "! LaTeX Error: Too deeply nested." even though plain LaTeX's own
% kernel itemize supports 4. Found 2026-07-15, direct instruction, on
% slides/hrg-Exchange-Rate-Quotation-Forward-current.tex -- genuinely
% 4 levels deep in the source pptx itself (confirmed via a PowerPoint COM
% render, not assumed), not a conversion artifact to design around.
% Since this deck already styles every existing level identically (same
% \gothamBulletDot, same \gothamItemStep margin -- see the comments right
% above/below), extending to a 4th level needs no new visual design, just
% patching beamer's own depth limit up by one and telling it what to call
% the new level ("subsubsub", the next name in beamer's own
% sub/subsub/... naming scheme) so its template/color lookups
% (\usebeamertemplate{itemize subsubsubitem}, \usebeamercolor{itemize/
% enumerate subsubsub body}, etc.) resolve to *something* instead of
% erroring -- undefined beamer templates/colors are expected to fall back
% to sensible defaults (no extra spacing hook, plain text color) rather
% than raising an error, same as any other beamer template/color that a
% given theme never bothers to set; the actual bullet symbol still needs
% its own \gothamBulletDot template below, added explicitly the same way
% as the existing 3 levels, not left to whatever fallback beamer has for
% an unset "itemize subsubsubitem" template.
% \beamer@computepref is patched wholesale (not via \patchcmd) since its
% body is one single \ifcase chain, not a stable string to match against.
% \itemize itself was *first* tried via \patchcmd{\itemize}{...}
% (matching this deck's own established practice for beamer/theme
% internals, e.g. beamerfontthemegotham.sty's own
% \patchcmd{\beamer@title}{...}) -- this failed outright, confirmed
% directly (\patchcmd's own failure branch fired, raising a clean
% "Patching failed" error rather than silently doing nothing): recent
% LaTeX2e kernels implement an optional-argument environment like
% beamer's own `\renewcommand{\itemize}[1][]{...}` via
% `\@protected@testopt`, which stores the *real* body not under `\itemize`
% itself but under a differently-named internal control sequence (`\show
% \itemize` confirms this directly: it prints as
% `\@protected@testopt \itemize \\itemize {}`, i.e. plain `\itemize` is
% only a thin dispatcher, and the actual `\ifnum \@itemdepth >2...` body
% \patchcmd was trying to match against was never there to find). Rather
% than chase that internal storage name (liable to differ across LaTeX
% kernel versions, which \patchcmd's own documented usage doesn't need to
% care about), \itemize is fully re-declared here instead, copied
% verbatim from beamer's own beamerbaselocalstructure.sty (this project's
% installed version) with only the depth bound itself changed from 2 to
% 3 -- the same declaration form beamer itself used
% (\renewcommand{\itemize}[1][]{...}), which regenerates both the
% dispatcher and the real body correctly under whatever internal name
% the current kernel expects, sidestepping the naming question entirely.
% Trade-off, stated plainly: if a future beamer release changes this
% macro's own body for an unrelated reason, this override will silently
% keep using today's (2026-07) copy instead of picking up that change --
% acceptable here since the override is one line different from upstream
% and clearly marked, but re-diff against beamerbaselocalstructure.sty's
% own \itemize if beamer is ever upgraded and itemize behavior looks off.
\makeatletter
\long\def\beamer@computepref#1{%
	\let\beameritemnestingprefix\@empty%
	\ifcase#1\or\or\def\beameritemnestingprefix{sub}\or\def\beameritemnestingprefix{subsub}\or\def\beameritemnestingprefix{subsubsub}\or\@toodeep\fi%
}
\renewcommand{\itemize}[1][]{%
	\ifblank{#1}{}{\def\beamer@defaultospec{#1}}%
	\ifnum \@itemdepth >3\relax\@toodeep\else
		\advance\@itemdepth\@ne
		\beamer@computepref\@itemdepth% sets \beameritemnestingprefix
		\usebeamerfont{itemize/enumerate \beameritemnestingprefix body}%
		\usebeamercolor[fg]{itemize/enumerate \beameritemnestingprefix body}%
		\usebeamertemplate{itemize/enumerate \beameritemnestingprefix body begin}%
		\list
			{\usebeamertemplate{itemize \beameritemnestingprefix item}}
			{\def\makelabel##1{%
					{%
						\hss\llap{{%
								\usebeamerfont*{itemize \beameritemnestingprefix item}%
								\usebeamercolor[fg]{itemize \beameritemnestingprefix item}##1}}%
					}%
				}%
			}
	\fi%
	\beamer@cramped%
	% \beamer@cramped (beamerbaselocalstructure.sty, unmodified upstream
	% code) unconditionally does `\itemsep0pt\topsep0.5pt`, every single
	% time -- called right here, in this project's own copy of beamer's
	% own \itemize body, same position as upstream. Setting \itemsep any
	% earlier (first tried as a separate \AtBeginEnvironment{itemize} hook,
	% run *before* this whole macro body even starts) gets silently wiped
	% out by this same line a moment later -- confirmed directly via
	% \typeout debugging (not assumed): the hook's own \setlength really
	% did run with the right values, \itemsep still read back as 0pt
	% immediately after. Setting it *here*, right after \beamer@cramped
	% already ran (not before), is what actually survives.
	\ifcase\@itemdepth
	\or\setlength{\itemsep}{\gothamItemsepRatioI em}% level 1
	\or\setlength{\itemsep}{\gothamItemsepRatioII em}% level 2
	\or\setlength{\itemsep}{\gothamItemsepRatioIII em}% level 3
	\or\setlength{\itemsep}{\gothamItemsepRatioIV em}% level 4
	\fi%
	\raggedright%
	\beamer@firstlineitemizeunskip%
}
\makeatother

% A nested \begin{enumerate} inherits the same \leftmarginN step as itemize
% above for *where its text starts* (matching itemize's own "one step past
% the parent level's text" rule -- confirmed directly against the compiled
% PDF: an enumerate's own body text, e.g. "Senior secured debt", already
% lands exactly one \gothamItemStep past its parent item's text), but its
% *label* ("1.", "2.", ...) doesn't follow \labelwidth the way plain LaTeX's
% list mechanism would -- beamer draws enumerate labels through its own
% template machinery, which turned out not to respect \labelwidth at all:
% \setlength{\labelwidth}{...} right after \begin{enumerate} measurably took
% effect (checked with \typeout{\the\labelwidth} immediately after) but the
% compiled numeral's own position never moved. What *does* move it: directly
% overriding \makelabel with a fixed-width \makebox -- empirically
% calibrated (render, measure the numeral's own x-position in the compiled
% PDF, adjust, repeat) since the box's own left edge turned out to sit at a
% fixed position unrelated to \labelwidth or \leftmarginN, with the
% right-aligned numeral inside it moving in lockstep with the box's width
% (confirmed linear: +1pt width -> +~1pt numeral position, across several
% test widths). 19.23pt (at this deck's default itemize font size) lands a
% single-digit numeral flush with the parent item's own text -- only valid
% at that font size and for single-digit lists (1-9); re-calibrate the same
% way (not by recomputing from first principles) for a differently-sized or
% 10+-item nested enumerate. Call right after \begin{enumerate}, before the
% first \item.
%
% 20.33pt (2026-07-08, direct instruction: "the numbers still extend a
% little farther to the left than the text above them does" --
% slides/f6m-Recovery-Rates-current.tex's own "Pecking order:" list).
% 19.23pt landed *close* but not quite flush: measured directly against
% the compiled PDF, every numeral's own left edge (20.55pt-20.95pt,
% varying slightly by digit width even though all are right-aligned to
% the same box) still sat up to ~1.05pt left of "Pecking order:"'s own
% text start (21.6pt) -- the widest digit tried, "6.", was the binding
% case. +1.1pt of box width (confirmed +1pt width -> +1pt numeral
% position, per this comment's own note above) covers that with a small
% margin, re-confirmed against the recompiled PDF.
\newcommand{\gothamEnumerateLabelWidth}{\renewcommand{\makelabel}[1]{\makebox[20.33pt][r]{##1}}}

% Automatic version of the above (2026-07-23, background-task follow-up
% flagged in slides/kfh-Asset-Allocation-Process-current.tex's own v1-1-0
% comment: "fix the enumerate behavior generally, if possible", task
% `task_bb3fbb66`). That comment's first attempt, an
% `\AtBeginEnvironment{enumerate}{...}` hook, failed because the hook fires
% *before* `\enumerate`'s own body runs, and that body's own internal
% `\list{...}{\def\makelabel...}` setup (the beamer copy directly below the
% `\itemize` override above) clobbers whatever a hook that early leaves
% behind -- the exact same "template runs its own setting, external pre-set
% gets overwritten" problem `\itemize`'s own `\itemsep` comment above
% already documents. The real fix, same idea `\itemize` itself used above:
% override beamer's own low-level piece directly rather than hooking
% \AtBeginEnvironment. Unlike `\itemize` (declared via
% `\@protected@testopt`, needing a full re-declaration to reach -- see the
% `\itemize` override's own comment for the full story), `\enumerate`'s
% relevant piece, `\beamer@enum@` (beamerbaselocalstructure.sty), is a
% plain `\def` with no optional-argument indirection, so it can be
% `\renewcommand`'d directly, verbatim from upstream except for the
% `\ifnum\@itemdepth>1` branch added below.
%
% Only fires for a *nested* enumerate (`\@itemdepth>1`, i.e. already inside
% at least one `\itemize`/`\enumerate` when this one opens) -- a top-level
% enumerate (`\@itemdepth=1`, e.g.
% slides/t9z-Outland-Steel-Hedging-Borrowing-Investing-current.tex/
% slides/g4c-Outland-Steel-Hedging-Forward-Sale-current.tex/
% slides/j84-Outland-Steel-No-Hedging-current.tex's own outermost numbered
% list, or slides/nkj-Agenda-current.tex's own \tableofcontents-driven one)
% is left at beamer's stock label position entirely untouched, matching
% this project's own established behavior for that shape (never reported
% as misaligned, unlike the nested case) -- and, for nkj specifically,
% doesn't matter either way, since that slide's own
% `\gothamTocLabelFlushLeft` call (right after its own `\begin{enumerate}`)
% overrides `\makelabel` again regardless of whatever this default sets.
%
% Reuses the exact same 20.33pt calibration as the manual
% `\gothamEnumerateLabelWidth` above, confirmed (pymupdf, compiled PDF) to
% land correctly not just at `\@itemdepth=2` (one `\itemize` above the
% enumerate -- slides/kfh-Asset-Allocation-Process-current.tex/
% slides/f6m-Recovery-Rates-current.tex/
% slides/6pz-Outland-Steel-Hedging-Strategies-Overview-current.tex, all
% previously needing the manual call) but also at `\@itemdepth=3` (two
% `\itemize`s above it -- slides/bdh-Working-Capital-Example-current.tex's
% own "Development of balance sheet:" list, also previously needing the
% manual call, with the identical 20.33pt value) -- this deck's itemize/
% enumerate templates don't shrink font size by nesting depth (only the
% bullet glyph/margin step does), so the numeral's own glyph width -- the
% only thing \gothamEnumerateLabelWidth's own box width actually
% compensates for -- stays the same at every depth, and one constant covers
% both cases without needing a depth-dependent formula. Not verified beyond
% depth 3 (no slide currently nests an enumerate that deep); if one ever
% does, re-measure before trusting this constant that far down, same as any
% other empirical calibration in this file.
% (slides/scs-Political-Risk-Stability-Index-current.tex also matches
% "`\gothamEnumerateLabelWidth`" in a grep -- but only inside its own prose
% comments describing the general technique; its actual call is the
% bespoke `\gothamEnumerateLabelFlushLeftShifted{16pt}{10.2pt}` for its own
% alphabetic, left-aligned, narrower-box labels, unaffected by this
% override either way since that call runs after `\begin{enumerate}` opens,
% same ordering as every other post-open `\makelabel` override here.)
%
% Former per-slide manual `\gothamEnumerateLabelWidth` calls removed from
% the 4 slides above now that this covers them automatically; the macro
% itself is kept defined, both for this override to reuse and as a manual
% escape hatch (e.g. a differently-calibrated width some future slide might
% need).
\makeatletter
\renewcommand{\beamer@enum@}{%
	\beamer@computepref\@itemdepth% sets \beameritemnestingprefix
	\usebeamerfont{itemize/enumerate \beameritemnestingprefix body}%
	\usebeamercolor[fg]{itemize/enumerate \beameritemnestingprefix body}%
	\usebeamertemplate{itemize/enumerate \beameritemnestingprefix body begin}%
	\expandafter
		\list
			{\usebeamertemplate{\beamer@enumtempl}}
			{\usecounter\@enumctr%
				\ifnum\@itemdepth>1
					\def\makelabel##1{\makebox[20.33pt][r]{##1}}%
				\else
					\def\makelabel##1{{\hss\llap{{%
							\usebeamerfont*{enumerate \beameritemnestingprefix item}%
							\usebeamercolor[fg]{enumerate \beameritemnestingprefix item}##1}}}}%
				\fi}%
	\beamer@cramped%
	\raggedright%
	\beamer@firstlineitemizeunskip%
}
\makeatother

% Extra breathing room when a nested itemize closes and control returns to
% its *parent* list, right before that parent's next \item -- e.g. sub-items
% under "Mid-Term Exam" followed by the next top-level "Final Exam" -- so a
% new top-level group reads as visually distinct from the one before it,
% rather than the same uniform gap used between every other item regardless
% of level. \@itemdepth (LaTeX's own nesting-level counter, already
% decremented back to the parent's level by the time \end{itemize}'s normal
% closing has run) says which level control returned to -- 0 means the
% outermost list itself just closed (nothing left to add space before,
% deliberately excluded), 1 means back to a level-1 item, 2 means back to a
% level-2 item.
%
% Level-dependent (2026-07-06, direct instruction) -- a flat, single
% constant regardless of *which* level control returned to (the original
% version of this hook) reads as "extra"/inconsistent precisely because a
% return-to-level-1 and a return-to-level-2 transition both got the exact
% same gap, rather than level 1 getting a visibly bigger gap than level 2 --
% found on slides/57r-Reasons-for-Leasing-current.tex, where a level-2
% sub-list closing back to a level-1 item looked like a random extra gap
% rather than "level 1 items are consistently more separated than level 2
% ones." \addvspace (not plain \vspace) matters too: it coalesces with
% whatever glue itemize's own \itemsep/\topsep already contributes at the
% closing boundary (taking whichever is larger) rather than *adding* on top
% of it unconditionally -- \vspace's own stacking behavior was the other
% half of why the old flat version could look like too much space.
\makeatletter
% Ratios, not lengths (2026-07-21, direct instruction: "all of these
% spacings should scale with font size") -- a \newlength+\setlength{...}{0.5em}
% pair, the way this used to be written, freezes to a fixed *pt* value the
% one time \setlength runs (at preamble-load time, using whatever font is
% active then), despite the "em" in its own definition -- it does NOT
% re-track whichever font size is actually active on a given slide/level
% later. Confirmed directly: a \typeout{\f@size} placed at the point each
% of these actually gets used (inside \AtEndEnvironment{itemize}'s own
% hook body, further below) showed the *closing* list's own font still
% active there (9pt/10pt/10.95pt for a level-4/3/2 list closing,
% respectively -- not yet restored to the parent's own font), confirming
% this hook fires early enough in `\end{itemize}`'s own processing that a
% live "0.5em" computed *there* would correctly track whichever level is
% actually involved. Fixed by storing the bare ratio as a plain macro
% instead, and writing the literal "em" back on at the point of use
% (`\addvspace{\gothamItemizeLevelDropSpaceRatioI em}`), so it's
% recomputed fresh, against whatever font is live, every time.
\newcommand{\gothamItemizeLevelDropSpaceRatioI}{0.5}
\newcommand{\gothamItemizeLevelDropSpaceRatioII}{0.3}
% Level 3's own drop-space, for when a closed *4th*-level list returns
% control to a level-3 item -- added 2026-07-15 alongside the 4th itemize
% level itself (search "subsubsub" further above). Interpolated from the
% existing I/II progression (0.5, 0.3 -- each step roughly 60% of the one
% before) rather than empirically re-measured against a compiled PDF the
% way I/II themselves were: no slide used a 4th level before
% slides/hrg-Exchange-Rate-Quotation-Forward-current.tex to measure
% against. Re-derive properly (render, measure, adjust) if this ever
% looks visibly off on a real slide.
\newcommand{\gothamItemizeLevelDropSpaceRatioIII}{0.2}
% Level 4's own drop-space, for when a closed level-4 list returns control
% to a level-3 item -- added 2026-07-21 alongside \gothamItemsepRatioIV.
% \@itemdepth at this hook's fire time is the depth of the list that *just
% closed*, confirmed directly via a temporary \typeout{\the\@itemdepth}
% (not assumed): closing the level-4 list on
% slides/ZZO-Itemize-Spacing-Reference-current.tex measured itemdepth=4
% there, closing level-3 measured 3, level-2 measured 2, and the outermost
% level-1 list itself closing measured 1 (never 0 -- \@itemdepth's own
% restore to the parent's depth happens *after* this hook runs, so the "0:
% outermost list itself closed" case below is dead code, kept only because
% removing a genuinely unreachable \ifcase branch isn't worth the risk of
% getting the numbering wrong a second time). So this is very much a real,
% reachable case -- not decorative symmetry -- and needed its own
% \addvspace branch (previously missing entirely, silently leaving the
% level-4-to-level-3 closing transition with no drop-space at all).
% Interpolated from the existing I/II/III progression as a starting guess,
% then re-measured directly against the compiled ZZO reference and bumped:
% 0.13 measured at 12.6pt, still *less* than level 4's own same-level gap
% of 13.4pt (violates rule III -- a level's own same-level gap must be
% smaller than that level's own closing-transition gap); raised until
% re-measurement cleared it with a comfortable margin again, matching the
% margin every other level's own pair already has.
\newcommand{\gothamItemizeLevelDropSpaceRatioIV}{0.5}
\AtEndEnvironment{itemize}{%
	\ifcase\@itemdepth
		% 0: dead code -- @itemdepth is never actually 0 when this hook
		% fires (see \gothamItemizeLevelDropSpaceRatioIV's own comment
		% just above); kept only for alignment with the \ifcase branches
		% below, each keyed on the depth of the list that just closed.
	\or\addvspace{\gothamItemizeLevelDropSpaceRatioI em}% level-1 (outermost) list itself closed -- a pre-existing quirk (not introduced by this comment's own edit): there's no parent item above level 1 for this to visually separate from, so this only actually matters if something follows the outermost list directly on the same frame; harmless on every slide seen so far since nothing does. Out of scope to redesign here.
	\or\addvspace{\gothamItemizeLevelDropSpaceRatioII em}% level-2 list closed, back to a level-1 item
	\or\addvspace{\gothamItemizeLevelDropSpaceRatioIII em}% level-3 list closed, back to a level-2 item
	\or\addvspace{\gothamItemizeLevelDropSpaceRatioIV em}% level-4 list closed, back to a level-3 item
	\fi
}

% Per-level \itemsep -- the item-to-item gap *within* one list, at one
% fixed nesting depth (as opposed to the drop-space mechanism just above,
% which only fires when a nested list *closes* back up to its parent).
% Ratios, not lengths -- same "scale with font size" reasoning and fix as
% \gothamItemizeLevelDropSpaceRatioI/II/III above (that comment has the
% full story on why a plain \newlength+\setlength pair here would freeze
% to one fixed pt value instead of actually tracking each level's own
% font size). Values only, defined here for discoverability alongside the
% drop-space ones above -- actually *applied* inside the `\itemize`
% redefinition itself (search "\beamer@cramped (beamerbaselocalstructure.sty"
% further above), not via a separate hook here; see that comment for why.
% Found missing entirely (2026-07-21, direct instruction, see
% slides/ZZO-Itemize-Spacing-Reference-current.tex's own header comment
% for the rules this and the drop-space mechanism together need to
% satisfy): this deck's own itemize never set \itemsep itself at any
% level -- confirmed directly, grepping this whole file -- so every depth
% shared whatever bare default beamer's own theme happens to leave
% \itemsep at, uniformly. That produced two visible problems, both
% measured directly (pymupdf against slides/2c2-Topics-Around-the-CAPM-Empirical-Validity-current.tex
% and an isolated bisection test, not guessed): (a) level-2 and level-3
% same-level item-to-item gaps came out *identical* (12.6pt each,
% confirmed against an isolated reproduction), since nothing distinguished
% them by depth at all -- only a closing transition ever got a
% level-dependent gap, never a same-level one; (b) that shared default
% \itemsep read as *smaller* than a single item's own wrapped-line
% \baselineskip (13.7-13.8pt within a wrap vs. 12.6pt between items),
% backwards from how a list should read -- an item boundary should never
% be tighter than a same-item line wrap.
% Ratios themselves are not arbitrary -- each chosen, then re-measured
% against the compiled ZZO reference slide and adjusted, specifically to
% satisfy all three of that slide's own header-comment rules together
% (extended to level 4 the same way, 2026-07-21), not just "level 1
% biggest" in isolation:
% - Rule I (level-1 same-level gap > level-2's > level-3's > level-4's)
%   falls straight out of \gothamItemsepRatioI > ...II > ...III > ...IV,
%   added on top of the same \baselineskip every level shares (itself
%   already smaller per level here, since this theme's own itemize/
%   enumerate body font shrinks with depth too -- confirmed directly,
%   10.95pt/10pt/9pt/9pt for levels 1-4 respectively, level 4 sharing
%   level 3's own font since this theme has no distinct "subsubsub" size
%   of its own).
% - Rule II (every level's own same-level gap > that level's own wrapped-
%   line gap) needs \gothamItemsepRatioIV alone (the smallest, hence the
%   binding case) to still be positive and large enough that it clearly
%   reads as more open than a plain wrapped line once rendered, not just
%   bigger on paper.
% - Rule III (a level's own same-level gap < that same level's own
%   *closing*-transition gap) is *not* automatic just because
%   \gothamItemizeLevelDropSpaceRatio* was already bigger than the old
%   shared default -- \addvspace's own coalesce-to-the-larger-value
%   behavior (documented above) means the closing gap is
%   max(natural glue, that *drop-space* ratio), not a sum of the two. Each
%   drop-space ratio's own \ifcase branch is keyed on the depth of the
%   list that *just closed* (see \gothamItemizeLevelDropSpaceRatioIV's own
%   comment for the full story on why, confirmed via \typeout, not
%   assumed), so it lines up with the same-numbered itemsep ratio despite
%   describing a different transition -- \gothamItemizeLevelDropSpaceRatioIV
%   governs the level-4-list-closing-back-to-level-3 transition and is the
%   one that has to individually stay bigger than \gothamItemsepRatioIV.
%   This branch was missing entirely at first (only I/II/III existed, a
%   leftover from before level 4 existed) -- confirmed directly on ZZO: the
%   level-4-to-level-3 closing transition measured 12.6pt with no drop-space
%   branch at all, less than level 4's own same-level gap of 13.4pt, a real
%   rule III violation. Adding the branch (\gothamItemizeLevelDropSpaceRatioIV,
%   starting at 0.5) fixed it. Re-verify rule III directly against ZZO's own
%   compiled PDF (don't assume it still holds on paper) if any of these ratios is
%   ever changed again.
\newcommand{\gothamItemsepRatioI}{0.4}
\newcommand{\gothamItemsepRatioII}{0.24}
\newcommand{\gothamItemsepRatioIII}{0.15}
\newcommand{\gothamItemsepRatioIV}{0.09}

% Per-level "opening" space -- the gap between a level-n item and the
% first item of a level-(n+1) list that immediately follows it, as
% opposed to \gothamItemizeLevelDropSpaceRatio* (a nested list *closing*
% back up) or \gothamItemsepRatio* (two items at the *same* depth).
% Added 2026-07-21, direct instruction, replacing the original three
% separate rules on slides/ZZO-Itemize-Spacing-Reference-current.tex's
% own header comment with one multi-part inequality stated per level n:
%   level-n -> level-(n-1) close  >  level-n same-level  >
%   level-n -> level-(n+1) open  >  level-n's own wrap gap
% The "open" term was entirely unhandled before this -- confirmed
% directly on ZZO: every opening transition (level 1 into level 2, level
% 2 into level 3, level 3 into level 4) measured *smaller* than that same
% level's own wrap gap, backwards from what the rule needs, since nothing
% but beamer's own bare default \topsep governed it.
% \AtBeginEnvironment{itemize}, not \topsep -- \topsep was tried first
% (set via the same \ifcase\@itemdepth block, right after \itemsep,
% inside the `\itemize` redefinition above) and measured *unreliable*:
% two of the three levels barely moved despite a sizeable ratio, while
% the third moved by roughly the right amount -- \topsep's own value is
% evidently consulted by `\list` at some point this override comes too
% late for, at least for some levels, not evenly. \addvspace via
% \AtBeginEnvironment instead measured consistently and predictably
% (each level's own change matched its own ratio's calculated pt value
% directly), the same mechanism \gothamItemizeLevelDropSpaceRatio* already
% relies on for the closing side. This hook fires with \@itemdepth
% reflecting the *parent* item's own depth (confirmed via \typeout, not
% assumed: opening the outermost level-1 list itself measured 0, opening
% a level-2 list from within a level-1 item measured 1, level-3-from-
% level-2 measured 2, level-4-from-level-3 measured 3) -- the opposite of
% \gothamItemizeLevelDropSpaceRatio*'s own hook, which keys on the depth
% of the list that *just closed*. That makes the \ifcase numbering here
% match the "n" in the rule directly: case 1 is level-1's own open-to-
% level-2 ratio, case 2 is level-2's own open-to-level-3, case 3 is
% level-3's own open-to-level-4; case 0 (opening the outermost list
% itself) is a genuine no-op, correctly -- there's no parent item above
% level 1 for that transition to separate from.
% Ratios chosen, then re-measured against the compiled ZZO reference and
% adjusted, the same way as every other ratio here: each has to
% individually land strictly between that level's own \gothamItemsepRatio*
% (the ceiling) and its own wrap gap (the floor). Re-verify directly
% against ZZO's own compiled PDF (don't assume it still holds on paper)
% if any of these, or the itemsep/drop-space ratios they're measured
% against, are ever changed again.
% \gothamItemizeLevelOpenSpaceRatioI trimmed from 0.28 to 0.02 (2026-07-21,
% direct instruction: "the level-1 to level-2 distance...seems a bit
% much"). Confirmed directly: even at 0, this transition still measured
% ~16.3pt -- \addvspace's own coalesce-to-the-larger-value behavior means
% the ratio can only ever push this *above* whatever glue the list
% mechanism inserts here regardless (that natural floor, not the ratio),
% so 0.28 down to 0 only ever moved the actual gap by ~2pt (18.5->16.3);
% 0.02 splits the difference, keeping the mechanism nominally "in effect"
% rather than reading as accidentally disabled. \gothamItemizeLevelOpenSpaceRatioIII
% (level-3-to-level-4) had to come down too, from 0.08 to 0.02, purely to
% restore its own margin -- changing RatioI shifted that unrelated
% transition's own measured gap closer to level 3's own same-level gap
% (12.1 vs. 12.0, too thin), confirmed directly on ZZO, not assumed;
% ratios elsewhere in this file have shown the same small-ripple behavior
% before (see \gothamItemizeLevelDropSpaceRatioIV's own comment) --
% re-measure the *whole* table after touching any single one of these,
% not just the transition being changed.
\newcommand{\gothamItemizeLevelOpenSpaceRatioI}{0.02}
\newcommand{\gothamItemizeLevelOpenSpaceRatioII}{0.15}
\newcommand{\gothamItemizeLevelOpenSpaceRatioIII}{0.02}
\AtBeginEnvironment{itemize}{%
	\ifcase\@itemdepth
		% 0: opening the outermost (level-1) list itself -- no parent item
		% precedes it, nothing to add.
	\or\addvspace{\gothamItemizeLevelOpenSpaceRatioI em}% opening a level-2 list from a level-1 item
	\or\addvspace{\gothamItemizeLevelOpenSpaceRatioII em}% opening a level-3 list from a level-2 item
	\or\addvspace{\gothamItemizeLevelOpenSpaceRatioIII em}% opening a level-4 list from a level-3 item
	\fi
}
\makeatother

% A \tableofcontents-driven agenda slide (see
% slides/nkj-Agenda-current.tex) needs the exact same top-margin/line-
% spacing as every other list in this deck, not whatever gotham's own
% "section in toc" templates use (both "gotham bullet" and "gotham simple"
% carry their own spacing/rule/font choices, tuned for gotham's generic
% look, not this project's own itemize metrics). Rather than re-tuning a
% *third* set of constants to match, this template just emits a plain
% \item -- wrapping the whole \tableofcontents call in a real
% \begin{enumerate}...\end{enumerate} then makes each section entry a
% genuine enumerate item, reusing this file's own \leftmargini/\itemsep/
% \gothamContentTopShift machinery exactly as any other list would, so it
% automatically stays consistent with the rest of the deck instead of
% drifting out of sync with it later.
\defbeamertemplate{section in toc}{gotham plain enumerate}{\item \inserttocsection}

% Even with the plain-enumerate template above, \tableofcontents still
% doesn't match this deck's own normal itemize spacing -- beamer's own
% \beamer@sectionintoc (beamerbasetoc.sty), which wraps *every* section
% entry regardless of which "section in toc" template is active,
% unconditionally inserts a fixed \vskip1.5em before each entry (an
% unbounded \vfill, even worse, if subsections aren't hidden -- neither is
% reachable via a public \tableofcontents option, just this show/hide
% binary), *and* wraps the entry itself in \hbox{\vbox{...}} -- box
% wrapping that (like \gothamMinipageTopCorrection's own minipage case
% elsewhere in this file) changes how surrounding baselineskip/topskip
% glue gets inserted compared to plain paragraph text. Both together added
% roughly 13pt of extra gap
% between entries beyond this deck's own normal ~18.6pt level-1 item
% spacing (measured directly against
% slides/57r-Reasons-for-Leasing-current.tex's own compiled PDF).
% Neither is controllable via a public \tableofcontents option, so this
% redefines \beamer@sectionintoc itself -- scoped safely since
% \tableofcontents is only ever called on the one Agenda slide in this
% deck -- keeping its numbering/hyperlink/part-number/hide-checking logic
% intact but dropping the vskip/vfill and the hbox/vbox wrapping, so an
% entry behaves exactly like a plain \item once printed.
\makeatletter
\def\beamer@sectionintoc#1#2#3#4#5{%
	\ifnum\c@tocdepth>0%
	\ifnum#4=\beamer@showpartnumber%
	{%
	\beamer@saveanother%
	\gdef\beamer@todo{}%
	\beamer@slideinframe=#1\relax%
	\expandafter\only\beamer@tocsections{\gdef\beamer@todo{%
			\beamer@tempcount=#5\relax%
			\advance\beamer@tempcount by\beamer@sectionadjust%
			\ifnum\beamer@tempcount>0
				\ifnum\beamer@tempcount>\beamer@toclastsection
					\def\inserttocsectionnumber{}%
				\else
					\edef\inserttocsectionnumber{\the\beamer@tempcount}%
				\fi
			\else
				\def\inserttocsectionnumber{}%
			\fi%
			\def\inserttocsection{\hyperlink{Navigation#3}{#2}}%
			\beamer@tocifnothide{\ifnum\c@section=#1\beamer@toc@cs\else\beamer@toc@os\fi}%
			{%
				\beamer@tocact{\ifnum\c@section=#1\beamer@toc@cs\else\beamer@toc@os\fi}{section in toc}%
			}%
		}%
	}%
	\beamer@restoreanother%
	}%
	\beamer@todo%
	\fi\fi%
}
\makeatother

% Left-aligns an enumerate's labels within a caller-chosen \labelwidth --
% used by the \tableofcontents Agenda slide (see
% slides/nkj-Agenda-current.tex's own comment for why its labels need
% this at 3.6pt). Defined here as a named helper rather than written as a
% literal \renewcommand{\makelabel}[1]{...#1...} inside the slide's own
% frame body, because that literal form -- specifically its raw `#1`
% parameter token -- breaks inside any frame that \tableofcontents forces
% through beamer's overlay reprocessing: beamer replays such a frame's
% collected body via \def-based machinery (\loop's own \iterate, and
% \beamer@doifinframe), and a `#` in the material of any \def must be
% doubled -- so the compile hits "! Illegal parameter number in definition
% of \iterate." (twice) plus the same for \beamer@doifinframe, once per
% compile, at that frame's own \end{frame}. TeX's error recovery ("I'll
% assume you meant ##") happens to produce the intended definition anyway,
% which is why the rendered PDF was always correct despite the errors --
% but lualatex still exits nonzero, which is exactly what kept breaking
% latexmk's automated lualatex/biber/rerun sequencing on the main deck
% (the "Known open issue" long tracked in CONVERSION_NOTES.md, originally
% misattributed to a biblatex interaction -- see its "Resolved" entry).
% Hoisting the definition out here means the frame body carries only this
% macro's *call* token, with no `#` in it, which sidesteps the whole
% mechanism. The inner \makelabel parameter is written `##1` below since
% it sits inside another \newcommand's body -- standard nesting, and also,
% incidentally, exactly the doubling rule the inline version fell foul of.
\newcommand{\gothamTocLabelFlushLeft}[1]{%
	\setlength{\labelwidth}{#1}%
	\renewcommand{\makelabel}[1]{\makebox[\labelwidth][l]{##1}}%
}

% Same idea as \gothamTocLabelFlushLeft above, plus an extra fixed
% \hspace shift before the label (for a list that needs to start further
% right than the plain nested-list indent, e.g. to line up with a
% sibling item's own text rather than its bullet -- see
% slides/scs-Political-Risk-Stability-Index-current.tex's own v1-5-0/
% v1-6-0 comments for the worked example this was hoisted out of).
% Found 2026-07-20 that the raw-`#1`-in-frame-body bug documented above
% is *not* actually \tableofcontents-specific, contrary to that comment's
% own framing -- slides/scs-Political-Risk-Stability-Index-current.tex has
% no \tableofcontents anywhere in the deck-wide overlay chain leading to
% it, yet hit the identical "Illegal parameter number in definition of
% \iterate" at its own \end{standardframe} from the exact same literal-
% \renewcommand{\makelabel}[1]{...#1...}-in-frame-body pattern. Whatever
% the precise trigger condition actually is, the same hoist-to-preamble
% fix (a bare macro *call*, no `#` token, left in the frame body) sidesteps
% it either way.
\newcommand{\gothamEnumerateLabelFlushLeftShifted}[2]{%
	\setlength{\labelwidth}{#2}%
	\renewcommand{\makelabel}[1]{\hspace{#1}\makebox[\labelwidth][l]{##1}}%
}

% Width of a single digit in this deck's normal content font, measured once
% here (safe, plain preamble-load-time \settowidth) rather than inline in
% the Agenda slide's own frame body -- \tableofcontents there puts the whole
% frame through beamer's overlay-replay machinery, which the comment above
% \gothamTocLabelFlushLeft already found chokes on raw tokens (like a bare
% \settowidth call) sitting directly in a replayed frame body, the same
% "Missing number"/"Illegal unit of measure" class of error as that `#`
% problem, confirmed directly the same way. Used by
% slides/nkj-Agenda-current.tex's own "widen the number-to-title gap"
% fix, which only ever needs to *reference* this precomputed value inside
% the frame, not measure anything itself there.
\newlength{\gothamTocLabelExtraGap}
\settowidth{\gothamTocLabelExtraGap}{0}

% For a two-column layout (e.g. bullets next to a chart/formula), nest
% beamer's own "columns" directly inside standardframe -- do NOT bake
% \begin{columns}/\end{columns} into a custom environment's begin/end code:
% beamer collects each frame's body as a token list to replay it once per
% overlay, and a columns environment that is opened in the "before" half and
% closed in the "after" half of a *different* environment confuses that
% collection (it never finds a properly matched \end{frame}). Nesting it in
% the document body, as below, avoids the issue entirely:
%
%   \begin{standardframe}{Title}
%   	\begin{columns}[t,onlytextwidth]
%   		\begin{column}{0.48\linewidth} ... \end{column}
%   		\begin{column}{0.48\linewidth} ... \end{column}
%   	\end{columns}
%   \end{standardframe}

% For automatic section dividers and single-statement "standout" slides, use
% what the gotham theme already provides instead of a custom environment:
%   \section{...}                                                (section divider, automatic)
%   \begin{frame}[standout]{Key takeaway}\end{frame}             (big single statement)
%   \begin{block}{...}/\begin{alertblock}{...}/\begin{exampleblock}{...}/\begin{referenceblock}{...}  (definitions, formulas, examples, references)
% For a course's own title slide, section-topic overview, quote, figure-only,
% and presenter-attention slides, see the five environments below.

% block/alertblock/exampleblock recolored to this project's own palette
% (2026-07-05, direct instruction) instead of gotham's own default colors
% -- gotham's own per-block-type color machinery (\gotham@block in
% beamerinnerthemegotham.sty) turns out to be dead code, never actually
% wired into any template, so these three blocks were rendering via plain
% beamer's own stock templates/colors the whole time; overriding the
% beamer color names directly here works regardless, since those stock
% templates read their colors from these same standard names (block~title/
% block~body and their ~alerted/~example variants) via beamercolorbox.
% Body background is the same c_tertiary_gray for all four
% block types (including referenceblock below) -- only each one's own
% title/header background differs. block's own header went
% gray→secondaryblue→secondaryviolet across three separate direct
% instructions in as many days -- secondaryviolet is the one actually in
% effect now.
\setbeamercolor{block title}{bg=c_secondary_violet, fg=c_primary_dark}
\setbeamercolor{block body}{bg=c_tertiary_gray, fg=c_primary_dark}
\setbeamercolor{block title example}{bg=c_secondary_green, fg=c_primary_dark}
\setbeamercolor{block body example}{bg=c_tertiary_gray, fg=c_primary_dark}
\setbeamercolor{block title alerted}{bg=c_secondary_red, fg=c_primary_dark}
\setbeamercolor{block body alerted}{bg=c_tertiary_gray, fg=c_primary_dark}

% All four block types' left/right margins widened to match a top-level
% itemize bullet's own text indent (2026-07-05, direct instruction: "the
% blocks should not have a smaller left-hand margin than the content
% text") -- a plain block spans \linewidth as-is, which lines up with the
% *frame's* own text margin, not with an itemize bullet's own text (which
% sits further right, past the bullet symbol and its gap) -- measured
% directly against a compiled slide with both an itemize and a block
% (slides/afa-Excursus-Price-Formation-current.tex): a bullet's own text
% starts 12.96pt further right than a plain block's left edge did.
%
% First attempt used changepage's \adjustwidth, on paper the "correct" tool
% for temporarily changing margins around a chunk of content -- it silently
% does *nothing* inside a beamer frame (beamer computes \linewidth/\hsize
% through its own internal frame-geometry machinery, not the page-geometry
% registers changepage adjusts, which are meant for article/report-style
% documents), and *loading the package at all* actively corrupted the
% title bar/logo on every frame (some interaction with gotham's own
% tikz/remember-picture-based title-bar drawing, not chased further since
% the package needed to go regardless) -- caught by actually compiling and
% looking, not assumed to work from the package's own documented purpose.
% Fixed instead with a plain \hspace{}+\begin{minipage}{\linewidth-2\times}
% wrapped around each block type's own begin/end template pair below --
% minipage genuinely resets \linewidth to its own local width for anything
% typeset inside it (already relied on elsewhere in this deck, e.g. the
% side-by-side text/image slides), which a bare beamercolorbox alone
% doesn't do on its own.
%
% 12.96pt matched a *second*-level bullet's own dot (which happens to
% coincide almost exactly with where a *first*-level bullet's own *text*
% starts, both ~21.6pt in from the frame's raw margin -- a coincidence of
% this deck's specific bullet-dot/gap/font metrics, not a general rule) --
% corrected the next day (direct instruction: align with the *first*-level
% bullet's own margin instead) to 4.20pt, matching a first-level bullet's
% own dot position (~12.84pt in from the frame's raw margin) instead.
\newlength{\gothamBlockSideMargin}
\setlength{\gothamBlockSideMargin}{4.20pt}

% Shared helper for a block title bar that carries a marker symbol flush
% right, opposite the title text (currently: infoblock's circled "i",
% definitionblock's $\triangleq$, exampleblock's $\in$ -- see each of
% those three's own template for the direct instruction behind its own
% marker choice). 2026-07-08, direct instruction, after infoblock's own
% "circled i nearly touching the title text" report on a two-line title
% (slides/u5h-Long-term-Bonds-Provisions-current.tex's own "Denomination,
% face value, or principal" step) -- generalized here so the same fix
% covers all three marker types at once, and any future one, rather than
% patching each block's own template separately.
%
% The marker itself is placed via a zero-width box *before* the title
% text, not via a trailing `\hfill<marker>` the way all three started out
% (definitionblock/exampleblock still do, until updated below) -- `\hfill`
% only fills space on whatever *line* it ends up typeset on, so for a
% title long enough to wrap onto a second line, `\hfill<marker>` lands on
% that second line instead of the first, printing the marker at the
% bottom right of the box instead of the top right (this is the same root
% cause infoblock's own marker already had fixed once already, see that
% template's own comment for the fuller writeup of *that* bug specifically
% -- this pushes the same fix out to the other two, and adds the part
% that bug's own fix didn't yet cover, below).
% `\makebox[0pt][l]{\hspace{\linewidth}\makebox[0pt][r]{<marker>}}` fixes
% *that* much on its own: the outer zero-width box doesn't displace the
% title text at all (still starts at the same point, still wraps however
% it needs to), but *inside* that zero-width box the marker is placed at
% a fixed horizontal offset (this box's own full \linewidth) from the
% title's own starting point, i.e. always the top right corner of the
% title bar regardless of how the title text itself wraps.
%
% What that alone still doesn't fix -- and what this round's actual
% report was about -- is the title text itself having no *awareness* that
% a marker will render on top of wherever it happens to end: a title just
% long enough to reach close to the marker's own fixed position, without
% actually needing to wrap, ends up visually touching or running under
% it. Fixed by also wrapping the title text itself in a
% `\begin{minipage}[t]{\linewidth-\gothamBlockTitleMarkerWidth}`, one
% \gothamBlockTitleMarkerWidth narrower than the full title bar, so the
% text itself wraps *before* reaching the marker's own reserved space
% instead of running into it -- confirmed directly against the compiled
% PDF on the exact case from the report. \gothamBlockTitleMarkerWidth is
% one shared reserved width for all three current marker types rather
% than a separate hand-measured width per symbol, which would need
% re-measuring by hand every time a new marker is added -- comfortably
% covers the widest of the three (checked directly) with a bit to spare.
% #1 = the marker's own content (e.g. \textcircled{\scriptsize i}, or a
% math-mode symbol like $\triangleq$); #2 = the title text
% (\insertblocktitle). Call as
% \gothamBlockTitleWithMarker{<marker>}{\insertblocktitle} in place of a
% bare \insertblocktitle, after the usual \usebeamerfont* call.
\newlength{\gothamBlockTitleMarkerWidth}
\setlength{\gothamBlockTitleMarkerWidth}{16pt}
% 2026-07-08 (direct instruction), two more fixes to this same helper:
% - A title long enough to wrap onto a second line
%   (slides/u5h-Long-term-Bonds-Provisions-current.tex's own
%   "Denomination, face value, or principal") read as *justified* rather
%   than left-aligned -- the minipage below had no explicit alignment, so
%   it fell back to LaTeX's default full justification, which only ever
%   visibly stretches a line short of the box's own right edge, i.e. only
%   shows up once a title is long enough to actually wrap. Every other
%   (single-line, shorter) title already looked left-aligned regardless,
%   since a line that's naturally shorter than the box is never stretched
%   whether or not justification is on -- confirmed directly, this is why
%   the bug went unnoticed until a title finally wrapped. `\raggedright`
%   fixes it without needing anything different for the single-line case.
% - Different titles' own title *bars* rendered at different heights
%   depending on whether that particular title's text happened to include
%   a descender (e.g. "Security"/"y", "Moody's rating at issue date"/"y")
%   or not (e.g. "Registered", "Trustee") -- the same class of bug already
%   fixed on slides/qxr-Introduction-to-R-Assignments-current.tex's own
%   colorboxes (see that slide's own comment): a box that tightly wraps its
%   own text naturally varies in height/depth with that text's own
%   ascenders/descenders. A leading `\strut` gives every title's first (and
%   for all but this one current title, only) line the same guaranteed
%   minimum height/depth, matching the qxr fix's own reasoning -- confirmed
%   directly, every current info/definition/example block title bar now
%   renders at the same height regardless of its own text.
\newcommand{\gothamBlockTitleWithMarker}[2]{%
	\makebox[0pt][l]{\hspace{\linewidth}\makebox[0pt][r]{#1}}%
	\begin{minipage}[t]{\dimexpr\linewidth-\gothamBlockTitleMarkerWidth\relax}\raggedright\strut#2\end{minipage}%
}

% A fourth block type, for a reference/citation-style callout (first used
% on slides/afa-Excursus-Price-Formation-current.tex and
% slides/4r7-IPO-Underpricing-International-current.tex, both replacing
% an earlier rotated hand-drawn callout box) -- beamer has no built-in
% "reference" block variant, so this defines one from scratch, mirroring
% exactly how beamer's own alertblock/exampleblock are built (compare
% beamerbaselocalstructure.sty's \newenvironment<>{alertblock} and
% beamerinnerthemedefault.sty's "block alerted begin/end" templates).
%
% First attempt (2026-07-05) copied these under the "rounded" style/
% beamerboxesrounded, on the wrong assumption that "rounded" was this
% deck's active block style generally -- it renders with visibly rounded
% corners, caught by the user, unlike block/alertblock/exampleblock's own
% square corners. Those three actually use beamerinnerthemedefault.sty's
% *starred* "default"-style templates (plain beamercolorbox, no rounding),
% which \defbeamertemplate*'s star selects immediately when that inner
% theme loads, regardless of beamer's own \defbeamertemplateparent{blocks}
% fallback to "rounded" -- gotham's own inner theme derives from
% beamerinnerthemedefault, so "default" (square) was the deck's real active
% style the whole time. Rewritten below to match that "default" style
% exactly (plain beamercolorbox instead of beamerboxesrounded) for
% consistent square corners across all four block types.
\setbeamercolor{block title reference}{bg=c_secondary_orange, fg=c_primary_dark}
\setbeamercolor{block body reference}{bg=c_tertiary_gray, fg=c_primary_dark}
% block/alertblock/exampleblock's own titles render bold via
% beamerfontthemegotham.sty's \setbeamerfont{block title}{series=\bfseries}
% (and an explicit ...alerted equivalent; example inherits it from
% beamer's own base font-fallback chain, verified directly -- both
% "Example"/"Caveat" titles on slides/ZZX-Definition-formula-current.tex
% render in the same heavier FiraSans-Regular as "Definition" despite
% gotham never setting "block title example" explicitly). "reference" is a
% brand new type name unknown to any of that inheritance, so it fell back
% to a plain, non-bold font with no color-analogous safety net (unlike
% \setbeamercolor{block title reference} above, colors for genuinely new
% names still resolve fine) -- caught only by comparing the actual
% rendered font name in the compiled PDF (FiraSans-Light, matching plain
% body text) against block/alertblock/exampleblock's own (FiraSans-Regular,
% this deck's de facto "bold" for Fira Sans), not by eye alone. Explicit
% setting needed, matching gotham's own two.
\setbeamerfont{block title reference}{size=\normalsize, series=\bfseries}
\defbeamertemplate*{block reference begin}{default}[1][]
{
	\hspace{\gothamBlockSideMargin}\begin{minipage}{\dimexpr\linewidth-2\gothamBlockSideMargin\relax}
	\par\vskip\medskipamount%
	\begin{beamercolorbox}[colsep*=.75ex,#1]{block title reference}
		\usebeamerfont*{block title reference}\strut\insertblocktitle%
	\end{beamercolorbox}\noframebreak%
	{\parskip0pt\par}%
	\ifbeamercolorempty[bg]{block title reference}
	{}
	{\ifbeamercolorempty[bg]{block body reference}{}{\nointerlineskip\vskip-0.5pt}}%
	\usebeamerfont{block body reference}%
	\begin{beamercolorbox}[colsep*=.75ex,vmode,#1]{block body reference}%
		\ifbeamercolorempty[bg]{block body reference}{\vskip-.25ex}{\vskip-.75ex}\vbox{}%
}
\defbeamertemplate*{block reference end}{default}
{\end{beamercolorbox}\vskip\smallskipamount\end{minipage}}
\newenvironment<>{referenceblock}[1]{%
	\begin{actionenv}#2%
		\def\insertblocktitle{#1}%
		\par%
		\usebeamertemplate{block reference begin}}
	{\par%
		\usebeamertemplate{block reference end}%
	\end{actionenv}}

% A sixth block type, for general background/context info that isn't
% itself a definition, example, or reference citation (2026-07-06, direct
% instruction: "use secondaryblue for a new type of info block") -- built
% the same from-scratch way as referenceblock just above (a brand-new
% style name needs \defbeamertemplate*'s registering form, not
% \setbeamertemplate -- see referenceblock's own comment for why). Marker
% is a circled "i" (\textcircled{\scriptsize i}, the generic LaTeX-kernel
% \textcircled -- not a Unicode "circled i" character: this deck's own
% Fira Sans/Fira Mono have no glyph for that codepoint either, the same
% missing-glyph problem \triangleq's own comment already found for
% definitionblock, and \textcircled sidesteps it entirely since it just
% draws a plain circle symbol around an ordinary letter rather than
% depending on one precomposed font glyph existing at all).
%
% The marker is placed via a zero-width box *before* \insertblocktitle,
% not via a trailing \hfill\textcircled{...} the way it first looked most
% natural to write -- \hfill only fills space on whatever line it's
% typeset on, and for a title long enough to wrap onto a second line (2026-07-08,
% found on the last reveal step of slides/ib3-Long-term-Bonds-Repayment-current.tex),
% \hfill\textcircled{...} then lands on that *second* line instead of the
% first, printing the marker at the bottom right of the box instead of the
% top right. Fixed originally with a one-off zero-width-box trick here;
% generalized the same day into the shared \gothamBlockTitleWithMarker
% helper (see its own comment, near \gothamBlockSideMargin above) once a
% *second*, related bug turned up on this exact block type (a long title
% visually touching the marker even without wrapping) -- that helper is
% what's actually called below now, not a copy of this trick inline.
\setbeamercolor{block title info}{bg=c_secondary_blue, fg=c_primary_dark}
\setbeamercolor{block body info}{bg=c_tertiary_gray, fg=c_primary_dark}
\setbeamerfont{block title info}{size=\normalsize, series=\bfseries}
\defbeamertemplate*{block info begin}{default}[1][]
{
	\hspace{\gothamBlockSideMargin}\begin{minipage}{\dimexpr\linewidth-2\gothamBlockSideMargin\relax}
	\par\vskip\medskipamount%
	\begin{beamercolorbox}[colsep*=.75ex,#1]{block title info}
		\usebeamerfont*{block title info}%
		\gothamBlockTitleWithMarker{\textcircled{\scriptsize i}}{\insertblocktitle}%
	\end{beamercolorbox}\noframebreak%
	{\parskip0pt\par}%
	\ifbeamercolorempty[bg]{block title info}
	{}
	{\ifbeamercolorempty[bg]{block body info}{}{\nointerlineskip\vskip-0.5pt}}%
	\usebeamerfont{block body info}%
	\begin{beamercolorbox}[colsep*=.75ex,vmode,#1]{block body info}%
		\ifbeamercolorempty[bg]{block body info}{\vskip-.25ex}{\vskip-.75ex}\vbox{}%
}
\defbeamertemplate*{block info end}{default}
{\end{beamercolorbox}\vskip\smallskipamount\end{minipage}}
\newenvironment<>{infoblock}[1]{%
	\begin{actionenv}#2%
		\def\insertblocktitle{#1}%
		\par%
		\usebeamertemplate{block info begin}}
	{\par%
		\usebeamertemplate{block info end}%
	\end{actionenv}}

% A fifth block type, for a "this introduces/defines a term" callout,
% visually distinguished from a plain \block (2026-07-06, direct
% instruction: "I would like to distinguish definition blocks from other
% blocks") by a $\triangleq$ symbol sitting flush right in the title bar,
% opposite the title text itself. Exact same colors/margins/font as plain
% \block (reuses "block title"/"block body", not a new color pair) --
% only the title-bar content differs, one extra `\hfill$\triangleq$`
% after `\insertblocktitle` -- so this mirrors referenceblock's own
% from-scratch pattern above (a brand new "definition" style name via
% \defbeamertemplate*, not \setbeamertemplate, for the same reason
% referenceblock's own comment gives: this style name has never been
% registered before, so the *-registration form is the correct one here,
% unlike plain block/alertblock/exampleblock below which are redefining
% already-registered "default"-style bodies). `\triangleq` is a genuine
% math symbol, so it needs `$...$` -- renders in this theme's math font
% (Computer Modern, see beamerfontthemegotham.sty's `no-math` option),
% a deliberate, visible one-glyph exception in an otherwise Fira Sans
% title bar, the same tradeoff every other math symbol in this deck
% already accepts.
\defbeamertemplate*{block definition begin}{default}[1][]
{
	\hspace{\gothamBlockSideMargin}\begin{minipage}{\dimexpr\linewidth-2\gothamBlockSideMargin\relax}
	\par\vskip\medskipamount%
	\begin{beamercolorbox}[colsep*=.75ex,#1]{block title}
		\usebeamerfont*{block title}\gothamBlockTitleWithMarker{$\triangleq$}{\insertblocktitle}%
	\end{beamercolorbox}\noframebreak%
	{\parskip0pt\par}%
	\ifbeamercolorempty[bg]{block title}
	{}
	{\ifbeamercolorempty[bg]{block body}{}{\nointerlineskip\vskip-0.5pt}}%
	\usebeamerfont{block body}%
	\begin{beamercolorbox}[colsep*=.75ex,vmode,#1]{block body}%
		\ifbeamercolorempty[bg]{block body}{\vskip-.25ex}{\vskip-.75ex}\vbox{}%
}
\defbeamertemplate*{block definition end}{default}
{\end{beamercolorbox}\vskip\smallskipamount\end{minipage}}
\newenvironment<>{definitionblock}[1]{%
	\begin{actionenv}#2%
		\def\insertblocktitle{#1}%
		\par%
		\usebeamertemplate{block definition begin}}
	{\par%
		\usebeamertemplate{block definition end}%
	\end{actionenv}}

% block/alertblock/exampleblock begin/end: same body as
% beamerinnerthemedefault.sty's own "default"-style templates (the ones
% actually in effect, per the comment above referenceblock), plus the
% minipage margin wrapper around each pair -- set via plain
% \setbeamertemplate{block begin}{...} (overriding the *body* of the
% existing "default" style directly) rather than
% \defbeamertemplate*{block begin}{default}{...} like referenceblock's own
% (new) templates above. The two look almost interchangeable but aren't:
% \defbeamertemplate* internally does a \newcommand-based *registration* of
% a style name, which errors outright ("... already defined") when that
% exact template+style pair already exists, as "block begin"/"default" and
% the other two pairs already do (registered once by
% beamerinnerthemedefault.sty itself) -- caught only by compiling and
% reading the log, not from any visible symptom in the rendered PDF (LaTeX
% recovers from the error by leaving the *original* stock template body in
% place, so these three blocks kept working and even picked up their new
% *colors* correctly via \setbeamercolor above, since color lookups happen
% at render time either way -- just silently without the margin fix, which
% *is* baked into the template body itself). \setbeamertemplate{name}{body}
% (no style-name argument) instead does a plain \def of that style's
% existing body storage, which is exactly "redefine what's already there"
% and carries no such restriction -- the correct tool whenever the target
% template+style pair is already registered, as opposed to
% \defbeamertemplate*, which is only for registering a style name for the
% very first time (as referenceblock's brand-new "reference" style name
% needs). Dropped the stock templates' own "#1" optional-argument slot
% (extra beamercolorbox options) in the process: nothing in this deck ever
% calls \usebeamertemplate{block begin}[...] with a bracket argument to
% populate it, so it was always empty in practice anyway.
% \strut before \insertblocktitle here, and in "block alerted"/"block
% reference" below: same title-bar-height-vs-descender fix as
% \gothamBlockTitleWithMarker's own comment above -- these three don't go
% through that helper at all (no marker symbol), so needed the same
% `\strut` added directly instead.
\setbeamertemplate{block begin}{
	\hspace{\gothamBlockSideMargin}\begin{minipage}{\dimexpr\linewidth-2\gothamBlockSideMargin\relax}
	\par\vskip\medskipamount%
	\begin{beamercolorbox}[colsep*=.75ex]{block title}
		\usebeamerfont*{block title}\strut\insertblocktitle%
	\end{beamercolorbox}\noframebreak%
	{\parskip0pt\par}%
	\ifbeamercolorempty[bg]{block title}
	{}
	{\ifbeamercolorempty[bg]{block body}{}{\nointerlineskip\vskip-0.5pt}}%
	\usebeamerfont{block body}%
	\begin{beamercolorbox}[colsep*=.75ex,vmode]{block body}%
		\ifbeamercolorempty[bg]{block body}{\vskip-.25ex}{\vskip-.75ex}\vbox{}%
}
\setbeamertemplate{block end}
{\end{beamercolorbox}\vskip\smallskipamount\end{minipage}}

% Warning-triangle marker for alertblock (direct instruction, 2026-07-12),
% same \gothamBlockTitleWithMarker treatment as definitionblock's
% $\triangleq$/exampleblock's $\in$/infoblock's \textcircled{\scriptsize i}
% above. No single built-in glyph renders a triangle-with-exclamation-mark
% in this deck's fonts (checked the same way \triangleq's own comment
% did), so it's hand-drawn: a plain tikz outline triangle with a bold "!"
% centered inside, sized by eye against \textcircled{\scriptsize i} until
% the two read as the same visual weight in a block title bar.
\newcommand{\gothamAlertMarker}{%
	\tikz[baseline=-0.6ex]{
		\draw[line width=1.1pt, line join=round] (-0.2,-0.11) -- (0.2,-0.11) -- (0,0.31) -- cycle;
		\node[font=\fontsize{7}{7}\selectfont\bfseries] at (0,0.01) {!};
	}%
}
\setbeamertemplate{block alerted begin}{
	\hspace{\gothamBlockSideMargin}\begin{minipage}{\dimexpr\linewidth-2\gothamBlockSideMargin\relax}
	\par\vskip\medskipamount%
	\begin{beamercolorbox}[colsep*=.75ex]{block title alerted}
		\usebeamerfont*{block title alerted}\gothamBlockTitleWithMarker{\gothamAlertMarker}{\insertblocktitle}%
	\end{beamercolorbox}\noframebreak%
	{\parskip0pt\par}%
	\ifbeamercolorempty[bg]{block title alerted}
	{}
	{\ifbeamercolorempty[bg]{block body alerted}{}{\nointerlineskip\vskip-0.5pt}}%
	\usebeamerfont{block body alerted}%
	\begin{beamercolorbox}[colsep*=.75ex,vmode]{block body alerted}%
		\ifbeamercolorempty[bg]{block body alerted}{\vskip-.25ex}{\vskip-.75ex}\vbox{}%
}
\setbeamertemplate{block alerted end}
{\end{beamercolorbox}\vskip\smallskipamount\end{minipage}}

% $\in$ ("is an instance/element of") flush right in the title bar
% (2026-07-06, direct instruction, mirroring definitionblock's own
% $\triangleq$ -- "what symbol would you propose ... to symbolize an
% example?"): a concrete example is naturally "an instance of" the general
% concept it illustrates, the same relationship $\in$ already denotes
% between an element and a set/class -- and, same as $\triangleq$, it's a
% math-mode symbol with a real Computer Modern glyph (no Fira Sans/Fira
% Mono text-mode substitute needed, unlike most Unicode symbol candidates
% -- see Slidemaster_preamble_uni-graz.tex's own "beamer/TeX gotchas" entry on
% $\triangleq$ for why that matters here). Applied to every \exampleblock/
% \gothamfloatingexampleblock in the deck via this shared template, not a
% one-off per slide, matching how $\triangleq$ applies to every
% definitionblock through its own shared template.
\setbeamertemplate{block example begin}{
	\hspace{\gothamBlockSideMargin}\begin{minipage}{\dimexpr\linewidth-2\gothamBlockSideMargin\relax}
	\par\vskip\medskipamount%
	\begin{beamercolorbox}[colsep*=.75ex]{block title example}
		\usebeamerfont*{block title example}\gothamBlockTitleWithMarker{$\in$}{\insertblocktitle}%
	\end{beamercolorbox}\noframebreak%
	{\parskip0pt\par}%
	\ifbeamercolorempty[bg]{block title example}
	{}
	{\ifbeamercolorempty[bg]{block body example}{}{\nointerlineskip\vskip-0.5pt}}%
	\usebeamerfont{block body example}%
	\begin{beamercolorbox}[colsep*=.75ex,vmode]{block body example}%
		\ifbeamercolorempty[bg]{block body example}{\vskip-.25ex}{\vskip-.75ex}\vbox{}%
}
\setbeamertemplate{block example end}
{\end{beamercolorbox}\vskip\smallskipamount\end{minipage}}

% A narrower, freely-positionable version of \begin{block}{...} (2026-07-06,
% direct instruction) -- a plain block always spans the *current*
% \linewidth (that's what block begin/end's own \hspace+minipage wrapper
% above is measured against), so it can't be dropped into a fixed-width
% spot (e.g. beside a table, inside a tikz overlay \node) without first
% narrowing \linewidth some other way. \gothamfloatingblock just wraps a
% plain block in one more, *outer* minipage of the caller's own chosen
% width -- block begin/end's own inner minipage then measures itself
% against that narrower width instead of the frame's full one, nesting
% correctly with no changes needed to the block templates themselves.
% Title/body colors, margins, and font are exactly \block's own (same
% "block title"/"block body" beamer color names) -- this only changes the
% overall width and (by taking it out of the normal document flow, left to
% the caller) the position, not the look. Font size is also left to the
% caller (wrap the whole environment in a size declaration, e.g.
% {\footnotesize \begin{gothamfloatingblock}{5cm}{Excel} ... \end{gothamfloatingblock}})
% rather than fixed here, since the right size depends entirely on how
% much room the block ends up placed in.
\newenvironment{gothamfloatingblock}[2]{%
	\begin{minipage}{#1}
	\begin{block}{#2}
}{%
	\end{block}
	\end{minipage}
}

% Same idea, but wrapping \exampleblock instead of \block (2026-07-06,
% direct instruction, first used on
% slides/fr3-Forms-of-Capital-Market-Efficiency-Strong-Semistrong-Weak-current.tex
% and slides/ha9-Monacorp-Truck-Lease-Breakeven-Rent-current.tex, both
% switching an existing \gothamfloatingblock to the "example" green header
% instead of the plain violet one) -- not worth generalizing into a single
% macro parameterized over the block name: \gothamfloatingblock's own body
% is only two lines, so a third variant is cheaper to just write out than
% to build a block-name-dispatch layer for.
\newenvironment{gothamfloatingexampleblock}[2]{%
	\begin{minipage}{#1}
	\begin{exampleblock}{#2}
}{%
	\end{exampleblock}
	\end{minipage}
}

% Same idea again, wrapping \infoblock (2026-07-06, direct instruction,
% first used on
% slides/fr3-Forms-of-Capital-Market-Efficiency-Strong-Semistrong-Weak-current.tex's
% own "Random walk" formula box, switched from the green
% \gothamfloatingexampleblock to this blue one the same day).
\newenvironment{gothamfloatinginfoblock}[2]{%
	\begin{minipage}{#1}
	\begin{infoblock}{#2}
}{%
	\end{infoblock}
	\end{minipage}
}

% Every gothamfloatinginfoblock anchored top-right against the institute
% logo (the deck's standard
% "(current page.north east)+(-\gothamFrametitleRightpading,-\gothamContentTopAlign)"
% corner pattern) lands its own visible right edge short of the logo's,
% by a fixed amount independent of the block's own chosen width (#1) --
% \infoblock's underlying template insets its colored fill by
% \gothamBlockSideMargin from *both* sides of whatever width it's given
% (that margin exists to match a first-level bullet's own left-hand
% position, see its own comment above; the fact that it also narrows the
% *right* side is just a side effect of reusing one symmetric
% hspace+minipage trick for both, not a deliberate right-margin choice),
% partly offset by beamer's own beamercolorbox "colsep*" background bleed
% (colsep* pads the colored fill outward past its own nominal box edge
% without also insetting the text, unlike unstarred "colsep" -- see
% beamerbasecolor.sty's own \beamercolorbox definition). Root-caused and
% measured directly (pymupdf against a compiled probe deck, not derived
% purely on paper) 2026-07-21, after
% slides/j8y-Asset-Allocation-Example-Weights-current.tex's own one-off
% "+2.2pt" nudge (measured that same day against just the institute
% logo's yellow chip) turned out to disagree with
% slides/fr3-Forms-of-Capital-Market-Efficiency-Strong-Semistrong-Weak-current.tex's
% own older "+3.165pt" nudge (2026-07-06, measured against the logo's
% *full* gray-bordered footprint, shadow/card included) -- both slides'
% own nudges were each individually correct for their own chosen
% reference edge, they just didn't agree on what "the logo's own edge"
% means. Direct instruction, 2026-07-21: the *full* gray-bordered
% footprint is the deck's intended reference, so 3.165pt is the value
% kept here; j8y's own newer 2.2pt nudge was the one revised to match,
% not this one.
%
% Use "-\gothamFloatingInfoblockRightAlign" (not the bare
% "-\gothamFrametitleRightpading") for a gothamfloatinginfoblock's own
% anchor node whenever its right edge is meant to align with the logo --
% every current use of this pattern across the deck has been updated to
% match; new slides using this same corner pattern should use this
% length from the start rather than re-measuring their own one-off nudge.
\newlength{\gothamFloatingInfoblockRightAlign}
\setlength{\gothamFloatingInfoblockRightAlign}{\dimexpr\gothamFrametitleRightpading-3.165pt\relax}

%%%%%%%%%%%%%%%%%%%%%%%%%%%%%%%%%%%%%%%%%%%%%%%%%%%%%%%%%%%%%%%%%%%%%%%%%%%%%
% Special frame types (converted from ExampleSlides.pptx; see README.md)
%
% All five below deliberately reuse as much of standardframe -- and of each
% other -- as possible, rather than reimplementing frame chrome from
% scratch: sectionframe, figureframe, and highlightframe are literally thin
% wrappers around standardframe itself (same title bar, logo, footer, bullet
% styling; only what's documented at each one differs). titleframe and
% quoteframe can't be standardframe wrappers -- both are title-less, and
% quoteframe replaces the whole body layout -- but still reuse its palette,
% fonts, and margins (\gothamFrametitleLeftpading, the institute logo
% drawing command, etc.) rather than inventing new ones.
%%%%%%%%%%%%%%%%%%%%%%%%%%%%%%%%%%%%%%%%%%%%%%%%%%%%%%%%%%%%%%%%%%%%%%%%%%%%%

% titleframe: the deck's own title slide (ExampleSlides.pptx slide 1) --
% typically used exactly once, at the very start of a course's first deck.
% One continuous yellow box -- like standardframe's own title bar -- flush
% with the page's top and left edges, stopping short of the institute logo
% on the right (the same margin already used around the logo elsewhere:
% \gothamLogoWidth + 2*\gothamFrametitleRightpading), holding the title and,
% below it, a second line (author/term/whatever) in the same black text as
% standardframe's own frametitle. Usage:
%
%   \titleframe{Financial Markets}{Stefan Palan\\Summer term 2026}
%
% A *command*, not an environment, taking the second (lower) line's content
% as its second argument -- plain text, use \\ for line breaks same as any
% normal LaTeX paragraph, not an itemize, since the source deck's own
% subtitle block is just stacked lines. This has to be a command rather
% than an environment with the body as its content: beamer's frame body is
% captured by a delimited-argument macro (\beamer@collect@@body, looked up
% by triggering the error below on purpose) hunting for a literal
% \end{frame}, and a beamercolorbox -- needed by an earlier version of this
% to wrap the subtitle text in its own colored band -- opened in a *custom
% environment's* begin code and closed in its end code confuses that hunt
% exactly like the \columns example above (same root cause, different
% environment): with
% \newenvironment{titleframe}{...\begin{beamercolorbox}{...}}{\end{beamercolorbox}...},
% compiling ends in "! File ended while scanning use of \beamer@collect@@body."
% every time, however short the body. Taking the content as a normal brace
% argument instead sidesteps the whole issue. No footer/frame number
% (matches beamer's own convention for title slides, e.g. the old
% \begin{frame}[plain,noframenumbering]\titlepage\end{frame} this replaces --
% see README.md) and no chapter-logo support: both would be meaningless on a
% slide that appears once, before any chapter exists.
%
% Unlike the frametitle hook further up, neither the logo nor the yellow box
% here need the white-patch-and-redraw trick: the box's own width already
% stops short of the logo (rather than being painted full-width and then
% partly painted back over), so there's no yellow left to paint over in the
% first place. Both still need \tikz's "remember picture" on their
% "current page"-anchored nodes, though, and so still need the usual second
% compile pass (see "Title bar and logo" in README.md) -- checked
% empirically: without it, "current page" doesn't resolve to the actual page
% box on the first pass, and the logo/yellow box end up positioned nowhere
% near the page corners they're meant to be anchored to.
%
% \gothamTitleframeBoxHeight is a fixed target for the box's *bottom* edge
% (matching what was actually asked for: "reaching down to where the lower
% grey bar currently ends" in the two-band design this replaces), not the
% blank space above the text -- that space is derived from it instead (box
% height minus the title+subtitle content's own measured height), so the
% bottom edge stays fixed even if a longer title ever wraps onto a second
% line; the blank space above would just shrink to compensate. Measured
% against the old two-band design's own rendered PDF coordinates (before
% this change): its second band's bottom edge sat 187.147pt below the page's
% top edge.
\newlength{\gothamTitleframeBoxHeight}
\setlength{\gothamTitleframeBoxHeight}{187.147pt}
\newlength{\gothamTitleframeBoxWidth}
\setlength{\gothamTitleframeBoxWidth}{\dimexpr\paperwidth-\gothamLogoWidth-2\gothamFrametitleRightpading\relax}
% Text inset from the box's own edges: \gothamFrametitleLeftpading, the same
% length standardframe's own title text uses, plus \gothamTitleframeTextSep
% (2.5ex, matching this command's own former per-band beamercolorbox "sep").
\newlength{\gothamTitleframeTextSep}
\setlength{\gothamTitleframeTextSep}{2.5ex}
\newlength{\gothamTitleframeBandGap}
\setlength{\gothamTitleframeBandGap}{0.8ex}
% Because the leading blank space above the text (further down) is itself
% computed as "box height minus the content's own natural height", the
% content's *bottom* padding is the only thing that determines how far the
% last line sits above the box's bottom edge -- the leading space silently
% absorbs whatever that padding is, keeping the bottom edge fixed either
% way. (The relationship is a direct 1:1 offset, checked empirically by
% comparing two rendered values, but with a nonzero intercept from beamer's
% own default top-of-frame vskip -- there's no clean closed-form for the
% right value, hence the empirical calibration here rather than a formula.)
% Calibrated (against the rendered PDF's own text coordinates, not just
% visually) to about 1.5 line-heights of the subtitle text's own \Large
% size (~20.6pt/line here) above the box's bottom edge; re-measure and
% adjust if the subtitle's font size, or beamer's own default top-of-frame
% vskip, ever change.
\newlength{\gothamTitleframeBottomPad}
\setlength{\gothamTitleframeBottomPad}{41pt}
\newbox\gothamTitleframeBox
\makeatletter
\newcommand{\titleframe}[2]{%
	\begin{frame}[plain,noframenumbering,t]%
	% A bare tikzpicture behaves like an hbox (same reasoning as the footer
	% box's own \vbox-to-0pt wrapper further up): sitting directly in the
	% frame's vertical list, it silently opens an implicit paragraph,
	% switching into horizontal mode. Everything placed after it -- the
	% \vspace* further down included -- would then run in horizontal mode
	% instead of vertical mode, which is what a first version of this command
	% got wrong: with two bare tikzpictures here, \vspace*'s glue ended up as
	% *horizontal* space on an implicit line together with them (visible as
	% the title/subtitle shifting sideways, not down, and the vspace amount
	% appearing to have no effect at all -- it was moving the wrong axis).
	% \vbox to 0pt{\vss ...} keeps each overlay at zero space in the normal
	% flow without ever leaving vertical mode.
	%
	% \gothamArtifactBegin/\gothamArtifactEnd (defined with the footer, "PDF
	% accessibility" comment): this box is purely decorative background, same
	% kind of element as the footer's own yellow box -- marked the same way,
	% for a screen reader to skip. titleframe has no footer at all to worry
	% about separately (it's a "plain" frame, no footline), so this box is
	% the only thing here that needs it; the title/subtitle text right below
	% is real content and stays unmarked.
	\gothamArtifactBegin{/Type /Background}%
	\vbox to 0pt{\vss
	\begin{tikzpicture}[overlay, remember picture]
		% \gothamTitleBoxColor, not a literal color token -- same setting the
		% standard frametitle bar (colorA, above) uses, so the title slide's
		% own box matches every other title bar in the deck (2026-07-16,
		% direct instruction).
		% outer sep=0pt here and on the logo below: without it, tikz's
		% default outer sep insets the fill/logo 0.2pt from the anchored
		% page corner -- a hairline white line above/left of this box, and
		% its right edge/the logo 0.2-0.4pt off every other frame's. See
		% the frametitle hook's own comment for the full story.
		\node[fill=\gothamTitleBoxColor, anchor=north west, inner sep=0pt, outer sep=0pt,
		      minimum width=\gothamTitleframeBoxWidth,
		      minimum height=\gothamTitleframeBoxHeight]
			at (current page.north west) {};
	\end{tikzpicture}%
	}%
	\gothamArtifactEnd
	\vbox to 0pt{\vss
	\begin{tikzpicture}[overlay, remember picture]
		% inner ysep derived, not \gothamFrametitleRightpading -- see the
		% frametitle hook's fixed-position logo redraw for why (matches the
		% inline logo's vertically-centered position on every other frame).
		\node[anchor=north east, inner xsep=\gothamFrametitleRightpading,
		      inner ysep=\dimexpr\gothamTitleBarBottomAlign-\gothamLogoHeight\relax,
		      outer sep=0pt]
			at (current page.north east) {\gothamInstituteLogoSquare};
	\end{tikzpicture}%
	}%
	\setbox\gothamTitleframeBox=\vbox{%
		\hsize=\gothamTitleframeBoxWidth
		\parindent=0pt
		\leftskip=\gothamFrametitleLeftpading
		\advance\leftskip by \gothamTitleframeTextSep
		\rightskip=\gothamFrametitleLeftpading
		\advance\rightskip by \gothamTitleframeTextSep
		% \gothamTitleTextColor, not a literal color token -- follows the
		% box color change right above (2026-07-16, direct instruction):
		% leaving this hardcoded dark would be unreadable once the box
		% itself switches to a dark blue for MCI/Innsbruck.
		\color{\gothamTitleTextColor}%
		% Beamer redefines \\ inside every frame to peek for an overlay spec
		% (\\<2->, for line-by-line reveals in itemize/enumerate); that
		% redefinition assumes it's running in beamer's own frame-body
		% machinery and chokes ("Undefined control sequence", silently
		% swallowing the following word) when \\ is used for a plain line
		% break in a raw \vbox like this one. Restoring plain LaTeX's own
		% \\ (\@normalcr) for this group only fixes it without affecting \\
		% anywhere else in the deck. (The @ here needs catcode 11, which is
		% why this whole command definition -- not just this line -- is
		% wrapped in \makeatletter/\makeatother: macro bodies are tokenized
		% at *definition* time, so a \makeatletter placed inside the body
		% itself would take effect too late, after "\@normalcr" had already
		% been mis-tokenized as "\@" followed by four ordinary-text letters.)
		\let\\\@normalcr
		\vskip\gothamTitleframeTextSep
		\noindent{\Huge\bfseries #1\par}%
		\vskip\gothamTitleframeTextSep
		\vskip\gothamTitleframeBandGap
		\vskip\gothamTitleframeTextSep
		\noindent{\Large #2\par}%
		\vskip\gothamTitleframeBottomPad
	}%
	\vspace*{\dimexpr\gothamTitleframeBoxHeight-\ht\gothamTitleframeBox-\dp\gothamTitleframeBox\relax}%
	\box\gothamTitleframeBox
	% This command's own \color{\gothamTitleTextColor} above (set inside a
	% stored box register) can leak into the *next* frame's body text --
	% see standardframe's own opening code ("Forces this frame's own body")
	% for the real fix and the full story (a trailing reset attempted
	% here first turned out to be a silent no-op, since nothing draws after
	% it within this same frame to actually force the color operator to be
	% emitted).
	\end{frame}%
}
\makeatother

% fullbleedframe: a slide that is *only* a full-screen image (FM2_2026s_ps_deck.pptx
% slide 13, a recruitment flyer pasted in as one embedded picture on an
% otherwise blank layout -- no title, no PPTX text at all) -- no title bar,
% no footer, no logo, matching that source's own blank layout intent, unlike
% figureframe (which is standardframe underneath and always keeps the usual
% title/footer chrome). The image is scaled to fit the slide while keeping
% its own aspect ratio (\includegraphics's keepaspectratio, capped by
% \paperwidth *and* \paperheight together), centered; whatever's left over
% on the binding axis (e.g. a widescreen slide holding a taller/narrower
% image) is filled by a solid black backdrop rather than left showing the
% frame's usual white canvas -- the point of a full-*bleed* image is that
% nothing reads as "empty margin" even when the aspect ratios don't match.
%
%   \fullbleedframe{media/photo-datacamp-recruitment-flyer}
%
% \gothamFullBleedBlack is a fixed HTML black, deliberately *not* a
% \colorlet of c_primary_dark: this frame type's backdrop should stay
% neutral true black regardless of whether c_primary_dark itself is ever
% restyled to something else (e.g. a softer near-black, the way
% c_custom_gray already is elsewhere in this file) for branding
% reasons -- those are two different jobs (brand color vs. letterbox
% backdrop) that only happen to want the same value today.
\definecolor{gothamFullBleedBlack}{HTML}{000000}
\newcommand{\fullbleedframe}[1]{%
	\begin{frame}[plain,noframenumbering,t]%
	% \gothamArtifactBegin/\gothamArtifactEnd (defined with the footer, "PDF
	% accessibility" comment): the black backdrop is purely decorative
	% letterboxing, marked the same way as titleframe's/quoteframe's own
	% background boxes so a screen reader skips it. The image itself is left
	% unmarked -- same as everywhere else in this file, since this project
	% has no alt-text infrastructure yet for imported images generally (a
	% separate, not-yet-made decision -- see README.md's accessibility
	% section for the same open question about the institute logo).
	\gothamArtifactBegin{/Type /Background}%
	\vbox to 0pt{\vss
	\begin{tikzpicture}[overlay, remember picture]
		% outer sep=0pt: see the frametitle hook's comment -- without it the
		% backdrop sits 0.2pt in from the anchored page corner.
		\node[fill=gothamFullBleedBlack, anchor=south west, inner sep=0pt, outer sep=0pt,
		      minimum width=\paperwidth, minimum height=\paperheight]
			at (current page.south west) {};
	\end{tikzpicture}%
	}%
	\gothamArtifactEnd
	\vbox to 0pt{\vss
	\begin{tikzpicture}[overlay, remember picture]
		\node[anchor=center] at (current page.center)
			{\includegraphics[width=\paperwidth, height=\paperheight, keepaspectratio]{#1}};
	\end{tikzpicture}%
	}%
	\end{frame}%
}

% sideimageframe: a standardframe that always carries one auto-placed
% "side filler image" (the user's own term, see
% slides/dfc-Angel-Investors-current.tex's own comment for the technique
% this formalizes into a reusable frame type) -- 2026-07-06, direct
% instruction. The image is positioned flush to three edges at once: its
% top touches the title bar's own bottom edge, its right touches the
% page's own right edge, and its bottom touches the footer box's own top
% edge -- so it's scaled by *height* to exactly
% \paperheight-\gothamTitleBarBottomAlign-\gothamFooterBoxHeight (the
% measured gap between those two), with `keepaspectratio` so the width
% simply follows from the image's own aspect ratio rather than being
% forced/distorted to some second, independent target -- anchoring the
% node at the page's own north-east corner then makes the top+right
% touches automatic, and the height exactly filling that gap makes the
% bottom touch automatic too, with no separate width constraint needed
% (only a *square* source image would ever visibly reach some further
% left boundary; this only guarantees top/right/bottom, matching what was
% actually asked for).
%
%   \begin{sideimageframe}{Title}{media/photo-example}
%   	\begin{itemize}
%   		\item ...
%   	\end{itemize}
%   	\source{...}
%   \end{sideimageframe}
%
% \sidefootnotes is switched on automatically (not left to the caller) --
% the whole point of this frame type is that its footnotes always end up
% rotated along the left edge of the image rather than colliding with it
% the way a full-width bottom footnote would (see
% slides/u29-Implications-of-Capital-Market-Efficiency-Prices-current.tex's
% own comment for exactly that collision, on a plain standardframe). Unlike
% every other \sidefootnotes use elsewhere in this deck, though, the strip
% can't sit at its usual fixed inset from the page's own right edge here --
% that would run it directly underneath the image, the same problem u29
% had. Fixed by *computing* the image's own rendered width (via
% \settowidth against the exact same \includegraphics call used to draw
% it) and using that to override \gothamSideFootnoteRightInset for just
% this one frame -- moving the whole rotated strip to sit immediately left
% of the image instead of at the page's raw edge, entirely automatically,
% with no manual x-coordinate needed from the caller (the direct
% instruction offered a manual-x-coordinate fallback "if it is possible
% [to do automatically], that would be great" -- turned out to be
% possible). \gothamSideImageFrameHeight/\gothamSideImageFrameWidth are
% declared once, globally, and merely *set* (not re-declared) inside the
% environment body, so the frame type can be used more than once per deck;
% the \setlength override of \gothamSideFootnoteRightInset is a plain,
% non-\global assignment, which -- like every other frame-local override in
% this file (e.g. slides/u29-...'s own \setbeamercolor{footnote}{fg=white})
% -- resets on its own once \begin{frame}...\end{frame}'s own TeX group
% closes, with no explicit reset needed.
%
% The extra gap added on top of the image's own measured width (2026-07-06,
% direct instruction: "replicate ... the precise margin we have on regular
% slides between footnote and footer box") was originally an unmeasured
% 4pt guess -- replaced with 3.1pt, back-solved from the deck's own real
% reference gap: measured directly on a plain bottom-footnote slide
% (slides/xhe-Active-vs-Passive-Management-2026-First-current.tex's own
% \source{} line) that the real gap between a footnote's own text and the
% footer box's top edge is 2.93pt, while the original 4pt guess here was
% actually landing the rotated text 3.8pt from the image -- 4pt - (3.8-2.93)
% = 3.13pt, rounded to 3.1pt.
\newlength{\gothamSideImageFrameHeight}
\newlength{\gothamSideImageFrameWidth}
\newenvironment{sideimageframe}[3][t]{%
	\setlength{\gothamSideImageFrameHeight}{\dimexpr\paperheight-\gothamTitleBarBottomAlign-\gothamFooterBoxHeight\relax}%
	\settowidth{\gothamSideImageFrameWidth}{\includegraphics[height=\gothamSideImageFrameHeight, keepaspectratio]{#3}}%
	\sidefootnotes
	\setlength{\gothamSideFootnoteRightInset}{\dimexpr\gothamSideImageFrameWidth+3.1pt\relax}%
	\begin{standardframe}[#1]{#2}%
	\vbox to 0pt{\vss
	\begin{tikzpicture}[overlay, remember picture]
		% outer sep=0pt: see the frametitle hook's comment -- without it the
		% image sits 0.2pt in from the page's right edge (and 0.2pt low).
		\node[anchor=north east, inner sep=0pt, outer sep=0pt] at ($(current page.north east)+(0,-\gothamTitleBarBottomAlign)$)
			{\includegraphics[height=\gothamSideImageFrameHeight, keepaspectratio]{#3}};
	\end{tikzpicture}%
	}%
}{%
	\end{standardframe}%
}

% sectionframe: a section-opening overview (ExampleSlides.pptx slide 3) --
% same layout as standardframe (footer, logo, bullet styling all completely
% unchanged), except the title bar itself is much taller (about a third of
% the page) and set in a bigger font, and the body text is a step bigger
% too, so it reads as "somewhat greater focus" and is unmistakable at a
% glance when flipping through slides -- while still keeping the deck's
% usual yellow-background/black-text look, unlike an earlier version of this
% that inverted to dark-background/yellow-text instead.
%
% The taller bar is NOT a separately-drawn box the way titleframe's is (see
% "Title slide" above): the gotham theme's own frametitle box already sizes
% itself from \gothamFrametitleToppading/\gothamFrametitleBottompading plus
% the title font's own strut height (beamerouterthemegotham.sty, the
% "frametitle" template), so overriding those two lengths locally -- the
% same trick already used below for the font/color overrides -- grows the
% *theme's own* box instead of fighting it with a second one. This also
% means the white-patch-behind-the-logo/chapter-logo machinery (tied to the
% frametitle template via \addtobeamertemplate) keeps working unmodified,
% since it measures the box's actual rendered bottom edge rather than
% assuming any particular height.
%
% The title text is bottom-anchored within the bar -- sitting a fixed "one
% line" above the bar's bottom edge, growing *upward* for a title long
% enough to wrap onto a second line -- rather than top-anchored the way
% gotham's own frametitle template flows by default (which would instead
% leave a large, growing gap *below* a wrapped title, between it and the
% bar's bottom edge). gotham's template has no built-in way to do this: its
% only two knobs are the toppading/bottompading struts, both anchored to
% *one* end of a single top-down paragraph flow. So instead of leaving
% toppading fixed (as an earlier version of this did) the whole title text
% is measured into \gothamSectionframeTitleBox first -- built with the same
% font and (approximately) the same usable width gotham's own template
% gives it, so its \ht+\dp closely track the title's real rendered size,
% including wrapping onto extra lines for a long title -- and toppading is
% then computed as "whatever's left" after that measured height and the
% fixed bottom gap are subtracted from the bar's total target height. A
% longer/wrapped title therefore automatically shrinks toppading (pushing
% the whole title block up, growing upward) rather than changing the gap
% below it. \gothamSectionframeTitleWidth is a slight underestimate of
% gotham's actual per-line width (which is only narrowed by the logo on
% whichever line ends up next to it, typically just the last one) --
% applying that narrower width to *every* line is a deliberate, harmless
% simplification rather than replicating gotham's per-line logic exactly.
%
% The institute logo moves down along with the title now (it's vertically
% centered against the title's own baseline, gotham's instituteLogo
% template), rather than staying pinned where it sits on every other frame
% type, as an earlier version of this kept it -- an intentional change from
% that version: with the title itself now sitting near the bar's bottom
% edge, keeping the logo up at the bar's top would visually separate it
% from the title it's meant to sit beside.
%
% \begingroup/\endgroup (rather than the \global-set-and-reset dance
% \chapterlogo{}/\sidefootnotes need) is enough here: this only ever touches
% \setbeamercolor/\setbeamerfont/\setlength/\setbox, all of which assign
% locally under the hood (checked against beamerbasecolor.sty/
% beamerbasefont.sty directly for the first two -- their key handlers are
% plain \def, not \gdef -- and \setlength/\setbox are plain TeX assignments,
% local by default), so a plain surrounding TeX group reverts them
% correctly on its own.
%
% \gothamSectionframeContentDrop adds a little extra room below the bar
% before the body content starts, on top of standardframe's own
% \gothamContentTopShift pull-up (applied inside \begin{standardframe}
% already, before this code regains control). All lengths here calibrated
% against the rendered PDF's own coordinates; re-measure and adjust if the
% title/body font sizes change.
\newlength{\gothamSectionframeContentDrop}
\setlength{\gothamSectionframeContentDrop}{10pt}
\newlength{\gothamSectionframeBoxHeight}
\setlength{\gothamSectionframeBoxHeight}{85.04pt}
% \gothamSectionframeBottomGap is the *visual* target -- one huge-font line
% height (baselineskip, measured directly from a deliberately 2-line test
% title: the gap between the two lines' own rendered top edges). It is NOT
% used directly as \gothamFrametitleBottompading: gotham's bottom strut
% (\gotham@frametitlestrut@end) consistently produces a visual gap about
% 5.58pt *smaller* than the padding value fed into it (checked empirically,
% same offset for both a one-line and a two-line test title, so it's a
% fixed per-box offset rather than something that scales with content) --
% \gothamSectionframeBottomPadCorrection adds that back.
\newlength{\gothamSectionframeBottomGap}
\setlength{\gothamSectionframeBottomGap}{14.32pt}
\newlength{\gothamSectionframeBottomPadCorrection}
\setlength{\gothamSectionframeBottomPadCorrection}{5.58pt}
% Similarly, the box's total rendered height comes out consistently ~11.4pt
% *less* than toppading+measuredTitleHeight+bottompading would suggest --
% the measured title box (built as a plain \vbox{...\par}) doesn't reproduce
% gotham's own leading strut (\gotham@frametitlestrut@start, a zero-width
% rule with an explicit height rather than natural interline glue) exactly,
% and this correction closes that gap. Checked the same way as the bottom
% one above: consistent across both the one-line and two-line test titles.
\newlength{\gothamSectionframeHeightCorrection}
\setlength{\gothamSectionframeHeightCorrection}{12.1pt}
\newlength{\gothamSectionframeTitleWidth}
\setlength{\gothamSectionframeTitleWidth}{\dimexpr\paperwidth-\gothamFrametitleLeftpading-\gothamLogoWidth-2\gothamFrametitleRightpading\relax}
\newbox\gothamSectionframeTitleBox
\newenvironment{sectionframe}[2][t]{%
	\begingroup
	% Disable hyphenation for the title -- both here, in the measurement
	% box below, and (since this is in effect for the rest of the group,
	% which \begin{standardframe} -- triggering the real \insertframetitle
	% -- runs inside) in the actual rendered title too, so a line break
	% never lands mid-word. Setting the penalties alone is the standard,
	% sufficient way to do this (TeX treats a penalty of 10000 at a
	% breakpoint as forbidden, not just heavily discouraged, so it's never
	% chosen even when the alternative is an overfull line); the hyphenmin
	% lengths are redundant insurance, not strictly needed on top of that.
	\hyphenpenalty=10000
	\exhyphenpenalty=10000
	\lefthyphenmin=100
	\righthyphenmin=100
	\setbeamerfont{frametitle}{size=\huge}
	\setbox\gothamSectionframeTitleBox=\vbox{%
		\hsize=\gothamSectionframeTitleWidth
		\usebeamerfont{frametitle}\usebeamercolor[fg]{frametitle}%
		\noindent #2\par
	}%
	\setlength{\gothamFrametitleBottompading}{\dimexpr\gothamSectionframeBottomGap+\gothamSectionframeBottomPadCorrection\relax}
	\setlength{\gothamFrametitleToppading}{\dimexpr\gothamSectionframeBoxHeight-\gothamFrametitleBottompading-\ht\gothamSectionframeTitleBox-\dp\gothamSectionframeTitleBox+\gothamSectionframeHeightCorrection\relax}
	\global\def\gothamFrametitleFullHeightYellow{1}%
	\global\def\gothamFrametitleLogoFixed{1}%
	\setbeamertemplate{frametitle}[gotham_no_inline_logo]%
	\begin{standardframe}[#1]{#2}%
	\vspace*{\gothamSectionframeContentDrop}%
	\large
}{%
	\end{standardframe}%
	\endgroup
}

% quoteframe: a quote slide (ExampleSlides.pptx slide 4) -- institute logo
% as usual, and the quote itself set large/italic/black on a yellow
% full-bleed block, same style as titleframe/sectionframe's own title boxes
% (flush top-left, stopping short of the logo on the right, same margin as
% elsewhere) rather than the dark block an earlier version of this used. The
% usual footer still shows (unlike titleframe, which drops it -- a title
% slide is a one-off outside any particular course content, while a quote
% is still one slide *among* a deck's others, worth being able to reference
% by its footer slide ID like any other). Opening/closing curly quotes are
% added automatically (write the quote text with none of its own); the
% second argument is the person it's from, set smaller/plain, below the
% yellow box and right-aligned to its own right edge (not the page's usual
% text margin) -- pass an empty second argument for an unattributed quote.
%
%   \quoteframe{One of the funny things about the stock market is that
%   	every time one person buys, another sells, and both think they
%   	are astute.}{William Feather, 1889-1981}
%
% \gothamQuoteBoxHeight is a plain percentage of the page, unlike
% titleframe's/sectionframe's own (calibrated-to-land-on-an-exact-fraction)
% box heights, since nothing here needs to line up with another element's
% measured position the way those do. \gothamTitleframeBoxWidth (defined
% above) is reused as-is for the box's width -- the same
% stop-short-of-the-logo formula, so no reason to duplicate it under a new
% name.
%
% Every piece here -- box, logo, quote text, attribution -- is drawn as a
% zero-space tikz overlay anchored to absolute page coordinates (the same
% \vbox to 0pt{\vss ...} pattern used elsewhere in this file for exactly
% this reason), rather than any of it flowing through the frame's normal
% vertical list. This sidesteps needing to know/counteract beamer's own
% default top-of-frame spacing (which standardframe explicitly measures and
% cancels via \gothamContentTopShift) -- there's no such frame-relative
% positioning to get right here, since literally everything is placed by
% explicit coordinate instead. Like the other page-corner overlays in this
% file, "current page" doesn't resolve correctly on a document's very first
% compile pass, so (like those) this needs the usual second pass too.
%
% A *command*, not an environment -- see the identically-reasoned comment on
% titleframe just above for why: the same beamercolorbox-spanning-a-custom-
% environment's-begin/end problem would apply here too, for the same reason.
\newlength{\gothamQuoteBoxHeight}
\setlength{\gothamQuoteBoxHeight}{0.8\paperheight}
\newlength{\gothamQuoteAttributionGap}
\setlength{\gothamQuoteAttributionGap}{1.8ex}
% \@empty (the empty-attribution check below) needs \makeatletter.
% \quotesource{text} -- opt-in \source{} for \quoteframe (call *before*
% \quoteframe{}{}, same set-then-consume pattern as \chapterlogo{}
% above), needed because \quoteframe is a *command* that auto-closes its
% own \end{frame} -- a bare \source{...} typed after \quoteframe{}{} in
% the calling slide file (as if it were still inside the frame, the way
% \source{} reads everywhere else in this deck) is actually issued
% *outside* any frame, and \source{}'s own \footnote{} then has no page to
% attach to yet. Confirmed directly (2026-07-12, direct instruction:
% traced a spurious extra blank leading overlay step on
% slides/c4c-Macbeth-Spot-Removers-Proposed-Structure-current.tex -- the
% very next multi-overlay frame after
% slides/w8w-Pie-Model-of-Capital-Structure-current.tex's own bare
% \source{} -- down to that orphaned footnote specifically, via a reduced
% two-slide test deck): the dangling footnote insert doesn't render
% anywhere visible on w8w itself, but silently corrupts beamer's own
% overlay-counting pre-pass on whichever frame happens to come right
% after -- reproduced with w8w+c4c alone (bug present) vs. c4c alone (bug
% absent), and confirmed it's specifically about the footnote being
% orphaned, not about w8w's content otherwise, by removing just its
% trailing \source{} line (bug also disappears). A plain \clearpage
% between the two frames also "fixes" c4c but only by giving the orphaned
% footnote a page of its own to attach to -- confirmed directly this
% splits w8w itself into two pages (real content, then a second page with
% nothing but the footnote/footer), not an acceptable trade. The only
% correct fix is keeping the footnote inside the frame that's actually
% citing something, i.e. before \quoteframe's own \end{frame} -- exactly
% what \quotesource{} below does.
\newcommand{\gothamQuoteSource}{}
\newcommand{\quotesource}[1]{\gdef\gothamQuoteSource{#1}}
\makeatletter
\newcommand{\quoteframe}[2]{%
	\begin{frame}[t]%
	% \gothamArtifactBegin/\gothamArtifactEnd (defined with the footer, "PDF
	% accessibility" comment): same reasoning as titleframe's own background
	% box above -- purely decorative, marked so a screen reader skips it. The
	% quote text and attribution below stay unmarked (real content); the
	% footer (this frame keeps one, unlike titleframe) is already marked via
	% the global footline hook.
	\gothamArtifactBegin{/Type /Background}%
	\vbox to 0pt{\vss
	\begin{tikzpicture}[overlay, remember picture]
		% \gothamTitleBoxColor, not a literal color token -- same setting
		% titleframe's own box and the standard frametitle bar (colorA) use
		% (2026-07-16, direct instruction).
		% outer sep=0pt here and on the logo/attribution below: same hairline-
		% inset-from-the-anchored-corner issue as titleframe's own box -- see
		% the frametitle hook's comment for the full story.
		\node[fill=\gothamTitleBoxColor, anchor=north west, inner sep=0pt, outer sep=0pt,
		      minimum width=\gothamTitleframeBoxWidth,
		      minimum height=\gothamQuoteBoxHeight]
			at (current page.north west) {};
	\end{tikzpicture}%
	}%
	\gothamArtifactEnd
	\vbox to 0pt{\vss
	\begin{tikzpicture}[overlay, remember picture]
		% inner ysep derived, not \gothamFrametitleRightpading -- see the
		% frametitle hook's fixed-position logo redraw for why (matches the
		% inline logo's vertically-centered position on every other frame).
		\node[anchor=north east, inner xsep=\gothamFrametitleRightpading,
		      inner ysep=\dimexpr\gothamTitleBarBottomAlign-\gothamLogoHeight\relax,
		      outer sep=0pt]
			at (current page.north east) {\gothamInstituteLogoSquare};
	\end{tikzpicture}%
	}%
	\vbox to 0pt{\vss
	\begin{tikzpicture}[overlay, remember picture]
		% \gothamTitleTextColor, not a literal color token -- follows the box
		% color change right above (2026-07-16, direct instruction). Set
		% *inside* the node's own {...} content group (not a stored box
		% register, unlike titleframe's own vbox above), so this doesn't hit
		% the same "\color leaks out of a stored \vbox" pitfall documented on
		% \titleframe's own closing code -- confirmed directly by checking
		% the frame right after this one for the same symptom.
		\node[anchor=center, inner sep=0pt,
		      text width=\dimexpr\gothamTitleframeBoxWidth-2\gothamFrametitleLeftpading\relax]
			at ($(current page.north west)+(0.5\gothamTitleframeBoxWidth,-0.5\gothamQuoteBoxHeight)$)
			{\color{\gothamTitleTextColor}\Large\itshape\textquotedblleft #1\textquotedblright};
	\end{tikzpicture}%
	}%
	\def\gothamQuoteAttributionTemp{#2}%
	\ifx\gothamQuoteAttributionTemp\@empty\else
		\vbox to 0pt{\vss
		\begin{tikzpicture}[overlay, remember picture]
			\node[anchor=north east, inner sep=0pt, outer sep=0pt]
				at ($(current page.north west)+(\gothamTitleframeBoxWidth,-\gothamQuoteBoxHeight-\gothamQuoteAttributionGap)$)
				{\normalfont\itshape\color{c_custom_gray} #2};
		\end{tikzpicture}%
		}%
	\fi
	\ifx\gothamQuoteSource\@empty\else
		\source{\gothamQuoteSource}%
		\global\let\gothamQuoteSource\@empty
	\fi
	\end{frame}%
}
\makeatother

% figureframe: a slide that is just a figure (ExampleSlides.pptx slide 5) --
% same standardframe layout (title, footer, \source{} for the usual
% "Source: ..."/"Image source: ..." line), but with the figure positioned
% via \figurecenter (call before the figure, inside the frame body) rather
% than left-aligned, which -- like the source deck's own figure slide -- is
% the default:
%
%   \begin{figureframe}{Correlation between DJIA stock returns over 15 years}
%   	\includegraphics[width=0.5\linewidth]{correlation-matrix}
%   	\source{Source: Own illustration, data from Yahoo! Finance, 24.02.2026.}
%   \end{figureframe}
%
%   \begin{figureframe}{Title}
%   	\figurecenter
%   	\includegraphics[width=0.6\linewidth]{diagram}
%   \end{figureframe}
%
% Vertical placement is standardframe's own [t]/[c] (its existing optional
% argument, passed straight through as figureframe's own -- so
% \begin{figureframe}[c]{Title} vertically centers exactly like
% \begin{standardframe}[c]{Title} already does); \figurecenter only ever
% affects the horizontal axis.
%
% \figurecenter is plain \centering under a more discoverable name, nothing
% more -- *not* an opt-in-before-the-frame flag like \chapterlogo{}/
% \sidefootnotes (those must run before \begin{standardframe}, since they
% configure something the frame's own opening code reads). This is called
% *inside* the body instead, exactly where it takes effect, so it can just
% take effect immediately like any other formatting declaration -- no flag,
% no state to reset afterwards, and (checked by first writing it as a flag
% read in figureframe's own opening code, mirroring \sidefootnotes) no way
% for the frame's opening code to see it anyway, since that code runs
% before the body -- and \figurecenter -- does.
\newcommand{\figurecenter}{\centering}
\newenvironment{figureframe}[2][t]{%
	\begin{standardframe}[#1]{#2}%
}{%
	\end{standardframe}%
}

% highlightframe: a standardframe with two things about it made unmistakably
% different -- the title bar's *and* the footer's background color, both
% together -- so a presenter can never flip to it without immediately
% noticing (ExampleSlides.pptx slide 6 was the inspiration for wanting a
% standardframe variant here, though the source slide itself has no such
% color cue; this is a new addition, not a conversion). Everything else --
% body, logo, bullet styling -- is exactly standardframe's own, via the same
% \begingroup/\endgroup wrapper sectionframe above uses for the title-color
% half of this. \gothamHighlightColor is a \colorlet, not a color baked
% directly into the environment, so retargeting every highlightframe in a
% deck to a different color -- title bar and footer both -- is still the
% same one-line change it was when this only touched the title bar.
%
% The footer box needs the \global\def dance (\gothamFooterBoxColor, defined
% above with the frametitle hook's other opt-in flags, reset in
% standardframe's own closing code) that \setbeamercolor{frametitle} doesn't:
% that assigns through beamer's own color-management system, which handles
% the frame-closing group correctly on its own (checked -- see
% \gothamChapterLogoPath's comment for why plain macros need \global\def but
% beamer's own set-commands don't), but the footer box's fill is read from a
% plain macro inside a tikz key, same category as \gothamChapterLogoPath.
\colorlet{gothamHighlightColor}{c_primary_gray}
\newenvironment{highlightframe}[2][t]{%
	\begingroup
	\setbeamercolor{frametitle}{bg=gothamHighlightColor}
	\global\def\gothamFooterBoxColor{gothamHighlightColor}%
	\begin{standardframe}[#1]{#2}%
}{%
	\end{standardframe}%
	\endgroup
}


) to replace the default marker set for every
% \footnote{}/\source{} from then on. \seq_item:Nn (expl3's indexed lookup;
% expl3 is already loaded by the gotham theme itself, see e.g.
% beamerfontthemegotham.sty) is what removes the old \ifcase's hard 9-symbol
% ceiling -- a custom list can be any length. Indexing wraps around (via
% \int_mod:nn) rather than erroring past the end of the list, so an unusually
% footnote-heavy slide degrades to repeating symbols instead of a hard
% failure; slides/ZZU-Footnote-symbol-options-current.tex is a reference
% sheet of candidate Unicode symbols to build a custom list from.
%
% Project-wide default (chosen deliberately, not the classic */dagger/etc.
% set): a currency-sign sequence, in this specific order. All 14 were
% confirmed against slides/ZZU-Footnote-symbol-options-current.tex to
% actually exist in Fira Sans.
\ExplSyntaxOn
\seq_new:N \g_gotham_footnote_symbols_seq
\seq_gset_from_clist:Nn \g_gotham_footnote_symbols_seq { £,¥,₹,₽,₺,₴,₪,₫,₦,₡,₩,ƒ,¢,€ }
\newcommand{\setfootnotesymbols}[1]{%
	\seq_gset_from_clist:Nn \g_gotham_footnote_symbols_seq { #1 }
}
\newcommand{\gothamFnSymbol}[1]{%
	\seq_item:Nn \g_gotham_footnote_symbols_seq
		{ \int_eval:n { \int_mod:nn { #1 - 1 } { \seq_count:N \g_gotham_footnote_symbols_seq } + 1 } }
}
\ExplSyntaxOff
\renewcommand{\thefootnote}{\gothamFnSymbol{\value{footnote}}}
% A \footnote{} called inside a minipage doesn't use \thefootnote/\c@footnote
% at all -- LaTeX's own minipage machinery silently switches to a separate
% counter/display command, \c@mpfootnote/\thempfootnote, normally displayed
% via \@alph (a, b, c, ...). Found 2026-07-14 converting
% slides/j6h-Foreign-Exchange-Market-current.tex (its first slide combining
% \footnote{} with a minipage, since a narrowed text column was added
% alongside a corner chart): the footnote *text* still correctly landed in
% this deck's usual bottom-of-frame collection (beamer's own footnote
% storage isn't minipage-scoped), but the *marker*, both inline and at the
% bottom, showed as a plain "a"/"b"/"c" instead of this project's own
% currency-symbol sequence. Redefined the same way as \thefootnote so any
% future \footnote{} inside a minipage picks up the right symbol too --
% \c@mpfootnote resets to 0 for every new minipage on its own (unlike
% \c@footnote, which needs standardframe's own explicit per-frame reset
% below), which happens to already match this deck's "restart every slide"
% footnote convention for the common case of one minipage per frame.
\renewcommand{\thempfootnote}{\gothamFnSymbol{\value{mpfootnote}}}
\setbeamerfont{footnote}{size=\fontsize{3.2}{1.9}\selectfont}
% Beamer's own \@makefnmark (used for the marker both inline in the body
% text and again next to the footnote itself) hard-codes the "footnote
% mark" beamerfont onto it -- an *absolute* size, unrelated to whatever
% text surrounds it -- which is why the inline markers looked tinier than
% a normal superscript when that was set to this footnote's own small size:
% a fixed size that looks fine amid 3.2pt footnote text looks minuscule
% amid much bigger body text. Reverting to plain LaTeX's own \@makefnmark
% instead makes it a normal \textsuperscript of whatever font size is
% already active at each call site -- proportionally scaled either way, so
% it looks like a normal superscript inline, and still comes out small next
% to the (small) footnote text since \usebeamerfont{footnote} above is
% already active there by the time it's invoked.
\makeatletter
\renewcommand{\@makefnmark}{\hbox{\@textsuperscript{\normalfont\@thefnmark}}}
\makeatother
% \raggedleft's own \rightskip (0pt) is flush with the frame's normal body
% text margin, not the logo's edge above it; this pushes further right by
% the difference between the two, empirically calibrated (like
% \gothamFooterRightShift above) against the rendered PDF coordinates of
% both -- re-measure and adjust if the text margins or logo inset change.
\newlength{\gothamFootnoteRightShift}
\setlength{\gothamFootnoteRightShift}{4.8pt}

% Side footnotes: \sidefootnotes (called *before* \begin{standardframe}, like
% \chapterlogo{}) moves that one frame's footnotes -- including \source{} lines,
% which are footnotes underneath -- from their usual spot above the footer to a
% vertical strip along the right edge of the page, rotated 90 degrees counter-
% clockwise: the text starts just above the footer box and reads bottom-to-top,
% as if the usual bottom-right footnote block had been turned on its side (with
% one deliberate difference: the lines are flush with the strip's bottom end,
% not the top -- see the \raggedright comment in the template below).
% standardframe resets the switch
% on \end{frame} (same \global mechanics as \chapterlogo{} -- see the comment
% there for why a plain local assignment doesn't survive beamer's frame group),
% so it opts in only the single frame that follows, and any frame type derived
% from standardframe inherits the behaviour automatically.
%
% How it works: the footnote *template* below branches on this switch. In side
% mode it typesets each footnote into \gothamSideFootnotesBox (a plain box
% register accumulated across the frame's footnotes) instead of contributing
% anything to beamer's own bottom-of-frame footnote insert, and the footline
% hook further down draws that box -- rotated -- as a tikz overlay once the
% whole frame body has run. The footline is the right place to draw from: it is
% typeset at page shipout, i.e. after every \footnote{} in the body has already
% executed, whereas e.g. the frametitle would run too early and find the box
% empty. (Beamer also typesets the footline once per frame just to *measure*
% it, at \begin{frame} time -- before any footnotes exist -- which is why the
% drawing hook guards on \ifvoid: the measurement pass must not consume, or
% try to draw, a box that doesn't exist yet.)
%
% Known quirk: a footnote with an overlay spec (\footnote<2->{...}) is
% collected -- and therefore shown -- on *every* overlay slide of its frame in
% side mode. Beamer implements covered footnotes by typesetting them invisibly
% in place, which relocating the text into a box register defeats.
\newif\ifgothamSideFootnotes
\newcommand{\sidefootnotes}{\global\gothamSideFootnotestrue}
\newbox\gothamSideFootnotesBox
% The rotated block's line width, i.e. the vertical run the strip may occupy
% down the right edge of the page: from just below the title bar to just above
% the footer's yellow box. Like \gothamFooterBoxHeight above, calibrated
% against the measured PDF coordinates: the title bar's bottom edge sits
% 219.57pt above the page's bottom edge; minus \gothamSideFootnoteTopSep (3pt),
% \gothamFooterBoxHeight (7.26pt), and ~2pt of clearance above the footer box
% leaves 207pt. Re-measure if the frametitle font or its padding ever changes.
\newlength{\gothamSideFootnoteHsize}
\setlength{\gothamSideFootnoteHsize}{207pt}
% Gap between the title bar's bottom edge and the top end of the strip. The
% text itself sits at the strip's *bottom* end (it is \raggedright, i.e.
% flush with the line start, which the rotation turns into bottom-aligned),
% so this gap only matters for footnotes long enough to fill the whole run.
\newlength{\gothamSideFootnoteTopSep}
\setlength{\gothamSideFootnoteTopSep}{3pt}
% Distance from the page's right edge to the strip's outer edge -- the rotated
% counterpart of the small gap the bottom footnotes keep above the footer box.
\newlength{\gothamSideFootnoteRightInset}
\setlength{\gothamSideFootnoteRightInset}{2pt}
% A real footnote's marker (e.g. "*") sits in an \hbox of its own whose
% *ink* -- unlike \source{}'s marker, which is emptied to nothing and so has
% no ink at all -- extends slightly left of that hbox's own nominal left
% edge (ordinary optical behaviour for a small centred glyph like an
% asterisk, exaggerated further by \@makefnmark's superscript raise). With
% every footnote's paragraph starting flush at the same \noindent left edge,
% that overhang made real footnotes' markers visibly poke out to the left of
% \source{}'s flush-left text above/below them. This shift nudges real
% footnotes' *entire* line (marker, comma, and text together) right by that
% overhang, so the marker's ink -- not the text after it -- lines up with
% \source{}'s left edge; \source{} itself is untouched; empirically
% calibrated against the rendered PDF, like the other shift constants here.
\newlength{\gothamSideFootnoteMarkerShift}
\setlength{\gothamSideFootnoteMarkerShift}{1.3pt}

% A new template variant (rather than reusing beamer's own "default", which
% it has already defined and then refuses to redefine) that we immediately
% activate below.
% Wrapped in \makeatletter (unlike the rest of this section) because the side
% branch below tests \@thefnmark against \@empty to tell a real footnote
% (non-empty marker) apart from a \source{} call (marker forced empty) --
% see \gothamSideFootnoteMarkerShift above.
\makeatletter
\defbeamertemplate{footnote}{gothamfn}{%
	\ifgothamSideFootnotes
		% Side mode (see \sidefootnotes above): append this footnote to the
		% collector box instead of typesetting it here. \insertfootnotetext
		% carries beamer's own \rule\z@\footnotesep lead-in and
		% \@finalstrut\strutbox tail along with it (beamer bakes both into the
		% text it hands the template), so the between-footnote spacing and the
		% strutbox rebuild work exactly as in the bottom branch below -- and,
		% equally, nothing of the footnote (not even the spacing rule) is left
		% behind in beamer's bottom-of-frame insert, which therefore stays
		% visually empty. No \gothamFootnoteRightShift here: that constant
		% compensates the *bottom* block's \raggedleft for the text-margin-vs-
		% logo-edge difference, while the strip's position is set directly by
		% the tikz node in the footline hook.
		\global\setbox\gothamSideFootnotesBox\vbox{%
			\unvbox\gothamSideFootnotesBox
			\hsize\gothamSideFootnoteHsize
			\usebeamerfont{footnote}%
			\setlength{\lineskip}{0.6pt}%
			\settoheight{\dimen0}{Xy}%
			\settodepth{\dimen1}{Xy}%
			\setbox\strutbox\hbox{\vrule height\dimen0 depth\dimen1 width0pt}%
			\parindent 1em\noindent%
			% \raggedright, not the bottom block's \raggedleft: after the
			% rotation, the line start is the *bottom* end of the strip, so
			% this makes the text begin just above the footer box and run
			% upward (rather than hang down from the title bar).
			\raggedright
			% \source{} forces \thefootnote (hence \@thefnmark) empty, so this
			% only fires for a real footnote's marker -- see
			% \gothamSideFootnoteMarkerShift above for why it needs the shift
			% and \source{} doesn't. The "\," below (a separator between the
			% marker and the footnote text) needs the same guard: it used to
			% run unconditionally, so an empty-marker \source{} still got a
			% thin space's worth of leading indent before its own text --
			% invisible under the *bottom* block's \raggedleft (nothing there
			% depends on the left edge), but a real, visible defect here under
			% \raggedright, where the left edge is exactly what's on display:
			% a wrapped \source{}'s first line sat measurably indented from
			% its own second line, which never picks up that one-time leading
			% space (2026-07-12, direct instruction: "the first line seems
			% indented compared to the second, which it shouldn't be").
			\ifx\@thefnmark\@empty
				\insertfootnotetext\par%
			\else
				\hspace*{\gothamSideFootnoteMarkerShift}\insertfootnotemark\,\insertfootnotetext\par%
			\fi
		}%
	\else
	\usebeamerfont{footnote}%
	\setlength{\lineskip}{0.6pt}%
	% beamer's footnote machinery resets to plain \footnotesize (a much
	% bigger font) before this template runs, and switching to our own
	% smaller font via \usebeamerfont doesn't update \strutbox to match --
	% so the invisible strut beamer appends after every footnote
	% (\@finalstrut\strutbox) is still sized for \footnotesize, silently
	% padding out extra depth below the last line. Rebuilding \strutbox
	% from this font's own measured metrics removes that padding, which is
	% most of the gap between the footnotes and the footer box below them.
	\settoheight{\dimen0}{Xy}%
	\settodepth{\dimen1}{Xy}%
	\setbox\strutbox\hbox{\vrule height\dimen0 depth\dimen1 width0pt}%
	\parindent 1em\noindent%
	\raggedleft
	\advance\rightskip by -\gothamFootnoteRightShift%
	% Same \@thefnmark emptiness guard as the side-footnote branch above,
	% for the same reason (an empty marker still leaving behind the "\,"
	% separator's own leading thin space) -- invisible here under
	% \raggedleft, where the left edge is never on display, but the same
	% underlying defect, so fixed the same way for consistency.
	\ifx\@thefnmark\@empty
		\insertfootnotetext\par%
	\else
		\insertfootnotemark\,\insertfootnotetext\par%
	\fi
	\fi
}
\makeatother
\setbeamertemplate{footnote}[gothamfn]

% Draws the side-footnote strip (see \sidefootnotes above). Lives in the
% footline template because that is typeset at page shipout, after the frame
% body -- and with it every \footnote{} call -- has run, so the collector box
% is complete by the time this executes. The \ifvoid guard doubles as the
% protection against beamer's measure-the-footline pass at \begin{frame} time
% (see the \sidefootnotes comment). \box (not \copy) empties the register
% globally as a side effect, which is exactly right: each overlay slide of a
% frame re-executes the body and re-collects its footnotes from scratch, so
% whatever this slide's footline draws must not leak into the next slide.
%
% Node geometry: tikz pins the node's *unrotated* anchor corner at the given
% coordinate and rotates the box counter-clockwise around it, so anchor=south
% east (the end of the collected block's last line) ends up as the strip's
% top-east corner: from there the block extends down the page (its lines run
% bottom-to-top, last footnote nearest the page edge), with the \raggedright
% text sitting at the strip's bottom end, just above the footer box (see the
% comment in the footnote template above). The y-coordinate
% comes from the same remembered gothamTitleBottom node the title-bar patches
% already use (same two-pass .aux mechanics -- see "Title bar and logo").
% The \ifcsname guard skips the drawing while the gothamTitleBottom shape is
% entirely unknown to pgf (i.e. no frametitle has been typeset yet in this
% run *and* nothing was inherited from a previous run's .aux) -- without it,
% referencing the node is a hard error rather than the usual harmless
% first-pass misplacement. In that situation the position would have been
% garbage anyway, and the frame-start clearing in standardframe below mops up
% the box this branch leaves unconsumed.
\makeatletter
\addtobeamertemplate{footline}{%
	\ifvoid\gothamSideFootnotesBox\else
		\ifcsname pgf@sh@ns@gothamTitleBottom\endcsname
		\vbox to 0pt{% zero-space trick, as in the yellow-box footline hook above
		\vss
		\begin{tikzpicture}[remember picture, overlay]
			\node[rotate=90, anchor=south east, inner sep=0pt]
				at ($(gothamTitleBottom -| current page.east)+(-\gothamSideFootnoteRightInset,-\gothamSideFootnoteTopSep)$)
				{\box\gothamSideFootnotesBox};
		\end{tikzpicture}%
		}%
		\fi
	\fi
}{}
% Empties the side-footnote collector (used by standardframe at frame start).
% Defined here, inside the \makeatletter block, because \voidb@x needs @ to be
% a letter at tokenization time -- standardframe's own definition below sits
% outside any \makeatletter region.
\newcommand{\gothamClearSideFootnotesBox}{\global\setbox\gothamSideFootnotesBox=\box\voidb@x}
\makeatother

% \source{text} -- for a "Source: ..."/"Image source: ..."/etc. citation line
% (common at the bottom/edge of a converted PowerPoint slide) that should look
% and sit *exactly* like the footnotes above -- same font, spacing,
% right-alignment, and the same stretch-to-sit-just-above-the-footer
% behaviour beamer already gives real footnotes -- but with no marker, neither
% inline in the body text nor next to the note itself, since a source line
% isn't anchored to one specific word. Implemented as a real \footnote{} (so
% it gets that positioning for free, rather than reimplementing it) with
% \thefootnote emptied locally; since \footnote{} still advances the footnote
% counter even with an empty marker, \addtocounter reverses that so a real
% \footnote{} elsewhere on the same slide doesn't skip a symbol in the
% \fnsymbol sequence. Unlike a real footnote, no label is added automatically
% here -- the label varies ("Source:", "Image source:", ...), so write it as
% part of #1.
\newcommand{\source}[1]{%
	\begingroup
	\renewcommand{\thefootnote}{}%
	\footnote{#1}%
	\endgroup
	\addtocounter{footnote}{-1}%
}

% A plain tabular with yellow \toprule/\midrule/\bottomrule, matching the
% frametitle/footer accent color. Takes the same column-spec argument as
% tabular itself, e.g. \begin{gothamtable}{lccc} ... \end{gothamtable}. Only
% booktabs rules are colored (via \arrayrulecolor, scoped to this table by
% \begingroup/\endgroup) -- if a table ever needs a plain black \hline too,
% that's unaffected since \arrayrulecolor only recolors rules drawn *after*
% it takes effect, and booktabs' rules read it fresh each time.
%
% RULE-BLEED WARNING (2026-07-12, direct instruction: "make sure the table
% lines do not extend beyond the text"), read this before writing a new
% gothamtable: a plain column spec like `{lccc}` above leaves the default
% \tabcolsep (6pt, or whatever a slide-local \setlength last set it to) as
% *implicit* padding on the table's own outermost edges -- invisible to the
% table's own text (which starts/ends flush with each column's actual
% content), but very visible to booktabs' \toprule/\midrule/\bottomrule and
% to colortbl's \rowcolor/\cellcolor, both of which span the *declared*
% column width including that implicit padding. The result, confirmed
% directly and independently on at least three slides so far
% (slides/u5h-Long-term-Bonds-Provisions-current.tex's own v1-1-0/v1-4-0,
% slides/c4c-Macbeth-Spot-Removers-Proposed-Structure-current.tex's own
% v1-2-2, slides/tw5-Exercise-Leverage-and-ROE-Solution-current.tex's own
% v1-4-0): every rule/highlight bleeds \tabcolsep past the visible text on
% both the left and right edges, most noticeable on a left-aligned
% (non-\centering) table where the mismatch is easy to see against the
% title's own left margin.
%
% Not fixed here in the environment definition itself: doing so would
% change every existing gothamtable's own column spacing out from under it
% (this environment has no way to tell "plain default spacing was intended"
% apart from "the bleed was never noticed"). Fix per table instead, the
% same three-part pattern every slide above uses -- worth reaching for by
% default on any *new* gothamtable, not just when the bleed is spotted:
%   {\setlength{\tabcolsep}{0pt}
%   \begin{gothamtable}{@{}l<{\hspace{<gap>pt}}c<{\hspace{<gap>pt}}c@{}}
%   ...
%   \end{gothamtable}}
% -- \tabcolsep zeroed (scoped locally, e.g. by wrapping the whole call in
% a `{...}` group, so it doesn't leak into later tables), explicit `@{}` at
% both outer edges of the column spec (so booktabs/colortbl have no
% implicit padding left to bleed into), and array's `<{\hspace{<gap>pt}}>`
% column modifier appended after every column but the last, to put back an
% explicit, *content-column* gap (matching whatever `2*\tabcolsep` the
% table used to have, if preserving its exact look, e.g. `<{\hspace{12pt}}>`
% for the ordinary 6pt-\tabcolsep default) -- `<{...}>` puts the gap
% *inside* the preceding column's own content, which `\rowcolor` (if used)
% colors along with everything else, unlike a plain `@{...}` separator,
% which `\rowcolor` deliberately leaves untinted (see u5h's own v1-6-0
% comment for that exact failure mode if a highlighted table ever needs
% this too).
%
% RE-CONFIRMED 2026-07-16, direct instruction, after this exact bleed
% reappeared on essentially every gothamtable in a freshly-converted
% chapter (Capital Budgeting Techniques) despite the warning above already
% existing at the time -- the rule above is not optional/only-if-noticed,
% it's the default way to write *any* new gothamtable, full stop. Check
% every new gothamtable against it before considering a slide done, not
% just when a rendered page happens to make the bleed visually obvious.
%
% BLOCK-RULE-COLOR CONVENTION (2026-07-16, direct instruction): a
% gothamtable placed *inside* any block environment (`infoblock`,
% `exampleblock`, `referenceblock`, `definitionblock`, plain `block`, or
% their `gothamfloating*` wrappers) should use black rules, not this
% environment's own default yellow -- yellow-on-the-block's-own-tinted-
% background reads poorly, unlike yellow-on-plain-white outside a block.
% Override locally, *inside* `\begin{gothamtable}{...}`, before `\toprule`
% (not before `\begin{gothamtable}` itself -- the environment's own
% opening code sets yellow *after* its `\begingroup`, so a color set
% before that point is simply overwritten; confirmed directly, first
% attempt this same day):
%   \begin{gothamtable}{...}
%   	\arrayrulecolor{c_primary_dark}
%   	\toprule
%   	...
%   \end{gothamtable}
% Booktabs rules read `\arrayrulecolor` fresh at each rule, and
% `\endgroup` (gothamtable's own closing code) resets it back to yellow
% afterward, so this never leaks into a later table. First applied on
% slides/6ux-Example-51-BMA-Payback-Table-current.tex's own
% `gothamfloatinginfoblock`.
\newenvironment{gothamtable}[1]{%
	\begingroup
	\arrayrulecolor{c_primary_highlight}%
	\begin{tabular}{#1}%
}{%
	\end{tabular}%
	\endgroup
}

% For a side-corner image overlay (the usual \vbox to 0pt{...} tikz overlay
% pattern used throughout this deck's converted slides) whose *top edge*
% should line up with the content text's own top edge -- distance from the
% page's own top edge to the cap-height top of a standardframe's first body
% line, measured directly against the compiled PDF's own text coordinates,
% not derived from \gothamContentTopShift or guessed. Previous per-slide
% image placements each guessed their own downward shift (e.g. a bare
% "-1.3cm") independently -- close, but not exact, landing the image's top
% a couple mm above the text's own top; use this constant instead of a
% fresh guess for any new "text top edge"-aligned image.
%
% 41.58pt (2026-07-22, re-measured): the previous value, 44.0pt, went stale
% the moment \gothamContentTopShift itself was increased (7.33pt -> 11.5964pt,
% "reduce the Top Margin", see that length's own comment above) without a
% matching update here -- that change pulls plain body text up, but this
% length is a separate, hand-measured constant with no formula tying it to
% \gothamContentTopShift, so it kept anchoring images at the *old*, lower
% position. Confirmed directly (pymupdf against this same freshly-compiled
% preview deck, not assumed from the \gothamContentTopShift delta alone):
% three independent plain-first-line standardframes with no per-slide
% top-shift correction (slides/5v4-CML-vs-SML-current.tex,
% slides/u5j-Exercise-CAPM-4-current.tex,
% slides/nt4-Exercise-CAPM-5-current.tex) all land their first line's own
% top edge at y=41.58pt post-change, not 44.0pt. Re-measure and adjust
% again if \gothamContentTopShift ever changes in the future.
\newlength{\gothamContentTopAlign}
\setlength{\gothamContentTopAlign}{41.58pt}

% For a "side filler image" -- a decorative image stuck flush to the page's
% own right edge (zero margin) with its top edge flush against the title
% bar's own bottom edge (zero gap), rather than sitting in the usual
% content-aligned corner inset -- distance from the page's own top edge
% down to the standardframe title bar's bottom edge, measured directly
% against the compiled PDF (a pixel scan for the yellow-to-white
% transition, at this deck's 16cm-wide/9cm-tall page size). First used on
% slides/dfc-Angel-Investors-current.tex.
%
% 25.43pt -- the raw measured value (originally 25.44pt; re-measured
% 2026-07-07 against the current compiled PDF via the title bar's own
% white-patch rectangle, which came back 25.427pt, close enough to the
% original measurement to be the same underlying edge, not a drift).
% Briefly set to 24.94pt instead (2026-07-06-ish), a deliberate 0.5pt
% overlap *into* the title bar, after a report of a visible hairline
% white gap right at this seam in one render -- not reproduced in this
% project's own pdftoppm-based rendering checks at the time (every pixel
% scan at the raw value showed a clean cut, no gap), so assumed to be a
% renderer-specific sub-pixel rounding artifact rather than a real error
% in the constant, and "fixed" by deliberately overlapping instead of
% chasing the specific viewer. That overlap turned out to be the wrong
% call: direct feedback 2026-07-07 reported the *opposite* symptom on
% `sideimageframe` (the image's own top edge visibly overlapping the
% title box) -- exactly what a real 0.5pt overlap looks like once it's
% actually noticed, rather than a rendering artifact. Reverted back to
% the precise raw measurement: a real, small, deliberate overlap chosen
% to paper over an unreproduced gap in an unspecified other renderer
% turns out to be a worse trade than just using the measured value and
% letting a genuine future gap report (if one ever recurs, in a
% specific, identified renderer) be diagnosed on its own terms.
%
% Re-measured 2026-07-16, 25.43pt -> 25.32pt: the frametitle hook's
% outer sep fix (see its own comment) moved the visible yellow bottom
% edge up by the 0.2pt the trim patch's fill had previously been inset
% below its gothamLogoAnchor.south anchor -- the yellow now genuinely
% ends flush with the logo's bottom edge, as always intended. New raw
% measurement against the rebuilt PDF: the trim patch's fill top edge
% (= the visible yellow bottom) sits 25.228 PDF units below the page
% top, i.e. 25.3226pt (earlier measurements in this comment quoted PDF
% coordinate values directly as "pt"; PDF units are big points, 1bp =
% 1.00375pt -- a systematic ~0.1pt understatement at this magnitude,
% harmless then, but corrected now that the constant is being reset
% anyway). Note this value is also what the overlay-logo nodes
% (titleframe/quoteframe/the hook's fixed-position redraw) derive their
% vertical inset from: visible-bar-bottom minus \gothamLogoHeight is
% exactly the inline logo's own vertically-centered top margin.
\newlength{\gothamTitleBarBottomAlign}
\setlength{\gothamTitleBarBottomAlign}{25.32pt}

% Note on the [t] default below: it top-aligns content within the frame
% (beamer's own default is vertical centering, "c"), so slides don't have
% their content floating in the middle. It also sidesteps a real beamer bug:
% \begin{frame}[]{Title} (a literal empty options bracket) fails outright,
% even in plain beamer with no theme at all -- so the optional argument must
% never default to an empty string. The title is always passed as an
% explicit argument (rather than left for beamer to pick up from a trailing
% brace group), since anything placed between \begin{frame} and its title --
% like our \setbeamertemplate calls -- breaks that trailing-brace mechanism.

% Pulls the frame's content up closer to the title bar: measured (via the
% PDF's own coordinates) at about 22pt of gap by default between the title
% bar's bottom edge and the first line of body content: this is a third of
% that, shrinking the gap to about two thirds of its original size.
%
% Called "Top Margin" in conversation from 2026-07-12 onward (direct
% instruction) -- the deck-wide name for the resulting gap between the
% title bar and a standardframe's first line of body content, i.e. the
% vertical position a plain-text frame's first baseline naturally lands at
% (measured directly against a reference slide's own compiled PDF, e.g.
% slides/aii-Leverage-Effects-current.tex's first line). Any first content
% that isn't plain text (a table, a minipage, an itemize with an unusual
% first item, ...) needs its own per-slide \gothamIncreaseContentTopShift/
% \gothamReduceContentTopShift correction to land on this same baseline --
% not portable between slides even for superficially similar content (see
% slides/zgz-EBIT-Definition-current.tex's own v1-2-0 comment for a case
% where an earlier, wrong correction was carried over uncritically).
\newlength{\gothamContentTopShift}
% 11.5964pt (2026-07-22, direct instruction: "reduce the Top Margin"
% generally, tested at 20%/30%/40% smaller against
% slides/8s8-Unit-Investment-Trusts-current.tex before settling on 40% --
% up from the original 7.33pt so the *resulting* gap (this length pulls
% content up by, so a bigger value means a smaller visible margin) comes
% out to 60% of its old size: measured directly on 8s8 at 10.67pt before,
% 6.42pt after (target was 6.40pt, within 0.02pt). Since every
% standardframe applies this same pull-up before any per-slide
% \gothamIncreaseContentTopShift/\gothamReduceContentTopShift correction,
% changing it here shifts the baseline project-wide, not just on 8s8.
\setlength{\gothamContentTopShift}{11.5964pt}

% "Top Margin" (above) is one of four deck-wide margin names introduced in
% conversation from 2026-07-12 onward (direct instruction) for the plain
% page-edge/title-bar/footer gaps every slide is expected to obey by
% default. The other three, for reference (each maps to a length already
% declared elsewhere, not a new one introduced here):
% - "Left Margin": the default gap from the page's own left edge to the
%   title text and body content -- \gothamFrametitleLeftpading
%   (beamerouterthemegotham.sty/beamerthemegotham.sty, 2.2ex, measured
%   directly against the compiled PDF at 12.552pt in this deck's usual
%   fonts -- see e.g. slides/tw5-Exercise-Leverage-and-ROE-Solution-current.tex's
%   own v1-1-0 comment for that same measured value).
% - "Right Margin": the default gap from the page's own right edge to the
%   institute logo (and, by the same convention, to anything else
%   deliberately flush-right, e.g. a corner-anchored image) --
%   \gothamFrametitleRightpading (0.28cm/7.93pt, this file's own
%   \gothamInstituteLogoSquare placement above uses it directly as the
%   logo's own inner sep).
% - "Bottom Margin": the default gap between a frame's own body content
%   and the \source{}/footnote line below it. Unlike the other three, this
%   one has no single settable length behind it -- \source{}'s own
%   vertical position isn't fixed the way the footer box is; it tracks the
%   frame's actual content height once that content gets tall enough (see
%   slides/u5h-Long-term-Bonds-Provisions-current.tex's own v1-3-0 comment,
%   "past a certain table height it starts tracking the table's own growth
%   roughly 1:1"). "Respecting Bottom Margin" in practice means leaving
%   comparable clearance to whatever a reference slide with unremarkable
%   content already shows, the same empirical, measure-against-a-reference-
%   slide approach "Top Margin" itself uses, rather than a fixed constant.

% Compensates a gap that opens up specifically when a standardframe's body
% starts with a \begin{minipage}{...} (used to width-limit text next to an
% image/diagram overlay, e.g. slides/igr-Budgeting-current.tex,
% slides/bdh-Working-Capital-Example-current.tex -- see README.md's
% "Available layouts" section for why minipage instead of columns there) --
% a minipage is a *box*, and TeX's vertical glue insertion (\topskip
% specifically) treats the first box in a vertical list differently from
% the first line of plain paragraph text, so the same \gothamContentTopShift
% pull-up above lands closer to the title for minipage-wrapped content than
% for plain content. Cannot be fixed by \vspace(*) right after
% \begin{minipage}{...}: both plain and starred forms were tried and
% measurably had zero effect (the star only protects glue from page-break
% discard, not from \topskip's own separate glue-insertion calculation,
% which absorbs leading glue at the start of a vertical list regardless).
% Fixed instead by shrinking standardframe's own pull-up for the affected
% frame -- see \gothamContentTopShiftReduction below. The value below is
% calibrated empirically, not derived from first principles: an isolated
% two-frame test (single-level itemize, one plain, one minipage-wrapped)
% measured only a 4.4pt gap, but real slides (nested itemize inside the
% minipage) showed a larger ~7.5-7.6pt gap against a real plain-itemize
% slide (slides/jyn-Categorization-current.tex) -- the extra few points
% likely come from itemize's own topsep/partopsep being computed differently
% depending on whether the list opens a box or a paragraph, on top of the
% topskip/box effect alone. Tuned so slides/igr-Budgeting-current.tex and
% slides/bdh-Working-Capital-Example-current.tex's first content line
% lands within 0.1pt of slides/jyn-Categorization-current.tex's own
% (200.917pt vs. 200.871pt/200.980pt in the compiled PDF).
\newlength{\gothamMinipageTopCorrection}
% 9.38pt (2026-07-22, re-measured): the previous 7.6pt value went stale
% the moment \gothamContentTopShift itself was increased (7.33pt ->
% 11.5964pt, "reduce the Top Margin", see that length's own comment above)
% without a matching update here -- confirmed directly (pymupdf against a
% freshly-compiled preview deck, not assumed from the \gothamContentTopShift
% delta alone, which doesn't translate 1:1 into a minipage's own \topskip
% behavior): four independent slides all using this constant completely
% unmodified (slides/fr3-Forms-of-Capital-Market-Efficiency-Strong-Semistrong-Weak-current.tex,
% slides/za2-Startup-Financing-Categorization-current.tex,
% slides/yke-Implications-of-Homemade-Dividends-current.tex,
% slides/6j4-Modigliani-Miller-Assumptions-current.tex) all land their
% own first line at an identical y=39.80pt post-change, 1.78pt short of
% this deck's current 41.58pt Top Margin reference -- not slide-specific
% noise, the same deterministic amount on every one of them. Re-measure
% and adjust again if \gothamContentTopShift ever changes in the future.
\setlength{\gothamMinipageTopCorrection}{9.38pt}

% A *third* minipage-top correction, on top of \gothamMinipageTopCorrection
% above rather than replacing it, for a minipage whose own first internal
% content is a zero-height `\vbox to 0pt{...}` tikz overlay (a
% corner-anchored decorative image, placed there specifically so it stays
% unconditionally visible even when the minipage's own later content
% includes a `\pause` -- see e.g.
% slides/y5z-CoCo-Bonds-current.tex's own "v1-2-0" note for the fuller
% pause-inheritance story behind why the image has to sit there at all).
% 2026-07-08, found while chasing "the top margin of the text on y5z/22s
% is too low" -- measured directly against the compiled PDF, y5h/22s both
% landed 16.5pt too far down despite already using
% \gothamMinipageTopCorrection, with the zero-height image overlay sitting
% *before* the minipage entirely as the actual root cause (that ordering
% makes the image, not the minipage, the frame body's true first box,
% which \topskip's own special first-box glue handling reacts to
% differently than a plain first minipage). Moving the image inside the
% minipage as its own first content closed most of the gap on its own
% (restoring the minipage as the frame's true first box) but left a
% further 10.35pt uncorrected -- confirmed the same on two independent
% slides (y5z, 22s) with identical structure, not assumed transferable
% from one to the other. Call via
% `\gothamReduceContentTopShift{\dimexpr\gothamMinipageTopCorrection-\gothamMinipageWithCornerImageExtraCorrection\relax}`
% for any slide combining this exact pattern (corner image as a minipage's
% own first content, ahead of a `\pause`) -- confirmed transferable as-is
% to two more slides with the same structure
% (slides/qbj-Bank-Loans-Credit-Lines-current.tex,
% slides/jkf-Bank-Loans-Term-Loan-current.tex), neither of which
% actually uses `\pause` itself, so this isn't specific to that -- the
% image-inside-minipage-as-first-content shape alone is what triggers it.
\newlength{\gothamMinipageWithCornerImageExtraCorrection}
\setlength{\gothamMinipageWithCornerImageExtraCorrection}{10.35pt}

% A *different* minipage-top correction, specifically for a minipage+
% itemize sitting as the first content of a `sideimageframe` (2026-07-06,
% found while chasing "the top margin of the text on ezz/acw is again
% off" -- see CONVERSION_NOTES.md's own "beamer/TeX gotchas" entry for the
% full story). `\gothamMinipageTopCorrection` above is calibrated for a
% *plain* standardframe, and simply does not transfer to `sideimageframe`:
% measured directly on three real sideimageframe slides
% (slides/xmp-Implications-of-Capital-Market-Efficiency-Considerations-current.tex,
% slides/dfc-Angel-Investors-current.tex,
% slides/ezz-Examples-of-Tests-for-Market-Efficiency-current.tex --
% one flat itemize, one flat itemize, one nested-first-item itemize
% respectively, so this isn't a flat-vs-nested effect either), all three
% landed within 0.6pt of each other at their own *uncorrected* natural
% position (52.03-52.57pt against this deck's usual 44.0pt first-line
% reference) regardless of which of the two *wrong* corrections
% (`\gothamReduceContentTopShift{\gothamMinipageTopCorrection}` or
% `\gothamIncreaseContentTopShift{\gothamFlatItemizeTopCorrection}`) had
% been applied to each beforehand -- i.e. `sideimageframe` itself pushes
% its own first-content minipage down by a consistent ~8.2pt relative to
% a plain standardframe, on top of whatever generic itemize-shape
% correction would otherwise apply, most likely from the extra
% \sidefootnotes/\settowidth setup work `sideimageframe`'s own opening
% code does before `\begin{standardframe}` ever runs. Use
% `\gothamIncreaseContentTopShift{\gothamSideImageFrameMinipageCorrection}`
% (note: *increase*, not reduce) as the *only* topshift call on any
% future sideimageframe slide whose first content is a minipage+itemize --
% not the generic `\gothamMinipageTopCorrection`/
% `\gothamFlatItemizeTopCorrection` this deck uses everywhere else.
\newlength{\gothamSideImageFrameMinipageCorrection}
\setlength{\gothamSideImageFrameMinipageCorrection}{8.2pt}

% 6.26pt (2026-07-22, re-measured): stale for the same reason
% \gothamMinipageTopCorrection above went stale -- \gothamContentTopShift
% itself changed (7.33pt -> 11.5964pt) without a matching update here.
% Re-measured directly against a freshly-compiled preview deck across six
% independent sideimageframe slides using this constant unmodified: four
% (ezz, xmp, hzx, u6n) land their own first line at an identical
% y=39.64pt, 1.94pt short of this deck's current 41.58pt Top Margin
% reference (i.e. already pulled up *too far*); the other two (c7w, dfc)
% are close but not identical (39.53pt/-2.05, 40.18pt/-1.4) -- no single
% value fixes all six exactly, so this recalibrates to the four-way
% cluster (the majority case) and accepts a small (<0.6pt) residual on the
% two outliers. Since this constant feeds
% `\gothamIncreaseContentTopShift{}` (more pull-up, i.e. *smaller* y), and
% the observed y was already short of target, the fix is a *decrease*
% (8.2pt - 1.94pt), not an increase -- caught by re-measuring after a
% first, wrong-direction attempt at 10.14pt drove the four-slide cluster
% to y=37.70pt, further from target instead of closer.
\setlength{\gothamSideImageFrameMinipageCorrection}{6.26pt}

% Opt-in override: reduces standardframe's \gothamContentTopShift pull-up by
% this amount for the frame currently being opened. Empty-default-style flag
% (mirrors \gothamChapterLogoPath etc.): 0pt by default (no effect), set via
% \gothamReduceContentTopShift{<length>} called *before* \begin{standardframe},
% reset back to 0pt in standardframe's own closing code so it never leaks
% into the next frame. Used by minipage-based slides to counteract
% \gothamMinipageTopCorrection above.
\newcommand{\gothamContentTopShiftReduction}{0pt}
\newcommand{\gothamReduceContentTopShift}[1]{\global\def\gothamContentTopShiftReduction{#1}}

% A *second*, separate first-line-position quirk, found and fixed
% 2026-07-05 after a run of slides (batch 4 plus several earlier ones) were
% found to land at two visibly different heights depending purely on
% content shape, with no minipage involved at all: a plain (non-minipage)
% standardframe whose top-level itemize's very *first* \item is immediately
% followed by its own nested \begin{itemize}/\begin{enumerate} lands 1.834pt
% higher than one whose first \item has no such immediate child (checked
% against every content slide converted so far via each compiled page's own
% first-line PDF baseline coordinate, not eyeballed -- 13 "nested-first-item"
% slides all landed within 0.05pt of each other, and 9 "flat-first-item"
% slides all landed within 0.05pt of each other, at a clean +1.834pt offset
% from the first group; isolated two-frame test cases, differing only in
% presence/absence of that immediate nested list, reproduced the exact same
% 1.834pt gap independent of which text/list type/item count was used).
% Not chased to one named TeX primitive (most likely some interaction of
% \addvspace's glue-coalescing behaviour with \partopsep/\topsep at the
% \begin{itemize}/\begin{frame} boundary, given the existing
% \gothamMinipageTopCorrection comment above already flags topsep/partopsep
% as varying with box vs. paragraph context) -- calibrated empirically like
% \gothamMinipageTopCorrection, not derived from first principles.
% Interacts with minipage differently, not additively: inside a minipage,
% this nested-vs-flat distinction disappears entirely (both land within
% 0.05pt of each other) -- so this correction is for *plain* itemize-first
% standardframes only, never combined with \gothamMinipageTopCorrection.
\newlength{\gothamFlatItemizeTopCorrection}
\setlength{\gothamFlatItemizeTopCorrection}{1.834pt}

% 0.01pt (2026-07-22, re-measured): stale for the same reason
% \gothamMinipageTopCorrection above went stale -- \gothamContentTopShift
% itself changed (7.33pt -> 11.5964pt) without a matching update here.
% Re-measured directly against three independent flat-itemize slides using
% this constant unmodified (574, 5cr, c8e): all three landed their own
% first line at an identical y=39.76pt, 1.82pt short of this deck's
% current 41.58pt Top Margin reference -- coincidentally, the recalibrated
% \gothamContentTopAlign reference now already accounts for almost all of
% what this correction used to supply on its own, leaving next to nothing
% left to add. Effectively zeroed out rather than removing the call
% itself, so every flat-itemize slide's own
% `\gothamIncreaseContentTopShift{\gothamFlatItemizeTopCorrection}` line
% keeps working unchanged.
\setlength{\gothamFlatItemizeTopCorrection}{0.01pt}

% Opposite-direction counterpart to \gothamReduceContentTopShift: *increases*
% the pull-up instead of reducing it (needed for \gothamFlatItemizeTopCorrection
% above, since a flat first item needs *more* pull-up, not less). Same
% storage variable, just negated -- \gothamContentTopShiftReduction is
% allowed to go negative, since the frame-opening code only ever computes
% \gothamContentTopShift-\gothamContentTopShiftReduction, never checks its
% sign.
\newcommand{\gothamIncreaseContentTopShift}[1]{\global\def\gothamContentTopShiftReduction{-#1}}

% Define custom frame type with en-dash bullets
\newenvironment{standardframe}[2][t]{%
	% Restart footnote numbering (the *, dagger, ... symbols) fresh on every
	% slide, rather than letting it run on cumulatively across the whole deck
	% the way beamer does by default. Two reasons: \fnsymbol only defines 9
	% symbols before erroring out, which a real multi-slide deck would run
	% through in no time if every \footnote{}/\source{} call across the whole
	% talk drew from one shared counter; and each slide's footnotes are read
	% and referenced only on that one slide anyway, so there's no reason for
	% their numbering to depend on how many footnotes earlier slides happened
	% to have. This must run *before* \begin{frame}: beamer saves the
	% counter's current value as that frame's baseline as part of
	% \begin{frame}'s own internal setup (restoring it before every overlay
	% replay, so overlays within one frame don't double-count) -- resetting
	% after \begin{frame} would be reset too late, after that baseline was
	% already captured.
	\setcounter{footnote}{0}%
	\begin{frame}[#1]{#2}%
	% Forces this frame's own body to start with the correct default text
	% color, regardless of what any *previous* frame left active
	% (2026-07-16, direct instruction, root-caused via an isolated minimal
	% reproduction). \titleframe's/\quoteframe's own
	% \color{\gothamTitleTextColor} calls -- set inside content built via a
	% stored box register or a zero-space overlay, used/shipped out *later*
	% rather than immediately -- don't reliably get "undone": xcolor's
	% \color doesn't emit a PDF color operator immediately when called, only
	% when the *next* real glyph/path actually needs it, so whatever color
	% one of those frame types last set can end up as the active PDF color
	% well into the *next* frame's own first draw call. Invisible for Graz
	% (the leaked color, c_primary_dark, happens to match ordinary body text
	% anyway) but catastrophic for MCI/Innsbruck (c_primary_white, rendering
	% plain non-itemize body text invisible against its own white
	% background -- itemize items escape it, since beamer's own itemize
	% template sets color fresh per bullet). An \AtBeginEnvironment{frame}
	% hook (etoolbox) was tried first and confirmed *not* to fire at all --
	% beamer's own "frame" isn't a plain \newenvironment etoolbox can patch
	% that way -- so this resets directly in standardframe's own opening
	% code instead, exactly where this project already has full control.
	\color{c_primary_dark}%
	% Defensive: the tokens after \begin{frame} are collected into the frame
	% body and replayed once per overlay slide, so this clears the side-
	% footnote collector at the start of every slide's pass -- normally a
	% no-op (the previous slide's footline already consumed the box), but it
	% stops stale footnotes leaking in if a previous frame ever collected
	% without shipping a footline.
	\gothamClearSideFootnotesBox%
	\setbeamertemplate{itemize item}{\gothamBulletDot}%
	\setbeamertemplate{itemize subitem}{\gothamBulletDot}%
	\setbeamertemplate{itemize subsubitem}{\gothamBulletDot}%
	\setbeamertemplate{itemize subsubsubitem}{\gothamBulletDot}% 4th level, see its own comment further below (search "subsubsub")
	% Only for top-aligned frames (the default, #1=t): a [c]-centered frame
	% (2026-07-22, direct instruction, found while centering
	% slides/bw9-Pearson-Correlation-Illustration-current.tex/agc/crj) was
	% landing its content almost flush against the title bar instead of
	% truly centered -- root cause was this same pull-up firing
	% unconditionally regardless of #1. beamer's own [c] option already
	% centers the frame body within the available space; stacking this
	% deck's *own* top-alignment pull-up on top of that just drags the
	% already-centered box further upward by the same fixed amount, which
	% is exactly the asymmetry that was measured (a [c] figureframe's image
	% landing ~0pt from the title bar and ~10pt from the footer instead of
	% symmetric). Guarded with \ifstrequal (etoolbox, already loaded above)
	% rather than \ifx, matching this file's own \ifstrequal{##1}{...}
	% pattern elsewhere (\ifx would need #1 fully expanded first).
	\ifstrequal{#1}{c}{}{\vspace*{-\dimexpr\gothamContentTopShift-\gothamContentTopShiftReduction\relax}}%
}{%
	\end{frame}%
	\global\def\gothamChapterLogoPath{}% opt-in only for the frame just closed
	\global\def\gothamContentTopShiftReduction{0pt}% opt-in only for the frame just closed
	\global\gothamSideFootnotesfalse% opt-in only for the frame just closed
	\global\def\gothamFrametitleFullHeightYellow{}% opt-in only for the frame just closed
	\global\def\gothamFrametitleLogoFixed{}% opt-in only for the frame just closed
	\global\def\gothamFooterBoxColor{\gothamFooterBoxColorDefault}% opt-in only for the frame just closed
}

% The bullet itself: a small filled circle, matching the source PowerPoint
% deck's own round bullets (a plain \textbullet, tried first, is too
% small/thin in Fira Sans at any standard text size to read as a solid dot
% -- a tikz-drawn circle gives full control over the actual rendered size
% instead of depending on the font's own glyph design). baseline=-0.65ex
% roughly centers it on the surrounding text's x-height, the same vertical
% position a font's own bullet glyph would normally sit at.
%
% Its fill color is \gothamBulletColor, not a literal color token here --
% set by the per-university preamble before \input{}ing this file, same
% "value supplied externally, referenced lazily" idiom as
% \gothamInstituteLogoPath/\gothamLogoHeight above (\gothamBulletDot's own
% body only expands \gothamBulletColor when a bullet actually gets drawn,
% long after this \newcommand runs). Graz's own bullets are its primary
% highlight color (matching this project's original, single-institution
% behavior); MCI's and Innsbruck's are their own secondary blue instead
% (2026-07-16, direct instruction) -- see each institution's own
% Slidemaster_preamble_*.tex for its \gothamBulletColor setting.
\newcommand{\gothamBulletDot}{\tikz[baseline=-0.65ex]{\fill[\gothamBulletColor] circle (1.7pt);}}

% "a)"/"b)"/"c)"-labeled list, for exercise sub-tasks that read better
% lettered than bulleted (2026-07-20, direct instruction, first used on
% slides/y3a-Exercise-Diversification-2-current.tex's own bottom three
% items).
%
% First attempt (superseded, kept as a documented wrong turn -- see
% README.md's own "lualatex vs latex" section for why this project keeps
% wrong write-ups rather than deleting them): reuse itemize's own top-level
% `\llap`-based label positioning by locally overriding
% `\setbeamertemplate{itemize item}` inside a plain `\begin{itemize}` --
% the reasoning being that `\llap` draws its content flush against the
% text's own left edge "regardless of the label's width." That reasoning
% turned out to be wrong in a way that only showed up once actually
% measured against the compiled PDF (not just visually eyeballed): a plain
% bullet's own `\llap`-drawn footprint (title/Left-Margin-relative-left-
% edge to text-start) measured only ~9.05pt end-to-end on this deck's real
% slides, matching the Left-Margin-to-text-start gap exactly -- but "a)"
% rendered through the *exact same* mechanism measured a ~13.97pt
% footprint for the same gap, not the ~8.5-9pt `\settowidth` reported for
% the bare glyphs in isolation. The extra ~5pt was never fully root-caused
% (this custom itemize's own `\list` call, further below, only overrides
% `\makelabel` -- it never sets `\leftmargin`/`\labelwidth` itself, so
% those are left at whatever this deck's global itemize defaults happen to
% be, and `\llap`'s actual box-width behavior inside that inherited,
% unexamined list-parameter state is where the extra width most likely
% comes from) -- rather than keep guessing at *why*, measured the actual
% failure directly (the compiled "a)" sat ~4.9pt left of the Left Margin,
% confirmed via exact PDF text-position extraction, not a visual guess)
% and abandoned the approach once even `\tiny` plus manual per-letter
% kerning only just barely closed the gap while reading as illegibly small
% -- a real dead end, not a tuning problem.
%
% Working fix: don't reuse itemize's `\llap` label mechanism at all -- use
% a real, width-reserving `\makebox[W][r]{...}` label instead (the same
% *kind* of mechanism `\gothamEnumerateLabelWidth` above already uses for
% a nested enumerate, just with an explicit `\leftmargin`/`\labelwidth`/
% `\labelsep` this environment sets itself rather than inheriting), inside
% its own plain `\list`, not a reused `\itemize`. `\leftmargin` (label box
% + gap) is set to 11.1pt -- 9.1pt label box (fits "a)"/"b)"/"c)" at this
% deck's normal `\small` itemize size with a little clearance, confirmed
% directly: measured label-left-edge stays right of the 12.552pt Left
% Margin for all three letters, closest case (widest letter, "c)") 1.37pt
% clear) plus 2pt `\labelsep` for a visible gap before the text. Trade-off,
% accepted deliberately rather than found by accident: item text lands
% ~2.2pt right of where a plain bulleted item's text would (23.8pt vs.
% 21.6pt) -- fitting a legible "a)"/"b)"/"c)" at normal reading size plus
% a real gap simply doesn't fit in the ~9pt a plain bullet's tiny dot
% leaves, so *exact* text-position parity with plain bullets and a legible
% label are mutually exclusive here; legibility won given the two-point
% shift is visually negligible and there's no plain-bulleted item left on
% the same slide to compare against directly anyway (y3a's own first item
% moved into its infoblock). Re-derive `\leftmargin`/`\labelwidth` the same
% way (render, measure with exact text-position extraction, adjust) if
% this is ever used at a different font size than this deck's normal
% itemize `\small`.
\newcounter{gothamalphaitem}
\newcommand{\gothamAlphaItem}{\stepcounter{gothamalphaitem}\alph{gothamalphaitem})}
\newenvironment{gothamalphaitemize}{%
	\setcounter{gothamalphaitem}{0}%
	\list{\gothamAlphaItem}{%
		% 4.75pt (2026-07-20, direct instruction: "half the width of the
		% 'C' in 'Calculate'" -- \settowidth{C} at this deck's normal
		% itemize \small measured 5.5pt, so 2pt original + 2.75pt half-C).
		% Only \labelsep changes the text's own start position -- the
		% label's own box stays exactly where it was (its left edge is
		% pinned to this list's origin regardless of \labelsep, since
		% \leftmargin=\labelwidth+\labelsep always cancels \labelsep back
		% out of the label box's own position), so this widens the gap
		% without reopening the Left-Margin question already solved above.
		\setlength{\labelsep}{4.75pt}%
		\setlength{\labelwidth}{9.1pt}%
		\setlength{\leftmargin}{\labelwidth}%
		\addtolength{\leftmargin}{\labelsep}%
		\setlength{\itemindent}{0pt}%
		% \itemsep/\parsep left at plain \list's own defaults otherwise --
		% since this environment is built directly on \list rather than
		% reusing beamer's own itemize template, it never got beamer's own
		% itemize spacing hooks at all, only whatever \list falls back to.
		% Measured directly (2026-07-21, direct instruction, found on
		% slides/nt4-Exercise-CAPM-5-current.tex): that default reads as
		% *less* space between two items (13.7pt, item-to-item) than
		% between two wrapped lines within one *single* item (15.6pt,
		% pure \baselineskip) -- backwards from how a list should read
		% (item boundaries should be at least as open as a same-item line
		% wrap, not tighter). \itemsep set explicitly so the total gap at
		% an item boundary (baselineskip + itemsep) reads as clearly more
		% open than a plain wrapped line, not accidentally tighter.
		% Reduced from the original 6pt to 2.5pt (2026-07-22, found during
		% a full-project check after the itemize-rules work, not part of
		% that work itself): 6pt overshot by more than needed to clear the
		% 1.9pt gap (15.6 vs. 13.7), and slides/kmf-Exercise-CAPM-6-current.tex's
		% own 5-item list -- longer than nt4's own -- compounded that
		% excess across 4 item boundaries into a ~25pt overfull frame, its
		% own last item cut off by the footer (confirmed directly,
		% rendered and compared against a pre-regression PDF). 2.5pt
		% keeps a real, visible margin over the 1.9pt needed (not just
		% barely crossing it) while re-measuring clean (no overfull) on
		% kmf; re-verify against both nt4 and kmf directly if this is
		% ever changed again -- they're this environment's own tightest
		% (most items) and originally-tuned-against cases respectively.
		\setlength{\itemsep}{2.5pt}%
	}
	\raggedright% plain \list defaults to justified+hyphenated text, unlike
	% this deck's own itemize (which calls this too, further below) --
	% missed on the first pass, only caught by comparing a rendered a)/b)/c)
	% item side by side against a plain bulleted one ("stan-\ndard deviation"
	% hyphenating where a plain-bulleted version of the same sentence
	% elsewhere in this deck doesn't).
}{%
	\endlist
}

% Itemize indentation: level 1 bullets should sit flush with the frame title
% (no extra left margin), and each deeper level should be indented just
% enough for its bullet to line up under the *text* (not the bullet) of the
% level above. A nested list's \leftmarginN is added on top of its parent's
% (indentation is cumulative), so giving every level the same "step" --
% label width + labelsep, the width a bullet plus its gap to text takes up --
% makes level N's text sit at N*step and level N's bullet at (N-1)*step, i.e.
% exactly at level (N-1)'s text (0 for level 1, flush with the title).
\makeatletter
\newlength{\gothamBulletWidth}
\settowidth{\gothamBulletWidth}{\gothamBulletDot}
\newlength{\gothamItemStep}
\setlength{\gothamItemStep}{\labelsep}
\addtolength{\gothamItemStep}{\gothamBulletWidth}
\setlength{\leftmargini}{\gothamItemStep}
\setlength{\leftmarginii}{\gothamItemStep}
\setlength{\leftmarginiii}{\gothamItemStep}
\setlength{\leftmarginiv}{\gothamItemStep}
\makeatother

% A 4th itemize level -- beamer's own \itemize (beamerbaselocalstructure.sty)
% hard-caps nesting at 3 levels (\ifnum \@itemdepth >2\relax\@toodeep),
% independently of anything this deck sets: a 4th \begin{itemize} always
% raised "! LaTeX Error: Too deeply nested." even though plain LaTeX's own
% kernel itemize supports 4. Found 2026-07-15, direct instruction, on
% slides/hrg-Exchange-Rate-Quotation-Forward-current.tex -- genuinely
% 4 levels deep in the source pptx itself (confirmed via a PowerPoint COM
% render, not assumed), not a conversion artifact to design around.
% Since this deck already styles every existing level identically (same
% \gothamBulletDot, same \gothamItemStep margin -- see the comments right
% above/below), extending to a 4th level needs no new visual design, just
% patching beamer's own depth limit up by one and telling it what to call
% the new level ("subsubsub", the next name in beamer's own
% sub/subsub/... naming scheme) so its template/color lookups
% (\usebeamertemplate{itemize subsubsubitem}, \usebeamercolor{itemize/
% enumerate subsubsub body}, etc.) resolve to *something* instead of
% erroring -- undefined beamer templates/colors are expected to fall back
% to sensible defaults (no extra spacing hook, plain text color) rather
% than raising an error, same as any other beamer template/color that a
% given theme never bothers to set; the actual bullet symbol still needs
% its own \gothamBulletDot template below, added explicitly the same way
% as the existing 3 levels, not left to whatever fallback beamer has for
% an unset "itemize subsubsubitem" template.
% \beamer@computepref is patched wholesale (not via \patchcmd) since its
% body is one single \ifcase chain, not a stable string to match against.
% \itemize itself was *first* tried via \patchcmd{\itemize}{...}
% (matching this deck's own established practice for beamer/theme
% internals, e.g. beamerfontthemegotham.sty's own
% \patchcmd{\beamer@title}{...}) -- this failed outright, confirmed
% directly (\patchcmd's own failure branch fired, raising a clean
% "Patching failed" error rather than silently doing nothing): recent
% LaTeX2e kernels implement an optional-argument environment like
% beamer's own `\renewcommand{\itemize}[1][]{...}` via
% `\@protected@testopt`, which stores the *real* body not under `\itemize`
% itself but under a differently-named internal control sequence (`\show
% \itemize` confirms this directly: it prints as
% `\@protected@testopt \itemize \\itemize {}`, i.e. plain `\itemize` is
% only a thin dispatcher, and the actual `\ifnum \@itemdepth >2...` body
% \patchcmd was trying to match against was never there to find). Rather
% than chase that internal storage name (liable to differ across LaTeX
% kernel versions, which \patchcmd's own documented usage doesn't need to
% care about), \itemize is fully re-declared here instead, copied
% verbatim from beamer's own beamerbaselocalstructure.sty (this project's
% installed version) with only the depth bound itself changed from 2 to
% 3 -- the same declaration form beamer itself used
% (\renewcommand{\itemize}[1][]{...}), which regenerates both the
% dispatcher and the real body correctly under whatever internal name
% the current kernel expects, sidestepping the naming question entirely.
% Trade-off, stated plainly: if a future beamer release changes this
% macro's own body for an unrelated reason, this override will silently
% keep using today's (2026-07) copy instead of picking up that change --
% acceptable here since the override is one line different from upstream
% and clearly marked, but re-diff against beamerbaselocalstructure.sty's
% own \itemize if beamer is ever upgraded and itemize behavior looks off.
\makeatletter
\long\def\beamer@computepref#1{%
	\let\beameritemnestingprefix\@empty%
	\ifcase#1\or\or\def\beameritemnestingprefix{sub}\or\def\beameritemnestingprefix{subsub}\or\def\beameritemnestingprefix{subsubsub}\or\@toodeep\fi%
}
\renewcommand{\itemize}[1][]{%
	\ifblank{#1}{}{\def\beamer@defaultospec{#1}}%
	\ifnum \@itemdepth >3\relax\@toodeep\else
		\advance\@itemdepth\@ne
		\beamer@computepref\@itemdepth% sets \beameritemnestingprefix
		\usebeamerfont{itemize/enumerate \beameritemnestingprefix body}%
		\usebeamercolor[fg]{itemize/enumerate \beameritemnestingprefix body}%
		\usebeamertemplate{itemize/enumerate \beameritemnestingprefix body begin}%
		\list
			{\usebeamertemplate{itemize \beameritemnestingprefix item}}
			{\def\makelabel##1{%
					{%
						\hss\llap{{%
								\usebeamerfont*{itemize \beameritemnestingprefix item}%
								\usebeamercolor[fg]{itemize \beameritemnestingprefix item}##1}}%
					}%
				}%
			}
	\fi%
	\beamer@cramped%
	% \beamer@cramped (beamerbaselocalstructure.sty, unmodified upstream
	% code) unconditionally does `\itemsep0pt\topsep0.5pt`, every single
	% time -- called right here, in this project's own copy of beamer's
	% own \itemize body, same position as upstream. Setting \itemsep any
	% earlier (first tried as a separate \AtBeginEnvironment{itemize} hook,
	% run *before* this whole macro body even starts) gets silently wiped
	% out by this same line a moment later -- confirmed directly via
	% \typeout debugging (not assumed): the hook's own \setlength really
	% did run with the right values, \itemsep still read back as 0pt
	% immediately after. Setting it *here*, right after \beamer@cramped
	% already ran (not before), is what actually survives.
	\ifcase\@itemdepth
	\or\setlength{\itemsep}{\gothamItemsepRatioI em}% level 1
	\or\setlength{\itemsep}{\gothamItemsepRatioII em}% level 2
	\or\setlength{\itemsep}{\gothamItemsepRatioIII em}% level 3
	\or\setlength{\itemsep}{\gothamItemsepRatioIV em}% level 4
	\fi%
	\raggedright%
	\beamer@firstlineitemizeunskip%
}
\makeatother

% A nested \begin{enumerate} inherits the same \leftmarginN step as itemize
% above for *where its text starts* (matching itemize's own "one step past
% the parent level's text" rule -- confirmed directly against the compiled
% PDF: an enumerate's own body text, e.g. "Senior secured debt", already
% lands exactly one \gothamItemStep past its parent item's text), but its
% *label* ("1.", "2.", ...) doesn't follow \labelwidth the way plain LaTeX's
% list mechanism would -- beamer draws enumerate labels through its own
% template machinery, which turned out not to respect \labelwidth at all:
% \setlength{\labelwidth}{...} right after \begin{enumerate} measurably took
% effect (checked with \typeout{\the\labelwidth} immediately after) but the
% compiled numeral's own position never moved. What *does* move it: directly
% overriding \makelabel with a fixed-width \makebox -- empirically
% calibrated (render, measure the numeral's own x-position in the compiled
% PDF, adjust, repeat) since the box's own left edge turned out to sit at a
% fixed position unrelated to \labelwidth or \leftmarginN, with the
% right-aligned numeral inside it moving in lockstep with the box's width
% (confirmed linear: +1pt width -> +~1pt numeral position, across several
% test widths). 19.23pt (at this deck's default itemize font size) lands a
% single-digit numeral flush with the parent item's own text -- only valid
% at that font size and for single-digit lists (1-9); re-calibrate the same
% way (not by recomputing from first principles) for a differently-sized or
% 10+-item nested enumerate. Call right after \begin{enumerate}, before the
% first \item.
%
% 20.33pt (2026-07-08, direct instruction: "the numbers still extend a
% little farther to the left than the text above them does" --
% slides/f6m-Recovery-Rates-current.tex's own "Pecking order:" list).
% 19.23pt landed *close* but not quite flush: measured directly against
% the compiled PDF, every numeral's own left edge (20.55pt-20.95pt,
% varying slightly by digit width even though all are right-aligned to
% the same box) still sat up to ~1.05pt left of "Pecking order:"'s own
% text start (21.6pt) -- the widest digit tried, "6.", was the binding
% case. +1.1pt of box width (confirmed +1pt width -> +1pt numeral
% position, per this comment's own note above) covers that with a small
% margin, re-confirmed against the recompiled PDF.
\newcommand{\gothamEnumerateLabelWidth}{\renewcommand{\makelabel}[1]{\makebox[20.33pt][r]{##1}}}

% Automatic version of the above (2026-07-23, background-task follow-up
% flagged in slides/kfh-Asset-Allocation-Process-current.tex's own v1-1-0
% comment: "fix the enumerate behavior generally, if possible", task
% `task_bb3fbb66`). That comment's first attempt, an
% `\AtBeginEnvironment{enumerate}{...}` hook, failed because the hook fires
% *before* `\enumerate`'s own body runs, and that body's own internal
% `\list{...}{\def\makelabel...}` setup (the beamer copy directly below the
% `\itemize` override above) clobbers whatever a hook that early leaves
% behind -- the exact same "template runs its own setting, external pre-set
% gets overwritten" problem `\itemize`'s own `\itemsep` comment above
% already documents. The real fix, same idea `\itemize` itself used above:
% override beamer's own low-level piece directly rather than hooking
% \AtBeginEnvironment. Unlike `\itemize` (declared via
% `\@protected@testopt`, needing a full re-declaration to reach -- see the
% `\itemize` override's own comment for the full story), `\enumerate`'s
% relevant piece, `\beamer@enum@` (beamerbaselocalstructure.sty), is a
% plain `\def` with no optional-argument indirection, so it can be
% `\renewcommand`'d directly, verbatim from upstream except for the
% `\ifnum\@itemdepth>1` branch added below.
%
% Only fires for a *nested* enumerate (`\@itemdepth>1`, i.e. already inside
% at least one `\itemize`/`\enumerate` when this one opens) -- a top-level
% enumerate (`\@itemdepth=1`, e.g.
% slides/t9z-Outland-Steel-Hedging-Borrowing-Investing-current.tex/
% slides/g4c-Outland-Steel-Hedging-Forward-Sale-current.tex/
% slides/j84-Outland-Steel-No-Hedging-current.tex's own outermost numbered
% list, or slides/nkj-Agenda-current.tex's own \tableofcontents-driven one)
% is left at beamer's stock label position entirely untouched, matching
% this project's own established behavior for that shape (never reported
% as misaligned, unlike the nested case) -- and, for nkj specifically,
% doesn't matter either way, since that slide's own
% `\gothamTocLabelFlushLeft` call (right after its own `\begin{enumerate}`)
% overrides `\makelabel` again regardless of whatever this default sets.
%
% Reuses the exact same 20.33pt calibration as the manual
% `\gothamEnumerateLabelWidth` above, confirmed (pymupdf, compiled PDF) to
% land correctly not just at `\@itemdepth=2` (one `\itemize` above the
% enumerate -- slides/kfh-Asset-Allocation-Process-current.tex/
% slides/f6m-Recovery-Rates-current.tex/
% slides/6pz-Outland-Steel-Hedging-Strategies-Overview-current.tex, all
% previously needing the manual call) but also at `\@itemdepth=3` (two
% `\itemize`s above it -- slides/bdh-Working-Capital-Example-current.tex's
% own "Development of balance sheet:" list, also previously needing the
% manual call, with the identical 20.33pt value) -- this deck's itemize/
% enumerate templates don't shrink font size by nesting depth (only the
% bullet glyph/margin step does), so the numeral's own glyph width -- the
% only thing \gothamEnumerateLabelWidth's own box width actually
% compensates for -- stays the same at every depth, and one constant covers
% both cases without needing a depth-dependent formula. Not verified beyond
% depth 3 (no slide currently nests an enumerate that deep); if one ever
% does, re-measure before trusting this constant that far down, same as any
% other empirical calibration in this file.
% (slides/scs-Political-Risk-Stability-Index-current.tex also matches
% "`\gothamEnumerateLabelWidth`" in a grep -- but only inside its own prose
% comments describing the general technique; its actual call is the
% bespoke `\gothamEnumerateLabelFlushLeftShifted{16pt}{10.2pt}` for its own
% alphabetic, left-aligned, narrower-box labels, unaffected by this
% override either way since that call runs after `\begin{enumerate}` opens,
% same ordering as every other post-open `\makelabel` override here.)
%
% Former per-slide manual `\gothamEnumerateLabelWidth` calls removed from
% the 4 slides above now that this covers them automatically; the macro
% itself is kept defined, both for this override to reuse and as a manual
% escape hatch (e.g. a differently-calibrated width some future slide might
% need).
\makeatletter
\renewcommand{\beamer@enum@}{%
	\beamer@computepref\@itemdepth% sets \beameritemnestingprefix
	\usebeamerfont{itemize/enumerate \beameritemnestingprefix body}%
	\usebeamercolor[fg]{itemize/enumerate \beameritemnestingprefix body}%
	\usebeamertemplate{itemize/enumerate \beameritemnestingprefix body begin}%
	\expandafter
		\list
			{\usebeamertemplate{\beamer@enumtempl}}
			{\usecounter\@enumctr%
				\ifnum\@itemdepth>1
					\def\makelabel##1{\makebox[20.33pt][r]{##1}}%
				\else
					\def\makelabel##1{{\hss\llap{{%
							\usebeamerfont*{enumerate \beameritemnestingprefix item}%
							\usebeamercolor[fg]{enumerate \beameritemnestingprefix item}##1}}}}%
				\fi}%
	\beamer@cramped%
	\raggedright%
	\beamer@firstlineitemizeunskip%
}
\makeatother

% Extra breathing room when a nested itemize closes and control returns to
% its *parent* list, right before that parent's next \item -- e.g. sub-items
% under "Mid-Term Exam" followed by the next top-level "Final Exam" -- so a
% new top-level group reads as visually distinct from the one before it,
% rather than the same uniform gap used between every other item regardless
% of level. \@itemdepth (LaTeX's own nesting-level counter, already
% decremented back to the parent's level by the time \end{itemize}'s normal
% closing has run) says which level control returned to -- 0 means the
% outermost list itself just closed (nothing left to add space before,
% deliberately excluded), 1 means back to a level-1 item, 2 means back to a
% level-2 item.
%
% Level-dependent (2026-07-06, direct instruction) -- a flat, single
% constant regardless of *which* level control returned to (the original
% version of this hook) reads as "extra"/inconsistent precisely because a
% return-to-level-1 and a return-to-level-2 transition both got the exact
% same gap, rather than level 1 getting a visibly bigger gap than level 2 --
% found on slides/57r-Reasons-for-Leasing-current.tex, where a level-2
% sub-list closing back to a level-1 item looked like a random extra gap
% rather than "level 1 items are consistently more separated than level 2
% ones." \addvspace (not plain \vspace) matters too: it coalesces with
% whatever glue itemize's own \itemsep/\topsep already contributes at the
% closing boundary (taking whichever is larger) rather than *adding* on top
% of it unconditionally -- \vspace's own stacking behavior was the other
% half of why the old flat version could look like too much space.
\makeatletter
% Ratios, not lengths (2026-07-21, direct instruction: "all of these
% spacings should scale with font size") -- a \newlength+\setlength{...}{0.5em}
% pair, the way this used to be written, freezes to a fixed *pt* value the
% one time \setlength runs (at preamble-load time, using whatever font is
% active then), despite the "em" in its own definition -- it does NOT
% re-track whichever font size is actually active on a given slide/level
% later. Confirmed directly: a \typeout{\f@size} placed at the point each
% of these actually gets used (inside \AtEndEnvironment{itemize}'s own
% hook body, further below) showed the *closing* list's own font still
% active there (9pt/10pt/10.95pt for a level-4/3/2 list closing,
% respectively -- not yet restored to the parent's own font), confirming
% this hook fires early enough in `\end{itemize}`'s own processing that a
% live "0.5em" computed *there* would correctly track whichever level is
% actually involved. Fixed by storing the bare ratio as a plain macro
% instead, and writing the literal "em" back on at the point of use
% (`\addvspace{\gothamItemizeLevelDropSpaceRatioI em}`), so it's
% recomputed fresh, against whatever font is live, every time.
\newcommand{\gothamItemizeLevelDropSpaceRatioI}{0.5}
\newcommand{\gothamItemizeLevelDropSpaceRatioII}{0.3}
% Level 3's own drop-space, for when a closed *4th*-level list returns
% control to a level-3 item -- added 2026-07-15 alongside the 4th itemize
% level itself (search "subsubsub" further above). Interpolated from the
% existing I/II progression (0.5, 0.3 -- each step roughly 60% of the one
% before) rather than empirically re-measured against a compiled PDF the
% way I/II themselves were: no slide used a 4th level before
% slides/hrg-Exchange-Rate-Quotation-Forward-current.tex to measure
% against. Re-derive properly (render, measure, adjust) if this ever
% looks visibly off on a real slide.
\newcommand{\gothamItemizeLevelDropSpaceRatioIII}{0.2}
% Level 4's own drop-space, for when a closed level-4 list returns control
% to a level-3 item -- added 2026-07-21 alongside \gothamItemsepRatioIV.
% \@itemdepth at this hook's fire time is the depth of the list that *just
% closed*, confirmed directly via a temporary \typeout{\the\@itemdepth}
% (not assumed): closing the level-4 list on
% slides/ZZO-Itemize-Spacing-Reference-current.tex measured itemdepth=4
% there, closing level-3 measured 3, level-2 measured 2, and the outermost
% level-1 list itself closing measured 1 (never 0 -- \@itemdepth's own
% restore to the parent's depth happens *after* this hook runs, so the "0:
% outermost list itself closed" case below is dead code, kept only because
% removing a genuinely unreachable \ifcase branch isn't worth the risk of
% getting the numbering wrong a second time). So this is very much a real,
% reachable case -- not decorative symmetry -- and needed its own
% \addvspace branch (previously missing entirely, silently leaving the
% level-4-to-level-3 closing transition with no drop-space at all).
% Interpolated from the existing I/II/III progression as a starting guess,
% then re-measured directly against the compiled ZZO reference and bumped:
% 0.13 measured at 12.6pt, still *less* than level 4's own same-level gap
% of 13.4pt (violates rule III -- a level's own same-level gap must be
% smaller than that level's own closing-transition gap); raised until
% re-measurement cleared it with a comfortable margin again, matching the
% margin every other level's own pair already has.
\newcommand{\gothamItemizeLevelDropSpaceRatioIV}{0.5}
\AtEndEnvironment{itemize}{%
	\ifcase\@itemdepth
		% 0: dead code -- @itemdepth is never actually 0 when this hook
		% fires (see \gothamItemizeLevelDropSpaceRatioIV's own comment
		% just above); kept only for alignment with the \ifcase branches
		% below, each keyed on the depth of the list that just closed.
	\or\addvspace{\gothamItemizeLevelDropSpaceRatioI em}% level-1 (outermost) list itself closed -- a pre-existing quirk (not introduced by this comment's own edit): there's no parent item above level 1 for this to visually separate from, so this only actually matters if something follows the outermost list directly on the same frame; harmless on every slide seen so far since nothing does. Out of scope to redesign here.
	\or\addvspace{\gothamItemizeLevelDropSpaceRatioII em}% level-2 list closed, back to a level-1 item
	\or\addvspace{\gothamItemizeLevelDropSpaceRatioIII em}% level-3 list closed, back to a level-2 item
	\or\addvspace{\gothamItemizeLevelDropSpaceRatioIV em}% level-4 list closed, back to a level-3 item
	\fi
}

% Per-level \itemsep -- the item-to-item gap *within* one list, at one
% fixed nesting depth (as opposed to the drop-space mechanism just above,
% which only fires when a nested list *closes* back up to its parent).
% Ratios, not lengths -- same "scale with font size" reasoning and fix as
% \gothamItemizeLevelDropSpaceRatioI/II/III above (that comment has the
% full story on why a plain \newlength+\setlength pair here would freeze
% to one fixed pt value instead of actually tracking each level's own
% font size). Values only, defined here for discoverability alongside the
% drop-space ones above -- actually *applied* inside the `\itemize`
% redefinition itself (search "\beamer@cramped (beamerbaselocalstructure.sty"
% further above), not via a separate hook here; see that comment for why.
% Found missing entirely (2026-07-21, direct instruction, see
% slides/ZZO-Itemize-Spacing-Reference-current.tex's own header comment
% for the rules this and the drop-space mechanism together need to
% satisfy): this deck's own itemize never set \itemsep itself at any
% level -- confirmed directly, grepping this whole file -- so every depth
% shared whatever bare default beamer's own theme happens to leave
% \itemsep at, uniformly. That produced two visible problems, both
% measured directly (pymupdf against slides/2c2-Topics-Around-the-CAPM-Empirical-Validity-current.tex
% and an isolated bisection test, not guessed): (a) level-2 and level-3
% same-level item-to-item gaps came out *identical* (12.6pt each,
% confirmed against an isolated reproduction), since nothing distinguished
% them by depth at all -- only a closing transition ever got a
% level-dependent gap, never a same-level one; (b) that shared default
% \itemsep read as *smaller* than a single item's own wrapped-line
% \baselineskip (13.7-13.8pt within a wrap vs. 12.6pt between items),
% backwards from how a list should read -- an item boundary should never
% be tighter than a same-item line wrap.
% Ratios themselves are not arbitrary -- each chosen, then re-measured
% against the compiled ZZO reference slide and adjusted, specifically to
% satisfy all three of that slide's own header-comment rules together
% (extended to level 4 the same way, 2026-07-21), not just "level 1
% biggest" in isolation:
% - Rule I (level-1 same-level gap > level-2's > level-3's > level-4's)
%   falls straight out of \gothamItemsepRatioI > ...II > ...III > ...IV,
%   added on top of the same \baselineskip every level shares (itself
%   already smaller per level here, since this theme's own itemize/
%   enumerate body font shrinks with depth too -- confirmed directly,
%   10.95pt/10pt/9pt/9pt for levels 1-4 respectively, level 4 sharing
%   level 3's own font since this theme has no distinct "subsubsub" size
%   of its own).
% - Rule II (every level's own same-level gap > that level's own wrapped-
%   line gap) needs \gothamItemsepRatioIV alone (the smallest, hence the
%   binding case) to still be positive and large enough that it clearly
%   reads as more open than a plain wrapped line once rendered, not just
%   bigger on paper.
% - Rule III (a level's own same-level gap < that same level's own
%   *closing*-transition gap) is *not* automatic just because
%   \gothamItemizeLevelDropSpaceRatio* was already bigger than the old
%   shared default -- \addvspace's own coalesce-to-the-larger-value
%   behavior (documented above) means the closing gap is
%   max(natural glue, that *drop-space* ratio), not a sum of the two. Each
%   drop-space ratio's own \ifcase branch is keyed on the depth of the
%   list that *just closed* (see \gothamItemizeLevelDropSpaceRatioIV's own
%   comment for the full story on why, confirmed via \typeout, not
%   assumed), so it lines up with the same-numbered itemsep ratio despite
%   describing a different transition -- \gothamItemizeLevelDropSpaceRatioIV
%   governs the level-4-list-closing-back-to-level-3 transition and is the
%   one that has to individually stay bigger than \gothamItemsepRatioIV.
%   This branch was missing entirely at first (only I/II/III existed, a
%   leftover from before level 4 existed) -- confirmed directly on ZZO: the
%   level-4-to-level-3 closing transition measured 12.6pt with no drop-space
%   branch at all, less than level 4's own same-level gap of 13.4pt, a real
%   rule III violation. Adding the branch (\gothamItemizeLevelDropSpaceRatioIV,
%   starting at 0.5) fixed it. Re-verify rule III directly against ZZO's own
%   compiled PDF (don't assume it still holds on paper) if any of these ratios is
%   ever changed again.
\newcommand{\gothamItemsepRatioI}{0.4}
\newcommand{\gothamItemsepRatioII}{0.24}
\newcommand{\gothamItemsepRatioIII}{0.15}
\newcommand{\gothamItemsepRatioIV}{0.09}

% Per-level "opening" space -- the gap between a level-n item and the
% first item of a level-(n+1) list that immediately follows it, as
% opposed to \gothamItemizeLevelDropSpaceRatio* (a nested list *closing*
% back up) or \gothamItemsepRatio* (two items at the *same* depth).
% Added 2026-07-21, direct instruction, replacing the original three
% separate rules on slides/ZZO-Itemize-Spacing-Reference-current.tex's
% own header comment with one multi-part inequality stated per level n:
%   level-n -> level-(n-1) close  >  level-n same-level  >
%   level-n -> level-(n+1) open  >  level-n's own wrap gap
% The "open" term was entirely unhandled before this -- confirmed
% directly on ZZO: every opening transition (level 1 into level 2, level
% 2 into level 3, level 3 into level 4) measured *smaller* than that same
% level's own wrap gap, backwards from what the rule needs, since nothing
% but beamer's own bare default \topsep governed it.
% \AtBeginEnvironment{itemize}, not \topsep -- \topsep was tried first
% (set via the same \ifcase\@itemdepth block, right after \itemsep,
% inside the `\itemize` redefinition above) and measured *unreliable*:
% two of the three levels barely moved despite a sizeable ratio, while
% the third moved by roughly the right amount -- \topsep's own value is
% evidently consulted by `\list` at some point this override comes too
% late for, at least for some levels, not evenly. \addvspace via
% \AtBeginEnvironment instead measured consistently and predictably
% (each level's own change matched its own ratio's calculated pt value
% directly), the same mechanism \gothamItemizeLevelDropSpaceRatio* already
% relies on for the closing side. This hook fires with \@itemdepth
% reflecting the *parent* item's own depth (confirmed via \typeout, not
% assumed: opening the outermost level-1 list itself measured 0, opening
% a level-2 list from within a level-1 item measured 1, level-3-from-
% level-2 measured 2, level-4-from-level-3 measured 3) -- the opposite of
% \gothamItemizeLevelDropSpaceRatio*'s own hook, which keys on the depth
% of the list that *just closed*. That makes the \ifcase numbering here
% match the "n" in the rule directly: case 1 is level-1's own open-to-
% level-2 ratio, case 2 is level-2's own open-to-level-3, case 3 is
% level-3's own open-to-level-4; case 0 (opening the outermost list
% itself) is a genuine no-op, correctly -- there's no parent item above
% level 1 for that transition to separate from.
% Ratios chosen, then re-measured against the compiled ZZO reference and
% adjusted, the same way as every other ratio here: each has to
% individually land strictly between that level's own \gothamItemsepRatio*
% (the ceiling) and its own wrap gap (the floor). Re-verify directly
% against ZZO's own compiled PDF (don't assume it still holds on paper)
% if any of these, or the itemsep/drop-space ratios they're measured
% against, are ever changed again.
% \gothamItemizeLevelOpenSpaceRatioI trimmed from 0.28 to 0.02 (2026-07-21,
% direct instruction: "the level-1 to level-2 distance...seems a bit
% much"). Confirmed directly: even at 0, this transition still measured
% ~16.3pt -- \addvspace's own coalesce-to-the-larger-value behavior means
% the ratio can only ever push this *above* whatever glue the list
% mechanism inserts here regardless (that natural floor, not the ratio),
% so 0.28 down to 0 only ever moved the actual gap by ~2pt (18.5->16.3);
% 0.02 splits the difference, keeping the mechanism nominally "in effect"
% rather than reading as accidentally disabled. \gothamItemizeLevelOpenSpaceRatioIII
% (level-3-to-level-4) had to come down too, from 0.08 to 0.02, purely to
% restore its own margin -- changing RatioI shifted that unrelated
% transition's own measured gap closer to level 3's own same-level gap
% (12.1 vs. 12.0, too thin), confirmed directly on ZZO, not assumed;
% ratios elsewhere in this file have shown the same small-ripple behavior
% before (see \gothamItemizeLevelDropSpaceRatioIV's own comment) --
% re-measure the *whole* table after touching any single one of these,
% not just the transition being changed.
\newcommand{\gothamItemizeLevelOpenSpaceRatioI}{0.02}
\newcommand{\gothamItemizeLevelOpenSpaceRatioII}{0.15}
\newcommand{\gothamItemizeLevelOpenSpaceRatioIII}{0.02}
\AtBeginEnvironment{itemize}{%
	\ifcase\@itemdepth
		% 0: opening the outermost (level-1) list itself -- no parent item
		% precedes it, nothing to add.
	\or\addvspace{\gothamItemizeLevelOpenSpaceRatioI em}% opening a level-2 list from a level-1 item
	\or\addvspace{\gothamItemizeLevelOpenSpaceRatioII em}% opening a level-3 list from a level-2 item
	\or\addvspace{\gothamItemizeLevelOpenSpaceRatioIII em}% opening a level-4 list from a level-3 item
	\fi
}
\makeatother

% A \tableofcontents-driven agenda slide (see
% slides/nkj-Agenda-current.tex) needs the exact same top-margin/line-
% spacing as every other list in this deck, not whatever gotham's own
% "section in toc" templates use (both "gotham bullet" and "gotham simple"
% carry their own spacing/rule/font choices, tuned for gotham's generic
% look, not this project's own itemize metrics). Rather than re-tuning a
% *third* set of constants to match, this template just emits a plain
% \item -- wrapping the whole \tableofcontents call in a real
% \begin{enumerate}...\end{enumerate} then makes each section entry a
% genuine enumerate item, reusing this file's own \leftmargini/\itemsep/
% \gothamContentTopShift machinery exactly as any other list would, so it
% automatically stays consistent with the rest of the deck instead of
% drifting out of sync with it later.
\defbeamertemplate{section in toc}{gotham plain enumerate}{\item \inserttocsection}

% Even with the plain-enumerate template above, \tableofcontents still
% doesn't match this deck's own normal itemize spacing -- beamer's own
% \beamer@sectionintoc (beamerbasetoc.sty), which wraps *every* section
% entry regardless of which "section in toc" template is active,
% unconditionally inserts a fixed \vskip1.5em before each entry (an
% unbounded \vfill, even worse, if subsections aren't hidden -- neither is
% reachable via a public \tableofcontents option, just this show/hide
% binary), *and* wraps the entry itself in \hbox{\vbox{...}} -- box
% wrapping that (like \gothamMinipageTopCorrection's own minipage case
% elsewhere in this file) changes how surrounding baselineskip/topskip
% glue gets inserted compared to plain paragraph text. Both together added
% roughly 13pt of extra gap
% between entries beyond this deck's own normal ~18.6pt level-1 item
% spacing (measured directly against
% slides/57r-Reasons-for-Leasing-current.tex's own compiled PDF).
% Neither is controllable via a public \tableofcontents option, so this
% redefines \beamer@sectionintoc itself -- scoped safely since
% \tableofcontents is only ever called on the one Agenda slide in this
% deck -- keeping its numbering/hyperlink/part-number/hide-checking logic
% intact but dropping the vskip/vfill and the hbox/vbox wrapping, so an
% entry behaves exactly like a plain \item once printed.
\makeatletter
\def\beamer@sectionintoc#1#2#3#4#5{%
	\ifnum\c@tocdepth>0%
	\ifnum#4=\beamer@showpartnumber%
	{%
	\beamer@saveanother%
	\gdef\beamer@todo{}%
	\beamer@slideinframe=#1\relax%
	\expandafter\only\beamer@tocsections{\gdef\beamer@todo{%
			\beamer@tempcount=#5\relax%
			\advance\beamer@tempcount by\beamer@sectionadjust%
			\ifnum\beamer@tempcount>0
				\ifnum\beamer@tempcount>\beamer@toclastsection
					\def\inserttocsectionnumber{}%
				\else
					\edef\inserttocsectionnumber{\the\beamer@tempcount}%
				\fi
			\else
				\def\inserttocsectionnumber{}%
			\fi%
			\def\inserttocsection{\hyperlink{Navigation#3}{#2}}%
			\beamer@tocifnothide{\ifnum\c@section=#1\beamer@toc@cs\else\beamer@toc@os\fi}%
			{%
				\beamer@tocact{\ifnum\c@section=#1\beamer@toc@cs\else\beamer@toc@os\fi}{section in toc}%
			}%
		}%
	}%
	\beamer@restoreanother%
	}%
	\beamer@todo%
	\fi\fi%
}
\makeatother

% Left-aligns an enumerate's labels within a caller-chosen \labelwidth --
% used by the \tableofcontents Agenda slide (see
% slides/nkj-Agenda-current.tex's own comment for why its labels need
% this at 3.6pt). Defined here as a named helper rather than written as a
% literal \renewcommand{\makelabel}[1]{...#1...} inside the slide's own
% frame body, because that literal form -- specifically its raw `#1`
% parameter token -- breaks inside any frame that \tableofcontents forces
% through beamer's overlay reprocessing: beamer replays such a frame's
% collected body via \def-based machinery (\loop's own \iterate, and
% \beamer@doifinframe), and a `#` in the material of any \def must be
% doubled -- so the compile hits "! Illegal parameter number in definition
% of \iterate." (twice) plus the same for \beamer@doifinframe, once per
% compile, at that frame's own \end{frame}. TeX's error recovery ("I'll
% assume you meant ##") happens to produce the intended definition anyway,
% which is why the rendered PDF was always correct despite the errors --
% but lualatex still exits nonzero, which is exactly what kept breaking
% latexmk's automated lualatex/biber/rerun sequencing on the main deck
% (the "Known open issue" long tracked in CONVERSION_NOTES.md, originally
% misattributed to a biblatex interaction -- see its "Resolved" entry).
% Hoisting the definition out here means the frame body carries only this
% macro's *call* token, with no `#` in it, which sidesteps the whole
% mechanism. The inner \makelabel parameter is written `##1` below since
% it sits inside another \newcommand's body -- standard nesting, and also,
% incidentally, exactly the doubling rule the inline version fell foul of.
\newcommand{\gothamTocLabelFlushLeft}[1]{%
	\setlength{\labelwidth}{#1}%
	\renewcommand{\makelabel}[1]{\makebox[\labelwidth][l]{##1}}%
}

% Same idea as \gothamTocLabelFlushLeft above, plus an extra fixed
% \hspace shift before the label (for a list that needs to start further
% right than the plain nested-list indent, e.g. to line up with a
% sibling item's own text rather than its bullet -- see
% slides/scs-Political-Risk-Stability-Index-current.tex's own v1-5-0/
% v1-6-0 comments for the worked example this was hoisted out of).
% Found 2026-07-20 that the raw-`#1`-in-frame-body bug documented above
% is *not* actually \tableofcontents-specific, contrary to that comment's
% own framing -- slides/scs-Political-Risk-Stability-Index-current.tex has
% no \tableofcontents anywhere in the deck-wide overlay chain leading to
% it, yet hit the identical "Illegal parameter number in definition of
% \iterate" at its own \end{standardframe} from the exact same literal-
% \renewcommand{\makelabel}[1]{...#1...}-in-frame-body pattern. Whatever
% the precise trigger condition actually is, the same hoist-to-preamble
% fix (a bare macro *call*, no `#` token, left in the frame body) sidesteps
% it either way.
\newcommand{\gothamEnumerateLabelFlushLeftShifted}[2]{%
	\setlength{\labelwidth}{#2}%
	\renewcommand{\makelabel}[1]{\hspace{#1}\makebox[\labelwidth][l]{##1}}%
}

% Width of a single digit in this deck's normal content font, measured once
% here (safe, plain preamble-load-time \settowidth) rather than inline in
% the Agenda slide's own frame body -- \tableofcontents there puts the whole
% frame through beamer's overlay-replay machinery, which the comment above
% \gothamTocLabelFlushLeft already found chokes on raw tokens (like a bare
% \settowidth call) sitting directly in a replayed frame body, the same
% "Missing number"/"Illegal unit of measure" class of error as that `#`
% problem, confirmed directly the same way. Used by
% slides/nkj-Agenda-current.tex's own "widen the number-to-title gap"
% fix, which only ever needs to *reference* this precomputed value inside
% the frame, not measure anything itself there.
\newlength{\gothamTocLabelExtraGap}
\settowidth{\gothamTocLabelExtraGap}{0}

% For a two-column layout (e.g. bullets next to a chart/formula), nest
% beamer's own "columns" directly inside standardframe -- do NOT bake
% \begin{columns}/\end{columns} into a custom environment's begin/end code:
% beamer collects each frame's body as a token list to replay it once per
% overlay, and a columns environment that is opened in the "before" half and
% closed in the "after" half of a *different* environment confuses that
% collection (it never finds a properly matched \end{frame}). Nesting it in
% the document body, as below, avoids the issue entirely:
%
%   \begin{standardframe}{Title}
%   	\begin{columns}[t,onlytextwidth]
%   		\begin{column}{0.48\linewidth} ... \end{column}
%   		\begin{column}{0.48\linewidth} ... \end{column}
%   	\end{columns}
%   \end{standardframe}

% For automatic section dividers and single-statement "standout" slides, use
% what the gotham theme already provides instead of a custom environment:
%   \section{...}                                                (section divider, automatic)
%   \begin{frame}[standout]{Key takeaway}\end{frame}             (big single statement)
%   \begin{block}{...}/\begin{alertblock}{...}/\begin{exampleblock}{...}/\begin{referenceblock}{...}  (definitions, formulas, examples, references)
% For a course's own title slide, section-topic overview, quote, figure-only,
% and presenter-attention slides, see the five environments below.

% block/alertblock/exampleblock recolored to this project's own palette
% (2026-07-05, direct instruction) instead of gotham's own default colors
% -- gotham's own per-block-type color machinery (\gotham@block in
% beamerinnerthemegotham.sty) turns out to be dead code, never actually
% wired into any template, so these three blocks were rendering via plain
% beamer's own stock templates/colors the whole time; overriding the
% beamer color names directly here works regardless, since those stock
% templates read their colors from these same standard names (block~title/
% block~body and their ~alerted/~example variants) via beamercolorbox.
% Body background is the same c_tertiary_gray for all four
% block types (including referenceblock below) -- only each one's own
% title/header background differs. block's own header went
% gray→secondaryblue→secondaryviolet across three separate direct
% instructions in as many days -- secondaryviolet is the one actually in
% effect now.
\setbeamercolor{block title}{bg=c_secondary_violet, fg=c_primary_dark}
\setbeamercolor{block body}{bg=c_tertiary_gray, fg=c_primary_dark}
\setbeamercolor{block title example}{bg=c_secondary_green, fg=c_primary_dark}
\setbeamercolor{block body example}{bg=c_tertiary_gray, fg=c_primary_dark}
\setbeamercolor{block title alerted}{bg=c_secondary_red, fg=c_primary_dark}
\setbeamercolor{block body alerted}{bg=c_tertiary_gray, fg=c_primary_dark}

% All four block types' left/right margins widened to match a top-level
% itemize bullet's own text indent (2026-07-05, direct instruction: "the
% blocks should not have a smaller left-hand margin than the content
% text") -- a plain block spans \linewidth as-is, which lines up with the
% *frame's* own text margin, not with an itemize bullet's own text (which
% sits further right, past the bullet symbol and its gap) -- measured
% directly against a compiled slide with both an itemize and a block
% (slides/afa-Excursus-Price-Formation-current.tex): a bullet's own text
% starts 12.96pt further right than a plain block's left edge did.
%
% First attempt used changepage's \adjustwidth, on paper the "correct" tool
% for temporarily changing margins around a chunk of content -- it silently
% does *nothing* inside a beamer frame (beamer computes \linewidth/\hsize
% through its own internal frame-geometry machinery, not the page-geometry
% registers changepage adjusts, which are meant for article/report-style
% documents), and *loading the package at all* actively corrupted the
% title bar/logo on every frame (some interaction with gotham's own
% tikz/remember-picture-based title-bar drawing, not chased further since
% the package needed to go regardless) -- caught by actually compiling and
% looking, not assumed to work from the package's own documented purpose.
% Fixed instead with a plain \hspace{}+\begin{minipage}{\linewidth-2\times}
% wrapped around each block type's own begin/end template pair below --
% minipage genuinely resets \linewidth to its own local width for anything
% typeset inside it (already relied on elsewhere in this deck, e.g. the
% side-by-side text/image slides), which a bare beamercolorbox alone
% doesn't do on its own.
%
% 12.96pt matched a *second*-level bullet's own dot (which happens to
% coincide almost exactly with where a *first*-level bullet's own *text*
% starts, both ~21.6pt in from the frame's raw margin -- a coincidence of
% this deck's specific bullet-dot/gap/font metrics, not a general rule) --
% corrected the next day (direct instruction: align with the *first*-level
% bullet's own margin instead) to 4.20pt, matching a first-level bullet's
% own dot position (~12.84pt in from the frame's raw margin) instead.
\newlength{\gothamBlockSideMargin}
\setlength{\gothamBlockSideMargin}{4.20pt}

% Shared helper for a block title bar that carries a marker symbol flush
% right, opposite the title text (currently: infoblock's circled "i",
% definitionblock's $\triangleq$, exampleblock's $\in$ -- see each of
% those three's own template for the direct instruction behind its own
% marker choice). 2026-07-08, direct instruction, after infoblock's own
% "circled i nearly touching the title text" report on a two-line title
% (slides/u5h-Long-term-Bonds-Provisions-current.tex's own "Denomination,
% face value, or principal" step) -- generalized here so the same fix
% covers all three marker types at once, and any future one, rather than
% patching each block's own template separately.
%
% The marker itself is placed via a zero-width box *before* the title
% text, not via a trailing `\hfill<marker>` the way all three started out
% (definitionblock/exampleblock still do, until updated below) -- `\hfill`
% only fills space on whatever *line* it ends up typeset on, so for a
% title long enough to wrap onto a second line, `\hfill<marker>` lands on
% that second line instead of the first, printing the marker at the
% bottom right of the box instead of the top right (this is the same root
% cause infoblock's own marker already had fixed once already, see that
% template's own comment for the fuller writeup of *that* bug specifically
% -- this pushes the same fix out to the other two, and adds the part
% that bug's own fix didn't yet cover, below).
% `\makebox[0pt][l]{\hspace{\linewidth}\makebox[0pt][r]{<marker>}}` fixes
% *that* much on its own: the outer zero-width box doesn't displace the
% title text at all (still starts at the same point, still wraps however
% it needs to), but *inside* that zero-width box the marker is placed at
% a fixed horizontal offset (this box's own full \linewidth) from the
% title's own starting point, i.e. always the top right corner of the
% title bar regardless of how the title text itself wraps.
%
% What that alone still doesn't fix -- and what this round's actual
% report was about -- is the title text itself having no *awareness* that
% a marker will render on top of wherever it happens to end: a title just
% long enough to reach close to the marker's own fixed position, without
% actually needing to wrap, ends up visually touching or running under
% it. Fixed by also wrapping the title text itself in a
% `\begin{minipage}[t]{\linewidth-\gothamBlockTitleMarkerWidth}`, one
% \gothamBlockTitleMarkerWidth narrower than the full title bar, so the
% text itself wraps *before* reaching the marker's own reserved space
% instead of running into it -- confirmed directly against the compiled
% PDF on the exact case from the report. \gothamBlockTitleMarkerWidth is
% one shared reserved width for all three current marker types rather
% than a separate hand-measured width per symbol, which would need
% re-measuring by hand every time a new marker is added -- comfortably
% covers the widest of the three (checked directly) with a bit to spare.
% #1 = the marker's own content (e.g. \textcircled{\scriptsize i}, or a
% math-mode symbol like $\triangleq$); #2 = the title text
% (\insertblocktitle). Call as
% \gothamBlockTitleWithMarker{<marker>}{\insertblocktitle} in place of a
% bare \insertblocktitle, after the usual \usebeamerfont* call.
\newlength{\gothamBlockTitleMarkerWidth}
\setlength{\gothamBlockTitleMarkerWidth}{16pt}
% 2026-07-08 (direct instruction), two more fixes to this same helper:
% - A title long enough to wrap onto a second line
%   (slides/u5h-Long-term-Bonds-Provisions-current.tex's own
%   "Denomination, face value, or principal") read as *justified* rather
%   than left-aligned -- the minipage below had no explicit alignment, so
%   it fell back to LaTeX's default full justification, which only ever
%   visibly stretches a line short of the box's own right edge, i.e. only
%   shows up once a title is long enough to actually wrap. Every other
%   (single-line, shorter) title already looked left-aligned regardless,
%   since a line that's naturally shorter than the box is never stretched
%   whether or not justification is on -- confirmed directly, this is why
%   the bug went unnoticed until a title finally wrapped. `\raggedright`
%   fixes it without needing anything different for the single-line case.
% - Different titles' own title *bars* rendered at different heights
%   depending on whether that particular title's text happened to include
%   a descender (e.g. "Security"/"y", "Moody's rating at issue date"/"y")
%   or not (e.g. "Registered", "Trustee") -- the same class of bug already
%   fixed on slides/qxr-Introduction-to-R-Assignments-current.tex's own
%   colorboxes (see that slide's own comment): a box that tightly wraps its
%   own text naturally varies in height/depth with that text's own
%   ascenders/descenders. A leading `\strut` gives every title's first (and
%   for all but this one current title, only) line the same guaranteed
%   minimum height/depth, matching the qxr fix's own reasoning -- confirmed
%   directly, every current info/definition/example block title bar now
%   renders at the same height regardless of its own text.
\newcommand{\gothamBlockTitleWithMarker}[2]{%
	\makebox[0pt][l]{\hspace{\linewidth}\makebox[0pt][r]{#1}}%
	\begin{minipage}[t]{\dimexpr\linewidth-\gothamBlockTitleMarkerWidth\relax}\raggedright\strut#2\end{minipage}%
}

% A fourth block type, for a reference/citation-style callout (first used
% on slides/afa-Excursus-Price-Formation-current.tex and
% slides/4r7-IPO-Underpricing-International-current.tex, both replacing
% an earlier rotated hand-drawn callout box) -- beamer has no built-in
% "reference" block variant, so this defines one from scratch, mirroring
% exactly how beamer's own alertblock/exampleblock are built (compare
% beamerbaselocalstructure.sty's \newenvironment<>{alertblock} and
% beamerinnerthemedefault.sty's "block alerted begin/end" templates).
%
% First attempt (2026-07-05) copied these under the "rounded" style/
% beamerboxesrounded, on the wrong assumption that "rounded" was this
% deck's active block style generally -- it renders with visibly rounded
% corners, caught by the user, unlike block/alertblock/exampleblock's own
% square corners. Those three actually use beamerinnerthemedefault.sty's
% *starred* "default"-style templates (plain beamercolorbox, no rounding),
% which \defbeamertemplate*'s star selects immediately when that inner
% theme loads, regardless of beamer's own \defbeamertemplateparent{blocks}
% fallback to "rounded" -- gotham's own inner theme derives from
% beamerinnerthemedefault, so "default" (square) was the deck's real active
% style the whole time. Rewritten below to match that "default" style
% exactly (plain beamercolorbox instead of beamerboxesrounded) for
% consistent square corners across all four block types.
\setbeamercolor{block title reference}{bg=c_secondary_orange, fg=c_primary_dark}
\setbeamercolor{block body reference}{bg=c_tertiary_gray, fg=c_primary_dark}
% block/alertblock/exampleblock's own titles render bold via
% beamerfontthemegotham.sty's \setbeamerfont{block title}{series=\bfseries}
% (and an explicit ...alerted equivalent; example inherits it from
% beamer's own base font-fallback chain, verified directly -- both
% "Example"/"Caveat" titles on slides/ZZX-Definition-formula-current.tex
% render in the same heavier FiraSans-Regular as "Definition" despite
% gotham never setting "block title example" explicitly). "reference" is a
% brand new type name unknown to any of that inheritance, so it fell back
% to a plain, non-bold font with no color-analogous safety net (unlike
% \setbeamercolor{block title reference} above, colors for genuinely new
% names still resolve fine) -- caught only by comparing the actual
% rendered font name in the compiled PDF (FiraSans-Light, matching plain
% body text) against block/alertblock/exampleblock's own (FiraSans-Regular,
% this deck's de facto "bold" for Fira Sans), not by eye alone. Explicit
% setting needed, matching gotham's own two.
\setbeamerfont{block title reference}{size=\normalsize, series=\bfseries}
\defbeamertemplate*{block reference begin}{default}[1][]
{
	\hspace{\gothamBlockSideMargin}\begin{minipage}{\dimexpr\linewidth-2\gothamBlockSideMargin\relax}
	\par\vskip\medskipamount%
	\begin{beamercolorbox}[colsep*=.75ex,#1]{block title reference}
		\usebeamerfont*{block title reference}\strut\insertblocktitle%
	\end{beamercolorbox}\noframebreak%
	{\parskip0pt\par}%
	\ifbeamercolorempty[bg]{block title reference}
	{}
	{\ifbeamercolorempty[bg]{block body reference}{}{\nointerlineskip\vskip-0.5pt}}%
	\usebeamerfont{block body reference}%
	\begin{beamercolorbox}[colsep*=.75ex,vmode,#1]{block body reference}%
		\ifbeamercolorempty[bg]{block body reference}{\vskip-.25ex}{\vskip-.75ex}\vbox{}%
}
\defbeamertemplate*{block reference end}{default}
{\end{beamercolorbox}\vskip\smallskipamount\end{minipage}}
\newenvironment<>{referenceblock}[1]{%
	\begin{actionenv}#2%
		\def\insertblocktitle{#1}%
		\par%
		\usebeamertemplate{block reference begin}}
	{\par%
		\usebeamertemplate{block reference end}%
	\end{actionenv}}

% A sixth block type, for general background/context info that isn't
% itself a definition, example, or reference citation (2026-07-06, direct
% instruction: "use secondaryblue for a new type of info block") -- built
% the same from-scratch way as referenceblock just above (a brand-new
% style name needs \defbeamertemplate*'s registering form, not
% \setbeamertemplate -- see referenceblock's own comment for why). Marker
% is a circled "i" (\textcircled{\scriptsize i}, the generic LaTeX-kernel
% \textcircled -- not a Unicode "circled i" character: this deck's own
% Fira Sans/Fira Mono have no glyph for that codepoint either, the same
% missing-glyph problem \triangleq's own comment already found for
% definitionblock, and \textcircled sidesteps it entirely since it just
% draws a plain circle symbol around an ordinary letter rather than
% depending on one precomposed font glyph existing at all).
%
% The marker is placed via a zero-width box *before* \insertblocktitle,
% not via a trailing \hfill\textcircled{...} the way it first looked most
% natural to write -- \hfill only fills space on whatever line it's
% typeset on, and for a title long enough to wrap onto a second line (2026-07-08,
% found on the last reveal step of slides/ib3-Long-term-Bonds-Repayment-current.tex),
% \hfill\textcircled{...} then lands on that *second* line instead of the
% first, printing the marker at the bottom right of the box instead of the
% top right. Fixed originally with a one-off zero-width-box trick here;
% generalized the same day into the shared \gothamBlockTitleWithMarker
% helper (see its own comment, near \gothamBlockSideMargin above) once a
% *second*, related bug turned up on this exact block type (a long title
% visually touching the marker even without wrapping) -- that helper is
% what's actually called below now, not a copy of this trick inline.
\setbeamercolor{block title info}{bg=c_secondary_blue, fg=c_primary_dark}
\setbeamercolor{block body info}{bg=c_tertiary_gray, fg=c_primary_dark}
\setbeamerfont{block title info}{size=\normalsize, series=\bfseries}
\defbeamertemplate*{block info begin}{default}[1][]
{
	\hspace{\gothamBlockSideMargin}\begin{minipage}{\dimexpr\linewidth-2\gothamBlockSideMargin\relax}
	\par\vskip\medskipamount%
	\begin{beamercolorbox}[colsep*=.75ex,#1]{block title info}
		\usebeamerfont*{block title info}%
		\gothamBlockTitleWithMarker{\textcircled{\scriptsize i}}{\insertblocktitle}%
	\end{beamercolorbox}\noframebreak%
	{\parskip0pt\par}%
	\ifbeamercolorempty[bg]{block title info}
	{}
	{\ifbeamercolorempty[bg]{block body info}{}{\nointerlineskip\vskip-0.5pt}}%
	\usebeamerfont{block body info}%
	\begin{beamercolorbox}[colsep*=.75ex,vmode,#1]{block body info}%
		\ifbeamercolorempty[bg]{block body info}{\vskip-.25ex}{\vskip-.75ex}\vbox{}%
}
\defbeamertemplate*{block info end}{default}
{\end{beamercolorbox}\vskip\smallskipamount\end{minipage}}
\newenvironment<>{infoblock}[1]{%
	\begin{actionenv}#2%
		\def\insertblocktitle{#1}%
		\par%
		\usebeamertemplate{block info begin}}
	{\par%
		\usebeamertemplate{block info end}%
	\end{actionenv}}

% A fifth block type, for a "this introduces/defines a term" callout,
% visually distinguished from a plain \block (2026-07-06, direct
% instruction: "I would like to distinguish definition blocks from other
% blocks") by a $\triangleq$ symbol sitting flush right in the title bar,
% opposite the title text itself. Exact same colors/margins/font as plain
% \block (reuses "block title"/"block body", not a new color pair) --
% only the title-bar content differs, one extra `\hfill$\triangleq$`
% after `\insertblocktitle` -- so this mirrors referenceblock's own
% from-scratch pattern above (a brand new "definition" style name via
% \defbeamertemplate*, not \setbeamertemplate, for the same reason
% referenceblock's own comment gives: this style name has never been
% registered before, so the *-registration form is the correct one here,
% unlike plain block/alertblock/exampleblock below which are redefining
% already-registered "default"-style bodies). `\triangleq` is a genuine
% math symbol, so it needs `$...$` -- renders in this theme's math font
% (Computer Modern, see beamerfontthemegotham.sty's `no-math` option),
% a deliberate, visible one-glyph exception in an otherwise Fira Sans
% title bar, the same tradeoff every other math symbol in this deck
% already accepts.
\defbeamertemplate*{block definition begin}{default}[1][]
{
	\hspace{\gothamBlockSideMargin}\begin{minipage}{\dimexpr\linewidth-2\gothamBlockSideMargin\relax}
	\par\vskip\medskipamount%
	\begin{beamercolorbox}[colsep*=.75ex,#1]{block title}
		\usebeamerfont*{block title}\gothamBlockTitleWithMarker{$\triangleq$}{\insertblocktitle}%
	\end{beamercolorbox}\noframebreak%
	{\parskip0pt\par}%
	\ifbeamercolorempty[bg]{block title}
	{}
	{\ifbeamercolorempty[bg]{block body}{}{\nointerlineskip\vskip-0.5pt}}%
	\usebeamerfont{block body}%
	\begin{beamercolorbox}[colsep*=.75ex,vmode,#1]{block body}%
		\ifbeamercolorempty[bg]{block body}{\vskip-.25ex}{\vskip-.75ex}\vbox{}%
}
\defbeamertemplate*{block definition end}{default}
{\end{beamercolorbox}\vskip\smallskipamount\end{minipage}}
\newenvironment<>{definitionblock}[1]{%
	\begin{actionenv}#2%
		\def\insertblocktitle{#1}%
		\par%
		\usebeamertemplate{block definition begin}}
	{\par%
		\usebeamertemplate{block definition end}%
	\end{actionenv}}

% block/alertblock/exampleblock begin/end: same body as
% beamerinnerthemedefault.sty's own "default"-style templates (the ones
% actually in effect, per the comment above referenceblock), plus the
% minipage margin wrapper around each pair -- set via plain
% \setbeamertemplate{block begin}{...} (overriding the *body* of the
% existing "default" style directly) rather than
% \defbeamertemplate*{block begin}{default}{...} like referenceblock's own
% (new) templates above. The two look almost interchangeable but aren't:
% \defbeamertemplate* internally does a \newcommand-based *registration* of
% a style name, which errors outright ("... already defined") when that
% exact template+style pair already exists, as "block begin"/"default" and
% the other two pairs already do (registered once by
% beamerinnerthemedefault.sty itself) -- caught only by compiling and
% reading the log, not from any visible symptom in the rendered PDF (LaTeX
% recovers from the error by leaving the *original* stock template body in
% place, so these three blocks kept working and even picked up their new
% *colors* correctly via \setbeamercolor above, since color lookups happen
% at render time either way -- just silently without the margin fix, which
% *is* baked into the template body itself). \setbeamertemplate{name}{body}
% (no style-name argument) instead does a plain \def of that style's
% existing body storage, which is exactly "redefine what's already there"
% and carries no such restriction -- the correct tool whenever the target
% template+style pair is already registered, as opposed to
% \defbeamertemplate*, which is only for registering a style name for the
% very first time (as referenceblock's brand-new "reference" style name
% needs). Dropped the stock templates' own "#1" optional-argument slot
% (extra beamercolorbox options) in the process: nothing in this deck ever
% calls \usebeamertemplate{block begin}[...] with a bracket argument to
% populate it, so it was always empty in practice anyway.
% \strut before \insertblocktitle here, and in "block alerted"/"block
% reference" below: same title-bar-height-vs-descender fix as
% \gothamBlockTitleWithMarker's own comment above -- these three don't go
% through that helper at all (no marker symbol), so needed the same
% `\strut` added directly instead.
\setbeamertemplate{block begin}{
	\hspace{\gothamBlockSideMargin}\begin{minipage}{\dimexpr\linewidth-2\gothamBlockSideMargin\relax}
	\par\vskip\medskipamount%
	\begin{beamercolorbox}[colsep*=.75ex]{block title}
		\usebeamerfont*{block title}\strut\insertblocktitle%
	\end{beamercolorbox}\noframebreak%
	{\parskip0pt\par}%
	\ifbeamercolorempty[bg]{block title}
	{}
	{\ifbeamercolorempty[bg]{block body}{}{\nointerlineskip\vskip-0.5pt}}%
	\usebeamerfont{block body}%
	\begin{beamercolorbox}[colsep*=.75ex,vmode]{block body}%
		\ifbeamercolorempty[bg]{block body}{\vskip-.25ex}{\vskip-.75ex}\vbox{}%
}
\setbeamertemplate{block end}
{\end{beamercolorbox}\vskip\smallskipamount\end{minipage}}

% Warning-triangle marker for alertblock (direct instruction, 2026-07-12),
% same \gothamBlockTitleWithMarker treatment as definitionblock's
% $\triangleq$/exampleblock's $\in$/infoblock's \textcircled{\scriptsize i}
% above. No single built-in glyph renders a triangle-with-exclamation-mark
% in this deck's fonts (checked the same way \triangleq's own comment
% did), so it's hand-drawn: a plain tikz outline triangle with a bold "!"
% centered inside, sized by eye against \textcircled{\scriptsize i} until
% the two read as the same visual weight in a block title bar.
\newcommand{\gothamAlertMarker}{%
	\tikz[baseline=-0.6ex]{
		\draw[line width=1.1pt, line join=round] (-0.2,-0.11) -- (0.2,-0.11) -- (0,0.31) -- cycle;
		\node[font=\fontsize{7}{7}\selectfont\bfseries] at (0,0.01) {!};
	}%
}
\setbeamertemplate{block alerted begin}{
	\hspace{\gothamBlockSideMargin}\begin{minipage}{\dimexpr\linewidth-2\gothamBlockSideMargin\relax}
	\par\vskip\medskipamount%
	\begin{beamercolorbox}[colsep*=.75ex]{block title alerted}
		\usebeamerfont*{block title alerted}\gothamBlockTitleWithMarker{\gothamAlertMarker}{\insertblocktitle}%
	\end{beamercolorbox}\noframebreak%
	{\parskip0pt\par}%
	\ifbeamercolorempty[bg]{block title alerted}
	{}
	{\ifbeamercolorempty[bg]{block body alerted}{}{\nointerlineskip\vskip-0.5pt}}%
	\usebeamerfont{block body alerted}%
	\begin{beamercolorbox}[colsep*=.75ex,vmode]{block body alerted}%
		\ifbeamercolorempty[bg]{block body alerted}{\vskip-.25ex}{\vskip-.75ex}\vbox{}%
}
\setbeamertemplate{block alerted end}
{\end{beamercolorbox}\vskip\smallskipamount\end{minipage}}

% $\in$ ("is an instance/element of") flush right in the title bar
% (2026-07-06, direct instruction, mirroring definitionblock's own
% $\triangleq$ -- "what symbol would you propose ... to symbolize an
% example?"): a concrete example is naturally "an instance of" the general
% concept it illustrates, the same relationship $\in$ already denotes
% between an element and a set/class -- and, same as $\triangleq$, it's a
% math-mode symbol with a real Computer Modern glyph (no Fira Sans/Fira
% Mono text-mode substitute needed, unlike most Unicode symbol candidates
% -- see Slidemaster_preamble_uni-graz.tex's own "beamer/TeX gotchas" entry on
% $\triangleq$ for why that matters here). Applied to every \exampleblock/
% \gothamfloatingexampleblock in the deck via this shared template, not a
% one-off per slide, matching how $\triangleq$ applies to every
% definitionblock through its own shared template.
\setbeamertemplate{block example begin}{
	\hspace{\gothamBlockSideMargin}\begin{minipage}{\dimexpr\linewidth-2\gothamBlockSideMargin\relax}
	\par\vskip\medskipamount%
	\begin{beamercolorbox}[colsep*=.75ex]{block title example}
		\usebeamerfont*{block title example}\gothamBlockTitleWithMarker{$\in$}{\insertblocktitle}%
	\end{beamercolorbox}\noframebreak%
	{\parskip0pt\par}%
	\ifbeamercolorempty[bg]{block title example}
	{}
	{\ifbeamercolorempty[bg]{block body example}{}{\nointerlineskip\vskip-0.5pt}}%
	\usebeamerfont{block body example}%
	\begin{beamercolorbox}[colsep*=.75ex,vmode]{block body example}%
		\ifbeamercolorempty[bg]{block body example}{\vskip-.25ex}{\vskip-.75ex}\vbox{}%
}
\setbeamertemplate{block example end}
{\end{beamercolorbox}\vskip\smallskipamount\end{minipage}}

% A narrower, freely-positionable version of \begin{block}{...} (2026-07-06,
% direct instruction) -- a plain block always spans the *current*
% \linewidth (that's what block begin/end's own \hspace+minipage wrapper
% above is measured against), so it can't be dropped into a fixed-width
% spot (e.g. beside a table, inside a tikz overlay \node) without first
% narrowing \linewidth some other way. \gothamfloatingblock just wraps a
% plain block in one more, *outer* minipage of the caller's own chosen
% width -- block begin/end's own inner minipage then measures itself
% against that narrower width instead of the frame's full one, nesting
% correctly with no changes needed to the block templates themselves.
% Title/body colors, margins, and font are exactly \block's own (same
% "block title"/"block body" beamer color names) -- this only changes the
% overall width and (by taking it out of the normal document flow, left to
% the caller) the position, not the look. Font size is also left to the
% caller (wrap the whole environment in a size declaration, e.g.
% {\footnotesize \begin{gothamfloatingblock}{5cm}{Excel} ... \end{gothamfloatingblock}})
% rather than fixed here, since the right size depends entirely on how
% much room the block ends up placed in.
\newenvironment{gothamfloatingblock}[2]{%
	\begin{minipage}{#1}
	\begin{block}{#2}
}{%
	\end{block}
	\end{minipage}
}

% Same idea, but wrapping \exampleblock instead of \block (2026-07-06,
% direct instruction, first used on
% slides/fr3-Forms-of-Capital-Market-Efficiency-Strong-Semistrong-Weak-current.tex
% and slides/ha9-Monacorp-Truck-Lease-Breakeven-Rent-current.tex, both
% switching an existing \gothamfloatingblock to the "example" green header
% instead of the plain violet one) -- not worth generalizing into a single
% macro parameterized over the block name: \gothamfloatingblock's own body
% is only two lines, so a third variant is cheaper to just write out than
% to build a block-name-dispatch layer for.
\newenvironment{gothamfloatingexampleblock}[2]{%
	\begin{minipage}{#1}
	\begin{exampleblock}{#2}
}{%
	\end{exampleblock}
	\end{minipage}
}

% Same idea again, wrapping \infoblock (2026-07-06, direct instruction,
% first used on
% slides/fr3-Forms-of-Capital-Market-Efficiency-Strong-Semistrong-Weak-current.tex's
% own "Random walk" formula box, switched from the green
% \gothamfloatingexampleblock to this blue one the same day).
\newenvironment{gothamfloatinginfoblock}[2]{%
	\begin{minipage}{#1}
	\begin{infoblock}{#2}
}{%
	\end{infoblock}
	\end{minipage}
}

% Every gothamfloatinginfoblock anchored top-right against the institute
% logo (the deck's standard
% "(current page.north east)+(-\gothamFrametitleRightpading,-\gothamContentTopAlign)"
% corner pattern) lands its own visible right edge short of the logo's,
% by a fixed amount independent of the block's own chosen width (#1) --
% \infoblock's underlying template insets its colored fill by
% \gothamBlockSideMargin from *both* sides of whatever width it's given
% (that margin exists to match a first-level bullet's own left-hand
% position, see its own comment above; the fact that it also narrows the
% *right* side is just a side effect of reusing one symmetric
% hspace+minipage trick for both, not a deliberate right-margin choice),
% partly offset by beamer's own beamercolorbox "colsep*" background bleed
% (colsep* pads the colored fill outward past its own nominal box edge
% without also insetting the text, unlike unstarred "colsep" -- see
% beamerbasecolor.sty's own \beamercolorbox definition). Root-caused and
% measured directly (pymupdf against a compiled probe deck, not derived
% purely on paper) 2026-07-21, after
% slides/j8y-Asset-Allocation-Example-Weights-current.tex's own one-off
% "+2.2pt" nudge (measured that same day against just the institute
% logo's yellow chip) turned out to disagree with
% slides/fr3-Forms-of-Capital-Market-Efficiency-Strong-Semistrong-Weak-current.tex's
% own older "+3.165pt" nudge (2026-07-06, measured against the logo's
% *full* gray-bordered footprint, shadow/card included) -- both slides'
% own nudges were each individually correct for their own chosen
% reference edge, they just didn't agree on what "the logo's own edge"
% means. Direct instruction, 2026-07-21: the *full* gray-bordered
% footprint is the deck's intended reference, so 3.165pt is the value
% kept here; j8y's own newer 2.2pt nudge was the one revised to match,
% not this one.
%
% Use "-\gothamFloatingInfoblockRightAlign" (not the bare
% "-\gothamFrametitleRightpading") for a gothamfloatinginfoblock's own
% anchor node whenever its right edge is meant to align with the logo --
% every current use of this pattern across the deck has been updated to
% match; new slides using this same corner pattern should use this
% length from the start rather than re-measuring their own one-off nudge.
\newlength{\gothamFloatingInfoblockRightAlign}
\setlength{\gothamFloatingInfoblockRightAlign}{\dimexpr\gothamFrametitleRightpading-3.165pt\relax}

%%%%%%%%%%%%%%%%%%%%%%%%%%%%%%%%%%%%%%%%%%%%%%%%%%%%%%%%%%%%%%%%%%%%%%%%%%%%%
% Special frame types (converted from ExampleSlides.pptx; see README.md)
%
% All five below deliberately reuse as much of standardframe -- and of each
% other -- as possible, rather than reimplementing frame chrome from
% scratch: sectionframe, figureframe, and highlightframe are literally thin
% wrappers around standardframe itself (same title bar, logo, footer, bullet
% styling; only what's documented at each one differs). titleframe and
% quoteframe can't be standardframe wrappers -- both are title-less, and
% quoteframe replaces the whole body layout -- but still reuse its palette,
% fonts, and margins (\gothamFrametitleLeftpading, the institute logo
% drawing command, etc.) rather than inventing new ones.
%%%%%%%%%%%%%%%%%%%%%%%%%%%%%%%%%%%%%%%%%%%%%%%%%%%%%%%%%%%%%%%%%%%%%%%%%%%%%

% titleframe: the deck's own title slide (ExampleSlides.pptx slide 1) --
% typically used exactly once, at the very start of a course's first deck.
% One continuous yellow box -- like standardframe's own title bar -- flush
% with the page's top and left edges, stopping short of the institute logo
% on the right (the same margin already used around the logo elsewhere:
% \gothamLogoWidth + 2*\gothamFrametitleRightpading), holding the title and,
% below it, a second line (author/term/whatever) in the same black text as
% standardframe's own frametitle. Usage:
%
%   \titleframe{Financial Markets}{Stefan Palan\\Summer term 2026}
%
% A *command*, not an environment, taking the second (lower) line's content
% as its second argument -- plain text, use \\ for line breaks same as any
% normal LaTeX paragraph, not an itemize, since the source deck's own
% subtitle block is just stacked lines. This has to be a command rather
% than an environment with the body as its content: beamer's frame body is
% captured by a delimited-argument macro (\beamer@collect@@body, looked up
% by triggering the error below on purpose) hunting for a literal
% \end{frame}, and a beamercolorbox -- needed by an earlier version of this
% to wrap the subtitle text in its own colored band -- opened in a *custom
% environment's* begin code and closed in its end code confuses that hunt
% exactly like the \columns example above (same root cause, different
% environment): with
% \newenvironment{titleframe}{...\begin{beamercolorbox}{...}}{\end{beamercolorbox}...},
% compiling ends in "! File ended while scanning use of \beamer@collect@@body."
% every time, however short the body. Taking the content as a normal brace
% argument instead sidesteps the whole issue. No footer/frame number
% (matches beamer's own convention for title slides, e.g. the old
% \begin{frame}[plain,noframenumbering]\titlepage\end{frame} this replaces --
% see README.md) and no chapter-logo support: both would be meaningless on a
% slide that appears once, before any chapter exists.
%
% Unlike the frametitle hook further up, neither the logo nor the yellow box
% here need the white-patch-and-redraw trick: the box's own width already
% stops short of the logo (rather than being painted full-width and then
% partly painted back over), so there's no yellow left to paint over in the
% first place. Both still need \tikz's "remember picture" on their
% "current page"-anchored nodes, though, and so still need the usual second
% compile pass (see "Title bar and logo" in README.md) -- checked
% empirically: without it, "current page" doesn't resolve to the actual page
% box on the first pass, and the logo/yellow box end up positioned nowhere
% near the page corners they're meant to be anchored to.
%
% \gothamTitleframeBoxHeight is a fixed target for the box's *bottom* edge
% (matching what was actually asked for: "reaching down to where the lower
% grey bar currently ends" in the two-band design this replaces), not the
% blank space above the text -- that space is derived from it instead (box
% height minus the title+subtitle content's own measured height), so the
% bottom edge stays fixed even if a longer title ever wraps onto a second
% line; the blank space above would just shrink to compensate. Measured
% against the old two-band design's own rendered PDF coordinates (before
% this change): its second band's bottom edge sat 187.147pt below the page's
% top edge.
\newlength{\gothamTitleframeBoxHeight}
\setlength{\gothamTitleframeBoxHeight}{187.147pt}
\newlength{\gothamTitleframeBoxWidth}
\setlength{\gothamTitleframeBoxWidth}{\dimexpr\paperwidth-\gothamLogoWidth-2\gothamFrametitleRightpading\relax}
% Text inset from the box's own edges: \gothamFrametitleLeftpading, the same
% length standardframe's own title text uses, plus \gothamTitleframeTextSep
% (2.5ex, matching this command's own former per-band beamercolorbox "sep").
\newlength{\gothamTitleframeTextSep}
\setlength{\gothamTitleframeTextSep}{2.5ex}
\newlength{\gothamTitleframeBandGap}
\setlength{\gothamTitleframeBandGap}{0.8ex}
% Because the leading blank space above the text (further down) is itself
% computed as "box height minus the content's own natural height", the
% content's *bottom* padding is the only thing that determines how far the
% last line sits above the box's bottom edge -- the leading space silently
% absorbs whatever that padding is, keeping the bottom edge fixed either
% way. (The relationship is a direct 1:1 offset, checked empirically by
% comparing two rendered values, but with a nonzero intercept from beamer's
% own default top-of-frame vskip -- there's no clean closed-form for the
% right value, hence the empirical calibration here rather than a formula.)
% Calibrated (against the rendered PDF's own text coordinates, not just
% visually) to about 1.5 line-heights of the subtitle text's own \Large
% size (~20.6pt/line here) above the box's bottom edge; re-measure and
% adjust if the subtitle's font size, or beamer's own default top-of-frame
% vskip, ever change.
\newlength{\gothamTitleframeBottomPad}
\setlength{\gothamTitleframeBottomPad}{41pt}
\newbox\gothamTitleframeBox
\makeatletter
\newcommand{\titleframe}[2]{%
	\begin{frame}[plain,noframenumbering,t]%
	% A bare tikzpicture behaves like an hbox (same reasoning as the footer
	% box's own \vbox-to-0pt wrapper further up): sitting directly in the
	% frame's vertical list, it silently opens an implicit paragraph,
	% switching into horizontal mode. Everything placed after it -- the
	% \vspace* further down included -- would then run in horizontal mode
	% instead of vertical mode, which is what a first version of this command
	% got wrong: with two bare tikzpictures here, \vspace*'s glue ended up as
	% *horizontal* space on an implicit line together with them (visible as
	% the title/subtitle shifting sideways, not down, and the vspace amount
	% appearing to have no effect at all -- it was moving the wrong axis).
	% \vbox to 0pt{\vss ...} keeps each overlay at zero space in the normal
	% flow without ever leaving vertical mode.
	%
	% \gothamArtifactBegin/\gothamArtifactEnd (defined with the footer, "PDF
	% accessibility" comment): this box is purely decorative background, same
	% kind of element as the footer's own yellow box -- marked the same way,
	% for a screen reader to skip. titleframe has no footer at all to worry
	% about separately (it's a "plain" frame, no footline), so this box is
	% the only thing here that needs it; the title/subtitle text right below
	% is real content and stays unmarked.
	\gothamArtifactBegin{/Type /Background}%
	\vbox to 0pt{\vss
	\begin{tikzpicture}[overlay, remember picture]
		% \gothamTitleBoxColor, not a literal color token -- same setting the
		% standard frametitle bar (colorA, above) uses, so the title slide's
		% own box matches every other title bar in the deck (2026-07-16,
		% direct instruction).
		% outer sep=0pt here and on the logo below: without it, tikz's
		% default outer sep insets the fill/logo 0.2pt from the anchored
		% page corner -- a hairline white line above/left of this box, and
		% its right edge/the logo 0.2-0.4pt off every other frame's. See
		% the frametitle hook's own comment for the full story.
		\node[fill=\gothamTitleBoxColor, anchor=north west, inner sep=0pt, outer sep=0pt,
		      minimum width=\gothamTitleframeBoxWidth,
		      minimum height=\gothamTitleframeBoxHeight]
			at (current page.north west) {};
	\end{tikzpicture}%
	}%
	\gothamArtifactEnd
	\vbox to 0pt{\vss
	\begin{tikzpicture}[overlay, remember picture]
		% inner ysep derived, not \gothamFrametitleRightpading -- see the
		% frametitle hook's fixed-position logo redraw for why (matches the
		% inline logo's vertically-centered position on every other frame).
		\node[anchor=north east, inner xsep=\gothamFrametitleRightpading,
		      inner ysep=\dimexpr\gothamTitleBarBottomAlign-\gothamLogoHeight\relax,
		      outer sep=0pt]
			at (current page.north east) {\gothamInstituteLogoSquare};
	\end{tikzpicture}%
	}%
	\setbox\gothamTitleframeBox=\vbox{%
		\hsize=\gothamTitleframeBoxWidth
		\parindent=0pt
		\leftskip=\gothamFrametitleLeftpading
		\advance\leftskip by \gothamTitleframeTextSep
		\rightskip=\gothamFrametitleLeftpading
		\advance\rightskip by \gothamTitleframeTextSep
		% \gothamTitleTextColor, not a literal color token -- follows the
		% box color change right above (2026-07-16, direct instruction):
		% leaving this hardcoded dark would be unreadable once the box
		% itself switches to a dark blue for MCI/Innsbruck.
		\color{\gothamTitleTextColor}%
		% Beamer redefines \\ inside every frame to peek for an overlay spec
		% (\\<2->, for line-by-line reveals in itemize/enumerate); that
		% redefinition assumes it's running in beamer's own frame-body
		% machinery and chokes ("Undefined control sequence", silently
		% swallowing the following word) when \\ is used for a plain line
		% break in a raw \vbox like this one. Restoring plain LaTeX's own
		% \\ (\@normalcr) for this group only fixes it without affecting \\
		% anywhere else in the deck. (The @ here needs catcode 11, which is
		% why this whole command definition -- not just this line -- is
		% wrapped in \makeatletter/\makeatother: macro bodies are tokenized
		% at *definition* time, so a \makeatletter placed inside the body
		% itself would take effect too late, after "\@normalcr" had already
		% been mis-tokenized as "\@" followed by four ordinary-text letters.)
		\let\\\@normalcr
		\vskip\gothamTitleframeTextSep
		\noindent{\Huge\bfseries #1\par}%
		\vskip\gothamTitleframeTextSep
		\vskip\gothamTitleframeBandGap
		\vskip\gothamTitleframeTextSep
		\noindent{\Large #2\par}%
		\vskip\gothamTitleframeBottomPad
	}%
	\vspace*{\dimexpr\gothamTitleframeBoxHeight-\ht\gothamTitleframeBox-\dp\gothamTitleframeBox\relax}%
	\box\gothamTitleframeBox
	% This command's own \color{\gothamTitleTextColor} above (set inside a
	% stored box register) can leak into the *next* frame's body text --
	% see standardframe's own opening code ("Forces this frame's own body")
	% for the real fix and the full story (a trailing reset attempted
	% here first turned out to be a silent no-op, since nothing draws after
	% it within this same frame to actually force the color operator to be
	% emitted).
	\end{frame}%
}
\makeatother

% fullbleedframe: a slide that is *only* a full-screen image (FM2_2026s_ps_deck.pptx
% slide 13, a recruitment flyer pasted in as one embedded picture on an
% otherwise blank layout -- no title, no PPTX text at all) -- no title bar,
% no footer, no logo, matching that source's own blank layout intent, unlike
% figureframe (which is standardframe underneath and always keeps the usual
% title/footer chrome). The image is scaled to fit the slide while keeping
% its own aspect ratio (\includegraphics's keepaspectratio, capped by
% \paperwidth *and* \paperheight together), centered; whatever's left over
% on the binding axis (e.g. a widescreen slide holding a taller/narrower
% image) is filled by a solid black backdrop rather than left showing the
% frame's usual white canvas -- the point of a full-*bleed* image is that
% nothing reads as "empty margin" even when the aspect ratios don't match.
%
%   \fullbleedframe{media/photo-datacamp-recruitment-flyer}
%
% \gothamFullBleedBlack is a fixed HTML black, deliberately *not* a
% \colorlet of c_primary_dark: this frame type's backdrop should stay
% neutral true black regardless of whether c_primary_dark itself is ever
% restyled to something else (e.g. a softer near-black, the way
% c_custom_gray already is elsewhere in this file) for branding
% reasons -- those are two different jobs (brand color vs. letterbox
% backdrop) that only happen to want the same value today.
\definecolor{gothamFullBleedBlack}{HTML}{000000}
\newcommand{\fullbleedframe}[1]{%
	\begin{frame}[plain,noframenumbering,t]%
	% \gothamArtifactBegin/\gothamArtifactEnd (defined with the footer, "PDF
	% accessibility" comment): the black backdrop is purely decorative
	% letterboxing, marked the same way as titleframe's/quoteframe's own
	% background boxes so a screen reader skips it. The image itself is left
	% unmarked -- same as everywhere else in this file, since this project
	% has no alt-text infrastructure yet for imported images generally (a
	% separate, not-yet-made decision -- see README.md's accessibility
	% section for the same open question about the institute logo).
	\gothamArtifactBegin{/Type /Background}%
	\vbox to 0pt{\vss
	\begin{tikzpicture}[overlay, remember picture]
		% outer sep=0pt: see the frametitle hook's comment -- without it the
		% backdrop sits 0.2pt in from the anchored page corner.
		\node[fill=gothamFullBleedBlack, anchor=south west, inner sep=0pt, outer sep=0pt,
		      minimum width=\paperwidth, minimum height=\paperheight]
			at (current page.south west) {};
	\end{tikzpicture}%
	}%
	\gothamArtifactEnd
	\vbox to 0pt{\vss
	\begin{tikzpicture}[overlay, remember picture]
		\node[anchor=center] at (current page.center)
			{\includegraphics[width=\paperwidth, height=\paperheight, keepaspectratio]{#1}};
	\end{tikzpicture}%
	}%
	\end{frame}%
}

% sideimageframe: a standardframe that always carries one auto-placed
% "side filler image" (the user's own term, see
% slides/dfc-Angel-Investors-current.tex's own comment for the technique
% this formalizes into a reusable frame type) -- 2026-07-06, direct
% instruction. The image is positioned flush to three edges at once: its
% top touches the title bar's own bottom edge, its right touches the
% page's own right edge, and its bottom touches the footer box's own top
% edge -- so it's scaled by *height* to exactly
% \paperheight-\gothamTitleBarBottomAlign-\gothamFooterBoxHeight (the
% measured gap between those two), with `keepaspectratio` so the width
% simply follows from the image's own aspect ratio rather than being
% forced/distorted to some second, independent target -- anchoring the
% node at the page's own north-east corner then makes the top+right
% touches automatic, and the height exactly filling that gap makes the
% bottom touch automatic too, with no separate width constraint needed
% (only a *square* source image would ever visibly reach some further
% left boundary; this only guarantees top/right/bottom, matching what was
% actually asked for).
%
%   \begin{sideimageframe}{Title}{media/photo-example}
%   	\begin{itemize}
%   		\item ...
%   	\end{itemize}
%   	\source{...}
%   \end{sideimageframe}
%
% \sidefootnotes is switched on automatically (not left to the caller) --
% the whole point of this frame type is that its footnotes always end up
% rotated along the left edge of the image rather than colliding with it
% the way a full-width bottom footnote would (see
% slides/u29-Implications-of-Capital-Market-Efficiency-Prices-current.tex's
% own comment for exactly that collision, on a plain standardframe). Unlike
% every other \sidefootnotes use elsewhere in this deck, though, the strip
% can't sit at its usual fixed inset from the page's own right edge here --
% that would run it directly underneath the image, the same problem u29
% had. Fixed by *computing* the image's own rendered width (via
% \settowidth against the exact same \includegraphics call used to draw
% it) and using that to override \gothamSideFootnoteRightInset for just
% this one frame -- moving the whole rotated strip to sit immediately left
% of the image instead of at the page's raw edge, entirely automatically,
% with no manual x-coordinate needed from the caller (the direct
% instruction offered a manual-x-coordinate fallback "if it is possible
% [to do automatically], that would be great" -- turned out to be
% possible). \gothamSideImageFrameHeight/\gothamSideImageFrameWidth are
% declared once, globally, and merely *set* (not re-declared) inside the
% environment body, so the frame type can be used more than once per deck;
% the \setlength override of \gothamSideFootnoteRightInset is a plain,
% non-\global assignment, which -- like every other frame-local override in
% this file (e.g. slides/u29-...'s own \setbeamercolor{footnote}{fg=white})
% -- resets on its own once \begin{frame}...\end{frame}'s own TeX group
% closes, with no explicit reset needed.
%
% The extra gap added on top of the image's own measured width (2026-07-06,
% direct instruction: "replicate ... the precise margin we have on regular
% slides between footnote and footer box") was originally an unmeasured
% 4pt guess -- replaced with 3.1pt, back-solved from the deck's own real
% reference gap: measured directly on a plain bottom-footnote slide
% (slides/xhe-Active-vs-Passive-Management-2026-First-current.tex's own
% \source{} line) that the real gap between a footnote's own text and the
% footer box's top edge is 2.93pt, while the original 4pt guess here was
% actually landing the rotated text 3.8pt from the image -- 4pt - (3.8-2.93)
% = 3.13pt, rounded to 3.1pt.
\newlength{\gothamSideImageFrameHeight}
\newlength{\gothamSideImageFrameWidth}
\newenvironment{sideimageframe}[3][t]{%
	\setlength{\gothamSideImageFrameHeight}{\dimexpr\paperheight-\gothamTitleBarBottomAlign-\gothamFooterBoxHeight\relax}%
	\settowidth{\gothamSideImageFrameWidth}{\includegraphics[height=\gothamSideImageFrameHeight, keepaspectratio]{#3}}%
	\sidefootnotes
	\setlength{\gothamSideFootnoteRightInset}{\dimexpr\gothamSideImageFrameWidth+3.1pt\relax}%
	\begin{standardframe}[#1]{#2}%
	\vbox to 0pt{\vss
	\begin{tikzpicture}[overlay, remember picture]
		% outer sep=0pt: see the frametitle hook's comment -- without it the
		% image sits 0.2pt in from the page's right edge (and 0.2pt low).
		\node[anchor=north east, inner sep=0pt, outer sep=0pt] at ($(current page.north east)+(0,-\gothamTitleBarBottomAlign)$)
			{\includegraphics[height=\gothamSideImageFrameHeight, keepaspectratio]{#3}};
	\end{tikzpicture}%
	}%
}{%
	\end{standardframe}%
}

% sectionframe: a section-opening overview (ExampleSlides.pptx slide 3) --
% same layout as standardframe (footer, logo, bullet styling all completely
% unchanged), except the title bar itself is much taller (about a third of
% the page) and set in a bigger font, and the body text is a step bigger
% too, so it reads as "somewhat greater focus" and is unmistakable at a
% glance when flipping through slides -- while still keeping the deck's
% usual yellow-background/black-text look, unlike an earlier version of this
% that inverted to dark-background/yellow-text instead.
%
% The taller bar is NOT a separately-drawn box the way titleframe's is (see
% "Title slide" above): the gotham theme's own frametitle box already sizes
% itself from \gothamFrametitleToppading/\gothamFrametitleBottompading plus
% the title font's own strut height (beamerouterthemegotham.sty, the
% "frametitle" template), so overriding those two lengths locally -- the
% same trick already used below for the font/color overrides -- grows the
% *theme's own* box instead of fighting it with a second one. This also
% means the white-patch-behind-the-logo/chapter-logo machinery (tied to the
% frametitle template via \addtobeamertemplate) keeps working unmodified,
% since it measures the box's actual rendered bottom edge rather than
% assuming any particular height.
%
% The title text is bottom-anchored within the bar -- sitting a fixed "one
% line" above the bar's bottom edge, growing *upward* for a title long
% enough to wrap onto a second line -- rather than top-anchored the way
% gotham's own frametitle template flows by default (which would instead
% leave a large, growing gap *below* a wrapped title, between it and the
% bar's bottom edge). gotham's template has no built-in way to do this: its
% only two knobs are the toppading/bottompading struts, both anchored to
% *one* end of a single top-down paragraph flow. So instead of leaving
% toppading fixed (as an earlier version of this did) the whole title text
% is measured into \gothamSectionframeTitleBox first -- built with the same
% font and (approximately) the same usable width gotham's own template
% gives it, so its \ht+\dp closely track the title's real rendered size,
% including wrapping onto extra lines for a long title -- and toppading is
% then computed as "whatever's left" after that measured height and the
% fixed bottom gap are subtracted from the bar's total target height. A
% longer/wrapped title therefore automatically shrinks toppading (pushing
% the whole title block up, growing upward) rather than changing the gap
% below it. \gothamSectionframeTitleWidth is a slight underestimate of
% gotham's actual per-line width (which is only narrowed by the logo on
% whichever line ends up next to it, typically just the last one) --
% applying that narrower width to *every* line is a deliberate, harmless
% simplification rather than replicating gotham's per-line logic exactly.
%
% The institute logo moves down along with the title now (it's vertically
% centered against the title's own baseline, gotham's instituteLogo
% template), rather than staying pinned where it sits on every other frame
% type, as an earlier version of this kept it -- an intentional change from
% that version: with the title itself now sitting near the bar's bottom
% edge, keeping the logo up at the bar's top would visually separate it
% from the title it's meant to sit beside.
%
% \begingroup/\endgroup (rather than the \global-set-and-reset dance
% \chapterlogo{}/\sidefootnotes need) is enough here: this only ever touches
% \setbeamercolor/\setbeamerfont/\setlength/\setbox, all of which assign
% locally under the hood (checked against beamerbasecolor.sty/
% beamerbasefont.sty directly for the first two -- their key handlers are
% plain \def, not \gdef -- and \setlength/\setbox are plain TeX assignments,
% local by default), so a plain surrounding TeX group reverts them
% correctly on its own.
%
% \gothamSectionframeContentDrop adds a little extra room below the bar
% before the body content starts, on top of standardframe's own
% \gothamContentTopShift pull-up (applied inside \begin{standardframe}
% already, before this code regains control). All lengths here calibrated
% against the rendered PDF's own coordinates; re-measure and adjust if the
% title/body font sizes change.
\newlength{\gothamSectionframeContentDrop}
\setlength{\gothamSectionframeContentDrop}{10pt}
\newlength{\gothamSectionframeBoxHeight}
\setlength{\gothamSectionframeBoxHeight}{85.04pt}
% \gothamSectionframeBottomGap is the *visual* target -- one huge-font line
% height (baselineskip, measured directly from a deliberately 2-line test
% title: the gap between the two lines' own rendered top edges). It is NOT
% used directly as \gothamFrametitleBottompading: gotham's bottom strut
% (\gotham@frametitlestrut@end) consistently produces a visual gap about
% 5.58pt *smaller* than the padding value fed into it (checked empirically,
% same offset for both a one-line and a two-line test title, so it's a
% fixed per-box offset rather than something that scales with content) --
% \gothamSectionframeBottomPadCorrection adds that back.
\newlength{\gothamSectionframeBottomGap}
\setlength{\gothamSectionframeBottomGap}{14.32pt}
\newlength{\gothamSectionframeBottomPadCorrection}
\setlength{\gothamSectionframeBottomPadCorrection}{5.58pt}
% Similarly, the box's total rendered height comes out consistently ~11.4pt
% *less* than toppading+measuredTitleHeight+bottompading would suggest --
% the measured title box (built as a plain \vbox{...\par}) doesn't reproduce
% gotham's own leading strut (\gotham@frametitlestrut@start, a zero-width
% rule with an explicit height rather than natural interline glue) exactly,
% and this correction closes that gap. Checked the same way as the bottom
% one above: consistent across both the one-line and two-line test titles.
\newlength{\gothamSectionframeHeightCorrection}
\setlength{\gothamSectionframeHeightCorrection}{12.1pt}
\newlength{\gothamSectionframeTitleWidth}
\setlength{\gothamSectionframeTitleWidth}{\dimexpr\paperwidth-\gothamFrametitleLeftpading-\gothamLogoWidth-2\gothamFrametitleRightpading\relax}
\newbox\gothamSectionframeTitleBox
\newenvironment{sectionframe}[2][t]{%
	\begingroup
	% Disable hyphenation for the title -- both here, in the measurement
	% box below, and (since this is in effect for the rest of the group,
	% which \begin{standardframe} -- triggering the real \insertframetitle
	% -- runs inside) in the actual rendered title too, so a line break
	% never lands mid-word. Setting the penalties alone is the standard,
	% sufficient way to do this (TeX treats a penalty of 10000 at a
	% breakpoint as forbidden, not just heavily discouraged, so it's never
	% chosen even when the alternative is an overfull line); the hyphenmin
	% lengths are redundant insurance, not strictly needed on top of that.
	\hyphenpenalty=10000
	\exhyphenpenalty=10000
	\lefthyphenmin=100
	\righthyphenmin=100
	\setbeamerfont{frametitle}{size=\huge}
	\setbox\gothamSectionframeTitleBox=\vbox{%
		\hsize=\gothamSectionframeTitleWidth
		\usebeamerfont{frametitle}\usebeamercolor[fg]{frametitle}%
		\noindent #2\par
	}%
	\setlength{\gothamFrametitleBottompading}{\dimexpr\gothamSectionframeBottomGap+\gothamSectionframeBottomPadCorrection\relax}
	\setlength{\gothamFrametitleToppading}{\dimexpr\gothamSectionframeBoxHeight-\gothamFrametitleBottompading-\ht\gothamSectionframeTitleBox-\dp\gothamSectionframeTitleBox+\gothamSectionframeHeightCorrection\relax}
	\global\def\gothamFrametitleFullHeightYellow{1}%
	\global\def\gothamFrametitleLogoFixed{1}%
	\setbeamertemplate{frametitle}[gotham_no_inline_logo]%
	\begin{standardframe}[#1]{#2}%
	\vspace*{\gothamSectionframeContentDrop}%
	\large
}{%
	\end{standardframe}%
	\endgroup
}

% quoteframe: a quote slide (ExampleSlides.pptx slide 4) -- institute logo
% as usual, and the quote itself set large/italic/black on a yellow
% full-bleed block, same style as titleframe/sectionframe's own title boxes
% (flush top-left, stopping short of the logo on the right, same margin as
% elsewhere) rather than the dark block an earlier version of this used. The
% usual footer still shows (unlike titleframe, which drops it -- a title
% slide is a one-off outside any particular course content, while a quote
% is still one slide *among* a deck's others, worth being able to reference
% by its footer slide ID like any other). Opening/closing curly quotes are
% added automatically (write the quote text with none of its own); the
% second argument is the person it's from, set smaller/plain, below the
% yellow box and right-aligned to its own right edge (not the page's usual
% text margin) -- pass an empty second argument for an unattributed quote.
%
%   \quoteframe{One of the funny things about the stock market is that
%   	every time one person buys, another sells, and both think they
%   	are astute.}{William Feather, 1889-1981}
%
% \gothamQuoteBoxHeight is a plain percentage of the page, unlike
% titleframe's/sectionframe's own (calibrated-to-land-on-an-exact-fraction)
% box heights, since nothing here needs to line up with another element's
% measured position the way those do. \gothamTitleframeBoxWidth (defined
% above) is reused as-is for the box's width -- the same
% stop-short-of-the-logo formula, so no reason to duplicate it under a new
% name.
%
% Every piece here -- box, logo, quote text, attribution -- is drawn as a
% zero-space tikz overlay anchored to absolute page coordinates (the same
% \vbox to 0pt{\vss ...} pattern used elsewhere in this file for exactly
% this reason), rather than any of it flowing through the frame's normal
% vertical list. This sidesteps needing to know/counteract beamer's own
% default top-of-frame spacing (which standardframe explicitly measures and
% cancels via \gothamContentTopShift) -- there's no such frame-relative
% positioning to get right here, since literally everything is placed by
% explicit coordinate instead. Like the other page-corner overlays in this
% file, "current page" doesn't resolve correctly on a document's very first
% compile pass, so (like those) this needs the usual second pass too.
%
% A *command*, not an environment -- see the identically-reasoned comment on
% titleframe just above for why: the same beamercolorbox-spanning-a-custom-
% environment's-begin/end problem would apply here too, for the same reason.
\newlength{\gothamQuoteBoxHeight}
\setlength{\gothamQuoteBoxHeight}{0.8\paperheight}
\newlength{\gothamQuoteAttributionGap}
\setlength{\gothamQuoteAttributionGap}{1.8ex}
% \@empty (the empty-attribution check below) needs \makeatletter.
% \quotesource{text} -- opt-in \source{} for \quoteframe (call *before*
% \quoteframe{}{}, same set-then-consume pattern as \chapterlogo{}
% above), needed because \quoteframe is a *command* that auto-closes its
% own \end{frame} -- a bare \source{...} typed after \quoteframe{}{} in
% the calling slide file (as if it were still inside the frame, the way
% \source{} reads everywhere else in this deck) is actually issued
% *outside* any frame, and \source{}'s own \footnote{} then has no page to
% attach to yet. Confirmed directly (2026-07-12, direct instruction:
% traced a spurious extra blank leading overlay step on
% slides/c4c-Macbeth-Spot-Removers-Proposed-Structure-current.tex -- the
% very next multi-overlay frame after
% slides/w8w-Pie-Model-of-Capital-Structure-current.tex's own bare
% \source{} -- down to that orphaned footnote specifically, via a reduced
% two-slide test deck): the dangling footnote insert doesn't render
% anywhere visible on w8w itself, but silently corrupts beamer's own
% overlay-counting pre-pass on whichever frame happens to come right
% after -- reproduced with w8w+c4c alone (bug present) vs. c4c alone (bug
% absent), and confirmed it's specifically about the footnote being
% orphaned, not about w8w's content otherwise, by removing just its
% trailing \source{} line (bug also disappears). A plain \clearpage
% between the two frames also "fixes" c4c but only by giving the orphaned
% footnote a page of its own to attach to -- confirmed directly this
% splits w8w itself into two pages (real content, then a second page with
% nothing but the footnote/footer), not an acceptable trade. The only
% correct fix is keeping the footnote inside the frame that's actually
% citing something, i.e. before \quoteframe's own \end{frame} -- exactly
% what \quotesource{} below does.
\newcommand{\gothamQuoteSource}{}
\newcommand{\quotesource}[1]{\gdef\gothamQuoteSource{#1}}
\makeatletter
\newcommand{\quoteframe}[2]{%
	\begin{frame}[t]%
	% \gothamArtifactBegin/\gothamArtifactEnd (defined with the footer, "PDF
	% accessibility" comment): same reasoning as titleframe's own background
	% box above -- purely decorative, marked so a screen reader skips it. The
	% quote text and attribution below stay unmarked (real content); the
	% footer (this frame keeps one, unlike titleframe) is already marked via
	% the global footline hook.
	\gothamArtifactBegin{/Type /Background}%
	\vbox to 0pt{\vss
	\begin{tikzpicture}[overlay, remember picture]
		% \gothamTitleBoxColor, not a literal color token -- same setting
		% titleframe's own box and the standard frametitle bar (colorA) use
		% (2026-07-16, direct instruction).
		% outer sep=0pt here and on the logo/attribution below: same hairline-
		% inset-from-the-anchored-corner issue as titleframe's own box -- see
		% the frametitle hook's comment for the full story.
		\node[fill=\gothamTitleBoxColor, anchor=north west, inner sep=0pt, outer sep=0pt,
		      minimum width=\gothamTitleframeBoxWidth,
		      minimum height=\gothamQuoteBoxHeight]
			at (current page.north west) {};
	\end{tikzpicture}%
	}%
	\gothamArtifactEnd
	\vbox to 0pt{\vss
	\begin{tikzpicture}[overlay, remember picture]
		% inner ysep derived, not \gothamFrametitleRightpading -- see the
		% frametitle hook's fixed-position logo redraw for why (matches the
		% inline logo's vertically-centered position on every other frame).
		\node[anchor=north east, inner xsep=\gothamFrametitleRightpading,
		      inner ysep=\dimexpr\gothamTitleBarBottomAlign-\gothamLogoHeight\relax,
		      outer sep=0pt]
			at (current page.north east) {\gothamInstituteLogoSquare};
	\end{tikzpicture}%
	}%
	\vbox to 0pt{\vss
	\begin{tikzpicture}[overlay, remember picture]
		% \gothamTitleTextColor, not a literal color token -- follows the box
		% color change right above (2026-07-16, direct instruction). Set
		% *inside* the node's own {...} content group (not a stored box
		% register, unlike titleframe's own vbox above), so this doesn't hit
		% the same "\color leaks out of a stored \vbox" pitfall documented on
		% \titleframe's own closing code -- confirmed directly by checking
		% the frame right after this one for the same symptom.
		\node[anchor=center, inner sep=0pt,
		      text width=\dimexpr\gothamTitleframeBoxWidth-2\gothamFrametitleLeftpading\relax]
			at ($(current page.north west)+(0.5\gothamTitleframeBoxWidth,-0.5\gothamQuoteBoxHeight)$)
			{\color{\gothamTitleTextColor}\Large\itshape\textquotedblleft #1\textquotedblright};
	\end{tikzpicture}%
	}%
	\def\gothamQuoteAttributionTemp{#2}%
	\ifx\gothamQuoteAttributionTemp\@empty\else
		\vbox to 0pt{\vss
		\begin{tikzpicture}[overlay, remember picture]
			\node[anchor=north east, inner sep=0pt, outer sep=0pt]
				at ($(current page.north west)+(\gothamTitleframeBoxWidth,-\gothamQuoteBoxHeight-\gothamQuoteAttributionGap)$)
				{\normalfont\itshape\color{c_custom_gray} #2};
		\end{tikzpicture}%
		}%
	\fi
	\ifx\gothamQuoteSource\@empty\else
		\source{\gothamQuoteSource}%
		\global\let\gothamQuoteSource\@empty
	\fi
	\end{frame}%
}
\makeatother

% figureframe: a slide that is just a figure (ExampleSlides.pptx slide 5) --
% same standardframe layout (title, footer, \source{} for the usual
% "Source: ..."/"Image source: ..." line), but with the figure positioned
% via \figurecenter (call before the figure, inside the frame body) rather
% than left-aligned, which -- like the source deck's own figure slide -- is
% the default:
%
%   \begin{figureframe}{Correlation between DJIA stock returns over 15 years}
%   	\includegraphics[width=0.5\linewidth]{correlation-matrix}
%   	\source{Source: Own illustration, data from Yahoo! Finance, 24.02.2026.}
%   \end{figureframe}
%
%   \begin{figureframe}{Title}
%   	\figurecenter
%   	\includegraphics[width=0.6\linewidth]{diagram}
%   \end{figureframe}
%
% Vertical placement is standardframe's own [t]/[c] (its existing optional
% argument, passed straight through as figureframe's own -- so
% \begin{figureframe}[c]{Title} vertically centers exactly like
% \begin{standardframe}[c]{Title} already does); \figurecenter only ever
% affects the horizontal axis.
%
% \figurecenter is plain \centering under a more discoverable name, nothing
% more -- *not* an opt-in-before-the-frame flag like \chapterlogo{}/
% \sidefootnotes (those must run before \begin{standardframe}, since they
% configure something the frame's own opening code reads). This is called
% *inside* the body instead, exactly where it takes effect, so it can just
% take effect immediately like any other formatting declaration -- no flag,
% no state to reset afterwards, and (checked by first writing it as a flag
% read in figureframe's own opening code, mirroring \sidefootnotes) no way
% for the frame's opening code to see it anyway, since that code runs
% before the body -- and \figurecenter -- does.
\newcommand{\figurecenter}{\centering}
\newenvironment{figureframe}[2][t]{%
	\begin{standardframe}[#1]{#2}%
}{%
	\end{standardframe}%
}

% highlightframe: a standardframe with two things about it made unmistakably
% different -- the title bar's *and* the footer's background color, both
% together -- so a presenter can never flip to it without immediately
% noticing (ExampleSlides.pptx slide 6 was the inspiration for wanting a
% standardframe variant here, though the source slide itself has no such
% color cue; this is a new addition, not a conversion). Everything else --
% body, logo, bullet styling -- is exactly standardframe's own, via the same
% \begingroup/\endgroup wrapper sectionframe above uses for the title-color
% half of this. \gothamHighlightColor is a \colorlet, not a color baked
% directly into the environment, so retargeting every highlightframe in a
% deck to a different color -- title bar and footer both -- is still the
% same one-line change it was when this only touched the title bar.
%
% The footer box needs the \global\def dance (\gothamFooterBoxColor, defined
% above with the frametitle hook's other opt-in flags, reset in
% standardframe's own closing code) that \setbeamercolor{frametitle} doesn't:
% that assigns through beamer's own color-management system, which handles
% the frame-closing group correctly on its own (checked -- see
% \gothamChapterLogoPath's comment for why plain macros need \global\def but
% beamer's own set-commands don't), but the footer box's fill is read from a
% plain macro inside a tikz key, same category as \gothamChapterLogoPath.
\colorlet{gothamHighlightColor}{c_primary_gray}
\newenvironment{highlightframe}[2][t]{%
	\begingroup
	\setbeamercolor{frametitle}{bg=gothamHighlightColor}
	\global\def\gothamFooterBoxColor{gothamHighlightColor}%
	\begin{standardframe}[#1]{#2}%
}{%
	\end{standardframe}%
	\endgroup
}

